% GENERATED FILE -- DO NOT EDIT.
%
% Produced by harness/assemble_contract.py from module.yaml and claims/*.tex.
% Edit the claim files and regenerate; `--check` fails if this file is stale.
\documentclass[11pt,a4paper]{article}
\pdftrailerid{554478a28dbecc23e233ee3d3ff51bdf}
%% --- 1. Mathematics (order matters: physics after braket wins \ket, \tr) ---
\usepackage{amsmath}
\usepackage{amssymb}
\usepackage{amsfonts}
\usepackage{amsthm}
\usepackage{mathtools}
\usepackage{dsfont}
\usepackage{braket}
\usepackage{physics}
%% --- 2. Page furniture ------------------------------------------------------
\usepackage{xcolor}
\usepackage{array}    % >{...} column prefixes in the wide tables of some claims
\usepackage{enumitem}
\usepackage{graphicx}
\usepackage{booktabs}
\usepackage{longtable}
\usepackage{xltabular}              % front- and back-matter tables across pages
\usepackage{needspace}
\usepackage{xurl}
\usepackage[noautomatic]{imakeidx}  % defined-terms index; latexmk runs makeindex
\makeindex[title=Index of defined terms,intoc]
\usepackage{tcolorbox}
\tcbuselibrary{skins,breakable}
%% --- 3. Links and navigation (hyperref late; bookmark, refcount after it) ---
\usepackage{hyperref}
\usepackage{bookmark}
\usepackage{refcount}
%% --- 4. Colours: every distinction also has a shape (rule, caps, size, weight)
\definecolor{nnGrey}{gray}{0.40}      % non-normative text
\definecolor{metaGrey}{gray}{0.30}    % the generated meta line
\definecolor{techGrey}{gray}{0.45}    % heading of a technical claim
\definecolor{mainRule}{RGB}{0,110,140}
\definecolor{reqAccent}{RGB}{165,70,0}
\definecolor{linkBlue}{RGB}{0,62,150}
\definecolor{extLink}{RGB}{120,40,140}  % links into another PDF
\hypersetup{colorlinks=true, linkcolor=linkBlue, citecolor=linkBlue,
  urlcolor=linkBlue, filecolor=extLink, bookmarksnumbered=true,
  bookmarksopen=true, bookmarksopenlevel=1, pdfdisplaydoctitle=true}
% A link into another PDF (\href to a .pdf, \hyperref{file}..., xr-hyper) is a
% relative URI, <file>#nameddest=<destination>, in the file-link colour:
% browser viewers (Edge, Chrome) ignore hyperref's default GoToR action.
\makeatletter
\renewcommand*{\hyper@linkfile}[3]{% anchor text, file name, destination
  \begingroup
    \ifx\\#3\\\def\CQ@target{#2}\else\def\CQ@target{#2\hyper@hash nameddest=#3}\fi
    \let\@urlcolor\@filecolor
    \hyper@linkurl{#1}{\CQ@target}%
  \endgroup}
\makeatother
%% --- 5. The mapped Unicode set (harness/contract_navigation.py, UNICODE_TEX)
\DeclareUnicodeCharacter{03B1}{\ensuremath{\alpha}}
\DeclareUnicodeCharacter{03B2}{\ensuremath{\beta}}
\DeclareUnicodeCharacter{03B3}{\ensuremath{\gamma}}
\DeclareUnicodeCharacter{03B4}{\ensuremath{\delta}}
\DeclareUnicodeCharacter{03B5}{\ensuremath{\varepsilon}}
\DeclareUnicodeCharacter{03F5}{\ensuremath{\epsilon}}
\DeclareUnicodeCharacter{03B6}{\ensuremath{\zeta}}
\DeclareUnicodeCharacter{03B7}{\ensuremath{\eta}}
\DeclareUnicodeCharacter{03B8}{\ensuremath{\theta}}
\DeclareUnicodeCharacter{03D1}{\ensuremath{\vartheta}}
\DeclareUnicodeCharacter{03B9}{\ensuremath{\iota}}
\DeclareUnicodeCharacter{03BA}{\ensuremath{\kappa}}
\DeclareUnicodeCharacter{03BB}{\ensuremath{\lambda}}
\DeclareUnicodeCharacter{03BC}{\ensuremath{\mu}}
\DeclareUnicodeCharacter{03BD}{\ensuremath{\nu}}
\DeclareUnicodeCharacter{03BE}{\ensuremath{\xi}}
\DeclareUnicodeCharacter{03C0}{\ensuremath{\pi}}
\DeclareUnicodeCharacter{03C1}{\ensuremath{\rho}}
\DeclareUnicodeCharacter{03F1}{\ensuremath{\varrho}}
\DeclareUnicodeCharacter{03C3}{\ensuremath{\sigma}}
\DeclareUnicodeCharacter{03C2}{\ensuremath{\varsigma}}
\DeclareUnicodeCharacter{03C4}{\ensuremath{\tau}}
\DeclareUnicodeCharacter{03C5}{\ensuremath{\upsilon}}
\DeclareUnicodeCharacter{03C6}{\ensuremath{\varphi}}
\DeclareUnicodeCharacter{03D5}{\ensuremath{\phi}}
\DeclareUnicodeCharacter{03C7}{\ensuremath{\chi}}
\DeclareUnicodeCharacter{03C8}{\ensuremath{\psi}}
\DeclareUnicodeCharacter{03C9}{\ensuremath{\omega}}
\DeclareUnicodeCharacter{0393}{\ensuremath{\Gamma}}
\DeclareUnicodeCharacter{0394}{\ensuremath{\Delta}}
\DeclareUnicodeCharacter{0398}{\ensuremath{\Theta}}
\DeclareUnicodeCharacter{039B}{\ensuremath{\Lambda}}
\DeclareUnicodeCharacter{039E}{\ensuremath{\Xi}}
\DeclareUnicodeCharacter{03A0}{\ensuremath{\Pi}}
\DeclareUnicodeCharacter{03A3}{\ensuremath{\Sigma}}
\DeclareUnicodeCharacter{03A5}{\ensuremath{\Upsilon}}
\DeclareUnicodeCharacter{03A6}{\ensuremath{\Phi}}
\DeclareUnicodeCharacter{03A8}{\ensuremath{\Psi}}
\DeclareUnicodeCharacter{03A9}{\ensuremath{\Omega}}
\DeclareUnicodeCharacter{00E1}{\'a}
\DeclareUnicodeCharacter{00E0}{\`a}
\DeclareUnicodeCharacter{00E2}{\^a}
\DeclareUnicodeCharacter{00E4}{\"a}
\DeclareUnicodeCharacter{00E5}{\aa{}}
\DeclareUnicodeCharacter{00E9}{\'e}
\DeclareUnicodeCharacter{00E8}{\`e}
\DeclareUnicodeCharacter{00EA}{\^e}
\DeclareUnicodeCharacter{00EB}{\"e}
\DeclareUnicodeCharacter{00ED}{\'{\i}}
\DeclareUnicodeCharacter{00EF}{\"{\i}}
\DeclareUnicodeCharacter{00F3}{\'o}
\DeclareUnicodeCharacter{00F2}{\`o}
\DeclareUnicodeCharacter{00F4}{\^o}
\DeclareUnicodeCharacter{00F6}{\"o}
\DeclareUnicodeCharacter{00F8}{\o{}}
\DeclareUnicodeCharacter{00FA}{\'u}
\DeclareUnicodeCharacter{00F9}{\`u}
\DeclareUnicodeCharacter{00FB}{\^u}
\DeclareUnicodeCharacter{00FC}{\"u}
\DeclareUnicodeCharacter{00F1}{\~n}
\DeclareUnicodeCharacter{00E7}{\c{c}}
\DeclareUnicodeCharacter{00DF}{\ss{}}
\DeclareUnicodeCharacter{00C4}{\"A}
\DeclareUnicodeCharacter{00C5}{\AA{}}
\DeclareUnicodeCharacter{00C9}{\'E}
\DeclareUnicodeCharacter{00D6}{\"O}
\DeclareUnicodeCharacter{00DC}{\"U}
\DeclareUnicodeCharacter{2013}{\textendash{}}
\DeclareUnicodeCharacter{2014}{\textemdash{}}
\DeclareUnicodeCharacter{2018}{`}
\DeclareUnicodeCharacter{2019}{'}
\DeclareUnicodeCharacter{201C}{``}
\DeclareUnicodeCharacter{201D}{''}
\DeclareUnicodeCharacter{2026}{\ldots{}}
\DeclareUnicodeCharacter{00A0}{~}
\DeclareUnicodeCharacter{00A7}{\S{}}
\DeclareUnicodeCharacter{2264}{\ensuremath{\le}}
\DeclareUnicodeCharacter{2265}{\ensuremath{\ge}}
\DeclareUnicodeCharacter{2260}{\ensuremath{\neq}}
\DeclareUnicodeCharacter{2248}{\ensuremath{\approx}}
\DeclareUnicodeCharacter{2261}{\ensuremath{\equiv}}
\DeclareUnicodeCharacter{223C}{\ensuremath{\sim}}
\DeclareUnicodeCharacter{2243}{\ensuremath{\simeq}}
\DeclareUnicodeCharacter{2245}{\ensuremath{\cong}}
\DeclareUnicodeCharacter{221D}{\ensuremath{\propto}}
\DeclareUnicodeCharacter{226A}{\ensuremath{\ll}}
\DeclareUnicodeCharacter{226B}{\ensuremath{\gg}}
\DeclareUnicodeCharacter{2192}{\ensuremath{\to}}
\DeclareUnicodeCharacter{2190}{\ensuremath{\leftarrow}}
\DeclareUnicodeCharacter{2194}{\ensuremath{\leftrightarrow}}
\DeclareUnicodeCharacter{21D2}{\ensuremath{\Rightarrow}}
\DeclareUnicodeCharacter{21D0}{\ensuremath{\Leftarrow}}
\DeclareUnicodeCharacter{21D4}{\ensuremath{\Leftrightarrow}}
\DeclareUnicodeCharacter{21A6}{\ensuremath{\mapsto}}
\DeclareUnicodeCharacter{2208}{\ensuremath{\in}}
\DeclareUnicodeCharacter{2209}{\ensuremath{\notin}}
\DeclareUnicodeCharacter{2286}{\ensuremath{\subseteq}}
\DeclareUnicodeCharacter{2282}{\ensuremath{\subset}}
\DeclareUnicodeCharacter{2287}{\ensuremath{\supseteq}}
\DeclareUnicodeCharacter{2283}{\ensuremath{\supset}}
\DeclareUnicodeCharacter{2229}{\ensuremath{\cap}}
\DeclareUnicodeCharacter{222A}{\ensuremath{\cup}}
\DeclareUnicodeCharacter{2216}{\ensuremath{\setminus}}
\DeclareUnicodeCharacter{2205}{\ensuremath{\emptyset}}
\DeclareUnicodeCharacter{2200}{\ensuremath{\forall}}
\DeclareUnicodeCharacter{2203}{\ensuremath{\exists}}
\DeclareUnicodeCharacter{2115}{\ensuremath{\mathbb{N}}}
\DeclareUnicodeCharacter{2124}{\ensuremath{\mathbb{Z}}}
\DeclareUnicodeCharacter{211A}{\ensuremath{\mathbb{Q}}}
\DeclareUnicodeCharacter{211D}{\ensuremath{\mathbb{R}}}
\DeclareUnicodeCharacter{2102}{\ensuremath{\mathbb{C}}}
\DeclareUnicodeCharacter{00B7}{\ensuremath{\cdot}}
\DeclareUnicodeCharacter{00D7}{\ensuremath{\times}}
\DeclareUnicodeCharacter{2212}{\ensuremath{-}}
\DeclareUnicodeCharacter{00B1}{\ensuremath{\pm}}
\DeclareUnicodeCharacter{2213}{\ensuremath{\mp}}
\DeclareUnicodeCharacter{221A}{\ensuremath{\surd}}
\DeclareUnicodeCharacter{2211}{\ensuremath{\sum}}
\DeclareUnicodeCharacter{220F}{\ensuremath{\prod}}
\DeclareUnicodeCharacter{222B}{\ensuremath{\int}}
\DeclareUnicodeCharacter{221E}{\ensuremath{\infty}}
\DeclareUnicodeCharacter{2202}{\ensuremath{\partial}}
\DeclareUnicodeCharacter{2207}{\ensuremath{\nabla}}
\DeclareUnicodeCharacter{2297}{\ensuremath{\otimes}}
\DeclareUnicodeCharacter{2295}{\ensuremath{\oplus}}
\DeclareUnicodeCharacter{2218}{\ensuremath{\circ}}
\DeclareUnicodeCharacter{2308}{\ensuremath{\lceil}}
\DeclareUnicodeCharacter{2309}{\ensuremath{\rceil}}
\DeclareUnicodeCharacter{230A}{\ensuremath{\lfloor}}
\DeclareUnicodeCharacter{230B}{\ensuremath{\rfloor}}
\DeclareUnicodeCharacter{27E8}{\ensuremath{\langle}}
\DeclareUnicodeCharacter{27E9}{\ensuremath{\rangle}}
\DeclareUnicodeCharacter{2016}{\ensuremath{\|}}
\DeclareUnicodeCharacter{2113}{\ensuremath{\ell}}
\DeclareUnicodeCharacter{00B0}{\ensuremath{^{\circ}}}
\DeclareUnicodeCharacter{00B2}{\ensuremath{^{2}}}
\DeclareUnicodeCharacter{00B3}{\ensuremath{^{3}}}
\DeclareUnicodeCharacter{00B9}{\ensuremath{^{1}}}
\DeclareUnicodeCharacter{2080}{\ensuremath{_{0}}}
\DeclareUnicodeCharacter{2081}{\ensuremath{_{1}}}
\DeclareUnicodeCharacter{2082}{\ensuremath{_{2}}}
\DeclareUnicodeCharacter{207F}{\ensuremath{^{n}}}
\DeclareUnicodeCharacter{2096}{\ensuremath{_{k}}}
%% --- 6. Numbering: one counter per section -----------------------
\numberwithin{equation}{section}
\setcounter{tocdepth}{1}
%% --- 7. The project macro set, then the claim environments -----------------
% spec/macros.tex, read relative to this module's directory (latexmk -cd).
% =============================================================================
%  spec/macros.tex -- the one macro set every generated contract inputs.
%
%  Two parts: part A, the macros spec/notation.md catalogues, and part B,
%  the macros local to one module, grouped by module.
%  harness/assemble_contract.py emits % =============================================================================
%  spec/macros.tex -- the one macro set every generated contract inputs.
%
%  Two parts: part A, the macros spec/notation.md catalogues, and part B,
%  the macros local to one module, grouped by module.
%  harness/assemble_contract.py emits % =============================================================================
%  spec/macros.tex -- the one macro set every generated contract inputs.
%
%  Two parts: part A, the macros spec/notation.md catalogues, and part B,
%  the macros local to one module, grouped by module.
%  harness/assemble_contract.py emits \input{../../macros} for it, so the
%  file is read from spec/modules/<id>/.
%
%  RULES
%    * A macro that spec/notation.md catalogues is a verbatim copy of it
%      (part A). Changing one is a notation change: bump spec/notation.md.
%    * Part B holds the module-local macros. A macro defined by two modules
%      with the same expansion appears once, under the first module (in the
%      order below) that defines it. Two macros with the same name and
%      DIFFERENT expansions are both kept under a module prefix, and the
%      claims write the prefixed names:
%        \Chain  -> \redChain (reductions, \mathfrak{C}: the reduction chain)
%                   \pcChain  (protocol-compiler, \mathfrak{R}: the two-segment
%                             chain; \Chsrc and \Chrec are built on it)
%        \Nmsm   -> \redNmsm  (reductions, \mathcal{N}: a POVM measurement
%                             channel)
%                   \pcNmsm   (\mathcal{N}_{\mathsf A}: Alice's measurement
%                             channel, written by
%                             transmitter/thm:certificate-sufficiency,
%                             transmitter/rem:xreg and
%                             reductions/lem:source-transfer)
%    * Deliberately omitted: a \provenance{tag}{details} macro. Provenance
%      is a header field of every claim file (spec/claim-format.md section
%      3); a claim body that calls such a macro fails to compile, which is
%      the intended signal.
%
%  REQUIRED PACKAGES, loaded by the assembler's preamble before this file:
%    amsmath, amssymb, amsfonts, mathtools, dsfont, braket, physics, xcolor
%    (braket before physics, so that physics wins for \ket, \bra, \tr, ...).
%
%  This file contains only definitions: no \usepackage, no document body.
% =============================================================================


% =============================================================================
%  PART A -- notation macros, the ones spec/notation.md catalogues, each a
%  verbatim copy of its catalogue row; grouped under numbered headings.
% =============================================================================

% -----------------------------------------------------------------------------
% 1. State spaces and operator sets (a "system label" -> spaces of operators)
% -----------------------------------------------------------------------------
\newcommand{\Hs}[1]{\mathcal{H}(#1)}        % Hilbert space of a system label
\newcommand{\Lin}[1]{\mathcal{L}(#1)}       % bounded linear operators on H(#1)
\newcommand{\Linn}[2]{\mathcal{L}(#1,#2)}   % linear maps H(#1) -> H(#2)
\newcommand{\Pos}[1]{\mathrm{Pos}(#1)}      % positive semidefinite operators on H(#1)
\newcommand{\Density}{\mathrm{D}}           % density operators: \Density(A) = D(A)
\newcommand{\Seq}[1]{S_{=}(#1)}             % normalized states (trace = 1)
\newcommand{\Sleq}[1]{S_{\leq}(#1)}         % subnormalized states (0 <= trace <= 1)
\newcommand{\id}{\mathds{1}}                % identity operator
\newcommand{\idon}[1]{\mathds{1}_{#1}}      % identity on a named system, e.g. \idon{C}

% DEPRECATED: applies \mathcal to both arguments itself, so it takes bare
% letters, not system labels. Prefer \Linn; do not use in new contracts.
\providecommand{\Lop}[2]{\mathcal{L}(\mathcal{#1}, \mathcal{#2})}


% -----------------------------------------------------------------------------
% 2. Composite systems
% -----------------------------------------------------------------------------
\newcommand{\dsum}{\oplus}                  % direct sum of system labels
% Tensor product: use \otimes directly.
% n-fold direct sum of an operator space: \Lin{A}^{\oplus n}


% -----------------------------------------------------------------------------
% 3. Events, conditioning and classical registers (for STATES)
%    Convention:  \wedge  -> subnormalized "and-event" component, e.g. \rho_{A\wedge x}
%                 |        -> normalized conditional state,        e.g. \rho_{A|x}
% (Write these directly in subscripts; no extra macros needed.)
% -----------------------------------------------------------------------------


% -----------------------------------------------------------------------------
% 4. Photon-number subspaces / direct-sum blocks
%    Convention:  operators carry the restriction "on top", in brackets;
%                 states carry it as a \wedge subscript (see Sec. 3).
% -----------------------------------------------------------------------------
\newcommand{\Pnum}[1]{P_{#1}}               % photon-number projector (system from context)
\newcommand{\Pnumon}[2]{P_{#1}^{#2}}        % with explicit system superscript, e.g. P_n^C
\newcommand{\blk}[2]{#1^{(#2)}}             % operator block, e.g. \blk{M}{\leqN} = M^{(<=N)}
\newcommand{\leqN}{\leq N}                  % use as \blk{M}{\leqN}
\newcommand{\gtN}{> N}                      % use as \blk{M}{\gtN}
\newcommand{\eqn}{= n}                      % use as \blk{M}{\eqn}


% -----------------------------------------------------------------------------
% 5. System labels used in the protocol
% -----------------------------------------------------------------------------
\newcommand{\Chan}{C}                       % quantum-channel (Eve) input to Bob
\newcommand{\Prb}{P}                        % probe input
\newcommand{\Flag}{{\mathrm{Flg}}}          % flag system (EXISTS ONLY AFTER squashing)
\newcommand{\Loss}{{\mathrm{Lss}}}          % loss system
\newcommand{\Vc}{{\mathrm{Vcm}}}            % vacuum system
\newcommand{\Nout}{N_{\mathrm{out}}}
\newcommand{\Nin}{N_{\mathrm{in}}}


% -----------------------------------------------------------------------------
% 6. Alphabets
% -----------------------------------------------------------------------------
\newcommand{\Yset}{\mathcal{Y}}             % Bob's (coarse-grained) detection alphabet
\newcommand{\Yraw}{\mathcal{Y}_{\mathrm{raw}}}% raw click patterns {0,1}^4
\newcommand{\Yscset}{\mathcal{Y}_{\singlec}}% single-click outcomes
\newcommand{\Ymcset}{\mathcal{Y}_{\mc}}     % multi-click outcomes
\newcommand{\Tset}{\mathcal{T}_{\mathrm{set}}}% private setting set (the probe settings in the optical instance)
\newcommand{\Xset}{\mathcal{X}}             % Alice's alphabet
\newcommand{\Dset}{\{H,V,+,-\}}             % polarization / detector labels
\newcommand{\pass}{\top}                    % pass sifting


% -----------------------------------------------------------------------------
% 7. Mode / outcome abbreviations
% -----------------------------------------------------------------------------
\newcommand{\on}{{\mathrm{on}}}
\newcommand{\off}{{\mathrm{off}}}
\newcommand{\dc}{{\mathrm{dc}}}             % dark count
\newcommand{\sq}{\mathrm{sq}}               % squashed
\newcommand{\mc}{\mathrm{mc}}               % multi-click
\newcommand{\nc}{\mathrm{nc}}               % no-click
\newcommand{\singlec}{\mathrm{sc}}          % single-click
\newcommand{\vac}{\mathrm{vac}}             % vacuum
\newcommand{\kout}{{k_\mathrm{out}}}
\newcommand{\kin}{{k_\mathrm{in}}}


% -----------------------------------------------------------------------------
% 8. Maps, channels and miscellaneous operators
% -----------------------------------------------------------------------------
\newcommand{\Prob}{\mathrm{Prob}}           % probability (use \Prob[\cdot] or \Prob(\cdot))
\newcommand{\lmin}{\lambda_{\min}}          % minimum eigenvalue
\newcommand{\Wflag}{W_{\Flag}}              % flag-space weight
\newcommand{\transpose}{\mathsf{T}}
\newcommand{\innerp}[2]{\left\langle #1 \vert #2 \right\rangle}
\renewcommand{\vec}[1]{\boldsymbol{#1}}     % NB: overrides the LaTeX arrow accent


% -----------------------------------------------------------------------------
% 9. Key-rate numerics
% -----------------------------------------------------------------------------
\newcommand{\cG}{\mathcal{G}}               % protocol map
\newcommand{\cZ}{\mathcal{Z}}               % pinching map
\newcommand{\Amap}{\mathcal{A}_{\mathrm{ann}}} % announcement + measurement map
\newcommand{\Amapnb}{\mathcal{A}_{\mathrm{ann}|\mathrm{nb}}}
\newcommand{\simpl}{}                       % empty suffix for simplified block maps
\newcommand{\acSet}{\mathcal{S}_{\mathrm{acc}}} % acceptance set
\newcommand{\testE}{\mathcal{E}_{\mathrm{test}}} % test-statistics map
\newcommand{\ecLeak}{\mathrm{leak}^{\mathrm{EC}}_{\mathrm{obs}}}
\newcommand{\pgen}{p_{\mathrm{gen}}}
\newcommand{\askeyr}{R_{\infty}}             % asymptotic key rate
\newcommand{\ppass}{p_{\mathrm{pass}}}       % sifting pass probability


% -----------------------------------------------------------------------------
% 10. Finite-size analysis / Renyi entropies
% -----------------------------------------------------------------------------

% --- Round types: test vs. key generation (register T, alphabet \Fset) --------
\newcommand{\tst}{\mathrm{test}}            % test round
\newcommand{\gen}{\mathrm{gen}}             % key-generation round
\newcommand{\Fset}{\mathcal{T}}             % round-type alphabet {test, gen} (not \Tset)
\newcommand{\Chat}{\widehat{C}}             % acceptance-test register
\newcommand{\Chatset}{\widehat{\mathcal{C}}}% its alphabet

% --- Basis / key extraction functions ----------------------------------------
\newcommand{\bA}{b_{A}}                     % Alice's basis-extraction function
\newcommand{\bB}{b_{B}}                     % Bob's basis-extraction function
\newcommand{\kB}{k_{B}}                     % Bob's key-extraction function

% --- Maps --------------------------------------------------------------------
\newcommand{\Mprot}{\mathcal{M}}            % single-round protocol map
\newcommand{\Phimap}{\boldsymbol{\Phi}}     % test-statistics map
\newcommand{\Pdeph}{\mathcal{P}}            % dephasing (pinching) map on \Blreg

% --- Isometries --------------------------------------------------------------
\newcommand{\Vmsm}{V_{\mathrm{msm}}}        % measurement isometry
\newcommand{\Vann}{V_{\mathrm{ann}}}        % announcement isometry
\newcommand{\VI}{V_{I}}                     % Eve's copy of the announcements
\newcommand{\VChat}{V_{\Chat}}              % acceptance-data isometry
\newcommand{\VR}{V_{R}}                     % sifting + key map -> raw key R
\newcommand{\VZ}{V_{Z}}                     % raw key -> key register Z
\newcommand{\Wprot}{W}                      % full protocol isometry

% --- Renyi entropies ---------------------------------------------------------
% SECURITY-CRITICAL: see spec/notation.md §4. Never change an arrow to
% make an inequality work out; never convert silently between the sandwiched
% (tilde) and Petz (overline) families.
\newcommand{\Hsandup}{\widetilde{H}^{\uparrow}_{\alpha}}   % sandwiched, up
\newcommand{\Hsanddown}{\widetilde{H}^{\downarrow}_{\alpha}}% sandwiched, down
\newcommand{\Hpetzup}{\overline{H}^{\uparrow}_{\alpha}}    % Petz, up
\newcommand{\Hpetzdown}{\overline{H}^{\downarrow}_{\alpha}}% Petz, down
\newcommand{\hup}{h^{\uparrow}_{\alpha}}                   % single-round rate
\newcommand{\hqkd}{h^{\uparrow,\mathrm{QKD}}_{\alpha}}     % single-round QKD rate
\newcommand{\hlow}{\underline{h}_{\alpha}}                 % lower bound on \hqkd
% --- spec/notation.md 3.5 and 4.5 -------------------------------------------
% \hlb  is the first-order linearization lower bound on \hlow: the headline
%       key rate (the `first-order` tier of foundations/def:tiers). NOT \hlow
%       itself.
% \tnum is the instance's one declared numerical tolerance. Its
%       near-collision with the declared marginal \tau_A is disclosed in
%       spec/notation.md 3.5; bare \tau is reserved and denotes neither.
\newcommand{\hlb}{\underline{h}^{\mathrm{FO}}_{\alpha}}    % first-order lower bound on \hlow
\newcommand{\tnum}{\tau_{\mathrm{num}}}                    % declared numerical tolerance


% =============================================================================
%  PART B -- module-local macros, grouped by module. Deprecated aliases,
%  defined so that a text written with them compiles, used in no new contract: \Lop,
%  \testE, \Pdeph, \pass, \ecLeak (spec/notation.md section 3.9).
% =============================================================================

% --- meat-core -----------------------------------------------------------------
\newcommand{\SecSet}{\mathbb{S}}            % sector index set of Q
\newcommand{\Qsys}{Q}                       % declared sector-carrying measured system
\newcommand{\Vset}{\mathcal{V}}             % declared joint outcome alphabet of the Q-POVM
\newcommand{\CAset}{\mathcal{C}_{A}}        % Alice announcement alphabet
\newcommand{\CBset}{\mathcal{C}_{B}}        % Bob announcement alphabet
\newcommand{\Party}{\mathfrak{P}}           % declared party partition
\newcommand{\Oset}[1]{\mathcal{O}_{#1}}     % readout alphabet of party #1
\newcommand{\rd}[1]{\pi_{#1}}               % readout map of party #1
\newcommand{\Sadm}{\mathcal{S}_{A}}         % admissible Alice subspace
\newcommand{\nomarg}{\mathsf{none}}         % "no trusted marginal declared"
\newcommand{\Zset}{\mathcal{Z}_{\perp}}     % raw-key alphabet incl. \perp
\newcommand{\Zval}{\mathcal{Z}}             % key alphabet without \perp
\newcommand{\Etil}{\widetilde{E}}           % single-round purifying register
\newcommand{\Ehat}{\widehat{E}}             % Eve's n-round register
\newcommand{\OmAT}{\Omega_{\mathrm{AT}}}    % acceptance event
\newcommand{\epa}{\varepsilon_{\mathrm{PA}}}
\newcommand{\eev}{\varepsilon_{\mathrm{EV}}}
\newcommand{\lEC}{\lambda_{\mathrm{EC}}}
\newcommand{\Decl}{\mathfrak{D}}            % protocol declaration
\newcommand{\CList}{\mathfrak{L}}           % declared constraint list
\newcommand{\PinS}{\mathcal{P}_{\SecSet}}   % sector pinching on Q
\newcommand{\Feas}{\mathcal{W}}             % feasible set of block states
\newcommand{\Chkey}{\Chatset_{\mathrm{key}}}% key-live acceptance values
\newcommand{\Kann}{\mathcal{K}}             % announcement-class alphabet
\newcommand{\hold}{\mathsf{h}}              % declared key holder (A or B)
\newcommand{\sfD}{D}                        % classical KL divergence (base 2)
\newcommand{\Dtwo}{D}                       % quantum relative entropy (base 2)
\newcommand{\Dsand}{\widetilde{D}_{\alpha}} % sandwiched Renyi divergence
\newcommand{\idmap}{\mathrm{id}}            % identity channel
\newcommand{\Aev}{\mathcal{A}}              % Eve's per-round attack channel
\newcommand{\AttSet}{\mathfrak{A}}          % licensed attack-channel set
\newcommand{\Chacc}{\Chatset_{\mathrm{acc}}}% acceptance support of \acSet
\newcommand{\epsstar}{\varepsilon^{\star}}      % distance to an exact POVM
\newcommand{\dfH}{\Delta_{\mathrm{herm}}}       % Hermiticity defect
\newcommand{\dfP}{\Delta_{\mathrm{pos}}}        % negative part
\newcommand{\dfC}{\Delta_{\mathrm{cpl}}}        % completeness defect
\newcommand{\lmoq}{g}                       % acceptance-side gradient block
\newcommand{\lmoJ}{G}                       % feasible-side gradient block
\newcommand{\bacc}{\beta_{\mathrm{acc}}}    % certified acceptance-side minimum
\newcommand{\Itset}{\mathfrak{I}}           % finite index set of iterates (C22)

% --- reductions ----------------------------------------------------------------
\newcommand{\Bsys}{B}                        % physical (pre-squashing) system
\newcommand{\Tcl}[1]{\mathcal{T}(#1)}        % trace-class operators on H(#1)
\newcommand{\Ins}{\mathcal{I}}               % generalized-instrument branch map
\newcommand{\redNmsm}{\mathcal{N}}              % POVM measurement channel
\newcommand{\Eprot}{\mathcal{E}}             % n-round protocol channel
\newcommand{\Grnd}{\mathsf{G}}               % one round map of a protocol
\newcommand{\Kfin}{\mathsf{K}}               % end-of-protocol processing
\newcommand{\Crit}{\mathsf{sec}}             % value of a security criterion
\newcommand{\Pflag}{\Pi_{\Flag}}             % flag-sector projector on Q
\newcommand{\Pleq}{\Pi_{\leqN}}              % preserved-block projector
\newcommand{\Pgt}{\Pi_{\gtN}}                % above-cutoff projector
\newcommand{\Sph}{S^{\mathrm{ph}}}           % family of physical inputs
\newcommand{\Srch}{S^{\mathrm{rch}}}         % reachable squashed states
\newcommand{\Ssq}{S^{\sq}}                   % squashed admissible set
\newcommand{\Ffss}{F^{\mathrm{FSS}}}         % flag-state-squasher target
\newcommand{\Fwp}{F^{\mathrm{WP}}}           % multi-bound WPFSS target
\newcommand{\Mwp}{M^{\mathrm{WP}}}           % renormalized above-cutoff POVM
\newcommand{\Mdec}{{M'}^{\mathrm{ph}}}       % declared physical family on B
\newcommand{\Ppvm}{P^{\Flag}}                % flag PVM
\newcommand{\Lmp}{{\Lambda_{\mathrm{mp}}}}   % measure-and-prepare channel
\newcommand{\Lwp}{{\Lambda^{\mathrm{WP}}}}   % multi-bound WPFSS channel
\newcommand{\Lfss}{{\Lambda^{\mathrm{FSS}}}} % flag-state-squasher channel
\newcommand{\lmax}{\lambda_{\max}}           % maximum eigenvalue
\newcommand{\Lnoise}{\widetilde{\Lambda}}    % flag-noise channel
\newcommand{\cg}{\varpi}                     % declared coarse-graining map
\newcommand{\MCK}{\mathsf{MC}}               % multiclick raw-pattern set
\newcommand{\redChain}{\mathfrak{C}}            % reduction chain
\newcommand{\Lnk}{\ell}                      % a single typed link
\newcommand{\dsq}{\preceq}                      % the squashing relation
\newcommand{\opn}[1]{\left\lVert #1\right\rVert}           % operator (spectral) norm
\newcommand{\mxn}[1]{\left\lVert #1\right\rVert_{\max}}    % entrywise max modulus
\newcommand{\Fld}{\mathbb{F}}                   % exact ordered substrate field
\newcommand{\Fldi}{\mathbb{F}_{i}}              % its Gaussian extension F + iF
\newcommand{\lat}[1]{\big\lfloor #1\big\rfloor_{\kappa}} % round DOWN onto the lattice

% --- protocol-compiler ---------------------------------------------------------
\newcommand{\Spec}{\mathfrak{S}}             % protocol specification (compiler input)
\newcommand{\pcChain}{\mathfrak{R}}            % reduction chain (two typed segments)
\newcommand{\Chsrc}{\pcChain_{\mathrm{src}}}   % its Alice-side source-map segment
\newcommand{\Chrec}{\pcChain_{\mathrm{rec}}}   % its receiver-side squashing segment
\newcommand{\rvirt}{\rho^{\mathrm{v}}}       % virtual (PRE-source-map) signal state
\newcommand{\Fph}{F^{\mathrm{ph}}}           % physical (pre-reduction) POVM element
\newcommand{\Vph}{\mathcal{V}^{\mathrm{ph}}} % physical outcome alphabet
\newcommand{\Tpub}{\mathcal{T}_{\mathrm{pub}}}% public-choice alphabet
\newcommand{\Uset}{\mathcal{U}}              % residual outcome alphabet
\newcommand{\idxm}{\mathsf{idx}}             % label-to-index map
\newcommand{\Asig}{A'}                       % transmitted (signal) system -- the REAL one
\newcommand{\Avirt}{A''}                     % virtual transmitted system (source descriptor of Xi_1)
\newcommand{\Abar}{\bar{A}}                  % Alice's measured register
\newcommand{\Ahat}{\hat{A}}                  % Alice's shield register
\newcommand{\pcNmsm}{\mathcal{N}_{\mathsf A}}  % Alice's measurement channel (see the RULES block)
\newcommand{\Phip}{\Phi^{+}}                 % maximally entangled vector label
\newcommand{\Smap}{\mathfrak{s}}             % the emitted sector map (fraktur: a compile-time object)

% --- receiver ------------------------------------------------------------------
\newcommand{\Modes}{\mathbb{M}}              % declared ordered mode set
\newcommand{\Inp}{\mathbb{I}}                % declared input modes
\newcommand{\InpC}{\mathbb{I}_{\Chan}}       % channel input modes
\newcommand{\InpS}{\mathbb{I}_{\Prb}}        % probe (source) input modes
\newcommand{\Env}{\mathbb{E}}                % dilation (environment) modes
\newcommand{\Port}[1]{\mathbb{P}_{#1}}       % mode set of detector port #1
\newcommand{\Emd}[1]{E_{#1}}                 % coordinate projector on modes #1
\newcommand{\Gmat}{G}                        % instrument (transfer) matrix
\newcommand{\Fock}{\mathfrak{F}}             % Fock / second-quantization functor
\newcommand{\Symn}[1]{\mathrm{Sym}^{#1}}     % symmetric tensor power
\newcommand{\normord}[1]{{:}\mspace{2mu}#1\mspace{2mu}{:}}   % normal ordering
\newcommand{\rk}{\operatorname{rank}}        % rank (physics.sty owns \rank)
\newcommand{\SCK}{\mathsf{SC}}               % single-click pattern set
\newcommand{\Recv}{\mathfrak{R}}             % receiver instrument tuple
\newcommand{\Cat}{\mathfrak{C}}              % component catalogue
\newcommand{\Grph}{\mathfrak{E}}             % declared element graph
\newcommand{\defic}{\mathbb{L}}              % loss deficiency operator
\newcommand{\Pvac}[1]{\Pi_{\vac}(#1)}        % joint no-click element on ports #1
\newcommand{\pmc}{p_{\mc}}                   % uniform high-photon mc weight
\newcommand{\Bnd}{\beta}                     % the scalar bound function beta(n)
\newcommand{\Rot}{R}                         % planar rotation matrix
\newcommand{\Wplate}{W}                      % retarder (wave-plate) matrix
\newcommand{\Bsp}{B}                         % beamsplitter core matrix

% --- statistics ----------------------------------------------------------------
\newcommand{\honest}{p^{\mathrm{hon}}}          % honest per-round distribution
\newcommand{\ecomp}{\varepsilon_{\mathrm{comp}}}% completeness budget
\newcommand{\simplex}[1]{\Delta(#1)}            % probability simplex
\newcommand{\Chatcon}{\Chatset_{\mathrm{con}}}  % constrained outcomes
\newcommand{\Ireg}[3]{\mathcal{I}_{#1}(#2,#3)}  % regularized incomplete beta
\newcommand{\Ncount}[1]{\mathsf{N}_{#1}}        % per-outcome count
\newcommand{\qlo}{\mathbf{q}^{-}}               % box lower bounds
\newcommand{\qhi}{\mathbf{q}^{+}}               % box upper bounds
 for it, so the
%  file is read from spec/modules/<id>/.
%
%  RULES
%    * A macro that spec/notation.md catalogues is a verbatim copy of it
%      (part A). Changing one is a notation change: bump spec/notation.md.
%    * Part B holds the module-local macros. A macro defined by two modules
%      with the same expansion appears once, under the first module (in the
%      order below) that defines it. Two macros with the same name and
%      DIFFERENT expansions are both kept under a module prefix, and the
%      claims write the prefixed names:
%        \Chain  -> \redChain (reductions, \mathfrak{C}: the reduction chain)
%                   \pcChain  (protocol-compiler, \mathfrak{R}: the two-segment
%                             chain; \Chsrc and \Chrec are built on it)
%        \Nmsm   -> \redNmsm  (reductions, \mathcal{N}: a POVM measurement
%                             channel)
%                   \pcNmsm   (\mathcal{N}_{\mathsf A}: Alice's measurement
%                             channel, written by
%                             transmitter/thm:certificate-sufficiency,
%                             transmitter/rem:xreg and
%                             reductions/lem:source-transfer)
%    * Deliberately omitted: a \provenance{tag}{details} macro. Provenance
%      is a header field of every claim file (spec/claim-format.md section
%      3); a claim body that calls such a macro fails to compile, which is
%      the intended signal.
%
%  REQUIRED PACKAGES, loaded by the assembler's preamble before this file:
%    amsmath, amssymb, amsfonts, mathtools, dsfont, braket, physics, xcolor
%    (braket before physics, so that physics wins for \ket, \bra, \tr, ...).
%
%  This file contains only definitions: no \usepackage, no document body.
% =============================================================================


% =============================================================================
%  PART A -- notation macros, the ones spec/notation.md catalogues, each a
%  verbatim copy of its catalogue row; grouped under numbered headings.
% =============================================================================

% -----------------------------------------------------------------------------
% 1. State spaces and operator sets (a "system label" -> spaces of operators)
% -----------------------------------------------------------------------------
\newcommand{\Hs}[1]{\mathcal{H}(#1)}        % Hilbert space of a system label
\newcommand{\Lin}[1]{\mathcal{L}(#1)}       % bounded linear operators on H(#1)
\newcommand{\Linn}[2]{\mathcal{L}(#1,#2)}   % linear maps H(#1) -> H(#2)
\newcommand{\Pos}[1]{\mathrm{Pos}(#1)}      % positive semidefinite operators on H(#1)
\newcommand{\Density}{\mathrm{D}}           % density operators: \Density(A) = D(A)
\newcommand{\Seq}[1]{S_{=}(#1)}             % normalized states (trace = 1)
\newcommand{\Sleq}[1]{S_{\leq}(#1)}         % subnormalized states (0 <= trace <= 1)
\newcommand{\id}{\mathds{1}}                % identity operator
\newcommand{\idon}[1]{\mathds{1}_{#1}}      % identity on a named system, e.g. \idon{C}

% DEPRECATED: applies \mathcal to both arguments itself, so it takes bare
% letters, not system labels. Prefer \Linn; do not use in new contracts.
\providecommand{\Lop}[2]{\mathcal{L}(\mathcal{#1}, \mathcal{#2})}


% -----------------------------------------------------------------------------
% 2. Composite systems
% -----------------------------------------------------------------------------
\newcommand{\dsum}{\oplus}                  % direct sum of system labels
% Tensor product: use \otimes directly.
% n-fold direct sum of an operator space: \Lin{A}^{\oplus n}


% -----------------------------------------------------------------------------
% 3. Events, conditioning and classical registers (for STATES)
%    Convention:  \wedge  -> subnormalized "and-event" component, e.g. \rho_{A\wedge x}
%                 |        -> normalized conditional state,        e.g. \rho_{A|x}
% (Write these directly in subscripts; no extra macros needed.)
% -----------------------------------------------------------------------------


% -----------------------------------------------------------------------------
% 4. Photon-number subspaces / direct-sum blocks
%    Convention:  operators carry the restriction "on top", in brackets;
%                 states carry it as a \wedge subscript (see Sec. 3).
% -----------------------------------------------------------------------------
\newcommand{\Pnum}[1]{P_{#1}}               % photon-number projector (system from context)
\newcommand{\Pnumon}[2]{P_{#1}^{#2}}        % with explicit system superscript, e.g. P_n^C
\newcommand{\blk}[2]{#1^{(#2)}}             % operator block, e.g. \blk{M}{\leqN} = M^{(<=N)}
\newcommand{\leqN}{\leq N}                  % use as \blk{M}{\leqN}
\newcommand{\gtN}{> N}                      % use as \blk{M}{\gtN}
\newcommand{\eqn}{= n}                      % use as \blk{M}{\eqn}


% -----------------------------------------------------------------------------
% 5. System labels used in the protocol
% -----------------------------------------------------------------------------
\newcommand{\Chan}{C}                       % quantum-channel (Eve) input to Bob
\newcommand{\Prb}{P}                        % probe input
\newcommand{\Flag}{{\mathrm{Flg}}}          % flag system (EXISTS ONLY AFTER squashing)
\newcommand{\Loss}{{\mathrm{Lss}}}          % loss system
\newcommand{\Vc}{{\mathrm{Vcm}}}            % vacuum system
\newcommand{\Nout}{N_{\mathrm{out}}}
\newcommand{\Nin}{N_{\mathrm{in}}}


% -----------------------------------------------------------------------------
% 6. Alphabets
% -----------------------------------------------------------------------------
\newcommand{\Yset}{\mathcal{Y}}             % Bob's (coarse-grained) detection alphabet
\newcommand{\Yraw}{\mathcal{Y}_{\mathrm{raw}}}% raw click patterns {0,1}^4
\newcommand{\Yscset}{\mathcal{Y}_{\singlec}}% single-click outcomes
\newcommand{\Ymcset}{\mathcal{Y}_{\mc}}     % multi-click outcomes
\newcommand{\Tset}{\mathcal{T}_{\mathrm{set}}}% private setting set (the probe settings in the optical instance)
\newcommand{\Xset}{\mathcal{X}}             % Alice's alphabet
\newcommand{\Dset}{\{H,V,+,-\}}             % polarization / detector labels
\newcommand{\pass}{\top}                    % pass sifting


% -----------------------------------------------------------------------------
% 7. Mode / outcome abbreviations
% -----------------------------------------------------------------------------
\newcommand{\on}{{\mathrm{on}}}
\newcommand{\off}{{\mathrm{off}}}
\newcommand{\dc}{{\mathrm{dc}}}             % dark count
\newcommand{\sq}{\mathrm{sq}}               % squashed
\newcommand{\mc}{\mathrm{mc}}               % multi-click
\newcommand{\nc}{\mathrm{nc}}               % no-click
\newcommand{\singlec}{\mathrm{sc}}          % single-click
\newcommand{\vac}{\mathrm{vac}}             % vacuum
\newcommand{\kout}{{k_\mathrm{out}}}
\newcommand{\kin}{{k_\mathrm{in}}}


% -----------------------------------------------------------------------------
% 8. Maps, channels and miscellaneous operators
% -----------------------------------------------------------------------------
\newcommand{\Prob}{\mathrm{Prob}}           % probability (use \Prob[\cdot] or \Prob(\cdot))
\newcommand{\lmin}{\lambda_{\min}}          % minimum eigenvalue
\newcommand{\Wflag}{W_{\Flag}}              % flag-space weight
\newcommand{\transpose}{\mathsf{T}}
\newcommand{\innerp}[2]{\left\langle #1 \vert #2 \right\rangle}
\renewcommand{\vec}[1]{\boldsymbol{#1}}     % NB: overrides the LaTeX arrow accent


% -----------------------------------------------------------------------------
% 9. Key-rate numerics
% -----------------------------------------------------------------------------
\newcommand{\cG}{\mathcal{G}}               % protocol map
\newcommand{\cZ}{\mathcal{Z}}               % pinching map
\newcommand{\Amap}{\mathcal{A}_{\mathrm{ann}}} % announcement + measurement map
\newcommand{\Amapnb}{\mathcal{A}_{\mathrm{ann}|\mathrm{nb}}}
\newcommand{\simpl}{}                       % empty suffix for simplified block maps
\newcommand{\acSet}{\mathcal{S}_{\mathrm{acc}}} % acceptance set
\newcommand{\testE}{\mathcal{E}_{\mathrm{test}}} % test-statistics map
\newcommand{\ecLeak}{\mathrm{leak}^{\mathrm{EC}}_{\mathrm{obs}}}
\newcommand{\pgen}{p_{\mathrm{gen}}}
\newcommand{\askeyr}{R_{\infty}}             % asymptotic key rate
\newcommand{\ppass}{p_{\mathrm{pass}}}       % sifting pass probability


% -----------------------------------------------------------------------------
% 10. Finite-size analysis / Renyi entropies
% -----------------------------------------------------------------------------

% --- Round types: test vs. key generation (register T, alphabet \Fset) --------
\newcommand{\tst}{\mathrm{test}}            % test round
\newcommand{\gen}{\mathrm{gen}}             % key-generation round
\newcommand{\Fset}{\mathcal{T}}             % round-type alphabet {test, gen} (not \Tset)
\newcommand{\Chat}{\widehat{C}}             % acceptance-test register
\newcommand{\Chatset}{\widehat{\mathcal{C}}}% its alphabet

% --- Basis / key extraction functions ----------------------------------------
\newcommand{\bA}{b_{A}}                     % Alice's basis-extraction function
\newcommand{\bB}{b_{B}}                     % Bob's basis-extraction function
\newcommand{\kB}{k_{B}}                     % Bob's key-extraction function

% --- Maps --------------------------------------------------------------------
\newcommand{\Mprot}{\mathcal{M}}            % single-round protocol map
\newcommand{\Phimap}{\boldsymbol{\Phi}}     % test-statistics map
\newcommand{\Pdeph}{\mathcal{P}}            % dephasing (pinching) map on \Blreg

% --- Isometries --------------------------------------------------------------
\newcommand{\Vmsm}{V_{\mathrm{msm}}}        % measurement isometry
\newcommand{\Vann}{V_{\mathrm{ann}}}        % announcement isometry
\newcommand{\VI}{V_{I}}                     % Eve's copy of the announcements
\newcommand{\VChat}{V_{\Chat}}              % acceptance-data isometry
\newcommand{\VR}{V_{R}}                     % sifting + key map -> raw key R
\newcommand{\VZ}{V_{Z}}                     % raw key -> key register Z
\newcommand{\Wprot}{W}                      % full protocol isometry

% --- Renyi entropies ---------------------------------------------------------
% SECURITY-CRITICAL: see spec/notation.md §4. Never change an arrow to
% make an inequality work out; never convert silently between the sandwiched
% (tilde) and Petz (overline) families.
\newcommand{\Hsandup}{\widetilde{H}^{\uparrow}_{\alpha}}   % sandwiched, up
\newcommand{\Hsanddown}{\widetilde{H}^{\downarrow}_{\alpha}}% sandwiched, down
\newcommand{\Hpetzup}{\overline{H}^{\uparrow}_{\alpha}}    % Petz, up
\newcommand{\Hpetzdown}{\overline{H}^{\downarrow}_{\alpha}}% Petz, down
\newcommand{\hup}{h^{\uparrow}_{\alpha}}                   % single-round rate
\newcommand{\hqkd}{h^{\uparrow,\mathrm{QKD}}_{\alpha}}     % single-round QKD rate
\newcommand{\hlow}{\underline{h}_{\alpha}}                 % lower bound on \hqkd
% --- spec/notation.md 3.5 and 4.5 -------------------------------------------
% \hlb  is the first-order linearization lower bound on \hlow: the headline
%       key rate (the `first-order` tier of foundations/def:tiers). NOT \hlow
%       itself.
% \tnum is the instance's one declared numerical tolerance. Its
%       near-collision with the declared marginal \tau_A is disclosed in
%       spec/notation.md 3.5; bare \tau is reserved and denotes neither.
\newcommand{\hlb}{\underline{h}^{\mathrm{FO}}_{\alpha}}    % first-order lower bound on \hlow
\newcommand{\tnum}{\tau_{\mathrm{num}}}                    % declared numerical tolerance


% =============================================================================
%  PART B -- module-local macros, grouped by module. Deprecated aliases,
%  defined so that a text written with them compiles, used in no new contract: \Lop,
%  \testE, \Pdeph, \pass, \ecLeak (spec/notation.md section 3.9).
% =============================================================================

% --- meat-core -----------------------------------------------------------------
\newcommand{\SecSet}{\mathbb{S}}            % sector index set of Q
\newcommand{\Qsys}{Q}                       % declared sector-carrying measured system
\newcommand{\Vset}{\mathcal{V}}             % declared joint outcome alphabet of the Q-POVM
\newcommand{\CAset}{\mathcal{C}_{A}}        % Alice announcement alphabet
\newcommand{\CBset}{\mathcal{C}_{B}}        % Bob announcement alphabet
\newcommand{\Party}{\mathfrak{P}}           % declared party partition
\newcommand{\Oset}[1]{\mathcal{O}_{#1}}     % readout alphabet of party #1
\newcommand{\rd}[1]{\pi_{#1}}               % readout map of party #1
\newcommand{\Sadm}{\mathcal{S}_{A}}         % admissible Alice subspace
\newcommand{\nomarg}{\mathsf{none}}         % "no trusted marginal declared"
\newcommand{\Zset}{\mathcal{Z}_{\perp}}     % raw-key alphabet incl. \perp
\newcommand{\Zval}{\mathcal{Z}}             % key alphabet without \perp
\newcommand{\Etil}{\widetilde{E}}           % single-round purifying register
\newcommand{\Ehat}{\widehat{E}}             % Eve's n-round register
\newcommand{\OmAT}{\Omega_{\mathrm{AT}}}    % acceptance event
\newcommand{\epa}{\varepsilon_{\mathrm{PA}}}
\newcommand{\eev}{\varepsilon_{\mathrm{EV}}}
\newcommand{\lEC}{\lambda_{\mathrm{EC}}}
\newcommand{\Decl}{\mathfrak{D}}            % protocol declaration
\newcommand{\CList}{\mathfrak{L}}           % declared constraint list
\newcommand{\PinS}{\mathcal{P}_{\SecSet}}   % sector pinching on Q
\newcommand{\Feas}{\mathcal{W}}             % feasible set of block states
\newcommand{\Chkey}{\Chatset_{\mathrm{key}}}% key-live acceptance values
\newcommand{\Kann}{\mathcal{K}}             % announcement-class alphabet
\newcommand{\hold}{\mathsf{h}}              % declared key holder (A or B)
\newcommand{\sfD}{D}                        % classical KL divergence (base 2)
\newcommand{\Dtwo}{D}                       % quantum relative entropy (base 2)
\newcommand{\Dsand}{\widetilde{D}_{\alpha}} % sandwiched Renyi divergence
\newcommand{\idmap}{\mathrm{id}}            % identity channel
\newcommand{\Aev}{\mathcal{A}}              % Eve's per-round attack channel
\newcommand{\AttSet}{\mathfrak{A}}          % licensed attack-channel set
\newcommand{\Chacc}{\Chatset_{\mathrm{acc}}}% acceptance support of \acSet
\newcommand{\epsstar}{\varepsilon^{\star}}      % distance to an exact POVM
\newcommand{\dfH}{\Delta_{\mathrm{herm}}}       % Hermiticity defect
\newcommand{\dfP}{\Delta_{\mathrm{pos}}}        % negative part
\newcommand{\dfC}{\Delta_{\mathrm{cpl}}}        % completeness defect
\newcommand{\lmoq}{g}                       % acceptance-side gradient block
\newcommand{\lmoJ}{G}                       % feasible-side gradient block
\newcommand{\bacc}{\beta_{\mathrm{acc}}}    % certified acceptance-side minimum
\newcommand{\Itset}{\mathfrak{I}}           % finite index set of iterates (C22)

% --- reductions ----------------------------------------------------------------
\newcommand{\Bsys}{B}                        % physical (pre-squashing) system
\newcommand{\Tcl}[1]{\mathcal{T}(#1)}        % trace-class operators on H(#1)
\newcommand{\Ins}{\mathcal{I}}               % generalized-instrument branch map
\newcommand{\redNmsm}{\mathcal{N}}              % POVM measurement channel
\newcommand{\Eprot}{\mathcal{E}}             % n-round protocol channel
\newcommand{\Grnd}{\mathsf{G}}               % one round map of a protocol
\newcommand{\Kfin}{\mathsf{K}}               % end-of-protocol processing
\newcommand{\Crit}{\mathsf{sec}}             % value of a security criterion
\newcommand{\Pflag}{\Pi_{\Flag}}             % flag-sector projector on Q
\newcommand{\Pleq}{\Pi_{\leqN}}              % preserved-block projector
\newcommand{\Pgt}{\Pi_{\gtN}}                % above-cutoff projector
\newcommand{\Sph}{S^{\mathrm{ph}}}           % family of physical inputs
\newcommand{\Srch}{S^{\mathrm{rch}}}         % reachable squashed states
\newcommand{\Ssq}{S^{\sq}}                   % squashed admissible set
\newcommand{\Ffss}{F^{\mathrm{FSS}}}         % flag-state-squasher target
\newcommand{\Fwp}{F^{\mathrm{WP}}}           % multi-bound WPFSS target
\newcommand{\Mwp}{M^{\mathrm{WP}}}           % renormalized above-cutoff POVM
\newcommand{\Mdec}{{M'}^{\mathrm{ph}}}       % declared physical family on B
\newcommand{\Ppvm}{P^{\Flag}}                % flag PVM
\newcommand{\Lmp}{{\Lambda_{\mathrm{mp}}}}   % measure-and-prepare channel
\newcommand{\Lwp}{{\Lambda^{\mathrm{WP}}}}   % multi-bound WPFSS channel
\newcommand{\Lfss}{{\Lambda^{\mathrm{FSS}}}} % flag-state-squasher channel
\newcommand{\lmax}{\lambda_{\max}}           % maximum eigenvalue
\newcommand{\Lnoise}{\widetilde{\Lambda}}    % flag-noise channel
\newcommand{\cg}{\varpi}                     % declared coarse-graining map
\newcommand{\MCK}{\mathsf{MC}}               % multiclick raw-pattern set
\newcommand{\redChain}{\mathfrak{C}}            % reduction chain
\newcommand{\Lnk}{\ell}                      % a single typed link
\newcommand{\dsq}{\preceq}                      % the squashing relation
\newcommand{\opn}[1]{\left\lVert #1\right\rVert}           % operator (spectral) norm
\newcommand{\mxn}[1]{\left\lVert #1\right\rVert_{\max}}    % entrywise max modulus
\newcommand{\Fld}{\mathbb{F}}                   % exact ordered substrate field
\newcommand{\Fldi}{\mathbb{F}_{i}}              % its Gaussian extension F + iF
\newcommand{\lat}[1]{\big\lfloor #1\big\rfloor_{\kappa}} % round DOWN onto the lattice

% --- protocol-compiler ---------------------------------------------------------
\newcommand{\Spec}{\mathfrak{S}}             % protocol specification (compiler input)
\newcommand{\pcChain}{\mathfrak{R}}            % reduction chain (two typed segments)
\newcommand{\Chsrc}{\pcChain_{\mathrm{src}}}   % its Alice-side source-map segment
\newcommand{\Chrec}{\pcChain_{\mathrm{rec}}}   % its receiver-side squashing segment
\newcommand{\rvirt}{\rho^{\mathrm{v}}}       % virtual (PRE-source-map) signal state
\newcommand{\Fph}{F^{\mathrm{ph}}}           % physical (pre-reduction) POVM element
\newcommand{\Vph}{\mathcal{V}^{\mathrm{ph}}} % physical outcome alphabet
\newcommand{\Tpub}{\mathcal{T}_{\mathrm{pub}}}% public-choice alphabet
\newcommand{\Uset}{\mathcal{U}}              % residual outcome alphabet
\newcommand{\idxm}{\mathsf{idx}}             % label-to-index map
\newcommand{\Asig}{A'}                       % transmitted (signal) system -- the REAL one
\newcommand{\Avirt}{A''}                     % virtual transmitted system (source descriptor of Xi_1)
\newcommand{\Abar}{\bar{A}}                  % Alice's measured register
\newcommand{\Ahat}{\hat{A}}                  % Alice's shield register
\newcommand{\pcNmsm}{\mathcal{N}_{\mathsf A}}  % Alice's measurement channel (see the RULES block)
\newcommand{\Phip}{\Phi^{+}}                 % maximally entangled vector label
\newcommand{\Smap}{\mathfrak{s}}             % the emitted sector map (fraktur: a compile-time object)

% --- receiver ------------------------------------------------------------------
\newcommand{\Modes}{\mathbb{M}}              % declared ordered mode set
\newcommand{\Inp}{\mathbb{I}}                % declared input modes
\newcommand{\InpC}{\mathbb{I}_{\Chan}}       % channel input modes
\newcommand{\InpS}{\mathbb{I}_{\Prb}}        % probe (source) input modes
\newcommand{\Env}{\mathbb{E}}                % dilation (environment) modes
\newcommand{\Port}[1]{\mathbb{P}_{#1}}       % mode set of detector port #1
\newcommand{\Emd}[1]{E_{#1}}                 % coordinate projector on modes #1
\newcommand{\Gmat}{G}                        % instrument (transfer) matrix
\newcommand{\Fock}{\mathfrak{F}}             % Fock / second-quantization functor
\newcommand{\Symn}[1]{\mathrm{Sym}^{#1}}     % symmetric tensor power
\newcommand{\normord}[1]{{:}\mspace{2mu}#1\mspace{2mu}{:}}   % normal ordering
\newcommand{\rk}{\operatorname{rank}}        % rank (physics.sty owns \rank)
\newcommand{\SCK}{\mathsf{SC}}               % single-click pattern set
\newcommand{\Recv}{\mathfrak{R}}             % receiver instrument tuple
\newcommand{\Cat}{\mathfrak{C}}              % component catalogue
\newcommand{\Grph}{\mathfrak{E}}             % declared element graph
\newcommand{\defic}{\mathbb{L}}              % loss deficiency operator
\newcommand{\Pvac}[1]{\Pi_{\vac}(#1)}        % joint no-click element on ports #1
\newcommand{\pmc}{p_{\mc}}                   % uniform high-photon mc weight
\newcommand{\Bnd}{\beta}                     % the scalar bound function beta(n)
\newcommand{\Rot}{R}                         % planar rotation matrix
\newcommand{\Wplate}{W}                      % retarder (wave-plate) matrix
\newcommand{\Bsp}{B}                         % beamsplitter core matrix

% --- statistics ----------------------------------------------------------------
\newcommand{\honest}{p^{\mathrm{hon}}}          % honest per-round distribution
\newcommand{\ecomp}{\varepsilon_{\mathrm{comp}}}% completeness budget
\newcommand{\simplex}[1]{\Delta(#1)}            % probability simplex
\newcommand{\Chatcon}{\Chatset_{\mathrm{con}}}  % constrained outcomes
\newcommand{\Ireg}[3]{\mathcal{I}_{#1}(#2,#3)}  % regularized incomplete beta
\newcommand{\Ncount}[1]{\mathsf{N}_{#1}}        % per-outcome count
\newcommand{\qlo}{\mathbf{q}^{-}}               % box lower bounds
\newcommand{\qhi}{\mathbf{q}^{+}}               % box upper bounds
 for it, so the
%  file is read from spec/modules/<id>/.
%
%  RULES
%    * A macro that spec/notation.md catalogues is a verbatim copy of it
%      (part A). Changing one is a notation change: bump spec/notation.md.
%    * Part B holds the module-local macros. A macro defined by two modules
%      with the same expansion appears once, under the first module (in the
%      order below) that defines it. Two macros with the same name and
%      DIFFERENT expansions are both kept under a module prefix, and the
%      claims write the prefixed names:
%        \Chain  -> \redChain (reductions, \mathfrak{C}: the reduction chain)
%                   \pcChain  (protocol-compiler, \mathfrak{R}: the two-segment
%                             chain; \Chsrc and \Chrec are built on it)
%        \Nmsm   -> \redNmsm  (reductions, \mathcal{N}: a POVM measurement
%                             channel)
%                   \pcNmsm   (\mathcal{N}_{\mathsf A}: Alice's measurement
%                             channel, written by
%                             transmitter/thm:certificate-sufficiency,
%                             transmitter/rem:xreg and
%                             reductions/lem:source-transfer)
%    * Deliberately omitted: a \provenance{tag}{details} macro. Provenance
%      is a header field of every claim file (spec/claim-format.md section
%      3); a claim body that calls such a macro fails to compile, which is
%      the intended signal.
%
%  REQUIRED PACKAGES, loaded by the assembler's preamble before this file:
%    amsmath, amssymb, amsfonts, mathtools, dsfont, braket, physics, xcolor
%    (braket before physics, so that physics wins for \ket, \bra, \tr, ...).
%
%  This file contains only definitions: no \usepackage, no document body.
% =============================================================================


% =============================================================================
%  PART A -- notation macros, the ones spec/notation.md catalogues, each a
%  verbatim copy of its catalogue row; grouped under numbered headings.
% =============================================================================

% -----------------------------------------------------------------------------
% 1. State spaces and operator sets (a "system label" -> spaces of operators)
% -----------------------------------------------------------------------------
\newcommand{\Hs}[1]{\mathcal{H}(#1)}        % Hilbert space of a system label
\newcommand{\Lin}[1]{\mathcal{L}(#1)}       % bounded linear operators on H(#1)
\newcommand{\Linn}[2]{\mathcal{L}(#1,#2)}   % linear maps H(#1) -> H(#2)
\newcommand{\Pos}[1]{\mathrm{Pos}(#1)}      % positive semidefinite operators on H(#1)
\newcommand{\Density}{\mathrm{D}}           % density operators: \Density(A) = D(A)
\newcommand{\Seq}[1]{S_{=}(#1)}             % normalized states (trace = 1)
\newcommand{\Sleq}[1]{S_{\leq}(#1)}         % subnormalized states (0 <= trace <= 1)
\newcommand{\id}{\mathds{1}}                % identity operator
\newcommand{\idon}[1]{\mathds{1}_{#1}}      % identity on a named system, e.g. \idon{C}

% DEPRECATED: applies \mathcal to both arguments itself, so it takes bare
% letters, not system labels. Prefer \Linn; do not use in new contracts.
\providecommand{\Lop}[2]{\mathcal{L}(\mathcal{#1}, \mathcal{#2})}


% -----------------------------------------------------------------------------
% 2. Composite systems
% -----------------------------------------------------------------------------
\newcommand{\dsum}{\oplus}                  % direct sum of system labels
% Tensor product: use \otimes directly.
% n-fold direct sum of an operator space: \Lin{A}^{\oplus n}


% -----------------------------------------------------------------------------
% 3. Events, conditioning and classical registers (for STATES)
%    Convention:  \wedge  -> subnormalized "and-event" component, e.g. \rho_{A\wedge x}
%                 |        -> normalized conditional state,        e.g. \rho_{A|x}
% (Write these directly in subscripts; no extra macros needed.)
% -----------------------------------------------------------------------------


% -----------------------------------------------------------------------------
% 4. Photon-number subspaces / direct-sum blocks
%    Convention:  operators carry the restriction "on top", in brackets;
%                 states carry it as a \wedge subscript (see Sec. 3).
% -----------------------------------------------------------------------------
\newcommand{\Pnum}[1]{P_{#1}}               % photon-number projector (system from context)
\newcommand{\Pnumon}[2]{P_{#1}^{#2}}        % with explicit system superscript, e.g. P_n^C
\newcommand{\blk}[2]{#1^{(#2)}}             % operator block, e.g. \blk{M}{\leqN} = M^{(<=N)}
\newcommand{\leqN}{\leq N}                  % use as \blk{M}{\leqN}
\newcommand{\gtN}{> N}                      % use as \blk{M}{\gtN}
\newcommand{\eqn}{= n}                      % use as \blk{M}{\eqn}


% -----------------------------------------------------------------------------
% 5. System labels used in the protocol
% -----------------------------------------------------------------------------
\newcommand{\Chan}{C}                       % quantum-channel (Eve) input to Bob
\newcommand{\Prb}{P}                        % probe input
\newcommand{\Flag}{{\mathrm{Flg}}}          % flag system (EXISTS ONLY AFTER squashing)
\newcommand{\Loss}{{\mathrm{Lss}}}          % loss system
\newcommand{\Vc}{{\mathrm{Vcm}}}            % vacuum system
\newcommand{\Nout}{N_{\mathrm{out}}}
\newcommand{\Nin}{N_{\mathrm{in}}}


% -----------------------------------------------------------------------------
% 6. Alphabets
% -----------------------------------------------------------------------------
\newcommand{\Yset}{\mathcal{Y}}             % Bob's (coarse-grained) detection alphabet
\newcommand{\Yraw}{\mathcal{Y}_{\mathrm{raw}}}% raw click patterns {0,1}^4
\newcommand{\Yscset}{\mathcal{Y}_{\singlec}}% single-click outcomes
\newcommand{\Ymcset}{\mathcal{Y}_{\mc}}     % multi-click outcomes
\newcommand{\Tset}{\mathcal{T}_{\mathrm{set}}}% private setting set (the probe settings in the optical instance)
\newcommand{\Xset}{\mathcal{X}}             % Alice's alphabet
\newcommand{\Dset}{\{H,V,+,-\}}             % polarization / detector labels
\newcommand{\pass}{\top}                    % pass sifting


% -----------------------------------------------------------------------------
% 7. Mode / outcome abbreviations
% -----------------------------------------------------------------------------
\newcommand{\on}{{\mathrm{on}}}
\newcommand{\off}{{\mathrm{off}}}
\newcommand{\dc}{{\mathrm{dc}}}             % dark count
\newcommand{\sq}{\mathrm{sq}}               % squashed
\newcommand{\mc}{\mathrm{mc}}               % multi-click
\newcommand{\nc}{\mathrm{nc}}               % no-click
\newcommand{\singlec}{\mathrm{sc}}          % single-click
\newcommand{\vac}{\mathrm{vac}}             % vacuum
\newcommand{\kout}{{k_\mathrm{out}}}
\newcommand{\kin}{{k_\mathrm{in}}}


% -----------------------------------------------------------------------------
% 8. Maps, channels and miscellaneous operators
% -----------------------------------------------------------------------------
\newcommand{\Prob}{\mathrm{Prob}}           % probability (use \Prob[\cdot] or \Prob(\cdot))
\newcommand{\lmin}{\lambda_{\min}}          % minimum eigenvalue
\newcommand{\Wflag}{W_{\Flag}}              % flag-space weight
\newcommand{\transpose}{\mathsf{T}}
\newcommand{\innerp}[2]{\left\langle #1 \vert #2 \right\rangle}
\renewcommand{\vec}[1]{\boldsymbol{#1}}     % NB: overrides the LaTeX arrow accent


% -----------------------------------------------------------------------------
% 9. Key-rate numerics
% -----------------------------------------------------------------------------
\newcommand{\cG}{\mathcal{G}}               % protocol map
\newcommand{\cZ}{\mathcal{Z}}               % pinching map
\newcommand{\Amap}{\mathcal{A}_{\mathrm{ann}}} % announcement + measurement map
\newcommand{\Amapnb}{\mathcal{A}_{\mathrm{ann}|\mathrm{nb}}}
\newcommand{\simpl}{}                       % empty suffix for simplified block maps
\newcommand{\acSet}{\mathcal{S}_{\mathrm{acc}}} % acceptance set
\newcommand{\testE}{\mathcal{E}_{\mathrm{test}}} % test-statistics map
\newcommand{\ecLeak}{\mathrm{leak}^{\mathrm{EC}}_{\mathrm{obs}}}
\newcommand{\pgen}{p_{\mathrm{gen}}}
\newcommand{\askeyr}{R_{\infty}}             % asymptotic key rate
\newcommand{\ppass}{p_{\mathrm{pass}}}       % sifting pass probability


% -----------------------------------------------------------------------------
% 10. Finite-size analysis / Renyi entropies
% -----------------------------------------------------------------------------

% --- Round types: test vs. key generation (register T, alphabet \Fset) --------
\newcommand{\tst}{\mathrm{test}}            % test round
\newcommand{\gen}{\mathrm{gen}}             % key-generation round
\newcommand{\Fset}{\mathcal{T}}             % round-type alphabet {test, gen} (not \Tset)
\newcommand{\Chat}{\widehat{C}}             % acceptance-test register
\newcommand{\Chatset}{\widehat{\mathcal{C}}}% its alphabet

% --- Basis / key extraction functions ----------------------------------------
\newcommand{\bA}{b_{A}}                     % Alice's basis-extraction function
\newcommand{\bB}{b_{B}}                     % Bob's basis-extraction function
\newcommand{\kB}{k_{B}}                     % Bob's key-extraction function

% --- Maps --------------------------------------------------------------------
\newcommand{\Mprot}{\mathcal{M}}            % single-round protocol map
\newcommand{\Phimap}{\boldsymbol{\Phi}}     % test-statistics map
\newcommand{\Pdeph}{\mathcal{P}}            % dephasing (pinching) map on \Blreg

% --- Isometries --------------------------------------------------------------
\newcommand{\Vmsm}{V_{\mathrm{msm}}}        % measurement isometry
\newcommand{\Vann}{V_{\mathrm{ann}}}        % announcement isometry
\newcommand{\VI}{V_{I}}                     % Eve's copy of the announcements
\newcommand{\VChat}{V_{\Chat}}              % acceptance-data isometry
\newcommand{\VR}{V_{R}}                     % sifting + key map -> raw key R
\newcommand{\VZ}{V_{Z}}                     % raw key -> key register Z
\newcommand{\Wprot}{W}                      % full protocol isometry

% --- Renyi entropies ---------------------------------------------------------
% SECURITY-CRITICAL: see spec/notation.md §4. Never change an arrow to
% make an inequality work out; never convert silently between the sandwiched
% (tilde) and Petz (overline) families.
\newcommand{\Hsandup}{\widetilde{H}^{\uparrow}_{\alpha}}   % sandwiched, up
\newcommand{\Hsanddown}{\widetilde{H}^{\downarrow}_{\alpha}}% sandwiched, down
\newcommand{\Hpetzup}{\overline{H}^{\uparrow}_{\alpha}}    % Petz, up
\newcommand{\Hpetzdown}{\overline{H}^{\downarrow}_{\alpha}}% Petz, down
\newcommand{\hup}{h^{\uparrow}_{\alpha}}                   % single-round rate
\newcommand{\hqkd}{h^{\uparrow,\mathrm{QKD}}_{\alpha}}     % single-round QKD rate
\newcommand{\hlow}{\underline{h}_{\alpha}}                 % lower bound on \hqkd
% --- spec/notation.md 3.5 and 4.5 -------------------------------------------
% \hlb  is the first-order linearization lower bound on \hlow: the headline
%       key rate (the `first-order` tier of foundations/def:tiers). NOT \hlow
%       itself.
% \tnum is the instance's one declared numerical tolerance. Its
%       near-collision with the declared marginal \tau_A is disclosed in
%       spec/notation.md 3.5; bare \tau is reserved and denotes neither.
\newcommand{\hlb}{\underline{h}^{\mathrm{FO}}_{\alpha}}    % first-order lower bound on \hlow
\newcommand{\tnum}{\tau_{\mathrm{num}}}                    % declared numerical tolerance


% =============================================================================
%  PART B -- module-local macros, grouped by module. Deprecated aliases,
%  defined so that a text written with them compiles, used in no new contract: \Lop,
%  \testE, \Pdeph, \pass, \ecLeak (spec/notation.md section 3.9).
% =============================================================================

% --- meat-core -----------------------------------------------------------------
\newcommand{\SecSet}{\mathbb{S}}            % sector index set of Q
\newcommand{\Qsys}{Q}                       % declared sector-carrying measured system
\newcommand{\Vset}{\mathcal{V}}             % declared joint outcome alphabet of the Q-POVM
\newcommand{\CAset}{\mathcal{C}_{A}}        % Alice announcement alphabet
\newcommand{\CBset}{\mathcal{C}_{B}}        % Bob announcement alphabet
\newcommand{\Party}{\mathfrak{P}}           % declared party partition
\newcommand{\Oset}[1]{\mathcal{O}_{#1}}     % readout alphabet of party #1
\newcommand{\rd}[1]{\pi_{#1}}               % readout map of party #1
\newcommand{\Sadm}{\mathcal{S}_{A}}         % admissible Alice subspace
\newcommand{\nomarg}{\mathsf{none}}         % "no trusted marginal declared"
\newcommand{\Zset}{\mathcal{Z}_{\perp}}     % raw-key alphabet incl. \perp
\newcommand{\Zval}{\mathcal{Z}}             % key alphabet without \perp
\newcommand{\Etil}{\widetilde{E}}           % single-round purifying register
\newcommand{\Ehat}{\widehat{E}}             % Eve's n-round register
\newcommand{\OmAT}{\Omega_{\mathrm{AT}}}    % acceptance event
\newcommand{\epa}{\varepsilon_{\mathrm{PA}}}
\newcommand{\eev}{\varepsilon_{\mathrm{EV}}}
\newcommand{\lEC}{\lambda_{\mathrm{EC}}}
\newcommand{\Decl}{\mathfrak{D}}            % protocol declaration
\newcommand{\CList}{\mathfrak{L}}           % declared constraint list
\newcommand{\PinS}{\mathcal{P}_{\SecSet}}   % sector pinching on Q
\newcommand{\Feas}{\mathcal{W}}             % feasible set of block states
\newcommand{\Chkey}{\Chatset_{\mathrm{key}}}% key-live acceptance values
\newcommand{\Kann}{\mathcal{K}}             % announcement-class alphabet
\newcommand{\hold}{\mathsf{h}}              % declared key holder (A or B)
\newcommand{\sfD}{D}                        % classical KL divergence (base 2)
\newcommand{\Dtwo}{D}                       % quantum relative entropy (base 2)
\newcommand{\Dsand}{\widetilde{D}_{\alpha}} % sandwiched Renyi divergence
\newcommand{\idmap}{\mathrm{id}}            % identity channel
\newcommand{\Aev}{\mathcal{A}}              % Eve's per-round attack channel
\newcommand{\AttSet}{\mathfrak{A}}          % licensed attack-channel set
\newcommand{\Chacc}{\Chatset_{\mathrm{acc}}}% acceptance support of \acSet
\newcommand{\epsstar}{\varepsilon^{\star}}      % distance to an exact POVM
\newcommand{\dfH}{\Delta_{\mathrm{herm}}}       % Hermiticity defect
\newcommand{\dfP}{\Delta_{\mathrm{pos}}}        % negative part
\newcommand{\dfC}{\Delta_{\mathrm{cpl}}}        % completeness defect
\newcommand{\lmoq}{g}                       % acceptance-side gradient block
\newcommand{\lmoJ}{G}                       % feasible-side gradient block
\newcommand{\bacc}{\beta_{\mathrm{acc}}}    % certified acceptance-side minimum
\newcommand{\Itset}{\mathfrak{I}}           % finite index set of iterates (C22)

% --- reductions ----------------------------------------------------------------
\newcommand{\Bsys}{B}                        % physical (pre-squashing) system
\newcommand{\Tcl}[1]{\mathcal{T}(#1)}        % trace-class operators on H(#1)
\newcommand{\Ins}{\mathcal{I}}               % generalized-instrument branch map
\newcommand{\redNmsm}{\mathcal{N}}              % POVM measurement channel
\newcommand{\Eprot}{\mathcal{E}}             % n-round protocol channel
\newcommand{\Grnd}{\mathsf{G}}               % one round map of a protocol
\newcommand{\Kfin}{\mathsf{K}}               % end-of-protocol processing
\newcommand{\Crit}{\mathsf{sec}}             % value of a security criterion
\newcommand{\Pflag}{\Pi_{\Flag}}             % flag-sector projector on Q
\newcommand{\Pleq}{\Pi_{\leqN}}              % preserved-block projector
\newcommand{\Pgt}{\Pi_{\gtN}}                % above-cutoff projector
\newcommand{\Sph}{S^{\mathrm{ph}}}           % family of physical inputs
\newcommand{\Srch}{S^{\mathrm{rch}}}         % reachable squashed states
\newcommand{\Ssq}{S^{\sq}}                   % squashed admissible set
\newcommand{\Ffss}{F^{\mathrm{FSS}}}         % flag-state-squasher target
\newcommand{\Fwp}{F^{\mathrm{WP}}}           % multi-bound WPFSS target
\newcommand{\Mwp}{M^{\mathrm{WP}}}           % renormalized above-cutoff POVM
\newcommand{\Mdec}{{M'}^{\mathrm{ph}}}       % declared physical family on B
\newcommand{\Ppvm}{P^{\Flag}}                % flag PVM
\newcommand{\Lmp}{{\Lambda_{\mathrm{mp}}}}   % measure-and-prepare channel
\newcommand{\Lwp}{{\Lambda^{\mathrm{WP}}}}   % multi-bound WPFSS channel
\newcommand{\Lfss}{{\Lambda^{\mathrm{FSS}}}} % flag-state-squasher channel
\newcommand{\lmax}{\lambda_{\max}}           % maximum eigenvalue
\newcommand{\Lnoise}{\widetilde{\Lambda}}    % flag-noise channel
\newcommand{\cg}{\varpi}                     % declared coarse-graining map
\newcommand{\MCK}{\mathsf{MC}}               % multiclick raw-pattern set
\newcommand{\redChain}{\mathfrak{C}}            % reduction chain
\newcommand{\Lnk}{\ell}                      % a single typed link
\newcommand{\dsq}{\preceq}                      % the squashing relation
\newcommand{\opn}[1]{\left\lVert #1\right\rVert}           % operator (spectral) norm
\newcommand{\mxn}[1]{\left\lVert #1\right\rVert_{\max}}    % entrywise max modulus
\newcommand{\Fld}{\mathbb{F}}                   % exact ordered substrate field
\newcommand{\Fldi}{\mathbb{F}_{i}}              % its Gaussian extension F + iF
\newcommand{\lat}[1]{\big\lfloor #1\big\rfloor_{\kappa}} % round DOWN onto the lattice

% --- protocol-compiler ---------------------------------------------------------
\newcommand{\Spec}{\mathfrak{S}}             % protocol specification (compiler input)
\newcommand{\pcChain}{\mathfrak{R}}            % reduction chain (two typed segments)
\newcommand{\Chsrc}{\pcChain_{\mathrm{src}}}   % its Alice-side source-map segment
\newcommand{\Chrec}{\pcChain_{\mathrm{rec}}}   % its receiver-side squashing segment
\newcommand{\rvirt}{\rho^{\mathrm{v}}}       % virtual (PRE-source-map) signal state
\newcommand{\Fph}{F^{\mathrm{ph}}}           % physical (pre-reduction) POVM element
\newcommand{\Vph}{\mathcal{V}^{\mathrm{ph}}} % physical outcome alphabet
\newcommand{\Tpub}{\mathcal{T}_{\mathrm{pub}}}% public-choice alphabet
\newcommand{\Uset}{\mathcal{U}}              % residual outcome alphabet
\newcommand{\idxm}{\mathsf{idx}}             % label-to-index map
\newcommand{\Asig}{A'}                       % transmitted (signal) system -- the REAL one
\newcommand{\Avirt}{A''}                     % virtual transmitted system (source descriptor of Xi_1)
\newcommand{\Abar}{\bar{A}}                  % Alice's measured register
\newcommand{\Ahat}{\hat{A}}                  % Alice's shield register
\newcommand{\pcNmsm}{\mathcal{N}_{\mathsf A}}  % Alice's measurement channel (see the RULES block)
\newcommand{\Phip}{\Phi^{+}}                 % maximally entangled vector label
\newcommand{\Smap}{\mathfrak{s}}             % the emitted sector map (fraktur: a compile-time object)

% --- receiver ------------------------------------------------------------------
\newcommand{\Modes}{\mathbb{M}}              % declared ordered mode set
\newcommand{\Inp}{\mathbb{I}}                % declared input modes
\newcommand{\InpC}{\mathbb{I}_{\Chan}}       % channel input modes
\newcommand{\InpS}{\mathbb{I}_{\Prb}}        % probe (source) input modes
\newcommand{\Env}{\mathbb{E}}                % dilation (environment) modes
\newcommand{\Port}[1]{\mathbb{P}_{#1}}       % mode set of detector port #1
\newcommand{\Emd}[1]{E_{#1}}                 % coordinate projector on modes #1
\newcommand{\Gmat}{G}                        % instrument (transfer) matrix
\newcommand{\Fock}{\mathfrak{F}}             % Fock / second-quantization functor
\newcommand{\Symn}[1]{\mathrm{Sym}^{#1}}     % symmetric tensor power
\newcommand{\normord}[1]{{:}\mspace{2mu}#1\mspace{2mu}{:}}   % normal ordering
\newcommand{\rk}{\operatorname{rank}}        % rank (physics.sty owns \rank)
\newcommand{\SCK}{\mathsf{SC}}               % single-click pattern set
\newcommand{\Recv}{\mathfrak{R}}             % receiver instrument tuple
\newcommand{\Cat}{\mathfrak{C}}              % component catalogue
\newcommand{\Grph}{\mathfrak{E}}             % declared element graph
\newcommand{\defic}{\mathbb{L}}              % loss deficiency operator
\newcommand{\Pvac}[1]{\Pi_{\vac}(#1)}        % joint no-click element on ports #1
\newcommand{\pmc}{p_{\mc}}                   % uniform high-photon mc weight
\newcommand{\Bnd}{\beta}                     % the scalar bound function beta(n)
\newcommand{\Rot}{R}                         % planar rotation matrix
\newcommand{\Wplate}{W}                      % retarder (wave-plate) matrix
\newcommand{\Bsp}{B}                         % beamsplitter core matrix

% --- statistics ----------------------------------------------------------------
\newcommand{\honest}{p^{\mathrm{hon}}}          % honest per-round distribution
\newcommand{\ecomp}{\varepsilon_{\mathrm{comp}}}% completeness budget
\newcommand{\simplex}[1]{\Delta(#1)}            % probability simplex
\newcommand{\Chatcon}{\Chatset_{\mathrm{con}}}  % constrained outcomes
\newcommand{\Ireg}[3]{\mathcal{I}_{#1}(#2,#3)}  % regularized incomplete beta
\newcommand{\Ncount}[1]{\mathsf{N}_{#1}}        % per-outcome count
\newcommand{\qlo}{\mathbf{q}^{-}}               % box lower bounds
\newcommand{\qhi}{\mathbf{q}^{+}}               % box upper bounds

\makeatletter
\newcommand*{\ClaimHeadFont}{\bfseries}   % switched by prominence
\newcommand*{\ClaimNoteFont}{\mdseries}
\newcommand*{\ClaimShortId}{}
\newcommand*{\SetShortId}[1]{\renewcommand*{\ClaimShortId}{#1}}
\newcommand*{\RestatedName}{}
% Upright bodies throughout: italic is reserved for non-normative text.
% \@endparpenalty\@M keeps the meta line on the statement's page.
\newtheoremstyle{claim}{\topsep}{\topsep}{\normalfont\@endparpenalty\@M}{}%
  {\ClaimHeadFont}{.}{.5em}{\thmname{#1}\thmnumber{ #2}\thmnote{ {\ClaimNoteFont(#3)}}}
\newtheoremstyle{claimass}{\topsep}{\topsep}{\normalfont\@endparpenalty\@M}{}%
  {\ClaimHeadFont}{.}{.5em}{\thmname{#1}\thmnumber{ #2}%
  \ifx\ClaimShortId\@empty\else\ [\textsf{\ClaimShortId}]\fi%
  \thmnote{ {\ClaimNoteFont(#3)}}}
\newtheoremstyle{claimremark}{.5\topsep}{.5\topsep}{\small\@endparpenalty\@M}{}%
  {\itshape}{.}{.5em}{\thmname{#1}\thmnumber{ #2}\thmnote{ {\upshape(#3)}}}
\newtheoremstyle{claimreq}{\topsep}{\topsep}{\normalfont\@endparpenalty\@M}{}%
  {\color{reqAccent}\scshape}{.}{.5em}%
  {\thmname{#1}\thmnumber{ #2}\thmnote{ {\upshape(#3)}}}
\makeatother
\theoremstyle{claim}
\newtheorem{theorem}{Theorem}[section]
\newtheorem{lemma}[theorem]{Lemma}
\newtheorem{proposition}[theorem]{Proposition}
\newtheorem{corollary}[theorem]{Corollary}
\newtheorem{fact}[theorem]{Fact}
\newtheorem{definition}[theorem]{Definition}
\newtheorem*{restatedclaim}{\RestatedName}
\theoremstyle{claimass}
\newtheorem{assumption}[theorem]{Assumption}
\newtheorem*{restatedassumption}{\RestatedName}
\theoremstyle{claimremark}
\newtheorem{remark}[theorem]{Remark}
\newtheorem{example}[theorem]{Example}
\theoremstyle{claimreq}
\newtheorem{requirement}[theorem]{Requirement}
\newtheorem*{restatedrequirement}{\RestatedName}
%% --- 8. Prominence: main results a margin rule and a bold heading ----------
\newtcolorbox{mainclaim}{enhanced jigsaw, breakable, blanker,
  borderline west={1.5pt}{-8pt}{mainRule}, left=0pt, right=0pt, top=0pt, bottom=0pt,
  boxsep=0pt, before skip=\topsep, after skip=\topsep, parbox=false,
  before upper={\renewcommand*{\ClaimNoteFont}{\bfseries}}}
\newenvironment{technicalclaim}{\renewcommand*{\ClaimHeadFont}{\mdseries\color{techGrey}}%
  \renewcommand*{\ClaimNoteFont}{}}{}
%% --- 9. Non-normative text, the meta line, navigation ------------------------
\makeatletter
% Tie what follows to this paragraph: no page break, no list penalty before
% the next theorem head (the \@nobreak route a section heading uses).
\newcommand{\keepwithnext}{\par\nobreak\vspace{3pt}\@nobreaktrue
  \everypar{\if@nobreak\@nobreakfalse\clubpenalty\@M\else\clubpenalty\@clubpenalty\everypar{}\fi}}
\newcommand{\claimsummary}[2]{% #1 claim id (named destination), #2 summary
  \par\addvspace{\bigskipamount}\Needspace{6\baselineskip}%
  {\small\itshape\color{nnGrey}\noindent\hypertarget{claim-#1}{}#2\par}\keepwithnext}
\newcommand{\claimanchor}[1]{% a claim without a summary: the destination alone
  \par\addvspace{\bigskipamount}\Needspace{6\baselineskip}%
  \phantomsection\hypertarget{claim-#1}{}\keepwithnext}
\newcommand{\claimleadin}[1]{{\color{nnGrey}\noindent#1\par}\keepwithnext}
\newcommand{\claimmeta}[1]{\par\nobreak\vspace{2pt}%
  {\footnotesize\color{metaGrey}\raggedright\interlinepenalty\@M\noindent#1\par}%
  \addvspace{\medskipamount}}
% After an assumption: the sentence generated from its verification field; after
% any claim with one: its plain-text reading (user_statement). Both
% non-normative, both grey.
\newcommand{\claimverification}[1]{\par\nobreak\vspace{2pt}%
  {\footnotesize\color{metaGrey}\noindent#1\par}\nobreak}
\newcommand{\claimreading}[1]{\par\nobreak\vspace{2pt}%
  {\footnotesize\color{nnGrey}\noindent\emph{Reading:}\enspace#1\par}\nobreak}
\makeatother
\newcommand{\claimbookmark}[4]{% #1 id, #2 kind word, #3 label, #4 plain title
  \bookmark[dest=claim-#1,level=2]{#2 \getrefnumber{#3} #4}}
\newcommand{\metasep}{\unskip\enspace\textperiodcentered\enspace\ignorespaces}
\newcommand{\metalabel}[1]{\emph{#1}\enspace}
\newcommand{\sectionintro}[1]{{\color{nnGrey}\noindent#1\par}\smallskip}
\newcommand{\sectionimports}[1]{{\footnotesize\color{metaGrey}\raggedright\noindent\metalabel{Imports}#1\par}\smallskip}
\newenvironment{overview}{\par\color{nnGrey}\itshape\noindent\ignorespaces}{\par}
% References: the whole "Def. 2.1" is the link; ids are the fallback text of
% every link into another PDF (a relative URI to a named destination).
\newcommand{\mref}[2]{\hyperref[#2]{#1~\ref*{#2}}}
\newcommand{\aref}[2]{\hyperref[#2]{\textsf{#1}}}
\DeclareUrlCommand{\claimid}{\urlstyle{sf}\def\UrlBreaks{\do\/\do\:\do\-\do\.\do\_}%
  \def\UrlBigBreaks{}\def\UrlNoBreaks{}}
\newcommand{\extref}[2]{\hyperref{#1}{}{claim-#2}{\claimid{#2}}}
\newcommand{\extaref}[3]{\hyperref{#1}{}{claim-#2}{\textsf{#3}}}
\newcommand{\coderef}[2]{\href{#1}{\claimid{#2}}}
% Claim-body commands (spec/claim-format.md section 4): \xref{<full id>} prints
% the kind, title and id of a claim of another contract, linked to its stable
% destination in that contract's PDF (one \DeclareXref per cited claim below);
% \term[gloss]{phrase} marks a defined phrase for the index (the gloss feeds
% the Terms table); \tref[phrase]{text} is a use of a defined term. The
% assembler rewrites both in the body into \termdef (the definition's named
% destination term-<slug>) and \termref / \exttermref (a link to it, in this
% or another contract's PDF); the fallbacks below serve text it copies
% elsewhere, such as the digest. \iface is an interface key.
\makeatletter
\newcommand{\DeclareXref}[4]{% id, kind word, title, contract PDF
  \@namedef{xref@#1}{\hyperref{#4}{}{claim-#1}{#2 ``#3'' ({\small\claimid{#1}})}}}
\newcommand{\xref}[1]{\@ifundefined{xref@#1}%
  {\textbf{[unresolved \nolinkurl{#1}]}\@latex@warning{Unresolved xref #1}}%
  {\@nameuse{xref@#1}}}
\makeatother
\DeclareUrlCommand{\iface}{\urlstyle{tt}}
\newcommand{\term}[2][]{\emph{#2}\index{#2}}
\newcommand{\tref}[2][]{#2}
\newcommand{\termdef}[2]{\hypertarget{#1}{}\emph{#2}\index{#2}}
\newcommand{\termref}[2]{\hyperlink{#1}{#2}}
\newcommand{\exttermref}[3]{\hyperref{#1}{}{#2}{#3}}
\hypersetup{pdftitle={Proof reductions by squashing (reductions) -- contract 0.22.0}}

% Claims of other contracts cited here (\xref, [[...]]).
\DeclareXref{foundations/ass:snapped-measurement}{Assumption}{The device measures the snapped family where the run reads the snap}{../foundations/contract.pdf}
\DeclareXref{foundations/def:binary64}{Definition}{Binary64 numbers, declared values and dyadic lattices}{../foundations/contract.pdf}
\DeclareXref{foundations/def:declared-data}{Definition}{Declared data and the parties to a module boundary}{../foundations/contract.pdf}
\DeclareXref{foundations/def:domination}{Definition}{Domination of one operator family by another at block weights}{../foundations/contract.pdf}
\DeclareXref{foundations/def:evaluated-off-snap}{Definition}{A reference handed beside a snap; a block evaluated off its snap}{../foundations/contract.pdf}
\DeclareXref{foundations/def:exact-admissibility}{Definition}{Exact admissibility of a declared operator family}{../foundations/contract.pdf}
\DeclareXref{foundations/def:exact-island}{Definition}{The exact island}{../foundations/contract.pdf}
\DeclareXref{foundations/def:exact-reference}{Definition}{Reference of a measurement family, exact or enclosed}{../foundations/contract.pdf}
\DeclareXref{foundations/def:factor-enclosure}{Definition}{Factor enclosure of a block}{../foundations/contract.pdf}
\DeclareXref{foundations/def:pinching}{Definition}{Operator conventions and pinching}{../foundations/contract.pdf}
\DeclareXref{foundations/def:round-homogeneity}{Definition}{Round homogeneity of a per-round family}{../foundations/contract.pdf}
\DeclareXref{foundations/def:round-model}{Definition}{Round model, attack sequences and attack sets}{../foundations/contract.pdf}
\DeclareXref{foundations/def:sector-map}{Definition}{Sector maps and block-list operators}{../foundations/contract.pdf}
\DeclareXref{foundations/def:snap}{Definition}{Snapping declared data onto a lattice}{../foundations/contract.pdf}
\DeclareXref{foundations/lem:absorbing-domination}{Lemma}{Domination of a family that absorbs its residue in one element}{../foundations/contract.pdf}
\DeclareXref{foundations/lem:binary64-lattice}{Lemma}{Binary64 numbers on dyadic lattices}{../foundations/contract.pdf}
\DeclareXref{foundations/lem:domination}{Lemma}{Elementary properties of domination}{../foundations/contract.pdf}
\DeclareXref{foundations/lem:exact-accumulation}{Lemma}{Exactness of a sum that rounds nowhere}{../foundations/contract.pdf}
\DeclareXref{foundations/lem:rational-obstruction}{Lemma}{Rationality obstruction to the domination of a lattice family}{../foundations/contract.pdf}
\DeclareXref{foundations/lem:recentred-enclosure}{Lemma}{Zeroing the centres an enclosure cannot tell from zero}{../foundations/contract.pdf}
\DeclareXref{foundations/rem:computed-completeness}{Remark}{A complete computed sum from an incomplete family}{../foundations/contract.pdf}
\DeclareXref{foundations/rem:noise-floor}{Remark}{What domination controls, and why it needs a noise floor}{../foundations/contract.pdf}
\DeclareXref{foundations/req:boundary-admission}{Requirement}{Exact admissibility at every module boundary}{../foundations/contract.pdf}
\DeclareXref{foundations/req:consumer-decides-exactly}{Requirement}{Exact decision by every consumer}{../foundations/contract.pdf}
\DeclareXref{foundations/req:declared-lattice}{Requirement}{The declared lattice condition}{../foundations/contract.pdf}
\DeclareXref{foundations/req:declared-tolerance}{Requirement}{One declared tolerance per instance}{../foundations/contract.pdf}
\DeclareXref{foundations/req:exact-emission}{Requirement}{Exact verification before emission}{../foundations/contract.pdf}
\DeclareXref{foundations/req:norm-conversion}{Requirement}{Explicit conversion between norms}{../foundations/contract.pdf}
\DeclareXref{foundations/req:rounding-direction}{Requirement}{Rounding direction of emitted bounds}{../foundations/contract.pdf}
\DeclareXref{foundations/req:snap-reference}{Requirement}{Snaps of a measurement family: kept, or moved and disclosed}{../foundations/contract.pdf}
\DeclareXref{foundations/req:support-safe-snap}{Requirement}{Support-safe snaps of a measurement family}{../foundations/contract.pdf}
\DeclareXref{foundations/req:tolerance-scope}{Requirement}{Scope of the declared tolerance}{../foundations/contract.pdf}
\DeclareXref{foundations/req:who-snaps}{Requirement}{Snapping by the declarer only}{../foundations/contract.pdf}
\DeclareXref{meat-core/cor:composed-security}{Corollary}{Security of the protocol and its net key}{../meat-core/contract.pdf}
\DeclareXref{meat-core/def:admissible-declaration}{Definition}{Admissible declaration}{../meat-core/contract.pdf}
\DeclareXref{meat-core/def:declaration}{Definition}{The protocol declaration as the MEAT core reads it}{../meat-core/contract.pdf}
\DeclareXref{meat-core/def:hqkd-list}{Definition}{Constrained target set and single-round rate}{../meat-core/contract.pdf}
\DeclareXref{meat-core/def:roundmodel}{Definition}{Round model and licensed attack sets}{../meat-core/contract.pdf}
\DeclareXref{meat-core/frank-wolfe/req:bracket-on-reference}{Requirement}{The first-order bound evaluated on the reference companion}{../meat-core/contract.pdf}
\DeclareXref{meat-core/frank-wolfe/req:certified-evaluation}{Requirement}{The certified evaluation of the first-order bound}{../meat-core/contract.pdf}
\DeclareXref{meat-core/lem:memory-lift}{Lemma}{Approximating an attack by all-at-once attacks with finite memory}{../meat-core/contract.pdf}
\DeclareXref{meat-core/lem:range}{Lemma}{Constraint transport and range inclusion}{../meat-core/contract.pdf}
\DeclareXref{meat-core/prop:optol-nobound}{Proposition}{Instability of the rate under measurement perturbations}{../meat-core/contract.pdf}
\DeclareXref{meat-core/rem:cgate}{Remark}{Degeneracy of the licensed-class half of the gate}{../meat-core/contract.pdf}
\DeclareXref{meat-core/rem:modes-da1}{Remark}{The two source modes on the program}{../meat-core/contract.pdf}
\DeclareXref{meat-core/req:companion-admission}{Requirement}{Admission of the reference companion beside its declaration}{../meat-core/contract.pdf}
\DeclareXref{meat-core/thm:protocol-channel}{Theorem}{The declaration generates a protocol channel}{../meat-core/contract.pdf}
\DeclareXref{protocol-compiler/ass:atk}{Assumption}{Attack-set licence}{../protocol-compiler/contract.pdf}
\DeclareXref{protocol-compiler/def:chaincond}{Definition}{Interface conditions on a chain}{../protocol-compiler/contract.pdf}
\DeclareXref{protocol-compiler/def:declaration}{Definition}{Protocol declaration}{../protocol-compiler/contract.pdf}
\DeclareXref{protocol-compiler/def:homog}{Definition}{Licensed constraint list}{../protocol-compiler/contract.pdf}
\DeclareXref{protocol-compiler/def:spec}{Definition}{Protocol specification}{../protocol-compiler/contract.pdf}
\DeclareXref{protocol-compiler/req:chain-refusals}{Requirement}{Refusals of the interface conditions on a chain}{../protocol-compiler/contract.pdf}
\DeclareXref{protocol-compiler/req:link-row-provenance}{Requirement}{Omission of the link-soundness row only on a proved chain}{../protocol-compiler/contract.pdf}
\DeclareXref{receiver/ass:eig}{Assumption}{Declared slack on the computed block eigenvalues}{../receiver/contract.pdf}
\DeclareXref{receiver/ass:trust}{Assumption}{Trusted and exactly characterised receiver}{../receiver/contract.pdf}
\DeclareXref{receiver/component-catalogue/thm:catalogue-admissibility}{Theorem}{Admissibility of the catalogue entries}{../receiver/contract.pdf}
\DeclareXref{receiver/cor:compiled-instrument}{Corollary}{Receiver instrument from a well-formed graph declaration}{../receiver/contract.pdf}
\DeclareXref{receiver/cor:unconditioned-family}{Corollary}{The unconditioned receiver is a setting family}{../receiver/contract.pdf}
\DeclareXref{receiver/def:complex-enclosure}{Definition}{Complex enclosures, and enclosures of square roots and angles}{../receiver/contract.pdf}
\DeclareXref{receiver/def:evaluation-model}{Definition}{The evaluation model: floating point, enclosures and rigour}{../receiver/contract.pdf}
\DeclareXref{receiver/def:graph}{Definition}{Declared element graph}{../receiver/contract.pdf}
\DeclareXref{receiver/def:instrument}{Definition}{Receiver instrument}{../receiver/contract.pdf}
\DeclareXref{receiver/def:pmcex}{Definition}{Exact-window weight of a setting with a declared depth}{../receiver/contract.pdf}
\DeclareXref{receiver/def:residuals}{Definition}{The six residuals}{../receiver/contract.pdf}
\DeclareXref{receiver/def:setting-family}{Definition}{Setting family}{../receiver/contract.pdf}
\DeclareXref{receiver/def:slot}{Definition}{The strengthening slot}{../receiver/contract.pdf}
\DeclareXref{receiver/detector-imperfections/def:dark-composed}{Definition}{The dark-composed family}{../receiver/contract.pdf}
\DeclareXref{receiver/detector-imperfections/def:darkadm}{Definition}{Dark-count channel and its independent product form}{../receiver/contract.pdf}
\DeclareXref{receiver/detector-imperfections/def:darktol}{Definition}{Receiver tolerances deciding dark-channel admissibility}{../receiver/contract.pdf}
\DeclareXref{receiver/detector-imperfections/def:normalized-station}{Definition}{Normalized station and factored loss}{../receiver/contract.pdf}
\DeclareXref{receiver/detector-imperfections/def:pure-loss}{Definition}{Pure-loss channel on the channel modes}{../receiver/contract.pdf}
\DeclareXref{receiver/detector-imperfections/prop:dark-slack-bound}{Proposition}{Multiclick bound under a dark-count channel admitted with slack}{../receiver/contract.pdf}
\DeclareXref{receiver/detector-imperfections/prop:efficiency-normalization}{Proposition}{A receiver as uniform loss followed by its normalized station}{../receiver/contract.pdf}
\DeclareXref{receiver/detector-imperfections/prop:noise-floor}{Proposition}{Loss and dark counts make every dark-composed element definite}{../receiver/contract.pdf}
\DeclareXref{receiver/detector-imperfections/rem:d2slack}{Remark}{Cost of admitting a dark-count channel with slack}{../receiver/contract.pdf}
\DeclareXref{receiver/detector-imperfections/req:dark-channel-record}{Requirement}{The emitted dark-count channel and the scalars of its composed bound}{../receiver/contract.pdf}
\DeclareXref{receiver/detector-imperfections/req:dark-composed-reference}{Requirement}{The emitted blocks are dark-composed and carry their reference}{../receiver/contract.pdf}
\DeclareXref{receiver/detector-imperfections/thm:dark-counts}{Theorem}{Dark counts on a photon-number block-diagonal POVM}{../receiver/contract.pdf}
\DeclareXref{receiver/lem:enclosure-arithmetic}{Lemma}{Arithmetic on complex enclosures}{../receiver/contract.pdf}
\DeclareXref{receiver/prop:slotvalid}{Proposition}{Range of the strengthening-slot weights}{../receiver/contract.pdf}
\DeclareXref{receiver/req:evaluation-label}{Requirement}{The evaluation label on every artifact}{../receiver/contract.pdf}
\DeclareXref{receiver/req:exact-route-decision}{Requirement}{Deciding the exact route and enclosing the port Grams}{../receiver/contract.pdf}
\DeclareXref{receiver/req:rigorous-weight}{Requirement}{Rigorous evaluation of the multiclick weight}{../receiver/contract.pdf}
\DeclareXref{receiver/thm:click-povm}{Theorem}{Click-pattern POVM}{../receiver/contract.pdf}
\DeclareXref{receiver/thm:coarse-graining}{Theorem}{Coarse-graining of the click-pattern POVM}{../receiver/contract.pdf}
\DeclareXref{receiver/thm:compilation-well-defined}{Theorem}{Well-definedness of compilation}{../receiver/contract.pdf}
\DeclareXref{receiver/thm:emitted-bound}{Theorem}{Multiclick bounds carried by the emitted weights}{../receiver/contract.pdf}
\DeclareXref{receiver/thm:enclosed-blocks}{Theorem}{Rigorous enclosures of the photon-number blocks}{../receiver/contract.pdf}
\DeclareXref{receiver/thm:enclosed-weight}{Theorem}{The weight on enclosed bound data bounds the multiclick element}{../receiver/contract.pdf}
\DeclareXref{receiver/thm:exact-and-evaluated}{Theorem}{Exact statements, floating-point and rigorous evaluation}{../receiver/contract.pdf}
\DeclareXref{receiver/thm:representation-independence}{Theorem}{Representation independence of the click-pattern POVM}{../receiver/contract.pdf}
\DeclareXref{transmitter/ass:ens}{Assumption}{Emission of the specified ensemble}{../transmitter/contract.pdf}
\DeclareXref{transmitter/ass:seq}{Assumption}{Sequential round structure}{../transmitter/contract.pdf}
\DeclareXref{transmitter/ass:src}{Assumption}{Retention of the system A in trusted-source mode}{../transmitter/contract.pdf}
\DeclareXref{transmitter/decoy/ass:wcp}{Assumption}{Emission of the declared weak-coherent-pulse states}{../transmitter/contract.pdf}
\DeclareXref{transmitter/def:source-link}{Definition}{Source link}{../transmitter/contract.pdf}
\DeclareXref{transmitter/def:source-segment}{Definition}{Source segment and its interface conditions}{../transmitter/contract.pdf}
\DeclareXref{transmitter/def:source-spec}{Definition}{Source specification}{../transmitter/contract.pdf}
\DeclareXref{transmitter/def:srcert}{Definition}{Source-replacement certificate}{../transmitter/contract.pdf}
\DeclareXref{transmitter/thm:certificate-sufficiency}{Theorem}{Checkability and sufficiency of the source-replacement certificate}{../transmitter/contract.pdf}

\begin{document}
\begin{center}
{\LARGE\bfseries Proof reductions by squashing (reductions)}\\[4pt]
{\large Contract, version 0.22.0}
\end{center}
\begin{overview}
This contract replaces a physical receiver, and in trusted-source mode a physical source as well, by finite models a security analysis can run on, and states exactly what each replacement transfers. A \termref{term-squashing-relation}{squashing relation} asks that the adjoint of a channel pull every \termref{term-target-povm}{target POVM} element back to the physical one, together with \termref{term-carried-operators}{carried operators} that move device facts into the squashed model; because the physical \termref{term-measurement-channel}{measurement channel} then factors through the \termref{term-squashing-channel}{squashing channel}, the physical and the \termref{term-squashed-protocol}{squashed protocol} have the same output state. The core module has four main results, in the order the contract states them. Security of the squashed model over analysed inputs that contain the squashed image of every physical input transfers to the physical protocol with the same \ensuremath{\varepsilon}, for every criterion (\hyperlink{claim-reductions/thm:transfer}{Theorem~\ref*{thm:transfer}}), and the attack-set licence places every squashed physical attack in the admissible set the consuming key-rate engine restricts to, whenever each carried operator's pull-back is bounded above or below by, or equal to, a multiple of the identity, a spectral certificate (\hyperlink{claim-reductions/thm:licence}{Theorem~\ref*{thm:licence}}). In trusted-source mode the source replacement, which exchanges the specified signal ensemble for its virtual one, and the squashing reduction compose, in either order, into one \termref{term-exact-operational-reduction}{exact operational reduction} whose image lies in the admissible attack set (\hyperlink{claim-reductions/thm:two-sided}{Theorem~\ref*{thm:two-sided}}), and the target a \termref{term-reduction-instance}{reduction instance} emits is exactly admissible and squashes the \termref{term-declared-model}{declared model} with no error term (\hyperlink{claim-reductions/thm:emit}{Theorem~\ref*{thm:emit}}). The two submodules, wpfss (flag-state and weight-preserving squashers) and domination (the link that relates a dominated family to the one dominating it, kept for a declared device weight), follow the core as parts of their own and add main results, listed with these four in the front matter. Every construction is a reduction instance, emitted as a \termref{term-reduction-chain}{reduction chain} of \termref{term-typed-link}{typed links} that each carry their \termref{term-channel-witness}{channel witness}, and every condition on declared numbers is decided exactly, every declared and emitted operator entry and declared real scalar lying on a declared lattice of binary64 (double-precision) values; where the emission has to move an instance's datum, exact arithmetic having left that lattice, the emission theorem is about the moved, declared model, which nothing proved relates to the handed one; a moved class is disclosed, and a run that reads it rests on the premise that what it concludes from the emission holds of the declared model, as it does where the device measures that model. Where the declaring instance hands a reference of its measurement beside its snap, the same construction run on the reference, with no normalizer, is emitted beside the artefact as its reference companion (\hyperlink{claim-reductions/cor:reference-companion}{Corollary~\ref*{cor:reference-companion}}). The supplier's device facts, the independence of the measuring party's private setting and the round structure are assumptions the caller carries, while the soundness of the emitted constraints is a conclusion of the chain, emission and licence theorems and not an assumption; the adversary's registers are finite-dimensional, as the round model's are, and the consuming engine carries the security statement to an adversary whose memory has any dimension (\xref{meat-core/cor:composed-security}).
\end{overview}
\medskip
\noindent{\small\textcolor{nnGrey}{\itshape Grey text}, here and throughout, is non-normative: it explains, and nothing in it is assumed or proved.}

\smallskip
\noindent{\small The \href{https://github.com/fraunhofersg/certiqs-key/blob/v0.1.0/spec/modules/reductions/design.md}{design notes} (non-normative) explain why the module is built as it is and draw its processing flowchart. The dependencies between the sections are drawn in \href{graphs/overview.svg}{the section dependency matrix}, and \href{https://github.com/fraunhofersg/certiqs-key/blob/v0.1.0/spec/modules/reductions/graphs/README.md}{the graph index} opens every section's dependency graph. The \href{digest.pdf}{statements digest} holds every definition, assumption, result and requirement without the proofs. \href{https://github.com/fraunhofersg/certiqs-key/blob/v0.1.0/harness/contract-reading.md}{How to read a contract} explains this layout and its links.}

{\small\tableofcontents}

\section*{Main results}\addcontentsline{toc}{section}{Main results}
\sectionintro{The module's most important results, those its overview names as what it establishes, with every assumption each rests on, directly or through the claims it uses. Each links to its statement and proof; the statements digest holds them in full.}
\par\noindent\hyperref[thm:transfer]{\textbf{Theorem~\ref*{thm:transfer}}}\enspace Security transfer from a constrained squashed model\par\nobreak
{\small\itshape\color{nnGrey}\noindent If the analysed squashed inputs contain the squashed image of every physical input, a bound \ensuremath{\varepsilon} on any criterion of the output state of the squashed protocol holds for the physical protocol with the same \ensuremath{\varepsilon}. No property of the criterion is used.\par}\nobreak
\claimmeta{\metalabel{Rests on}\aref{A-MEAS}{ass:meas}, \aref{A-ROUND-SQ}{ass:round}, \extaref{../transmitter/contract.pdf}{transmitter/ass:seq}{A-SEQ}}
\par\noindent\hyperref[thm:licence]{\textbf{Theorem~\ref*{thm:licence}}}\enspace The attack-set licence\par\nobreak
{\small\itshape\color{nnGrey}\noindent If each operator the squashing carries pulls back to a physical operator \ensuremath{\le}, = or \ensuremath{\ge} a multiple of the identity, every squashed physical attack respects the resulting constraint list in either source mode, factor by factor for per-factor bounds. Security against such attacks transfers with the same \ensuremath{\varepsilon}.\par}\nobreak
\claimmeta{\metalabel{Rests on}\aref{A-CONST}{ass:constdecl}, \aref{A-MEAS}{ass:meas}, \aref{A-ROUND-SQ}{ass:round}, \extaref{../transmitter/contract.pdf}{transmitter/ass:seq}{A-SEQ}}
\par\noindent\hyperref[thm:two-sided]{\textbf{Theorem~\ref*{thm:two-sided}}}\enspace The two-sided reduction in trusted-source mode\par\nobreak
{\small\itshape\color{nnGrey}\noindent In trusted-source mode the source map and the squashing reduction compose, in either order, into one exact operational reduction into the admissible attack set. Security of the virtual squashed model thus transfers with the same \ensuremath{\varepsilon} to the specified physical protocol, against attacks of its type.\par}\nobreak
\claimmeta{\metalabel{Rests on}\aref{A-CONST}{ass:constdecl}, \aref{A-MEAS}{ass:meas}, \aref{A-ROUND-SQ}{ass:round}, \extaref{../transmitter/contract.pdf}{transmitter/ass:seq}{A-SEQ}, \extaref{../transmitter/contract.pdf}{transmitter/ass:src}{A-SRC}}
\par\noindent\hyperref[thm:emit]{\textbf{Theorem~\ref*{thm:emit}}}\enspace Exact emission of a reduction instance\par\nobreak
{\small\itshape\color{nnGrey}\noindent The emitted target is computed exactly, decided exactly admissible before emission, and squashes the declared model exactly; nearness to the handed model bounds nothing. On blocks the record names as resting on \hyperref{../foundations/contract.pdf}{}{claim-foundations/ass:snapped-measurement}{A-SNAP}, the run's conclusions about the device rest on that premise.\par}\nobreak
\claimmeta{\metalabel{Rests on}\aref{A-CONST}{ass:constdecl}, \aref{A-MEAS}{ass:meas}, \aref{A-ROUND-SQ}{ass:round}, \extaref{../transmitter/contract.pdf}{transmitter/ass:seq}{A-SEQ}, \extaref{../foundations/contract.pdf}{foundations/ass:snapped-measurement}{A-SNAP}}
\par\noindent\hyperref[cor:reference-companion]{\textbf{Corollary~\ref*{cor:reference-companion}}}\enspace The reference companion of a reduction instance\par\nobreak
{\small\itshape\color{nnGrey}\noindent The same instance, run on the reference in place of the snap, or on its own snapped blocks where none is handed, is a reduction instance of the emitted shape that squashes that model exactly. Its entries are enclosed from the reference, requested ones where their certificates hold there.\par}\nobreak
\claimmeta{\metalabel{Rests on}\aref{A-CONST}{ass:constdecl}, \aref{A-MEAS}{ass:meas}, \aref{A-ROUND-SQ}{ass:round}, \extaref{../transmitter/contract.pdf}{transmitter/ass:seq}{A-SEQ}, \extaref{../foundations/contract.pdf}{foundations/ass:snapped-measurement}{A-SNAP}}
\par\noindent\hyperref[prop:chain]{\textbf{Proposition~\ref*{prop:chain}}}\enspace Flag-state squashers: existence, chain and encoded weights\par\nobreak
{\small\itshape\color{nnGrey}\noindent With total certified weight below one, the physical POVM squashes to both targets, and the squashers compose exactly through a flag-noise channel. A weight the weight-preserving target encodes holds on every squashed state; under the flag-state target it holds only on reachable ones.\par}\nobreak
\claimmeta{\metalabel{Rests on}no assumption}
\par\noindent\hyperref[cor:emitted-flat]{\textbf{Corollary~\ref*{cor:emitted-flat}}}\enspace Exact emission of the flat instance\par\nobreak
{\small\itshape\color{nnGrey}\noindent The flat instance's emitted family is exactly admissible and pulls back exactly to the declared model on the channel system, its weights rounded down yet certified. Where the normalizer moved a block below the cutoff, nearness to the unrounded family bounds nothing.\par}\nobreak
\claimmeta{\metalabel{Rests on}\aref{A-CONST}{ass:constdecl}, \extaref{../receiver/contract.pdf}{receiver/ass:eig}{A-EIG}, \extaref{../receiver/contract.pdf}{receiver/ass:exact}{A-EXACT}, \aref{A-MEAS}{ass:meas}, \extaref{../receiver/contract.pdf}{receiver/ass:modes}{A-MODES}, \extaref{../receiver/contract.pdf}{receiver/ass:nosel}{A-NOSEL}, \extaref{../receiver/contract.pdf}{receiver/ass:passive}{A-PASSIVE}, \extaref{../receiver/contract.pdf}{receiver/ass:port}{A-PORT}, \aref{A-PROBE}{ass:probe}, \aref{A-RECV}{ass:recv}, \extaref{../receiver/contract.pdf}{receiver/ass:round}{A-ROUND}, \aref{A-ROUND-SQ}{ass:round}, \extaref{../transmitter/contract.pdf}{transmitter/ass:seq}{A-SEQ}, \extaref{../foundations/contract.pdf}{foundations/ass:snapped-measurement}{A-SNAP}, \extaref{../receiver/contract.pdf}{receiver/ass:thresh}{A-THRESH}, \extaref{../receiver/contract.pdf}{receiver/ass:trust}{A-TRUST}, \aref{A-TSM}{ass:tsm}, \extaref{../receiver/contract.pdf}{receiver/ass:vac}{A-VAC}}
\par\noindent\hyperref[thm:noise-channel]{\textbf{Theorem~\ref*{thm:noise-channel}}}\enspace The noise channel of a domination\par\nobreak
{\small\itshape\color{nnGrey}\noindent If G dominates F sector by sector, G squashes exactly to F with added flag sectors, or to F itself where F has flag sectors, through a noise channel that keeps each sector with probability one minus its weight. The adjoint returns each sector projector scaled by that probability.\par}\nobreak
\claimmeta{\metalabel{Rests on}\aref{A-CONST}{ass:constdecl}, \aref{A-MEAS}{ass:meas}, \aref{A-ROUND-SQ}{ass:round}, \extaref{../transmitter/contract.pdf}{transmitter/ass:seq}{A-SEQ}}

\section*{Assumptions at a glance}\addcontentsline{toc}{section}{Assumptions at a glance}
\sectionintro{Every assumption this contract states or its claims assume, with its category and the claims of this contract that assume it directly. An assumption of another contract links into that contract's PDF.}
{\footnotesize\begin{xltabular}{\linewidth}{@{}>{\raggedright\arraybackslash}p{1.6cm}>{\raggedright\arraybackslash}p{1.9cm}>{\raggedright\arraybackslash}X>{\raggedright\arraybackslash}p{2.3cm}>{\raggedright\arraybackslash}p{2.8cm}@{}}
\toprule
Assumption & Category & Title & Stated as & Assumed directly by\\
\midrule
\endhead
\aref{A-MEAS}{ass:meas} & receiver & Each measuring party's device as one POVM per round & \mref{Assumption}{ass:meas} & \mref{Lem.}{lem:opreduction-transfer}, \mref{Thm.}{thm:factorization}, \mref{Rem.}{rem:coherent}, \mref{Rem.}{rem:criteria}\\
\addlinespace
\aref{A-ROUND-SQ}{ass:round} & receiver & Round homogeneity of the physical measurement & \mref{Assumption}{ass:round} & \mref{Thm.}{thm:factorization}, \mref{Rem.}{rem:variants}, \mref{Rem.}{rem:coherent}\\
\addlinespace
\extaref{../foundations/contract.pdf}{foundations/ass:snapped-measurement}{A-SNAP} & receiver & The device measures the snapped family where the run reads the snap & \extref{../foundations/contract.pdf}{foundations/ass:snapped-measurement} & \mref{Thm.}{thm:emit}, \mref{Cor.}{cor:reference-companion}, \mref{Req.}{req:emission-record}\\
\addlinespace
\extaref{../transmitter/contract.pdf}{transmitter/ass:seq}{A-SEQ} & protocol structure & Sequential round structure & \extref{../transmitter/contract.pdf}{transmitter/ass:seq} & \mref{Thm.}{thm:factorization}, \mref{Lem.}{lem:source-transfer}\\
\addlinespace
\aref{A-TSM}{ass:tsm} & protocol structure & Output-state test statistics & \mref{Assumption}{ass:tsm} & \mref{Rem.}{rem:intertwining}, \mref{Prop.}{prop:flat-instance}, \mref{Cor.}{cor:emitted-flat}, \mref{Cor.}{cor:flat-companion}, \mref{Cor.}{cor:normalized-flat}\\
\addlinespace
\aref{A-PROBE}{ass:probe} & local randomness & The station's private setting drawn from its own declared randomness & \mref{Assumption}{ass:probe} & \mref{Prop.}{prop:flat-instance}, \mref{Cor.}{cor:emitted-flat}, \mref{Cor.}{cor:flat-companion}, \mref{Cor.}{cor:normalized-flat}\\
\addlinespace
\aref{A-CONST}{ass:constdecl} & supplied artefact & Declared constants of the constraint list & \mref{Assumption}{ass:constdecl} & \mref{Rem.}{rem:unsound}, \mref{Thm.}{thm:licence}, \mref{Rem.}{rem:notlicensed}, \mref{Cor.}{cor:entries}, \mref{Rem.}{rem:claimscope}, \mref{Rem.}{rem:published}\\
\addlinespace
\aref{A-RECV}{ass:recv} & supplied artefact & Supplier artefact whose claims hold of the device & \mref{Assumption}{ass:recv} & \mref{Rem.}{rem:notlicensed}, \mref{Prop.}{prop:route}, \mref{Rem.}{rem:absorbneeded}, \mref{Rem.}{rem:minweight}, \mref{Rem.}{rem:darkcompose}, \mref{Rem.}{rem:receiver-imports}, \mref{Prop.}{prop:flat-instance}, \mref{Cor.}{cor:emitted-flat}, \mref{Cor.}{cor:two-sector-decided}, \mref{Cor.}{cor:flat-companion}, \mref{Cor.}{cor:normalized-flat}\\
\bottomrule
\end{xltabular}}

\section*{Notation}\addcontentsline{toc}{section}{Notation}
\sectionintro{The symbols of the repository's notation catalogue (\href{../../notation.md}{spec/notation.md}) that this contract uses, in the catalogue's order, with their meaning.}
{\small\begin{xltabular}{\linewidth}{@{}>{\raggedright\arraybackslash}p{3.6cm}>{\raggedright\arraybackslash}X@{}}
\toprule
Symbol & Meaning\\
\midrule
\endhead
$\Hs{A}$ & Hilbert space of the system label $A$\\
$\Lin{A}$ & Bounded linear operators on $\mathcal{H}(A)$\\
$\Linn{A}{B}$ & Linear maps $\mathcal{H}(A)\to\mathcal{H}(B)$\\
$\Pos{A}$ & Positive semidefinite operators\\
$\Seq{A}$ & Normalized states\\
$\Sleq{A}$ & Subnormalized states\\
$\id$, $\idon{C}$ & Identity operator (contextual / named)\\
$\idmap$ & Identity channel\\
$\dsum$ & Direct sum of system labels\\
$\Pnum{n}$, $\Pnumon{n}{C}$ & Photon-number projector\\
$\blk{M}{\leqN}$ & Operator block restriction\\
$\leqN$, $\gtN$, $\eqn$ & Block indices for \texttt{\textbackslash{}blk}\\
$\Decl$ & Protocol declaration $\big(A,\tau_A;\,Q,\mathbb{S};\,\{M_x\};\,\{F_v\};\,\mathfrak{P};\,a_A,a_B;\,\phi;\,r,\mathsf{h};\,\mathfrak{L}\big)$ \textemdash{} the tuple of \extref{../meat-core/contract.pdf}{meat-core/def:declaration}: the party partition $\mathfrak{P}$ is a slot, and the marginal slot $\tau_A$ takes a state or the no-marginal value (trusted- vs untrusted-source mode)\\
$\SecSet$ & Sector index set of $Q$\\
$\Qsys$ & The declared sector-carrying measured system\\
$\Party$ & The declared party partition: the readout maps $\pi_{\mathsf A}$, $\pi_{\mathsf B}$ onto the alphabets $\mathcal{O}_{\mathsf A}$, $\mathcal{O}_{\mathsf B}$, with a party label for $A$ and, for each sector $s$, a declared tensor factorisation of $\mathcal{H}(Q_s)$ with a party label per factor\\
$\Xset$ & Outcome alphabet of the measurement on $A$ \textemdash{} Alice's under the default partition\\
$\Vset$ & The measuring parties' declared joint outcome alphabet \textemdash{} Bob's alone under the default partition\\
$\Chat$, $\Chatset$ & The acceptance-statistic register $\widehat{C}$ and its alphabet $\widehat{\mathcal{C}}$\\
$\CList$ & Declared constraint list $\{(\hat O_k,\bowtie_k,\Theta_k)\}$\\
$\Mprot$ & Single-round protocol channel $AQ\to Z\widehat{C}I$\\
$\PinS$ & Sector pinching on $Q$\\
$\Chan$, $\Prb$ & Channel input; probe input\\
$\Flag$ & Flag system \textemdash{} post-squashing only\\
$\Yset$ & Bob's coarse-grained detection alphabet, $\vert\mathcal{Y}\vert=6$\\
$\Tset$ & Bob's private setting set, weights $q_t$ (\mref{Def.}{def:joint-family}); in the optical instance the probe settings (register $T_{\mathrm{set}}$)\\
$\on$, $\off$, $\dc$, $\sq$, $\mc$, $\nc$, $\singlec$, $\vac$ & Mode / outcome abbreviations\\
$\lmin$ & Minimum eigenvalue\\
$\transpose$ & Transpose, $X^{\mathsf T}$\\
$\innerp{a}{b}$ & Inner product\\
$\vec{v}$ & Vector (bold)\\
\bottomrule
\end{xltabular}}

\section*{Terms}\addcontentsline{toc}{section}{Terms}
\sectionintro{The phrases this contract defines, and the terms of other contracts it uses, each with a one-line gloss and the definition that fixes it. In the text, a linked term opens its definition.}
{\small\begin{xltabular}{\linewidth}{@{}>{\raggedright\arraybackslash}p{3.6cm}>{\raggedright\arraybackslash}X>{\raggedright\arraybackslash}p{2.6cm}@{}}
\toprule
Term & Meaning & Defined in\\
\midrule
\endhead
\exttermref{../foundations/contract.pdf}{term-admissible-attack-sequence}{admissible attack sequence} & an attack sequence whose every round lies in the attack set of a given constraint list & \extref{../foundations/contract.pdf}{foundations/def:round-model}\\
\addlinespace[2pt]
\termref{term-admissible-datum}{admissible datum} & A datum that is Hermitian, positive semidefinite and class complete in every data sector; the set D\_adm. & \mref{Definition}{def:reduction-instance}\\
\addlinespace[2pt]
\termref{term-admitted}{admitted} & Said of a declaration with a non-empty source segment and a constraint list that meets the admission rule: licensed entries only, in trusted-source mode. & \mref{Definition}{def:chain}\\
\addlinespace[2pt]
\termref{term-analysed-strategy-set}{analysed strategy set} & The strategy set of the presentation an exact operational reduction lands in, which the analysis runs over; also called admissible. & \mref{Definition}{def:opreduction}\\
\addlinespace[2pt]
\exttermref{../foundations/contract.pdf}{term-attack-sequence}{attack sequence} & one channel per round whose output is the per-round system, with Alice's system where the adversary supplies it, and Eve's new memory & \extref{../foundations/contract.pdf}{foundations/def:round-model}\\
\addlinespace[2pt]
\exttermref{../foundations/contract.pdf}{term-attack-set}{attack set} & the channels whose adjoint maps every listed operator into the stated relation with its constant & \extref{../foundations/contract.pdf}{foundations/def:round-model}\\
\addlinespace[2pt]
\termref{term-attack-strategy}{attack strategy} & An element of the strategy set of a protocol presentation; a state, a channel or any other object. & \mref{Definition}{def:opreduction}\\
\addlinespace[2pt]
\termref{term-bare-form}{bare form} & The input form of the weight route in which the supplied preserved blocks are those of the conditional POVM before the dark-count channel, so that step 1 acts on them. & \mref{Definition}{def:weight-route}\\
\addlinespace[2pt]
\exttermref{../foundations/contract.pdf}{term-binary64-numbers}{binary64 numbers} & the finite values of the IEEE 754 binary64 format, each a dyadic rational & \extref{../foundations/contract.pdf}{foundations/def:binary64}\\
\addlinespace[2pt]
\exttermref{../foundations/contract.pdf}{term-block-list}{block list} & one square complex matrix per block of a sector map, of that block's dimension, identified with their direct sum & \extref{../foundations/contract.pdf}{foundations/def:sector-map}\\
\addlinespace[2pt]
\termref{term-carried-operators}{carried operators} & Hermitian operators on the squashed system, listed beside the target POVM, whose pull-backs a squashing relation prescribes. & \mref{Definition}{def:squashing}\\
\addlinespace[2pt]
\termref{term-certificate}{certificate} & The hypothesis that makes a transported constraint hold: an operator identity or inequality on the pull-back, or an assumption on the input states; not a source-replacement certificate. & \mref{Definition}{def:certificate}\\
\addlinespace[2pt]
\termref{term-certified-entries}{certified entries} & The constraint entries of a reduction instance, each a lattice combination of transported operators whose certificate is discharged at every admissible datum. & \mref{Definition}{def:reduction-instance}\\
\addlinespace[2pt]
\termref{term-certified-weights}{certified weights} & Non-negative numbers, one per outcome, each an operator lower bound on that outcome's POVM element above the cutoff. & \mref{Definition}{def:targets}\\
\addlinespace[2pt]
\exttermref{../receiver/contract.pdf}{term-channel-modes}{channel modes} & The declared input modes through which the adversary's signal enters. & \extref{../receiver/contract.pdf}{receiver/def:instrument}\\
\addlinespace[2pt]
\termref{term-channel-witness}{channel witness} & The squashing channel a link carries, relating its source family to its target together with the carried operators. & \mref{Definition}{def:link}\\
\addlinespace[2pt]
\termref{term-class-complete}{class complete} & Said of a datum whose elements in each class sum exactly to the class weight times the identity, in every data sector. & \mref{Definition}{def:reduction-instance}\\
\addlinespace[2pt]
\termref{term-class-weights}{class weights} & The declared probabilities of the settings: positive lattice values that sum to one exactly. & \mref{Definition}{def:joint-family}\\
\addlinespace[2pt]
\termref{term-classes}{classes} & The blocks of the joint alphabet, one per setting: every outcome paired with that setting, in the order of the outcome alphabet. & \mref{Definition}{def:joint-family}\\
\addlinespace[2pt]
\termref{term-class-marginal-flag}{class-marginal flag} & The flag of a reduction instance recording whether its construction establishes the class-marginal identity of the model at every admissible datum. & \mref{Definition}{def:reduction-instance}\\
\addlinespace[2pt]
\termref{term-class-marginal-identity}{class-marginal identity} & The operator identity that the elements of each class of a joint family sum to the class weight times the identity. & \mref{Definition}{def:joint-family}\\
\addlinespace[2pt]
\exttermref{../receiver/contract.pdf}{term-complex-enclosure}{complex enclosure} & a Gaussian-rational centre and a non-negative rational radius that bounds the distance of a complex number from the centre; a matrix is enclosed entry by entry & \extref{../receiver/contract.pdf}{receiver/def:complex-enclosure}\\
\addlinespace[2pt]
\termref{term-composable}{composable} & Said of a receiver segment in which each link's target is the next link's source family. & \mref{Definition}{def:chain}\\
\addlinespace[2pt]
\termref{term-composite}{composite} & The pair of a chain's source composite and receiver composite. & \mref{Definition}{def:chain}\\
\addlinespace[2pt]
\exttermref{../foundations/contract.pdf}{term-consumer}{consumer} & a module other than a supplier that reads declared data & \extref{../foundations/contract.pdf}{foundations/def:declared-data}\\
\addlinespace[2pt]
\exttermref{../receiver/contract.pdf}{term-dark-composed-family}{dark-composed family} & the click patterns of each setting composed with a dark-count channel, and their sums over the fibres of the declared coarse-graining & \extref{../receiver/contract.pdf}{receiver/detector-imperfections/def:dark-composed}\\
\addlinespace[2pt]
\termref{term-dark-composed-form}{dark-composed form} & The input form of the weight route in which the supplied preserved blocks are a declared complete family that the model takes below the cutoff after the dark counts, so that step 1 is not applied to them. & \mref{Definition}{def:weight-route}\\
\addlinespace[2pt]
\exttermref{../receiver/contract.pdf}{term-dark-count-channel}{dark-count channel} & A column-stochastic matrix on click patterns that can add clicks and never removes one. & \extref{../receiver/contract.pdf}{receiver/detector-imperfections/def:darkadm}\\
\addlinespace[2pt]
\termref{term-data-sector}{data sector} & A sector of a reduction instance's target whose blocks are the datum's matrices. & \mref{Definition}{def:reduction-instance}\\
\addlinespace[2pt]
\termref{term-datum}{datum} & A family of square matrices, one per data sector and outcome, with entries in the substrate and zero off each sector's support. & \mref{Definition}{def:reduction-instance}\\
\addlinespace[2pt]
\exttermref{../receiver/contract.pdf}{term-declared-coarse-graining}{declared coarse-graining} & A surjection from the raw click patterns onto the declared emitted alphabet, supplied as data by the protocol declaration. & \extref{../receiver/contract.pdf}{receiver/def:coarse}\\
\addlinespace[2pt]
\exttermref{../receiver/contract.pdf}{term-declared-dark-count-route}{declared dark-count route} & The dark-count channel is a matrix R declared as such. & \extref{../receiver/contract.pdf}{receiver/detector-imperfections/def:darkadm}\\
\addlinespace[2pt]
\exttermref{../foundations/contract.pdf}{term-declared-data}{declared data} & the data an instance states and the data a producer emits as its own declaration, with the sets such data define & \extref{../foundations/contract.pdf}{foundations/def:declared-data}\\
\addlinespace[2pt]
\termref{term-declared-datum}{declared datum} & The datum the emission normalizer outputs: exactly admissible on the declared lattice, and carried entry by entry in the emitted target. & \mref{Definition}{def:emitnorm}\\
\addlinespace[2pt]
\exttermref{../transmitter/contract.pdf}{term-declared-descriptor-family}{declared descriptor family} & The family of source descriptors for which a source link's channel identity is asserted. & \extref{../transmitter/contract.pdf}{transmitter/def:source-link}\\
\addlinespace[2pt]
\exttermref{../receiver/contract.pdf}{term-declared-element-graph}{declared element graph} & A finite directed graph whose vertices are catalogue elements, each with an entry type, a parameter tuple and a support. & \extref{../receiver/contract.pdf}{receiver/def:graph}\\
\addlinespace[2pt]
\exttermref{../receiver/contract.pdf}{term-declared-mode-space}{declared mode space} & A finite, totally ordered set of optical mode labels together with the Hilbert space whose orthonormal basis they index. & \extref{../receiver/contract.pdf}{receiver/def:modes}\\
\addlinespace[2pt]
\termref{term-declared-model}{declared model} & The handed model with its datum replaced by the declared datum, as the instance's model map defines it. & \mref{Definition}{def:emitnorm}\\
\addlinespace[2pt]
\exttermref{../foundations/contract.pdf}{term-declared-order}{declared order} & a total order on a finite alphabet, fixed together with the alphabet & \extref{../foundations/contract.pdf}{foundations/def:declared-order}\\
\addlinespace[2pt]
\exttermref{../receiver/contract.pdf}{term-declared-source-modes}{declared source modes} & The declared input modes that carry the receiver's own probe state. & \extref{../receiver/contract.pdf}{receiver/def:instrument}\\
\addlinespace[2pt]
\exttermref{../foundations/contract.pdf}{term-declared-tolerance}{declared tolerance} & a binary64 number, at least zero and below one, that an instance states in its run configuration & \extref{../foundations/contract.pdf}{foundations/def:tolerance}\\
\addlinespace[2pt]
\exttermref{../foundations/contract.pdf}{term-declared-value}{declared value} & a binary64 number stated as a datum, which denotes exactly that dyadic rational and approximates nothing & \extref{../foundations/contract.pdf}{foundations/def:binary64}\\
\addlinespace[2pt]
\termref{term-default-party-partition}{default party partition} & The assignment of every factor of every target sector to Bob, the party whose system the handed model lives on. & \mref{Definition}{def:reduction-instance}\\
\addlinespace[2pt]
\termref{term-domination-link}{domination link} & A receiver-stage typed link of kind family-dominating from a finite target family to a family it dominates at the domination weights, whose channel witness is the noise channel with added or shared flag sectors. & \mref{Definition}{def:domination-link}\\
\addlinespace[2pt]
\termref{term-domination-weights}{domination weights} & The weights, one per sector of a finite target system, at which the source family of a domination link dominates its target family. & \mref{Definition}{def:domination-link}\\
\addlinespace[2pt]
\exttermref{../receiver/contract.pdf}{term-eigenvalue-slack}{eigenvalue slack} & The declared number subtracted from each computed block eigenvalue of an exact window. & \extref{../receiver/contract.pdf}{receiver/def:pmcex}\\
\addlinespace[2pt]
\termref{term-emission-displacement}{emission displacement} & Per data sector, the largest entrywise modulus of the difference between the declared datum and the handed datum. & \mref{Definition}{def:emitnorm}\\
\addlinespace[2pt]
\termref{term-emission-normalizer}{emission normalizer} & The deterministic procedure that puts a reduction instance's exactly formed datum back on the declared lattice: it keeps the classes already there, snaps every other class for fidelity without reviving a zero row or column or, failing that, by a flagged last resort, and decides the six emission conditions last. & \mref{Definition}{def:emitnorm}\\
\addlinespace[2pt]
\exttermref{../receiver/contract.pdf}{term-enclosed-weight}{enclosed weight} & The per-setting weight formed on enclosed bound data and, at an exact window, on rigorous window bounds, with no computed eigenvalue and no eigenvalue slack. & \extref{../receiver/contract.pdf}{receiver/def:enclosed-weight}\\
\addlinespace[2pt]
\exttermref{../foundations/contract.pdf}{term-entrywise-enclosure}{entrywise enclosure} & a Gaussian-rational Hermitian centre together with symmetric non-negative rational radii that bound, entry by entry, the modulus of the difference between a matrix and that centre & \extref{../foundations/contract.pdf}{foundations/def:exact-reference}\\
\addlinespace[2pt]
\exttermref{../foundations/contract.pdf}{term-evaluated-off-its-snap}{evaluated off its snap} & said of a block of a snap in a run: no number from which the run computes its key length reads the snapped block, save the points at which those numbers are evaluated and the values other than the output length that the run fixes before its first round & \extref{../foundations/contract.pdf}{foundations/def:evaluated-off-snap}\\
\addlinespace[2pt]
\exttermref{../receiver/contract.pdf}{term-evaluation-model}{evaluation model} & The arithmetic of this contract: its exact claims evaluated in binary64 with diagnostic residuals and no rigorous enclosure of a floating-point number beyond the eigenvalue-slack assumption, beside exact rational enclosures of the blocks and a rigorous evaluation of the weights. & \extref{../receiver/contract.pdf}{receiver/def:evaluation-model}\\
\addlinespace[2pt]
\exttermref{../receiver/contract.pdf}{term-exact-flag}{exact flag} & The flag of a setting, exact or general; it may read exact only where every prefactor is one by construction. & \extref{../receiver/contract.pdf}{receiver/def:setting-family}\\
\addlinespace[2pt]
\exttermref{../foundations/contract.pdf}{term-exact-island}{exact island} & the computations of the engine, its instances' declarations included, by which a consumer of data or the declarer about to declare them decides whether they satisfy the conditions their declaration states, a declarer computes and decides the reference of a measurement family, a consumer's replacement of an emitted scalar bound is decided to lie on the bound's weakening side, a producer constructs the operator data it emits, the acceptance side of the optimisation is constructed and certified, or a lower bound on a linear minimum over the feasible side of the optimisation is certified & \extref{../foundations/contract.pdf}{foundations/def:exact-island}\\
\addlinespace[2pt]
\termref{term-exact-operational-reduction}{exact operational reduction} & A map of attack strategies between two presentations over the same output registers under which every strategy and its image produce the same full output. & \mref{Definition}{def:opreduction}\\
\addlinespace[2pt]
\exttermref{../foundations/contract.pdf}{term-exact-reference}{exact reference} & a reference all of whose radii are zero, so that each block is its Gaussian-rational centre & \extref{../foundations/contract.pdf}{foundations/def:exact-reference}\\
\addlinespace[2pt]
\exttermref{../foundations/contract.pdf}{term-exactly-admissible}{exactly admissible} & said of a declared operator family whose sector map agrees with the given one, whose blocks are exactly Hermitian and positive semidefinite, and whose sum in the declared order equals the constant times the identity on every block and rounds nowhere; the constant is one where none is named & \extref{../foundations/contract.pdf}{foundations/def:exact-admissibility}\\
\addlinespace[2pt]
\termref{term-exactly-admissible-datum}{exactly admissible datum} & An admissible datum every entry of which satisfies the lattice condition of the number model; the set D\_\ensuremath{\kappa}. & \mref{Definition}{def:reduction-instance}\\
\addlinespace[2pt]
\termref{term-exactly-admissible-on-the-declared-lattice}{exactly admissible on the declared lattice} & A family that is exactly admissible, in the sense of foundations, with the constant the invoking rule names, and whose every entry satisfies the lattice condition of this definition. & \mref{Definition}{def:exactadm}\\
\addlinespace[2pt]
\exttermref{../receiver/contract.pdf}{term-exact-window-depth}{exact-window depth} & The declared integer up to which the weight of a setting uses, block by block, the larger of the analytic bound and a computed block eigenvalue less the eigenvalue slack. & \extref{../receiver/contract.pdf}{receiver/def:pmcex}\\
\addlinespace[2pt]
\exttermref{../foundations/contract.pdf}{term-factor-enclosure}{factor enclosure} & a matrix box of columns together with an exact positive diagonal core, such that the matrix is the product of a matrix L the box contains, the core and the adjoint of L & \extref{../foundations/contract.pdf}{foundations/def:factor-enclosure}\\
\addlinespace[2pt]
\exttermref{../receiver/contract.pdf}{term-factored-loss}{factored loss} & The factor t, between the largest declared efficiency and one, by which a normalized station divides every efficiency. & \extref{../receiver/contract.pdf}{receiver/detector-imperfections/def:normalized-station}\\
\addlinespace[2pt]
\termref{term-flag-state-target}{flag-state target} & The target that keeps the POVM below the cutoff and replaces everything above it by one flag per outcome: the target of the flag-state squasher. & \mref{Definition}{def:targets}\\
\addlinespace[2pt]
\exttermref{../receiver/contract.pdf}{term-fock-space}{Fock space} & The direct sum over n of the symmetric n-photon subspaces of a finite-dimensional mode space. & \extref{../receiver/contract.pdf}{receiver/def:fock}\\
\addlinespace[2pt]
\termref{term-full-output}{full output} & The state on every output register that a protocol produces under one attack strategy, and no marginal of it. & \mref{Definition}{def:opreduction}\\
\addlinespace[2pt]
\termref{term-generalized-instrument}{generalized instrument} & A finite family of completely positive maps on one input system whose traces add up to the trace; each map is a branch, and its output system may depend on the branch. & \mref{Definition}{def:generalized-instrument}\\
\addlinespace[2pt]
\exttermref{../receiver/contract.pdf}{term-generated-dark-count-route}{generated dark-count route} & The dark-count channel is the independent product form of the declared per-port probabilities. & \extref{../receiver/contract.pdf}{receiver/detector-imperfections/def:darkadm}\\
\addlinespace[2pt]
\termref{term-handed-datum}{handed datum} & The admissible datum a reduction instance is handed, formed exactly from the numbers it was given. & \mref{Definition}{def:reduction-instance}\\
\addlinespace[2pt]
\termref{term-handed-model}{handed model} & The model map of a reduction instance at its handed datum. & \mref{Definition}{def:reduction-instance}\\
\addlinespace[2pt]
\exttermref{../foundations/contract.pdf}{term-handed-reference}{handed reference} & a reference of a snapped measurement family that its declarer passes, with its centres and radii and unaltered, to every consumer of the snap & \extref{../foundations/contract.pdf}{foundations/def:evaluated-off-snap}\\
\addlinespace[2pt]
\exttermref{../receiver/contract.pdf}{term-hermiticity-and-positivity-defect}{Hermiticity-and-positivity defect} & The larger of the spectral norm of X minus its adjoint and the negative part of the smallest eigenvalue of the Hermitian part of X: one number, measured in the spectral norm. & \extref{../receiver/contract.pdf}{receiver/def:residuals}\\
\addlinespace[2pt]
\termref{term-hypotheses-h}{hypotheses (H)} & The finite list of hypotheses a reduction instance makes about the parts of the model its datum does not describe. & \mref{Definition}{def:reduction-instance}\\
\addlinespace[2pt]
\termref{term-image}{image} & The set of analysed strategies that an exact operational reduction reaches from the physical strategy set. & \mref{Definition}{def:opreduction}\\
\addlinespace[2pt]
\exttermref{../receiver/contract.pdf}{term-independent-product-form}{independent product form} & The dark-count channel built from per-port dark-count probabilities, each port firing independently. & \extref{../receiver/contract.pdf}{receiver/detector-imperfections/def:darkadm}\\
\addlinespace[2pt]
\termref{term-indicator-sectors}{indicator sectors} & One one-dimensional sector per outcome, on which every element of a target family is one at its own outcome and zero at every other. & \mref{Definition}{def:domination-link}\\
\addlinespace[2pt]
\exttermref{../foundations/contract.pdf}{term-instance}{instance} & the party that fixes one protocol's device data, protocol data and run configuration & \extref{../foundations/contract.pdf}{foundations/def:declared-data}\\
\addlinespace[2pt]
\exttermref{../transmitter/contract.pdf}{term-interface-conditions}{interface conditions} & The conditions (S1) to (S4) a source segment satisfies: descriptor match, index preservation, the source identity, and the slots each link carries. & \extref{../transmitter/contract.pdf}{transmitter/def:source-segment}\\
\addlinespace[2pt]
\termref{term-joint-alphabet}{joint alphabet} & The product of the outcome alphabet and the settings, ordered with the outcome varying fastest. & \mref{Definition}{def:joint-family}\\
\addlinespace[2pt]
\termref{term-joint-family}{joint family} & The POVM over the joint alphabet whose element at an outcome and a setting is the class weight times the conditional element of that outcome given that setting. & \mref{Definition}{def:joint-family}\\
\addlinespace[2pt]
\exttermref{../foundations/contract.pdf}{term-lattice-exponent}{lattice exponent} & an integer k with 1 \ensuremath{\le} k \ensuremath{\le} 53 & \extref{../foundations/contract.pdf}{foundations/def:binary64}\\
\addlinespace[2pt]
\exttermref{../foundations/contract.pdf}{term-lattice-value}{lattice value} & at a lattice exponent k, a real number that is an integer multiple of two to the power \ensuremath{-}k and has modulus at most one; a complex number or a matrix is one when the real and the imaginary part of every entry are & \extref{../foundations/contract.pdf}{foundations/def:binary64}\\
\addlinespace[2pt]
\termref{term-licence-record}{licence record} & The record of the attack-set licence a link carries exactly when its constraint list is non-empty. & \mref{Definition}{def:link}\\
\addlinespace[2pt]
\termref{term-licensed}{licensed} & Said of a constraint list every entry of which has a real constant as its right-hand side, either untagged in the form its source mode admits or in the party-local form. & \mref{Definition}{def:licensed}\\
\addlinespace[2pt]
\exttermref{../foundations/contract.pdf}{term-matrix-box}{matrix box} & a matrix of Gaussian-rational centres together with non-negative rational radii that bound, entry by entry, the modulus of the difference between a rectangular complex matrix and that centre & \extref{../foundations/contract.pdf}{foundations/def:factor-enclosure}\\
\addlinespace[2pt]
\termref{term-measured-system}{measured system} & The quantum system the measuring parties receive in one round, and no other system; Bob's alone where he alone measures. & \mref{Definition}{def:protocol}\\
\addlinespace[2pt]
\termref{term-measurement-channel}{measurement channel} & The channel that measures a POVM on a quantum input and writes the outcome into a classical register. & \mref{Definition}{def:msm}\\
\addlinespace[2pt]
\termref{term-measuring-parties}{measuring parties} & The parties a declaration names as measuring; their measurements are the factors of the round's measurement channel. & \mref{Definition}{def:protocol}\\
\addlinespace[2pt]
\termref{term-model-map}{model map} & The map assigning to every admissible datum a POVM on the instance's system, under its hypotheses (H). & \mref{Definition}{def:reduction-instance}\\
\addlinespace[2pt]
\exttermref{../foundations/contract.pdf}{term-module-boundary}{module boundary} & the passage of declared data from their declarer to a consumer & \extref{../foundations/contract.pdf}{foundations/def:declared-data}\\
\addlinespace[2pt]
\exttermref{../foundations/contract.pdf}{term-moved}{moved} & a block of a snap of an operator family read as a measurement that is not kept: the snap is not decided exactly to equal on it the data it is made from & \extref{../foundations/contract.pdf}{foundations/def:snap}\\
\addlinespace[2pt]
\exttermref{../receiver/contract.pdf}{term-multiclick-set}{multiclick set} & The outcomes in which at least two ports click. & \extref{../receiver/contract.pdf}{receiver/def:clickalphabet}\\
\addlinespace[2pt]
\exttermref{../receiver/contract.pdf}{term-normalized-station}{normalized station} & A receiver instrument with the same channel modes and ports whose port Grams are those of the receiver divided by the factored loss t. & \extref{../receiver/contract.pdf}{receiver/detector-imperfections/def:normalized-station}\\
\addlinespace[2pt]
\termref{term-n-round-protocol}{$n$-round protocol} & A channel that in each of n rounds applies one measurement channel to the round's measured system and then processing that touches no measured system, followed by end-of-protocol processing. & \mref{Definition}{def:protocol}\\
\addlinespace[2pt]
\exttermref{../transmitter/contract.pdf}{term-obligation-record}{obligation record} & What a theorem-supplied source link carries in place of a decided identity: its kind, the soundness claim, the assumptions the obligation carries, and the family member its descriptors are. & \extref{../transmitter/contract.pdf}{transmitter/def:source-link}\\
\addlinespace[2pt]
\termref{term-outcome-alphabet}{outcome alphabet} & The finite ordered set of device outcomes that a private setting is joined to. & \mref{Definition}{def:joint-family}\\
\addlinespace[2pt]
\termref{term-output-registers}{output registers} & The tuple of registers on which the full outputs of a protocol presentation live. & \mref{Definition}{def:opreduction}\\
\addlinespace[2pt]
\termref{term-outside-the-sector-support}{outside the sector support} & Said of the elements of a datum at outcomes outside a data sector's support: exactly zero, and read, rounded or repaired by no step. & \mref{Definition}{def:reduction-instance}\\
\addlinespace[2pt]
\termref{term-party-local-form}{party-local form} & The form of a constraint entry scoped to one factor of a product measured system: tagged party-local and read as an operator inequality against the marginal actually present. & \mref{Definition}{def:licensed}\\
\addlinespace[2pt]
\exttermref{../receiver/contract.pdf}{term-per-setting-weight}{per-setting weight} & The high-photon multiclick weight of one setting: its collapse value, its window and tail, or its exact window. & \extref{../receiver/contract.pdf}{receiver/def:pmc}\\
\addlinespace[2pt]
\termref{term-physical-inputs}{physical inputs} & The input states the physical protocol is analysed over, one set per admissible adversary system. & \mref{Definition}{def:sound}\\
\addlinespace[2pt]
\termref{term-physical-strategy-set}{physical strategy set} & The strategy set of the presentation an exact operational reduction starts from. & \mref{Definition}{def:opreduction}\\
\addlinespace[2pt]
\termref{term-post-processed-datum}{post-processed datum} & The preserved photon-number blocks of the joint family formed by steps 1 and 2 of the weight route. & \mref{Definition}{def:weight-route}\\
\addlinespace[2pt]
\termref{term-protocol-presentation}{protocol presentation} & A set of attack strategies, a tuple of output registers, and the full output each strategy produces. & \mref{Definition}{def:opreduction}\\
\addlinespace[2pt]
\termref{term-pull-back}{pull-back} & The image of an operator on the squashed system under the adjoint of the squashing channel. & \mref{Definition}{def:squashing}\\
\addlinespace[2pt]
\exttermref{../receiver/contract.pdf}{term-pure-loss-channel}{pure-loss channel} & The channel on the channel system that sends every channel mode through a beamsplitter of transmittance t with vacuum in its other input and discards the reflected light. & \extref{../receiver/contract.pdf}{receiver/detector-imperfections/def:pure-loss}\\
\addlinespace[2pt]
\termref{term-reachable}{reachable} & Said of a squashed state that is the image of some physical input under the n-fold squashing channel. & \mref{Definition}{def:sound}\\
\addlinespace[2pt]
\termref{term-reachable-set-statement}{reachable-set statement} & The set of reachable squashed states of a link, for its channel witness and the declared family of inputs. & \mref{Definition}{def:link}\\
\addlinespace[2pt]
\termref{term-receiver-composite}{receiver composite} & The composition of the channel witnesses of a chain's receiver segment. & \mref{Definition}{def:chain}\\
\addlinespace[2pt]
\exttermref{../receiver/contract.pdf}{term-receiver-instrument}{receiver instrument} & A declared mode space with channel and probe inputs, disjoint detector ports and a contractive instrument matrix. & \extref{../receiver/contract.pdf}{receiver/def:instrument}\\
\addlinespace[2pt]
\termref{term-receiver-segment}{receiver segment} & The list of receiver-stage typed links of a reduction chain. & \mref{Definition}{def:chain}\\
\addlinespace[2pt]
\exttermref{../receiver/contract.pdf}{term-receiver-tolerances}{receiver tolerances} & The eight declared numbers at which the receiver decides its exact conditions on declared data, bounds its residuals and spends its eigenvalue slack; not the instance's declared tolerance. & \extref{../receiver/contract.pdf}{receiver/def:tolmap}\\
\addlinespace[2pt]
\termref{term-receiver-stage-typed-link}{receiver-stage typed link} & A tuple of ten slots relating a source family to a target through a channel witness, with its carried operators, reachable-set statement, transported entries, licence record and status. & \mref{Definition}{def:link}\\
\addlinespace[2pt]
\termref{term-reduction-chain}{reduction chain} & A pair of a source segment and a receiver segment of typed links. & \mref{Definition}{def:chain}\\
\addlinespace[2pt]
\termref{term-reduction-instance}{reduction instance} & A handed model with a finite datum and hypotheses about the rest, a target built from the datum and scalar values, a squashing relation with a transported set, and certified entries, all established for every admissible datum. & \mref{Definition}{def:reduction-instance}\\
\addlinespace[2pt]
\exttermref{../foundations/contract.pdf}{term-reference}{reference} & a family of Hermitian positive semidefinite block lists complete at the constant on every block, read together with an entrywise enclosure of every block of every element & \extref{../foundations/contract.pdf}{foundations/def:exact-reference}\\
\addlinespace[2pt]
\exttermref{../foundations/contract.pdf}{term-residual-family}{residual family} & on a block with weight q, the family (G \ensuremath{-} (1\ensuremath{-}q)F)/q where q is positive, and G itself where q is zero & \extref{../foundations/contract.pdf}{foundations/def:domination}\\
\addlinespace[2pt]
\exttermref{../transmitter/contract.pdf}{term-retained-register-sector-map}{retained-register sector map} & An optional matched pair of sector structures, on Alice's retained system and on the descriptor the certificate is written against. & \extref{../transmitter/contract.pdf}{transmitter/def:source-link}\\
\addlinespace[2pt]
\exttermref{../foundations/contract.pdf}{term-round-model}{round model} & n rounds of a per-round system S, with the attack sequences and attack sets below & \extref{../foundations/contract.pdf}{foundations/def:round-model}\\
\addlinespace[2pt]
\exttermref{../foundations/contract.pdf}{term-round-homogeneous}{round-homogeneous} & one fixed list of objects describes every round, and every round behaves as they describe whatever the other rounds do & \extref{../foundations/contract.pdf}{foundations/def:round-homogeneity}\\
\addlinespace[2pt]
\exttermref{../foundations/contract.pdf}{term-rounds-nowhere}{rounds nowhere} & said of an accumulation every one of whose binary64 additions is exact & \extref{../foundations/contract.pdf}{foundations/def:declared-order}\\
\addlinespace[2pt]
\termref{term-routed-weights}{routed weights} & The weights the weight route assigns to the joint outcomes: at a bounded letter, the class weight times that setting's certified weight sharpened by the dark counts, rounded down onto the lattice; zero elsewhere. & \mref{Definition}{def:weight-route}\\
\addlinespace[2pt]
\termref{term-scalar-sector}{scalar sector} & A sector of a reduction instance's target on which every target element is a scalar value times the identity. & \mref{Definition}{def:reduction-instance}\\
\addlinespace[2pt]
\termref{term-scalar-value}{scalar value} & The number in the unit interval, satisfying the lattice condition, that a scalar sector assigns to an outcome; the values of a sector sum exactly to one. & \mref{Definition}{def:reduction-instance}\\
\addlinespace[2pt]
\termref{term-scoped-to-a-factor}{scoped to a factor} & Said of a constraint entry whose operator acts on one factor of a product measured system and which carries the party-local tag. & \mref{Definition}{def:licensed}\\
\addlinespace[2pt]
\exttermref{../foundations/contract.pdf}{term-sector-map}{sector map} & a block kind together with a non-empty, ordered list of block identifiers, each with a dimension & \extref{../foundations/contract.pdf}{foundations/def:sector-map}\\
\addlinespace[2pt]
\termref{term-sector-set}{sector set} & The finite ordered list of a reduction instance's target sectors, data and scalar, each with a declared dimension. & \mref{Definition}{def:reduction-instance}\\
\addlinespace[2pt]
\termref{term-security-criterion}{security criterion} & A family of functions, one per admissible adversary system, of the state of the protocol's output registers together with that system. & \mref{Definition}{def:criterion}\\
\addlinespace[2pt]
\exttermref{../receiver/contract.pdf}{term-setting-family}{setting family} & A finite set of settings, each with a photon-number block-diagonal POVM whose joint no-click elements have the prefactor form, and an exact flag. & \extref{../receiver/contract.pdf}{receiver/def:setting-family}\\
\addlinespace[2pt]
\termref{term-settings}{settings} & The finite ordered set of values of a private per-round setting of the receiver. & \mref{Definition}{def:joint-family}\\
\addlinespace[2pt]
\termref{term-single-bound-case}{single-bound case} & The weight-preserving target in which only one certified weight is non-zero. & \mref{Definition}{def:targets}\\
\addlinespace[2pt]
\exttermref{../foundations/contract.pdf}{term-snap}{snap} & at a lattice exponent k, data of the same shape as the given data, over the same alphabet in the same order, with every entry a lattice value at k, that satisfy exactly the conditions the declaration of the given data states & \extref{../foundations/contract.pdf}{foundations/def:snap}\\
\addlinespace[2pt]
\termref{term-sound}{sound} & Said of a family of squashed states, or of a constraint, that contains every reachable state; not the soundness of a transmitter source link. & \mref{Definition}{def:sound}\\
\addlinespace[2pt]
\exttermref{../transmitter/contract.pdf}{term-soundness-registry}{soundness registry} & The closed list, emitted with every artefact, of admitted source-link kinds with their soundness claims and carried assumptions. & \extref{../transmitter/contract.pdf}{transmitter/def:source-segment}\\
\addlinespace[2pt]
\termref{term-source-composite}{source composite} & The composition of a chain's source links, the identity when the source segment is empty. & \mref{Definition}{def:chain}\\
\addlinespace[2pt]
\termref{term-source-family}{source family} & The system and POVM a receiver-stage link starts from. & \mref{Definition}{def:link}\\
\addlinespace[2pt]
\exttermref{../transmitter/contract.pdf}{term-source-link}{source link} & A link of a reduction chain's source stage that maps one Alice-side ensemble to another with the same labels and weights. & \extref{../transmitter/contract.pdf}{transmitter/def:source-link}\\
\addlinespace[2pt]
\exttermref{../transmitter/contract.pdf}{term-source-model}{source model} & A submodule's named definition of the signal ensemble as mathematical objects, possibly infinite dimensional, rather than as matrices. & \extref{../transmitter/contract.pdf}{transmitter/def:source-spec}\\
\addlinespace[2pt]
\exttermref{../transmitter/contract.pdf}{term-source-segment}{source segment} & The ordered, possibly empty list of source links leading from the virtual ensemble to the specified one. & \extref{../transmitter/contract.pdf}{transmitter/def:source-segment}\\
\addlinespace[2pt]
\exttermref{../transmitter/contract.pdf}{term-source-specification}{source specification} & The declared description of Alice's side of a protocol: alphabet, retained system with sectors, measurement, signal ensemble and source mode. & \extref{../transmitter/contract.pdf}{transmitter/def:source-spec}\\
\addlinespace[2pt]
\exttermref{../transmitter/contract.pdf}{term-source-replacement-certificate}{source-replacement certificate} & An operator on Alice's retained system and the virtual transmitted system whose marginal on her side is her declared marginal and which, measured on her side by her POVM, leaves the weighted virtual members, with retention attested. & \extref{../transmitter/contract.pdf}{transmitter/def:srcert}\\
\addlinespace[2pt]
\termref{term-squashed-protocol}{squashed protocol} & The protocol with the physical system and POVM of each round replaced by the squashed system and the target POVM, all other processing unchanged. & \mref{Definition}{def:protocol}\\
\addlinespace[2pt]
\termref{term-squashed-sequence}{squashed sequence} & The attack sequence followed, in every round, by the squashing channel on the measured system and the identity on everything else. & \mref{Definition}{def:roundmodel}\\
\addlinespace[2pt]
\termref{term-squashing-channel}{squashing channel} & A channel from the physical system to the squashed system whose adjoint realises a squashing relation. & \mref{Definition}{def:squashing}\\
\addlinespace[2pt]
\termref{term-squashing-relation}{squashing relation} & The relation in which the adjoint of a channel takes every target POVM element to the physical one and every carried operator to its prescribed physical operator. & \mref{Definition}{def:squashing}\\
\addlinespace[2pt]
\termref{term-stability-clause}{stability clause} & The requirement that no hypothesis in (H) read an entry of the datum and that parts (I2) to (I4) hold at every admissible datum. & \mref{Definition}{def:reduction-instance}\\
\addlinespace[2pt]
\termref{term-stage-ordered}{stage-ordered} & Said of a chain in which every source-stage link precedes every receiver-stage link. & \mref{Definition}{def:chain}\\
\addlinespace[2pt]
\exttermref{../receiver/contract.pdf}{term-strengthening-slot}{strengthening slot} & The list of disjoint raw-outcome sets with per-setting high-photon lower-bound weights that the artifact carries. & \extref{../receiver/contract.pdf}{receiver/def:slot}\\
\addlinespace[2pt]
\exttermref{../foundations/contract.pdf}{term-structural-support}{structural support} & of a block of an element: the set of indices whose row or column holds an entry that is not structurally zero & \extref{../foundations/contract.pdf}{foundations/def:snap}\\
\addlinespace[2pt]
\termref{term-substrate}{substrate} & An ordered field F with \ensuremath{\mathbb{Q}} \ensuremath{\subseteq} F \ensuremath{\subseteq} \ensuremath{\mathbb{R}}, together with its Gaussian extension F + iF, in which every exact decision of this contract is taken. & \mref{Definition}{def:substrate}\\
\addlinespace[2pt]
\termref{term-support}{support} & The subset of the alphabet on which a data sector's elements may be nonzero; it meets every class. & \mref{Definition}{def:reduction-instance}\\
\addlinespace[2pt]
\termref{term-system-setting}{system setting} & The setting of this contract: a separable, possibly infinite-dimensional physical system and finite-dimensional, declared target spaces and side registers other than the adversary's. & \mref{Definition}{def:system-setting}\\
\addlinespace[2pt]
\termref{term-target-family}{target family} & The target POVM of a reduction instance at a datum, assigned one sector block at a time from the datum and the scalar values. & \mref{Definition}{def:reduction-instance}\\
\addlinespace[2pt]
\termref{term-target-povm}{target POVM} & The POVM on the squashed system that a squashing relation pulls back to the physical POVM. & \mref{Definition}{def:squashing}\\
\addlinespace[2pt]
\exttermref{../transmitter/contract.pdf}{term-theorem-constructor}{theorem constructor} & A named claim of a submodule that supplies a source link's channel for each member of a declared family of descriptors. & \extref{../transmitter/contract.pdf}{transmitter/def:source-link}\\
\addlinespace[2pt]
\termref{term-total-weight}{total weight} & The sum of the certified weights over all outcomes. & \mref{Definition}{def:targets}\\
\addlinespace[2pt]
\exttermref{../transmitter/contract.pdf}{term-transmitted-system}{transmitted system} & The system Alice sends into the channel, carrying the specified signal ensemble. & \extref{../transmitter/contract.pdf}{transmitter/def:source-spec}\\
\addlinespace[2pt]
\exttermref{../transmitter/contract.pdf}{term-transported-constraint-entries}{transported constraint entries} & Entries of a consumer's constraint list carried across a link, each with the certificate case that discharges it. & \extref{../transmitter/contract.pdf}{transmitter/def:source-link}\\
\addlinespace[2pt]
\termref{term-transported-set}{transported set} & The finite set of sector-block-diagonal Hermitian operators on the target, every sector projector among them, that a reduction instance transports with their pull-backs. & \mref{Definition}{def:reduction-instance}\\
\addlinespace[2pt]
\exttermref{../transmitter/contract.pdf}{term-trusted-source-mode}{trusted-source mode} & The source mode in which Alice's retained system never reaches the adversary, with a declared marginal and, above dimension one, a certificate. & \extref{../transmitter/contract.pdf}{transmitter/def:source-spec}\\
\addlinespace[2pt]
\termref{term-tuple-chained}{tuple-chained} & Said of a receiver segment in which each link's target-side carried operators are the next link's source-side ones. & \mref{Definition}{def:chain}\\
\addlinespace[2pt]
\termref{term-typed-link}{typed link} & A link of a reduction chain, at the source stage or the receiver stage, carrying what is needed to re-check it. & \mref{Definition}{def:link}\\
\addlinespace[2pt]
\exttermref{../receiver/contract.pdf}{term-uniform-weight-of-an-entry}{uniform weight of an entry} & The minimum over the settings of a slot entry's weights; it bounds the entry at every setting. & \extref{../receiver/contract.pdf}{receiver/def:slot}\\
\addlinespace[2pt]
\exttermref{../transmitter/contract.pdf}{term-untrusted-source-mode}{untrusted-source mode} & The source mode with no certificate and no marginal, in which the adversary may have prepared Alice's retained system. & \extref{../transmitter/contract.pdf}{transmitter/def:source-spec}\\
\addlinespace[2pt]
\exttermref{../transmitter/contract.pdf}{term-virtual-ensemble}{virtual ensemble} & The finite-dimensional ensemble the source segment starts from and the certificate is written against. & \extref{../transmitter/contract.pdf}{transmitter/def:source-segment}\\
\addlinespace[2pt]
\exttermref{../transmitter/contract.pdf}{term-virtual-transmitted-system}{virtual transmitted system} & The transmitted system of the virtual ensemble, where the source segment starts. & \extref{../transmitter/contract.pdf}{transmitter/def:source-segment}\\
\addlinespace[2pt]
\termref{term-weight-route}{weight route} & The three ordered steps that carry a receiver artifact's blocks and certified weights through the dark counts and the private setting into the flat instance. & \mref{Definition}{def:weight-route}\\
\addlinespace[2pt]
\termref{term-weight-preserving-target}{weight-preserving target} & The target that keeps the POVM below the cutoff and gives each outcome a flag block carrying its certified weight on every flag and, in addition, one minus the total weight on its own flag: the target of the multi-bound weight-preserving flag-state squasher. & \mref{Definition}{def:targets}\\
\bottomrule
\end{xltabular}}

\clearpage
\part{Mathematics}

\section{Squashing relations and their statistics}\label{group-reductions/def:system-setting}
\sectionintro{A \termref{term-squashing-relation}{squashing relation} replaces the physical measurement on Bob's system by a target measurement on a finite system, through a channel whose adjoint pulls every target element back to the physical one, and a tuple of \termref{term-carried-operators}{carried operators} travels along with prescribed pull-backs. The relation is equivalent to agreement of all statistics, side system included (\hyperlink{claim-reductions/lem:statistics}{Lemma~\ref*{lem:statistics}}), and it turns a device fact about a carried operator's pull-back into an operator constraint on every state the \termref{term-squashing-channel}{squashing channel} produces, the fact being the constraint's certificate (\hyperlink{claim-reductions/cor:pullback}{Corollary~\ref*{cor:pullback}}). The setting comes first, a separable physical system with finite declared targets and side registers (\hyperlink{claim-reductions/def:system-setting}{Definition~\ref*{def:system-setting}}), and with it the infinite-dimensional bookkeeping behind adjoints and tensor channels (\hyperlink{claim-reductions/lem:sep}{Lemma~\ref*{lem:sep}}); the section ends with the \termref{term-measurement-channel}{measurement channel}, through which the relation becomes a factorisation of channels (\hyperlink{claim-reductions/lem:msmfact}{Lemma~\ref*{lem:msmfact}}).}

% claim reductions/def:system-setting
\claimsummary{reductions/def:system-setting}{The physical system Bob's receiver acts on has a separable Hilbert space, possibly infinite-dimensional, while every target space and every side register other than the adversary's is finite-dimensional with declared dimensions. The squashing relations and their transport are stated in it.}
\claimbookmark{reductions/def:system-setting}{Definition}{def:system-setting}{The system setting of the reductions}
\claimleadin{%
Bob's physical receiver acts on a system whose Hilbert space may be infinite
dimensional; in every optical application it is a \exttermref{../receiver/contract.pdf}{term-fock-space}{Fock space}. Every
other system of the analysis, apart from the adversary's registers, is
finite-dimensional and declared. Separability is no restriction on such a
receiver: a Fock space over countably many modes has the countable
occupation-number basis. It is used only to type the systems, and for a
branch register whose dimension is nowhere declared it is a convention.
}
\begin{definition}[The system setting of the reductions]\label{def:system-setting}
The
\termdef{term-system-setting}{system setting}
of this contract is the following:
\begin{enumerate}[label=(\alph*),leftmargin=2.4em]
\item the Hilbert space $\Hs{\Bsys}$ of the system $\Bsys$ that Bob's
  physical receiver acts on is separable and may be infinite-dimensional;
\item every target space $\Hs{\Qsys}$ is finite-dimensional, with declared
  dimensions;
\item every side register other than the adversary's registers, that is,
  Alice's retained system and the classical registers, is finite-dimensional,
  with declared dimensions.
\end{enumerate}
\end{definition}
\claimmeta{\metalabel{Used by}\mref{Lem.}{lem:sep}, \mref{Thm.}{thm:two-sided}, \mref{Cor.}{cor:identity-instance}, \mref{Rem.}{rem:exactboundary}, \mref{Prop.}{prop:flat-instance}, \mref{Cor.}{cor:normalized-flat}
\metasep newly derived
\metasep accepted}

% claim reductions/lem:sep
\claimsummary{reductions/lem:sep}{States separate bounded operators on a separable space; trace-preserving completely positive maps of any output dimension have unital completely positive adjoints; and tensor products of channels into finite targets, or from finite inputs, are channels. Transport arguments rest on this.}
\claimbookmark{reductions/lem:sep}{Lemma}{lem:sep}{Separation and tensor adjoints in the infinite-dimensional case}
\claimleadin{%
Bob's physical receiver acts on a system whose Hilbert space may be infinite
dimensional; in every optical application it is a \exttermref{../receiver/contract.pdf}{term-fock-space}{Fock space}. Adjoints
are therefore set up for a separable input and an output of any dimension;
tensor products of channels are set up for separable inputs and
finite-dimensional outputs, the case of the squashing channels, and for
finite-dimensional inputs and separable outputs, the case of an attack on
finite registers and of a source channel on a finite virtual system. The \termref{term-system-setting}{system setting} and the finite
dimension of the adversary's registers are what make part (ii) available and
what make the trace pairing on a side system separate points; the statistics
equivalence of Lemma~\ref{lem:statistics},
Corollaries~\ref{cor:pullback} and~\ref{cor:transported} and the licence,
Theorem~\ref{thm:licence}, use them. Besides its hypotheses and notation, the
Setting paragraph carries assertions of the lemma: that the adjoint it names
exists, is unique and is completely positive, unital where the map is trace
preserving and sub-unital where it is trace non-increasing. Steps 2 to 4 of
the proof prove them, beside conclusions (i) to (iii).
}
\begin{lemma}[Separation and tensor adjoints in the infinite-dimensional case]\label{lem:sep}
\emph{Setting.} $\Bsys$ is the system Bob's physical receiver acts on and
$\Qsys$ the system it is squashed to. $\Hs{\Bsys}$ is separable and may be
infinite-dimensional, while $\Hs{\Qsys}$ and every side register are
finite-dimensional (Definition~\ref{def:system-setting}, and, for the
adversary's registers, the \exttermref{../foundations/contract.pdf}{term-round-model}{round model} of
\xref{foundations/def:round-model}, every register of which is
finite-dimensional). $\Tcl{\Bsys}$ denotes the
trace-class operators on $\Hs{\Bsys}$; a \emph{channel}
$\Lambda\colon\Tcl{\Bsys}\to\Tcl{\Qsys}$ is a completely positive
trace-preserving map, and its \emph{adjoint}
$\Lambda^{\dagger}\colon\Lin{\Qsys}\to\Lin{\Bsys}$ is the unital completely
positive map fixed by
\begin{equation}\label{eq:pairing}
  \tr\big[X\,\Lambda(\rho)\big]=\tr\big[\Lambda^{\dagger}(X)\,\rho\big]
  \qquad\text{for all }X\in\Lin{\Qsys},\ \rho\in\Tcl{\Bsys}.
\end{equation}
More generally, if $\Lambda\colon\Tcl{\Bsys}\to\Tcl{\Bsys'}$, for a Hilbert
space $\Hs{\Bsys'}$ of any dimension, is only completely positive and trace
non-increasing ($\tr\Lambda(\rho)\le\tr\rho$ for $\rho\succeq0$), then
\eqref{eq:pairing} for all $X\in\Lin{\Bsys'}$ and $\rho\in\Tcl{\Bsys}$ fixes
a completely positive map $\Lambda^{\dagger}\colon\Lin{\Bsys'}\to\Lin{\Bsys}$,
again called its \emph{adjoint}, and $\Lambda^{\dagger}$ is sub-unital:
$\idon{\Bsys}-\Lambda^{\dagger}(\idon{\Bsys'})\succeq0$. Where such a
$\Lambda$ is moreover trace preserving, $\Lambda^{\dagger}$ is unital,
$\Lambda^{\dagger}(\idon{\Bsys'})=\idon{\Bsys}$, whatever the dimension of
$\Hs{\Bsys'}$.

Let $\Hs{\Bsys_1},\dots,\Hs{\Bsys_n}$ be separable and
$\Hs{\Qsys_1},\dots,\Hs{\Qsys_n}$ finite-dimensional.
\begin{enumerate}[label=(\roman*),leftmargin=3.0em]
\item \emph{(Separation.)} If $X\in\Lin{\Bsys_1}$ satisfies $\tr[X\rho]=0$ for
  every $\rho\in\Seq{\Bsys_1}$, then $X=0$.
\item \emph{(Tensor channels and their adjoints.)} If $\Lambda_i\colon
  \Tcl{\Bsys_i}\to\Tcl{\Qsys_i}$ are channels, then exactly one trace-norm
  continuous linear map $\Lambda_1\otimes\dots\otimes\Lambda_n$ from
  $\Tcl{\Bsys_1\cdots\Bsys_n}$, the trace class of the completed tensor
  product of the $\Hs{\Bsys_i}$, to $\Tcl{\Qsys_1\cdots\Qsys_n}$ takes every
  $\rho_1\otimes\dots\otimes\rho_n$ with $\rho_i\in\Tcl{\Bsys_i}$ to
  $\Lambda_1(\rho_1)\otimes\dots\otimes\Lambda_n(\rho_n)$; it is a channel,
  and $(\Lambda_1\otimes\dots\otimes\Lambda_n)^{\dagger}
  =\Lambda_1^{\dagger}\otimes\dots\otimes\Lambda_n^{\dagger}$ on all of
  $\Lin{\Qsys_1\cdots\Qsys_n}$. The same holds with any of the $\Lambda_i$
  replaced by $\idmap_{R_i}$ for a finite-dimensional $R_i$.
\item \emph{(Tensor channels from finite inputs.)} Let
  $\Hs{R_1},\dots,\Hs{R_n}$ be finite-dimensional,
  $\Hs{\Bsys'_1},\dots,\Hs{\Bsys'_n}$ separable, of any dimension, and
  $\Phi_i\colon\Lin{R_i}\to\Tcl{\Bsys'_i}$ completely positive and trace
  preserving. Then exactly one linear map $\Phi_1\otimes\dots\otimes\Phi_n$
  from $\Lin{R_1\cdots R_n}$ to $\Tcl{\Bsys'_1\cdots\Bsys'_n}$, the trace
  class of the completed tensor product of the $\Hs{\Bsys'_i}$, takes every
  $X_1\otimes\dots\otimes X_n$ with $X_i\in\Lin{R_i}$ to
  $\Phi_1(X_1)\otimes\dots\otimes\Phi_n(X_n)$, and it is completely positive
  and trace preserving. The identity map $\idmap_{R}$ of a finite-dimensional
  $R$ is such a $\Phi_i$.
\end{enumerate}
\end{lemma}
\claimmeta{\metalabel{Uses}\mref{Def.}{def:system-setting}
\metasep \metalabel{Used by}\mref{Def.}{def:squashing}, \mref{Lem.}{lem:statistics}, \mref{Lem.}{lem:msmfact}, \mref{Prop.}{prop:platform}, \mref{Cor.}{cor:transported}, \mref{Prop.}{prop:preorder}, \mref{Prop.}{prop:tensor}, \mref{Lem.}{lem:master}, \mref{Thm.}{thm:mixture}, \mref{Lem.}{lem:party-local}, \mref{Thm.}{thm:licence}, \mref{Rem.}{rem:publishedlicence}, \mref{Lem.}{lem:precomposition}, \mref{Lem.}{lem:source-transfer}, \mref{Thm.}{thm:instrument-instance}, \mref{Lem.}{lem:precomposed-instance}, \mref{Thm.}{thm:emitted-product}, \mref{Lem.}{lem:row-relaxation}, \mref{Prop.}{prop:chain}
\metasep newly derived
\metasep accepted}

\begin{proof}
\emph{Idea.} Separation is tested on rank-one operators. The adjoint is built
on rank-one operators by the Riesz representation theorem and extended to the
whole trace class because finite-rank operators are dense in trace norm. For
(ii) the output is finite-dimensional, so the tensor map is defined from the
tensor product of the adjoints on a basis of matrix units; it acts on products
as required, and products are dense. For (iii) the input is
finite-dimensional, so the tensor map is defined on a basis of matrix units
of the input, and its complete positivity is read off its Choi operator, a
tensor product of positive operators.

\emph{Step 1 (separation, part (i)).} Every $\tau\in\Tcl{\Bsys_1}$ is a
complex linear combination of at most four positive trace-class operators
(split $\tau$ into Hermitian and anti-Hermitian parts, then each into positive
and negative parts), and every non-zero positive trace-class operator is a
positive multiple of a state. Hence $\tr[X\tau]=0$ for every
$\tau\in\Tcl{\Bsys_1}$. If $X\ne0$ there are unit vectors
$\ket{\psi},\ket{\varphi}$ with $\bra{\varphi}X\ket{\psi}\ne0$; take
$\tau=\ketbra{\psi}{\varphi}$, which is trace class of rank one, and
$\tr[X\tau]=\bra{\varphi}X\ket{\psi}\ne0$, a contradiction.

\emph{Step 2 (the adjoint exists and is unique).} Let
$\Lambda\colon\Tcl{\Bsys}\to\Tcl{\Bsys'}$ be completely positive and trace
non-increasing; a channel is the case $\Bsys'=\Qsys$ with the trace preserved.
It is bounded in trace norm. Write $\tau\in\Tcl{\Bsys}$ as $\tau=H_1+iH_2$ with
the Hermitian $H_1=\tfrac12(\tau+\tau^{\dagger})$ and
$H_2=\tfrac1{2i}(\tau-\tau^{\dagger})$, each of trace norm at most
$\|\tau\|_1$ since $\|\tau^{\dagger}\|_1=\|\tau\|_1$, and each $H_l$ as
$H_l^+-H_l^-$ with positive trace-class $H_l^{\pm}$ and
$\|H_l^+\|_1+\|H_l^-\|_1=\|H_l\|_1$ (its spectral decomposition). A positive
operator has trace norm equal to its trace, and $\Lambda$ is positive (as every
completely positive map is) and trace non-increasing, so
\begin{equation*}
  \|\Lambda(\tau)\|_1\le\sum_{l,\pm}\tr\Lambda(H_l^{\pm})
  \le\sum_{l,\pm}\tr H_l^{\pm}=\|H_1\|_1+\|H_2\|_1\le2\|\tau\|_1 .
\end{equation*}
Now fix $X\in\Lin{\Bsys'}$. For $\psi,\chi\in\Hs{\Bsys}$ the rank-one operator
$\ketbra{\psi}{\chi}$ has trace norm $\|\psi\|\,\|\chi\|$, so
$s(\chi,\psi)\coloneqq\tr[X\,\Lambda(\ketbra{\psi}{\chi})]$ is linear in
$\psi$, conjugate-linear in $\chi$, and bounded:
$|s(\chi,\psi)|\le\|X\|\,\|\Lambda(\ketbra{\psi}{\chi})\|_1
\le2\|X\|\,\|\psi\|\,\|\chi\|$. For fixed $\psi$ the map
$\chi\mapsto\overline{s(\chi,\psi)}$ is a bounded linear functional on
$\Hs{\Bsys}$, so the Riesz representation theorem gives exactly one vector
$\eta_\psi\in\Hs{\Bsys}$ with $s(\chi,\psi)=\innerp{\chi}{\eta_\psi}$ for
every $\chi$. Set $Y\psi\coloneqq\eta_\psi$: uniqueness makes $Y$ linear, and
$\chi=\eta_\psi$ gives $\|Y\psi\|\le2\|X\|\,\|\psi\|$, so $Y\in\Lin{\Bsys}$.
Then $\tr[Y\ketbra{\psi}{\chi}]=\bra{\chi}Y\ket{\psi}
=\tr[X\,\Lambda(\ketbra{\psi}{\chi})]$, which is \eqref{eq:pairing} for $Y$ at
every rank-one $\rho$, hence by linearity at every finite-rank $\rho$.
Finite-rank operators are trace-norm dense in $\Tcl{\Bsys}$ (truncate the
singular value decomposition $\rho=\sum_js_j\ketbra{u_j}{w_j}$, whose tail
beyond $J$ terms has trace norm at most $\sum_{j>J}s_j\to0$ since
$\sum_js_j=\|\rho\|_1$), and both sides of \eqref{eq:pairing} are trace-norm
continuous in $\rho$, bounded by $2\|X\|\,\|\rho\|_1$ and
$\|Y\|\,\|\rho\|_1$; so \eqref{eq:pairing} holds for $Y$ at every
$\rho\in\Tcl{\Bsys}$. Two operators satisfying \eqref{eq:pairing} for the same
$X$ have a difference that (i), read on $\Hs{\Bsys}$, shows is zero, so
$\Lambda^{\dagger}(X)\coloneqq Y$ defines the one map on $\Lin{\Bsys'}$ that
satisfies \eqref{eq:pairing}. It is linear, since
$c\Lambda^{\dagger}(X)+c'\Lambda^{\dagger}(X')$ satisfies \eqref{eq:pairing}
for $cX+c'X'$.

\emph{Step 3 (the adjoint is completely positive).} Let $k\ge1$, let
$\{\ket{a}\}_{a=1}^{k}$ be the standard basis of $\mathbb{C}^k$, let
$Z=\sum_{a,b=1}^{k}\ketbra{a}{b}\otimes Z_{ab}\succeq0$ on
$\mathbb{C}^k\otimes\Hs{\Bsys'}$, so that
$(\idmap_k\otimes\Lambda^{\dagger})(Z)=\sum_{a,b}\ketbra{a}{b}\otimes
\Lambda^{\dagger}(Z_{ab})$, and let
$\xi=\sum_{a=1}^{k}\ket{a}\otimes\xi_a\in\mathbb{C}^k\otimes\Hs{\Bsys}$, so
that $\dyad{\xi}=\sum_{a,b}\ketbra{a}{b}\otimes\ketbra{\xi_a}{\xi_b}$. By
\eqref{eq:pairing} at the rank-one operators $\ketbra{\xi_b}{\xi_a}$,
\begin{equation*}
  \bra{\xi}(\idmap_k\otimes\Lambda^{\dagger})(Z)\ket{\xi}
  =\sum_{a,b}\bra{\xi_a}\Lambda^{\dagger}(Z_{ab})\ket{\xi_b}
  =\sum_{a,b}\tr\big[Z_{ab}\,\Lambda(\ketbra{\xi_b}{\xi_a})\big]
  =\tr\big[Z\,(\idmap_k\otimes\Lambda)(\dyad{\xi})\big],
\end{equation*}
which is non-negative: $T\coloneqq(\idmap_k\otimes\Lambda)(\dyad{\xi})$ is
positive by complete positivity of $\Lambda$, and for its spectral
decomposition $T=\sum_mt_m\dyad{e_m}$, with $t_m\ge0$ and $\{e_m\}$
orthonormal, $\tr[ZT]=\sum_mt_m\bra{e_m}Z\ket{e_m}\ge0$, the trace being taken
in an orthonormal basis that extends $\{e_m\}$ by vectors of the kernel of
$T$. As $\xi$ is arbitrary, $(\idmap_k\otimes\Lambda^{\dagger})(Z)\succeq0$.

\emph{Step 4 (sub-unitality, and unitality where the trace is preserved).}
Sub-unitality is trace non-increase read through \eqref{eq:pairing}: for
every $\psi\in\Hs{\Bsys}$,
$\bra{\psi}\Lambda^{\dagger}(\idon{\Bsys'})\ket{\psi}
=\tr[\Lambda(\dyad{\psi})]\le\tr\dyad{\psi}=\|\psi\|^2$, so
$\idon{\Bsys}-\Lambda^{\dagger}(\idon{\Bsys'})\succeq0$. Where $\Lambda$ is
trace preserving, unitality is trace preservation read the same way: we have
$\tr[\Lambda^{\dagger}(\idon{\Bsys'})\rho]=\tr[\Lambda(\rho)]=\tr\rho$ for
every $\rho\in\Tcl{\Bsys}$, so $\Lambda^{\dagger}(\idon{\Bsys'})=\idon{\Bsys}$
by (i), read on $\Hs{\Bsys}$. Nothing in this step reads the dimension of
$\Hs{\Bsys'}$; a channel is the case $\Bsys'=\Qsys$.

\emph{Step 5 (the tensor map, part (ii)).} Fix orthonormal bases of the
finite-dimensional $\Hs{\Qsys_i}$. The tensor products of their matrix units
form a basis of $\Lin{\Qsys_1\cdots\Qsys_n}$, so there is exactly one linear
map $\Psi\colon\Lin{\Qsys_1\cdots\Qsys_n}\to\Lin{\Bsys_1\cdots\Bsys_n}$ with
$\Psi(X_1\otimes\dots\otimes X_n)=\Lambda_1^{\dagger}(X_1)\otimes\dots\otimes
\Lambda_n^{\dagger}(X_n)$ for every product operator, both sides being
multilinear in $(X_1,\dots,X_n)$; it is the map the statement writes
$\Lambda_1^{\dagger}\otimes\dots\otimes\Lambda_n^{\dagger}$. For an
orthonormal basis $\{e_r\}$ of $\Hs{\Qsys_1\cdots\Qsys_n}$ and
$\rho\in\Tcl{\Bsys_1\cdots\Bsys_n}$ put
\begin{equation*}
  \Lambda(\rho)\coloneqq\sum_{r,s}\tr\big[\Psi(\ketbra{e_r}{e_s})\,\rho\big]\,
  \ketbra{e_s}{e_r}.
\end{equation*}
Then $\tr[X\Lambda(\rho)]=\tr[\Psi(X)\rho]$ for all
$X\in\Lin{\Qsys_1\cdots\Qsys_n}$: the identity holds at every
$X=\ketbra{e_r}{e_s}$, and both sides are linear in $X$. $\Lambda$ is linear,
and bounded in trace norm,
$\|\Lambda(\rho)\|_1\le\sum_{r,s}\|\Psi(\ketbra{e_r}{e_s})\|\,\|\rho\|_1$.

\emph{Step 6 (action on products, and uniqueness).} For
$\rho=\rho_1\otimes\dots\otimes\rho_n$ and $X=X_1\otimes\dots\otimes X_n$ the
traces below factor, and the pairing of Step~5 with \eqref{eq:pairing} for
each $\Lambda_i$ gives
\begin{equation*}
\begin{split}
  \tr\big[X\,\Lambda(\rho)\big]&=\tr\big[\Psi(X)\,\rho\big]
  =\prod_{i}\tr\big[\Lambda_i^{\dagger}(X_i)\,\rho_i\big]\\
  &=\prod_{i}\tr\big[X_i\,\Lambda_i(\rho_i)\big]
  =\tr\big[X\,\big(\Lambda_1(\rho_1)\otimes\dots\otimes\Lambda_n(\rho_n)\big)\big].
\end{split}
\end{equation*}
Product operators span $\Lin{\Qsys_1\cdots\Qsys_n}$, and an operator $\sigma$
on this finite-dimensional space with $\tr[X\sigma]=0$ for every $X$ is zero
(take $X=\sigma^{\dagger}$), so $\Lambda(\rho_1\otimes\dots\otimes\rho_n)
=\Lambda_1(\rho_1)\otimes\dots\otimes\Lambda_n(\rho_n)$. Product trace-class
operators span a trace-norm dense subspace of $\Tcl{\Bsys_1\cdots\Bsys_n}$:
finite-rank operators are dense there (Step~2); a rank-one
$\ketbra{\psi}{\chi}$ is the trace-norm limit of $\ketbra{\psi_m}{\chi_m}$ for
finite sums of product vectors $\psi_m\to\psi$ and $\chi_m\to\chi$, which are
dense in the completed tensor product, since
$\|\ketbra{\psi}{\chi}-\ketbra{\psi_m}{\chi_m}\|_1
\le\|\psi-\psi_m\|\,\|\chi\|+\|\psi_m\|\,\|\chi-\chi_m\|$; and each
$\ketbra{\psi_m}{\chi_m}$ is a finite sum of products
$\ketbra{\alpha_1}{\beta_1}\otimes\dots\otimes\ketbra{\alpha_n}{\beta_n}$. So
$\Lambda$ is the only trace-norm continuous linear map that acts on products
as the statement says, and $\Lambda_1\otimes\dots\otimes\Lambda_n\coloneqq
\Lambda$.

\emph{Step 7 (the tensor map is a channel).} Every $\Lambda_i^{\dagger}$ is
unital, so $\Psi(\idon{})=\idon{}$ and
$\tr\Lambda(\rho)=\tr[\Psi(\idon{})\rho]=\tr\rho$. For complete positivity let
$k\ge1$, let $\omega=\sum_{a,b=1}^{k}\ketbra{a}{b}\otimes\omega_{ab}\succeq0$
be trace class on $\mathbb{C}^k\otimes\Hs{\Bsys_1\cdots\Bsys_n}$, and let
$z=\sum_{a=1}^{k}\ket{a}\otimes z_a\in\mathbb{C}^k\otimes
\Hs{\Qsys_1\cdots\Qsys_n}$. With $P\coloneqq(\idmap_k\otimes\Psi)(\dyad{z})$
the pairing of Step~5 gives
\begin{equation*}
  \bra{z}(\idmap_k\otimes\Lambda)(\omega)\ket{z}
  =\sum_{a,b}\bra{z_a}\Lambda(\omega_{ab})\ket{z_b}
  =\sum_{a,b}\tr\big[\Psi(\ketbra{z_b}{z_a})\,\omega_{ab}\big]
  =\tr\big[P\,\omega\big],
\end{equation*}
and $\tr[P\omega]\ge0$ once $P\succeq0$, by the spectral decomposition of
$\omega$ as in Step~3. $P$ is bounded, so $P\succeq0$ follows from
$\bra{\xi}P\ket{\xi}\ge0$ for every $\xi$ in the dense subspace of finite sums
of product vectors $\ket{a}\otimes\beta_1\otimes\dots\otimes\beta_n$. Such a
$\xi$ lies in $\mathbb{C}^k\otimes S_1\otimes\dots\otimes S_n$ for
finite-dimensional subspaces $S_i\subseteq\Hs{\Bsys_i}$. Let $J_i$ be the
inclusion of $S_i$ into $\Hs{\Bsys_i}$ and
$J\coloneqq\idon{}\otimes J_1\otimes\dots\otimes J_n$, with $\idon{}$ on
$\mathbb{C}^k$. The map $\Psi_i(X)\coloneqq J_i^{\dagger}\Lambda_i^{\dagger}(X)J_i$
from $\Lin{\Qsys_i}$ to $\Lin{S_i}$ is completely positive, since
$(\idmap_m\otimes\Psi_i)(Z)=(\idon{}\otimes J_i)^{\dagger}
(\idmap_m\otimes\Lambda_i^{\dagger})(Z)(\idon{}\otimes J_i)$ for every $m$,
and expanding $\dyad{z}$ in product operators gives
$J^{\dagger}PJ=(\idmap_k\otimes\Psi_1\otimes\dots\otimes\Psi_n)(\dyad{z})$.
Between these finite-dimensional spaces,
$\idmap_k\otimes\Psi_1\otimes\dots\otimes\Psi_n$ is the composite, over
$i=1,\dots,n$, of the maps that apply $\Psi_i$ to the $i$-th of the last $n$
tensor factors and the identity to the others; after a unitary reordering of
the factors each of these is $\idmap_m\otimes\Psi_i$, for $m$ the finite
dimension of the other factors, and so positive. Hence
$\bra{\xi}P\ket{\xi}=\bra{\xi}J^{\dagger}PJ\ket{\xi}\ge0$, with $\xi$ read as a
vector of $\mathbb{C}^k\otimes S_1\otimes\dots\otimes S_n$ on the right.

\emph{Step 8 (its adjoint).} By the pairing of Step~5, $\Psi$ satisfies
\eqref{eq:pairing} for the channel $\Lambda_1\otimes\dots\otimes\Lambda_n$,
whose input space is separable and whose output space is finite-dimensional,
so by Step~2 it is the one map that does, the adjoint, and it is unital and
completely positive (Steps~3 and~4). The $\idmap_{R_i}$ case is the same
computation with $\Lambda_i=\idmap$.

\emph{Step 9 (tensor products of trace-class operators).} Let
$T_i\in\Tcl{\Bsys'_i}$ have the singular value decompositions
$T_i=\sum_ks_{ik}\ketbra{u_{ik}}{w_{ik}}$, with $s_{ik}\ge0$,
$\sum_ks_{ik}=\|T_i\|_1$ and orthonormal families $\{u_{ik}\}_k$,
$\{w_{ik}\}_k$. The series
$\sum_{k_1,\dots,k_n}\big(\prod_is_{ik_i}\big)\,
\ketbra{u_{1k_1}\otimes\dots\otimes u_{nk_n}}{w_{1k_1}\otimes\dots\otimes w_{nk_n}}$
of rank-one operators on the completed tensor product, each of trace norm
$\prod_is_{ik_i}$, converges absolutely in trace norm, its sum of norms being
$\prod_i\|T_i\|_1$; its sum is the operator $T_1\otimes\dots\otimes T_n$,
which is therefore trace class, with
$\tr(T_1\otimes\dots\otimes T_n)=\sum_{k_1,\dots,k_n}\prod_is_{ik_i}
\innerp{w_{ik_i}}{u_{ik_i}}=\prod_i\tr T_i$. Where every $T_i\succeq0$ its
spectral decomposition is such an expansion with $u_{ik}=w_{ik}$, so the
partial sums are positive and so is their trace-norm limit.

\emph{Step 10 (the tensor map from finite inputs, part (iii)).} Fix
orthonormal bases of the $\Hs{R_i}$ and write $E^{(i)}_{ab}$ for the matrix
units. Their tensor products $E_{\mathbf{a}\mathbf{b}}\coloneqq
E^{(1)}_{a_1b_1}\otimes\dots\otimes E^{(n)}_{a_nb_n}$ form a basis of
$\Lin{R_1\cdots R_n}$, so exactly one linear map $\Phi$ takes each
$E_{\mathbf{a}\mathbf{b}}$ to
$\Phi_1(E^{(1)}_{a_1b_1})\otimes\dots\otimes\Phi_n(E^{(n)}_{a_nb_n})$, a
trace-class operator by Step~9; both sides of the product rule being
multilinear in $(X_1,\dots,X_n)$, $\Phi$ takes every product
$X_1\otimes\dots\otimes X_n$ to $\Phi_1(X_1)\otimes\dots\otimes\Phi_n(X_n)$,
and a linear map is fixed by its values on a basis, which gives uniqueness.
It is trace preserving: by Step~9 and trace preservation of each $\Phi_i$,
$\tr\Phi(E_{\mathbf{a}\mathbf{b}})=\prod_i\tr E^{(i)}_{a_ib_i}
=\tr E_{\mathbf{a}\mathbf{b}}$, and both sides are linear. For complete
positivity, let $\Omega_i\coloneqq\sum_a\ket{a}\otimes\ket{a}$ on two copies
of $\Hs{R_i}$; complete positivity of $\Phi_i$ makes
$J_i\coloneqq\sum_{a,b}E^{(i)}_{ab}\otimes\Phi_i(E^{(i)}_{ab})
=(\idmap_{R_i}\otimes\Phi_i)(\dyad{\Omega_i})$ positive, and after a unitary
reordering of the factors
$J\coloneqq\sum_{\mathbf{a},\mathbf{b}}E_{\mathbf{a}\mathbf{b}}\otimes
\Phi(E_{\mathbf{a}\mathbf{b}})$ is $J_1\otimes\dots\otimes J_n$, positive by
Step~9. Let $k\ge1$ and let $x=\sum_{\mathbf{a}}\ket{y_{\mathbf{a}}}\otimes
\ket{\mathbf{a}}$ be a vector of $\mathbb{C}^k\otimes\Hs{R_1\cdots R_n}$, with
$y_{\mathbf{a}}\in\mathbb{C}^k$, and put
$V\coloneqq\sum_{\mathbf{a}}\ketbra{y_{\mathbf{a}}}{\mathbf{a}}$. Then
\begin{equation*}
  (\idmap_k\otimes\Phi)(\dyad{x})
  =\sum_{\mathbf{a},\mathbf{b}}\ketbra{y_{\mathbf{a}}}{y_{\mathbf{b}}}\otimes
   \Phi(E_{\mathbf{a}\mathbf{b}})
  =(V\otimes\idon{})\,J\,(V\otimes\idon{})^{\dagger}\;\succeq\;0 .
\end{equation*}
Every positive operator on the finite-dimensional
$\mathbb{C}^k\otimes\Hs{R_1\cdots R_n}$ is a finite sum of such $\dyad{x}$, so
$\idmap_k\otimes\Phi$ is positive for every $k$, and $\Phi$ is completely
positive. The identity map $\idmap_R$ is linear, completely positive and
trace preserving from $\Lin{R}$ to $\Tcl{R}=\Lin{R}$.
\end{proof}

% claim reductions/def:squashing
\claimsummary{reductions/def:squashing}{A squashing relation asks that the adjoint of a channel pull every target POVM element back to the physical one, and every carried operator back to its prescribed physical partner. The carried operators let one reduction move device facts into the squashed model as well as outcome statistics.}
\claimbookmark{reductions/def:squashing}{Definition}{def:squashing}{Squashing relation}
\claimleadin{%
Squashing replaces one \emph{model} of the receiver by another. A model is what
the analysis runs on: the system Bob receives and the POVM he measures on it.
The physical model is $(\Bsys,M)$, the one the security statement is wanted
for; the squashed model $(\Qsys,F)$ lives on a system a proof technique can
handle. The squashed model must reproduce the physical model's outcome
statistics \emph{exactly}, since a protocol reads Bob's device through nothing
else (Definition~\ref{def:protocol}), and it must be at least as generous to
the adversary. Handing $\Lambda$ to the adversary accomplishes that: the
squashing channel is counted as part of her attack rather than of Bob's
device, so every attack on the physical model, followed by $\Lambda$, is an
attack on the squashed model. Both verbs the definition introduces, squashing
and transporting, run from the physical side to the target side, in the
direction of the symbol $\dsq$ and of $\Lambda$; the identities that define
them are written with the adjoint $\Lambda^{\dagger}$, which runs the other
way.
}
\begin{definition}[Squashing relation]\label{def:squashing}
Let $M=\{M_v\}_{v\in\Vset}$ be a POVM on $\Bsys$ and $F=\{F_v\}_{v\in\Vset}$ a
POVM on $\Qsys$ over the same finite alphabet $\Vset$, and let
$\hat O=(\hat O_k)_{k=1}^{m}$ and $O=(O_k)_{k=1}^{m}$ be tuples of Hermitian
operators $\hat O_k\in\Lin{\Qsys}$, $O_k\in\Lin{\Bsys}$. For a channel
$\Lambda\colon\Tcl{\Bsys}\to\Tcl{\Qsys}$, with adjoint $\Lambda^{\dagger}$ as
in Lemma~\ref{lem:sep}, write
\begin{equation}\label{eq:squashing}
  (M;O)\;\dsq_{\Lambda}\;(F;\hat O)
  \quad:\Longleftrightarrow\quad
  \Lambda^{\dagger}(F_v)=M_v\ \ \forall v\in\Vset
  \ \ \text{and}\ \
  \Lambda^{\dagger}(\hat O_k)=O_k\ \ \forall k\in\{1,\dots,m\},
\end{equation}
and call \eqref{eq:squashing} a
\termdef{term-squashing-relation}{squashing relation}:
$M$ \emph{squashes to} the
\termdef{term-target-povm}{target POVM}
$F$ via the
\termdef{term-squashing-channel}{squashing channel}
$\Lambda$, which \emph{transports} $O$ to $\hat O$. The operators
$\hat O_k$ are the relation's
\termdef{term-carried-operators}{carried operators}.
The empty case $m=0$ is written $M\dsq_{\Lambda}F$, and $\dsq$ replaces
$\dsq_{\Lambda}$ when only the existence of a squashing channel is asserted.
We call $\Lambda^{\dagger}(\hat O_k)$ the
\termdef{term-pull-back}{pull-back}
of $\hat O_k$.
\end{definition}
\claimmeta{\metalabel{Uses}\mref{Lem.}{lem:sep}
\metasep \metalabel{Used by}\mref{Lem.}{lem:statistics}, \mref{Rem.}{rem:dictionary}, \mref{Rem.}{rem:carried-operators}, \mref{Rem.}{rem:channel-choice}, \mref{Prop.}{prop:preorder}, \mref{Lem.}{lem:linear}, \mref{Prop.}{prop:pinching}, \mref{Prop.}{prop:tensor}, \mref{Lem.}{lem:obstructions}, \mref{Lem.}{lem:top}, \mref{Prop.}{prop:obstruction}, \mref{Def.}{def:reduction-instance}, \mref{Lem.}{lem:precomposed-instance}, \mref{Cor.}{cor:reference-companion}, \mref{Lem.}{lem:row-relaxation}, \mref{Def.}{def:domination-link}, \mref{Thm.}{thm:noise-channel}
\metasep newly derived
\metasep accepted}

% claim reductions/lem:statistics
\claimsummary{reductions/lem:statistics}{A squashing relation holds exactly when the physical and the squashed model give the same statistics for every POVM element and carried operator, correlations with a finite side system included. Two channels realising one relation therefore have adjoints that agree on the span of those operators.}
\claimbookmark{reductions/lem:statistics}{Lemma}{lem:statistics}{Squashing as statistics equivalence}
\claimleadin{%
The two lists in \eqref{eq:squashing} contain the same kind of object, an
operator on $\Qsys$ whose pull-back is prescribed, and differ in role, not in
type: the $F_v$ are what Bob measures and must form a POVM, while the
$\hat O_k$ are never measured and serve only to carry statements about the
physical device into the squashed model.
}
\begin{lemma}[Squashing as statistics equivalence]\label{lem:statistics}
Let $M,F,O,\hat O$ be as in Definition~\ref{def:squashing} and let
$\Lambda\colon\Tcl{\Bsys}\to\Tcl{\Qsys}$ be a channel. The following are
equivalent.
\begin{enumerate}[label=(\roman*),leftmargin=3.0em]
\item $(M;O)\dsq_{\Lambda}(F;\hat O)$.
\item For every $\rho\in\Seq{\Bsys}$, every $v$ and every $k$,
  $\tr[M_v\rho]=\tr[F_v\Lambda(\rho)]$ and
  $\tr[O_k\rho]=\tr[\hat O_k\Lambda(\rho)]$.
\item For every finite-dimensional side system $R$, every
  $\rho\in\Seq{\Bsys R}$, every $X\in\Lin{R}$, every $v$ and every $k$, writing
  $\sigma\coloneqq(\Lambda\otimes\idmap_R)(\rho)$,
  \begin{equation*}
    \tr\big[(M_v\otimes X)\rho\big]=\tr\big[(F_v\otimes X)\sigma\big]
    \qquad\text{and}\qquad
    \tr\big[(O_k\otimes X)\rho\big]=\tr\big[(\hat O_k\otimes X)\sigma\big].
  \end{equation*}
\end{enumerate}
Moreover, two channels $\Lambda,\Lambda'$ that both realize the relation
satisfy $\Lambda^{\dagger}=\Lambda'^{\dagger}$ on
$\operatorname{span}\big(\{F_v\}_{v\in\Vset}\cup\{\hat O_k\}_{k=1}^{m}\big)$.
\end{lemma}
\claimmeta{\metalabel{Uses}\mref{Lem.}{lem:sep}, \mref{Def.}{def:squashing}
\metasep \metalabel{Used by}\mref{Cor.}{cor:pullback}, \mref{Lem.}{lem:msmfact}, \mref{Rem.}{rem:completeness}, \mref{Rem.}{rem:unsound}, \mref{Rem.}{rem:channel-choice}
\metasep newly derived
\metasep accepted}

\begin{proof}
(i)$\Rightarrow$(iii). By the tensor-adjoint identity of
Lemma~\ref{lem:sep}\,(ii),
$(\Lambda\otimes\idmap_R)^{\dagger}(Y\otimes X)=\Lambda^{\dagger}(Y)\otimes X$.
Taking $Y=F_v$ and $Y=\hat O_k$ and inserting \eqref{eq:squashing} turns the
two right-hand sides of (iii) into the two left-hand sides, by
\eqref{eq:pairing} for $\Lambda\otimes\idmap_R$.

(iii)$\Rightarrow$(ii). Take $\Hs{R}=\mathbb{C}$ and $X=1$.

(ii)$\Rightarrow$(i). By \eqref{eq:pairing}, the first identity in (ii) reads
$\tr\big[\big(\Lambda^{\dagger}(F_v)-M_v\big)\rho\big]=0$ for every
$\rho\in\Seq{\Bsys}$, so $\Lambda^{\dagger}(F_v)=M_v$ by the separation
property, Lemma~\ref{lem:sep}\,(i). The same argument applied to the second
identity gives $\Lambda^{\dagger}(\hat O_k)=O_k$.

For the last claim, $\Lambda^{\dagger}$ and $\Lambda'^{\dagger}$ take the same
prescribed values on the operators $F_v$ and $\hat O_k$, hence agree on their
span by linearity.
\end{proof}

% claim reductions/rem:dictionary
\claimsummary{reductions/rem:dictionary}{The physical POVM stands on the left of the squashing symbol, as in the squashing-model literature, while the measurement-theory literature writes the same relation the other way round. The direction is to be checked before a statement is imported from there.}
\claimbookmark{reductions/rem:dictionary}{Remark}{rem:dictionary}{Direction of the squashing symbol in the literature}
\claimleadin{%
The side system $R$ in Lemma~\ref{lem:statistics}\,(iii) is not a
technicality: in a security proof it holds Alice's retained register and Eve's
side information, both of which are present whenever a squashed constraint is
evaluated.
}
\begin{remark}[Direction of the squashing symbol in the literature]\label{rem:dictionary}
In Definition~\ref{def:squashing} the \emph{physical} POVM stands on the left
and the target on the right, in the direction of
$\Lambda\colon\Tcl{\Bsys}\to\Tcl{\Qsys}$ and of the QKD usage in which
$\Lambda$ maps physical states to squashed states. The identity
\eqref{eq:squashing} is the one the squashing-model literature writes: it is
\cite[Eqs.~(1)--(2)]{beaudry2008squashing}, which states the statistical
equivalence $\tr[\rho_{\mathrm{in}}F_M(i)]
=\tr[\Lambda(\rho_{\mathrm{in}})F_Q(i)]$ and observes that it holds for all
inputs \emph{if and only if} $F_M(i)=\Lambda^{\dagger}(F_Q(i))$, with
$\Lambda^{\dagger}$ completely positive and unital but not necessarily trace
preserving;
and it is \cite[Eq.~(146)]{tupkary2026decoy} in measurement-channel form.
The measurement-theory literature writes the same relation the \emph{other way
round}: \cite[Def.~6.1]{buscemi2005clean} defines $P\succ Q$, ``$P$ is cleaner
than $Q$'', to mean $Q=\mathcal{E}(P)$ for a channel $\mathcal{E}$. The
paper's channels acting on effects are its identity-preserving completely
positive maps, so $\mathcal{E}$ is what this contract writes
$\Lambda^{\dagger}$ for a trace-preserving $\Lambda$, and the left-hand
argument $P$ is our \emph{target}. The direction is worth checking
before importing a statement from there; this contract imports none.
\end{remark}
\claimmeta{\metalabel{Uses}\mref{Def.}{def:squashing}
\metasep newly derived
\metasep accepted}

% claim reductions/rem:carried-operators
\claimsummary{reductions/rem:carried-operators}{The pull-back identity of the POVM alone does not confine a squashed analysis to the image of the physical inputs. Operators with a prescribed pull-back carry device facts across, which is why the squashing relation carries a tuple of them from the start.}
\claimbookmark{reductions/rem:carried-operators}{Remark}{rem:carried-operators}{Carried operators and the image of the squashing channel}
\begin{remark}[Carried operators and the image of the squashing channel]\label{rem:carried-operators}
The relation of Definition~\ref{def:squashing} carries more than the POVM.
Once $\Lambda$ is handed to the adversary, that is, counted as part of her
attack, nothing in the identity
$\Lambda^{\dagger}(F_v)=M_v$ confines the squashed analysis to states a
physical input can produce, and optimizing over all of $\Seq{\Qsys}$ instead
of over the image is where a squashed analysis loses. What recovers the loss
is that the identity keeps its meaning for operators that are \emph{not} POVM
elements: if $\Lambda^{\dagger}(\hat O)=O$, the expectation of $\hat O$ on a
squashed state equals that of $O$ on its physical preimage, so a device fact
about $O$ transfers to a statement about $\hat O$. The relation therefore
carries a designated tuple of such operators, its carried operators, from the
start.
\end{remark}
\claimmeta{\metalabel{Uses}\mref{Def.}{def:squashing}
\metasep newly derived
\metasep accepted}

% claim reductions/cor:pullback
\claimsummary{reductions/cor:pullback}{A carried operator's partial expectation on a squashed state equals its pull-back's on the physical input, as an operator on the side system. An identity, positivity or spectral bound on the pull-back therefore becomes an operator constraint on every state the squashing channel produces.}
\claimbookmark{reductions/cor:pullback}{Corollary}{cor:pullback}{Pull-back of operator bounds}
\claimleadin{%
The following corollary is the recipe every constraint of this contract is
built by: pick an operator $\hat O$ on $\Qsys$, compute its pull-back, and
read off what the physical device already guarantees about it.
}
\begin{corollary}[Pull-back of operator bounds]\label{cor:pullback}
Let $(M;O)\dsq_{\Lambda}(F;\hat O)$, let $R$ be a finite-dimensional side
system, let $\rho\in\Seq{\Bsys R}$ and let
$\sigma=(\Lambda\otimes\idmap_R)(\rho)$. Then for every $k$,
\begin{equation}\label{eq:transport}
  \tr_{\Qsys}\big[(\hat O_k\otimes\idon{R})\,\sigma\big]
  =\tr_{\Bsys}\big[(O_k\otimes\idon{R})\,\rho\big]
\end{equation}
as operators on $R$. Consequently, on every such $\sigma$:
\begin{enumerate}[label=(\roman*),leftmargin=3.0em]
\item \emph{Exact values.} If $O_k=c\,\idon{\Bsys}$ then
  $\tr_{\Qsys}[(\hat O_k\otimes\idon{R})\sigma]=c\,\sigma_R$; in particular
  $\tr[\hat O_k\sigma]=c$, and at $c=0$,
  $\tr[(\hat O_k\otimes X)\sigma]=0$ for every $X\in\Lin{R}$.
\item \emph{Positivity.} If $O_k\succeq0$ then
  $\tr_{\Qsys}[(\hat O_k\otimes\idon{R})\sigma]\succeq0$; in particular
  $\tr[\hat O_k\sigma]\ge0$.
\item \emph{Spectral bounds.} If
  $a\,\idon{\Bsys}\preceq O_k\preceq b\,\idon{\Bsys}$ then
  $a\,\sigma_R\preceq\tr_{\Qsys}[(\hat O_k\otimes\idon{R})\sigma]
  \preceq b\,\sigma_R$; in particular
  $a\le\tr[\hat O_k\sigma]\le b$.
\end{enumerate}
\end{corollary}
\claimmeta{\metalabel{Uses}\mref{Lem.}{lem:statistics}
\metasep \metalabel{Used by}\mref{Def.}{def:certificate}, \mref{Cor.}{cor:transported}, \mref{Rem.}{rem:unsound}, \mref{Lem.}{lem:master}, \mref{Thm.}{thm:mixture}, \mref{Thm.}{thm:licence}, \mref{Def.}{def:link}, \mref{Def.}{def:reduction-instance}, \mref{Thm.}{thm:instrument-instance}, \mref{Cor.}{cor:entries}, \mref{Req.}{req:transported-set-certificates}, \mref{Prop.}{prop:chain}, \mref{Rem.}{rem:published}, \mref{Cor.}{cor:dom-chain}
\metasep newly derived
\metasep accepted}

\begin{proof}
Test \eqref{eq:transport} against an arbitrary $X\in\Lin{R}$. By the defining
property of the partial trace and the side-system form of the statistics
equivalence, Lemma~\ref{lem:statistics}\,(iii),
\begin{equation*}
  \tr_R\!\Big\{X\,\tr_{\Qsys}\big[(\hat O_k\otimes\idon{R})\sigma\big]\Big\}
  =\tr\big[(\hat O_k\otimes X)\sigma\big]
  =\tr\big[(O_k\otimes X)\rho\big]
  =\tr_R\!\Big\{X\,\tr_{\Bsys}\big[(O_k\otimes\idon{R})\rho\big]\Big\},
\end{equation*}
and since $R$ is finite-dimensional the trace pairing on $\Lin{R}$ separates
points, so the two operators on $R$ coincide.

(i) With $O_k=c\,\idon{\Bsys}$ the right-hand side of \eqref{eq:transport} is
$c\,\tr_{\Bsys}\rho=c\,\rho_R$, and $\rho_R=\sigma_R$ because $\Lambda$ is
trace preserving and acts only on $\Bsys$. (ii) For $X\succeq0$ we have
$O_k\otimes X\succeq0$, so $\tr[(O_k\otimes X)\rho]\ge0$; hence
$\tr_R\{X\,\tr_{\Qsys}[(\hat O_k\otimes\idon{R})\sigma]\}\ge0$ for every
$X\succeq0$, which for a Hermitian operator on a finite-dimensional space is
positive semidefiniteness. (iii) Apply (ii) to $O_k-a\,\idon{\Bsys}\succeq0$
and to $b\,\idon{\Bsys}-O_k\succeq0$, then use (i) on the identity terms; the
scalar forms follow by taking traces and using $\tr\sigma_R=1$.
\end{proof}

% claim reductions/def:certificate
\claimsummary{reductions/def:certificate}{The certificate of a transported constraint is the hypothesis that makes it hold: for a case of \hyperlink{claim-reductions/cor:pullback}{Corollary~\ref*{cor:pullback}} an operator fact about the device and the channel, valid whatever the input, and for the assumed case of \hyperlink{claim-reductions/cor:transported}{Corollary~\ref*{cor:transported}} an assumption on the input states.}
\claimbookmark{reductions/def:certificate}{Definition}{def:certificate}{Certificates of transported constraints}
\claimleadin{%
A constraint is a condition on squashed states that an analysis imposes.
Corollary~\ref{cor:pullback} derives one from an identity, a positivity or a
spectral bound on a carried operator's pull-back, and
Corollary~\ref{cor:transported}, in the section on
output-state equality, reads the same derivation over $n$ rounds and also
admits a bound assumed to hold on the \termref{term-physical-inputs}{physical inputs}.
}
\begin{definition}[Certificates of transported constraints]\label{def:certificate}
The
\termdef{term-certificate}{certificate}
of a transported constraint is the hypothesis that makes it hold: for a
constraint from Corollary~\ref{cor:pullback}, the hypothesis of the case that
yields it, an operator identity in its case~(i) and an operator inequality in
its cases (ii)--(iii), both facts about the device and the channel, valid
whatever the input; and, for an assumed constraint of
Corollary~\ref{cor:transported}\,(ii), an assumption on the input states.
\end{definition}
\claimmeta{\metalabel{Uses}\mref{Cor.}{cor:pullback}
\metasep \metalabel{Used by}\mref{Cor.}{cor:transported}, \mref{Def.}{def:link}, \mref{Def.}{def:reduction-instance}, \mref{Lem.}{lem:row-relaxation}
\metasep newly derived
\metasep accepted}

% claim reductions/rem:notchannel
\claimsummary{reductions/rem:notchannel}{The universal squash model establishes a statistical equivalence and cutoff-based dimension reduction controls the effect of a cutoff, neither by exhibiting a squashing channel. Without such a channel the output-state equality behind \hyperlink{claim-reductions/thm:transfer}{Theorem~\ref*{thm:transfer}} is not available.}
\claimbookmark{reductions/rem:notchannel}{Remark}{rem:notchannel}{Reductions that are not squashings}
\begin{remark}[Reductions that are not squashings]\label{rem:notchannel}
Not every dimension reduction used in QKD comes with a squashing channel, and
one that does not cannot be fed into the security transfer,
Theorem~\ref{thm:transfer}. The universal squash model of
\cite[Sec.~II]{fung2011universal} establishes a statistical
equivalence between the threshold-detector setup and an ideal squashed one by
routing both through classical post-processing of a photon-number-resolving
measurement, and the dimension-reduction method of
\cite[Sec.~III]{upadhyaya2021dimension} controls the effect of a finite cutoff
rather than exhibiting a channel. Neither proceeds through the existence of a
$\Lambda$ with $\Lambda^{\dagger}(F_v)=M_v$, which is exactly the hypothesis the
output-state equality of Theorem~\ref{thm:factorization} runs on and the
hypothesis \cite[Lemma~6]{nahar2024postselection} states for a narrower model.
The rule that follows for the emitted artefact, a \termref{term-reduction-chain}{reduction chain} of
links, is Requirement~\ref{req:channel-witness}: an explicit
\termref{term-channel-witness}{channel witness} for every link. The companion rule for constraint
entries, exactly one named certificate each, is
Requirement~\ref{req:certified-entries-only}.
\end{remark}
\claimmeta{\metalabel{Used by}\mref{Req.}{req:channel-witness}
\metasep newly derived
\metasep accepted}

% claim reductions/def:msm
\claimsummary{reductions/def:msm}{The measurement channel of a POVM takes a quantum input to the classical register holding its outcome. It separates Bob's measurement from the protocol's later processing.}
\claimbookmark{reductions/def:msm}{Definition}{def:msm}{Measurement channel}
\claimleadin{%
Replacing Bob's physical POVM $M$ by a squashed target $F$ can only
overestimate the adversary, for a short reason. First, $M\dsq_{\Lambda}F$ says
exactly that Bob's physical measurement \emph{channel} is the squashing channel
followed by the squashed one. Second, Bob's device enters the protocol only
through that channel. Third, therefore, the physical protocol and the
\termref{term-squashed-protocol}{squashed protocol} have the same output state, not merely close ones,
which is why the transfer
costs no penalty term, no correction factor, and the same $\varepsilon$ on both
sides. What remains to check is that the squashed analysis runs over a set of
inputs containing the images of the physical ones, and that is
Theorem~\ref{thm:transfer}.
}
\begin{definition}[Measurement channel]\label{def:msm}
For a POVM $M=\{M_v\}_{v\in\Vset}$ on $\Bsys$ and a classical register $Y$ with
alphabet $\Vset$, the
\termdef{term-measurement-channel}{measurement channel}
of $M$ is
\begin{equation}\label{eq:msm}
  \redNmsm_M\colon\Tcl{\Bsys}\to\Tcl{Y},\qquad
  \redNmsm_M(\rho)\coloneqq\sum_{v\in\Vset}\tr\big[M_v\rho\big]\,\dyad{v}_{Y}.
\end{equation}
\end{definition}
\claimmeta{\metalabel{Used by}\mref{Def.}{def:protocol}, \mref{Prop.}{prop:platform}
\metasep newly derived
\metasep accepted}

% claim reductions/lem:msmfact
\claimsummary{reductions/lem:msmfact}{A POVM squashes to a target via a channel exactly when its measurement channel is the target's measurement channel after that channel. This is the channel identity the protocol-level argument inserts round by round.}
\claimbookmark{reductions/lem:msmfact}{Lemma}{lem:msmfact}{Measurement-channel factorization}
\begin{lemma}[Measurement-channel factorization]\label{lem:msmfact}
Let $M,F$ be POVMs on $\Bsys,\Qsys$ over $\Vset$ and let
$\Lambda\colon\Tcl{\Bsys}\to\Tcl{\Qsys}$ be a channel. Then
\begin{equation}\label{eq:msmfact}
  M\dsq_{\Lambda}F
  \qquad\Longleftrightarrow\qquad
  \redNmsm_M=\redNmsm_F\circ\Lambda .
\end{equation}
\end{lemma}
\claimmeta{\metalabel{Uses}\mref{Lem.}{lem:sep}, \mref{Lem.}{lem:statistics}
\metasep \metalabel{Used by}\mref{Thm.}{thm:factorization}, \mref{Cor.}{cor:sequential}
\metasep newly derived
\metasep accepted}

\begin{proof}
($\Rightarrow$) For every $\rho\in\Tcl{\Bsys}$,
$\redNmsm_M(\rho)=\sum_v\tr[\Lambda^{\dagger}(F_v)\rho]\dyad{v}
=\sum_v\tr[F_v\Lambda(\rho)]\dyad{v}=\redNmsm_F(\Lambda(\rho))$ by
\eqref{eq:pairing}. ($\Leftarrow$) Comparing the coefficient of $\dyad{v}$ on
both sides gives $\tr[M_v\rho]=\tr[F_v\Lambda(\rho)]$ for every state $\rho$,
which is condition (ii) of the statistics equivalence,
Lemma~\ref{lem:statistics}, with $m=0$.
\end{proof}

\section{Protocols, security criteria and exact operational reductions}\label{group-reductions/def:protocol}
\sectionintro{The protocols of this contract pass each round's \termref{term-measured-system}{measured system} through one \termref{term-measurement-channel}{measurement channel} and process only its classical outcome afterwards. A \termref{term-security-criterion}{security criterion} is any family of functions of the protocol's final output together with the adversary's system, and an \termref{term-exact-operational-reduction}{exact operational reduction} is a map of attack strategies under which every strategy and its image produce the same \termref{term-full-output}{full output}. Such a map transfers every bound on every criterion with the same \ensuremath{\varepsilon}, and reductions restrict, enlarge and compose (\hyperlink{claim-reductions/lem:opreduction-transfer}{Lemma~\ref*{lem:opreduction-transfer}}). The section closes with the reason both reductions of this contract are exact rather than approximate.}

% claim reductions/def:protocol
\claimsummary{reductions/def:protocol}{An n-round protocol passes each round's measured system, held by the measuring parties, through one measurement channel, the same in every round, and no later processing touches that system. The squashed protocol is the same channel with the target POVM in place of the physical one.}
\claimbookmark{reductions/def:protocol}{Definition}{def:protocol}{Measurement-factorizing n-round protocol}
\claimleadin{%
The protocols this contract reduces touch each round's measured system only
through a measurement; what follows reads its classical outcome and never the
system itself. The definition fixes that structure. The declaration it speaks
of is the protocol declaration that defines the consuming engine's protocol
(\xref{protocol-compiler/def:declaration}); it names the parties that measure. The
structure rests on a round order, every public message of a round sent after
every measurement of that round, which is Assumption A-SEQ of the
transmitter's contract. This definition places the
\termref{term-measurement-channel}{measurement channel} before the round's remaining processing and fixes
one POVM and one channel for all rounds; the claims that rest on that
structure declare A-SEQ and A-ROUND-SQ (Assumption~\ref{ass:round}), a
definition assuming nothing. Assumption~\ref{ass:meas} and
Assumption~\ref{ass:round} leave the round order out. The embedding by which
a declaration names Alice as a measuring party is that of
\xref{meat-core/rem:modes-da1}\,(b).
}
\begin{definition}[Measurement-factorizing $n$-round protocol]
\label{def:protocol}
An
\termdef{term-n-round-protocol}{$n$-round protocol}
with receiver POVM $M$ is a channel
\begin{equation}\label{eq:protocol}
  \Eprot_M\;=\;\Kfin\circ\Grnd_n\circ\dots\circ\Grnd_1,
  \qquad
  \Grnd_j\;=\;\bar{\Grnd}_j\circ\big(\idmap\otimes\redNmsm_M\big),
\end{equation}
in which the \termref{term-measurement-channel}{measurement channel} $\redNmsm_M$ of
Definition~\ref{def:msm} acts on the round-$j$
\termdef{term-measured-system}{measured system}
$\Bsys_j$, the quantum system that the
\termdef{term-measuring-parties}{measuring parties},
the parties the declaration names as measuring, receive in round $j$, and no
other system, and the identity acts on everything else. $\bar{\Grnd}_j$ is the
round's remaining processing (Alice's measurement on $A_j$, the round's
announcements appended to the public record, and the retained per-round data),
and $\Kfin$ is the end-of-protocol processing: sifting, parameter estimation
and the accept/abort decision, error correction, verification, privacy
amplification and the key maps. The
\termdef{term-squashed-protocol}{squashed protocol}
$\Eprot_F$ is \eqref{eq:protocol} with $\Bsys_j$ replaced by $\Qsys_j$, $M$ by
$F$, and $\bar{\Grnd}_j$, $\Kfin$ unchanged. One POVM and one channel serve
all $n$ rounds: the same $M$, and in the squashed protocol the same $F$ and
the same \termref{term-squashing-channel}{squashing channel}, stands in every $\Grnd_j$, with no
dependence on $j$. That the measuring devices measure by the same $M$ in
every round, with no memory between rounds, is Assumption~\ref{ass:round};
$F$ and the squashing channel are objects of the analysis, chosen once for all
rounds. That is what makes the composite of the per-round channels a tensor
power, such as $\Lambda^{\otimes n}$ for the squashing channel, rather than a
product of $n$ different maps.

\emph{One measuring party, or several.} $M$ is one POVM on the measured
system, over the alphabet of the measuring parties' joint outcome. Where one
party measures, which is the one-party case, $\Bsys_j$ is that party's
system, Bob's in a prepare-and-measure declaration. Where several parties
measure, $\Bsys_j=\bigotimes_f\Bsys_{f,j}$ is the product of their
party-scoped parts, each factor $\Bsys_{f,j}$ belonging to one of them, and
\eqref{eq:protocol} places \emph{all} of their measurements at the one step
$\redNmsm_M$, before the round's remaining processing $\bar{\Grnd}_j$ and so
before every public message of the round.
The declaration says on which side Alice's measurement sits, and each reading
is covered by exactly one clause of the first paragraph: the measurement on
$A_j$ listed in $\bar{\Grnd}_j$ is that of a party the declaration does
\emph{not} name as measuring (Alice's, in a prepare-and-measure declaration),
while where she \emph{is} named as measuring, her system is a factor
$\Bsys_{f,j}$ of the measured system and her measurement is a factor of
$\redNmsm_M$, not a step of $\bar{\Grnd}_j$. Either way $\bar{\Grnd}_j$ and
$\Kfin$ touch no $\Bsys_j$, which is all the theorems of this contract read.
A declaration realizes the second reading by an embedding: it declares
Alice's system $A$
one-dimensional, $d_A=1$, folds her measured system into the sectors of
$\Qsys$ and her outcome components into $v$, and names the components she
holds in its party partition (\xref{protocol-compiler/def:declaration}). That the
physical protocol has that order, every public message of round $j$ sent after
every measurement of round $j$ whichever party sends it and whichever party
measured, is Assumption~A-SEQ of the transmitter
(\xref{transmitter/ass:seq}). The definition says nothing about which party
sends which message or about a message's dependence on earlier messages of
its own round, and a quantum operation of round $j$ that depended on a
message of round $j$ is outside it, as it is outside A-SEQ.
\end{definition}
\claimmeta{\metalabel{Uses}\mref{Def.}{def:msm}
\metasep \metalabel{Used by}\mref{Thm.}{thm:factorization}, \mref{Def.}{def:roundmodel}, \mref{Cor.}{cor:sequential}, \mref{Prop.}{prop:platform}, \mref{Thm.}{thm:transfer}, \mref{Lem.}{lem:source-transfer}, \mref{Ass.}{ass:meas}, \mref{Ass.}{ass:round}, \extref{../meat-core/contract.pdf}{meat-core/lem:honest-attacks}, \extref{../meat-core/contract.pdf}{meat-core/lem:memory-lift}
\metasep newly derived
\metasep accepted}

% claim reductions/def:criterion
\claimsummary{reductions/def:criterion}{A security criterion is any family of functions, one per admissible adversary system, of the joint state of the protocol's output registers and that system. Equality of those states therefore transfers every criterion at once, with no proof for its particular formula.}
\claimbookmark{reductions/def:criterion}{Definition}{def:criterion}{Security criterion}
\claimleadin{%
The protocol $\Eprot_M$ of Definition~\ref{def:protocol} on which a criterion
is evaluated factors in each round as
$\Grnd_j=\bar{\Grnd}_j\circ(\idmap\otimes\redNmsm_M)$, and that factorization
rests on Assumption~\ref{ass:meas}, together with the round order and the one
POVM for all rounds that Definition~\ref{def:protocol} declares: each
measuring party's interaction with the quantum system it receives is
described by a single POVM, and $M$ is their product on the round's measured
system, Bob's POVM where he alone measures. The measurement emits the
classical outcome $v$ and nothing else; no post-measurement quantum system is
retained or reused, and all later processing reads the measured system only
through that outcome.
}
\begin{definition}[Security criterion]\label{def:criterion}
A
\termdef{term-security-criterion}{security criterion}
is a family $f=(f_E)_E$ of functions, one for every admissible adversary
system $E$, defined on states of the protocol's output registers together with
$E$; an \emph{admissible adversary system} is a finite-dimensional system,
the form the \exttermref{../foundations/contract.pdf}{term-round-model}{round model} of \xref{foundations/def:round-model} gives
the adversary's registers. Its value
on the protocol $\Eprot_M$ at an input $\rho$, a state of the systems the
protocol acts on together with $E$, is
\begin{equation}\label{eq:criterion}
  \Crit(\Eprot_M,\rho)\coloneqq
  f_E\big((\Eprot_M\otimes\idmap_E)(\rho)\big).
\end{equation}
\end{definition}
\claimmeta{\metalabel{Used by}\mref{Lem.}{lem:opreduction-transfer}, \mref{Rem.}{rem:why-exact}, \mref{Thm.}{thm:factorization}, \mref{Def.}{def:roundmodel}, \mref{Rem.}{rem:criteria}, \mref{Rem.}{rem:completeness}, \mref{Thm.}{thm:transfer}, \mref{Cor.}{cor:sound}, \mref{Thm.}{thm:two-sided}
\metasep newly derived
\metasep accepted}

% claim reductions/def:opreduction
\claimsummary{reductions/def:opreduction}{A protocol presentation is a set of attack strategies with the full output each produces; an exact operational reduction maps strategies between two presentations over the same registers so that every strategy and its image give the same output. Bounds carry back with no property of the criterion.}
\claimbookmark{reductions/def:opreduction}{Definition}{def:opreduction}{Exact operational reduction}
\claimleadin{%
The two reductions this contract performs, replacing a physical receiver by a
finite model and replacing a physical source by a virtual one, are the same
kind of object, and everything either of them transfers is transferred by that
kind of object alone. This definition states the object once and
Lemma~\ref{lem:opreduction-transfer} states what it buys;
Corollary~\ref{cor:sequential} and Theorem~\ref{thm:transfer} are the
receiver-side instance, Lemma~\ref{lem:source-transfer} the source-side one,
and Theorem~\ref{thm:two-sided} their composition. Nothing in the definition
is about squashing, about measurement, or about any particular protocol: it
speaks of a set of strategies, of the state each of them produces, and of
nothing else.
}
\begin{definition}[Exact operational reduction]\label{def:opreduction}
A
\termdef{term-protocol-presentation}{protocol presentation}
is a triple $\mathsf{P}=(\mathbb{A},W,\mathsf{out})$: a set $\mathbb{A}$, each
element of which is called an
\termdef{term-attack-strategy}{attack strategy},
a tuple $W$ of
\termdef{term-output-registers}{output registers},
and a map $\mathsf{out}\colon\mathbb{A}\to\Sleq{W}$ assigning to each
strategy the
\termdef{term-full-output}{full output}
the protocol produces under it. The full output is the state on \emph{every}
register a security criterion may read (Definition~\ref{def:criterion}): both
parties' key registers, the public transcript with the acceptance and abort
data, and the adversary's final register, and no marginal of them.

Let $\mathsf{P}=(\mathbb{A},W,\mathsf{out})$ and
$\mathsf{P}'=(\mathbb{A}',W,\mathsf{out}')$ be presentations \emph{over the
same output registers} $W$. An
\termdef{term-exact-operational-reduction}{exact operational reduction}
from $\mathsf{P}$ to $\mathsf{P}'$ is a map $r\colon\mathbb{A}\to\mathbb{A}'$
with
\begin{equation}\label{eq:opreduction}
  \mathsf{out}(a)\;=\;\mathsf{out}'\big(r(a)\big)
  \qquad\text{for every }a\in\mathbb{A} .
\end{equation}
We call $\mathbb{A}$ the
\termdef{term-physical-strategy-set}{physical strategy set}
and $\mathbb{A}'$ the
\termdef{term-analysed-strategy-set}{analysed strategy set},
equivalently the \emph{admissible} strategy set, and
$r(\mathbb{A})\subseteq\mathbb{A}'$ the reduction's
\termdef{term-image}{image}.

\emph{The conditions, in full.}
\begin{enumerate}[label=(R\arabic*),leftmargin=3.0em]
\item \emph{(one output tuple)} the two presentations carry literally the same
  output registers, and which register of $\mathsf{P}'$ is which register of
  $\mathsf{P}$ is part of the data rather than inferred. Where a model
  protocol's adversary register is a \emph{different} system, an
  identification of it with the physical one has to be declared and the
  criterion has to be invariant under that identification; this contract
  declares none and needs none, the adversary's registers being literally the
  same system in both of its instances;
\item \emph{(exact equality)} \eqref{eq:opreduction} is equality of the two
  states on the nose: not of a marginal, not up to a distance, and not up to
  any map on $W$;
\item \emph{(totality and typing)} $r$ is defined at every strategy of
  $\mathbb{A}$ and takes its values in $\mathbb{A}'$. That its values lie in
  $\mathbb{A}'$ is a condition on $r$ and on $\mathbb{A}'$ together, and it is
  the only place where a restriction of the analysed set has to be earned.
\end{enumerate}

\emph{What is not required, each item being used by an instance in this
contract.} $r$ need not be injective, need not be surjective onto
$\mathbb{A}'$, and need not preserve any structure of $\mathbb{A}$; the
analysed set $\mathbb{A}'$ may be, and in this contract typically is, strictly
larger than the image, because what a numerical analysis can optimize over is
a set cut out by finitely many conditions and not an image
(Definition~\ref{def:attset}); no metric, topology, order, linear or convex
structure on $\mathbb{A}$, on $\mathbb{A}'$ or on the output states is used;
$\mathsf{out}$ and $\mathsf{out}'$ need not be linear, affine, continuous or
given by channels, and a strategy need not be a channel or a tuple of
channels (in Theorem~\ref{thm:transfer} a strategy is a state); and nothing
is asserted in the other direction, a bound over $\mathbb{A}$ transferring to
$\mathbb{A}'$ under no hypothesis of this definition.

\emph{Two presentations, one per adversary system.} A security criterion in
the sense of Definition~\ref{def:criterion} is a \emph{family} $(f_E)_E$, one
function per admissible adversary system. Fixing $E$ fixes $W$ and turns each
$f_E$ into one function on $\Sleq{W}$, so the quantification over $E$ in that
definition is a quantification over presentations: a statement of this
contract that holds ``for every admissible $E$'' exhibits one presentation
pair and one reduction per $E$, and the transfer of
Lemma~\ref{lem:opreduction-transfer} is applied to each.
\end{definition}
\claimmeta{\metalabel{Used by}\mref{Lem.}{lem:opreduction-transfer}, \mref{Rem.}{rem:why-exact}, \mref{Cor.}{cor:sequential}, \mref{Thm.}{thm:transfer}, \mref{Cor.}{cor:sound}, \mref{Lem.}{lem:source-transfer}, \mref{Thm.}{thm:two-sided}
\metasep newly derived
\metasep accepted}

% claim reductions/lem:opreduction-transfer
\claimsummary{reductions/lem:opreduction-transfer}{An exact operational reduction carries every bound on every function of the full output from the analysed to the physical strategy set with the same \ensuremath{\varepsilon}, and reductions restrict, enlarge and compose. A bound over an intersection of analysed sets needs one map into all of them.}
\claimbookmark{reductions/lem:opreduction-transfer}{Lemma}{lem:opreduction-transfer}{Transfer and composition of exact operational reductions}
\claimleadin{%
The lemma is elementary set arithmetic. Its part (iv) matters in this
contract because two channels realising one and the same
\termref{term-squashing-relation}{squashing relation}
give two different strategy maps, whose reachable sets are in general
incomparable; that is why Corollary~\ref{cor:sound}\,(ii)--(iii) are stated for
one common channel and one common family of \termref{term-physical-inputs}{physical inputs}
(Remark~\ref{rem:channel-choice}). Part (vi) says that an analysed set larger
than the image costs tightness and never soundness, the same trade the
reachable sets of Definition~\ref{def:sound} make.
}
\begin{lemma}[Transfer and composition of exact operational reductions]\label{lem:opreduction-transfer}
Let $r$ be an exact operational reduction from
$\mathsf{P}=(\mathbb{A},W,\mathsf{out})$ to
$\mathsf{P}'=(\mathbb{A}',W,\mathsf{out}')$ in the sense of
Definition~\ref{def:opreduction}. Write $\mathcal{F}(W)$ for the set of all
functions $\Sleq{W}\to[-\infty,+\infty]$; a security criterion in the sense of
Definition~\ref{def:criterion}, at a fixed adversary system, is one such
function.
\begin{enumerate}[label=(\roman*),leftmargin=3.0em]
\item \emph{(transfer, with the same $\varepsilon$ and for every criterion at
  once)} for every $f\in\mathcal{F}(W)$ and every $a\in\mathbb{A}$,
  $f(\mathsf{out}(a))=f(\mathsf{out}'(r(a)))$, and consequently
  \begin{equation}\label{eq:opsup}
    \sup_{a\in\mathbb{A}}f\big(\mathsf{out}(a)\big)
    \;=\;\sup_{b\in r(\mathbb{A})}f\big(\mathsf{out}'(b)\big)
    \;\le\;\sup_{b\in\mathbb{A}'}f\big(\mathsf{out}'(b)\big) .
  \end{equation}
  In particular, if $f(\mathsf{out}'(b))\le\varepsilon$ for every
  $b\in\mathbb{A}'$, then $f(\mathsf{out}(a))\le\varepsilon$ for every
  $a\in\mathbb{A}$, with the \emph{same} $\varepsilon$, and simultaneously
  for every $f\in\mathcal{F}(W)$, no property of $f$ being used;
\item \emph{(enlargement of the analysed set)} if
  $\mathbb{A}'\subseteq\mathbb{A}''$ and
  $\mathsf{out}''|_{\mathbb{A}'}=\mathsf{out}'$, then the same $r$ is an
  exact operational reduction from $\mathsf{P}$ to
  $\mathsf{P}''=(\mathbb{A}'',W,\mathsf{out}'')$. So it is enough that the
  image $r(\mathbb{A})$ lie in the analysed set, and the last inequality of
  \eqref{eq:opsup} is where that slack is spent;
\item \emph{(restriction of the physical set)} if
  $\mathbb{A}_0\subseteq\mathbb{A}$, then $r|_{\mathbb{A}_0}$ is an exact
  operational reduction from $(\mathbb{A}_0,W,\mathsf{out}|_{\mathbb{A}_0})$ to
  $\mathsf{P}'$;
\item \emph{(one map, many analysed sets)} let $(\mathbb{A}'_k)_{k\in K}$ be a
  family of subsets of $\mathbb{A}'$ such that \emph{one and the same} $r$
  satisfies $r(\mathbb{A})\subseteq\mathbb{A}'_k$ for every $k\in K$. Then $r$
  is an exact operational reduction from $\mathsf{P}$ to
  $\big(\bigcap_{k\in K}\mathbb{A}'_k,W,\mathsf{out}'\big)$, and by (i) a bound
  over that intersection transfers. The hypothesis cannot be weakened to the
  existence, for each $k$, of \emph{some} exact operational reduction into
  $\mathbb{A}'_k$: let $W$ be one classical bit, $\mathbb{A}=\{a\}$,
  $\mathbb{A}'=\{b_1,b_2\}$ with $b_1\ne b_2$ and
  $\mathsf{out}(a)=\mathsf{out}'(b_1)=\mathsf{out}'(b_2)=\dyad{0}$, and
  $\mathbb{A}'_k=\{b_k\}$. Each $r_k\colon a\mapsto b_k$ is an exact
  operational reduction into $\mathbb{A}'_k$, the intersection is empty, so
  ``$f(\mathsf{out}'(b))\le0$ for every $b$ in the intersection'' holds
  vacuously; yet $f(\mathsf{out}(a))=1$ for
  $f(\rho)\coloneqq\tr[\dyad{0}\rho]$, so the conclusion of (i) is false for
  the intersection;
\item \emph{(composition, and how the analysed sets compose)} let
  $\mathbb{A}'_0\subseteq\mathbb{A}'$ contain $r(\mathbb{A})$ and let $r'$ be an
  exact operational reduction from
  $\big(\mathbb{A}'_0,W,\mathsf{out}'|_{\mathbb{A}'_0}\big)$ to
  $\mathsf{P}''=(\mathbb{A}'',W,\mathsf{out}'')$. Then $r'\circ r$ is an exact
  operational reduction from $\mathsf{P}$ to $\mathsf{P}''$, with image
  $r'(r(\mathbb{A}))$. The analysed set the composite needs is the second
  reduction's, or any subset of $\mathbb{A}''$ containing
  $r'(r(\mathbb{A}))$; the first reduction's analysed set enters only as a
  domain the second must cover, which is what
  $r(\mathbb{A})\subseteq\mathbb{A}'_0$ says. Composition is associative and
  the identity map on $\mathbb{A}$ is an exact operational reduction from
  $\mathsf{P}$ to itself;
\item \emph{(the direction, and what the slack costs)} nothing above bounds
  $\sup_{\mathbb{A}'}f\circ\mathsf{out}'$ by
  $\sup_{\mathbb{A}}f\circ\mathsf{out}$, and the first can be strictly larger:
  take $\mathbb{A}=\{a\}$, $\mathbb{A}'=\{r(a),b\}$ with
  $\mathsf{out}'(b)\ne\mathsf{out}(a)$ and an $f$ separating the two states.
  An analysed set larger than the image therefore costs \emph{tightness} and
  never soundness.
\end{enumerate}
\end{lemma}
\claimmeta{\metalabel{Uses}\mref{Def.}{def:criterion}, \mref{Def.}{def:opreduction}
\metasep \metalabel{Assumes}\aref{A-MEAS}{ass:meas}
\metasep \metalabel{Used by}\mref{Rem.}{rem:why-exact}, \mref{Cor.}{cor:sequential}, \mref{Thm.}{thm:transfer}, \mref{Cor.}{cor:sound}, \mref{Lem.}{lem:source-transfer}, \mref{Thm.}{thm:two-sided}
\metasep newly derived
\metasep accepted}

\begin{proof}
(i) The pointwise identity is $f$ applied to both sides of
\eqref{eq:opreduction}. Hence
$\{\mathsf{out}(a):a\in\mathbb{A}\}
=\{\mathsf{out}'(r(a)):a\in\mathbb{A}\}
=\{\mathsf{out}'(b):b\in r(\mathbb{A})\}$
as subsets of $\Sleq{W}$, so the two suprema of $f$ over them agree, which is
the equality in \eqref{eq:opsup}; the inequality is
$r(\mathbb{A})\subseteq\mathbb{A}'$ and monotonicity of the supremum under
enlarging the index set. The last sentence is \eqref{eq:opsup} read with the
right-hand supremum bounded by $\varepsilon$, and $f$ was arbitrary.

(ii) The map is unchanged and \eqref{eq:opreduction} is unchanged, because
$\mathsf{out}''(r(a))=\mathsf{out}'(r(a))$ for every $a$ by
$r(\mathbb{A})\subseteq\mathbb{A}'$ and the hypothesis on $\mathsf{out}''$;
the typing condition (R3) of Definition~\ref{def:opreduction} holds because
$\mathbb{A}'\subseteq\mathbb{A}''$, and (R1), (R2) are conditions on $W$ and on
the values, which did not move.

(iii) \eqref{eq:opreduction} is a condition at each strategy separately, so it
survives restricting the set of strategies it is required at.

(iv) $r(\mathbb{A})\subseteq\bigcap_k\mathbb{A}'_k$ is exactly the hypothesis,
and \eqref{eq:opreduction} does not read the analysed set, so $r$ is an exact
operational reduction into the intersection; the transfer is (i). For the
counterexample: each $r_k$ satisfies \eqref{eq:opreduction} by construction and
$r_k(\mathbb{A})=\{b_k\}=\mathbb{A}'_k$, so each is an exact operational
reduction into $\mathbb{A}'_k$; $\mathbb{A}'_1\cap\mathbb{A}'_2=\varnothing$
because $b_1\ne b_2$; a supremum over the empty set is $-\infty$, so every
bound holds over it; and $f(\mathsf{out}(a))=\tr[\dyad{0}\dyad{0}]=1$. What
fails is not (i) but its hypothesis: no map has its image in the intersection.

(v) For every $a\in\mathbb{A}$, $r(a)\in\mathbb{A}'_0$ by hypothesis, so
$r'(r(a))$ is defined and lies in $\mathbb{A}''$, and
$\mathsf{out}(a)=\mathsf{out}'(r(a))=\mathsf{out}''(r'(r(a)))$ by
\eqref{eq:opreduction} for $r$ and then for $r'$. The image statement is the
definition of a composite's image. Associativity and the identity are those of
maps of sets; the identity satisfies \eqref{eq:opreduction} trivially.

(vi) The two suprema are unrelated by anything proved above, and the displayed
instance has $\sup_{\mathbb{A}'}f\circ\mathsf{out}'
=\max\{f(\mathsf{out}(a)),f(\mathsf{out}'(b))\}
>f(\mathsf{out}(a))=\sup_{\mathbb{A}}f\circ\mathsf{out}$ for an $f$ with
$f(\mathsf{out}'(b))>f(\mathsf{out}(a))$, which exists because the two states
differ.
\end{proof}

% claim reductions/rem:why-exact
\claimsummary{reductions/rem:why-exact}{Equality of full outputs lets one argument transfer every criterion with the same \ensuremath{\varepsilon}, while a distance bound would need a modulus of continuity no criterion is given. Every reduction and link of this contract, the domination link included, is exact.}
\claimbookmark{reductions/rem:why-exact}{Remark}{rem:why-exact}{Exactness of the reductions and generality of the transfer}
\begin{remark}[Exactness of the reductions and generality of the transfer]\label{rem:why-exact}
Definition~\ref{def:opreduction} demands equality of the \emph{full} output
state, and that demand is what pays for everything
Lemma~\ref{lem:opreduction-transfer} delivers. A security criterion in the
sense of Definition~\ref{def:criterion} is an arbitrary function of the output
state: no linearity, no monotonicity and no continuity is assumed of it
anywhere in this contract. Under exact equality that costs nothing, because
part (i) of the lemma is $f$ applied to two equal states; under anything
weaker it costs everything, because a hypothesis
$\|\mathsf{out}(a)-\mathsf{out}'(r(a))\|_1\le\delta$ bounds
$f(\mathsf{out}(a))$ in terms of $f(\mathsf{out}'(r(a)))$ only through a
modulus of continuity for $f$, which the definition of a criterion does not
give and which the consuming engine's criterion would have to supply case by
case. So the generality of the transfer, one argument for every criterion with
the same $\varepsilon$, and the exactness are one fact, not two.

\emph{Both instances are exact, for a checkable reason.} On the receiver side
the substitution rests on the channel identity \eqref{eq:factorization}, an
identity of maps with no error term (Theorem~\ref{thm:factorization}). On the
source side it rests on the source identity (S3) of
\xref{transmitter/def:source-segment}, an exact matrix identity on declared
numbers: it is decided entry by entry where every \exttermref{../transmitter/contract.pdf}{term-source-link}{source link} supplies
its channel as data, and concluded from a named claim where a link supplies
its channel through a \exttermref{../transmitter/contract.pdf}{term-theorem-constructor}{theorem constructor}
(Lemma~\ref{lem:source-transfer}).

\emph{What an approximate reduction would need.} It would need a theorem
converting a distance between output states into a security penalty, and this
contract has none. The security transfer, Theorem~\ref{thm:transfer}, rests on
an equality and has no penalty slot, so there is no parameter for such a term
to enter; accordingly a declared relation carrying an error term is refused
(Requirement~\ref{req:no-approximate-relations}), and a \termref{term-reduction-chain}{reduction chain}
mixing an exact link with one carrying an error term is refused rather than
given a composite status, a status for the chain as a whole
(Theorem~\ref{thm:chain}\,(iv), Requirement~\ref{req:mixed-status-chain}).
The nearest thing to a continuity statement in these contracts is a refusal of
one: \xref{meat-core/prop:optol-nobound} exhibits two
\exttermref{../foundations/contract.pdf}{term-exactly-admissible}{exactly admissible} target families within a declared distance of each
other whose rate bounds per round differ by an amount limited only by the
number model. The limit is that proposition's own representability condition,
which binary64 arithmetic imposes and which caps the integer $N$ measuring the
gap at $N\le52$; so proximity of models licenses nothing about rates there
either.

\emph{The \termref{term-emission-displacement}{emission displacement} is not an error term of a reduction.}
The displacement $\Delta_s$ of the \termref{term-emission-normalizer}{emission normalizer}, the largest
change it makes to an entry of an instance's datum in the \termref{term-data-sector}{data sector}
$s$, does not make a reduction approximate. It moves the \termref{term-declared-model}{declared model}
the reduction is stated about, and the reduction's identity then holds of that
model with no error term at all (Theorem~\ref{thm:emit},
Remark~\ref{rem:emitclaim}). Nor does the step back to the
\termref{term-handed-model}{handed model}, the construction before the normalizer's \exttermref{../foundations/contract.pdf}{term-snap}{snap},
need an error term. No bound relates a key length computed from the declared
model alone to one computed from the handed model, and the displacement is
not what relates them (Theorem~\ref{thm:emit}\,(v)). The gap is not bridged
at all where a run can avoid it: a run that forms the number its key length
is formed from on the reference companion
(Corollary~\ref{cor:reference-companion}), built with no normalizer from the
\exttermref{../foundations/contract.pdf}{term-handed-reference}{handed reference}, or from the instance's own snapped blocks where none
is handed, reads the declared model only to find points. Where a run does read it, the gap is closed by a disclosed
premise about the device, never by an error term. The \termref{term-domination-link}{domination link},
which no requirement emits, is exact as well: a noise channel followed by its
target with classical flags (Theorem~\ref{thm:noise-channel}), so that a
squashing link followed by it is exact as a whole
(Corollary~\ref{cor:dom-chain}).
\end{remark}
\claimmeta{\metalabel{Uses}\mref{Def.}{def:criterion}, \mref{Def.}{def:opreduction}, \mref{Lem.}{lem:opreduction-transfer}, \mref{Thm.}{thm:factorization}, \mref{Thm.}{thm:transfer}, \mref{Thm.}{thm:chain}, \mref{Thm.}{thm:emit}
\metasep \metalabel{Used by}\mref{Req.}{req:no-approximate-relations}
\metasep newly derived
\metasep accepted}

\section{Output-state equality and security transfer}\label{group-reductions/thm:factorization}
\sectionintro{Replacing each round's physical measurement by the \termref{term-squashing-channel}{squashing channel} followed by the target measurement leaves the protocol's whole output state unchanged (\hyperlink{claim-reductions/thm:factorization}{Theorem~\ref*{thm:factorization}}), also against adaptive attack sequences in either source mode (\hyperlink{claim-reductions/cor:sequential}{Corollary~\ref*{cor:sequential}}). The main result, \hyperlink{claim-reductions/thm:transfer}{Theorem~\ref*{thm:transfer}}, is that security of the \termref{term-squashed-protocol}{squashed protocol} over analysed inputs that contain the squashed image of every physical input transfers to the physical protocol with the same \ensuremath{\varepsilon}, for every criterion. A constraint is sound when every reachable squashed state satisfies it; sound constraints may be intersected when they are sound for one common channel and one family of \termref{term-physical-inputs}{physical inputs}, and \termref{term-carried-operators}{carried operators} supply them (\hyperlink{claim-reductions/cor:sound}{Corollary~\ref*{cor:sound}}, \hyperlink{claim-reductions/cor:transported}{Corollary~\ref*{cor:transported}}). The remarks give the variants of the factorisation, the criteria that transfer, and the standard way a transplanted operator identity makes a constraint unsound.}
\sectionimports{\extref{../foundations/contract.pdf}{foundations/def:round-model}}

% claim reductions/thm:factorization
\claimsummary{reductions/thm:factorization}{Replacing each round's physical measurement by the squashing channel followed by the target measurement leaves the protocol's whole output state unchanged. Every security criterion therefore takes the same value on the physical and the squashed protocol, with no error term.}
\claimbookmark{reductions/thm:factorization}{Theorem}{thm:factorization}{Equality of the physical and squashed output states}
\claimleadin{%
A \termref{term-squashing-relation}{squashing relation} is a statement about one round's measurement.
Through the factorization of the \termref{term-measurement-channel}{measurement channel}
(Lemma~\ref{lem:msmfact}) it
becomes a statement about the whole protocol, and the proof uses no property
of the criterion $f_E$ beyond its being a function of the output state: not
linearity, not monotonicity, not continuity. The one structural hypothesis is
the factorization $\Grnd_j=\bar{\Grnd}_j\circ(\idmap\otimes\redNmsm_M)$ of
Definition~\ref{def:protocol}, and it rests on the measuring devices
entering only through one POVM per round (Assumption~\ref{ass:meas})
together with the round order and the one POVM for all rounds that the
definition declares. The round's order is not part of
Assumption~\ref{ass:meas}: that every public message of the round follows
every measurement of the round is A-SEQ (\xref{transmitter/ass:seq}).
}
\begin{theorem}[Equality of the physical and squashed output states]
\label{thm:factorization}
Let $M\dsq_{\Lambda}F$ over $\Vset$ and let $\Eprot_M,\Eprot_F$ be as in
Definition~\ref{def:protocol}. Then
\begin{equation}\label{eq:factorization}
  \Eprot_M\;=\;\Eprot_F\circ\big(\idmap\otimes\Lambda^{\otimes n}\big),
\end{equation}
the identity acting on every register other than $\Bsys_1\dots\Bsys_n$; and
consequently, for every finite-dimensional adversary system $E$ and every input
$\rho$,
\begin{equation}\label{eq:sameoutput}
  (\Eprot_M\otimes\idmap_E)(\rho)
  =(\Eprot_F\otimes\idmap_E)\big(\sigma\big),
  \qquad
  \sigma\coloneqq\big(\idmap\otimes\Lambda^{\otimes n}\otimes\idmap_E\big)(\rho).
\end{equation}
In particular $\Crit(\Eprot_M,\rho)=\Crit(\Eprot_F,\sigma)$ for every
\termref{term-security-criterion}{security criterion} in the sense of Definition~\ref{def:criterion}.
\end{theorem}
\claimmeta{\metalabel{Uses}\mref{Lem.}{lem:msmfact}, \mref{Def.}{def:protocol}, \mref{Def.}{def:criterion}
\metasep \metalabel{Assumes}\aref{A-MEAS}{ass:meas}, \aref{A-ROUND-SQ}{ass:round}, \extaref{../transmitter/contract.pdf}{transmitter/ass:seq}{A-SEQ}
\metasep \metalabel{Used by}\mref{Rem.}{rem:why-exact}, \mref{Def.}{def:roundmodel}, \mref{Rem.}{rem:variants}, \mref{Rem.}{rem:coherent}, \mref{Rem.}{rem:completeness}, \mref{Def.}{def:sound}, \mref{Thm.}{thm:transfer}, \mref{Cor.}{cor:monotone}, \mref{Thm.}{thm:chain}, \mref{Cor.}{cor:dom-chain}
\metasep newly derived
\metasep accepted
\metasep \metalabel{Oracles}\coderef{https://github.com/fraunhofersg/certiqs-key/blob/v0.1.0/tests/oracles/registry_reductions.py}{registry_reductions.py}}

\begin{proof}
By Lemma~\ref{lem:msmfact}, $\redNmsm_M=\redNmsm_F\circ\Lambda$. Substituting into
\eqref{eq:protocol},
\begin{equation*}
  \Grnd^{M}_j
  =\bar{\Grnd}_j\circ\big(\idmap\otimes\redNmsm_F\circ\Lambda\big)
  =\bar{\Grnd}_j\circ\big(\idmap\otimes\redNmsm_F\big)\circ\big(\idmap\otimes\Lambda_j\big)
  =\Grnd^{F}_j\circ\big(\idmap\otimes\Lambda_j\big),
\end{equation*}
where $\Lambda_j$ denotes $\Lambda$ acting on $\Bsys_j$ and the identity
elsewhere. For $i\ne j$ the maps $\Grnd^{F}_i$ and $\Lambda_j$ act on disjoint
tensor factors --- $\Grnd^{F}_i$ touches $A_i$, $\Qsys_i$ and classical
registers, while $\Lambda_j$ touches only $\Bsys_j$ --- so they commute.
Composing the $n$ rounds therefore moves every $\Lambda_j$ to the front, and
the $\Lambda_j$ assemble into $\Lambda^{\otimes n}$ because they act on
disjoint systems. Applying $\Kfin$, which touches no $\Bsys_j$, gives
\eqref{eq:factorization}. Equation~\eqref{eq:sameoutput} is
\eqref{eq:factorization} tensored with $\idmap_E$, and the equality of the
criteria is \eqref{eq:criterion} evaluated on the two equal states.
\end{proof}

% claim reductions/def:roundmodel
\claimsummary{reductions/def:roundmodel}{An attack sequence is the round model's, at the measured system; its squashed sequence applies the squashing channel to the measured system alone, leaving Alice's system and Eve's memory untouched. Where several parties measure, that channel is a product of one channel per party.}
\claimbookmark{reductions/def:roundmodel}{Definition}{def:roundmodel}{Round model and attack sequences, typed by source mode}
\claimleadin{%
Equation~\eqref{eq:sameoutput} is stronger than it may look: the two protocols
do not merely have close output states, they have the \emph{same} one. This is
why the transfer costs nothing and holds simultaneously for every criterion of
the form \eqref{eq:criterion}, with the same $\varepsilon$ on both sides. To
state the equality against an adversary who acts round by round, we need the
\exttermref{../foundations/contract.pdf}{term-attack-sequence}{attack sequences} of the \exttermref{../foundations/contract.pdf}{term-round-model}{round model} and their
squashed counterparts.
}
\begin{definition}[Round model and attack sequences, typed by source mode]
\label{def:roundmodel}
The \exttermref{../foundations/contract.pdf}{term-round-model}{round model} of the protocol \eqref{eq:protocol} is that of
\xref{foundations/def:round-model} with per-round system $\Bsys$, Alice's
system $A$ and the declared source mode. An \exttermref{../foundations/contract.pdf}{term-attack-sequence}{attack sequence} for the
protocol is thus a family $(\Aev_j)_{j=1}^{n}$ of channels whose output carries
the systems the adversary supplies in the declared source mode, the display
being that definition's at $S=\Bsys$:
\begin{equation}\label{eq:attackchannel}
  \Aev_j\in
  \begin{cases}
    \mathrm{CPTP}\big(E_{j-1},\ \Bsys_j\,E'_j\big),
      & \text{trusted-source mode},\\[0.3ex]
    \mathrm{CPTP}\big(E_{j-1},\ A_j\,\Bsys_j\,E'_j\big),
      & \text{untrusted-source mode},
  \end{cases}
\end{equation}
where $E_0$ is Eve's initial register, $E'_j$ her outgoing memory, and $E_j$
collects $E'_j$ together with whatever public data of round $j$ she retains. In
\exttermref{../transmitter/contract.pdf}{term-trusted-source-mode}{trusted-source mode} $A_j$ is \emph{not} an output of $\Aev_j$: it is
supplied by Alice and carried into the round as a separate input. The
protocol's output under the sequence is obtained from \eqref{eq:protocol} by
supplying $\Bsys_j$ --- and, in \exttermref{../transmitter/contract.pdf}{term-untrusted-source-mode}{untrusted-source mode}, $A_j$ as well
--- as the output of $\Aev_j$ in round $j$. The corresponding
\termdef{term-squashed-sequence}{squashed sequence} is
\begin{equation}\label{eq:composite}
  \Aev^{\sq}_j\;\coloneqq\;
  \begin{cases}
    \big(\Lambda\otimes\idmap_{E'_j}\big)\circ\Aev_j
      \in\mathrm{CPTP}\big(E_{j-1},\Qsys_jE'_j\big),
      & \text{trusted-source mode},\\[0.3ex]
    \big(\idmap_{A}\otimes\Lambda\otimes\idmap_{E'_j}\big)\circ\Aev_j
      \in\mathrm{CPTP}\big(E_{j-1},A_j\Qsys_jE'_j\big),
      & \text{untrusted-source mode},
  \end{cases}
\end{equation}
the \termref{term-squashing-channel}{squashing channel} acting on the $\Bsys_j$ factor in both cases and
on nothing else.

\emph{Several \termref{term-measuring-parties}{measuring parties}.} Where the \termref{term-measured-system}{measured system} of a
round is the product $\Bsys_j=\bigotimes_f\Bsys_{f,j}$ of party-scoped parts
(Definition~\ref{def:protocol}), the squashing channel is the product
$\Lambda=\bigotimes_f\Lambda_f$ of one channel per part,
$\Qsys_j=\bigotimes_f\Qsys_{f,j}$, and the alphabet $\Vset$ of $M$ and $F$ is
the product of the parts' alphabets. Neither line of \eqref{eq:attackchannel} changes:
$\Bsys_j$ is one output system of $\Aev_j$ whatever its factorisation, so the
adversary supplies all of its parts in one channel and is given the same
freedom as with one measuring party; in particular the parts may be
arbitrarily correlated, and the adversary is not required to act on them
independently. Nor does \eqref{eq:composite} change: the product channel
still acts on the $\Bsys_j$ factor alone, leaving $A_j$ and $E'_j$ untouched.
That the product channel realizes a squashing relation for the product family
is Proposition~\ref{prop:tensor}, which this definition does not assert. With
one measuring party the product has one factor.
\end{definition}
\claimmeta{\metalabel{Uses}\mref{Def.}{def:protocol}, \mref{Def.}{def:criterion}, \mref{Thm.}{thm:factorization}, \extref{../foundations/contract.pdf}{foundations/def:round-model}
\metasep \metalabel{Used by}\mref{Cor.}{cor:sequential}, \mref{Thm.}{thm:licence}, \mref{Lem.}{lem:source-transfer}, \mref{Thm.}{thm:two-sided}
\metasep newly derived
\metasep accepted}

% claim reductions/cor:sequential
\claimsummary{reductions/cor:sequential}{In either source mode, the physical protocol under an attack sequence and the squashed protocol under its squashed sequence have the same output state. Security against squashed attack families containing every squashed physical attack therefore transfers with the same \ensuremath{\varepsilon}.}
\claimbookmark{reductions/cor:sequential}{Corollary}{cor:sequential}{Output-state equality against adaptive attack sequences}
\claimleadin{%
Theorem~\ref{thm:factorization} fixes the adversary's input before the
protocol runs; an adversary who acts round by round, with a memory, is an
attack sequence of Definition~\ref{def:roundmodel}. The same substitution
handles it, one squashing channel per round. The corollary is the
receiver-side instance of an \termref{term-exact-operational-reduction}{exact operational reduction}, which
Theorem~\ref{thm:two-sided} composes with the source-side one in
trusted-source mode.
}
\begin{corollary}[Output-state equality against adaptive attack sequences]\label{cor:sequential}
Let $M\dsq_{\Lambda}F$ and fix a source mode. For every attack sequence
$(\Aev_j)_{j=1}^{n}$ of the corresponding type \eqref{eq:attackchannel}, the
output state of the physical protocol $\Eprot_M$ under $(\Aev_j)_j$ equals the
output state of the \termref{term-squashed-protocol}{squashed protocol} $\Eprot_F$ under $(\Aev^{\sq}_j)_j$ of
the same branch of \eqref{eq:composite}. Consequently every security criterion
takes the same value on the two, and if $\Eprot_F$ is $\varepsilon$-secure
against every sequence drawn from families $\AttSet^{\sq}_j$ with
\begin{equation}\label{eq:seqcover}
  \AttSet^{\sq}_j\;\supseteq\;
  \begin{cases}
    \big(\Lambda\otimes\idmap\big)\circ
      \mathrm{CPTP}\big(E_{j-1},\Bsys_jE'_j\big),
      & \text{trusted-source mode},\\[0.3ex]
    \big(\idmap_{A}\otimes\Lambda\otimes\idmap\big)\circ
      \mathrm{CPTP}\big(E_{j-1},A_j\Bsys_jE'_j\big),
      & \text{untrusted-source mode},
  \end{cases}
\end{equation}
then $\Eprot_M$ is $\varepsilon$-secure against \emph{all} attack sequences of
that type.
\end{corollary}
\claimmeta{\metalabel{Uses}\mref{Lem.}{lem:msmfact}, \mref{Def.}{def:protocol}, \mref{Def.}{def:opreduction}, \mref{Lem.}{lem:opreduction-transfer}, \mref{Def.}{def:roundmodel}
\metasep \metalabel{Used by}\mref{Prop.}{prop:platform}, \mref{Rem.}{rem:coherent}, \mref{Thm.}{thm:licence}, \mref{Rem.}{rem:publishedlicence}, \mref{Thm.}{thm:chain}, \mref{Thm.}{thm:two-sided}
\metasep newly derived
\metasep accepted}

\begin{proof}
\emph{Step 1 (equal output states).} Both branches are the same computation;
only the register the identity acts on differs. In either mode, each
occurrence of $\redNmsm_M$ in \eqref{eq:protocol} acts on the $\Bsys_j$-output
of $\Aev_j$ and on nothing else that $\Aev_j$ produces --- in untrusted-source
mode it does not touch the $A_j$-output either, which is measured by Alice's
POVM inside $\bar{\Grnd}_j$ (Definition~\ref{def:protocol}). Writing $\iota_j$
for $\idmap_{E'_j}$ in trusted mode and for $\idmap_{A}\otimes\idmap_{E'_j}$
in untrusted mode, Lemma~\ref{lem:msmfact} gives
\begin{equation*}
  \big(\idmap\otimes\redNmsm_M\otimes\idmap\big)\circ\Aev_j
  =\big(\idmap\otimes\redNmsm_F\otimes\idmap\big)\circ
   \big(\Lambda\ \text{on }\Bsys_j,\ \iota_j\ \text{elsewhere}\big)\circ\Aev_j
  =\big(\idmap\otimes\redNmsm_F\otimes\idmap\big)\circ\Aev^{\sq}_j ,
\end{equation*}
the middle map being exactly the branch of \eqref{eq:composite} for that mode.
Substituting round by round replaces every $(\redNmsm_M,\Aev_j)$ pair by
$(\redNmsm_F,\Aev^{\sq}_j)$ and changes nothing else in \eqref{eq:protocol}, so
the two output states are equal.

\emph{Step 2 (security, as the transfer of an exact operational reduction).}
Fix an admissible adversary system and let $\mathsf{P}$ be the
\termref{term-protocol-presentation}{presentation} whose strategies are the attack
sequences of the declared type against $\Eprot_M$, with the output map of
Step~1, and $\mathsf{P}'$ the presentation whose strategies are the sequences
drawn from the families $\AttSet^{\sq}_j$ against $\Eprot_F$. The
\termref{term-output-registers}{output registers} are the same in both --- those of
$\Kfin$ together with the adversary's final register --- and the map
$r_{\sq}\colon(\Aev_j)_j\mapsto(\Aev^{\sq}_j)_j$ of \eqref{eq:composite} takes
its values in $\mathsf{P}'$'s strategy set precisely because
\eqref{eq:seqcover} holds round by round. Step~1 is then
\eqref{eq:opreduction} for $r_{\sq}$, so $r_{\sq}$ is an exact operational
reduction (Definition~\ref{def:opreduction}), and the security statement is
Lemma~\ref{lem:opreduction-transfer}\,(i): the supremum over the physical
sequences equals the supremum over the image and is at most the supremum over
the families $\AttSet^{\sq}_j$, which is where the slack of
\eqref{eq:seqcover} is spent. That any larger families serve as well is part
(ii) of the same lemma. No property of the criterion is used, which is why the
conclusion holds for all of them at once.
\end{proof}

% claim reductions/prop:platform
\claimsummary{reductions/prop:platform}{For a well-formed declaration with one sector on A, one measured system, POVM form, two messages neither reading the other and no party-scoped entry, the round channel and its physical counterpart factor as in \hyperlink{claim-reductions/def:protocol}{Definition~\ref*{def:protocol}}. A squashing relation carries the one output to the other.}
\claimbookmark{reductions/prop:platform}{Proposition}{prop:platform}{Measurement factorization of the platform's protocol channel}
\claimleadin{%
The consuming engine's protocol channel is generated from a declaration. For a
declaration in the format its hypothesis fixes, the following proposition
shows that this channel, and the physical one defined by the same closed form
on the physical system, have the measurement factorization of
Definition~\ref{def:protocol}, so that Corollary~\ref{cor:sequential} applies
to them: an adversary who chooses round $j$ after seeing the announcements of
rounds $1,\dots,j-1$ is covered, since $E_{j-1}$ may carry the public record,
and the substitution is unchanged --- one $\Lambda$ per round, absorbed round
by round.
}
\begin{proposition}[Measurement factorization of the platform's protocol channel]\label{prop:platform}
Let $\Decl$ be a protocol declaration in the sense of
\xref{meat-core/def:declaration} that satisfies the conditions of
\xref{meat-core/def:admissible-declaration}, in the format in which $A$ carries one sector, one
system is measured, Alice's measurement is declared as a POVM
$\{M_x\}_{x\in\Xset}$, the public discussion is two messages neither of which
reads the other, and no entry is scoped to a party. Write $a$ for its joint
announcement map, $\phi$ for its acceptance-statistic map and $r$ for its
raw-key map, and let $\Mprot^{\sq}$ be the single-round protocol channel it
generates, on $A\Qsys$ with the declared POVM $F=\{F_v\}_{v\in\Vset}$ on the
measured system $\Qsys$, the system the declaration's measuring parties hold
(Definition~\ref{def:protocol}). Let $M=\{M_v\}_{v\in\Vset}$ be a POVM on a
separable $\Bsys$, and let $\Mprot^{\mathrm{ph}}$ be the map that the same
components define with $(\Bsys,M)$ in place of $(\Qsys,F)$,
\begin{equation}\label{eq:platformph}
  \Mprot^{\mathrm{ph}}(\omega)\coloneqq\sum_{x,v}
  \dyad{r(x,v)}_{Z}\otimes\dyad{\phi(a(x,v))}_{\Chat}\otimes
  \dyad{a(x,v)}_{I}\;\tr\big[(M_x\otimes M_v)\,\omega\big],
  \qquad\omega\in\Tcl{A\Bsys}.
\end{equation}
Then:
\begin{enumerate}[label=(\roman*),leftmargin=3.0em]
\item both maps satisfy the hypothesis of Definition~\ref{def:protocol} with
  one remaining processing $\bar{\Grnd}$,
  \begin{equation*}
    \Mprot^{\sq}=\bar{\Grnd}\circ\big(\idmap_A\otimes\redNmsm_F\big),
    \qquad
    \Mprot^{\mathrm{ph}}=\bar{\Grnd}\circ\big(\idmap_A\otimes\redNmsm_M\big),
  \end{equation*}
  where $\bar{\Grnd}$ measures $A$ with Alice's POVM $\{M_x\}$ and applies to
  the pair $(x,v)$ the classical map $\dyad{x,v}\mapsto
  \dyad{r(x,v)}_{Z}\otimes\dyad{\phi(a(x,v))}_{\Chat}\otimes\dyad{a(x,v)}_{I}$.
  Alice's measurement is thus a step of $\bar{\Grnd}$, where
  Definition~\ref{def:protocol} places the measurement of a party the
  declaration does not name as measuring, and not of the measurement channel;
\item if $M\dsq_{\Lambda}F$, then for every finite-dimensional $E$ and every
  $\omega\in\Sleq{A\Bsys E}$
  \begin{equation}\label{eq:platformfact}
    \big(\Mprot^{\mathrm{ph}}\otimes\idmap_E\big)(\omega)
    =\big(\Mprot^{\sq}\otimes\idmap_E\big)
     \big((\idmap_A\otimes\Lambda\otimes\idmap_E)(\omega)\big).
  \end{equation}
\end{enumerate}
\end{proposition}
\claimmeta{\metalabel{Uses}\mref{Lem.}{lem:sep}, \mref{Def.}{def:msm}, \mref{Def.}{def:protocol}, \mref{Cor.}{cor:sequential}
\metasep newly derived
\metasep accepted}

\begin{proof}
\emph{Step 1 (the two closed forms).} Under the conditions of
\xref{meat-core/def:admissible-declaration}, the closed form of \xref{meat-core/thm:protocol-channel}\,(ii)
reads, verbatim, with $\Mprot=\Mprot^{\sq}$,
\begin{equation}\label{eq:povmform}
\begin{split}
  (\Mprot\otimes\idmap_E)(\omega)
  ={}&\sum_{x,v}
  \dyad{r(x,v)}_{Z}\otimes\dyad{\phi(a(x,v))}_{\Chat}\otimes{}\\
  &\qquad\dyad{a(x,v)}_{I}\otimes
  \tr_{A\Qsys}\!\big[(M_x\otimes F_v\otimes\idon{E})\,\omega\big]
\end{split}
\end{equation}
for every system $E$ and every $\omega\in\Sleq{A\Qsys E}$. Index letters
carry the whole distinction here: in \eqref{eq:povmform},
\eqref{eq:platformph} and \eqref{eq:jointtransport}, $M_x$ is Alice's element
on $A$, quoted from the declaration, while $M_v$ is the physical element on
the measured system $\Bsys$ as defined in this contract; they act on
different tensor factors and the squashing channel touches only the second.
The map \eqref{eq:platformph} has the same form with $(\Bsys,M)$ in place of
$(\Qsys,F)$: for a finite-dimensional $E$ with orthonormal basis
$\{\ket{i}\}_i$, every $\omega\in\Tcl{A\Bsys E}$ is
$\sum_{i,j}\omega_{ij}\otimes\ketbra{i}{j}_E$ with $\omega_{ij}\in\Tcl{A\Bsys}$,
and $\sum_{i,j}\tr[(M_x\otimes M_v)\,\omega_{ij}]\,\ketbra{i}{j}_E
=\tr_{A\Bsys}[(M_x\otimes M_v\otimes\idon{E})\,\omega]$, so by linearity
\begin{equation}\label{eq:povmformph}
\begin{split}
  (\Mprot^{\mathrm{ph}}\otimes\idmap_E)(\omega)
  ={}&\sum_{x,v}
  \dyad{r(x,v)}_{Z}\otimes\dyad{\phi(a(x,v))}_{\Chat}\otimes{}\\
  &\qquad\dyad{a(x,v)}_{I}\otimes
  \tr_{A\Bsys}\!\big[(M_x\otimes M_v\otimes\idon{E})\,\omega\big].
\end{split}
\end{equation}

\emph{Step 2 (the factorization).} Let $(S,G)$ be $(\Qsys,F)$ or
$(\Bsys,M)$. For $\omega\in\Tcl{AS}$, \eqref{eq:msm} gives
$(\idmap_A\otimes\redNmsm_G)(\omega)
=\sum_v\tr_S\big[(\idon{A}\otimes G_v)\,\omega\big]\otimes\dyad{v}_Y$, and
$\bar{\Grnd}$ maps $\tau\otimes\dyad{v}_Y$, $\tau\in\Tcl{A}$, to
$\sum_x\tr[M_x\tau]\,\dyad{r(x,v)}_{Z}\otimes\dyad{\phi(a(x,v))}_{\Chat}
\otimes\dyad{a(x,v)}_{I}$. Hence
\begin{equation*}
  \bar{\Grnd}\big((\idmap_A\otimes\redNmsm_G)(\omega)\big)
  =\sum_{x,v}\dyad{r(x,v)}_{Z}\otimes\dyad{\phi(a(x,v))}_{\Chat}\otimes
  \dyad{a(x,v)}_{I}\;\tr\big[(M_x\otimes G_v)\,\omega\big],
\end{equation*}
which for $(\Bsys,M)$ is \eqref{eq:platformph} and for $(\Qsys,F)$ is
\eqref{eq:povmform} at trivial $E$, extended from $\Sleq{A\Qsys}$ to
$\Tcl{A\Qsys}$ by linearity, the finite-dimensional $A\Qsys$ being spanned by
its states. $\bar{\Grnd}$ is a channel, a measurement of $A$ followed by a
classical map, and it touches neither $\Qsys$ nor $\Bsys$; that is the
hypothesis of Definition~\ref{def:protocol} for both maps.

\emph{Step 3 (the substitution).} For \eqref{eq:platformfact} it therefore
suffices, by \eqref{eq:povmformph} and \eqref{eq:povmform}, that for all $x$
and $v$
\begin{equation}\label{eq:jointtransport}
  \tr_{A\Bsys}\!\big[(M_x\otimes M_v\otimes\idon{E})\,\omega\big]
  =\tr_{A\Qsys}\!\big[(M_x\otimes F_v\otimes\idon{E})\,
    (\idmap_A\otimes\Lambda\otimes\idmap_E)(\omega)\big]
\end{equation}
as operators on $E$, the state $(\idmap_A\otimes\Lambda\otimes\idmap_E)(\omega)$
lying in $\Sleq{A\Qsys E}$. Test both sides against an arbitrary
$X\in\Lin{E}$. By \eqref{eq:pairing} applied to
$\idmap_A\otimes\Lambda\otimes\idmap_E$ and Lemma~\ref{lem:sep}\,(ii),
\begin{equation*}
  \tr\big[(M_x\otimes F_v\otimes X)(\idmap_A\otimes\Lambda\otimes\idmap_E)(\omega)\big]
  =\tr\big[\big(M_x\otimes\Lambda^{\dagger}(F_v)\otimes X\big)\,\omega\big]
  =\tr\big[(M_x\otimes M_v\otimes X)\,\omega\big],
\end{equation*}
and operators on the finite-dimensional $E$ separate, so the two partial traces
agree. Substituting into \eqref{eq:povmformph} and \eqref{eq:povmform} gives
\eqref{eq:platformfact}.
\end{proof}

% claim reductions/rem:platform-published
\claimsummary{reductions/rem:platform-published}{The reduction of \hyperlink{claim-reductions/cor:sequential}{Corollary~\ref*{cor:sequential}} and \hyperlink{claim-reductions/prop:platform}{Proposition~\ref*{prop:platform}} is published for narrower protocol models, with a restricted or an unrestricted squashed attack set, and in a non-adaptive single-channel version. No proof here uses those statements.}
\claimbookmark{reductions/rem:platform-published}{Remark}{rem:platform-published}{Published forms of the sequential reduction}
\begin{remark}[Published forms of the sequential reduction]\label{rem:platform-published}
The published statement of the reduction of Corollary~\ref{cor:sequential}
and Proposition~\ref{prop:platform}, for a narrower protocol model, is
\cite[Lemma~9.3, Eqs.~(146)--(148)]{tupkary2026decoy}; the case in which the
squashed \exttermref{../foundations/contract.pdf}{term-attack-set}{attack set} is left unrestricted is its Rem.~9.1, and the
non-adaptive, single-channel version with an unrestricted set is
\cite[Lemma~6]{nahar2024postselection}. Neither is used in the proofs of
Corollary~\ref{cor:sequential} and Proposition~\ref{prop:platform}.
\end{remark}
\claimmeta{newly derived
\metasep accepted}

% claim reductions/rem:variants
\claimsummary{reductions/rem:variants}{The factorization holds verbatim for round-dependent squashers, for a private per-round setting squashed by one common channel, and for a correlated n-round device. A setting-dependent squashing channel does not fit, and hands the setting register to the adversary.}
\claimbookmark{reductions/rem:variants}{Remark}{rem:variants}{Variants of the factorization}
\begin{remark}[Variants of the factorization]\label{rem:variants}
The proof of Theorem~\ref{thm:factorization} uses only linearity and the
single-round identity, so it applies verbatim to round-dependent squashers
($M^{(j)}\dsq_{\Lambda_j}F^{(j)}$, composite channel
$\bigotimes_j\Lambda_j$); to a private per-round setting drawn from a family
$\{M^{(s)}\}_s$ squashed by a \emph{common} channel with
$\Lambda^{\dagger}(F^{(s)}_v)=M^{(s)}_v$, where the controlled measurement
channel factorizes through $\idmap_S\otimes\Lambda$; and to a genuinely
correlated $n$-round device, where a joint POVM $\{M_{\vec v}\}$ on
$\Bsys^{n}$ and a channel $\Lambda^{(n)}$ with
$\Lambda^{(n)\dagger}(F_{v_1}\otimes\dots\otimes F_{v_n})=M_{\vec v}$ give the
same conclusion. A device with detector memory across rounds, which
Assumption~\ref{ass:round} excludes, is a joint POVM on $\Bsys^{n}$, so
where a channel $\Lambda^{(n)}$ as above exists, \eqref{eq:sameoutput} holds
for it at a fixed input; nothing is asserted of it in the round model. What
Assumption~\ref{ass:round} fixes is that one physical POVM describes every
round, with no memory across rounds; that the declared inputs describe one
round and are applied to all $n$, with one squashing channel for every round,
is the case the engine emits
(Requirement~\ref{req:round-homogeneous-emission}).
A setting-\emph{dependent} family $\Lambda^{(s)}$ does
\emph{not} fit this scheme and yields only the weaker model in which the
setting register is handed to the adversary; that is why the probe setting
enters this contract as a joint outcome under one common channel
(Definition~\ref{def:joint-family}, Lemma~\ref{lem:linear}\,(ii)) and never as
a setting-dependent squasher. If the setting is not drawn as
Assumption~\ref{ass:probe} says, the \termref{term-joint-family}{joint family} $\{q_tM_{y|t}\}$ is not a
POVM with a common squashing channel and the setting register would have to
be handed to the adversary; a probe setting the adversary can predict or
influence also breaks the countermeasure it models, outside this contract.
\end{remark}
\claimmeta{\metalabel{Uses}\mref{Thm.}{thm:factorization}
\metasep \metalabel{Assumes}\aref{A-ROUND-SQ}{ass:round}
\metasep \metalabel{Used by}\mref{Rem.}{rem:coherent}, \mref{Req.}{req:round-homogeneous-emission}
\metasep newly derived
\metasep accepted}

% claim reductions/rem:coherent
\claimsummary{reductions/rem:coherent}{The output-state equality holds for an arbitrary adversary input and an arbitrary attack sequence, so it covers an adversary entangling all n rounds. The tensor power is a hypothesis about Bob's device, not about the adversary.}
\claimbookmark{reductions/rem:coherent}{Remark}{rem:coherent}{Coherent attacks covered by the output-state equality}
\begin{remark}[Coherent attacks covered by the output-state equality]\label{rem:coherent}
Nothing in Theorem~\ref{thm:factorization} or
Corollary~\ref{cor:sequential} refers to the structure of the adversary's
input: $\rho$ is an arbitrary state on $A^{n}\Bsys^{n}E$ and $(\Aev_j)_j$ an
arbitrary sequence of channels, so the conclusion covers Eve entangling all $n$
rounds and holding $E$ correlated with everything. The tensor power in
\eqref{eq:factorization} is a hypothesis about Bob's \emph{device} --- the same
POVM every round, with no memory between rounds
(Assumptions~\ref{ass:meas}, \ref{ass:round}) --- and is itself relaxed in
Remark~\ref{rem:variants}. It is not a hypothesis about the adversary.
\end{remark}
\claimmeta{\metalabel{Uses}\mref{Thm.}{thm:factorization}, \mref{Cor.}{cor:sequential}, \mref{Rem.}{rem:variants}
\metasep \metalabel{Assumes}\aref{A-MEAS}{ass:meas}, \aref{A-ROUND-SQ}{ass:round}
\metasep newly derived
\metasep accepted}

% claim reductions/rem:criteria
\claimsummary{reductions/rem:criteria}{Correctness and secrecy in the composable sense are functions of the output state, so equality of output states transfers them. Completeness compares the protocol against an honest reference channel, is not of that form and is not transferred.}
\claimbookmark{reductions/rem:criteria}{Remark}{rem:criteria}{Correctness, secrecy and completeness as criteria}
\begin{remark}[Correctness, secrecy and completeness as criteria]\label{rem:criteria}
Correctness and secrecy in the composable sense ---
$\Pr[K_A\neq K_B]\le\varepsilon_{\rm cor}$ and
$(1-p_{\mathrm{abort}})\,D(\rho_{AE},\pi_A\otimes\rho_E)\le
\varepsilon_{\rm sec}$ of \cite[Eqs.~(13)--(14)]{portmann2022security}, in
which $\pi_A$ is the \emph{fully mixed key state} (that source writes $\tau_A$
for it; it is \emph{not} the trusted marginal $\tau_A$ of
Definition~\ref{def:licensed}, and the two never occur in one formula here)
--- are both functions of the output state in the sense of
Definition~\ref{def:criterion}, because the ideal object each compares
against is a fixed channel applied to the protocol's own output. What is
\emph{not} of that form, and is therefore not transferred by
Theorem~\ref{thm:transfer}, is a statement comparing the protocol against an
\emph{honest reference channel}: the completeness criterion of
\cite[Eq.~(9)]{portmann2022security} --- the second construction condition,
evaluated against a blank adversarial interface, from which that paper derives
robustness at its \cite[Eq.~(15)]{portmann2022security} --- and completeness
generally. Those are the concern of a different module, and
Remark~\ref{rem:completeness} records why no gap opens.
\end{remark}
\claimmeta{\metalabel{Uses}\mref{Def.}{def:criterion}
\metasep \metalabel{Assumes}\aref{A-MEAS}{ass:meas}
\metasep newly derived
\metasep accepted}

% claim reductions/rem:completeness
\claimsummary{reductions/rem:completeness}{A completeness criterion compares the protocol against an honest reference channel and is outside the form that transfers. No gap opens: the squashed model reproduces the physical statistics exactly, so a completeness statement computed in one model holds in the other.}
\claimbookmark{reductions/rem:completeness}{Remark}{rem:completeness}{Completeness statements in the physical and squashed models}
\begin{remark}[Completeness statements in the physical and squashed models]
\label{rem:completeness}
A criterion comparing the protocol against an \emph{honest reference channel}
--- ``on the honest channel the protocol aborts with probability at most
$\varepsilon$'', the completeness criterion of
\cite[Eq.~(9)]{portmann2022security} (robustness is its phrasing at that
paper's Eq.~(15)) --- is not of the form
\eqref{eq:criterion} evaluated on an adversarial input, and
Theorem~\ref{thm:factorization} says nothing about it.
No gap opens: the squashed model reproduces the physical outcome statistics
\emph{exactly} (Lemma~\ref{lem:statistics}\,(ii)), so the honest distribution
on $\Vset$ is the same object in both models, and a completeness statement
computed in one is literally the statement in the other. Completeness budgets
belong to the statistics module, and nothing in this contract carries one.
\end{remark}
\claimmeta{\metalabel{Uses}\mref{Lem.}{lem:statistics}, \mref{Def.}{def:criterion}, \mref{Thm.}{thm:factorization}
\metasep newly derived
\metasep accepted}

% claim reductions/def:sound
\claimsummary{reductions/def:sound}{A squashed state is reachable when it is the squashed image of a physical input, and a family of squashed states, or a constraint, is sound when it contains every reachable state. Soundness is what keeps a constrained analysis from excluding a physical attack.}
\claimbookmark{reductions/def:sound}{Definition}{def:sound}{Reachable states and sound constraints}
\claimleadin{%
Theorem~\ref{thm:factorization} leaves exactly one thing to check, and the
following definition names it: the squashed analysis must run over a set of
inputs containing the squashed images of the physical ones. Everything a
constrained squashed analysis can legitimately assume is the supply of such
sets, and every way such an analysis can fail is a constraint some
\emph{reachable} state violates. Both notions depend on the channel and on
the physical inputs; Remark~\ref{rem:channel-choice} says what has to be
watched when these are not fixed.
}
\begin{definition}[Reachable states and sound constraints]\label{def:sound}
Let $M\dsq_{\Lambda}F$ and let $\Sph=(\Sph_E)_E$ be a family of
\termdef{term-physical-inputs}{physical inputs}, one set
$\Sph_E\subseteq\Seq{A^{n}\Bsys^{n}E}$ per admissible adversary system $E$. A
state $\sigma\in\Seq{A^{n}\Qsys^{n}E}$ is \termdef{term-reachable}{reachable} if
\begin{equation}\label{eq:reachable}
  \sigma=\big(\idmap_{A^{n}}\otimes\Lambda^{\otimes n}\otimes\idmap_E\big)(\rho)
  \qquad\text{for some }\rho\in\Sph_E ,
\end{equation}
and $\Srch_E$ denotes the set of reachable states. A family $C=(C_E)_E$ of
subsets $C_E\subseteq\Seq{A^{n}\Qsys^{n}E}$, and any constraint whose solution
set it is, is \termdef{term-sound}{sound} if $\Srch_E\subseteq C_E$ for every admissible $E$. Both notions
depend on $\Lambda$ and on $\Sph$, which are suppressed in the notation and
fixed throughout any one analysis.
\end{definition}
\claimmeta{\metalabel{Uses}\mref{Thm.}{thm:factorization}
\metasep \metalabel{Used by}\mref{Cor.}{cor:sound}, \mref{Cor.}{cor:transported}, \mref{Cor.}{cor:blockdiag}, \mref{Def.}{def:link}, \mref{Cor.}{cor:normalized-flat}, \mref{Def.}{def:domination-link}
\metasep newly derived
\metasep accepted}

% claim reductions/thm:transfer
\claimsummary{reductions/thm:transfer}{If the analysed squashed inputs contain the squashed image of every physical input, a bound \ensuremath{\varepsilon} on any criterion of the output state of the squashed protocol holds for the physical protocol with the same \ensuremath{\varepsilon}. No property of the criterion is used.}
\claimbookmark{reductions/thm:transfer}{Theorem}{thm:transfer}{Security transfer from a constrained squashed model}
\claimleadin{%
With the output state unchanged (Theorem~\ref{thm:factorization}), security
of the squashed protocol carries over as soon as the squashed analysis runs
over enough inputs. Theorem~\ref{thm:two-sided} composes the receiver-side
reduction with the source map in trusted-source mode.
}
\begin{mainclaim}
\begin{theorem}[Security transfer from a constrained squashed model]\label{thm:transfer}
Let $M\dsq_{\Lambda}F$ and let $\Eprot_M,\Eprot_F$ be as in
Definition~\ref{def:protocol}. For every admissible adversary system $E$ let
$\Sph_E\subseteq\Seq{A^{n}\Bsys^{n}E}$ and
$\Ssq_E\subseteq\Seq{A^{n}\Qsys^{n}E}$ satisfy
\begin{equation}\label{eq:inclusion}
  \big(\idmap_{A^{n}}\otimes\Lambda^{\otimes n}\otimes\idmap_E\big)
  \big(\Sph_E\big)\;\subseteq\;\Ssq_E .
\end{equation}
If $\Crit(\Eprot_F,\sigma)\le\varepsilon$ for all $\sigma\in\Ssq_E$ and all
$E$, then $\Crit(\Eprot_M,\rho)\le\varepsilon$ for all $\rho\in\Sph_E$ and all
$E$ --- with the \emph{same} $\varepsilon$, and simultaneously for every
security criterion in the sense of Definition~\ref{def:criterion}.
\end{theorem}
\end{mainclaim}
\claimmeta{\metalabel{Uses}\mref{Def.}{def:protocol}, \mref{Def.}{def:criterion}, \mref{Def.}{def:opreduction}, \mref{Lem.}{lem:opreduction-transfer}, \mref{Thm.}{thm:factorization}
\metasep \metalabel{Used by}\mref{Rem.}{rem:why-exact}, \mref{Cor.}{cor:sound}, \mref{Rem.}{rem:unsound}, \mref{Rem.}{rem:channel-choice}, \mref{Thm.}{thm:chain}, \mref{Thm.}{thm:emit}, \mref{Cor.}{cor:reference-companion}, \mref{Rem.}{rem:tworoutes}, \mref{Cor.}{cor:dom-chain}, \extref{../protocol-compiler/contract.pdf}{protocol-compiler/req:link-row-provenance}
\metasep newly derived
\metasep accepted
\metasep \metalabel{Oracles}\coderef{https://github.com/fraunhofersg/certiqs-key/blob/v0.1.0/tests/oracles/registry_reductions.py}{registry_reductions.py}}

\begin{proof}
Fix an admissible adversary system $E$ and put
$r(\rho)\coloneqq
(\idmap_{A^{n}}\otimes\Lambda^{\otimes n}\otimes\idmap_E)(\rho)$. Let
$\mathsf{P}$ be the presentation whose strategies are the states
$\rho\in\Sph_E$ --- a strategy here is an input state and not a channel, which
Definition~\ref{def:opreduction} allows and says --- with
$\mathsf{out}(\rho)=(\Eprot_M\otimes\idmap_E)(\rho)$, and $\mathsf{P}'$ the
presentation with strategies $\Ssq_E$ and
$\mathsf{out}'(\sigma)=(\Eprot_F\otimes\idmap_E)(\sigma)$; both have the output
registers of $\Kfin$ together with $E$, which is condition (R1). The
hypothesis \eqref{eq:inclusion} says exactly that $r$ takes its values in
$\Ssq_E$, which is condition (R3), and \eqref{eq:sameoutput} of
Theorem~\ref{thm:factorization} is \eqref{eq:opreduction} for $r$. So $r$ is an
exact operational reduction from $\mathsf{P}$ to $\mathsf{P}'$, and
Lemma~\ref{lem:opreduction-transfer}\,(i), applied to the criterion $f_E$ at
this $E$, gives
$\Crit(\Eprot_M,\rho)=\Crit(\Eprot_F,r(\rho))\le\varepsilon$ for every
$\rho\in\Sph_E$, with no property of $f_E$ used. The quantification over $E$ is
one presentation pair and one reduction per admissible $E$
(Definition~\ref{def:opreduction}, closing paragraph).
\end{proof}

% claim reductions/cor:sound
\claimsummary{reductions/cor:sound}{Analysed squashed inputs satisfy the inclusion of \hyperlink{claim-reductions/thm:transfer}{Theorem~\ref*{thm:transfer}} exactly when they form a sound family, and any intersection of sound families for one channel and one physical input family is sound. Enlarging a sound family costs tightness, never soundness.}
\claimbookmark{reductions/cor:sound}{Corollary}{cor:sound}{Imposing and intersecting sound constraints}
\begin{corollary}[Imposing and intersecting sound constraints]
\label{cor:sound}
Let $M\dsq_{\Lambda}F$ and fix a family $\Sph$ of \termref{term-physical-inputs}{physical inputs}.
\begin{enumerate}[label=(\roman*),leftmargin=3.0em]
\item A family $\Ssq$ satisfies \eqref{eq:inclusion} if and only if it is
  sound.
\item An arbitrary intersection of sound families is sound.
\item Consequently, if $(C^{(i)})_{i\in I}$ is any collection of sound families
  and $\Crit(\Eprot_F,\sigma)\le\varepsilon$ for every admissible $E$ and every
  $\sigma\in\bigcap_{i\in I}C^{(i)}_E$, then $\Crit(\Eprot_M,\rho)\le
  \varepsilon$ for every admissible $E$ and every $\rho\in\Sph_E$, with the
  same $\varepsilon$ and for every criterion of the form \eqref{eq:criterion}.
\item The tightest sound family is $\Srch$ itself and the loosest is
  $\Ssq_E=\Seq{A^{n}\Qsys^{n}E}$; enlarging a sound family costs tightness and
  never soundness.
\end{enumerate}
\end{corollary}
\claimmeta{\metalabel{Uses}\mref{Def.}{def:criterion}, \mref{Def.}{def:opreduction}, \mref{Lem.}{lem:opreduction-transfer}, \mref{Def.}{def:sound}, \mref{Thm.}{thm:transfer}
\metasep \metalabel{Used by}\mref{Rem.}{rem:channel-choice}, \mref{Cor.}{cor:blockdiag}
\metasep newly derived
\metasep accepted}

\begin{proof}
(i) The left-hand side of \eqref{eq:inclusion} is $\Srch_E$ by
\eqref{eq:reachable}, so the inclusion says exactly that $\Ssq_E$ contains
every reachable state --- equivalently, that the strategy map
$r\colon\rho\mapsto(\idmap_{A^{n}}\otimes\Lambda^{\otimes n}\otimes
\idmap_E)(\rho)$ of the proof of Theorem~\ref{thm:transfer}, the one belonging
to \emph{this} $\Lambda$ and \emph{this} $\Sph$, has its image in $\Ssq_E$,
which is condition (R3) of Definition~\ref{def:opreduction}. (ii) A reachable
state lies in every member of the collection, hence in the intersection.
(iii) By (i), soundness of each $C^{(i)}$ says that \emph{one and the same}
map $r$ has its image in $C^{(i)}_E$, for every $i$ and every admissible $E$;
Lemma~\ref{lem:opreduction-transfer}\,(iv) therefore makes $r$ an exact
operational reduction into $\bigcap_iC^{(i)}_E$, and its part (i) gives the
bound with the same $\varepsilon$ for every criterion. The hypothesis is one
map and not one map per $i$, and this matters: part (iv) of that lemma
exhibits analysed sets, each of them reached by some exact operational
reduction, whose intersection is empty and over which every bound holds
vacuously; the concrete form of that failure here is two channels realizing
one squashing relation, which is why this corollary fixes one $\Lambda$ and
one $\Sph$ (Remark~\ref{rem:channel-choice}). (iv) Immediate from (i); the
loosest family satisfies \eqref{eq:inclusion} trivially, which is the
unconstrained squashing statement.
\end{proof}

% claim reductions/cor:transported
\claimsummary{reductions/cor:transported}{On the n-round relation, a transported operator tested on a squashed image gives what its pull-back gives on the physical input. Certified operator facts, and bounds assumed on the physical input family, therefore give sound constraints on the squashed states.}
\claimbookmark{reductions/cor:transported}{Corollary}{cor:transported}{Sound constraints from transported operators}
\claimleadin{%
Part (iii) of Corollary~\ref{cor:sound} is what a constrained analysis runs
on: soundness is established once per constraint and never again for the
combination, so constraints from unrelated sources --- a device property, a
calibration bound, an assumption on Eve's input --- may be imposed
simultaneously, \emph{provided they are sound for one and the same squashing
channel} (Remark~\ref{rem:channel-choice}). What a squashed analysis can do is
therefore fixed by the supply of sound constraints, and that supply is
Corollary~\ref{cor:pullback} read on the $n$-round relation.
}
\begin{corollary}[Sound constraints from transported operators]
\label{cor:transported}
Let $(M;O)\dsq_{\Lambda}(F;\hat O)$, fix a family $\Sph$ of physical inputs and
write $R\coloneqq A^{n}E$. Let $\hat O^{(n)}\in\Lin{\Qsys^{n}}$ be Hermitian
and let $O^{(n)}\coloneqq(\Lambda^{\dagger})^{\otimes n}(\hat O^{(n)})$ be its
pull-back. Then
\begin{equation}\label{eq:transportn}
  \tr_{\Qsys^{n}}\big[(\hat O^{(n)}\otimes\idon{R})\,\sigma\big]
  =\tr_{\Bsys^{n}}\big[(O^{(n)}\otimes\idon{R})\,\rho\big]
\end{equation}
for every $\rho\in\Sph_E$ and its image $\sigma$, and each of the following
families of squashed states is sound:
\begin{enumerate}[label=(\roman*),leftmargin=3.0em]
\item \emph{(certified)} for each \termref{term-certificate}{certificate} case of
  Corollary~\ref{cor:pullback} --- identity, positivity, spectral bounds ---
  read with $\Bsys^{n},\Qsys^{n},\Lambda^{\otimes n},R$ in place of
  $\Bsys,\Qsys,\Lambda,R$, the family of squashed states satisfying its
  conclusion;
\item \emph{(assumed)}
  $\{\sigma:\tr[(\hat O^{(n)}\otimes X)\sigma]\bowtie c\}$ for a fixed
  $\bowtie\,\in\{\le,=,\ge\}$, a fixed $X\in\Lin{R}$, Hermitian unless
  $\bowtie$ is $=$, and a fixed $c$, provided
  $\tr[(O^{(n)}\otimes X)\rho]\bowtie c$ holds for every $\rho\in\Sph_E$.
\end{enumerate}
Operators with a prescribed pull-back are supplied round by round: if
$\hat O^{(n)}$ is $\hat O_k$ on $\Qsys_j$ and $\idon{\Qsys}$ on the remaining
rounds, then $O^{(n)}$ is $O_k$ on $\Bsys_j$ and $\idon{\Bsys}$ elsewhere; and
a product $\hat O_{k_1}\otimes\dots\otimes\hat O_{k_n}$ pulls back to
$O_{k_1}\otimes\dots\otimes O_{k_n}$.
\end{corollary}
\claimmeta{\metalabel{Uses}\mref{Lem.}{lem:sep}, \mref{Cor.}{cor:pullback}, \mref{Def.}{def:certificate}, \mref{Def.}{def:sound}
\metasep \metalabel{Used by}\mref{Rem.}{rem:channel-choice}, \mref{Rem.}{rem:reading-master}, \mref{Thm.}{thm:licence}, \mref{Thm.}{thm:chain}, \mref{Req.}{req:certified-entries-only}, \mref{Rem.}{rem:published}
\metasep newly derived
\metasep accepted}

\begin{proof}
$\Lambda^{\otimes n}$ is a channel $\Tcl{\Bsys^{n}}\to\Tcl{\Qsys^{n}}$ whose
adjoint is $(\Lambda^{\dagger})^{\otimes n}$ by Lemma~\ref{lem:sep}\,(ii).
Reordering $A^{n}\Bsys^{n}E$ as $\Bsys^{n}R$, the channel of
\eqref{eq:reachable} is $\Lambda^{\otimes n}\otimes\idmap_R$, so
Corollary~\ref{cor:pullback} applies with
$\Bsys^{n},\Qsys^{n},\Lambda^{\otimes n},R$ in place of
$\Bsys,\Qsys,\Lambda,R$ and with the single carried pair
$(O^{(n)},\hat O^{(n)})$. Equation~\eqref{eq:transportn} is its
\eqref{eq:transport}, and each conclusion of
Corollary~\ref{cor:pullback}\,(i)--(iii) holds on \emph{every} reachable
$\sigma$, which is soundness of the families in (i). For (ii), test
\eqref{eq:transport} against $X$ and bound the right-hand side by hypothesis.
For the round-by-round supply, $(\Lambda^{\dagger})^{\otimes n}$ acts
factorwise and $\Lambda^{\dagger}$ is unital.
\end{proof}

% claim reductions/rem:unsound
\claimsummary{reductions/rem:unsound}{Imposed as an operator identity on a target that fails it, the physical class-marginal identity would exclude reachable states and void \hyperlink{claim-reductions/thm:transfer}{Theorem~\ref*{thm:transfer}}. Its sound squashed form is the expectation identity, whose certificate is available exactly where the reduction instance discharges it.}
\claimbookmark{reductions/rem:unsound}{Remark}{rem:unsound}{The standard way to make a squashed constraint unsound}
\claimleadin{%
Case (i) of Corollary~\ref{cor:transported} carries the identity and
inequality certificates: they hold on every reachable state whatever the
adversary is allowed to send, so the resulting families are sound for
\emph{every} choice of $\Sph$ and may be adjoined to any analysis. Its case
(ii) is an assumption about the input and is sound only for the $\Sph$ it was
assumed on. The following remark describes the standard way a constraint
fails to be sound at all.
}
\begin{remark}[The standard way to make a squashed constraint unsound]
\label{rem:unsound}
Equation~\eqref{eq:inclusion} is the whole of soundness, and the standard way
to violate it is to transplant an operator identity that holds in the
\emph{physical} model. Take a private setting of a measuring party's. Where the physical
model \emph{is} a \termref{term-joint-family}{joint family} (Definition~\ref{def:joint-family}\,(b)) ---
one branch, $\Bsys$ the system the conditional families live on --- the
physical POVM satisfies the operator identity
\begin{equation}\label{eq:probeidentity}
  \sum_{v\in\Vset_t}M_v\;=\;q_t\,\idon{\Bsys}
  \qquad\text{for every }t\in\Tset ,
\end{equation}
which is the \termref{term-class-marginal-identity}{class-marginal identity} \eqref{eq:classmarginal}: the
weight $q_t$ is that party's own and
\eqref{eq:jointfamily} \emph{builds} $M_{(y,t)}=q_t\widetilde M_{y|t}$ out of
a conditional family that is complete at each fixed $t$
(Lemma~\ref{lem:linear}\,(ii), a statement about conditional families).

\emph{When \eqref{eq:probeidentity} is available, and when it is not.} It is
a theorem of the construction where the model is a joint family outright, and
it is a property of a \termref{term-reduction-instance}{reduction instance}'s model only where the instance
discharges it (Definition~\ref{def:reduction-instance}\,(I4)). An instance built by the
branching rule inherits it from its branches and from nothing else
(Theorem~\ref{thm:instrument-instance}, Corollary~\ref{cor:register-instance}):
a branch whose POVM is an arbitrary one --- the flag instance of an
uncharacterised branch, Corollary~\ref{cor:flag-instance}, on which nothing
is assumed and which may be undeclared --- contributes to the class sum
merely \emph{some} operator between $0$ and its branch weight, and the
identity on $\Bsys$ is then neither a theorem of this contract nor a
validated datum unless the branch's declaration establishes it. Whether a
declaration can establish it for such a branch is a question about what the
declaration presents (Remark~\ref{rem:exactboundary}). The identity is
discharged only where every branch establishes it; an undeclared branch, a
branch missing an element and a branch that does not establish the identity
alike leave it undischarged, and the entry that rests on it is then refused
and never emitted (Requirement~\ref{req:class-marginal-entry}). Where it is
not discharged, \eqref{eq:probeidentity} may still hold of the device, but
nothing declared establishes it. The squashed operators, where the identity
is discharged, satisfy
$\Lambda^{\dagger}\big(\sum_{v\in\Vset_t}F_v\big)=q_t\,\idon{\Bsys}$, while
$\sum_{v\in\Vset_t}F_v\ne q_t\,\idon{\Qsys}$ on every \termref{term-scalar-sector}{scalar sector} of a
flag-shaped target whose values do not sum to $q_t$ over $\Vset_t$. The sound squashed
constraint is therefore the \emph{expectation} identity
$\tr[\sum_{v\in\Vset_t}F_v\,\sigma]=q_t$ on reachable $\sigma$ --- a
pull-back constraint in the sense of Corollary~\ref{cor:pullback}\,(i) ---
and \emph{not} the operator identity on $\Qsys$. Imposing the latter would
exclude reachable states and void Theorem~\ref{thm:transfer}, which is a
security break rather than a rate loss.

\emph{The two forms, and the certificate of the expectation form.} The
operator form is refused because it is unsound
(Requirement~\ref{req:class-marginal-operator-form}). The expectation form is
sound but is not emitted by default either: the default constraint list is
empty on the flag-state submodule's encoded route and carries an instance's
own entries, at their constants, on its licensed route
(Corollary~\ref{cor:entries}\,(iii)), and adding an entry on the encoded route
would force a licence there and contradict that route's definition. The sound
content of the physical identity is instead carried by the target elements
themselves, which reproduce the statistics exactly
(Lemma~\ref{lem:statistics}\,(ii)). The expectation form remains available on
request \emph{under the condition that discharges its certificate, and not
otherwise}. The entry is $\hat O=\sum_{v\in\Vset_t}F_v$, $\bowtie\,=\,=$,
$c=q_t$ of Theorem~\ref{thm:licence}. Its hypothesis \eqref{eq:certificate}
needs two steps, and only the first is unconditional: by the relation and
linearity, $\Lambda^{\dagger}\big(\sum_{v\in\Vset_t}F_v\big)=\sum_{v\in\Vset_t}M_v$
always (the fourth consequence of Definition~\ref{def:reduction-instance});
that this in turn equals $q_t\idon{\Bsys}$ is \eqref{eq:probeidentity}, which
holds by construction where the model is a joint family and otherwise exactly
where the instance discharges it. \emph{The certificate is therefore
available exactly when the instance's \termref{term-class-marginal-flag}{class-marginal flag} says the
class-marginal identity is discharged}; for a model that is a joint family outright the flag is set
by Lemma~\ref{lem:linear}\,(ii) alone. The constant $q_t$ is declared before
the run in either case (Assumption~\ref{ass:constdecl}). The single equality
entry at $c=q_t$ is what Corollary~\ref{cor:entries}\,(ii) emits, the
certificate being discharged for the emitted numbers themselves. A request
for the entry that meets the condition obtains it, under the licence a
non-empty constraint list requires; a request that does not meet the
condition is refused (Requirement~\ref{req:class-marginal-entry}), because
emitting the entry there would impose $\tr[\sum_{v\in\Vset_t}F_v\,\sigma]=q_t$
on states an uncharacterised branch can reach and for which it is false ---
the security break this remark is about, arrived at from the other direction.
\end{remark}
\claimmeta{\metalabel{Uses}\mref{Lem.}{lem:statistics}, \mref{Cor.}{cor:pullback}, \mref{Thm.}{thm:transfer}, \mref{Lem.}{lem:linear}, \mref{Def.}{def:joint-family}, \mref{Def.}{def:reduction-instance}
\metasep \metalabel{Assumes}\aref{A-CONST}{ass:constdecl}
\metasep \metalabel{Used by}\mref{Req.}{req:class-marginal-operator-form}, \mref{Req.}{req:class-marginal-entry}, \mref{Rem.}{rem:tworoutes}
\metasep newly derived
\metasep accepted}

% claim reductions/rem:channel-choice
\claimsummary{reductions/rem:channel-choice}{Channels realizing one squashing relation can have different reachable sets, so constraints sound for different channels may not be combined. A constraint built from an operator in the carried tuple has the same pull-back for every realizing channel.}
\claimbookmark{reductions/rem:channel-choice}{Remark}{rem:channel-choice}{Soundness relative to the channel, fixed by the carried tuple}
\begin{remark}[Soundness relative to the channel, fixed by the carried tuple]\label{rem:channel-choice}
The relation $M\dsq F$ asserts that \emph{some} squashing channel exists, and
different channels realizing it have different reachable sets: by
Lemma~\ref{lem:statistics}, two channels realizing
$(M;O)\dsq(F;\hat O)$ agree only on
$V\coloneqq\operatorname{span}\big(\{F_v\}\cup\{\hat O_k\}\big)$, and nothing
forces $\Srch_E(\Lambda)=\Srch_E(\Lambda')$. Three consequences follow.
\emph{(a)} Theorem~\ref{thm:transfer} needs only one channel, and every
realizing channel gives a valid conclusion, since $\Lambda$ does not appear in
it. \emph{(b)} There is something to choose between them for
\emph{tightness}: the reachable sets are in general incomparable, so which
channel is used decides which constraints are available. \emph{(c)}
Corollary~\ref{cor:sound}\,(ii)--(iii) therefore require a \emph{common}
channel: constraints sound for $\Lambda$ and constraints sound for $\Lambda'$
may have an intersection excluding every state reachable under either, and
imposing them together would void Theorem~\ref{thm:transfer} exactly as an
unsound constraint does.

What removes the hazard is the tuple of \termref{term-carried-operators}{carried operators} of
Definition~\ref{def:squashing}. If $\hat O^{(n)}$ lies in the span of the
$n$-fold products of elements of $V$, then
$(\Lambda^{\dagger})^{\otimes n}(\hat O^{(n)})
=(\Lambda'^{\dagger})^{\otimes n}(\hat O^{(n)})$, since the adjoints act
factorwise and agree on $V$; the pull-back, and with it every constraint
Corollary~\ref{cor:transported} builds from it, is then the same for every
realizing channel. This is what the tuple is for: listing an operator in the
relation promotes a constraint that would otherwise be tied to one construction
into one that belongs to the relation itself, which is why the emitted
constraints are built from the tuple
(Requirement~\ref{req:carried-tuple-emitted}).
\end{remark}
\claimmeta{\metalabel{Uses}\mref{Def.}{def:squashing}, \mref{Lem.}{lem:statistics}, \mref{Thm.}{thm:transfer}, \mref{Cor.}{cor:sound}, \mref{Cor.}{cor:transported}
\metasep \metalabel{Used by}\mref{Req.}{req:pinched-channel-witness}, \mref{Req.}{req:carried-tuple-emitted}, \mref{Rem.}{rem:tworoutes}
\metasep newly derived
\metasep accepted}

\section{Operations on squashing relations}\label{group-reductions/prop:preorder}
\sectionintro{\termref{term-squashing-relation}{Squashing relations} compose along a pre-order and survive common classical post-processing, a privately drawn setting under one common channel, mixtures and tensor products. Pinching onto the sectors of a block-diagonal target is itself a squashing; where the \termref{term-carried-operators}{carried operators} are block diagonal too, the squashed analysis may be restricted to block-diagonal states for the pinched channel (\hyperlink{claim-reductions/cor:blockdiag}{Corollary~\ref*{cor:blockdiag}}). A private setting with declared \termref{term-class-weights}{class weights} enters as part of a joint outcome, the weight inside the element, and the \termref{term-joint-family}{joint family} is defined once here (\hyperlink{claim-reductions/def:joint-family}{Definition~\ref*{def:joint-family}}). Two necessary conditions can show a declared relation impossible, and both are decided exactly on the declared numbers (\hyperlink{claim-reductions/lem:obstructions}{Lemma~\ref*{lem:obstructions}}, \hyperlink{claim-reductions/rem:obstructiondecide}{Remark~\ref*{rem:obstructiondecide}}).}

% claim reductions/prop:preorder
\claimsummary{reductions/prop:preorder}{Every POVM squashes to itself through the identity channel, and two successive squashing relations compose into one, their carried operator tuples included. A chain of squashers is therefore again a squashing relation that keeps its constraint information.}
\claimbookmark{reductions/prop:preorder}{Proposition}{prop:preorder}{The squashing relation as a pre-order}
\claimleadin{%
This proposition, Lemma~\ref{lem:linear} and Propositions~\ref{prop:pinching}
and~\ref{prop:tensor} are the operations that build
\termref{term-squashing-relation}{squashing relations} out of squashing relations. With the branching rule of Lemma~\ref{lem:master}
they generate every construction this contract emits, so a squasher
submodule assembles its constructions from them rather than proving each
afresh. Where a rule also produces a constraint it is stated for one round,
and Corollary~\ref{cor:transported} lifts every such statement to the
$n$-round soundness of Definition~\ref{def:sound}.
}
\begin{proposition}[The squashing relation as a pre-order]\label{prop:preorder}
$M\dsq_{\idmap}M$, and $M\dsq_{\Lambda_1}F_1$ together with
$F_1\dsq_{\Lambda_2}F_2$ imply $M\dsq_{\Lambda_2\circ\Lambda_1}F_2$. The same
holds for the relation with \termref{term-carried-operators}{carried operators}, the operator tuples
composing the same way: $(M;O)\dsq_{\Lambda_1}(F_1;\hat O^{1})$ and
$(F_1;\hat O^{1})\dsq_{\Lambda_2}(F_2;\hat O^{2})$ give
$(M;O)\dsq_{\Lambda_2\circ\Lambda_1}(F_2;\hat O^{2})$.
\end{proposition}
\claimmeta{\metalabel{Uses}\mref{Lem.}{lem:sep}, \mref{Def.}{def:squashing}
\metasep \metalabel{Used by}\mref{Cor.}{cor:blockdiag}, \mref{Thm.}{thm:chain}, \mref{Ex.}{ex:identity}
\metasep newly derived
\metasep accepted}

\begin{proof}
The identity is a channel and $\idmap^{\dagger}=\idmap$. Compositions of
channels are channels and
$(\Lambda_2\circ\Lambda_1)^{\dagger}=\Lambda_1^{\dagger}\circ
\Lambda_2^{\dagger}$, which is \eqref{eq:pairing} applied twice. Hence
$\Lambda_1^{\dagger}\big(\Lambda_2^{\dagger}(F_{2,v})\big)
=\Lambda_1^{\dagger}(F_{1,v})=M_v$, and identically for each $\hat O_k$.
\end{proof}

% claim reductions/cor:monotone
\claimsummary{reductions/cor:monotone}{If one target squashes to a cleaner one, its worst-case criterion over a family of inputs equals the cleaner target's over the family's image under the connecting channel, and is at most the cleaner target's over any larger family. Moving up the pre-order can only raise the worst case.}
\claimbookmark{reductions/cor:monotone}{Corollary}{cor:monotone}{Monotonicity of the worst-case criterion along the pre-order}
\claimleadin{%
Composing squashers costs nothing in soundness; the corollary says in which
direction it costs tightness. Moving \emph{up} the pre-order, to a cleaner
target, hands the adversary more.
}
\begin{corollary}[Monotonicity of the worst-case criterion along the pre-order]\label{cor:monotone}
For a POVM $G$ and a family $S$ of inputs write
$\varepsilon^{*}(G,S)\coloneqq\sup_E\sup_{\sigma\in S_E}
\Crit(\Eprot_G,\sigma)$. Let $F_2\dsq_{\Lambda_{12}}F_1$ and let both
\termref{term-squashed-protocol}{squashed protocols} use the same processing $\bar{\Grnd}_j$ and $\Kfin$. Then for
every family $S$,
\begin{equation*}
  \varepsilon^{*}(F_2,S)
  =\varepsilon^{*}\big(F_1,(\idmap\otimes\Lambda_{12}^{\otimes n}\otimes\idmap_E)(S)\big)
  \;\le\;\varepsilon^{*}(F_1,S_1)
  \quad\text{for every }S_1\supseteq
  (\idmap\otimes\Lambda_{12}^{\otimes n}\otimes\idmap_E)(S).
\end{equation*}
\end{corollary}
\claimmeta{\metalabel{Uses}\mref{Thm.}{thm:factorization}
\metasep \metalabel{Used by}\mref{Rem.}{rem:flagcost}, \mref{Rem.}{rem:tworoutes}, \mref{Prop.}{prop:weight-monotone}, \mref{Rem.}{rem:weightfree}, \mref{Rem.}{rem:minweight}
\metasep newly derived
\metasep accepted}

\begin{proof}
The equality is Theorem~\ref{thm:factorization} applied with $(F_2,F_1)$ in
place of $(M,F)$, pointwise in $\sigma$ and then in the suprema; the inequality
is monotonicity of a supremum in the set.
\end{proof}

% claim reductions/lem:linear
\claimsummary{reductions/lem:linear}{A squashing relation survives a stochastic post-processing of both sides, and conditional relations with one common channel give relations for the weighted joint family and for its mixture. A private setting therefore enters as part of the outcome and opens no block of the target space.}
\claimbookmark{reductions/lem:linear}{Lemma}{lem:linear}{Linearity: post-processing, private settings and mixtures}
\claimleadin{%
Linearity of $\Lambda^{\dagger}$ lets both sides of a relation be reprocessed
together. It does so in three ways, which differ only in what happens to a
classical variable: it is post-processed, kept as an outcome, or averaged out.
A comma in a subscript denotes a joint outcome and a bar a conditional POVM
(\href{../../../spec/notation.md\#26-comma-vs-bar-on-povm-subscripts}{\nolinkurl{spec/notation.md}\texttt{\#26-comma-vs-bar-on-povm-subscripts}}).
}
\begin{lemma}[Linearity: post-processing, private settings and mixtures]\label{lem:linear}
Let $\Lambda\colon\Tcl{\Bsys}\to\Tcl{\Qsys}$ be a channel.
\begin{enumerate}[label=(\roman*),leftmargin=3.0em]
\item \emph{(post-processing)} If $M\dsq_{\Lambda}F$ and $P(z|v)$ is a
  stochastic matrix from $\Vset$ to a finite alphabet $\mathcal{U}$, then the
  post-processed families $(PM)_z\coloneqq\sum_vP(z|v)M_v$ and
  $(PF)_z\coloneqq\sum_vP(z|v)F_v$ are POVMs and
  $PM\dsq_{\Lambda}PF$.
\item \emph{(a private classical variable)} For each $t\in\Tset$ let
  $\{M_{y|t}\}_{y\in\Yset}$ be a \emph{conditional} POVM on $\Bsys$ ---
  complete for each fixed $t$, $\sum_yM_{y|t}=\idon{\Bsys}$ --- and let
  $\{F_{y|t}\}_{y\in\Yset}$ be conditional POVMs
  on the same $\Qsys$ with $M_{|t}\dsq_{\Lambda}F_{|t}$ for \emph{one common}
  $\Lambda$. Let $(q_t)_{t\in\Tset}$ be a probability distribution. Then, over
  the \emph{joint} alphabet $\Vset=\Yset\times\Tset$,
  \begin{equation}\label{eq:private}
    \big\{q_tM_{y|t}\big\}_{(y,t)}\;\dsq_{\Lambda}\;
    \big\{q_tF_{y|t}\big\}_{(y,t)} ,
  \end{equation}
  which in the comma notation reads
  $\{M_{y,t}\}\dsq_{\Lambda}\{F_{y,t}\}$ with $M_{y,t}=q_tM_{y|t}$ and
  $F_{y,t}=q_tF_{y|t}$.
\item \emph{(mixtures)} In the setting of (ii),
  $\big\{\sum_tq_tM_{y|t}\big\}_y\dsq_{\Lambda}
   \big\{\sum_tq_tF_{y|t}\big\}_y$.
\end{enumerate}
\end{lemma}
\claimmeta{\metalabel{Uses}\mref{Def.}{def:squashing}
\metasep \metalabel{Used by}\mref{Rem.}{rem:unsound}, \mref{Def.}{def:joint-family}, \mref{Rem.}{rem:intertwining}, \mref{Lem.}{lem:master}, \mref{Rem.}{rem:master-scope}, \mref{Thm.}{thm:instrument-instance}, \mref{Def.}{def:weight-route}, \mref{Rem.}{rem:colstoch}, \mref{Prop.}{prop:flat-instance}, \extref{../protocol-compiler/contract.pdf}{protocol-compiler/req:link-row-provenance}
\metasep newly derived
\metasep accepted}

\begin{proof}
(i) $(PM)_z\succeq0$ and $\sum_z(PM)_z=\sum_v\big(\sum_zP(z|v)\big)M_v
=\sum_vM_v=\idon{\Bsys}$, likewise for $PF$; and
$\Lambda^{\dagger}\big((PF)_z\big)=\sum_vP(z|v)\Lambda^{\dagger}(F_v)
=\sum_vP(z|v)M_v=(PM)_z$.
(ii) $\Lambda^{\dagger}\big(q_tF_{y|t}\big)=q_t\Lambda^{\dagger}(F_{y|t})
=q_tM_{y|t}$ by linearity and the relation at each fixed $t$; and both joint
families are POVMs since
$\sum_{y,t}q_tM_{y|t}=\sum_tq_t\big(\sum_yM_{y|t}\big)
=\sum_tq_t\idon{\Bsys}=\idon{\Bsys}$, using completeness of each conditional
POVM, and likewise for $F$.
(iii) is (ii) followed by (i) for the stochastic matrix that forgets $t$.
\end{proof}

% claim reductions/def:joint-family
\claimsummary{reductions/def:joint-family}{A private setting with exact class weights joins the outcome in one joint alphabet, partitioned into one class per setting, and its conditional POVMs form one joint family with each weight inside the elements. The family satisfies the class-marginal identity and opens no target block.}
\claimbookmark{reductions/def:joint-family}{Definition}{def:joint-family}{Joint alphabet, classes and the joint family of a private setting}
\claimleadin{%
A receiver with a private per-round setting, such as a probe setting or an
internal coin of known bias, is described by one conditional POVM per value
of the setting, and the analysis reads the pair (outcome, setting) as one
joint outcome with the setting's weight inside the element
(\href{../../../spec/notation.md\#26-comma-vs-bar-on-povm-subscripts}{\nolinkurl{spec/notation.md}\texttt{\#26-comma-vs-bar-on-povm-subscripts}}). The
definition names that object once and in two layers. The combinatorial layer
is the joint alphabet with its classes and their weights; the abstract
\termref{term-reduction-instance}{reduction instance} of Definition~\ref{def:reduction-instance} is stated over
it, and the branching instance of Theorem~\ref{thm:instrument-instance}
extends it by a branch label. The POVM layer is what a device with a private
setting hands this module. The weights are the measuring party's own
(Assumption~\ref{ass:probe}) and enter only through the statistics of the
joint outcome (Assumption~\ref{ass:tsm}). The identity
\eqref{eq:classmarginal} is a fact about the \emph{physical} model; which
squashed form of it is sound is the subject of Remark~\ref{rem:unsound}, and
why the joint alphabet is built \emph{before} the first squashing link, and
what may still be relabelled after one, is Remark~\ref{rem:intertwining}.
}
\begin{definition}[Joint alphabet, classes and the joint family of a private setting]\label{def:joint-family}
Let $\Yset$ be a finite non-empty ordered
\termdef{term-outcome-alphabet}{outcome alphabet},
let $\Tset$ be a finite non-empty ordered set of
\termdef{term-settings}{settings},
and let $(q_t)_{t\in\Tset}$ be
\termdef{term-class-weights}{class weights},
each a declared dyadic rational on the declared lattice
(Definition~\ref{def:exactadm}\,(b)), with $q_t>0$ and
$\sum_{t\in\Tset}q_t=1$ \emph{exactly}.
\begin{enumerate}[label=(\alph*),leftmargin=2.4em]
\item \emph{(the joint alphabet, its classes and weights)} The
  \termdef{term-joint-alphabet}{joint alphabet}
  is $\Vset\coloneqq\Yset\times\Tset$, ordered with $y$ varying
  fastest, which is its \exttermref{../foundations/contract.pdf}{term-declared-order}{declared order}; its
  \termdef{term-classes}{classes}
  are $\Vset_t\coloneqq\Yset\times\{t\}$, $t\in\Tset$, a partition
  $\Vset=\bigsqcup_{t\in\Tset}\Vset_t$, each class ordered as $\Yset$ is; and
  the triple $\big(\Vset,\{\Vset_t\}_{t\in\Tset},\{q_t\}_{t\in\Tset}\big)$ is
  the \emph{joint alphabet with classes and weights}. When no private setting
  is declared, $\Tset$ is a singleton $\{\ast\}$ with $q_{\ast}=1$, and
  $\Vset=\Yset\times\{\ast\}$ is identified with $\Yset$: one class, of weight
  one.
\item \emph{(the joint family)} Let $\Bsys$ be a system and, for each
  $t\in\Tset$, let $\{\widetilde M_{y|t}\}_{y\in\Yset}$ be a
  \emph{conditional} POVM on $\Bsys$ --- complete at each fixed $t$,
  $\sum_{y}\widetilde M_{y|t}=\idon{\Bsys}$. Their
  \termdef{term-joint-family}{joint family}
  is
  \begin{equation}\label{eq:jointfamily}
    M_{(y,t)}\;\coloneqq\;q_t\,\widetilde M_{y|t},
    \qquad (y,t)\in\Vset ,
  \end{equation}
  formed from the conditional elements and the weights by one product each,
  exactly. It is a POVM on $\Bsys$ over $\Vset$ by
  Lemma~\ref{lem:linear}\,(ii) --- $\sum_{(y,t)}q_t\widetilde M_{y|t}
  =\sum_tq_t\idon{\Bsys}=\idon{\Bsys}$ --- and it satisfies the
  \termdef{term-class-marginal-identity}{class-marginal identity}
  \begin{equation}\label{eq:classmarginal}
    \sum_{v\in\Vset_t}M_v\;=\;q_t\,\idon{\Bsys}
    \qquad\text{for every }t\in\Tset ,
  \end{equation}
  an operator identity on $\Bsys$, by completeness of the conditional POVM at
  $t$.
\end{enumerate}
Two consequences are immediate and are used throughout. First, given
conditional relations $\widetilde M_{|t}\dsq_{\Lambda}\widetilde F_{|t}$,
$t\in\Tset$, with one common channel $\Lambda$, Lemma~\ref{lem:linear}\,(ii)
gives the relation $\{q_t\widetilde M_{y|t}\}\dsq_{\Lambda}\{q_t\widetilde
F_{y|t}\}$ for the joint families: the setting never enters the target space,
no block is opened for it, and the weights are carried by the target elements
themselves. A relation for a joint family need not have a target of this
form. Second, where every $\widetilde M_{y|t}$ is block
diagonal for an orthogonal decomposition of $\Hs{\Bsys}$, so is every
$M_{(y,t)}$, and \eqref{eq:classmarginal} holds block by block; where such a
joint family is itself a target, the decomposition being finite and the
carried operators block diagonal for it as well,
Corollary~\ref{cor:blockdiag} applies to it.
\end{definition}
\claimmeta{\metalabel{Uses}\mref{Lem.}{lem:linear}, \mref{Def.}{def:exactadm}
\metasep \metalabel{Used by}\mref{Rem.}{rem:unsound}, \mref{Rem.}{rem:intertwining}, \mref{Def.}{def:reduction-instance}, \mref{Cor.}{cor:identity-instance}, \mref{Cor.}{cor:flag-instance}, \mref{Thm.}{thm:instrument-instance}, \mref{Def.}{def:emitnorm}, \mref{Thm.}{thm:emit}, \mref{Cor.}{cor:entries}, \mref{Def.}{def:weight-route}, \mref{Prop.}{prop:route}, \mref{Rem.}{rem:colstoch}, \mref{Rem.}{rem:minweight}, \mref{Rem.}{rem:darkcompose}, \mref{Prop.}{prop:flat-instance}
\metasep newly derived
\metasep accepted}

% claim reductions/rem:intertwining
\claimsummary{reductions/rem:intertwining}{A noise or coarse-graining map changes the physical family a squasher is built from, so it acts on the physical POVM before the first link. A stochastic relabelling of both sides by the private setting after a link keeps the channel, the sectors, every pull-back and every certificate.}
\claimbookmark{reductions/rem:intertwining}{Remark}{rem:intertwining}{Classical maps before a squasher, and relabelling after one}
\claimleadin{%
Lemma~\ref{lem:linear} post-processes both sides of a relation at once
(part (i)) and folds a private variable into the outcome (part (ii)). This
remark says where in a chain of squashers each of the two may be used.
}
\begin{remark}[Classical maps before a squasher, and relabelling after one]\label{rem:intertwining}
\emph{A private variable inside the alphabet.} In
Lemma~\ref{lem:linear}\,(ii) the private variable never enters $\Qsys$: the
target space is unchanged, no block is opened, and no constraint is produced,
the weights $q_t$ being carried by the target effects themselves. A variable
the measuring party draws privately, such as a probe setting or an internal coin of known
bias, is therefore squashed into the outcome alphabet at no cost \emph{by (ii)
alone}. That is the treatment the output-state test-statistics model of
Assumption~\ref{ass:tsm} requires, and the reason
Definition~\ref{def:joint-family} builds the joint alphabet before squashing
rather than after. The no-cost statement concerns (ii) alone and need not
survive a squasher composed with it: once (ii) is composed with a flag-shaped
squasher, every value of the private variable opens flag sectors of its own.
A variable the squasher's channel does not depend on, such as a test-round designation drawn against the setting, is
then carried more cheaply \emph{after} the link, by (i), than inside the
alphabet before it (the relabelling \eqref{eq:relabel}); only a
variable the channel depends on, as a probe setting does through the
conditional families, has to be in the alphabet before the first link. The
branching rule of Lemma~\ref{lem:master}, by contrast, files its classical
variable in $\Qsys$ and does cost something (Remark~\ref{rem:master-scope}).

Part (i) applies a post-processing to \emph{both} sides at once, and two uses
of it differ: one is refused across a link and the other is not.

\emph{A noise or coarse-graining map changes the family a squasher is built
from.} A \exttermref{../receiver/contract.pdf}{term-dark-count-channel}{dark-count channel} or a coarse-graining applied to the
physical POVM produces a different physical family, and a squasher
constructed for the original one is not a squasher to the same target for it:
$M\dsq_{\Lambda}F$ gives $\Lambda^{\dagger}(F_v)=M_v$, which is not $(PM)_v$
in general. That is why such a map is applied to the physical POVM
\emph{before} the first squashing link (in the \termref{term-weight-route}{weight route} of the
flag-state submodule, the descended dark-count matrix acts before the joint
family of Definition~\ref{def:joint-family} is formed). Part (i) does move a
stochastic $P$ across a fixed squasher, but only together with the target:
from $M\dsq_{\Lambda}F$ it gives $PM\dsq_{\Lambda}PF$, with the same channel
and the target $PF$ in place of $F$, whose elements, and with them its
\termref{term-scalar-value}{scalar values}, its class sums and the entries built from
them, are not the ones the chain was built and certified for. Where the map
acts on a finer alphabet than the squasher's --- a dark-count channel $P$ on
raw click patterns, applied before the \exttermref{../receiver/contract.pdf}{term-declared-coarse-graining}{declared coarse-graining} $K$
that produces the
squasher's family $KM_{\mathrm{raw}}$ from the raw click-pattern POVM
$M_{\mathrm{raw}}$ --- moving it across the squasher asks for a stochastic
$P'$ on the coarse alphabet with $P'KM_{\mathrm{raw}}=KPM_{\mathrm{raw}}$, a
linear feasibility problem in $P'$ with no solution in general. For two
detectors whose single-click patterns $K$ merges into one letter, with
dark-count probabilities $d_1\ne d_2$ acting independently, the double-click
element of $KPM_{\mathrm{raw}}$ carries the element of the pattern in which
only detector~$1$ clicked with the coefficient $d_2$ and that of the pattern
in which only detector~$2$ clicked with the coefficient $d_1$, while every
element of $P'KM_{\mathrm{raw}}$ carries the two with one common
coefficient; so no $P'$ exists when the four elements of $M_{\mathrm{raw}}$
are linearly independent. This contract does neither: the order is fixed
(Requirement~\ref{req:post-processing-order}).

\emph{A stochastic relabelling applied to both sides after a link is part (i)
itself, and it is sound.} Let $M\dsq_{\Lambda}F$ be a link's relation over
the joint alphabet $\Vset$ of Definition~\ref{def:joint-family}, and let
$P(f\,|\,t)$ be a stochastic matrix from the settings $\Tset$ to a finite
alphabet of values $f$, depending on the private component $t=t(v)$ of the
outcome only and never on the device outcome $y$, since otherwise $f$ would
be inferred from an outcome. Then, over the support of $P$,
\begin{equation}\label{eq:relabel}
  \big\{P(f\,|\,t(v))\,M_v\big\}_{(v,f)}\ \dsq_{\Lambda}\
  \big\{P(f\,|\,t(v))\,F_v\big\}_{(v,f)}
\end{equation}
by Lemma~\ref{lem:linear}\,(i) for the stochastic matrix from $\Vset$ to
$\Vset\times\{f\}$ with entries $\delta_{vv'}P(f\,|\,t(v'))$ --- every column
sums to $\sum_fP(f|t(v'))=1$ --- with the \emph{same} channel. Nothing of the
link changes: $\Qsys$, its \exttermref{../foundations/contract.pdf}{term-sector-map}{sector map} and its dimension are unchanged;
every element $P(f|t(v))F_v$ is a scalar multiple of an existing one, so
block diagonality holds by construction; every carried operator has the
same pull-back; and every entry's certificate is untouched. On a
\termref{term-scalar-sector}{scalar sector} the value becomes $P(f|t(v))\phi_s(v)$, and no sector is indexed by
$f$. Such a relabelling is therefore not a link in the sense of
Definition~\ref{def:link}: a link composes a new channel, while a relabelling
changes the alphabet at both ends of an existing relation. This contract
declares no kind, no register and no position in a chain
(Definition~\ref{def:chain}) for it: it is the \emph{consumer's} to apply,
after the chain, by the same lemma, which the module that compiles the
protocol may use directly.
\end{remark}
\claimmeta{\metalabel{Uses}\mref{Lem.}{lem:linear}, \mref{Def.}{def:joint-family}
\metasep \metalabel{Assumes}\aref{A-TSM}{ass:tsm}
\metasep \metalabel{Used by}\mref{Req.}{req:post-processing-order}
\metasep newly derived
\metasep accepted}

% claim reductions/prop:pinching
\claimsummary{reductions/prop:pinching}{Pinching onto finitely many orthogonal sectors is a self-adjoint channel, and a target whose elements are block diagonal for those sectors squashes to itself through it. The pinching can therefore follow a squasher whose target and carried operators are block diagonal.}
\claimbookmark{reductions/prop:pinching}{Proposition}{prop:pinching}{Sector pinching as a squashing}
\claimleadin{%
The sector pinching below is the pinching of
\xref{foundations/def:pinching} on the target $\Qsys$: it deletes the
coherences between sectors. The proposition proves the two facts the
squashing constructions use, that it is a self-adjoint channel and that a
block-diagonal target squashes to itself through it.
}
\begin{proposition}[Sector pinching as a squashing]\label{prop:pinching}
Let $\{\Pi_s\}_{s\in\SecSet}$ be finitely many mutually orthogonal projections
on $\Hs{\Qsys}$ with $\sum_s\Pi_s=\idon{\Qsys}$, and let
$\PinS(\sigma)\coloneqq\sum_s\Pi_s\sigma\Pi_s$. Then $\PinS$ is a channel with
$\PinS^{\dagger}=\PinS$, and if every $F_v$ is block diagonal,
$F_v=\sum_s\Pi_sF_v\Pi_s$, then $F\dsq_{\PinS}F$.
\end{proposition}
\claimmeta{\metalabel{Uses}\mref{Def.}{def:squashing}
\metasep \metalabel{Used by}\mref{Cor.}{cor:blockdiag}, \mref{Ex.}{ex:identity}, \mref{Def.}{def:reduction-instance}, \mref{Cor.}{cor:identity-instance}, \mref{Cor.}{cor:flag-instance}, \mref{Lem.}{lem:precomposed-instance}, \mref{Thm.}{thm:emitted-product}, \mref{Cor.}{cor:emitted-identity}
\metasep newly derived
\metasep accepted}

\begin{proof}
$\PinS$ is CP as a sum of conjugations and trace preserving by cyclicity and
$\sum_s\Pi_s=\idon{\Qsys}$; it is self-adjoint because
$\tr[X\sum_s\Pi_s\sigma\Pi_s]=\tr[(\sum_s\Pi_sX\Pi_s)\sigma]$. Block
diagonality is $\PinS(F_v)=F_v$, which is
$\PinS^{\dagger}(F_v)=F_v$.
\end{proof}

% claim reductions/cor:blockdiag
\claimsummary{reductions/cor:blockdiag}{If every target element and carried operator is block diagonal for a sector decomposition, the squashing channel followed by the pinching realises the same relation and makes every reachable state block diagonal. For that channel the squashed analysis may be restricted to block-diagonal states.}
\claimbookmark{reductions/cor:blockdiag}{Corollary}{cor:blockdiag}{Restriction to block-diagonal squashed states}
\claimleadin{%
Composing a squasher with the pinching of Proposition~\ref{prop:pinching}
costs nothing when the target and the carried operators are block diagonal,
and it makes block diagonality of the squashed states a sound constraint.
}
\begin{corollary}[Restriction to block-diagonal squashed states]\label{cor:blockdiag}
Let $(M;O)\dsq_{\Lambda}(F;\hat O)$ and suppose every $F_v$ and every
$\hat O_k$ is block diagonal for a finite orthogonal decomposition
$\{\Pi_s\}_{s\in\SecSet}$ of $\Hs{\Qsys}$. Then
\begin{enumerate}[label=(\roman*),leftmargin=3.0em]
\item $(M;O)\dsq_{\PinS\circ\Lambda}(F;\hat O)$, with the same target, the same
  tuple and the same pull-backs;
\item every state reachable through $\PinS\circ\Lambda$ is block diagonal in
  each round, so ``$\sigma$ is block diagonal in every round's sector
  decomposition'' is a sound constraint for that channel
  (Definition~\ref{def:sound});
\item consequently the squashed analysis may be restricted to block-diagonal
  states without excluding a reachable one.
\end{enumerate}
\end{corollary}
\claimmeta{\metalabel{Uses}\mref{Def.}{def:sound}, \mref{Cor.}{cor:sound}, \mref{Prop.}{prop:preorder}, \mref{Prop.}{prop:pinching}
\metasep \metalabel{Used by}\mref{Thm.}{thm:mixture}, \mref{Def.}{def:reduction-instance}, \mref{Cor.}{cor:identity-instance}, \mref{Thm.}{thm:instrument-instance}, \mref{Thm.}{thm:emit}, \mref{Thm.}{thm:emitted-product}, \mref{Cor.}{cor:reference-companion}, \mref{Req.}{req:pinched-channel-witness}, \mref{Prop.}{prop:flat-instance}
\metasep newly derived
\metasep accepted}

\begin{proof}
(i) By Proposition~\ref{prop:pinching}, $F\dsq_{\PinS}F$ carrying $\hat O$ as
well, since $\PinS^{\dagger}(\hat O_k)=\hat O_k$; compose with
Proposition~\ref{prop:preorder}. (ii) The output of $\PinS$ is block diagonal
by construction, and $(\PinS\circ\Lambda)^{\otimes n}$ is block diagonal in
every factor. (iii) is (ii) with Corollary~\ref{cor:sound}\,(i).
\end{proof}

% claim reductions/prop:tensor
\claimsummary{reductions/prop:tensor}{Two squashing relations give one for the product measurement, through the tensor product of their channels. Independent components can therefore be reduced separately and assembled without a fresh existence proof.}
\claimbookmark{reductions/prop:tensor}{Proposition}{prop:tensor}{Tensor products of squashing relations}
\claimleadin{%
Squashing relations for independent components combine by tensor product,
because the adjoint of a tensor product of channels is the tensor product of
their adjoints, also when the physical systems are infinite-dimensional
(Lemma~\ref{lem:sep}\,(ii)).
}
\begin{proposition}[Tensor products of squashing relations]\label{prop:tensor}
If $M\dsq_{\Lambda}F$ and $M'\dsq_{\Lambda'}F'$ then
$\{M_v\otimes M'_{v'}\}_{(v,v')}\dsq_{\Lambda\otimes\Lambda'}
\{F_v\otimes F'_{v'}\}_{(v,v')}$.
\end{proposition}
\claimmeta{\metalabel{Uses}\mref{Lem.}{lem:sep}, \mref{Def.}{def:squashing}
\metasep \metalabel{Used by}\mref{Thm.}{thm:licence}, \mref{Thm.}{thm:emitted-product}
\metasep newly derived
\metasep accepted}

\begin{proof}
$(\Lambda\otimes\Lambda')^{\dagger}(F_v\otimes F'_{v'})
=\Lambda^{\dagger}(F_v)\otimes\Lambda'^{\dagger}(F'_{v'})$ by
Lemma~\ref{lem:sep}\,(ii).
\end{proof}

% claim reductions/lem:obstructions
\claimsummary{reductions/lem:obstructions}{If M squashes to F, the spectrum of each real combination of physical elements lies between the extreme spectral values of the same combination of target elements, and the physical span is no larger and inherits every linear relation with the identity. A relation violating either is impossible.}
\claimbookmark{reductions/lem:obstructions}{Lemma}{lem:obstructions}{Necessary conditions for a squashing relation}
\claimleadin{%
Not every pair $(M,F)$ admits a \termref{term-squashing-channel}{squashing channel}, and two necessary
conditions are cheap enough to check by inspection. They serve to
\emph{refuse} a declared relation, never to establish one;
Remark~\ref{rem:obstructiondecide} says how both are decided on declared
numbers.
}
\begin{lemma}[Necessary conditions for a squashing relation]\label{lem:obstructions}
Let $M\dsq F$, let $X=\sum_{v}a_{v}F_{v}$ be a real linear combination of
target effects and let $X'\coloneqq\sum_{v}a_{v}M_{v}$. Then:
\begin{enumerate}[label=(\roman*),leftmargin=3.0em]
\item \emph{(spectral inclusion)}
  $\operatorname{spec}(X')\subseteq[\inf\operatorname{spec}X,\
  \sup\operatorname{spec}X]$; in particular
  $\operatorname{spec}(M_v)\subseteq[\inf\operatorname{spec}F_v,\
  \sup\operatorname{spec}F_v]$ and $\|M_v\|\le\|F_v\|$ for every $v$;
\item \emph{(span monotonicity)}
  $\dim\operatorname{span}\{M_v\}\le\dim\operatorname{span}\{F_v\}
  \le(\dim\Hs{\Qsys})^{2}$, and every linear relation
  $\sum_va_vF_v=c\,\idon{\Qsys}$ forces $\sum_va_vM_v=c\,\idon{\Bsys}$.
\end{enumerate}
\end{lemma}
\claimmeta{\metalabel{Uses}\mref{Def.}{def:squashing}
\metasep \metalabel{Used by}\mref{Req.}{req:declared-relation-check}, \mref{Prop.}{prop:chain}
\metasep newly derived
\metasep accepted}

\begin{proof}
Let $\Lambda$ realize the relation, so $X'=\Lambda^{\dagger}(X)$ by linearity.
(i) With $a\coloneqq\inf\operatorname{spec}X$ and
$b\coloneqq\sup\operatorname{spec}X$ we have
$a\,\idon{\Qsys}\preceq X\preceq b\,\idon{\Qsys}$; applying the positive unital
map $\Lambda^{\dagger}$ preserves both inequalities, so
$a\,\idon{\Bsys}\preceq X'\preceq b\,\idon{\Bsys}$. Taking $X=F_v$ gives the
inclusion for $M_v$, and the norm bound follows since
$\|Y\|=\sup\operatorname{spec}Y$ for a positive operator.
(ii) $M_v=\Lambda^{\dagger}(F_v)$, so
$\operatorname{span}\{M_v\}=\Lambda^{\dagger}\big(\operatorname{span}\{F_v\}\big)$
and the dimension cannot increase; the bound
$(\dim\Hs{\Qsys})^{2}$ is $\dim\Lin{\Qsys}$. Unitality gives the last
statement.
\end{proof}

% claim reductions/rem:obstructiondecide
\claimsummary{reductions/rem:obstructiondecide}{On declared numbers the spectral condition is decided exactly in operator form on rational enclosures, and the span condition by exact ranks and linear relations. A loose decision would refuse more or less than the lemma says, which is safe: nothing emitted rests on the check having passed.}
\claimbookmark{reductions/rem:obstructiondecide}{Remark}{rem:obstructiondecide}{Exact decision of the necessary conditions}
\begin{remark}[Exact decision of the necessary conditions]\label{rem:obstructiondecide}
Lemma~\ref{lem:obstructions} is an exact statement about an arbitrary
relation. Applied to \emph{declared} numbers, both of its parts are decided
\emph{exactly} in the substrate of Lemma~\ref{lem:exactpsd}, each in a form
that differs from the lemma's wording.
\emph{(i)} The spectral enclosure is decided in its \emph{operator} form,
which is what a refusal needs and what \eqref{eq:psdtest} decides: for
rationals $a\le b$ the pair $a\idon{\Qsys}\preceq X\preceq b\idon{\Qsys}$ is
decided exactly, by any of the exact procedures of
Remark~\ref{rem:psdalgorithm}, and where it holds the same pair is decided
for $X'$. \emph{The lemma's $\inf\operatorname{spec}X$ and
$\sup\operatorname{spec}X$ are in general algebraic and not rational}, so the
decision tests the rational enclosures the declaration exhibits rather than
the extreme values themselves; it refuses a subset of what the lemma forbids,
and no eigenvalue is computed anywhere.
\emph{(ii)} The span dimension is the \emph{exact} rank over $\Fldi$ of the
matrix whose columns are the vectorized declared elements, decided by
elimination in the field. A floating-point rank is a threshold decision, and
this one is not. The linear relations are decided the same way: the
solutions $(a,c)$ of $\sum_va_vF_v=c\,\idon{\Qsys}$ form a subspace, an exact
basis of which elimination in the field produces, and testing
$\sum_va_vM_v=c\,\idon{\Bsys}$ on each basis solution tests it on every
solution, the condition being linear in $(a,c)$.

\emph{A loose decision would have been safe here.} A rank read at a threshold
may count a nearly dependent family as independent, and an enclosure read at
a window may fail to notice a violation just inside it. Made on the target's
elements, either error refuses \emph{less} than the lemma forbids; made on the
physical elements, a loose decision can refuse a relation that exists.
Neither can produce a wrong artefact, because of what the check is for: it
is a cheap early refusal, not a warrant. Nothing this module emits is
proved from the check's having passed; what the emitted artefact rests on are
the constructions themselves, which either produce a channel or do not. A
missed refusal costs the construction's time and produces no artefact, and a
spurious one produces none either; neither can produce a wrong one. The check is exact all the same
(Requirement~\ref{req:declared-relation-check}). \emph{It is also vacuous on
this module's own emission path}, where the request declares no relation of
its own and the only relation in play is the one the reduction instance
constructs (Definition~\ref{def:reduction-instance}\,(I3)).
\end{remark}
\claimmeta{\metalabel{Used by}\mref{Req.}{req:declared-relation-check}
\metasep newly derived
\metasep accepted}

\section{Branching and flag targets}\label{group-reductions/lem:top}
\sectionintro{Every measurement with a finite alphabet squashes to the flag PVM on its outcomes (\hyperlink{claim-reductions/lem:top}{Lemma~\ref*{lem:top}}), so the existence of a squasher is never the difficulty; its tightness is. A \termref{term-generalized-instrument}{generalized instrument} followed by a squasher in each branch gives one \termref{term-squashing-relation}{squashing relation} whose target stacks the branch targets as blocks, with each branch weight transported as an operator (\hyperlink{claim-reductions/lem:master}{Lemma~\ref*{lem:master}}). Its instances cover a known mixture, a measurement block diagonal for an orthogonal decomposition and a register read by a POVM (\hyperlink{claim-reductions/cor:instances}{Corollary~\ref*{cor:instances}}, \hyperlink{claim-reductions/thm:mixture}{Theorem~\ref*{thm:mixture}}); common targets that merge the blocks again, flag-extended targets and their rate cost, and the price of a general POVM on the register follow.}

% claim reductions/lem:top
\claimsummary{reductions/lem:top}{Every POVM with a finite alphabet squashes to the flag PVM on its outcomes through a measure-and-prepare channel, and to any PVM on the same alphabet that vanishes only where the POVM does. A squasher therefore always exists; the work lies in finding a target tight enough to leave a key rate.}
\claimbookmark{reductions/lem:top}{Lemma}{lem:top}{Flag PVM as a top element}
\claimleadin{%
A receiver is rarely characterized as a single POVM on a single space. It is
characterized \emph{conditionally}: below a photon-number cutoff we know what
it does and above it we do not; on one value of a classical branch register it
behaves as specified and on another we promise nothing. The lemma shows first
that existence of a squasher is never the difficulty. The work of this
module's constructions is therefore not to find a squasher but to find one
whose target is tight enough to leave a key rate.
}
\begin{lemma}[Flag PVM as a top element]\label{lem:top}
Let $M$ be any POVM on any $\Bsys$ with finite alphabet $\Vset$, and let
$\Ppvm=\{\dyad{v}\}_{v\in\Vset}$ be the \emph{flag PVM} on
$\Hs{\Flag}=\mathbb{C}^{|\Vset|}$ with orthonormal basis
$\{\ket{v}\}_{v\in\Vset}$. Then
\begin{equation}\label{eq:mpsquasher}
  M\;\dsq_{\Lmp}\;\Ppvm,
  \qquad
  \Lmp(\rho)\coloneqq\sum_{v\in\Vset}\tr[M_v\rho]\,\dyad{v} .
\end{equation}
More generally, if $P$ is any PVM on a space $\Qsys$ with the same alphabet and
$P_v=0\Rightarrow M_v=0$, then $M\dsq P$.
\end{lemma}
\claimmeta{\metalabel{Uses}\mref{Def.}{def:squashing}
\metasep \metalabel{Used by}\mref{Cor.}{cor:instances}, \mref{Cor.}{cor:flagsenough}, \mref{Rem.}{rem:flagcost}, \mref{Cor.}{cor:flag-instance}, \mref{Prop.}{prop:chain}, \mref{Thm.}{thm:noise-channel}
\metasep newly derived
\metasep accepted}

\begin{proof}
$\Lmp$ is a measure-and-prepare channel: CP, and trace preserving because
$\sum_v\tr[M_v\rho]=\tr\rho$. Its adjoint is
$\Lmp^{\dagger}(X)=\sum_v\bra{v}X\ket{v}\,M_v$, since
$\tr[X\Lmp(\rho)]=\sum_v\tr[M_v\rho]\bra{v}X\ket{v}$; evaluating at
$X=\dyad{v'}$ gives $M_{v'}$. For the general case pick, for each $v$ with
$M_v\neq0$, a state $\omega_v$ with $P_v\omega_vP_v=\omega_v$ (possible since
$P_v\neq0$), and set
$\Lambda(\rho)\coloneqq\sum_{v:M_v\neq0}\tr[M_v\rho]\,\omega_v$. This is trace
preserving because the omitted terms vanish and $\sum_vM_v=\idon{\Bsys}$, and
$\Lambda^{\dagger}(P_{v'})=\sum_{v:M_v\neq0}\tr[P_{v'}\omega_v]\,M_v=M_{v'}$,
because $\tr[P_{v'}\omega_v]=\delta_{vv'}$ by orthogonality of the ranges of a
PVM, and because the identity reads $0=0$ when $M_{v'}=0$.
\end{proof}

% claim reductions/def:generalized-instrument
\claimsummary{reductions/def:generalized-instrument}{A generalized instrument is a finite family of completely positive maps on one input system whose traces add up to the trace. Each map is a branch, occurring with the probability its trace gives and leaving a subnormalized state on a system that may differ from branch to branch.}
\claimbookmark{reductions/def:generalized-instrument}{Definition}{def:generalized-instrument}{Generalized instrument}
\claimleadin{%
By Lemma~\ref{lem:top} every POVM squashes to a flag PVM, so a
\termref{term-squashing-relation}{squashing relation} on its own is worth nothing: the target $\Ppvm$
lets the adversary prepare a flag state and fix Bob's outcome outright. What
buys a key rate is a target whose elements are not all rank-one projectors, or
a \termref{term-sound}{sound} constraint that excludes the flag states an adversary would
use; the squasher constructions of this module's submodules and the
attack-set licence of
Theorem~\ref{thm:licence} are the two ways of getting one. The squasher
constructions describe a receiver branch by branch, and the branching step is
the following object.
}
\begin{definition}[Generalized instrument]\label{def:generalized-instrument}
A \termdef{term-generalized-instrument}{generalized instrument} on $\Bsys$ is a finite family
$\{\Ins_j\}_{j\in J}$ of CP maps $\Ins_j\colon\Tcl{\Bsys}\to\Tcl{\Bsys'_j}$
whose traces add up to the trace, $\sum_j\tr\circ\Ins_j=\tr$: branch $j$
occurs with probability $\tr\Ins_j(\rho)$ and leaves the subnormalized state
$\Ins_j(\rho)$ on the system $\Bsys'_j$, which may differ from branch to
branch.
\end{definition}
\claimmeta{\metalabel{Used by}\mref{Lem.}{lem:master}, \mref{Thm.}{thm:mixture}, \mref{Thm.}{thm:instrument-instance}, \mref{Cor.}{cor:register-instance}, \mref{Thm.}{thm:noise-channel}
\metasep newly derived
\metasep accepted}

% claim reductions/lem:master
\claimsummary{reductions/lem:master}{A generalized instrument followed by a squasher in each branch squashes to the direct sum of the branch targets, with or without the branch label in the outcome. Branch operators are transported, and a branch weight is a fixed number only where its pull-back is a multiple of the identity.}
\claimbookmark{reductions/lem:master}{Lemma}{lem:master}{Instrument-conditioned squashing}
\claimleadin{%
A receiver described branch by branch squashes as a whole once each branch
does. The lemma combines a generalized instrument with one squasher per branch
into a single squashing relation whose target stacks the branch targets as
diagonal blocks, and carries each branch weight along as an operator.
}
\begin{lemma}[Instrument-conditioned squashing]\label{lem:master}
Let $J$ be a finite set and let $\Ins_j\colon\Tcl{\Bsys}\to\Tcl{\Bsys'_j}$ be
CP maps with $\sum_{j}\tr\circ\Ins_j=\tr$, that is, a generalized instrument
(Definition~\ref{def:generalized-instrument}). For each $j$ let $M^{(j)}$ be a
POVM on $\Bsys'_j$ over a common finite alphabet $\Vset$ with
$(M^{(j)};O^{(j)})\dsq_{\Lambda_j}(F^{(j)};\hat O^{(j)})$, where $F^{(j)}$ is a
POVM on a finite-dimensional $\Qsys_j$. Each $\Ins_j$ is trace
non-increasing; write $\Ins_j^{\dagger}$ for its adjoint in the sense of
Lemma~\ref{lem:sep}, the completely positive map fixed by
$\tr[X\Ins_j(\rho)]=\tr[\Ins_j^{\dagger}(X)\rho]$ for $X\in\Lin{\Bsys'_j}$ and
$\rho\in\Tcl{\Bsys}$. Then:
\begin{enumerate}[label=(\roman*),leftmargin=3.0em]
\item $M_{v,j}\coloneqq\Ins_j^{\dagger}\big(M^{(j)}_v\big)$ defines a POVM on
  $\Bsys$ with joint outcome $(v,j)\in\Vset\times J$;
\item with $\Qsys\coloneqq\bigoplus_{j\in J}\Qsys_j$ and $F_{v,j}$ the operator
  equal to $F^{(j)}_v$ in the $j$-th block and zero elsewhere,
  \begin{equation}\label{eq:master}
    \{M_{v,j}\}\;\dsq_{\Lambda}\;\{F_{v,j}\},
    \qquad
    \Lambda(\rho)\coloneqq\bigoplus_{j\in J}\Lambda_j\big(\Ins_j(\rho)\big);
  \end{equation}
\item \emph{(coarse graining)} with the same channel,
  $\big\{\sum_jM_{v,j}\big\}_v\dsq_{\Lambda}
   \big\{\bigoplus_jF^{(j)}_v\big\}_v$;
\item \emph{(branch weights)} $\Lambda^{\dagger}(\Pi_j)
  =\Ins_j^{\dagger}(\idon{\Bsys'_j})$, where $\Pi_j$ projects onto the $j$-th
  block. In particular, if $\Ins_j^{\dagger}(\idon{\Bsys'_j})=q_j\,\idon{\Bsys}$
  for constants $q_j$, then for every finite-dimensional side system $R$ and
  every \termref{term-reachable}{reachable} $\sigma$,
  $\tr_{\Qsys}[(\Pi_j\otimes\idon{R})\sigma]=q_j\,\sigma_R$, and in particular
  $\tr[\Pi_j\sigma]=q_j$;
\item \emph{(transported operators)} the composite carries the operator equal
  to $\hat O^{(j)}_k$ in block $j$ and zero elsewhere, with pull-back
  $\Ins_j^{\dagger}\big(O^{(j)}_k\big)$.
\end{enumerate}
\end{lemma}
\claimmeta{\metalabel{Uses}\mref{Lem.}{lem:sep}, \mref{Cor.}{cor:pullback}, \mref{Lem.}{lem:linear}, \mref{Def.}{def:generalized-instrument}
\metasep \metalabel{Used by}\mref{Rem.}{rem:reading-master}, \mref{Rem.}{rem:master-scope}, \mref{Rem.}{rem:branch-blocks}, \mref{Rem.}{rem:branch-probability}, \mref{Cor.}{cor:instances}, \mref{Prop.}{prop:common}, \mref{Thm.}{thm:mixture}, \mref{Thm.}{thm:instrument-instance}, \mref{Cor.}{cor:register-instance}, \mref{Prop.}{prop:chain}, \mref{Thm.}{thm:noise-channel}
\metasep newly derived
\metasep accepted}

\begin{proof}
For $\rho\succeq0$ every term of $\sum_{j'}\tr\Ins_{j'}(\rho)=\tr\rho$ is
non-negative, so $\tr\Ins_j(\rho)\le\tr\rho$: each $\Ins_j$ is trace
non-increasing, and Lemma~\ref{lem:sep} gives its adjoint $\Ins_j^{\dagger}$,
completely positive and so positive. The hypothesis
$\sum_j\tr\circ\Ins_j=\tr$ says exactly
$\sum_j\Ins_j^{\dagger}(\idon{\Bsys'_j})=\idon{\Bsys}$, by
Lemma~\ref{lem:sep}\,(i) applied to
$\sum_j\Ins_j^{\dagger}(\idon{})-\idon{\Bsys}$.

(i) $M_{v,j}\succeq0$, and
$\sum_{v,j}\Ins_j^{\dagger}(M^{(j)}_v)
=\sum_j\Ins_j^{\dagger}(\idon{\Bsys'_j})=\idon{\Bsys}$.

(ii) $\Lambda$ is CP as a direct sum of compositions of CP maps, and trace
preserving: $\tr\Lambda(\rho)=\sum_j\tr\Lambda_j(\Ins_j(\rho))
=\sum_j\tr\Ins_j(\rho)=\tr\rho$. Its output is block diagonal, so only the
diagonal blocks of an operator on $\Qsys$ contribute to the pairing
\eqref{eq:pairing}; for $X=\bigoplus_jX_j$,
\begin{equation}\label{eq:masteradjoint}
  \tr\big[X\,\Lambda(\rho)\big]
  =\sum_j\tr\big[X_j\,\Lambda_j(\Ins_j(\rho))\big]
  =\sum_j\tr\Big[\Ins_j^{\dagger}\big(\Lambda_j^{\dagger}(X_j)\big)\,\rho\Big],
  \quad\text{so}\quad
  \Lambda^{\dagger}\Big(\bigoplus_jX_j\Big)
  =\sum_j\Ins_j^{\dagger}\big(\Lambda_j^{\dagger}(X_j)\big).
\end{equation}
Taking $X=F_{v,j}$, with only the $j$-th block non-zero, gives
$\Lambda^{\dagger}(F_{v,j})=\Ins_j^{\dagger}\big(\Lambda_j^{\dagger}
(F^{(j)}_v)\big)=\Ins_j^{\dagger}(M^{(j)}_v)=M_{v,j}$. That
$\{F_{v,j}\}$ is a POVM on $\Qsys$ follows blockwise from
$\sum_vF^{(j)}_v=\idon{\Qsys_j}$.

(iii) is Lemma~\ref{lem:linear}\,(i) for the stochastic map that forgets $j$.

(iv) Take $X=\Pi_j$, that is $X_j=\idon{\Qsys_j}$ and zero elsewhere, in
\eqref{eq:masteradjoint}: $\Lambda^{\dagger}(\Pi_j)
=\Ins_j^{\dagger}\big(\Lambda_j^{\dagger}(\idon{\Qsys_j})\big)
=\Ins_j^{\dagger}(\idon{\Bsys'_j})$ by unitality of $\Lambda_j^{\dagger}$. The
constraint claim is Corollary~\ref{cor:pullback}\,(i).

(v) is \eqref{eq:masteradjoint} evaluated on the stated operator.
\end{proof}

% claim reductions/rem:reading-master
\claimsummary{reductions/rem:reading-master}{The branch weights of \hyperlink{claim-reductions/lem:master}{Lemma~\ref*{lem:master}} (iv) are the only constraint the branching rule produces, and the operator form is the one to impose: it fixes each block's weight conditioned on everything Alice and Eve hold, while its trace alone admits a block index perfectly correlated with Eve.}
\claimbookmark{reductions/rem:reading-master}{Remark}{rem:reading-master}{The operator form of the branch-weight constraint}
\begin{remark}[The operator form of the branch-weight constraint]
\label{rem:reading-master}
Part~(iv) of Lemma~\ref{lem:master} is the only place the branching rule
produces a constraint, and the operator form is the one worth imposing: it
says the weight of block $j$ is $q_j$ \emph{conditioned on everything Alice
and Eve hold}, and Corollary~\ref{cor:transported}\,(i) turns it into a sound
constraint, one per round. The scalar $\tr[\Pi_j\sigma]=q_j$ is only its trace
and is strictly weaker: it still admits a squashed state whose block index is
perfectly correlated with Eve's register. Nothing is averaged anywhere in the
rule, because averaging over $j$ before squashing would lose the branch
structure and with it the relation.
\end{remark}
\claimmeta{\metalabel{Uses}\mref{Cor.}{cor:transported}, \mref{Lem.}{lem:master}
\metasep newly derived
\metasep accepted}

% claim reductions/rem:master-scope
\claimsummary{reductions/rem:master-scope}{The branching rule files the branch label as a block index of the system handed to the adversary, so it models branches Eve may know. A setting the measuring party draws privately belongs to the common-channel rule of \hyperlink{claim-reductions/lem:linear}{Lemma~\ref*{lem:linear}} (ii), which opens no block and produces no constraint.}
\claimbookmark{reductions/rem:master-scope}{Remark}{rem:master-scope}{Branch labels as information the adversary may hold}
\begin{remark}[Branch labels as information the adversary may hold]\label{rem:master-scope}
The branch label ends up as a block index of $\Qsys$, and $\Qsys$ is the system
the squashed model hands to the adversary. Lemma~\ref{lem:master} therefore
models a branch Eve may as well know: a detector-response mode she selected, a
test on the light she sent, a register that arrived with it. A setting the measuring
party draws \emph{privately} can be pushed through the rule, taking
$\Ins_j=q_j\idmap$, but the resulting model files that setting in Eve's state,
and only part~(iv) of Lemma~\ref{lem:master} in its operator form ties it
back. Settings meant to stay on the measuring party's side belong to
Lemma~\ref{lem:linear}\,(ii) instead: one common channel, no block opened, no
constraint produced. The probe setting therefore enters this contract
through Lemma~\ref{lem:linear}\,(ii), and a mode register that arrived with
the light through Lemma~\ref{lem:master}; swapping the two would be a
modelling error, giving Eve a register she does not have or withholding one
she does.
\end{remark}
\claimmeta{\metalabel{Uses}\mref{Lem.}{lem:linear}, \mref{Lem.}{lem:master}
\metasep newly derived
\metasep accepted}

% claim reductions/rem:branch-blocks
\claimsummary{reductions/rem:branch-blocks}{The branching rule squashes a conditionally characterized receiver as soon as each branch squashes, stacking the branch targets as blocks. A branch that cannot be modelled gets a block of its own, where the adversary fixes the outcome only in proportion to the branch's probability.}
\claimbookmark{reductions/rem:branch-blocks}{Remark}{rem:branch-blocks}{Conditionally characterized receivers under the branching rule}
\begin{remark}[Conditionally characterized receivers under the branching rule]\label{rem:branch-blocks}
The branching rule of Lemma~\ref{lem:master} turns any conditional
description of a receiver into a squashing relation. The device is squashable
as soon as \emph{each branch} is, and the squashed target is the branch
targets stacked as diagonal blocks. A branch one cannot model is not fatal: it
gets a block of its own, on which the squashed adversary rather than Bob's
detector fixes the outcome, but only in proportion to how likely that branch
was.
\end{remark}
\claimmeta{\metalabel{Uses}\mref{Lem.}{lem:master}
\metasep newly derived
\metasep accepted}

% claim reductions/rem:branch-probability
\claimsummary{reductions/rem:branch-probability}{In the branching rule only the branch label is classical: the probability of a branch is in general an operator, not a number. That operator is what the squashed model retains about the branching, and it is the only constraint the rule produces.}
\claimbookmark{reductions/rem:branch-probability}{Remark}{rem:branch-probability}{The branch probability as an operator}
\begin{remark}[The branch probability as an operator]\label{rem:branch-probability}
In Lemma~\ref{lem:master} the branch label is classical, but nothing else
about the branching is: the step that produces the label is an arbitrary CP
map, and the probability of branch $j$ is in general an operator, not a
number. That operator is what the squashed model retains about the branching,
and it is the only constraint the rule produces.
\end{remark}
\claimmeta{\metalabel{Uses}\mref{Lem.}{lem:master}
\metasep newly derived
\metasep accepted}

% claim reductions/cor:instances
\claimsummary{reductions/cor:instances}{The branching rule applies to a convex mixture with a known weight, to a measurement block diagonal for an orthogonal decomposition and to a register read by a POVM. Only the known mixture fixes the branch weight as a number without further information.}
\claimbookmark{reductions/cor:instances}{Corollary}{cor:instances}{Three special cases of instrument-conditioned squashing}
\claimleadin{%
Three choices of the generalized instrument recur in this module's
constructions: a known mixture, a split of the input space into orthogonal
blocks, and a POVM measurement of a register that arrives with the light.
}
\begin{corollary}[Three special cases of instrument-conditioned squashing]\label{cor:instances}
Let $\Lambda_j$ and $F^{(j)}$ be branch channels and branch targets as in
Lemma~\ref{lem:master}.
\begin{enumerate}[label=(\roman*),leftmargin=3.0em]
\item \emph{(convex mixture with a known weight)} If
  $M_v=(1-q)G_v+qR_v$ for POVMs $G,R$ on $\Bsys$ and $q\in[0,1]$, with
  $G\dsq_{\Lambda_G}F^{(G)}$ and $R$ \emph{arbitrary}, then
  $M\dsq\big\{F^{(G)}_v\dsum\dyad{v}_{\Flag}\big\}_v$, with the sound branch
  constraints $\tr[\Pi_1\sigma]=1-q$ and $\tr[\Pflag\sigma]=q$, in the operator
  form of Lemma~\ref{lem:master}\,(iv).
\item \emph{(block-wise squashing)} Let $\Hs{\Bsys}=\bigoplus_j\Hs{\Bsys_j}$ be
  a finite orthogonal decomposition and let every $M_v$ be block diagonal,
  $M_v=\bigoplus_jM^{(j)}_v$, with $M^{(j)}\dsq_{\Lambda_j}F^{(j)}$. Then
  $M\dsq\bigoplus_jF^{(j)}$.
\item \emph{(POVM-conditioned squashing)} If $\Bsys=\Bsys_0D$ with a POVM
  $\{N_j\}_{j\in J}$ on $D$ and $M_v=\sum_jM_{v|j}\otimes N_j$ with
  $M_{|j}\dsq_{\Lambda_j}F_{|j}$ on $\Qsys_j$, then
  $\{M_v\}_v\dsq_{\Lambda}\big\{\bigoplus_jF_{v|j}\big\}_v$ with
  \begin{equation}\label{eq:povmcond}
    \Lambda(\rho)=\bigoplus_{j\in J}\Lambda_j\Big(
      \tr_{D}\big[(\idon{\Bsys_0}\otimes N_j^{1/2})\,\rho\,
      (\idon{\Bsys_0}\otimes N_j^{1/2})\big]\Big).
  \end{equation}
\end{enumerate}
\end{corollary}
\claimmeta{\metalabel{Uses}\mref{Lem.}{lem:top}, \mref{Lem.}{lem:master}
\metasep \metalabel{Used by}\mref{Thm.}{thm:mixture}, \mref{Prop.}{prop:chain}
\metasep newly derived
\metasep accepted}

\begin{proof}
Apply Lemma~\ref{lem:master} with, respectively:

(i) $\Ins_1=(1-q)\idmap$ and $\Ins_2=q\,\idmap$, branch targets $F^{(G)}$ and
the flag PVM of Lemma~\ref{lem:top} (which exists for the unknown $R$), so that
$M_{v,1}=(1-q)G_v$ and $M_{v,2}=qR_v$ and part~(iii) gives
$\sum_jM_{v,j}=M_v$. Here $\Ins_j^{\dagger}(\idon{})$ is a multiple of
$\idon{\Bsys}$, which is the hypothesis of part~(iv).

(ii) $\Ins_j(\rho)\coloneqq P_j\rho P_j$ read as an operator on $\Bsys_j$, for
$P_j$ the projection onto $\Hs{\Bsys_j}$; then $\Ins_j^{\dagger}(X_j)$ is
$X_j$ embedded in the $j$-th block, so $M_{v,j}$ is $M^{(j)}_v$ so embedded and
part~(iii) gives $\sum_jM_{v,j}=M_v$. The instrument condition is
$\sum_jP_j=\idon{\Bsys}$.

(iii) $\Ins_j(\rho)\coloneqq
\tr_{D}\big[(\idon{\Bsys_0}\otimes N_j^{1/2})\rho(\idon{\Bsys_0}\otimes
N_j^{1/2})\big]$, which is CP as a Kraus map followed by a partial trace, and
satisfies $\sum_j\tr\Ins_j(\rho)=\sum_j\tr[(\idon{}\otimes N_j)\rho]=\tr\rho$.
Its adjoint is
$\Ins_j^{\dagger}(X)=(\idon{}\otimes N_j^{1/2})(X\otimes\idon{D})
(\idon{}\otimes N_j^{1/2})=X\otimes N_j$, the two factors commuting because
they act on different tensor factors; hence
$M_{v,j}=M_{v|j}\otimes N_j$ and part~(iii) gives
$\sum_jM_{v,j}=\sum_jM_{v|j}\otimes N_j=M_v$.
\end{proof}

% claim reductions/prop:common
\claimsummary{reductions/prop:common}{If every branch target squashes to one POVM, the coarse-grained measurement of \hyperlink{claim-reductions/lem:master}{Lemma~\ref*{lem:master}} squashes to that POVM directly, through the sum of the composed branch channels. The branch blocks then disappear from the target.}
\claimbookmark{reductions/prop:common}{Proposition}{prop:common}{Common and dominating targets}
\claimleadin{%
Only in Corollary~\ref{cor:instances}\,(i) is the branch weight state
independent, hence the constraint a fixed number. In its case~(ii) the weight
is the projection onto the block the input fell in, and that operator-valued
case is the one a flag-state squasher runs on. Its case~(iii) is stated for a
general POVM $\{N_j\}$ rather than a PVM, and it is what the
register-conditioned squashing of Theorem~\ref{thm:mixture} runs on: a
\emph{classical branch register} read by a POVM rather than a PVM, which
Corollary~\ref{cor:register-instance} lifts to \termref{term-reduction-instance}{reduction instances}. The price
of the generality is that each branch gets its own block of $\Qsys$ rather
than sharing one register, which is exactly the block structure of $\Qsys$
the consuming declaration wants anyway. Where every branch target squashes to
one common target, the blocks can be merged again.
}
\begin{proposition}[Common and dominating targets]\label{prop:common}
In the setting of Lemma~\ref{lem:master}, suppose additionally
$F^{(j)}\dsq_{\Psi_j}F$ for a single POVM $F$ on a space $\Qsys_{*}$ (a
\emph{dominating target}; $\Psi_j=\idmap$ if all branch targets already
coincide). Then
\begin{equation*}
  \Big\{\sum_{j}M_{v,j}\Big\}_v\;\dsq_{\Lambda'}\;F,
  \qquad
  \Lambda'\coloneqq\sum_{j\in J}\Psi_j\circ\Lambda_j\circ\Ins_j .
\end{equation*}
\end{proposition}
\claimmeta{\metalabel{Uses}\mref{Lem.}{lem:master}
\metasep \metalabel{Used by}\mref{Rem.}{rem:dominating}, \mref{Req.}{req:reserved-link-kinds}, \mref{Thm.}{thm:noise-channel}
\metasep newly derived
\metasep accepted}

\begin{proof}
$\Lambda'$ is CP as a sum of compositions of CP maps and trace preserving by
the computation in Lemma~\ref{lem:master}\,(ii). Then
$\Lambda'^{\dagger}(F_v)=\sum_j\Ins_j^{\dagger}\Lambda_j^{\dagger}
\Psi_j^{\dagger}(F_v)
=\sum_j\Ins_j^{\dagger}\Lambda_j^{\dagger}\big(F^{(j)}_v\big)
=\sum_j\Ins_j^{\dagger}\big(M^{(j)}_v\big)=\sum_jM_{v,j}$.
\end{proof}

% claim reductions/rem:dominating
\claimsummary{reductions/rem:dominating}{A detector whose POVM dominates a characterized one is the characterized device plus a weighted flag, as in the published block-resolved construction. By \hyperlink{claim-reductions/prop:common}{Proposition~\ref*{prop:common}}, a family of devices whose members squash to a common target becomes one squashed model with the label a branch register.}
\claimbookmark{reductions/rem:dominating}{Remark}{rem:dominating}{Dominated detectors and device families}
\begin{remark}[Dominated detectors and device families]\label{rem:dominating}
A device has the form (i) of Corollary~\ref{cor:instances} whenever its POVM
dominates a characterized one: if $M_v\succeq(1-q)G_v$ for all $v$ with
$q\in(0,1]$, then $R_v\coloneqq\big(M_v-(1-q)G_v\big)/q$ is positive and sums
to $\idon{\Bsys}$, so any detector that is $q$-close to a characterized
detector \emph{in this operator sense} may be analysed as the characterized
device plus a $q$-weighted flag. The published block-resolved form of that
construction is \cite[Thm.~3, Eqs.~(20)--(23)]{nahar2025imperfect}, which
builds exactly this channel from the block dominations
$\Gamma_{m,\text{noise}}-(1-q_m)\Gamma_{m,\text{ideal}}\succeq0$ and records
the weight $\Lambda^{\dagger}[\Pi_m]=(1-q_m)\Pi_m$ kept on block $m$, the flag
receiving $q_m$.

By Proposition~\ref{prop:common}, a family of devices indexed by a label $\xi$
that the adversary selects becomes one squashed model as soon as each member
squashes to a \emph{common} target, with the label a branch register carrying
its own constraints from Lemma~\ref{lem:master}\,(iv). The operator
factorization such a construction rests on (range inclusion, majorization and
factorization are equivalent for bounded operators) is
\cite[Thm.~1]{douglas1966majorization}; no proof of this contract uses it.
\end{remark}
\claimmeta{\metalabel{Uses}\mref{Prop.}{prop:common}
\metasep newly derived
\metasep accepted}

% claim reductions/cor:flagsenough
\claimsummary{reductions/cor:flagsenough}{In a flag-extended target, the block the squashing channel never populates may carry any POVM, and the target does not depend on the device. An uncharacterized branch can therefore be represented without a device model for it.}
\claimbookmark{reductions/cor:flagsenough}{Corollary}{cor:flagsenough}{Arbitrary unpopulated block in a flag-extended target}
\claimleadin{%
Lemma~\ref{lem:top} says every POVM squashes to a flag PVM. The following
sharpening is what a construction needs: in a \emph{flag-extended} target, the
block the \termref{term-squashing-channel}{squashing channel} never populates may be replaced by any POVM
at all. An uncharacterized branch therefore needs no model, neither an
approximate nor a conservative one, and one target covers every device with
the same flag structure. The branch blocks of a composite target \emph{may}
accordingly be padded to a common size; a reduction instance emits them as its
branches supply them, and a padded variant only where a submodule's interface
declares it (Requirement~\ref{req:unpadded-branch-blocks}).
}
\begin{corollary}[Arbitrary unpopulated block in a flag-extended target]\label{cor:flagsenough}
Let $\{G_v\}_{v\in\Vset}$ be any POVM on a finite-dimensional space
$\Bsys_0$ and let $\{H_v\}_{v\in\Vset}$ be any POVM on $\Bsys_0$. Then, on
$\Hs{\Bsys_0}\dsum\Hs{\Flag}$,
\begin{equation}\label{eq:flagsenough}
  \big\{G_v\dsum\dyad{v}_{\Flag}\big\}_v
  \;\dsq_{\Lambda}\;
  \big\{H_v\dsum\dyad{v}_{\Flag}\big\}_v,
  \qquad
  \Lambda(\rho)\coloneqq
   0_{\Bsys_0}\dsum\sum_v\tr\big[(G_v\dsum\dyad{v})\rho\big]\,\dyad{v}_{\Flag}.
\end{equation}
In particular $H_v=\tfrac{1}{|\Vset|}\idon{\Bsys_0}$ is admissible, and the
squashing channel depends on $G$ while the target does not.
\end{corollary}
\claimmeta{\metalabel{Uses}\mref{Lem.}{lem:top}
\metasep \metalabel{Used by}\mref{Rem.}{rem:flagcost}, \mref{Req.}{req:unpadded-branch-blocks}
\metasep newly derived
\metasep accepted
\metasep \metalabel{Oracles}\coderef{https://github.com/fraunhofersg/certiqs-key/blob/v0.1.0/tests/oracles/registry_reductions.py}{registry_reductions.py}}

\begin{proof}
$\Lambda$ is a measure-and-prepare channel, hence CP; it is trace preserving
because $\sum_v(G_v\dsum\dyad{v})=\idon{\Bsys_0}\dsum\idon{\Flag}$. Its output
is supported in the flag block, so
$\Lambda^{\dagger}\big(X_{\Bsys_0}\dsum X_{\Flag}\big)
=\sum_v\bra{v}X_{\Flag}\ket{v}\,\big(G_v\dsum\dyad{v}\big)$, the
$\Bsys_0$-part of the observable contributing nothing, and evaluating at
$X_{\Bsys_0}=H_v$, $X_{\Flag}=\dyad{v}$ returns $G_v\dsum\dyad{v}$. That
$\{H_v\dsum\dyad{v}\}$ is a POVM follows from $\sum_vH_v=\idon{\Bsys_0}$.
\end{proof}

% claim reductions/rem:flagcost
\claimsummary{reductions/rem:flagcost}{Replacing the unpopulated block by an arbitrary POVM gives a cleaner target, which hands the adversary more and costs rate. What it buys is one target, independent of the device, for a branch that is unknown or adversarially chosen; an infinite-dimensional branch takes the flag PVM.}
\claimbookmark{reductions/rem:flagcost}{Remark}{rem:flagcost}{Rate cost of a flag-extended target}
\begin{remark}[Rate cost of a flag-extended target]\label{rem:flagcost}
The replacement of Corollary~\ref{cor:flagsenough} is not free in rate: by
Corollary~\ref{cor:monotone} the target of \eqref{eq:flagsenough} is cleaner
than the source, so it hands the adversary more, and every squashed state it
admits is a flag state as far as the $\Bsys_0$ block is concerned. What it
buys is that $G$ appears nowhere in the target: a branch whose device is
unknown or adversarially chosen is covered by one target that does not depend
on the device, and an infinite-dimensional branch by the finite flag PVM of
Lemma~\ref{lem:top}. The instance of such a branch is the flag instance of
Corollary~\ref{cor:flag-instance}, on which by construction nothing is
assumed and nothing may be.
\end{remark}
\claimmeta{\metalabel{Uses}\mref{Cor.}{cor:monotone}, \mref{Lem.}{lem:top}, \mref{Cor.}{cor:flagsenough}
\metasep newly derived
\metasep accepted}

% claim reductions/thm:mixture
\claimsummary{reductions/thm:mixture}{A receiver conditioned by a POVM measurement of an incoming register squashes to the direct sum of its branch targets, with branch operators and branch weights transported as operators. A branch weight constrains the squashed states only where it is a declared multiple of the side system's state.}
\claimbookmark{reductions/thm:mixture}{Theorem}{thm:mixture}{Register-conditioned squashing}
\claimleadin{%
When the branch is selected by a measurement on a register that arrives with
the light, the per-branch squashers combine into one map. The theorem is
Corollary~\ref{cor:instances}\,(iii) with the transported operators and the
branch weights carried along, and it is stated for a general POVM on the
register rather than a projective measurement.
}
\begin{theorem}[Register-conditioned squashing]\label{thm:mixture}
Let $\Bsys=\Bsys_0D$, let $\{N_j\}_{j\in J}$ be a POVM on $D$ with $J$ finite,
and let
\begin{equation}\label{eq:regmodel}
  M_v=\sum_{j\in J}M_{v|j}\otimes N_j ,
  \qquad v\in\Vset ,
\end{equation}
where $\{M_{v|j}\}_{v}$ is a POVM on $\Bsys_0$ for each $j$. Suppose
$(M_{|j};O^{(j)})\dsq_{\Lambda_j}(F_{|j};\hat O^{(j)})$ for each $j$, with
$F_{|j}$ a POVM on a finite-dimensional $\Qsys_j$. Put
$\Hs{\Qsys}=\bigoplus_{j\in J}\Hs{\Qsys_j}$ with block projections $\Pi_j$ and
\begin{equation}\label{eq:regchannel}
  \Lambda(\rho)\coloneqq\bigoplus_{j\in J}\Lambda_j\Big(
    \tr_{D}\big[(\idon{\Bsys_0}\otimes N_j^{1/2})\,\rho\,
    (\idon{\Bsys_0}\otimes N_j^{1/2})\big]\Big).
\end{equation}
Then $\Lambda$ is a channel and:
\begin{enumerate}[label=(\roman*),leftmargin=3.0em]
\item $\{M_v\}_v\dsq_{\Lambda}\big\{\bigoplus_jF_{v|j}\big\}_v$; in particular
  $\{M_v\}$ is a POVM on $\Bsys$;
\item the composite transports the operator $\hat O^{(j)}_k$ placed in block
  $j$ to $O^{(j)}_k\otimes N_j$;
\item $\Lambda^{\dagger}(\Pi_j)=\idon{\Bsys_0}\otimes N_j$, so on every
  reachable $\sigma$ and every finite-dimensional side system $R$,
  \begin{equation}\label{eq:branchweight}
    \tr_{\Qsys}\big[(\Pi_j\otimes\idon{R})\,\sigma\big]
    =\tr_{\Bsys}\big[\big((\idon{\Bsys_0}\otimes N_j)\otimes\idon{R}\big)\,\rho\big];
  \end{equation}
  this is an identity, not a constraint, unless the right-hand side is a
  declared multiple of $\rho_R$;
\item if every $F_{|j}$ and every $\hat O^{(j)}_k$ is block diagonal for a
  finite orthogonal decomposition of $\Hs{\Qsys_j}$, then every
  $\bigoplus_jF_{v|j}$ and every carried operator is block diagonal for the
  union decomposition of $\Hs{\Qsys}$, and $\PinS\circ\Lambda$ realizes the
  same relation with block-diagonal reachable states
  (Corollary~\ref{cor:blockdiag}).
\end{enumerate}
\end{theorem}
\claimmeta{\metalabel{Uses}\mref{Lem.}{lem:sep}, \mref{Cor.}{cor:pullback}, \mref{Cor.}{cor:blockdiag}, \mref{Def.}{def:generalized-instrument}, \mref{Lem.}{lem:master}, \mref{Cor.}{cor:instances}
\metasep \metalabel{Used by}\mref{Rem.}{rem:povmreg}, \mref{Cor.}{cor:register-instance}
\metasep newly derived
\metasep accepted
\metasep \metalabel{Oracles}\coderef{https://github.com/fraunhofersg/certiqs-key/blob/v0.1.0/tests/oracles/registry_reductions.py}{registry_reductions.py}}

\begin{proof}
Put $\Ins_j(\rho)\coloneqq\tr_{D}[(\idon{\Bsys_0}\otimes N_j^{1/2})\rho
(\idon{\Bsys_0}\otimes N_j^{1/2})]$. Each $\Ins_j$ is CP, being a
single-Kraus-operator map followed by a partial trace, and
$\sum_j\tr\Ins_j(\rho)=\sum_j\tr[(\idon{}\otimes N_j)\rho]=\tr\rho$, so
$\{\Ins_j\}$ is a generalized instrument
(Definition~\ref{def:generalized-instrument}). Its adjoint is
\begin{equation*}
  \Ins_j^{\dagger}(X)
  =(\idon{}\otimes N_j^{1/2})(X\otimes\idon{D})(\idon{}\otimes N_j^{1/2})
  =X\otimes N_j
  \qquad(X\in\Lin{\Bsys_0}),
\end{equation*}
the two factors commuting because they act on different tensor factors; the
first equality is \eqref{eq:pairing} for $\Ins_j$, since
$\tr[X\,\tr_D(K\rho K)]=\tr[(X\otimes\idon{D})\,K\rho K]
=\tr[K(X\otimes\idon{D})K\,\rho]$ for the Hermitian
$K=\idon{\Bsys_0}\otimes N_j^{1/2}$, by cyclicity.

Now apply Lemma~\ref{lem:master} with these $\Ins_j$, branch POVMs
$M^{(j)}_v\coloneqq M_{v|j}$ and branch channels $\Lambda_j$. Part (i) of that
lemma gives that $\{M_{v,j}=M_{v|j}\otimes N_j\}$ is a POVM on $\Bsys$ over
$\Vset\times J$, and part (iii), which coarse-grains the branch label away,
gives $\sum_jM_{v,j}=M_v$ and the relation (i) here, with the channel
\eqref{eq:regchannel}, which is \eqref{eq:master}. Part (ii) here is
Lemma~\ref{lem:master}\,(v) with $\Ins_j^{\dagger}(O^{(j)}_k)=O^{(j)}_k\otimes
N_j$. Part (iii) here is Lemma~\ref{lem:master}\,(iv) with
$\Ins_j^{\dagger}(\idon{\Bsys_0})=\idon{\Bsys_0}\otimes N_j$, followed by
Corollary~\ref{cor:pullback} read on the carried operator $\Pi_j$, whose
pull-back is that operator. Part (iv) is immediate: a direct sum of
block-diagonal operators is block diagonal for the union of the
decompositions, and Corollary~\ref{cor:blockdiag} applies.
\end{proof}

% claim reductions/rem:povmreg
\claimsummary{reductions/rem:povmreg}{Keeping the register intact in the target would need a projective measurement of it; filing each branch as a block of the squashed system removes that hypothesis. The price is bookkeeping only: a direct sum instead of a tensor product, the branch weight still transported as an operator.}
\claimbookmark{reductions/rem:povmreg}{Remark}{rem:povmreg}{A general POVM on the register, and its price}
\begin{remark}[A general POVM on the register, and its price]
\label{rem:povmreg}
Theorem~\ref{thm:mixture} assumes only that $\{N_j\}$ is a POVM. A
construction that keeps the register $D$ intact in the target, squashing
$\sum_jM_{v|j}\otimes N_j$ to $\sum_jF_{v|j}\otimes N_j$ on $\Qsys_0D$, does
need $\{N_j\}$ to be projective, because its proof turns on
$\mathcal{P}_i(X\otimes N_j)=\delta_{ij}X\otimes N_j$ for the projections
$\mathcal{P}_i=(\idon{}\otimes N_i)(\cdot)(\idon{}\otimes N_i)$, and that
identity fails for a general POVM. Filing each branch as a \emph{block} of
$\Qsys$ instead of a value of a retained register removes the hypothesis. The
price is bookkeeping and nothing else: the target is a direct sum rather than
a tensor product, which is the \exttermref{../foundations/contract.pdf}{term-block-list}{block-list} form of
\xref{foundations/def:sector-map} and the form every reduction instance emits
(Definition~\ref{def:reduction-instance}\,(I2)); and
Theorem~\ref{thm:mixture}\,(iii) shows the branch weight is still
transported, as the operator $\idon{\Bsys_0}\otimes N_j$ rather than as the
number it is not.
\end{remark}
\claimmeta{\metalabel{Uses}\mref{Thm.}{thm:mixture}
\metasep newly derived
\metasep accepted}

\section{Event-level device facts, attack sets and the licence}\label{group-reductions/prop:obstruction}
\sectionintro{A device fact about a union of outcomes cannot in general be kept by a classical target (\hyperlink{claim-reductions/prop:obstruction}{Proposition~\ref*{prop:obstruction}}), so it travels as a constraint entry, which is sound only on the reachable set. The consuming engine restricts the squashed attacks to the admissible set the entries cut out. The main result, the attack-set licence \hyperlink{claim-reductions/thm:licence}{Theorem~\ref*{thm:licence}}, proves that every squashed physical attack lies in that set whenever each carried operator's pull-back satisfies a spectral certificate, in either source mode and, for party-local entries, on a product \termref{term-measured-system}{measured system} with a product \termref{term-squashing-channel}{squashing channel} and each certificate on its factor's own channel. The remarks say how the relation symbols read at the channel and the declaration level, what the licence does and does not license, and where its format is published.}
\sectionimports{\extref{../foundations/contract.pdf}{foundations/def:round-model}}

% claim reductions/prop:obstruction
\claimsummary{reductions/prop:obstruction}{Some three-outcome qubit POVM with a lower bound on the sum of two of its elements has no squashing to a classical target that keeps that bound as an operator inequality. A device fact about a union of outcomes therefore travels as a constraint entry rather than inside a flag target.}
\claimbookmark{reductions/prop:obstruction}{Proposition}{prop:obstruction}{Obstruction to encoding event bounds in classical targets}
\claimleadin{%
The \termref{term-weight-preserving-target}{weight-preserving flag-state target} of the flag-state submodule encodes
one certified weight \emph{per outcome}. For an \emph{event}
$\mathcal{O}_{*}\subseteq\Vset$, a union of outcomes that must not be merged
(for instance because they carry key information), a device may certify only
the sum, $\sum_{v\in\mathcal{O}_{*}}M_v\succeq p_{*}\idon{\Bsys}$. The
following proposition shows that no classical target can absorb such a
bound, whatever the construction.
}
\begin{proposition}[Obstruction to encoding event bounds in classical targets]
\label{prop:obstruction}
Call a \termref{term-target-povm}{target POVM} $F$ on $\Qsys$ \emph{classical} if every $F_v$ is diagonal
in one fixed orthonormal basis $\{\ket{f}\}_f$ of $\Hs{\Qsys}$. There exist a
POVM $M$ on $\Hs{\Bsys}=\mathbb{C}^{2}$, an event
$\mathcal{O}_{*}\subseteq\Vset$ and a constant $p_{*}>0$ with
$\sum_{v\in\mathcal{O}_{*}}M_v\succeq p_{*}\idon{\Bsys}$ such that \emph{no}
pair $(F,\Lambda)$ exists with $F$ classical, $M\dsq_{\Lambda}F$ and
$\sum_{v\in\mathcal{O}_{*}}F_v\succeq p_{*}\idon{\Qsys}$. Concretely,
\begin{equation}\label{eq:obstruction}
  M_1=\tfrac12\dyad{0},\quad
  M_2=\tfrac12\dyad{+},\quad
  M_3=\idon{\Bsys}-M_1-M_2,\quad
  \mathcal{O}_{*}=\{1,2\},\quad
  p_{*}=\tfrac12\Big(1-\tfrac{1}{\sqrt2}\Big),
\end{equation}
with $\ket{+}=(\ket{0}+\ket{1})/\sqrt2$.
\end{proposition}
\claimmeta{\metalabel{Uses}\mref{Def.}{def:squashing}
\metasep \metalabel{Used by}\mref{Rem.}{rem:obstruction-scope}, \mref{Req.}{req:constraint-and-licence-slots}
\metasep newly derived
\metasep accepted
\metasep \metalabel{Oracles}\coderef{https://github.com/fraunhofersg/certiqs-key/blob/v0.1.0/tests/oracles/registry_reductions.py}{registry_reductions.py}}

\begin{proof}
\emph{The data are admissible.} For two rank-one projections with overlap
$|\innerp{0}{+}|=1/\sqrt2$, the operator $\dyad{0}+\dyad{+}$ has trace $2$ and
eigenvalues $1\pm1/\sqrt2$; hence
$\lmin(M_1+M_2)=\tfrac12(1-1/\sqrt2)=p_{*}>0$ and
$\lmax(M_1+M_2)=\tfrac12(1+1/\sqrt2)<1$, so $M_3\succeq0$ and $\{M_1,M_2,M_3\}$
is a POVM with $M_1+M_2\succeq p_{*}\idon{\Bsys}$.

\emph{No classical target encodes the bound.} Suppose $(F,\Lambda)$ existed.
Put $T_f\coloneqq\Lambda^{\dagger}(\dyad{f})$. Then $\{T_f\}_f$ is a POVM on
$\Bsys$: each $T_f\succeq0$ because $\Lambda^{\dagger}$ is positive, and
$\sum_fT_f=\Lambda^{\dagger}(\idon{\Qsys})=\idon{\Bsys}$ by unitality. Writing
$F_v=\sum_fF_v(f)\dyad{f}$ with $F_v(f)\ge0$ and $\sum_vF_v(f)=1$ (POVM
completeness read on the diagonal), linearity gives
\begin{equation*}
  M_v=\Lambda^{\dagger}(F_v)=\sum_fF_v(f)\,T_f .
\end{equation*}
The encoded event bound $\sum_{v\in\mathcal{O}_{*}}F_v\succeq p_{*}\idon{\Qsys}$
reads, on the diagonal, $F_1(f)+F_2(f)\ge p_{*}>0$ for every $f$, so every $f$
has $F_1(f)>0$ or $F_2(f)>0$.

Suppose $F_1(f)>0$. Every summand of $M_1=\sum_{f'}F_1(f')T_{f'}$ is positive
semidefinite, so $0\preceq F_1(f)T_f\preceq M_1=\tfrac12\dyad{0}$. A positive
operator dominated by a rank-one positive operator is a non-negative multiple
of it: if $X\preceq c\dyad{0}$ with $X\succeq0$, then for every
$\ket{\psi}\perp\ket{0}$ we get $\bra{\psi}X\ket{\psi}\le0$, hence
$\bra{\psi}X\ket{\psi}=0$ and $X\ket{\psi}=0$, so $\operatorname{supp}X
\subseteq\operatorname{span}\ket{0}$. Thus $T_f\propto\dyad{0}$. Symmetrically
$F_2(f)>0$ forces $T_f\propto\dyad{+}$. A non-zero operator cannot be
proportional to both, since $\dyad{0}$ and $\dyad{+}$ are distinct rank-one
projections; so $F_1(f)F_2(f)>0$ forces $T_f=0$, and in every case each $T_f$
is $0$, or a non-negative multiple of $\dyad{0}$, or a non-negative multiple of
$\dyad{+}$.

Summing, $\idon{\Bsys}=\sum_fT_f=\alpha\dyad{0}+\beta\dyad{+}$ for some
$\alpha,\beta\ge0$ (two non-negative scalars local to this proof, unrelated
to the R\'enyi order $\alpha$). Testing against $\ket{1}$ gives
$1=\bra{1}\big(\alpha\dyad{0}+\beta\dyad{+}\big)\ket{1}=\beta/2$, so
$\beta=2$; testing the off-diagonal entry gives
$0=\bra{0}\big(\alpha\dyad{0}+\beta\dyad{+}\big)\ket{1}=\beta/2$, so
$\beta=0$ --- a contradiction.
\end{proof}

% claim reductions/rem:obstruction-scope
\claimsummary{reductions/rem:obstruction-scope}{A general quantum target encodes an event bound trivially; the obstruction concerns targets classical on the whole target space. Such a bound travels as a constraint entry under the licence, and merging the outcomes is no remedy when they carry key information.}
\claimbookmark{reductions/rem:obstruction-scope}{Remark}{rem:obstruction-scope}{Scope of the obstruction, and what it forces}
\begin{remark}[Scope of the obstruction, and what it forces]
\label{rem:obstruction-scope}
With a general \emph{quantum} target the encoding is trivial: take $F=M$ and
$\Lambda=\idmap$, and the physical bound is itself an operator inequality on
the target. The content of Proposition~\ref{prop:obstruction} is that a
target classical on the whole of $\Qsys$ cannot in general carry an
event-level bound, whatever the channel.
Per-outcome bounds, by contrast, are encodable: a
\termref{term-weight-preserving-target}{weight-preserving flag-state target} of the flag-state submodule carries one weight per outcome
as an operator inequality on the target. Event-level information that the
target cannot carry must travel as a pull-back constraint, and a pull-back constraint is
sound only on the reachable set, which is what the licence of
Theorem~\ref{thm:licence} certifies. Three consequences follow.
\begin{enumerate}[label=(\alph*),leftmargin=2.4em]
\item A record without a constraint list and a licence slot excludes a class
  of device knowledge that no classical target can carry; this is why the
  artefact carries both even when both are empty
  (Requirement~\ref{req:constraint-and-licence-slots}).
\item Re-declaring the coarse-graining does not help. Merging the outcomes of
  $\mathcal{O}_{*}$ into one letter turns the event bound into a per-outcome
  bound and makes it encodable, but only by destroying the distinction
  between the merged outcomes, which is precisely what ``must not be merged''
  forbids when they carry key information (the fibre condition of the
  flag-state submodule's \termref{term-weight-route}{weight route}).
\item The counterexample is two-dimensional and uses only three outcomes, so it
  is not an artefact of infinite dimension, of photon number, or of the optical
  setting. It applies to every squashed model whose target is classical on
  the whole of $\Qsys$, and it says nothing of a target with a block on which
  it is not classical, which can carry the bound as the trivial encoding
  above does.
\end{enumerate}
\end{remark}
\claimmeta{\metalabel{Uses}\mref{Prop.}{prop:obstruction}
\metasep \metalabel{Used by}\mref{Req.}{req:constraint-and-licence-slots}
\metasep newly derived
\metasep accepted}

% claim reductions/def:attset
\claimsummary{reductions/def:attset}{For a constraint list with real constants, the admissible attack channels are the round model's attack sets: adjoint conditions in each source mode and, for party-local entries, a condition with a reference system. The licence \hyperlink{claim-reductions/thm:licence}{Theorem~\ref*{thm:licence}} says when every squashed physical attack lies there.}
\claimbookmark{reductions/def:attset}{Definition}{def:attset}{Admissible attack sets in the consumer's format}
\claimleadin{%
The consuming engine accepts a declared constraint list inside its security
chain only under an assumption it never verifies: that security of the
physical protocol reduces to the squashed \exttermref{../foundations/contract.pdf}{term-round-model}{round model} with attacks
restricted to the declared admissible set. This definition fixes that set in
the consumer's format, and Theorem~\ref{thm:licence} proves the reduction
from one hypothesis.
}
\begin{definition}[Admissible attack sets in the consumer's format]
\label{def:attset}
Let $\CList=\{(\hat O_k,\bowtie_k,\Theta_k)\}_{k=1}^{m}$ be a declared list of
Hermitian $\hat O_k\in\Lin{\Qsys}$, relations
$\bowtie_k\in\{\preceq,=,\succeq\}$ and right-hand sides typed by the declared
source mode. For a list whose entries carry real constants $c_k$, as every
list this module emits does (Definition~\ref{def:licensed}), the admissible
sets are the \exttermref{../foundations/contract.pdf}{term-attack-set}{attack sets} of $\CList$ of
\xref{foundations/def:round-model} at the per-round system $S=\Qsys$, which the
consuming contract states at the same system
(\xref{meat-core/def:roundmodel}); their displays at that system are,
verbatim,
\begin{equation}\label{eq:attset}
  \AttSet_{\CList}(E',\,\Qsys E'')\coloneqq
  \Big\{\Aev\in\mathrm{CPTP}(E',\Qsys E''):\
  \Aev^{\dagger}\big(\hat O_k\otimes\idon{E''}\big)
  \bowtie_k c_k\,\idon{E'}\ \ \forall k\Big\}
  \qquad\text{(trusted-source mode)},
\end{equation}
and, with the constraint operator lifted to the larger supplied system,
\begin{equation}\label{eq:attsetun}
  \AttSet_{\CList}(E',\,A\Qsys E'')\coloneqq
  \Big\{\Aev\in\mathrm{CPTP}(E',A\Qsys E''):\
  \Aev^{\dagger}\big(\idon{A}\otimes\hat O_k\otimes\idon{E''}\big)
  \bowtie_k c_k\,\idon{E'}\ \ \forall k\Big\}
\end{equation}
in \exttermref{../transmitter/contract.pdf}{term-untrusted-source-mode}{untrusted-source mode}, the two distinguished by the output system
named in the argument. This contract writes $\bowtie_k$ in the operator
spelling of \eqref{eq:attset} and \eqref{eq:attsetun} in both source modes,
and everything it emits is admitted under either spelling. An entry that declares a Hermitian
right-hand side $\Theta_k$ on Alice's system, rather than a real constant, is
a \emph{general} entry: it is admissible input to the consuming contract but
is never licensed (\xref{protocol-compiler/def:homog}). Every entry this module emits
carries a real scalar $c_k$ and is written $(\hat O_k,\bowtie_k,c_k)$
(Definition~\ref{def:licensed}, Theorem~\ref{thm:licence}\,(iv)).

\emph{Party-local entries.} Where the \termref{term-measured-system}{measured system} is a product
$\Qsys=\bigotimes_f\Qsys_f$ and $\CList$ carries entries tagged party-local
(Definition~\ref{def:licensed}), the admissible set is cut out, for each such
entry $(\hat O_k,\bowtie_k,c_k;f)$, by the condition
\begin{equation}\label{eq:attsetpl}
  \big(\idmap_{R'}\otimes\Aev\big)^{\dagger}
  \Big(X\otimes\big(\hat O_k-c_k\,\idon{\Qsys_f}\big)\otimes\idon{E''}\Big)
  \;\bowtie_k\;0
  \qquad\text{for every finite-dimensional }R'\text{ and every }
  X\in\Pos{R'R\,\Qsys_{f^c}} ,
\end{equation}
in place of the scalar condition of \eqref{eq:attset} or
\eqref{eq:attsetun}, where $\Qsys_{f^c}$ is the product of the other
factors and $R$ is $A$ in untrusted-source mode and trivial in
\exttermref{../transmitter/contract.pdf}{term-trusted-source-mode}{trusted-source mode}; the untagged entries keep their conditions. The
reference system $R'$ is what carries Alice's retained system in
trusted-source mode, where it is no output of $\Aev$. Taking $R'$ trivial
and $X=\idon{}$ in \eqref{eq:attsetpl} gives the scalar condition for the
lifted operator $\idon{\Qsys_{f^c}}\otimes\hat O_k$, so the set with
party-local entries is contained in the set with the same entries read as
scalars. The condition \eqref{eq:attsetpl} is that of an entry tagged with the
factor $f$ in \xref{foundations/def:round-model}\,(iii) at $S=\Qsys$, which the
consuming contract states at the same system. An \exttermref{../foundations/contract.pdf}{term-attack-sequence}{attack sequence} is
$\CList$-admissible (an \exttermref{../foundations/contract.pdf}{term-admissible-attack-sequence}{admissible attack sequence} for $\CList$) if
$\Aev_j$ lies in the corresponding set with
$E'=E_{j-1}$ and $E''=E'_j$ for every $j$.
\end{definition}
\claimmeta{\metalabel{Uses}\extref{../foundations/contract.pdf}{foundations/def:round-model}
\metasep \metalabel{Used by}\mref{Rem.}{rem:relation-symbols}, \mref{Thm.}{thm:licence}, \mref{Rem.}{rem:notlicensed}, \mref{Rem.}{rem:publishedlicence}, \mref{Def.}{def:chain}, \mref{Rem.}{rem:chain-admission}, \mref{Lem.}{lem:precomposition}, \mref{Thm.}{thm:two-sided}, \mref{Rem.}{rem:ordering}, \mref{Lem.}{lem:row-relaxation}
\metasep newly derived
\metasep accepted}

% claim reductions/rem:relation-symbols
\claimsummary{reductions/rem:relation-symbols}{The consuming contract writes the relations operator-valued at the channel level and as comparisons of real numbers at the untrusted declaration level; both spellings denote the same three relations. Every entry this module emits carries a real constant and is admitted under either.}
\claimbookmark{reductions/rem:relation-symbols}{Remark}{rem:relation-symbols}{Relation symbols at the channel and declaration levels}
\begin{remark}[Relation symbols at the channel and declaration levels]\label{rem:relation-symbols}
The consuming contract splits the relation symbols between two
\emph{levels}, not between the two source modes. At the \emph{channel} level,
the adjoint-form sets themselves, \xref{meat-core/def:roundmodel} writes
$\bowtie_k$ operator-valued in \emph{both} modes, comparing
$\Aev^{\dagger}(\cdot)$ against $c_k\idon{E'}$; that is the spelling the two
displays of Definition~\ref{def:attset} reproduce, with
$\bowtie_k\in\{\preceq,=,\succeq\}$. At the \emph{declaration and state}
level, \xref{meat-core/def:admissible-declaration}\,(vi) gives $\bowtie_k\in\{\preceq,=,\succeq\}$
with $\Theta_k$ Hermitian on $A$ in trusted-source mode, and
$\bowtie_k\in\{\le,=,\ge\}$ with $\Theta_k=c_k\in\mathbb{R}$ in
untrusted-source mode, where the entry is a comparison of real numbers;
\xref{protocol-compiler/def:homog} is on the same side, with the clause letters
\emph{(a) trusted-source} and \emph{(b) untrusted-source}. That definition
also fixes what \emph{licensed} means there: a licensed trusted entry declares
a real $c_k$ and the right-hand side $\Theta_k\coloneqq c_k\tau_A$ is
\emph{constructed} from it, while an entry that declares a Hermitian
$\Theta_k$ is a general entry, admissible \emph{input} but never licensed.
Nothing this module emits is touched by that: every entry it emits carries a
real scalar $c_k$ and is written $(\hat O_k,\bowtie_k,c_k)$
(Definition~\ref{def:licensed}, Theorem~\ref{thm:licence}\,(iv)). The two
spellings denote the same three relations at the two levels; this contract
states operator conditions throughout and so uses the operator spelling, and
everything it emits is admitted under either.

The entries it emits are those of Corollary~\ref{cor:entries}:
$\bowtie\,=\,\succeq$ with $c=0$ for every sector projector (its part~(i)),
the single equality entry $(\hat O,=,q_t)$ for a class-marginal entry where
the instance discharges it (part~(ii)), and the constant at which its
certificate is discharged for an instance's own or a caller's entry
(part~(iii)). Every one of those constants is exact as declared, and the only
rounding that can touch one is a caller's own constant off the declared
lattice, which is rounded in the direction that \emph{enlarges} the
admissible set (Requirement~\ref{req:caller-constant-rounding},
\xref{foundations/req:rounding-direction}).
\end{remark}
\claimmeta{\metalabel{Uses}\mref{Def.}{def:attset}
\metasep newly derived
\metasep accepted}

% claim reductions/def:licensed
\claimsummary{reductions/def:licensed}{A constraint list is licensed when every right-hand side is a real constant: untagged, as a multiple of Alice's declared marginal in trusted mode and a real number in untrusted mode, or party-local. A party-local entry acts on one factor of a product measured system, against the marginal present.}
\claimbookmark{reductions/def:licensed}{Definition}{def:licensed}{Licensed entries and the party-local form}
\begin{definition}[Licensed entries and the party-local form]\label{def:licensed}
The list $\CList$ is
\termdef{term-licensed}{licensed}
for the declared source mode when each of its entries either is untagged
with its right-hand side in the form that mode admits, or is in the
party-local form below. The untagged forms are $\Theta_k=c_k\,\tau_A$ with
$c_k\in\mathbb{R}$ in trusted-source mode, the state-level entry reading
$\tr_{\Qsys}[(\idon{A}\otimes\hat O_k)\,\omega]\bowtie_kc_k\,\tau_A$; and
$\Theta_k=c_k\in\mathbb{R}$ in untrusted-source mode, the entry reading
$\tr[(\idon{A}\otimes\hat O_k)\,\omega]\bowtie_kc_k$, a comparison of real
numbers (\xref{protocol-compiler/def:homog}). In both cases the untagged
entries are written $(\hat O_k,\bowtie_k,c_k)$.

The
\termdef{term-party-local-form}{party-local form}.
Where the measured system is declared as a product
$\Qsys=\bigotimes_{f}\Qsys_f$ of party-scoped factors, an entry may in
addition be
\termdef{term-scoped-to-a-factor}{scoped to a factor}
$f$. It is written $(\hat O_k,\bowtie_k,c_k;f)$ with
$\hat O_k\in\Lin{\Qsys_f}$ Hermitian and $c_k\in\mathbb{R}$, it carries the
party-local tag, and its state-level reading, in either source mode, is
\begin{equation}\label{eq:partylocal-entry}
  \tr_{\Qsys_f}\!\Big[\big(\idon{}\otimes(\hat O_k-c_k\,\idon{\Qsys_f})\big)
  \,\omega\Big]\;\bowtie_k\;0
  \qquad\text{as an operator on the systems of }\omega\text{ other than }
  \Qsys_f ,
\end{equation}
those systems being $A$, where the declaration carries one, and the factors
$\Qsys_{f'}$, $f'\ne f$. Equivalently
$\tr_{\Qsys_f}[(\idon{}\otimes\hat O_k)\omega]\bowtie_kc_k\tr_{\Qsys_f}\omega$:
the entry is homogeneous with respect to the marginal that is actually
there. The attack-set licence of Theorem~\ref{thm:licence} covers a
party-local entry only through a certificate on the factor's own
\termref{term-squashing-channel}{squashing channel} (its part (v)); the product structure that makes it
sound is Lemma~\ref{lem:party-local}.

\emph{What the tag does not change.} An entry without the tag is read in
the form of its source mode, exactly as above. With one factor, the
party-local reading in trusted-source mode is
$\tr_{\Qsys}[(\idon{A}\otimes\hat O_k)\omega]\bowtie_kc_k\,\omega_A$, which
under the consumer's marginal condition $\omega_A=\tau_A$ is the
trusted-source form; with one factor and $d_A=1$ it is the untrusted-source
comparison of real numbers. \emph{With one factor, in untrusted-source mode
and $d_A>1$ it is neither}: it reads
$\tr_{\Qsys}[(\idon{A}\otimes\hat O_k)\omega]\bowtie_kc_k\,\omega_A$, an
operator inequality on $A$ against the marginal that is actually there, and
the untrusted-source scalar entry is its trace
(Lemma~\ref{lem:party-local}\,(iii)), implied by it and not equivalent to
it. With more than one factor it implies the untagged reading of the same
$(\hat O_k,\bowtie_k,c_k)$ for the lifted operator $\idon{}\otimes\hat O_k$,
which is obtained from it by the partial trace over the other factors (and
over $A$ in untrusted-source mode; in trusted-source mode under the
consumer's marginal condition $\omega_A=\tau_A$), and in general it is not
implied by it.
\end{definition}
\claimmeta{\metalabel{Used by}\mref{Thm.}{thm:licence}, \mref{Def.}{def:chain}, \mref{Lem.}{lem:precomposition}, \mref{Thm.}{thm:two-sided}, \mref{Rem.}{rem:ordering}, \mref{Cor.}{cor:entries}, \mref{Rem.}{rem:entry-tags}, \mref{Thm.}{thm:emitted-product}
\metasep newly derived
\metasep accepted}

% claim reductions/lem:party-local
\claimsummary{reductions/lem:party-local}{For a product squashing channel, a spectral bound certified on one part's channel holds of every squashed attack as an operator inequality on all other systems, whose trace is the scalar entry. Without the product structure the operator form can fail while the scalar form holds.}
\claimbookmark{reductions/lem:party-local}{Lemma}{lem:party-local}{Party-local transport}
\claimleadin{%
A party-local entry (Definition~\ref{def:licensed}) is read against the
marginal that is actually present rather than a declared one. The lemma shows
that a certificate on the part's own channel makes that reading hold of every
squashed attack, and that it can fail once the channel correlates the parts.
}
\begin{lemma}[Party-local transport]\label{lem:party-local}
Let the measured system be a tensor product $\Bsys=\Bsys_f\otimes\Bsys_{f^c}$
of a \emph{part} $\Bsys_f$ and the product $\Bsys_{f^c}$ of all other parts,
let $\Lambda_f\colon\Tcl{\Bsys_f}\to\Tcl{\Qsys_f}$ and
$\Lambda_{f^c}\colon\Tcl{\Bsys_{f^c}}\to\Tcl{\Qsys_{f^c}}$ be channels, and
put $\Lambda\coloneqq\Lambda_f\otimes\Lambda_{f^c}$ and
$\Qsys\coloneqq\Qsys_f\otimes\Qsys_{f^c}$. Let $\hat O\in\Lin{\Qsys_f}$ be
Hermitian with a spectral pull-back certificate \emph{on the part},
\begin{equation}\label{eq:partcert}
  \Lambda_f^{\dagger}(\hat O)\;\bowtie\;c\,\idon{\Bsys_f},
  \qquad\bowtie\,\in\{\preceq,=,\succeq\},\quad c\in\mathbb{R} .
\end{equation}
Let $R$ and $R'$ be finite-dimensional systems, $E'$, $E''$
finite-dimensional, and $\Aev\in\mathrm{CPTP}(E',R\,\Bsys\,E'')$ \emph{any}
channel, with
$\Aev^{\sq}\coloneqq(\idmap_R\otimes\Lambda\otimes\idmap_{E''})\circ\Aev$.
Write $\hat D\coloneqq\hat O-c\,\idon{\Qsys_f}$.
\begin{enumerate}[label=(\roman*),leftmargin=2.2em]
\item \emph{(channel level)} for every $X\in\Pos{R'R\,\Qsys_{f^c}}$,
\begin{equation}\label{eq:partylocal-channel}
  \big(\idmap_{R'}\otimes\Aev^{\sq}\big)^{\dagger}
  \big(X\otimes\hat D\otimes\idon{E''}\big)\;\bowtie\;0
  \qquad\text{as an operator on }R'E' ,
\end{equation}
the tensor factors written in the order $R'R\,\Qsys_{f^c}$, $\Qsys_f$, $E''$.
\item \emph{(state level)} for every $\varrho\in\Pos{R'E'}$, the operator
$\omega\coloneqq\tr_{E''}\big[(\idmap_{R'}\otimes\Aev^{\sq})(\varrho)\big]$
on $R'R\,\Qsys$ satisfies
\begin{equation}\label{eq:partylocal-state}
  \tr_{\Qsys_f}\!\big[(\idon{R'R\,\Qsys_{f^c}}\otimes\hat D)\,\omega\big]
  \;\bowtie\;0
  \qquad\text{as an operator on }R'R\,\Qsys_{f^c},
\end{equation}
equivalently
$\tr_{\Qsys_f}[(\idon{}\otimes\hat O)\,\omega]\bowtie
c\,\tr_{\Qsys_f}\omega$: the right-hand side is $c$ times the marginal of
$\omega$ \emph{that is actually there}, and no marginal is assumed.
\item \emph{(the scalar entry is its trace)} taking the trace of
\eqref{eq:partylocal-state} gives
$\tr[(\idon{}\otimes\hat O)\,\omega]\bowtie c\,\tr\omega$, the scalar form.
With $R'$ trivial and $R=A$, $\Qsys_{f^c}$ trivial (one factor), \eqref{eq:partylocal-state}
reads $\tr_{\Qsys}[(\idon{A}\otimes\hat O)\omega]\bowtie c\,\omega_A$.
\item \emph{(the product structure is needed)} if $\Lambda$ is not a product
the conclusion can fail while the scalar one holds. Let $\Qsys_f$ and
$\Qsys_{f^c}$ be qubits, $R$, $R'$ trivial, $\hat O=\dyad{0}$, $c=\tfrac12$,
and let $\Lambda$ discard its input and prepare
$\tfrac12(\dyad{00}+\dyad{11})$. Then
$\Lambda^{\dagger}(\idon{}\otimes\hat O)=\tfrac12\idon{\Bsys}\succeq c\,\idon{}$,
so every squashed attack satisfies the scalar entry with $\succeq$; but
$\tr_{\Qsys_f}[(\idon{}\otimes\hat D)\omega]
=\tfrac14\dyad{0}-\tfrac14\dyad{1}$, which is not positive semidefinite.
\end{enumerate}
\end{lemma}
\claimmeta{\metalabel{Uses}\mref{Lem.}{lem:sep}
\metasep \metalabel{Used by}\mref{Thm.}{thm:licence}, \mref{Rem.}{rem:entry-tags}, \mref{Thm.}{thm:emitted-product}
\metasep newly derived
\metasep accepted}

\begin{proof}
(i) Let $\bowtie$ be $\succeq$. By Lemma~\ref{lem:sep}\,(ii), applied to the
product $\idmap_{R'R}\otimes\Lambda_{f^c}\otimes\Lambda_f\otimes\idmap_{E''}$,
\begin{equation*}
  \big(\idmap_{R'R}\otimes\Lambda\otimes\idmap_{E''}\big)^{\dagger}
  \big(X\otimes\hat D\otimes\idon{E''}\big)
  =\big(\idmap_{R'R}\otimes\Lambda_{f^c}\big)^{\dagger}(X)\;\otimes\;
   \Lambda_f^{\dagger}(\hat D)\;\otimes\;\idon{E''} .
\end{equation*}
The first factor is positive semidefinite: $X\succeq0$ and the adjoint of
the channel $\idmap_{R'R}\otimes\Lambda_{f^c}$ is a positive map, the
adjoint of a completely positive map being completely positive. The second
is $\Lambda_f^{\dagger}(\hat O)-c\,\Lambda_f^{\dagger}(\idon{\Qsys_f})
=\Lambda_f^{\dagger}(\hat O)-c\,\idon{\Bsys_f}\succeq0$ by unitality of
$\Lambda_f^{\dagger}$ and \eqref{eq:partcert}. A tensor product of positive
semidefinite operators is positive semidefinite. Finally
$(\idmap_{R'}\otimes\Aev^{\sq})^{\dagger}
=(\idmap_{R'}\otimes\Aev)^{\dagger}\circ
(\idmap_{R'R}\otimes\Lambda\otimes\idmap_{E''})^{\dagger}$, and
$(\idmap_{R'}\otimes\Aev)^{\dagger}$ is a positive map because
$\idmap_{R'}\otimes\Aev$ is a channel; this gives
\eqref{eq:partylocal-channel}. For $\preceq$ apply the argument to
$-\hat D$, and for $=$ both hold, the left-hand side of
\eqref{eq:partylocal-channel} then vanishing for every $X\succeq0$ and
hence, by linearity, for every $X$.

(ii) For $X\in\Pos{R'R\,\Qsys_{f^c}}$, by the definition of the adjoint,
\begin{equation*}
  \tr\!\Big[X\;\tr_{\Qsys_f}\!\big[(\idon{}\otimes\hat D)\,\omega\big]\Big]
  =\tr\!\big[(X\otimes\hat D\otimes\idon{E''})\,
     (\idmap_{R'}\otimes\Aev^{\sq})(\varrho)\big]
  =\tr\!\Big[\big(\idmap_{R'}\otimes\Aev^{\sq}\big)^{\dagger}
     \big(X\otimes\hat D\otimes\idon{E''}\big)\,\varrho\Big] ,
\end{equation*}
which by (i) and $\varrho\succeq0$ is $\ge0$ for $\succeq$, $\le0$ for
$\preceq$ and $0$ for $=$. A Hermitian operator $Y$ with $\tr[XY]\ge0$ for
every $X\succeq0$ is positive semidefinite, and
$\tr_{\Qsys_f}[(\idon{}\otimes\hat D)\omega]$ is Hermitian because $\hat D$
is and acts on the traced factor alone. The equivalent form is linearity:
$\tr_{\Qsys_f}[(\idon{}\otimes c\,\idon{\Qsys_f})\omega]
=c\,\tr_{\Qsys_f}\omega$.

(iii) The trace is a positive functional, and
$\tr\big[\tr_{\Qsys_f}[(\idon{}\otimes\hat O)\omega]\big]
=\tr[(\idon{}\otimes\hat O)\omega]$. The special case is a substitution.

(iv) $\Lambda^{\dagger}(Y)=\tr[Y\varsigma]\,\idon{\Bsys}$ for the prepared
state $\varsigma=\tfrac12(\dyad{00}+\dyad{11})$, so
$\Lambda^{\dagger}(\idon{}\otimes\dyad{0})=\tfrac12\idon{\Bsys}$. Every
squashed attack outputs $\omega=\varsigma$ on $\Qsys$, and
$\tr_{\Qsys_f}[(\idon{}\otimes(\dyad{0}-\tfrac12\idon{}))\varsigma]
=\tfrac12\big(\tfrac12\dyad{0}\big)+\tfrac12\big(-\tfrac12\dyad{1}\big)$,
with the traced qubit written second.
\end{proof}

% claim reductions/thm:licence
\claimsummary{reductions/thm:licence}{If each operator the squashing carries pulls back to a physical operator \ensuremath{\le}, = or \ensuremath{\ge} a multiple of the identity, every squashed physical attack respects the resulting constraint list in either source mode, factor by factor for per-factor bounds. Security against such attacks transfers with the same \ensuremath{\varepsilon}.}
\claimbookmark{reductions/thm:licence}{Theorem}{thm:licence}{The attack-set licence}
\claimleadin{%
The consuming engine may restrict the squashed attacks to the admissible set
of Definition~\ref{def:attset} only if every squashed physical attack lies in
it. The theorem proves this from one hypothesis about the device, a spectral
bound on each carried operator's physical pull-back, and uses nothing about
the attack beyond its being a channel.
}
\begin{mainclaim}
\begin{theorem}[The attack-set licence]\label{thm:licence}
Let $(M;O)\dsq_{\Lambda}(F;\hat O)$ with $\Lambda\colon\Tcl{\Bsys}\to
\Tcl{\Qsys}$, and suppose each carried operator has a \emph{spectral pull-back
certificate}: for each $k$ there are $\bowtie_k\in\{\preceq,=,\succeq\}$ and
$c_k\in\mathbb{R}$ with
\begin{equation}\label{eq:certificate}
  O_k\;=\;\Lambda^{\dagger}(\hat O_k)\;\bowtie_k\;c_k\,\idon{\Bsys}
  \qquad\text{as operators on }\Hs{\Bsys}.
\end{equation}
Then, with $\CList\coloneqq\{(\hat O_k,\bowtie_k,c_k)\}_{k=1}^{m}$:
\begin{enumerate}[label=(\roman*),leftmargin=3.0em]
\item \emph{(trusted-source mode)} for all finite-dimensional $E',E''$ and
  \emph{every} $\Aev\in\mathrm{CPTP}(E',\Bsys E'')$, the composite
  $\Aev^{\sq}\coloneqq(\Lambda\otimes\idmap_{E''})\circ\Aev$ lies in
  $\AttSet_{\CList}(E',\Qsys E'')$ of \eqref{eq:attset};
\item \emph{(untrusted-source mode)} for every
  $\Aev\in\mathrm{CPTP}(E',A\Bsys E'')$, the composite
  $\Aev^{\sq}\coloneqq(\idmap_A\otimes\Lambda\otimes\idmap_{E''})\circ\Aev$
  lies in $\AttSet_{\CList}(E',A\Qsys E'')$ of \eqref{eq:attsetun};
\item \emph{(the licensed reduction, in the declared mode)} consequently, for
  every attack sequence $(\Aev_j)_{j=1}^{n}$ against the physical protocol of
  the type \eqref{eq:attackchannel} for the declared source mode, the sequence
  $(\Aev^{\sq}_j)_j$ of the \emph{same} branch of \eqref{eq:composite} is
  $\CList$-admissible --- by (i) with the trusted branch of
  \eqref{eq:attackchannel}, by (ii) with the untrusted branch --- and produces
  the \emph{same} protocol output state, by the correspondingly named branch of
  Corollary~\ref{cor:sequential}; so security of the squashed round model
  against $\CList$-admissible sequences implies security of the physical
  protocol against all sequences of that type, with the same $\varepsilon$;
\item \emph{(licensed entries)} $\CList$ is licensed for either source mode in
  the sense of Definition~\ref{def:licensed}, with the same constants $c_k$;
  on every reachable state $\sigma$ the transported entries hold in the form
  with the marginal that is present,
  $\tr_{\Qsys}[(\idon{A}\otimes\hat O_k)\sigma]\bowtie_kc_k\,\sigma_A$, and
  so in the untrusted-source form
  $\tr[(\idon{A}\otimes\hat O_k)\sigma]\bowtie_kc_k$; and the trusted-source
  form $\tr_{\Qsys}[(\idon{A}\otimes\hat O_k)\sigma]\bowtie_kc_k\,\tau_A$
  holds on every reachable state whose marginal on $A$ is $\tau_A$.
\item \emph{(party-local entries, product measured system)} \emph{in place
  of the hypotheses above}, which type $\hat O_k$ on $\Qsys$ and certify it
  against $\Lambda$, let
  $\Bsys=\bigotimes_f\Bsys_f$, let $M^{f}\dsq_{\Lambda_f}F^{f}$ for every
  factor with $\Lambda_f\colon\Tcl{\Bsys_f}\to\Tcl{\Qsys_f}$, and put
  $M\coloneqq\{\bigotimes_fM^{f}_{v_f}\}$, $F\coloneqq\{\bigotimes_fF^{f}_{v_f}\}$
  and $\Lambda\coloneqq\bigotimes_f\Lambda_f$, so that $M\dsq_{\Lambda}F$
  (Proposition~\ref{prop:tensor}). Let each entry of $\CList$ be scoped to a
  factor $f_k$ with $\hat O_k\in\Lin{\Qsys_{f_k}}$ and a certificate
  \emph{on that factor's channel},
  $\Lambda_{f_k}^{\dagger}(\hat O_k)\bowtie_kc_k\,\idon{\Bsys_{f_k}}$. Then
  for every attack $\Aev$ of either type of \eqref{eq:attackchannel}, the
  squashed attack of the same branch of \eqref{eq:composite} satisfies
  \eqref{eq:attsetpl} for every $k$, and also the scalar condition of
  \eqref{eq:attset} or \eqref{eq:attsetun} for the lifted operator
  $\idon{\Qsys_{f_k^c}}\otimes\hat O_k$ with the same $c_k$; every
  \termref{term-squashed-sequence}{squashed sequence} is admissible for $\CList$ read with any subset of its entries
  tagged party-local, it produces the same protocol output state
  as the physical sequence, and the security implication of (iii) holds
  with the same $\varepsilon$. On every reachable state the entries hold in
  the form \eqref{eq:partylocal-entry}.
\end{enumerate}
\end{theorem}
\end{mainclaim}
\claimmeta{\metalabel{Uses}\mref{Lem.}{lem:sep}, \mref{Cor.}{cor:pullback}, \mref{Def.}{def:roundmodel}, \mref{Cor.}{cor:sequential}, \mref{Cor.}{cor:transported}, \mref{Prop.}{prop:tensor}, \mref{Def.}{def:attset}, \mref{Def.}{def:licensed}, \mref{Lem.}{lem:party-local}
\metasep \metalabel{Assumes}\aref{A-CONST}{ass:constdecl}
\metasep \metalabel{Used by}\mref{Rem.}{rem:notlicensed}, \mref{Rem.}{rem:gateweight}, \mref{Rem.}{rem:publishedlicence}, \mref{Thm.}{thm:chain}, \mref{Thm.}{thm:two-sided}, \mref{Rem.}{rem:ordering}, \mref{Def.}{def:reduction-instance}, \mref{Thm.}{thm:instrument-instance}, \mref{Cor.}{cor:entries}, \mref{Rem.}{rem:entry-tags}, \mref{Thm.}{thm:emitted-product}, \mref{Cor.}{cor:reference-companion}, \mref{Lem.}{lem:row-relaxation}, \mref{Req.}{req:licence-record}, \mref{Req.}{req:constraint-list-by-request}, \mref{Cor.}{cor:fssinstance}, \mref{Cor.}{cor:two-sector-decided}, \mref{Cor.}{cor:flat-companion}, \mref{Cor.}{cor:dom-chain}, \extref{../protocol-compiler/contract.pdf}{protocol-compiler/req:link-row-provenance}
\metasep newly derived
\metasep accepted
\metasep \metalabel{Oracles}\coderef{https://github.com/fraunhofersg/certiqs-key/blob/v0.1.0/tests/oracles/registry_reductions.py}{registry_reductions.py}}

\begin{proof}
(i) Fix $k$ and write $\Aev^{\sq\dagger}
=\Aev^{\dagger}\circ(\Lambda\otimes\idmap_{E''})^{\dagger}$, which is
\eqref{eq:pairing} applied twice. By Lemma~\ref{lem:sep}\,(ii),
\begin{equation}\label{eq:compositeadjoint}
  (\Lambda\otimes\idmap_{E''})^{\dagger}
  \big(\hat O_k\otimes\idon{E''}\big)
  =\Lambda^{\dagger}(\hat O_k)\otimes\idon{E''}
  =O_k\otimes\idon{E''} .
\end{equation}
Suppose $\bowtie_k$ is $\succeq$. From $O_k\succeq c_k\idon{\Bsys}$ we get
$O_k\otimes\idon{E''}\succeq c_k\,\idon{\Bsys E''}$, and $\Aev^{\dagger}$ is
positive and unital because $\Aev$ is completely positive and trace preserving;
hence
\begin{equation*}
  \Aev^{\sq\dagger}\big(\hat O_k\otimes\idon{E''}\big)-c_k\,\idon{E'}
  =\Aev^{\dagger}\big(O_k\otimes\idon{E''}-c_k\,\idon{\Bsys E''}\big)
  \;\succeq\;0 ,
\end{equation*}
using $\Aev^{\dagger}(\idon{\Bsys E''})=\idon{E'}$. For
$\bowtie_k\,=\,\preceq$ apply the same argument to
$c_k\idon{\Bsys}-O_k\succeq0$; for $\bowtie_k\,=\,=$ both directions hold, or
directly $\Aev^{\dagger}(c_k\idon{})=c_k\idon{E'}$. Since $k$ was arbitrary,
$\Aev^{\sq}\in\AttSet_{\CList}(E',\Qsys E'')$.

(ii) Identical, with \eqref{eq:compositeadjoint} replaced by
$(\idmap_A\otimes\Lambda\otimes\idmap_{E''})^{\dagger}
\big(\idon{A}\otimes\hat O_k\otimes\idon{E''}\big)
=\idon{A}\otimes O_k\otimes\idon{E''}$, again Lemma~\ref{lem:sep}\,(ii), and
with $\idon{A}\otimes O_k\otimes\idon{E''}\bowtie_k c_k\,\idon{A\Bsys E''}$
following from \eqref{eq:certificate} by tensoring with identities.

(iii) In trusted-source mode, admissibility is (i) applied per round with
$E'=E_{j-1}$, $E''=E'_j$ and the trusted branch of \eqref{eq:composite}; in
untrusted-source mode it is (ii) with the untrusted branch, the lifted operator
$\idon{A}\otimes\hat O_k\otimes\idon{E'_j}$ of \eqref{eq:attsetun} being
exactly the one that branch's composite is tested against. The equality of
output states is the correspondingly typed branch of
Corollary~\ref{cor:sequential}, whose hypothesis is exactly $M\dsq_{\Lambda}F$
in both modes; and the security implication is the last sentence of that
corollary, the supremum over the $\CList$-admissible sequences dominating the
value at $(\Aev^{\sq}_j)_j$. The two branches of \eqref{eq:seqcover} are
$\AttSet_{\CList}(E_{j-1},\Qsys E'_j)$ and
$\AttSet_{\CList}(E_{j-1},A\Qsys E'_j)$ respectively, by (i) and (ii).

(iv) The right-hand sides are the scalars $c_k$ of \eqref{eq:certificate},
which are declared constants of the reduction and not functions of any observed
quantity (Assumption~\ref{ass:constdecl}); reading $c_k$ as $c_k\tau_A$ in
trusted-source mode and as the real number $c_k$ in untrusted-source mode is
Definition~\ref{def:licensed} in the two modes. That the entries hold on every
reachable state $\sigma$ in the form with the marginal present is
Corollary~\ref{cor:transported}\,(i) applied to the certificate
\eqref{eq:certificate}, which is case (i) or (iii) of
Corollary~\ref{cor:pullback}, read on the side it bounds, with $R=A$: that
corollary delivers
$\tr_{\Qsys}[(\idon{A}\otimes\hat O_k)\sigma]\bowtie_kc_k\,\sigma_A$, and its
scalar form, the trace of both sides with $\tr\sigma_A=1$, is the
untrusted-source form. In trusted-source mode the identification of the
right-hand side with $c_k\tau_A$ is the \emph{declaration's} and not this
contract's: it holds where $\sigma_A=\tau_A$, which is the consumer's
marginal condition $\omega_A=\tau_A$ (defined in
\xref{meat-core/def:hqkd-list} and asserted by the consuming contract under
its own source-replacement assumption), and nothing in the hypotheses of this
theorem restricts the marginal on $A$.

(v) Fix $k$, write $f=f_k$, $\Lambda_{f^c}\coloneqq\bigotimes_{f'\ne f}
\Lambda_{f'}$, and reorder tensor factors so that
$\Lambda=\Lambda_f\otimes\Lambda_{f^c}$. Lemma~\ref{lem:party-local}\,(i),
with $R=A$ in untrusted-source mode and $R$ trivial in trusted-source mode,
and with the certificate on the factor as its hypothesis
\eqref{eq:partcert}, is \eqref{eq:attsetpl} for the squashed attack, for
every reference system $R'$. Its part (i) at $R'$ trivial and $X=\idon{}$,
with unitality of the adjoint of the squashed attack, a channel, gives the scalar
condition for the lifted operator; alternatively
$\Lambda^{\dagger}(\idon{}\otimes\hat O_k)
=\idon{\Bsys_{f^c}}\otimes\Lambda_f^{\dagger}(\hat O_k)\bowtie_k
c_k\,\idon{\Bsys}$ by Lemma~\ref{lem:sep}\,(ii) and unitality, which is a
certificate \eqref{eq:certificate} for $\Lambda$, so (i) and (ii) apply to
it. Admissibility of the squashed sequence is this per round, with
$E'=E_{j-1}$ and $E''=E'_j$. The equality of output states and the security
implication are Corollary~\ref{cor:sequential} for the channel $\Lambda$,
whose hypothesis $M\dsq_{\Lambda}F$ is Proposition~\ref{prop:tensor}; the
supremum over the admissible sequences dominates the value at the squashed
one, whichever entries are tagged. The state-level form is
Lemma~\ref{lem:party-local}\,(ii), with $R'=A$ carrying Alice's retained
system in trusted-source mode.
\end{proof}

% claim reductions/rem:notlicensed
\claimsummary{reductions/rem:notlicensed}{The licence permits one thing, restricting the squashed attacks to the declared admissible set. It licenses no trusted marginal, no constraint on Eve's side information, no data-dependent or multi-round constraint, no reuse under another proof technique, and nothing about the device.}
\claimbookmark{reductions/rem:notlicensed}{Remark}{rem:notlicensed}{Scope of the attack-set licence}
\begin{remark}[Scope of the attack-set licence]
\label{rem:notlicensed}
Theorem~\ref{thm:licence} licenses exactly one thing: restricting the squashed
round model's attacks to the declared admissible set of
Definition~\ref{def:attset}. It does \emph{not} license, and this contract
asserts none of the following.
\begin{enumerate}[label=(\alph*),leftmargin=2.4em]
\item \emph{A trusted marginal.} Nothing here licenses declaring $\tau_A$. That
  is a separate input with a separate attestation on the consumer's side, and a
  licensed $\CList$ never substitutes for it: a constraint entry restricts what
  \emph{Eve's channel} may do, whereas $\tau_A$ asserts that a register never
  reached her at all.
\item \emph{Constraints on the purifying register.} Every entry is an operator
  on $\Qsys$ (resp.\ $A\Qsys$); no constraint on Eve's side information is
  expressible in this format or admitted by it.
\item \emph{Data-dependent constraints.} The constants $c_k$ are declared
  before the run (Assumption~\ref{ass:constdecl}); an entry whose $c_k$ encoded
  a statistic of the current run would be unsound here regardless of
  provenance, and such an entry is refused
  (Requirement~\ref{req:declared-constants-only}).
\item \emph{Joint multi-round constraints.} Admissibility is per round
  (Definition~\ref{def:attset}); Theorem~\ref{thm:licence} is proved per round,
  and no constraint on the joint $n$-round state is licensed or expressible.
\item \emph{Reuse under another proof technique.} The restriction is compatible
  with the accumulation route the consuming contract imports. It is \emph{not}
  automatically compatible with a postselection-based reduction: the published
  framework records that the analogous restriction ``cannot be
  straightforwardly used with the postselection technique''
  (\cite[Rem.~9.4]{tupkary2026decoy}). A module importing $\CList$ into a
  different technique must re-establish the licence there; the artefact names
  the technique the licence was issued for
  (Requirement~\ref{req:licence-record}).
\item \emph{Anything about the device.} Whether the receiver actually satisfies
  the hypothesis \eqref{eq:certificate} --- concretely, whether
  $\blk{M}{\gtN}_o\succeq p_o\idon{}$ holds of the physical detector --- is
  \emph{not} proved here and cannot be: it is
  Assumption~\ref{ass:recv}, discharged by the accepted receiver contract for
  the device that contract describes, under its own ledger.
\end{enumerate}
\end{remark}
\claimmeta{\metalabel{Uses}\mref{Def.}{def:attset}, \mref{Thm.}{thm:licence}
\metasep \metalabel{Assumes}\aref{A-CONST}{ass:constdecl}, \aref{A-RECV}{ass:recv}
\metasep \metalabel{Used by}\mref{Req.}{req:licence-record}, \mref{Req.}{req:declared-constants-only}
\metasep newly derived
\metasep accepted}

% claim reductions/rem:gateweight
\claimsummary{reductions/rem:gateweight}{In untrusted-source mode, or where Alice's system has dimension one, the consumer's mechanical check of the licensed class degenerates and its gate rests entirely on the supplier's licence. The licence is therefore a theorem with a named hypothesis rather than an attestation.}
\claimbookmark{reductions/rem:gateweight}{Remark}{rem:gateweight}{Weight of the licence in the consumer's gate}
\begin{remark}[Weight of the licence in the consumer's gate]
\label{rem:gateweight}
The consuming contract gates a key-length claim on a non-empty $\CList$ behind
two checks: (i) that the list lies in the licensed class of the declared source
mode, which is mechanical; and (ii) that the supplier's attack-set licence is
present, which no program can verify. It also records that check (i)
\emph{degenerates}: at $d_A=1$ every $1\times1$ Hermitian right-hand side is a
multiple of $\tau_A$, and in untrusted-source mode the type of an entry is
scalar by construction, so part (i) rejects nothing in untrusted-source mode
and nothing a producer cannot trivially avoid at $d_A=1$, and in both cases
the gate rests entirely on part (ii) (\xref{meat-core/rem:cgate}). This is the reason
Theorem~\ref{thm:licence} is stated as a theorem with an explicit hypothesis
rather than as an attestation: in the cases that matter, the attestation
\emph{is} the security argument, and the only defensible form for it is a
proof plus a named hypothesis that some other accepted contract discharges.
The \termref{term-licence-record}{licence record} names both, and the artefact that discharges the
hypothesis (Requirement~\ref{req:licence-record}).
\end{remark}
\claimmeta{\metalabel{Uses}\mref{Thm.}{thm:licence}
\metasep \metalabel{Used by}\mref{Req.}{req:licence-record}
\metasep newly derived
\metasep accepted}

% claim reductions/rem:publishedlicence
\claimsummary{reductions/rem:publishedlicence}{The adjoint format of the admissible set is the published squashed-attack-set format, and the published inclusion is the special case of \hyperlink{claim-reductions/thm:licence}{Theorem~\ref*{thm:licence}} (i) at one flag-state operator with constant zero. The licence generalises it to any certified carried operator in either source mode.}
\claimbookmark{reductions/rem:publishedlicence}{Remark}{rem:publishedlicence}{The published squashed-attack-set format}
\begin{remark}[The published squashed-attack-set format]
\label{rem:publishedlicence}
The adjoint format of \eqref{eq:attset} is the published squashed-attack-set
format: \cite[Lemma~9.4, Eq.~(151)]{tupkary2026decoy} constructs, for a
flag-state-squashed measurement, the set
$\AttSet_j\coloneqq\{\Aev^{\mathrm{Sq}}_j:
\Aev^{\mathrm{Sq}\dagger}_j(F^{(Q)}_o\otimes\idon{E_j})\succeq
\lambda_{\min}\Aev^{\mathrm{Sq}\dagger}_j(\Pi_{Sq}\otimes\idon{E_j})\}$ from
the physical operator inequality
$M^{(B)}_o\succeq\lambda_{\min}\Pi_{m>N_{\mathrm{FSS}}}$, the projector being
onto the photon numbers above the cutoff $N_{\mathrm{FSS}}$ of the flag-state
squasher, and
proves the inclusion $\AttSet_j\supseteq\Lambda\circ
\mathrm{CPTP}(E_{j-1},\Bsys_jE_j)$; its Eq.~(152) is the dimension-independent
marginal form of the same set, and the modelling step of restricting the
attack-channel set at all is its Def.~8.3 with Rem.~8.4. That inclusion is the
special case of Theorem~\ref{thm:licence}\,(i) at
$\hat O=F_o-\lambda_{\min}\Pi_{Sq}$, $c=0$; that theorem is the general
statement, for an arbitrary carried operator with a spectral pull-back
certificate, in either source mode, with the round-model bookkeeping of
Corollary~\ref{cor:sequential} attached. It is proved there from
\eqref{eq:pairing}, positivity and unitality, and appeals to no published
result.
\end{remark}
\claimmeta{\metalabel{Uses}\mref{Lem.}{lem:sep}, \mref{Cor.}{cor:sequential}, \mref{Def.}{def:attset}, \mref{Thm.}{thm:licence}
\metasep newly derived
\metasep accepted}

\section{Typed links and reduction chains}\label{group-reductions/def:link}
\sectionintro{The artefact is a chain of \termref{term-typed-link}{typed links}, each carrying enough to be re-checked: its \termref{term-channel-witness}{channel witness}, \termref{term-carried-operators}{carried operators}, \termref{term-reachable-set-statement}{reachable-set statement}, transported constraint entries and, where there are entries, a \termref{term-licence-record}{licence record}. A chain is a source segment of the transmitter followed by composable receiver-stage links, and its \termref{term-receiver-composite}{receiver composite} is again a \termref{term-squashing-relation}{squashing relation} to which output-state equality, security transfer and the licence apply (\hyperlink{claim-reductions/thm:chain}{Theorem~\ref*{thm:chain}}). A declaration with a non-empty source segment and constraint entries is admitted only with licensed entries and in trusted-source mode. An already-finite receiver reduces by one identity link.}
\sectionimports{\extref{../transmitter/contract.pdf}{transmitter/def:source-link}, \extref{../transmitter/contract.pdf}{transmitter/def:source-segment}}

% claim reductions/def:link
\claimsummary{reductions/def:link}{A typed link is a source link of the transmitter's type or a receiver-stage tuple of ten slots: stage, kind, source family, target, channel witness, carried operators, reachable-set statement, transported entries, licence record and status. Each link carries enough to be re-checked.}
\claimbookmark{reductions/def:link}{Definition}{def:link}{Typed link}
\claimleadin{%
What this module emits is not a squashed POVM but a chain of typed links,
each carrying enough to re-check it, in an order that is data rather than
convention. This definition states a link. How
\exttermref{../transmitter/contract.pdf}{term-source-link}{source links} and receiver-stage links compose is
Definition~\ref{def:chain}, with its condition (C3) for the order the record
carries and its admission rule for which such chains are admitted at all; Remark~\ref{rem:ordering} gives the
reason, Theorem~\ref{thm:chain} what the receiver stage composes to, and
Theorem~\ref{thm:two-sided} what the two stages compose to.
}
\begin{definition}[Typed link]\label{def:link}
A
\termdef{term-typed-link}{typed link}
carries a stage $\textsf{stage}\in\{\textsf{source},\textsf{receiver}\}$.
A link at the stage \textsf{source} is a \exttermref{../transmitter/contract.pdf}{term-source-link}{source link} in the sense of
\xref{transmitter/def:source-link}: its descriptors are Alice-side, a
\exttermref{../transmitter/contract.pdf}{term-transmitted-system}{transmitted system} with a signal ensemble on either end and a channel
witness between them, and no condition on it mentions the receiver's
systems, \termref{term-outcome-alphabet}{outcome alphabets} or measurements.
\emph{This module reads that type as the transmitter states it}: the channel
witness is supplied either
\emph{as data}, in which case the source identity is decided on declared
numbers, or \emph{as a \exttermref{../transmitter/contract.pdf}{term-theorem-constructor}{theorem constructor}}, in which case the link
carries an \exttermref{../transmitter/contract.pdf}{term-obligation-record}{obligation record}, namely the triple of the kind, the id of
the soundness claim and the assumptions the obligation carries, which
\xref{transmitter/def:source-segment}\,(A1) defines and which this module
neither computes nor re-decides, together with the member of its
\exttermref{../transmitter/contract.pdf}{term-declared-descriptor-family}{declared descriptor family} that its descriptors are; and a source link
carries in addition that family, an optional
\exttermref{../transmitter/contract.pdf}{term-retained-register-sector-map}{retained-register sector map}, its
\exttermref{../transmitter/contract.pdf}{term-transported-constraint-entries}{transported constraint entries} with its licence slot, and its status.
Which kinds are admitted at all, and which components a consumer decides
\emph{on data} (the emitted \exttermref{../transmitter/contract.pdf}{term-soundness-registry}{soundness registry}, the triple compared
against it entry by entry with those assumptions demanded in the ledger, and
the provenance rule for a non-empty segment), is
\xref{transmitter/def:source-segment}, conditions (A1) to (A3); this module
re-decides none of them, adds none, and adds no source-stage kind. A
\termdef{term-receiver-stage-typed-link}{receiver-stage typed link}
$\Lnk$ is a tuple
\begin{equation}\label{eq:link}
  \Lnk=\Big(\textsf{stage},\ \textsf{kind},\
  \big(\Bsys^{\mathrm{in}},M^{\mathrm{in}}\big),\
  \big(\Bsys^{\mathrm{out}},M^{\mathrm{out}}\big),\
  \Lambda_{\Lnk},\
  \big(\hat O^{\mathrm{out}},O^{\mathrm{in}}\big),\
  \Srch_{\Lnk},\
  \CList_{\Lnk},\
  \textsf{licence},\ \textsf{status}\Big)
\end{equation}
in which: $\textsf{stage}=\textsf{receiver}$;
\textsf{kind} is one of the receiver-stage link kinds the module's interface
declares (the core's kinds, and those a submodule's interface adds);
$(\Bsys^{\mathrm{in}},M^{\mathrm{in}})$ is the
\termdef{term-source-family}{source family}
--- the system and POVM the link starts from --- and
$(\Bsys^{\mathrm{out}},M^{\mathrm{out}})$ its \emph{target};
$\Lambda_{\Lnk}$ is the
\termdef{term-channel-witness}{channel witness},
with
$\big(M^{\mathrm{in}};O^{\mathrm{in}}\big)\dsq_{\Lambda_{\Lnk}}
\big(M^{\mathrm{out}};\hat O^{\mathrm{out}}\big)$;
$\Srch_{\Lnk}$ is the
\termdef{term-reachable-set-statement}{reachable-set statement}
of Definition~\ref{def:sound} for $\Lambda_{\Lnk}$ and the declared family of
inputs; $\CList_{\Lnk}$ is the list of transported constraint entries
$(\hat O_k,\bowtie_k,c_k)$ with, per entry, the \termref{term-certificate}{certificate} case of
Corollary~\ref{cor:pullback} that discharges it; \textsf{licence} is the
\termdef{term-licence-record}{licence record}
of the module's outputs, present exactly when $\CList_{\Lnk}\ne\varnothing$;
and $\textsf{status}\in\{\textsf{exact},\textsf{certified\_outer},
\textsf{exploratory\_numeric}\}$, with \textsf{exact} meaning that
$\Lambda_{\Lnk}^{\dagger}(M^{\mathrm{out}}_v)=M^{\mathrm{in}}_v$ holds with no
error term.

\emph{The count.} A receiver-stage link has ten slots. Six of
them are the required fields --- source family, target, channel witness,
reachable-set statement, transported constraint entries, licence record ---
and the remaining four are \textsf{stage}, \textsf{kind}, the carried
operator pair, and \textsf{status}. Wherever this contract says ``the six
required fields'' it means the first group; the tuple \eqref{eq:link} is the
whole of what an implementer binds.
\end{definition}
\claimmeta{\metalabel{Uses}\mref{Cor.}{cor:pullback}, \mref{Def.}{def:certificate}, \mref{Def.}{def:sound}, \extref{../transmitter/contract.pdf}{transmitter/def:source-link}, \extref{../transmitter/contract.pdf}{transmitter/def:source-segment}
\metasep \metalabel{Used by}\mref{Def.}{def:chain}, \mref{Thm.}{thm:chain}, \mref{Cor.}{cor:identity-instance}, \mref{Thm.}{thm:emit}, \mref{Cor.}{cor:reference-companion}, \mref{Def.}{def:domination-link}, \mref{Thm.}{thm:noise-channel}, \mref{Cor.}{cor:dom-chain}
\metasep newly derived
\metasep accepted}

% claim reductions/def:chain
\claimsummary{reductions/def:chain}{A reduction chain pairs a source segment with composable, tuple-chained receiver-stage links, and its composite is the pair of their composites. A declaration with a non-empty source segment and constraints is admitted only with licensed entries and in trusted-source mode.}
\claimbookmark{reductions/def:chain}{Definition}{def:chain}{Reduction chain}
\claimleadin{%
Remark~\ref{rem:ordering} gives the reason for the stage order (C3), and for
the admission rule of the definition the reason it gives is
Lemma~\ref{lem:precomposition} and its counterexample.
}
\begin{definition}[Reduction chain]\label{def:chain}
A
\termdef{term-reduction-chain}{reduction chain}
is a pair $\redChain=(\redChain_{\mathrm{src}},\redChain_{\mathrm{rec}})$ of
a \exttermref{../transmitter/contract.pdf}{term-source-segment}{source segment} and a
\termdef{term-receiver-segment}{receiver segment}.
The source segment
$\redChain_{\mathrm{src}}=(\Xi_1,\dots,\Xi_{L_{\mathrm s}})$, $L_{\mathrm s}\ge0$,
is a source segment in the sense of \xref{transmitter/def:source-segment}
(written $\Chsrc$ there): a list of \exttermref{../transmitter/contract.pdf}{term-source-link}{source links}
(\xref{transmitter/def:source-link}) satisfying that definition's
\exttermref{../transmitter/contract.pdf}{term-interface-conditions}{interface conditions} (S1)--(S4) --- descriptor match, index
preservation, the source identity where checkable, and the slots --- every
descriptor of which is Alice-side. The receiver segment
$\redChain_{\mathrm{rec}}=(\Lnk_1,\dots,\Lnk_k)$, $k\ge1$, is a list of
receiver-stage typed links (Definition~\ref{def:link}) that is
\begin{enumerate}[label=(C\arabic*),leftmargin=3.0em]
\item
  \termdef{term-composable}{composable}:
  $(\Bsys^{\mathrm{out}}_i,M^{\mathrm{out}}_i)
  =(\Bsys^{\mathrm{in}}_{i+1},M^{\mathrm{in}}_{i+1})$ for $1\le i<k$;
\item
  \termdef{term-tuple-chained}{tuple-chained}:
  $\hat O^{\mathrm{out}}_i=O^{\mathrm{in}}_{i+1}$ for
  $1\le i<k$, so that each declared operator on the final target has a
  pull-back at every intermediate stage;
\item
  \termdef{term-stage-ordered}{stage-ordered}:
  every \textsf{source}-stage link precedes every \textsf{receiver}-stage
  link, which the pair form makes structural.
\end{enumerate}
Its
\termdef{term-composite}{composite}
is the \emph{pair}
\begin{equation}\label{eq:chaincomposite}
  \big(\Xi_{\redChain},\ \Lambda_{\redChain}\big),
  \qquad
  \Xi_{\redChain}\coloneqq\Xi_{L_{\mathrm s}}\circ\dots\circ\Xi_1 ,
  \qquad
  \Lambda_{\redChain}\coloneqq\Lambda_{\Lnk_k}\circ\dots\circ\Lambda_{\Lnk_1},
\end{equation}
of the
\termdef{term-source-composite}{source composite}
--- the composite the transmitter's (S3) writes $\Xi$, the identity when
$L_{\mathrm s}=0$ --- and the
\termdef{term-receiver-composite}{receiver composite}.
The two act on different systems and no condition above relates them:
$\Xi_{\redChain}$ carries the \exttermref{../transmitter/contract.pdf}{term-virtual-transmitted-system}{virtual transmitted system} to the
transmitted one, on Alice's side, and $\Lambda_{\redChain}$ the
\termref{term-measured-system}{measured system} to the target, on the receiver's. What the pair
composes to \emph{as a reduction} is Theorem~\ref{thm:two-sided}, and
Theorem~\ref{thm:chain} is what the receiver composite alone delivers.

\emph{Which chains are admitted.} The rule is a condition on a
\emph{declaration} that carries both a non-empty source segment and a
constraint list: such a declaration is
\termdef{term-admitted}{admitted}
only where every entry of the list is \termref{term-licensed}{licensed} in the sense of
Definition~\ref{def:licensed} --- a real constant $c_k$ as the right-hand
side, the entry tagged or untagged --- \emph{and only in
\exttermref{../transmitter/contract.pdf}{term-trusted-source-mode}{trusted-source mode}}, Theorem~\ref{thm:two-sided} being stated there
alone, so that a declaration pairing a non-empty source segment with
\exttermref{../transmitter/contract.pdf}{term-untrusted-source-mode}{untrusted-source mode} is refused rather than admitted in silence. An
entry declaring a Hermitian right-hand side on Alice's retained system, a
\emph{general} entry in the sense of Definition~\ref{def:attset}, is
\emph{refused} beside one. The rule
is decided by the consumer that assembles those two.

\emph{Where the source segment is empty}, $L_{\mathrm s}=0$, the source
composite is the identity, $\Avirt=\Asig$, the \exttermref{../transmitter/contract.pdf}{term-virtual-ensemble}{virtual ensemble} is the
specified one, and the pair \eqref{eq:chaincomposite} carries exactly the
receiver composite: Alice's declared states are the ones the security
statement is about, and the receiver segment is the whole of the chain's
content (Theorem~\ref{thm:chain}).
\end{definition}
\claimmeta{\metalabel{Uses}\mref{Def.}{def:attset}, \mref{Def.}{def:licensed}, \mref{Def.}{def:link}, \extref{../transmitter/contract.pdf}{transmitter/def:source-link}, \extref{../transmitter/contract.pdf}{transmitter/def:source-segment}
\metasep \metalabel{Used by}\mref{Rem.}{rem:chain-admission}, \mref{Thm.}{thm:chain}, \mref{Ex.}{ex:identity}, \mref{Rem.}{rem:ordering}, \mref{Cor.}{cor:identity-instance}, \mref{Thm.}{thm:emit}, \mref{Cor.}{cor:reference-companion}, \mref{Req.}{req:identity-link-recorded}, \mref{Cor.}{cor:dom-chain}
\metasep newly derived
\metasep implemented (effective: accepted)
\metasep \metalabel{Code}\coderef{https://github.com/fraunhofersg/certiqs-key/blob/v0.1.0/src/certiqs_key/reductions/links.py}{certiqs_key.reductions.links:Chain}}

% claim reductions/rem:chain-admission
\claimsummary{reductions/rem:chain-admission}{This module states the admission rule of a chain and decides no instance of it; the protocol compiler decides both its halves. Every entry this module emits carries a real constant and every artefact an empty source segment, so no emitted artefact can violate the rule.}
\claimbookmark{reductions/rem:chain-admission}{Remark}{rem:chain-admission}{The admission rule and what this module emits}
\begin{remark}[The admission rule and what this module emits]\label{rem:chain-admission}
This module states the admission rule of Definition~\ref{def:chain} and
decides no instance of it. It declares no general entry and emits none, every
entry it emits carrying a real scalar (Definition~\ref{def:attset}), and every
artefact it emits has an empty source segment; so no artefact of this module
can violate the rule. The rule is decided where a declaration's source
segment is assembled with its constraint list, which is the protocol
compiler's chain condition (\xref{protocol-compiler/def:chaincond}). It
decides both halves of the rule: the \emph{entry} half, a general entry
beside a non-empty source segment, and the \emph{mode} half, a declaration
carrying a non-empty source segment in untrusted-source mode, each refused by
name there (\xref{protocol-compiler/req:chain-refusals}).
\end{remark}
\claimmeta{\metalabel{Uses}\mref{Def.}{def:attset}, \mref{Def.}{def:chain}
\metasep newly derived
\metasep accepted}

% claim reductions/thm:chain
\claimsummary{reductions/thm:chain}{The receiver composite of a chain is again a squashing relation with the composed pull-backs, so output-state equality, security transfer and the licence apply to it at every length of the source segment. Its status is exact whenever every link's is.}
\claimbookmark{reductions/thm:chain}{Theorem}{thm:chain}{Composition of a chain}
\claimleadin{%
The order of the source and receiver stages, which this theorem does not
read, is the subject of Remark~\ref{rem:ordering}.
}
\begin{theorem}[Composition of a chain]\label{thm:chain}
Let $\redChain=(\redChain_{\mathrm{src}},\redChain_{\mathrm{rec}})$ be a
reduction chain with receiver segment
$\redChain_{\mathrm{rec}}=(\Lnk_1,\dots,\Lnk_k)$ and receiver composite
$\Lambda_{\redChain}$ of \eqref{eq:chaincomposite}. Then:
\begin{enumerate}[label=(\roman*),leftmargin=3.0em]
\item $\big(M^{\mathrm{in}}_1;O^{\mathrm{in}}_1\big)
  \dsq_{\Lambda_{\redChain}}
  \big(M^{\mathrm{out}}_k;\hat O^{\mathrm{out}}_k\big)$, and
  $\Lambda_{\redChain}^{\dagger}
  =\Lambda_{\Lnk_1}^{\dagger}\circ\dots\circ\Lambda_{\Lnk_k}^{\dagger}$;
\item Theorem~\ref{thm:factorization}, Theorem~\ref{thm:transfer} and
  Corollary~\ref{cor:sequential} apply to the composite: the physical protocol
  with $M^{\mathrm{in}}_1$ and the final \termref{term-squashed-protocol}{squashed protocol} with
  $M^{\mathrm{out}}_k$ have the same output state, round by round and
  adaptively;
\item an entry $(\hat O_k,\bowtie_k,c_k)$ declared on the \emph{final} target
  is sound for the receiver composite when the certificate
  \eqref{eq:certificate} holds for its pull-back
  $\Lambda_{\redChain}^{\dagger}(\hat O_k)$, which (C2) computes link by link; and
  Theorem~\ref{thm:licence} then applies with $\Lambda_{\redChain}$;
\item the composite status is \textsf{exact} whenever every link's is, and no
  other composite value is asserted.
\end{enumerate}
Nothing above reads the source segment: (i)--(iv) are statements about the
receiver composite $\Lambda_{\redChain}$ and hold at every length
$L_{\mathrm s}\ge0$. What the \emph{pair} \eqref{eq:chaincomposite} composes to
as a reduction is Theorem~\ref{thm:two-sided}, under the source segment's own
interface conditions and the admission rule of Definition~\ref{def:chain}, and
this theorem does not restate it. Where the source segment is empty the source
composite is the identity and the receiver composite is the whole of the
chain's reduction.
\end{theorem}
\claimmeta{\metalabel{Uses}\mref{Thm.}{thm:factorization}, \mref{Cor.}{cor:sequential}, \mref{Thm.}{thm:transfer}, \mref{Cor.}{cor:transported}, \mref{Prop.}{prop:preorder}, \mref{Thm.}{thm:licence}, \mref{Def.}{def:link}, \mref{Def.}{def:chain}
\metasep \metalabel{Used by}\mref{Rem.}{rem:why-exact}, \mref{Ex.}{ex:identity}, \mref{Cor.}{cor:identity-instance}, \mref{Thm.}{thm:emit}, \mref{Cor.}{cor:emitted-identity}, \mref{Cor.}{cor:reference-companion}, \mref{Req.}{req:mixed-status-chain}, \mref{Req.}{req:identity-link-recorded}, \mref{Cor.}{cor:dom-chain}, \extref{../protocol-compiler/contract.pdf}{protocol-compiler/req:link-row-provenance}
\metasep newly derived
\metasep implemented (effective: accepted)
\metasep \metalabel{Code}\coderef{https://github.com/fraunhofersg/certiqs-key/blob/v0.1.0/src/certiqs_key/reductions/links.py}{certiqs_key.reductions.links:Chain}
\metasep \metalabel{Oracles}\coderef{https://github.com/fraunhofersg/certiqs-key/blob/v0.1.0/tests/oracles/registry_reductions.py}{registry_reductions.py}}

\begin{proof}
(i) is Proposition~\ref{prop:preorder} applied $k-1$ times, the tuple chaining
(C2) being exactly the hypothesis that makes the operator tuples compose. (ii)
is those results applied to the relation of (i), which is all they need. (iii)
is Corollary~\ref{cor:transported} and Theorem~\ref{thm:licence} applied to the
composite relation; the pull-back is computed as
$\Lambda_{\Lnk_1}^{\dagger}\big(\dots\Lambda_{\Lnk_k}^{\dagger}
(\hat O_k)\dots\big)$ by (i). (iv): a composition of exact identities is an
exact identity, by (i) applied to identities carrying no error term. The
remaining half is not proved and is not claimed: a chain containing a link
that is not \textsf{exact} is declined rather than assigned a composite
status, by the rules of Requirements~\ref{req:no-approximate-relations}
and~\ref{req:mixed-status-chain}.
The closing paragraph is Definition~\ref{def:chain}: the composite is defined
there as the pair \eqref{eq:chaincomposite}, whose second member is the
receiver segment's and is the only member any of (i)--(iv) reads --- each of
them quantifies over the links $\Lnk_1,\dots,\Lnk_k$ and over nothing else ---
and with $L_{\mathrm s}=0$ its first member is the identity.
\end{proof}

% claim reductions/ex:identity
\claimsummary{reductions/ex:identity}{An already-finite receiver reduces by one identity link (the identity channel or the sector pinching) with the same POVM, no constraint entries and exact status; the chain theorem returns the physical model unchanged. No input the module's interface declares carries such a receiver.}
\claimbookmark{reductions/ex:identity}{Example}{ex:identity}{The identity link}
\begin{example}[The identity link]\label{ex:identity}
For an already-finite receiver (an ideal qubit protocol, or any fixture whose
POVM can be consumed directly), the chain has one link with
$\textsf{stage}=\textsf{receiver}$, $\textsf{kind}=\textsf{identity}$,
$\Lambda_{\Lnk}=\idmap$ --- or the sector pinching $\PinS$ of the declared
decomposition, which is the identity on every block-diagonal state and is the
channel of the identity instance, Corollary~\ref{cor:identity-instance} ---,
$M^{\mathrm{out}}=M^{\mathrm{in}}$, the carried tuple
empty or the identity on any declared operators (block-diagonal ones where
the channel is $\PinS$),
$\Srch_{\Lnk}$ the whole declared input family, pinched where the channel is
$\PinS$,
$\CList_{\Lnk}=\varnothing$, no licence, and
$\textsf{status}=\textsf{exact}$. It is a \termref{term-squashing-relation}{squashing relation}: by
reflexivity (Proposition~\ref{prop:preorder}) for $\idmap$, and by
Proposition~\ref{prop:pinching} for $\PinS$, every declared element being
block diagonal for the decomposition. So Theorem~\ref{thm:chain} applies and
returns the physical model unchanged. The link is recorded rather than
omitted (Requirement~\ref{req:identity-link-recorded}). It is the chain of
the identity instance (Corollary~\ref{cor:identity-instance}), whose emitted
family is Corollary~\ref{cor:emitted-identity}. No input the module's
interface declares carries an already-finite receiver, so no declaration
reaches this link; admitting one is a new declared input with a new refusal
rule.
\end{example}
\claimmeta{\metalabel{Uses}\mref{Prop.}{prop:preorder}, \mref{Prop.}{prop:pinching}, \mref{Def.}{def:chain}, \mref{Thm.}{thm:chain}
\metasep newly derived
\metasep accepted}

\section{The source-side reduction and the two-sided theorem}\label{group-reductions/lem:precomposition}
\sectionintro{Admissibility for a list of entries with real constants survives precomposition with any channel, in either source mode (\hyperlink{claim-reductions/lem:precomposition}{Lemma~\ref*{lem:precomposition}}). Replacing the specified signal ensemble of a trusted-source specification by its virtual ensemble is itself an \termref{term-exact-operational-reduction}{exact operational reduction}, and the entanglement-based description a source-replacement certificate presents has the same \termref{term-full-output}{full output} (\hyperlink{claim-reductions/lem:source-transfer}{Lemma~\ref*{lem:source-transfer}}). The main result, \hyperlink{claim-reductions/thm:two-sided}{Theorem~\ref*{thm:two-sided}}, is that the source map and the squashing reduction then compose, in either order, into one exact operational reduction whose image lies in the admissible attack set, so security of the virtual squashed model transfers to the specified physical protocol with the same \ensuremath{\varepsilon}, against attacks that receive every signal in the first round and keep finite-dimensional registers; every other attack is reached by a limit in the consuming engine (\xref{meat-core/cor:composed-security}). In untrusted-source mode there is no source segment to transfer, and the receiver-side reduction stands alone.}
\sectionimports{\extref{../transmitter/contract.pdf}{transmitter/def:source-segment}, \extref{../transmitter/contract.pdf}{transmitter/def:source-spec}, \extref{../transmitter/contract.pdf}{transmitter/def:srcert}, \extref{../transmitter/contract.pdf}{transmitter/thm:certificate-sufficiency}}

% claim reductions/lem:precomposition
\claimsummary{reductions/lem:precomposition}{An attack channel admissible for a list of entries with real constants stays admissible when any channel is composed before it, in either source mode and for party-local entries. A Hermitian right-hand side that is not a real multiple of the identity is not preserved.}
\claimbookmark{reductions/lem:precomposition}{Lemma}{lem:precomposition}{Licensed admissibility survives precomposition with a source channel}
\claimleadin{%
A \termref{term-receiver-stage-typed-link}{receiver-stage link} hands the analysis an
\exttermref{../foundations/contract.pdf}{term-attack-set}{attack set} that is a subset of the channels, in general a proper one
(Theorem~\ref{thm:licence}), while a \exttermref{../transmitter/contract.pdf}{term-source-link}{source-stage link}
re-parameterises the first-round attack as a composition
with a channel; applying the source-stage link after the receiver-stage link
needs an inclusion (Remark~\ref{rem:ordering}). The lemma proves that inclusion for every list
whose right-hand sides are real constants, which is the form of a
\termref{term-licensed}{licensed} list in either source mode (Definition~\ref{def:licensed}) and of every entry
this contract emits (Theorem~\ref{thm:licence}\,(iv)), and exhibits a
right-hand side for which it fails.
}
\begin{lemma}[Licensed admissibility survives precomposition with a source channel]\label{lem:precomposition}
Let $\CList=\{(\hat O_k,\bowtie_k,c_k)\}_{k=1}^{m}$ be a list of Hermitian
$\hat O_k$, relations $\bowtie_k\in\{\preceq,=,\succeq\}$ and \emph{real}
constants $c_k$. Let $\Hs{E_{\mathrm{pre}}}$ and $\Hs{E'}$ be separable, let
$\Qsys$, $E''$, $A$, $R$ and $R'$ be finite-dimensional, and let
$\Gamma\in\mathrm{CPTP}(E_{\mathrm{pre}},E')$ be \emph{any} channel. The sets
below are the ones the displays \eqref{eq:attset}, \eqref{eq:attsetun} and
\eqref{eq:attsetpl} of Definition~\ref{def:attset} cut out, read at these
registers with the adjoints of Lemma~\ref{lem:sep}; where $E_{\mathrm{pre}}$
and $E'$ are finite-dimensional they are the attack sets that definition
names.
\begin{enumerate}[label=(\roman*),leftmargin=3.0em]
\item \emph{(\exttermref{../transmitter/contract.pdf}{term-trusted-source-mode}{trusted-source mode})} if
  $\Aev\in\AttSet_{\CList}(E',\Qsys E'')$ of \eqref{eq:attset}, then
  $\Aev\circ\Gamma\in\AttSet_{\CList}(E_{\mathrm{pre}},\Qsys E'')$;
\item \emph{(\exttermref{../transmitter/contract.pdf}{term-untrusted-source-mode}{untrusted-source mode})} if
  $\Aev\in\AttSet_{\CList}(E',A\Qsys E'')$ of \eqref{eq:attsetun}, then
  $\Aev\circ\Gamma\in\AttSet_{\CList}(E_{\mathrm{pre}},A\Qsys E'')$;
\item \emph{(party-local entries)} let $\Qsys=\bigotimes_f\Qsys_f$ and let the
  entry $(\hat O_k,\bowtie_k,c_k;f_k)$ carry the party-local tag
  (Definition~\ref{def:licensed}), so that its condition on $\Aev$ is
  \eqref{eq:attsetpl}. If $\Aev$ satisfies \eqref{eq:attsetpl} for every
  finite-dimensional $R'$ and every $X\in\Pos{R'R\,\Qsys_{f_k^c}}$, then so does
  $\Aev\circ\Gamma$, with $E_{\mathrm{pre}}$ in place of $E'$. Consequently the
  admissible set of Definition~\ref{def:attset}, with \emph{any} subset of its
  entries tagged and the rest read as scalars, is closed under precomposition
  with a channel into its input register;
\item \emph{(the real constant is what carries it)} replace the right-hand side
  $c_k\,\idon{E'}$ of \eqref{eq:attset} by a Hermitian $\Theta\in\Lin{E'}$ that
  is not a real multiple of $\idon{E'}$. Then the condition
  $\Aev^{\dagger}(\hat O\otimes\idon{E''})\succeq\Theta$ is not preserved by
  precomposition: take $E_{\mathrm{pre}}=E'=\Qsys=\mathbb{C}^{2}$ with declared
  basis $\{\ket{0},\ket{1}\}$, $E''$ one-dimensional, $\Aev=\idmap$ the
  identity channel $E'\to\Qsys$, $\hat O=\Theta=\dyad{0}$, and
  $\Gamma(\rho)\coloneqq\tr[\rho]\,\dyad{1}$. The condition holds for $\Aev$
  with equality, and
  $(\Aev\circ\Gamma)^{\dagger}(\hat O)=0\not\succeq\dyad{0}=\Theta$.
\end{enumerate}
\end{lemma}
\claimmeta{\metalabel{Uses}\mref{Lem.}{lem:sep}, \mref{Def.}{def:attset}, \mref{Def.}{def:licensed}
\metasep \metalabel{Used by}\mref{Rem.}{rem:precomposition-scope}, \mref{Thm.}{thm:two-sided}, \mref{Rem.}{rem:ordering}
\metasep newly derived
\metasep accepted}

\begin{proof}
\emph{Two facts are used throughout, and nothing else.} First, for a channel
$\Phi$ whose output system is finite-dimensional the adjoint $\Phi^{\dagger}$
is the unital positive map fixed by \eqref{eq:pairing}; the composite
$\Aev\circ\Gamma$ has the same output system as $\Aev$, so
$(\Aev\circ\Gamma)^{\dagger}$ is defined on the same operators as
$\Aev^{\dagger}$, and no adjoint of $\Gamma$ is formed in (i) to (iii), which
is what lets $\Hs{E'}$ be separable and infinite-dimensional. Second, a
bounded self-adjoint $H$ on a separable $\Hs{K}$ with $\tr[H\varrho]\ge0$ for
every $\varrho\in\Pos{K}$ of finite trace satisfies $H\succeq0$: take
$\varrho=\dyad{\phi}$ for a unit vector $\phi$, so that
$\bra{\phi}H\ket{\phi}=\tr[H\dyad{\phi}]\ge0$.

(i) Fix $k$, write $Z_k\coloneqq\hat O_k\otimes\idon{E''}$ and suppose
$\bowtie_k$ is $\succeq$, so that $\Aev^{\dagger}(Z_k)\succeq c_k\,\idon{E'}$.
Let $\varrho\in\Pos{E_{\mathrm{pre}}}$ have finite trace. By
\eqref{eq:pairing} applied to $\Aev\circ\Gamma$, then to $\Aev$,
\begin{equation*}
  \tr\!\big[(\Aev\circ\Gamma)^{\dagger}(Z_k)\,\varrho\big]
  =\tr\!\big[Z_k\,\Aev\big(\Gamma(\varrho)\big)\big]
  =\tr\!\big[\Aev^{\dagger}(Z_k)\,\Gamma(\varrho)\big]
  \;\ge\;c_k\,\tr\!\big[\Gamma(\varrho)\big]
  \;=\;c_k\,\tr\varrho ,
\end{equation*}
the inequality because $\Aev^{\dagger}(Z_k)-c_k\idon{E'}$ is a positive bounded
operator and $\Gamma(\varrho)$ a positive trace-class one, so their product has
non-negative trace, and the last step because $\Gamma$ is trace preserving.
Hence $\tr\!\big[\big((\Aev\circ\Gamma)^{\dagger}(Z_k)
-c_k\idon{E_{\mathrm{pre}}}\big)\varrho\big]\ge0$ for every such $\varrho$, and
the second fact gives
$(\Aev\circ\Gamma)^{\dagger}(Z_k)\succeq c_k\,\idon{E_{\mathrm{pre}}}$.
The right-hand side's being a real multiple of the identity is used exactly
once, in the last step of the display: it makes the right-hand side a
multiple of $\tr\varrho$, and therefore the same constant on the two
registers. For $\bowtie_k\,=\,\preceq$ apply the argument to $-Z_k$ and
$-c_k$; for $\bowtie_k\,=\,=$ both directions hold. Since $k$ was arbitrary and
$\Aev\circ\Gamma$ is a channel from $E_{\mathrm{pre}}$ with the output system of
$\Aev$, it lies in $\AttSet_{\CList}(E_{\mathrm{pre}},\Qsys E'')$.

(ii) Identical, with $Z_k\coloneqq\idon{A}\otimes\hat O_k\otimes\idon{E''}$:
the tested operator is the only thing that differs between \eqref{eq:attset}
and \eqref{eq:attsetun}, and $\Gamma$ acts on the input register in both.

(iii) Fix $k$, $R'$ and $X\in\Pos{R'R\,\Qsys_{f_k^c}}$, write
$\hat D_k\coloneqq\hat O_k-c_k\,\idon{\Qsys_{f_k}}$ and
$Z\coloneqq X\otimes\hat D_k\otimes\idon{E''}$, and suppose $\bowtie_k$ is
$\succeq$, so that $(\idmap_{R'}\otimes\Aev)^{\dagger}(Z)\succeq0$. Since
$\idmap_{R'}\otimes(\Aev\circ\Gamma)
=(\idmap_{R'}\otimes\Aev)\circ(\idmap_{R'}\otimes\Gamma)$, the same two
applications of \eqref{eq:pairing} give, for every
$\varrho\in\Pos{R'E_{\mathrm{pre}}}$ of finite trace,
\begin{equation*}
  \tr\!\Big[\big(\idmap_{R'}\otimes(\Aev\circ\Gamma)\big)^{\dagger}(Z)\,
    \varrho\Big]
  =\tr\!\Big[(\idmap_{R'}\otimes\Aev)^{\dagger}(Z)\;
    \big(\idmap_{R'}\otimes\Gamma\big)(\varrho)\Big]\;\ge\;0 ,
\end{equation*}
because $(\idmap_{R'}\otimes\Gamma)(\varrho)$ is positive and trace class. The
second fact gives \eqref{eq:attsetpl} for $\Aev\circ\Gamma$ at that $R'$ and
that $X$, and $R'$, $X$ were arbitrary. For $\preceq$ apply this to $-Z$; for
$=$ both directions hold. No unitality and no constant is used here, the
right-hand side of \eqref{eq:attsetpl} being $0$. The closure statement is (i)
or (ii) for the untagged entries and this for the tagged ones, each entry's
condition being a separate condition on the same channel.

(iv) $\Aev=\idmap$ has $\Aev^{\dagger}=\idmap$, so
$\Aev^{\dagger}(\hat O)=\dyad{0}=\Theta$ and the condition holds with equality,
hence in the direction $\succeq$. The map $\Gamma$ is completely positive and
trace preserving, and $\tr[Y\,\Gamma(\rho)]=\tr[\rho]\,\bra{1}Y\ket{1}
=\tr\big[\big(\bra{1}Y\ket{1}\,\idon{E_{\mathrm{pre}}}\big)\rho\big]$ for every
$Y$ and $\rho$, so $\Gamma^{\dagger}(Y)=\bra{1}Y\ket{1}\,\idon{}$; this is the
one place in the proof where an adjoint of $\Gamma$ is written, and here both
spaces are two-dimensional. Hence
$(\Aev\circ\Gamma)^{\dagger}(\hat O)=\Gamma^{\dagger}(\dyad{0})
=\braket{1}{0}\braket{0}{1}\,\idon{}=0$, while $\Theta=\dyad{0}$, and
$0-\dyad{0}$ is not positive semidefinite. The scalar entry built from the
same operator is not broken by the same $\Gamma$:
$\Aev^{\dagger}(\dyad{0})=\dyad{0}\succeq0\cdot\idon{}$ is the licensed entry
at $c=0$, and $\Gamma^{\dagger}(\dyad{0})=0\succeq0\cdot\idon{}$ holds.
\end{proof}

% claim reductions/rem:precomposition-scope
\claimsummary{reductions/rem:precomposition-scope}{The precomposition lemma needs only membership in the set the attack-set displays cut out, at a separable input, and gives the inclusion of the specified-source attack set, composed with any channel, in the virtual-source one. A general entry has no reading beside a non-empty source segment.}
\claimbookmark{reductions/rem:precomposition-scope}{Remark}{rem:precomposition-scope}{Scope of the precomposition lemma}
\begin{remark}[Scope of the precomposition lemma]\label{rem:precomposition-scope}
\emph{The general entry.} The right-hand side $\Theta$ of
Lemma~\ref{lem:precomposition}\,(iv), Hermitian and not a real multiple of the
identity, is the channel-level shape that the Hermitian right-hand side of a
\emph{general} entry asks for, an entry Definition~\ref{def:attset} records as
admissible input to the consuming contract and never licensed there. At the state level
such an entry beside a \exttermref{../transmitter/contract.pdf}{term-source-link}{source link} fails earlier still, and this is
separate from the counterexample. There its right-hand side is a Hermitian
operator on Alice's retained system (Definition~\ref{def:attset};
\xref{meat-core/def:admissible-declaration}\,(vi)), and the retained system a
\exttermref{../transmitter/contract.pdf}{term-source-replacement-certificate}{source-replacement certificate} is written against is the one of the
\emph{virtual} description, at every length of the \exttermref{../transmitter/contract.pdf}{term-source-segment}{source segment}
(\xref{transmitter/def:srcert}, ``One reading at every $L_{\mathrm s}$''); no
contract relates that system or its marginal to the specified description's.
So such an entry does not merely fail to be preserved: beside a non-empty
source segment it has no reading to preserve. That, and not the
counterexample alone, is why the admission rule of
Definition~\ref{def:chain} is a refusal.

\emph{Membership, and where it comes from.} The hypothesis of
Lemma~\ref{lem:precomposition}\,(i) to (iii) is \emph{membership} in the set
the displays of Definition~\ref{def:attset} cut out. The attack sets that
definition names are those of the \exttermref{../foundations/contract.pdf}{term-round-model}{round model}, whose input registers are
finite-dimensional; the displays themselves carry no condition on the
dimension of $E'$, and the lemma reads them at a separable input register,
so (i) to (iii) hold there.
Obtaining that membership from Theorem~\ref{thm:licence}\,(i), (ii) or (v) is
a different matter, because that theorem types its attack on a
finite-dimensional input register. Theorem~\ref{thm:two-sided}\,(ii)
therefore does not route the first round through the lemma: it applies
Theorem~\ref{thm:licence} to the pre-composed attack, whose input register is
the \exttermref{../transmitter/contract.pdf}{term-virtual-transmitted-system}{virtual transmitted system}. The lemma is the inclusion
Remark~\ref{rem:ordering} needs, and the second of the two routes in the
proof of Theorem~\ref{thm:two-sided}\,(ii).

\emph{The inclusion.} Write
$\AttSet^{\mathrm{real}}_1\coloneqq\AttSet_{\CList}(E',\cdot)$ for the set the
first-round attack of the specified-source protocol is admissible in and
$\AttSet^{\mathrm{virtual}}_1\coloneqq
\AttSet_{\CList}(E_{\mathrm{pre}},\cdot)$ for the set its precomposition with
$\Gamma$ must land in. Then Lemma~\ref{lem:precomposition}\,(i) to (iii) say
exactly
$\AttSet^{\mathrm{real}}_1\circ\Gamma\subseteq\AttSet^{\mathrm{virtual}}_1$,
\emph{for every channel $\Gamma$}. That is the inclusion
Remark~\ref{rem:ordering} names, and Theorem~\ref{thm:two-sided} is where it
is spent.
\end{remark}
\claimmeta{\metalabel{Uses}\mref{Lem.}{lem:precomposition}
\metasep newly derived
\metasep accepted}

% claim reductions/lem:source-transfer
\claimsummary{reductions/lem:source-transfer}{In trusted-source mode, replacing the specified signal ensemble by the virtual one is an exact operational reduction, and the entanglement-based description a certificate presents has the same full output. One source channel common to every letter is essential.}
\claimbookmark{reductions/lem:source-transfer}{Lemma}{lem:source-transfer}{Exactness of the source map in trusted-source mode}
\claimleadin{%
A source map replaces \emph{Alice's} preparation, and only a
\exttermref{../transmitter/contract.pdf}{term-source-specification}{source specification} in trusted-source mode has one. In
untrusted-source mode the specification carries ``no certificate and no
marginal'' (\xref{transmitter/def:source-spec}), the retained system $A_j$ is
an \emph{output} of the attack (Definition~\ref{def:roundmodel}, untrusted
branch), and both (SR3), that no operation of the adversary acts on $A$, and
the register discipline that \xref{transmitter/thm:certificate-sufficiency}\,(iii)
requires of its channel fail there; correspondingly A-SRC
(\xref{transmitter/ass:src}) is stated of a specification in trusted-source
mode. So
the lemma is stated in trusted-source mode alone; in untrusted-source mode
there is no source segment to transfer, and what remains is the receiver-side
instance, Corollary~\ref{cor:sequential}. The specified-source input
$\rho^{\mathrm{sp}}$ of the lemma is the state in which the physical source prepares
the letters and signals of the $n$ rounds, under the transmitter's
physical-source premise that makes the rounds independent: A-ENS
(\xref{transmitter/ass:ens}) or, for a specification declared by a
\exttermref{../transmitter/contract.pdf}{term-source-model}{source model}, the premise that model declares in its place, A-WCP
(\xref{transmitter/decoy/ass:wcp}) for the decoy model. The lemma is a
statement about $\rho^{\mathrm{sp}}$ and does not assume that premise.
}
\begin{lemma}[Exactness of the source map in trusted-source mode]\label{lem:source-transfer}
Let a source specification in \emph{trusted-source mode}
(\xref{transmitter/def:source-spec}) declare an alphabet $\Xset$ with weights
$(p_x)_{x\in\Xset}$, the \emph{specified} ensemble $\{p_x,\rho_x\}$ on the
\exttermref{../transmitter/contract.pdf}{term-transmitted-system}{transmitted system} $\Asig$ and the \exttermref{../transmitter/contract.pdf}{term-virtual-ensemble}{virtual ensemble}
$\{p_x,\rvirt_x\}$ on the virtual transmitted system $\Avirt$, with
$\Hs{\Asig}$ separable and $\Hs{\Avirt}$ finite-dimensional, as
\xref{transmitter/def:source-spec} declares the virtual transmitted system,
and let its source segment satisfy
the \exttermref{../transmitter/contract.pdf}{term-interface-conditions}{interface conditions} (S1)--(S4) of
\xref{transmitter/def:source-segment}, so that its composite
$\Xi\colon\Tcl{\Avirt}\to\Tcl{\Asig}$ satisfies the \emph{source identity}
\begin{equation}\label{eq:sourceidentity}
  \Xi\big(\rvirt_x\big)=\rho_x\qquad\text{for every }x\in\Xset .
\end{equation}
Fix $n$, a finite-dimensional adversary register $E_0$ and a protocol in the
sense of Definition~\ref{def:protocol}, and put
\begin{equation}\label{eq:prepstates}
  \rho^{\mathrm{sp}}\coloneqq\bigotimes_{j=1}^{n}
    \Big(\sum_{x\in\Xset}p_x\,\dyad{x}_{X_j}\otimes\rho_x\Big),
  \qquad
  \rho^{\mathrm{v}}\coloneqq\bigotimes_{j=1}^{n}
    \Big(\sum_{x\in\Xset}p_x\,\dyad{x}_{X_j}\otimes\rvirt_x\Big),
\end{equation}
the $j$-th factor living on $X_j\Asig_j$ and on $X_j\Avirt_j$ respectively.

\emph{The two \termref{term-protocol-presentation}{presentations}.} A \emph{strategy}
is a pair
$a=(\zeta,(\Aev_j)_{j=1}^{n})$ of an initial adversary state
$\zeta\in\Seq{E_0}$ and channels
\begin{equation}\label{eq:srcstrategy}
  \Aev_1\in\mathrm{CPTP}\big(\Asig^{n}E_0,\ \Bsys_1E'_1\big),
  \qquad
  \Aev_j\in\mathrm{CPTP}\big(E_{j-1},\ \Bsys_jE'_j\big)\ \ (2\le j\le n),
\end{equation}
with $E'_j$, $E_j$ as in Definition~\ref{def:roundmodel}: the \emph{trusted}
branch of \eqref{eq:attackchannel}, the only branch this lemma is about, in
which $A_j$ is Alice's and is no output of $\Aev_j$. A strategy of the
\emph{virtual} type is the same with $\Avirt^{n}$ in place of $\Asig^{n}$ in
\eqref{eq:srcstrategy}. The \emph{specified-source} presentation
$\mathsf{P}^{\mathrm{sp}}$ has these strategies and the output map
$\mathsf{out}^{\mathrm{sp}}(a)\coloneqq$ the state \eqref{eq:protocol} produces
on the \termref{term-output-registers}{output registers} of $\Kfin$ together with $E_n$, from the input
$\rho^{\mathrm{sp}}\otimes\zeta$, with $\Bsys_j$ supplied by $\Aev_j$ and the
classical $X_j$ carried into $\bar{\Grnd}_j$; the \emph{virtual-source}
presentation $\mathsf{P}^{\mathrm{v}}$ has the strategies of the virtual type
and the output map built the same way from $\rho^{\mathrm{v}}\otimes\zeta$, the
two protocols agreeing in \emph{every} other component: the
\termref{term-measured-system}{measured system} and its POVM, $\bar{\Grnd}_j$ and $\Kfin$. The output registers are literally
the same, which is condition (R1) of Definition~\ref{def:opreduction}.
\begin{enumerate}[label=(\roman*),leftmargin=3.0em]
\item \emph{(the source map)} define
  $r_{\mathrm{src}}(\zeta,(\Aev_j)_j)\coloneqq(\zeta,(\Aev^{\mathrm{v}}_j)_j)$
  by
  \begin{equation}\label{eq:srcmap-strategy}
    \Aev^{\mathrm{v}}_1\coloneqq
      \Aev_1\circ\big(\Xi^{\otimes n}\otimes\idmap_{E_0}\big),
    \qquad
    \Aev^{\mathrm{v}}_j\coloneqq\Aev_j\ \ (2\le j\le n).
  \end{equation}
  Then $r_{\mathrm{src}}$ is an \termref{term-exact-operational-reduction}{exact operational reduction} from
  $\mathsf{P}^{\mathrm{sp}}$ to $\mathsf{P}^{\mathrm{v}}$ in the sense of
  Definition~\ref{def:opreduction}, with the whole of the virtual type as the
  \termref{term-analysed-strategy-set}{admissible strategy set};
\item \emph{(the two descriptions of the virtual protocol)} let the virtual
  protocol be described \emph{entanglement-based} instead: Alice's retained
  systems $A_1,\dots,A_n$ carry, jointly with the virtual transmitted systems,
  the preparation $\Psi^{\otimes n}$, where $\Psi$ is a source-replacement
  certificate for the \exttermref{../foundations/contract.pdf}{term-declared-data}{declared data} $(A,\mathbb{T},\tau_A,\{M_x\})$
  against $\{p_x,\rvirt_x\}$ in any of the declared forms of
  \xref{transmitter/def:srcert}, satisfying (SR1)--(SR3); and $x_j$ is
  produced by her declared \termref{term-measurement-channel}{measurement channel} $\pcNmsm$, the
  measurement map of part (iii) of \xref{transmitter/thm:certificate-sufficiency}, applied to
  $A_j$ at its declared time in round $j$. Write $\mathsf{out}^{\mathrm{eb}}$
  for the resulting output map on the virtual type. Then, under A-SEQ
  (\xref{transmitter/ass:seq}),
  $\mathsf{out}^{\mathrm{eb}}(a)=\mathsf{out}^{\mathrm{v}}(a)$ for every
  strategy $a$ of that type, so the identity map on strategies is an exact
  operational reduction between the two presentations in either direction;
\item \emph{(the composite)} consequently $r_{\mathrm{src}}$ is an exact
  operational reduction from $\mathsf{P}^{\mathrm{sp}}$ to the virtual protocol
  in its entanglement-based description, the object the rest of this
  contract and the consuming engine analyse;
\item \emph{(one channel for every label is essential)} \eqref{eq:sourceidentity}
  with \emph{one} $\Xi$ cannot be weakened to a family $(\Xi_x)_x$ with
  $\Xi_x(\rvirt_x)=\rho_x$. Take $n=1$, $\Xset=\{0,1\}$,
  $p_0=p_1=\tfrac12$, $\Asig=\Avirt=\mathbb{C}^{2}$,
  $\rvirt_0=\rvirt_1=\dyad{0}$, $\rho_0=\dyad{0}$, $\rho_1=\dyad{1}$, and
  $\Xi_x(\cdot)\coloneqq\tr[\cdot]\,\rho_x$. Every identity
  $\Xi_x(\rvirt_x)=\rho_x$ holds, yet for the attack that measures the received
  signal in the basis $\{\ket{0},\ket{1}\}$ and keeps the outcome in $E'_1$,
  the marginal on $X_1E'_1$ is perfectly correlated in the specified-source
  protocol and, for \emph{every} attack of the virtual type, a product in the
  virtual one, so no map of strategies can make the two
  \termref{term-full-output}{full outputs} equal; the protocol exhibited in the proof announces nothing in round one, so that
  $E_1$ is $E'_1$ there.
\end{enumerate}
\end{lemma}
\claimmeta{\metalabel{Uses}\mref{Lem.}{lem:sep}, \mref{Def.}{def:protocol}, \mref{Def.}{def:opreduction}, \mref{Lem.}{lem:opreduction-transfer}, \mref{Def.}{def:roundmodel}, \extref{../transmitter/contract.pdf}{transmitter/def:source-segment}, \extref{../transmitter/contract.pdf}{transmitter/def:source-spec}, \extref{../transmitter/contract.pdf}{transmitter/def:srcert}, \extref{../transmitter/contract.pdf}{transmitter/thm:certificate-sufficiency}
\metasep \metalabel{Assumes}\extaref{../transmitter/contract.pdf}{transmitter/ass:seq}{A-SEQ}
\metasep \metalabel{Used by}\mref{Rem.}{rem:source-transfer-registers}, \mref{Rem.}{rem:source-transfer-scope}, \mref{Thm.}{thm:two-sided}
\metasep newly derived
\metasep accepted}

\begin{proof}
(i) The map $\Gamma\coloneqq\Xi^{\otimes n}\otimes\idmap_{E_0}$, and with it
$\idmap_{X^{n}}\otimes\Gamma$, is a channel that acts factor by factor on
products: its inputs $\Avirt$, $E_0$ and $X_j$ are finite-dimensional, so
Lemma~\ref{lem:sep}\,(iii) gives it, with $\Phi_i=\Xi$ on each $\Avirt_j$
and the identity maps on $E_0$ and on the classical registers. First,
$\big(\idmap_{X^{n}}\otimes\Xi^{\otimes n}\big)
(\rho^{\mathrm{v}})=\rho^{\mathrm{sp}}$. The map acts factor by factor on
\eqref{eq:prepstates}, and on the $j$-th factor
$\sum_xp_x\dyad{x}\otimes\Xi(\rvirt_x)=\sum_xp_x\dyad{x}\otimes\rho_x$ by
\eqref{eq:sourceidentity}, linearity and continuity of $\Xi$ on the trace-class
operators; a term with $p_x=0$ is the zero operator on both sides whatever
$\rvirt_x$ and $\rho_x$ are, which is why \eqref{eq:sourceidentity} is needed
only where the weight is positive. Tensoring with $\idmap_{E_0}$ gives
\begin{equation}\label{eq:srcinput}
  \big(\idmap_{X^{n}}\otimes\Xi^{\otimes n}\otimes\idmap_{E_0}\big)
  \big(\rho^{\mathrm{v}}\otimes\zeta\big)
  =\rho^{\mathrm{sp}}\otimes\zeta .
\end{equation}
Second, write $\Grnd^{\Aev_1}_1$ for the round-one map of \eqref{eq:protocol}
with $\Bsys_1$ supplied by $\Aev_1$, and
$\mathsf{C}\coloneqq\Kfin\circ\Grnd^{\Aev_n}_n\circ\dots\circ\Grnd^{\Aev_2}_2$
for everything that follows it. Since $\Xi^{\otimes n}$ acts on $\Avirt^{n}$
alone and $\Aev^{\mathrm{v}}_1=\Aev_1\circ(\Xi^{\otimes n}\otimes\idmap_{E_0})$,
\begin{equation*}
  \Grnd^{\Aev^{\mathrm{v}}_1}_1
  =\Grnd^{\Aev_1}_1\circ
   \big(\idmap_{X^{n}}\otimes\Xi^{\otimes n}\otimes\idmap_{E_0}\big),
\end{equation*}
the inserted map being the identity on every system $\Grnd^{\Aev_1}_1$ reads
apart from the transmitted ones. And $\mathsf{C}$ is \emph{literally the same}
composite in the two executions: $\Aev^{\mathrm{v}}_j=\Aev_j$ for $j\ge2$ by
\eqref{eq:srcmap-strategy}, the two protocols agree in every component other
than the released signal system, and no map of $\mathsf{C}$ touches that
system, the transmitted systems having been consumed by $\Aev_1$. Therefore
\begin{equation*}
  \mathsf{out}^{\mathrm{v}}\big(r_{\mathrm{src}}(a)\big)
  =\mathsf{C}\circ\Grnd^{\Aev^{\mathrm{v}}_1}_1
     \big(\rho^{\mathrm{v}}\otimes\zeta\big)
  =\mathsf{C}\circ\Grnd^{\Aev_1}_1\big(\rho^{\mathrm{sp}}\otimes\zeta\big)
  =\mathsf{out}^{\mathrm{sp}}(a),
\end{equation*}
using \eqref{eq:srcinput} in the middle step. This is
\eqref{eq:opreduction}. Condition (R3) holds because
$\Aev^{\mathrm{v}}_1$ is a composition of channels and has the input
$\Avirt^{n}E_0$ and the output $\Bsys_1E'_1$ required by
\eqref{eq:srcstrategy} of the virtual type, and (R2) is the display.

(ii) By (SR3) no operation of the protocol and no operation of the adversary
acts on $A_j$ except Alice's declared measurement, and $\pcNmsm$ acts on $A_j$
alone and writes into the fresh classical register $X_j$. Every map of the
execution that \emph{reads} $X_j$, namely the announcement of round $j$ inside
$\bar{\Grnd}_j$ and, through the public record, every later map, comes after
Alice's round-$j$ measurement, by A-SEQ: every public message of round $j$
follows every measurement of round $j$, and no operation on a quantum system of
round $j$ depends on a message of round $j$. So each map of the execution that
precedes the first map reading $X_j$ acts on systems complementary to $A_jX_j$,
and part (iii)(a) of \xref{transmitter/thm:certificate-sufficiency}, the
interchange law, which reads no property of the channel it is stated for
beyond the register discipline, lets Alice's round-$j$ measurement be
performed before it. Applying this for $j=1,\dots,n$ moves all $n$
measurements to the start of the execution, where the state is
$\Psi^{\otimes n}\otimes\zeta$. There
\begin{equation*}
  \big(\pcNmsm^{\otimes n}\otimes\idmap_{E_0}\big)
  \big(\Psi^{\otimes n}\otimes\zeta\big)
  =\bigotimes_{j=1}^{n}\Big(\sum_{x\in\Xset}\dyad{x}_{X_j}\otimes
    \tr_{A}\big[(M_x\otimes\idon{\Avirt})\,\Psi\big]\Big)\otimes\zeta
  =\rho^{\mathrm{v}}\otimes\zeta ,
\end{equation*}
by the definition of $\pcNmsm$ and by (SR2), which is the source
identification of that theorem's part (iii)(b) at the trivial attack. From that
point the two executions are the same composite of the same maps applied to the
same state, so the full outputs agree; the output registers agree because
$\pcNmsm$ consumes $A_j$ and produces $X_j$ in both descriptions. Neither
direction of the equality reads any property of $\Psi$ beyond (SR2) and (SR3),
so the conclusion holds in whichever declared form the certificate was given;
(SR1) enters as a hypothesis of the theorem cited and as what the consuming
engine reads as $\tau_A$, and no step of this proof uses it.

(iii) By (ii) the identity map is an exact operational reduction from
$\mathsf{P}^{\mathrm{v}}$ to the entanglement-based presentation, whose strategy
set is the same; composing it with $r_{\mathrm{src}}$ is
Lemma~\ref{lem:opreduction-transfer}\,(v).

(iv) Each $\Xi_x$ is a channel and $\Xi_x(\rvirt_x)=\tr[\rvirt_x]\rho_x=\rho_x$.
In the specified-source protocol the state on $X_1\Asig_1$ is
$\tfrac12\dyad{0}_{X_1}\otimes\dyad{0}+\tfrac12\dyad{1}_{X_1}\otimes\dyad{1}$,
so the stated attack leaves $X_1$ and $E'_1$ perfectly correlated; in the
virtual protocol the state is
$\big(\tfrac12\dyad{0}_{X_1}+\tfrac12\dyad{1}_{X_1}\big)\otimes\dyad{0}$, a
product, so the same attack and, the signal carrying no dependence on $x$,
\emph{every} attack of the virtual type leaves $E'_1$ independent of $X_1$.
Take the protocol whose $\bar{\Grnd}_1$ retains $x_1$ as its per-round datum and
\emph{appends nothing to the public record}, and whose $\Kfin$ is the identity
on the retained data and on $E_1$; it is a protocol of the form
\eqref{eq:protocol}, and with no round-one announcement $E_1$ is $E'_1$, the
adversary retaining no public datum of that round. The two outputs are then
compared on the marginal on $X_1E'_1$: for the stated attack it is
perfectly correlated in the specified-source protocol and, for \emph{every}
attack of the virtual type, a product in the virtual one, so the two full
outputs differ and no map of strategies can equate them. The argument reads
the marginal on $X_1E'_1$, and that is why the exhibited protocol announces
nothing in round one: at a protocol whose round-one announcement is a
function of $x_1$, $E_1$ would collect that datum and would not be
independent of $X_1$ in either protocol.
\end{proof}

% claim reductions/rem:source-transfer-registers
\claimsummary{reductions/rem:source-transfer-registers}{The transmitted systems enter the first-round attack as systems in transit, not as registers of the adversary, and the virtual one is finite-dimensional. A round-by-round attack with a memory of any dimension is reached from this type by a limit (\xref{meat-core/lem:memory-lift}).}
\claimbookmark{reductions/rem:source-transfer-registers}{Remark}{rem:source-transfer-registers}{The released signals as systems in transit}
\begin{remark}[The released signals as systems in transit]\label{rem:source-transfer-registers}
\emph{What the released signals are.} The first round's input register in
\eqref{eq:srcstrategy} is $\Asig^{n}\otimes E_0$, of which \emph{only $E_0$}
is one of the adversary's registers, finite-dimensional as every register of
the round model of Definition~\ref{def:roundmodel} is. The
transmitted systems are named beside it because they are systems \emph{in
transit}: each is an intermediate system of one channel, produced by the
trusted preparation, consumed by $\Aev_1$, and held by no party at the end of
the round. That is the standing the measured system $\Bsys$ has on the
receiver's side, which is why Definition~\ref{def:system-setting} names
$\Bsys$ separately as separable, and neither it nor the round model names the
transmitted system at all. Neither is read more widely than it is written:
every finiteness that
Lemma~\ref{lem:source-transfer} and Theorem~\ref{thm:two-sided} use, of
$E_0$, of each $E'_j$, of $\Qsys$ and of $\Avirt$, is hypothesised where it is
used, and every claim they apply at a register is applied only where that
claim's own typing holds.

\emph{Which attacks the type covers.} Where $\Hs{\Asig}$ is
finite-dimensional, a strategy that receives all $n$ signals before the first
round subsumes one that receives them round by round: the adversary may hold
the later ones untouched in a finite-dimensional $E'_1$ large enough to hold
them and pass them on, through $E_1,\dots,E_{j-1}$, to the round $j$ that
releases them. A-SEQ
(\xref{transmitter/ass:seq}) orders the round's \emph{messages} after the
round's \emph{measurements} and says nothing about when a signal is released.
Where $\Hs{\Asig}$ is infinite-dimensional, as the transmitted system of a
\exttermref{../transmitter/contract.pdf}{term-source-model}{source model} is, a finite-dimensional $E'_1$ cannot hold the later
signals untouched, and a round-by-round attack in which $\Asig_j$ arrives in
transit at round $j$ is not a strategy of the type \eqref{eq:srcstrategy}
with finite registers. The type is the one the published source map makes:
its first attack channel receives all $n$ signals
\cite[Lemma~9.1, Eq.~(140)]{tupkary2026decoy}, and each $E'_j$ has a fixed
finite dimension \cite[Rem.~8.4]{tupkary2026decoy}. Every other attack is
reached by a limit over the dimensions of the adversary's registers, taken
in the consuming engine (\xref{meat-core/lem:memory-lift}), as the published
analysis takes it \cite[Thm.~A.1, Rem.~A.2]{tupkary2026decoy}: compressing
the adversary's memory, and each signal released after the first round, onto
finite-dimensional subspaces turns a round-by-round attack with a memory of
any dimension into a strategy of this type with finite registers whose output
state is as close as wanted, and the trace distances of the security
statement pass to the limit. Lemma~\ref{lem:source-transfer} and
Theorem~\ref{thm:two-sided} are exact about the type they name and transfer
every criterion; the limit is taken for the trace distances alone.

\emph{The virtual transmitted system is finite-dimensional.} It is so in both
cases of \xref{transmitter/def:source-spec}: trivially where the ensemble is
declared as matrices, and by its clause (b) where the specified system is
separable, ``what a consumer receives is never the model but the virtual
ensemble, which is finite dimensional''. So the virtual presentation's
first-round input register is finite-dimensional even where the specified one
is not, which is the hypothesis Theorem~\ref{thm:two-sided}\,(ii) spends.
\end{remark}
\claimmeta{\metalabel{Uses}\mref{Lem.}{lem:source-transfer}
\metasep newly derived
\metasep accepted}

% claim reductions/rem:source-transfer-scope
\claimsummary{reductions/rem:source-transfer-scope}{The source-map lemma gives the specified-source protocol no entanglement-based description, uses no purification and asserts no equality of marginals. With an empty source segment the two descriptions coincide, and parts (i) and (iii) of the lemma add nothing to the definition of one protocol.}
\claimbookmark{reductions/rem:source-transfer-scope}{Remark}{rem:source-transfer-scope}{Scope of the source-map lemma}
\begin{remark}[Scope of the source-map lemma]\label{rem:source-transfer-scope}
Lemma~\ref{lem:source-transfer} never gives the specified-source protocol an
entanglement-based description. So no purification of a mixed $\rho_x$ is a
system of either presentation, nothing is asserted about its dimension, its
marginal or its correlations, and no statement relating the purifications of
the two preparations is used: where $\rho_x$ is mixed such a purification
exists, and it is a system of neither protocol. Nor is any equality of
marginals asserted: $\tau_A$ is the marginal of the \emph{virtual}
description, which is the pair a certificate is written against at every
length of the source segment (\xref{transmitter/def:srcert}, ``One reading at
every $L_{\mathrm s}$''), and nothing in the lemma equates it with a marginal
of the specified description. At $L_{\mathrm s}=0$ the two descriptions
coincide, $\Xi$ is the identity, $\Avirt=\Asig$ and $\rvirt_x=\rho_x$, and
Lemma~\ref{lem:source-transfer}\,(i) and (iii) say nothing beyond what
Definition~\ref{def:protocol} already says of one protocol.
\end{remark}
\claimmeta{\metalabel{Uses}\mref{Lem.}{lem:source-transfer}
\metasep newly derived
\metasep accepted}

% claim reductions/thm:two-sided
\claimsummary{reductions/thm:two-sided}{In trusted-source mode the source map and the squashing reduction compose, in either order, into one exact operational reduction into the admissible attack set. Security of the virtual squashed model thus transfers with the same \ensuremath{\varepsilon} to the specified physical protocol, against attacks of its type.}
\claimbookmark{reductions/thm:two-sided}{Theorem}{thm:two-sided}{The two-sided reduction in trusted-source mode}
\claimleadin{%
The theorem is about trusted-source mode, as Lemma~\ref{lem:source-transfer}
is and for its reasons: in untrusted-source mode there is no source segment to
transfer, the specification carrying no certificate and no marginal and the
retained system being an output of the attack
(Definition~\ref{def:roundmodel}, untrusted branch), and what remains there is
the receiver-side instance alone, Corollary~\ref{cor:sequential} with
Theorem~\ref{thm:licence}\,(ii), (iii); nothing in the theorem is asserted of
that mode. In part (iii) the specified-source physical protocol pairs the
specified source with the physical receiver. That the physical source is the
specified one is the transmitter's physical-source premise, A-ENS (\xref{transmitter/ass:ens}), or A-WCP
(\xref{transmitter/decoy/ass:wcp}) for the decoy model, and the theorem does
not assume it.
}
\begin{mainclaim}
\begin{theorem}[The two-sided reduction in trusted-source mode]\label{thm:two-sided}
Assume the hypotheses of Lemma~\ref{lem:source-transfer}: a trusted-source
specification with virtual ensemble $\{p_x,\rvirt_x\}$, a source segment
whose composite $\Xi$ satisfies the source identity
\eqref{eq:sourceidentity}, a source-replacement certificate satisfying
(SR1)--(SR3), and A-SEQ. Assume in addition $M\dsq_{\Lambda}F$ together with
a spectral pull-back certificate \eqref{eq:certificate} for the list
$\CList=\{(\hat O_k,\bowtie_k,c_k)\}$, or the factorwise hypotheses of
Theorem~\ref{thm:licence}\,(v) for a list with entries carrying the
party-local tag. Let $\mathsf{P}^{\mathrm{sp}}$ be the specified-source
presentation of Lemma~\ref{lem:source-transfer}, its receiver measuring $M$
on $\Bsys_j$, and let $\mathsf{P}^{\mathrm{v},\sq}_{\CList}$ be the
presentation whose protocol has the virtual source and the target $F$ on
$\Qsys_j$ and whose strategies are the $\CList$-admissible sequences of the
virtual type (Definition~\ref{def:attset}). Define
\begin{equation}\label{eq:twosidedmap}
  r_{\mathrm{two}}\big(\zeta,(\Aev_j)_j\big)
  \coloneqq\big(\zeta,(\Aev^{\mathrm{v},\sq}_j)_j\big),
  \qquad
  \Aev^{\mathrm{v},\sq}_1
    \coloneqq\lambda_1\circ\Aev_1\circ
      \big(\Xi^{\otimes n}\otimes\idmap_{E_0}\big),
  \quad
  \Aev^{\mathrm{v},\sq}_j\coloneqq\lambda_j\circ\Aev_j\ \ (j\ge2),
\end{equation}
where $\lambda_j\coloneqq\Lambda\otimes\idmap_{E'_j}$ is the trusted-source
branch of \eqref{eq:composite}, the only branch the strategies
\eqref{eq:srcstrategy} of $\mathsf{P}^{\mathrm{sp}}$ are typed in.
\begin{enumerate}[label=(\roman*),leftmargin=3.0em]
\item \emph{(one map, and the order of the two replacements is immaterial)}
  $r_{\mathrm{two}}=r_{\sq}\circ r_{\mathrm{src}}
  =r_{\mathrm{src}}\circ r_{\sq}$, where $r_{\mathrm{src}}$ is the source map
  \eqref{eq:srcmap-strategy} and $r_{\sq}$ the squashing map
  $(\Aev_j)_j\mapsto(\Aev^{\sq}_j)_j$ of \eqref{eq:composite}; both composites
  are defined on the same strategies and agree there. And $r_{\mathrm{two}}$ is
  an exact operational reduction from $\mathsf{P}^{\mathrm{sp}}$ to the
  virtual, \termref{term-squashed-protocol}{squashed protocol} taken with \emph{all} strategies of its type;
\item \emph{(the image is admissible)} let $E_0$, every $E'_j$ and $\Qsys$ be
  finite-dimensional, as the round model of Definition~\ref{def:roundmodel}
  makes the adversary's registers, every register of it being
  finite-dimensional, and Definition~\ref{def:system-setting} makes the
  target space, and let
  $\Hs{\Avirt}$ be finite-dimensional, as it is in both
  cases of \xref{transmitter/def:source-spec}: trivially for an ensemble
  declared as matrices, and by its clause (b) for a separable specified
  system. Then every $\Aev^{\mathrm{v},\sq}_j$ of \eqref{eq:twosidedmap}
  lies in the set Definition~\ref{def:attset} assigns to its round, so
  $r_{\mathrm{two}}$ has its image in the strategy set of
  $\mathsf{P}^{\mathrm{v},\sq}_{\CList}$ and is an exact operational reduction
  into it. \emph{No hypothesis on the specified transmitted system is needed}:
  $\Hs{\Asig}$ may be separable and infinite-dimensional, as the transmitted
  system of a source model is, and it is not a register
  Definition~\ref{def:system-setting} or the round model ranges over
  but a system in transit;
\item \emph{(transfer, with the same $\varepsilon$)} consequently, if the
  virtual, squashed protocol is $\varepsilon$-secure against every
  $\CList$-admissible sequence of the virtual type, a \termref{term-security-criterion}{security criterion} in
  the sense of Definition~\ref{def:criterion} taking a value at most
  $\varepsilon$ at each of them, then the specified-source physical protocol
  is $\varepsilon$-secure against \emph{every} \exttermref{../foundations/contract.pdf}{term-attack-sequence}{attack sequence} of its
  type, with the \emph{same} $\varepsilon$ and simultaneously for every such
  criterion;
\item \emph{(in the description the certificate presents)} (i) to (iii) hold
  with the virtual protocol read entanglement-based, Alice retaining $A_j$,
  which trusted-source mode is the mode that gives her, and measuring
  $\{M_x\}$, this being the description the consuming engine analyses;
\item \emph{(the empty source segment)} where the source segment is empty,
  $\Xi$ is the identity, $\Avirt=\Asig$, $\rvirt_x=\rho_x$,
  $r_{\mathrm{src}}$ is the identity map on strategies and $r_{\mathrm{two}}$ is
  $r_{\sq}$; (i) to (iii) are then Corollary~\ref{cor:sequential} and
  Theorem~\ref{thm:licence}\,(iii) word for word, and the theorem states
  nothing new about the case this module emits.
\end{enumerate}
\end{theorem}
\end{mainclaim}
\claimmeta{\metalabel{Uses}\mref{Def.}{def:system-setting}, \mref{Def.}{def:criterion}, \mref{Def.}{def:opreduction}, \mref{Lem.}{lem:opreduction-transfer}, \mref{Def.}{def:roundmodel}, \mref{Cor.}{cor:sequential}, \mref{Def.}{def:attset}, \mref{Def.}{def:licensed}, \mref{Thm.}{thm:licence}, \mref{Lem.}{lem:precomposition}, \mref{Lem.}{lem:source-transfer}, \extref{../transmitter/contract.pdf}{transmitter/def:source-segment}
\metasep \metalabel{Used by}\mref{Rem.}{rem:two-sided-scope}, \mref{Rem.}{rem:ordering}
\metasep newly derived
\metasep accepted}

\begin{proof}
(i) Both composites are defined: $r_{\sq}$ is defined at an attack sequence of
either the specified or the virtual type, being the insertion of $\lambda_j$
after $\Aev_j$ and reading nothing about the attack's input register; and
$r_{\mathrm{src}}$ is defined at a sequence with either receiver, being the
insertion of $\Xi^{\otimes n}\otimes\idmap_{E_0}$ before $\Aev_1$ and reading
nothing about the attack's output systems. Applying them in either order gives,
in round one,
$\lambda_1\circ\big(\Aev_1\circ(\Xi^{\otimes n}\otimes\idmap_{E_0})\big)
=\big(\lambda_1\circ\Aev_1\big)\circ(\Xi^{\otimes n}\otimes\idmap_{E_0})$ by
associativity of composition, and in every later round $\lambda_j\circ\Aev_j$;
so both equal $r_{\mathrm{two}}$. The two insertions are at disjoint places
of one composite: $\Xi^{\otimes n}$ acts on the transmitted systems, which are
inputs of $\Aev_1$, and $\lambda_j$ on the measured system, which is an output
of $\Aev_j$, and neither acts on a system the other acts on.

For exactness, apply Lemma~\ref{lem:opreduction-transfer}\,(v) to
$r_{\mathrm{src}}$, an exact operational reduction from
$\mathsf{P}^{\mathrm{sp}}$ to the virtual-source protocol with the physical
receiver (Lemma~\ref{lem:source-transfer}\,(i)), and to $r_{\sq}$ on that
protocol. That $r_{\sq}$ is an exact operational reduction there is
Corollary~\ref{cor:sequential}: its first sentence is the equality of the two
output states, at every attack sequence of the type of
\eqref{eq:attackchannel}, and the strategies of the virtual-source protocol are
sequences of that type, \eqref{eq:srcstrategy} naming the first round's input
register and nothing else. Both reductions are over the same output registers,
which are those of $\Kfin$ together with $E_n$ in all three presentations, so
(R1) holds along the composition.

(ii) Fix $j\ge2$. Theorem~\ref{thm:licence}\,(i), which is its trusted-source
part, applied at $E'=E_{j-1}$ and $E''=E'_j$, gives
$\lambda_j\circ\Aev_j\in\AttSet_{\CList}(E_{j-1},\Qsys E'_j)$; for a list with
tagged entries it is its part (v), read at the trusted branch of
\eqref{eq:attackchannel}, which gives both the party-local condition
\eqref{eq:attsetpl} and the scalar one.

For $j=1$ the same parts are applied \emph{to the pre-composed attack}, and
that is why no hypothesis on $\Hs{\Asig}$ is needed. Write
$\Gamma\coloneqq\Xi^{\otimes n}\otimes\idmap_{E_0}$ and
$\Aev^{\mathrm{v}}_1\coloneqq\Aev_1\circ\Gamma$, so that
$\Aev^{\mathrm{v},\sq}_1=\lambda_1\circ\Aev^{\mathrm{v}}_1$ by
\eqref{eq:twosidedmap} and associativity of composition. Being a composition of
two channels, $\Aev^{\mathrm{v}}_1$ is a channel
$\Tcl{\Avirt^{n}E_0}\to\Tcl{\Bsys_1E'_1}$, of exactly the type
Theorem~\ref{thm:licence}\,(i) and (v) quantify over, and its
\emph{input} register $\Avirt^{n}E_0$ is finite-dimensional by hypothesis.
Applying that theorem at $E'=\Avirt^{n}E_0$ and $E''=E'_1$ therefore gives
$\Aev^{\mathrm{v},\sq}_1\in\AttSet_{\CList}(\Avirt^{n}E_0,\cdot)$ outright.
The separable system $\Asig^{n}$ occurs nowhere in this step: it is an
intermediate system of one channel (Remark~\ref{rem:source-transfer-registers}),
and the trusted branch of \eqref{eq:attackchannel} types an attack by the
systems it \emph{outputs}, the measured system and the outgoing memory,
placing no condition on the register the round receives beyond its being that
register.

\emph{The other order lands on the same set and needs a hypothesis this one
does not.} Applying Theorem~\ref{thm:licence} to $\Aev_1$ first requires
$E'=\Asig^{n}E_0$, which that theorem types finite-dimensional, and then
Lemma~\ref{lem:precomposition} at $E_{\mathrm{pre}}=\Avirt^{n}E_0$,
$E'=\Asig^{n}E_0$ and the same $\Gamma$, its part (i) and its part (iii) for
the tagged entries, gives the same membership. That lemma asks only that the
two input registers be separable, so what makes the second route the narrower
one is Theorem~\ref{thm:licence}'s own typing and nothing else; that the two
routes agree is (i) read at the level of admissible sets. The image of
$r_{\mathrm{two}}$ therefore lies in the strategy set of
$\mathsf{P}^{\mathrm{v},\sq}_{\CList}$. Equation~\eqref{eq:opreduction} holds
at every strategy by (i) and does not read the analysed set, and condition (R3)
of Definition~\ref{def:opreduction} now holds for the smaller set, so
$r_{\mathrm{two}}$ is an exact operational reduction into it.

(iii) Lemma~\ref{lem:opreduction-transfer}\,(i) applied to the reduction of
(ii), with $f$ the criterion at the adversary system of the presentation:
$\sup$ over the physical strategies equals $\sup$ over the image and is at most
$\sup$ over the $\CList$-admissible sequences, which is at most $\varepsilon$.
No property of the criterion is used, which is why the conclusion is
simultaneous in all of them.

(iv) By Lemma~\ref{lem:source-transfer}\,(ii) the identity map on strategies is
an exact operational reduction between the prepare-and-measure and the
entanglement-based presentations of the virtual protocol, in either direction,
and the squashing replacement touches only $\Bsys_j$, a system neither
description of the source names. Composing with the identity map changes
neither the strategies nor their images, so
Lemma~\ref{lem:opreduction-transfer}\,(v) carries (i) to (iii) across
unchanged.

(v) The closing paragraph of \xref{transmitter/def:source-segment} states
that when $L_{\mathrm s}=0$ the segment is empty, $\Xi$ is the identity,
$\Avirt=\Asig$ and $\rvirt_x=\rho_x$, so
$\Aev^{\mathrm{v}}_1=\Aev_1$ in \eqref{eq:srcmap-strategy} and
$r_{\mathrm{src}}$ is the identity; then \eqref{eq:twosidedmap} is
$\Aev^{\sq}_j=\lambda_j\circ\Aev_j$ of \eqref{eq:composite}. Part (i) is then
the output equality of Corollary~\ref{cor:sequential}, part (ii) is
Theorem~\ref{thm:licence}\,(i) or (v) per round, and part (iii) is that
corollary's last sentence as Theorem~\ref{thm:licence}\,(iii) states it. In
\emph{untrusted}-source mode there is no source segment to transfer and this
theorem asserts nothing, so what is left there is exactly those two claims,
read at the untrusted branch.
\end{proof}

% claim reductions/rem:two-sided-scope
\claimsummary{reductions/rem:two-sided-scope}{The two-sided reduction relates no marginals of the specified and virtual descriptions and licenses no entry whose right-hand side is not a real constant. Its order-independence concerns soundness; the order a chain records stays data.}
\claimbookmark{reductions/rem:two-sided-scope}{Remark}{rem:two-sided-scope}{Scope of the two-sided reduction}
\begin{remark}[Scope of the two-sided reduction]\label{rem:two-sided-scope}
Theorem~\ref{thm:two-sided} uses and implies no relation between a marginal
of the specified description and one of the virtual description, and none
would follow (Remark~\ref{rem:source-transfer-scope}). It licenses no entry
whose right-hand side is not a real constant:
Lemma~\ref{lem:precomposition}\,(iv) is a counterexample for such an entry,
and the admission rule of Definition~\ref{def:chain} refuses one beside a
non-empty source segment. And part (i) does not say that a chain may be
\emph{recorded} in either order: the record's order is data, stated as
condition (C3) of Definition~\ref{def:chain}, and what (i) says is that for a
licensed list the order is not what carries the soundness
(Remark~\ref{rem:ordering}).
\end{remark}
\claimmeta{\metalabel{Uses}\mref{Thm.}{thm:two-sided}
\metasep newly derived
\metasep accepted}

% claim reductions/rem:ordering
\claimsummary{reductions/rem:ordering}{Condition (C3) records the order of the source and receiver stages as data; for a list with real constants that order does not carry soundness (\hyperlink{claim-reductions/lem:precomposition}{Lemma~\ref*{lem:precomposition}}, \hyperlink{claim-reductions/thm:two-sided}{Theorem~\ref*{thm:two-sided}}). A general entry is not preserved and is refused beside a non-empty source segment.}
\claimbookmark{reductions/rem:ordering}{Remark}{rem:ordering}{Recorded stage order, immaterial to soundness for licensed lists}
\claimleadin{%
One of the artefact shapes the module's design fixes is that composition order
is data (\href{../../../spec/modules/reductions/design.md\#structural-design-decisions}{\nolinkurl{spec/modules/reductions/design.md}\texttt{\#structural-design-decisions}},
the decisions on the artefact's shape and on the two-sided chain), and
condition (C3) of Definition~\ref{def:chain} states it. This remark says what
that condition carries and what it does not.
}
\begin{remark}[Recorded stage order, immaterial to soundness for licensed lists]\label{rem:ordering}
\emph{The record.} Condition (C3) is a statement about the \emph{record}: the
order in which a reduction was performed is carried as data, so that a later
reader re-checks the path instead of inferring it from when an artefact was
built.

\emph{The inclusion a source link after a receiver link needs.} A \emph{source}-stage
link changes what Alice emits: it asserts that the states actually prepared
are the image, under a declared channel $\Gamma$, of simpler virtual ones, and
its security argument re-parameterises the first-round attack as
$\Aev_1=\Aev'_1\circ\Gamma$ ranging over \emph{all} CPTP maps. A
\emph{receiver}-stage link, by Theorem~\ref{thm:licence}, hands the analysis
an attack set that is a subset of the CPTP maps, in general a \emph{proper}
one. Applying a
source link after one is therefore not free: what must additionally hold is
the inclusion
$\AttSet^{\mathrm{real}}_1\circ\Gamma\subseteq\AttSet^{\mathrm{virtual}}_1$.

\emph{For a licensed list the inclusion holds, and the order is immaterial to
soundness.}
Lemma~\ref{lem:precomposition} establishes exactly that inclusion, for
\emph{every} channel $\Gamma$ and for every list whose right-hand sides are
real constants, which is every list this module emits
(Theorem~\ref{thm:licence}\,(iv)) and every list the consuming contract calls
licensed (Definition~\ref{def:licensed}), tagged entries included. For such a
list the order of the two stages is immaterial to soundness, and that is
proved: Theorem~\ref{thm:two-sided}\,(i), in trusted-source mode, the only
mode in which a source stage exists, shows that replacing the source first
and replacing the receiver first give one and the same map of
\termref{term-attack-strategy}{attack strategies}, because the two replacements act
on different systems, the \termref{term-source-composite}{source composite} at the \emph{input} of the
first round's attack and the \termref{term-squashing-channel}{squashing channel} at the \emph{output} of
each round's attack. Condition (C3) stays as the shape of the record and the
order in which the artefact presents its two segments; it does not carry the
soundness.

\emph{The published framework names both resolutions, and this contract takes
the second.} \cite[Rem.~9.2]{tupkary2026decoy} records the same interaction and
resolves it by applying its source-map lemma \emph{before} its squashing
lemma, which is the ordering resolution; the same remark names the
alternative, reproving the source-map statement for restricted attack sets,
and Lemma~\ref{lem:precomposition} is that alternative carried out here. No
step of the lemma appeals to the paper: it is proved from positivity and trace
preservation.

\emph{What is refused, and why.} The inclusion is established for a licensed
list and for no other. An entry whose right-hand side is a Hermitian operator
rather than a real constant, a \emph{general} entry, which
Definition~\ref{def:attset} records as admissible input to the consuming
contract and never licensed there, is not preserved:
Lemma~\ref{lem:precomposition}\,(iv) exhibits a channel and an entry that hold
before the composition and fail after it. At the state level such an entry
has no reading beside a non-empty source segment at all, its right-hand side
being an operator on a retained system the two descriptions do not share
(Remark~\ref{rem:precomposition-scope}). A chain with a non-empty source
segment carrying such an entry is therefore \emph{refused}, which is the
admission rule of Definition~\ref{def:chain}.

\emph{Where the rule is decided.} The rule is a condition on a declaration
that carries both a non-empty source segment and a constraint list, and it is
decided where the declaration is assembled, by the chain conditions of the
protocol compiler (\xref{protocol-compiler/def:chaincond}); this module states
the rule and decides no instance of it. It declares no general entry and
emits none, every entry it emits carrying a real scalar, so no artefact of
this module can violate the rule. The rule is not vacuous: a presented
declaration may carry a non-empty source segment. Both of its halves are
decided by the protocol compiler's chain conditions. The \emph{entry} half
refuses a general entry beside a non-empty source segment; the \emph{mode} half admits a non-empty source
segment only in trusted-source mode, and refuses a declaration that carries
one while its marginal slot holds the no-marginal value $\nomarg$. So neither
half rests on the provenance of the artefact that carried the segment.
\end{remark}
\claimmeta{\metalabel{Uses}\mref{Def.}{def:attset}, \mref{Def.}{def:licensed}, \mref{Thm.}{thm:licence}, \mref{Def.}{def:chain}, \mref{Lem.}{lem:precomposition}, \mref{Thm.}{thm:two-sided}
\metasep newly derived
\metasep accepted}

\section{The number model and exact decisions}\label{group-reductions/def:exactadm}
\sectionintro{Every number this module reads or emits, except the recorded \termref{term-emission-displacement}{emission displacement}, is a binary64 value read as the dyadic rational it is, every declared or emitted operator entry and declared real scalar satisfies one declared lattice condition, and every condition on declared data is decided exactly; the section states the strengthening of exact admissibility this module decides (\hyperlink{claim-reductions/def:exactadm}{Definition~\ref*{def:exactadm}}). The decisions are taken in an ordered field that the contract leaves open, since no admissible choice changes a decision (\hyperlink{claim-reductions/def:substrate}{Definition~\ref*{def:substrate}}). Positivity and operator-norm bounds reduce to sign comparisons of characteristic-polynomial coefficients (\hyperlink{claim-reductions/lem:exactpsd}{Lemma~\ref*{lem:exactpsd}}), and partial sums on the lattice are exact while the alphabet is small enough (\hyperlink{claim-reductions/lem:latsum}{Lemma~\ref*{lem:latsum}}).}
\sectionimports{\extref{../foundations/contract.pdf}{foundations/def:binary64}, \extref{../foundations/contract.pdf}{foundations/def:exact-admissibility}, \extref{../foundations/contract.pdf}{foundations/lem:binary64-lattice}}

% claim reductions/def:exactadm
\claimsummary{reductions/def:exactadm}{Every number this module reads or emits, except the recorded emission displacement, is a binary64 value read as its dyadic rational; every declared or emitted operator entry and declared real scalar meets one lattice condition. Exact admissibility here is the foundations core plus that condition.}
\claimbookmark{reductions/def:exactadm}{Definition}{def:exactadm}{The number model, the declared lattice and exact admissibility}
\claimleadin{%
Every operator datum this module consumes arrives as a finite list of
\exttermref{../foundations/contract.pdf}{term-binary64-numbers}{binary64 numbers}, and under the number model those numbers are the
model: a \exttermref{../foundations/contract.pdf}{term-declared-value}{declared value} \emph{is} its binary64 datum
(\xref{foundations/def:binary64}), declared operator data are
\exttermref{../foundations/contract.pdf}{term-exactly-admissible}{exactly admissible} at every \exttermref{../foundations/contract.pdf}{term-module-boundary}{module boundary}
(\xref{foundations/req:boundary-admission}), and every consumer decides that
exactly and refuses otherwise
(\xref{foundations/req:consumer-decides-exactly}). This definition states the
module's side of that policy, for what it reads and for what it emits: the
common core of exact admissibility, restated in this module's terms so that
both ends of one boundary name one condition, and strengthened by one
declared lattice. On operator entries the lattice condition
\eqref{eq:latinv} of (b) is condition (D-LAT) of
\xref{foundations/req:declared-lattice} at $k=\kappa$; this module states it
of the declared real scalars of (b) as well, as a condition of what it
admits. Together with \eqref{eq:latinv}, conditions (a)--(e) of (c) are, for
the family this module emits, the six conditions of
\xref{foundations/req:exact-emission}. Condition (e) is stated on the
additions because \xref{foundations/rem:computed-completeness} is a family
on which a vanishing completeness defect and completeness of the numbers
differ, and Remark~\ref{rem:exactboundary} discusses the objects that scope
clause (d) leaves to the hypotheses of an instance.
}
\begin{definition}[The number model, the declared lattice and exact admissibility]\label{def:exactadm}
\emph{(a) The model.} Every numerical datum this module reads or emits,
the displacement record below aside, is a binary64 value on the wire and
\emph{is} that value: no exact object is posited behind it and none is
approximated by it. A binary64 value is a dyadic
rational, and \emph{this contract reads it as that rational}
(\xref{foundations/def:binary64}\,(iii)). Every predicate this contract
decides is decided, and every emitted number but one is formed, by finitely
many operations in the substrate of Definition~\ref{def:substrate} and
Lemma~\ref{lem:exactpsd}. \emph{The exception is the \termref{term-emission-displacement}{emission displacement}
$\Delta_s$ of \eqref{eq:displacement}}: an entrywise maximum \emph{modulus},
hence the square root of an exact rational and in general irrational, emitted
as a binary64 rendering of a record of what the normalizer of
Definition~\ref{def:emitnorm} did, in a direction nothing records; it decides
nothing, and no statement of this contract reads it. That is the reading the
\exttermref{../foundations/contract.pdf}{term-exact-island}{exact island} (\xref{foundations/def:exact-island}) gives, which
contains both the admission decisions on declared blocks and a producer's
construction of what it emits. The only non-numerical data are combinatorial
(alphabets, orders, block identifiers, a coarse-graining map, a branch index
set, a cutoff) and carry no arithmetic.

\emph{(b) The declared lattice.} The request declares one integer exponent
$\kappa$, and every declared and every emitted operator entry, together with
every \emph{real scalar} the request declares of the following kinds --- the
\termref{term-class-weights}{class weights} $q_t$, any certified weight, the entries of a declared
stochastic matrix, and any declared device constant --- satisfies
\begin{equation}\label{eq:latinv}
  2^{\kappa}x\in\mathbb{Z},
  \qquad |x|\le1,
  \qquad 1\le\kappa\le53 .
\end{equation}
Such an $x$ is exactly a binary64 number
(\xref{foundations/lem:binary64-lattice}\,(i)), its significand carrying at
most $\kappa\le53$ bits. The list is exhaustive: the constant a caller declares for a
requested constraint entry is not among them, and where it is off the
lattice it is rounded onto it or refused (Corollary~\ref{cor:entries}\,(iii)). It is a condition on a \emph{set} of values: it says which numbers
may be declared and nothing about how near one number is to another. A
complex entry satisfies \eqref{eq:latinv} when its real and imaginary parts
do.

\emph{(c) Exact admissibility on the declared lattice.} Let
$\{A_v\}_{v\in I}$ be a finite family of matrices presented by its blocks
over declared bases, indexed by the finite set $I$ the invoking rule names
(the \termref{term-outcome-alphabet}{outcome alphabet} at a fixed setting for a declared block family, a class
$\Vset_t$ of the datum of a \termref{term-reduction-instance}{reduction instance}
(Definition~\ref{def:reduction-instance}), $\Vset$ for the family this module
emits), and let $c$ be the constant the rule names for its completeness
identity: $1$ for a POVM, the class weight $q_t$ for a
class of a datum. The family is
\termdef{term-exactly-admissible-on-the-declared-lattice}{exactly admissible on the declared lattice}
when
\begin{enumerate}[label=(\alph*),leftmargin=2.4em]
\item \emph{(the \exttermref{../foundations/contract.pdf}{term-sector-map}{sector map})} the block kind, the block identifiers and their
  order are the ones the consumer holds, a comparison of declared labels,
  decided and never measured;
\item \emph{(Hermiticity)} $A_v=A_v^{\dagger}$, entry for entry;
\item \emph{(positivity)} $A_v\succeq0$, decided by \eqref{eq:psdtest};
\item \emph{(completeness)} $\sum_{v\in I}A_v=c\,\idon{}$ on each block, as an
  identity of the entries;
\item \emph{(the sum \exttermref{../foundations/contract.pdf}{term-rounds-nowhere}{rounds nowhere})} every entrywise binary64
  addition of that accumulation, taken in the \exttermref{../foundations/contract.pdf}{term-declared-order}{declared order} of $I$,
  is exact, so that the \emph{computed} sum \emph{is} the sum of (d);
\end{enumerate}
and every entry satisfies \eqref{eq:latinv}. Conditions (a)--(e) say that
the family is exactly admissible with constant $c$ for the sector map the
consumer holds (\xref{foundations/def:exact-admissibility}), each block being
presented at the dimension the consumer's map declares for it. Condition (e)
is a condition in its own right and is stated on the \emph{additions}: a
vanishing completeness \emph{defect} says that the computed sum is the
identity, not that $\sum_vA_v=c\,\idon{}$ holds of the numbers.
Lemma~\ref{lem:latsum} shows that (e) is satisfiable on this
contract's own emission, rather than a condition that would refuse every
family.

\emph{(d) Scope, stated once and binding wherever this definition is
invoked.} This definition and the rules that invoke it are about data this
contract requires to be presented as a finite list of matrices of declared
dimension: the \termref{term-data-sector}{data sectors} of a reduction instance
(Definition~\ref{def:reduction-instance}\,(I1)) and the declared blocks they
are formed from; the declared classical scalars and matrices a submodule's
route reads (a stochastic matrix and its descended form, the class weights,
each certified weight, any declared device constant); and the family this
module emits. They are \emph{not} available for the objects an instance's
\termref{term-hypotheses-h}{hypotheses (H)} are about: operators on a part of $\Hs{\Bsys}$ that no
finite list of matrices presents, an above-cutoff family, a register POVM, an
uncharacterised branch. There is nothing on those to decide, and their
conditions remain hypotheses of the instance
(Definition~\ref{def:reduction-instance}\,(I1)).
\end{definition}
\claimmeta{\metalabel{Uses}\extref{../foundations/contract.pdf}{foundations/def:binary64}, \extref{../foundations/contract.pdf}{foundations/def:exact-admissibility}, \extref{../foundations/contract.pdf}{foundations/lem:binary64-lattice}
\metasep \metalabel{Used by}\mref{Def.}{def:joint-family}, \mref{Def.}{def:substrate}, \mref{Lem.}{lem:latsum}, \mref{Def.}{def:reduction-instance}, \mref{Cor.}{cor:identity-instance}, \mref{Cor.}{cor:flag-instance}, \mref{Thm.}{thm:instrument-instance}, \mref{Def.}{def:emitnorm}, \mref{Thm.}{thm:emit}, \mref{Cor.}{cor:entries}, \mref{Cor.}{cor:emitted-identity}, \mref{Def.}{def:weight-route}, \mref{Prop.}{prop:flat-instance}, \mref{Cor.}{cor:emitted-flat}
\metasep newly derived
\metasep accepted}

% claim reductions/def:substrate
\claimsummary{reductions/def:substrate}{Every exact decision is taken in an ordered field F between \ensuremath{\mathbb{Q}} and \ensuremath{\mathbb{R}}, with its Gaussian extension F + iF, and the field is left open. No admissible choice changes a decision, every number a rule reads being rational, so \ensuremath{\mathbb{Q}} always suffices.}
\claimbookmark{reductions/def:substrate}{Definition}{def:substrate}{The arithmetic substrate}
\begin{definition}[The arithmetic substrate]
\label{def:substrate}
The
\termdef{term-substrate}{substrate}
of this contract is an ordered field $\Fld$ with
$\mathbb{Q}\subseteq\Fld\subseteq\mathbb{R}$ together with its Gaussian
extension $\Fldi\coloneqq\Fld+i\Fld$, which is closed under complex
conjugation. What is fixed is the \emph{predicate} each rule decides and the
requirement that it be decided exactly; which such field carries the
arithmetic is not. The choice cannot change any decision. Binary64 values are
dyadic rationals (Definition~\ref{def:exactadm}\,(a)), and every number a rule
reads is formed from the declared entries and from rational constants by
field operations, floors of rationals and roundings of rationals to a
lattice, so it is rational; since $\mathbb{Q}$ is an ordered subfield of every
admissible $\Fld$ with the induced order, and $\mathbb{Q}+i\mathbb{Q}$ a
subfield of $\Fldi$, each decision, a comparison or an equality of such
numbers, is the same decision taken over $\mathbb{Q}$. In particular
$\Fld=\mathbb{Q}$ always suffices. The latitude is bounded: it is a choice of
arithmetic substrate, not of mathematics.
\end{definition}
\claimmeta{\metalabel{Uses}\mref{Def.}{def:exactadm}
\metasep \metalabel{Used by}\mref{Lem.}{lem:exactpsd}, \mref{Thm.}{thm:emit}, \mref{Cor.}{cor:reference-companion}, \mref{Req.}{req:exact-admission-refusal}, \mref{Def.}{def:weight-route}, \mref{Prop.}{prop:route}, \mref{Cor.}{cor:emitted-flat}, \mref{Cor.}{cor:two-sector-decided}, \mref{Cor.}{cor:flat-companion}
\metasep newly derived
\metasep accepted}

% claim reductions/lem:exactpsd
\claimsummary{reductions/lem:exactpsd}{A Hermitian matrix with entries in the substrate is positive semidefinite exactly when, for every k, the sum of its principal minors of order k is non-negative, and an operator-norm bound reduces to such a test. Both are decided by field operations and sign comparisons, with no eigenvalue computed.}
\claimbookmark{reductions/lem:exactpsd}{Lemma}{lem:exactpsd}{Exact decision of positivity and of an operator-norm bound}
\claimleadin{%
An eigenvalue of a matrix with rational entries is in general irrational, so
an exact decision of positivity cannot go through one. The lemma replaces it
by the signs of the coefficients of the characteristic polynomial, which lie
in the substrate.
}
\begin{lemma}[Exact decision of positivity and of an operator-norm bound]
\label{lem:exactpsd}
Let $\Fld,\Fldi$ be the substrate of Definition~\ref{def:substrate} and let
$d,d'\ge1$.
\begin{enumerate}[label=(\roman*),leftmargin=3.0em]
\item Let $X=X^{\dagger}\in\Lin{\mathbb{C}^{d}}$ have all entries in $\Fldi$ and
  write
  \begin{equation}\label{eq:charpoly}
    \det\big(\lambda\,\idon{d}+X\big)=\sum_{k=0}^{d}e_k(X)\,\lambda^{d-k},
    \qquad e_0(X)=1 ,
  \end{equation}
  so that $e_k(X)$ is the sum of the $k\times k$ principal minors of $X$. Then
  every $e_k(X)$ lies in $\Fld$, and
  \begin{equation}\label{eq:psdtest}
    X\succeq0
    \quad\Longleftrightarrow\quad
    e_k(X)\ge0\ \text{ for every }k=1,\dots,d .
  \end{equation}
\item Let $Z\in\Linn{\mathbb{C}^{d'}}{\mathbb{C}^{d}}$ have all entries in $\Fldi$ and
  let $\delta\in\Fld$, $\delta\ge0$. Then
  \begin{equation}\label{eq:normtest}
    \opn{Z}\le\delta
    \quad\Longleftrightarrow\quad
    \delta^{2}\,\idon{d'}-Z^{\dagger}Z\;\succeq\;0 ,
  \end{equation}
  the right-hand side being a test of the form (i) on a Hermitian matrix with
  entries in $\Fldi$.
\end{enumerate}
Both predicates are therefore decided by finitely many field operations and
sign comparisons in $\Fld$: no rounding, no threshold, and \emph{no
eigenvalue is computed}.
\end{lemma}
\claimmeta{\metalabel{Uses}\mref{Def.}{def:substrate}
\metasep \metalabel{Used by}\mref{Rem.}{rem:psdalgorithm}, \mref{Def.}{def:reduction-instance}, \mref{Def.}{def:emitnorm}, \mref{Thm.}{thm:emit}, \mref{Cor.}{cor:emitted-identity}, \mref{Rem.}{rem:exactboundary}, \mref{Def.}{def:weight-route}, \mref{Prop.}{prop:route}, \mref{Cor.}{cor:oneabsorbing}, \mref{Cor.}{cor:emitted-flat}, \mref{Cor.}{cor:two-sector-decided}
\metasep newly derived
\metasep accepted}

\begin{proof}
(i) A principal submatrix of a Hermitian matrix is Hermitian, so each principal
minor is real; it is also a polynomial with integer coefficients in entries of
$\Fldi$, hence lies in $\Fldi$. Being both real and in $\Fldi$, it lies in
$\Fld$, and so does each $e_k(X)$. The identification of $e_k(X)$ with the
coefficient in \eqref{eq:charpoly} is the standard expansion of the
characteristic polynomial by principal minors.

Let $\lambda_1,\dots,\lambda_d\in\mathbb{R}$ be the eigenvalues of $X$ with
multiplicity, so that
$\det(\lambda\idon{d}+X)=\prod_{j=1}^{d}(\lambda+\lambda_j)$ and $e_k(X)$ is
the $k$-th elementary symmetric polynomial in them. If $X\succeq0$ then every
$\lambda_j\ge0$ and every $e_k(X)$, a sum of products of $\lambda_j$'s, is
$\ge0$. Conversely suppose $e_k(X)\ge0$ for every $k$ and that some
$\lambda_{j_0}<0$. Put $\lambda\coloneqq-\lambda_{j_0}>0$. On the one hand
$\prod_j(\lambda+\lambda_j)=0$, because the factor $j_0$ vanishes. On the other
hand
$\sum_{k=0}^{d}e_k(X)\lambda^{d-k}\ge e_0(X)\lambda^{d}=\lambda^{d}>0$, every
omitted term being non-negative. The two sides of \eqref{eq:charpoly} at that
$\lambda$ are therefore $0$ and something strictly positive, a
contradiction. Hence no eigenvalue is negative and $X\succeq0$.

(ii) $Z^{\dagger}Z$ is Hermitian with entries in $\Fldi$, and
$\opn{Z}^{2}=\lmax(Z^{\dagger}Z)$ by definition of the operator norm. So
$\opn{Z}\le\delta$ iff $\lmax(Z^{\dagger}Z)\le\delta^{2}$ iff
$\delta^{2}\idon{d'}-Z^{\dagger}Z\succeq0$.
\end{proof}

% claim reductions/rem:psdalgorithm
\claimsummary{reductions/rem:psdalgorithm}{Any procedure that decides the positivity test of \hyperlink{claim-reductions/lem:exactpsd}{Lemma~\ref*{lem:exactpsd}} exactly in the substrate conforms, such as the Faddeev\textendash{}Le Verrier recursion or a rational LDL factorisation with the zero-pivot test. A procedure that decides it approximately does not, whatever its error bound.}
\claimbookmark{reductions/rem:psdalgorithm}{Remark}{rem:psdalgorithm}{Algorithms that decide the positivity test}
\begin{remark}[Algorithms that decide the positivity test]
\label{rem:psdalgorithm}
Lemma~\ref{lem:exactpsd} fixes \emph{what} is decided, not \emph{how}. The
coefficients $e_k(X)$ are produced exactly by the Faddeev--Le~Verrier
recursion or by any division-free characteristic-polynomial method. A
rational $LDL^{\dagger}$ factorisation with symmetric pivoting and the
zero-pivot row test decides the same predicate, often faster, in $O(d^{3})$
field operations. Enumerating all $2^{d}-1$ principal minors also decides it
and is never necessary. Any procedure that decides \eqref{eq:psdtest} exactly
in $\Fld$ conforms to this contract; \emph{one that decides it approximately
does not, whatever its error bound}. A floating-point eigenvalue is such a
procedure, which is why \xref{foundations/req:exact-emission}\,(iii) excludes
one by name.
\end{remark}
\claimmeta{\metalabel{Uses}\mref{Lem.}{lem:exactpsd}
\metasep newly derived
\metasep accepted}

% claim reductions/lem:latsum
\claimsummary{reductions/lem:latsum}{If \ensuremath{\kappa} + \ensuremath{\lceil}log\ensuremath{_{2}} L\ensuremath{\rceil} \ensuremath{\le} 53, every partial sum of L entries on the declared lattice is a binary64 number and every addition forming it is exact. A family on the lattice over such an alphabet therefore meets the rounds-nowhere condition by construction.}
\claimbookmark{reductions/lem:latsum}{Lemma}{lem:latsum}{Partial sums on the declared lattice are exact}
\claimleadin{%
A sum \exttermref{../foundations/contract.pdf}{term-rounds-nowhere}{rounds nowhere} exactly when every partial sum is a binary64
number (\xref{foundations/lem:exact-accumulation}). On the declared lattice
that holds as long as the alphabet is small enough for the partial sums to
fit in the significand.
}
\begin{technicalclaim}
\begin{lemma}[Partial sums on the declared lattice are exact]\label{lem:latsum}
Let $x_1,\dots,x_L$ satisfy \eqref{eq:latinv} and let
$\Sigma_j\coloneqq\sum_{i\le j}x_i$, $1\le j\le L$, be the partial sums of the
accumulation in that order. If
\begin{equation}\label{eq:latsumcond}
  \kappa+\big\lceil\log_2L\big\rceil\ \le\ 53 ,
\end{equation}
then every $\Sigma_j$ is exactly a binary64 number and every one of the $L-1$
additions forming it is exact in binary64. Consequently a family whose entries
satisfy \eqref{eq:latinv}, over an alphabet satisfying
\eqref{eq:latsumcond}, meets condition~(e) of Definition~\ref{def:exactadm} by
construction.
\end{lemma}
\end{technicalclaim}
\claimmeta{\metalabel{Uses}\mref{Def.}{def:exactadm}
\metasep \metalabel{Used by}\mref{Thm.}{thm:emit}
\metasep newly derived
\metasep accepted}

\begin{proof}
Work one real component of one entry position at a time; the accumulation is
entrywise and the real and imaginary parts are accumulated independently. Each
$x_i$ is an integer multiple of $2^{-\kappa}$, so each $\Sigma_j$ is one too,
and $|\Sigma_j|\le\sum_{i\le j}|x_i|\le L$. Hence $2^{\kappa}\Sigma_j$ is an
integer with $\big|2^{\kappa}\Sigma_j\big|\le L\,2^{\kappa}\le2^{53}$ by
\eqref{eq:latsumcond}. Every integer of modulus at most $2^{53}$ is exactly
representable in binary64, and multiplication by the power of two
$2^{-\kappa}$ is exact and changes no significand, so $\Sigma_j$ is a binary64
number; and $2^{-\kappa}\ge2^{-53}$ with $|\Sigma_j|\le L$ keeps it far
inside the normal range, so no subnormal or overflow case arises. An IEEE-754
addition returns the correctly rounded value of the exact sum of its operands;
here that exact sum is $\Sigma_j$, which is representable, so the rounding is
the identity and the addition is exact.
\end{proof}

\section{Reduction instances}\label{group-reductions/def:reduction-instance}
\sectionintro{A \termref{term-reduction-instance}{reduction instance} packages what every construction emits: a \termref{term-handed-model}{handed model} with a finite datum and hypotheses about the rest, a target built block by block from the datum and \termref{term-scalar-value}{scalar values}, a \termref{term-squashing-relation}{squashing relation} with a \termref{term-transported-set}{transported set}, and \termref{term-certified-entries}{certified entries}, all holding at every \termref{term-admissible-datum}{admissible datum} (\hyperlink{claim-reductions/def:reduction-instance}{Definition~\ref*{def:reduction-instance}}). The core proves the three instances a protocol needs regardless of device: the identity instance, the flag instance of an uncharacterised branch, and the branching instance of a \termref{term-generalized-instrument}{generalized instrument} with its register form (\hyperlink{claim-reductions/thm:instrument-instance}{Theorem~\ref*{thm:instrument-instance}}). An instance precomposed with a channel on its system, such as a known loss in front of the device, is again an instance with the same datum, target, entries and emission (\hyperlink{claim-reductions/lem:precomposed-instance}{Lemma~\ref*{lem:precomposed-instance}}).}

% claim reductions/def:reduction-instance
\claimsummary{reductions/def:reduction-instance}{A reduction instance is a handed model with a finite datum and hypotheses about the rest, a target built block by block from the datum and scalar values, a squashing relation with a transported set, and certified entries, all holding at every admissible datum. Exact emission is proved once over it.}
\claimbookmark{reductions/def:reduction-instance}{Definition}{def:reduction-instance}{Reduction instance}
\claimleadin{%
The emission theorem of this module (Theorem~\ref{thm:emit}) is stated over
the object defined here, in the way the receiver states its emitted bound
over a \exttermref{../receiver/contract.pdf}{term-setting-family}{setting family} (\xref{receiver/thm:emitted-bound}). The object
packages what every construction this module or a submodule emits has in
common: a handed model on $\Bsys$ with a finite datum, a finite target built
from that datum and from scalar values, an existence theorem giving the
\termref{term-squashing-relation}{squashing relation} with a transported set, and a list of certified
entries.
The obligation of exact emission is then proved once, and each construction
is shown to be an instance where it lives. Because every part is established
at every admissible datum, the emission theorem can apply the existence
theorem to the datum the normalizer declares
(Definition~\ref{def:emitnorm}, Theorem~\ref{thm:emit}\,(iii)) without
re-proving anything, the role the \exttermref{../receiver/contract.pdf}{term-exact-flag}{exact flag} plays for the receiver's
setting family. The core establishes three instances a protocol needs
regardless of device: the identity instance
(Corollary~\ref{cor:identity-instance}), the flag instance of an
uncharacterised branch (Corollary~\ref{cor:flag-instance}) and the branching
instance (Theorem~\ref{thm:instrument-instance}) with its register form
(Corollary~\ref{cor:register-instance}); a squasher submodule proves that its
construction is an instance and applies Theorem~\ref{thm:emit}. The downward
rounding of a scalar value in (I2) is the rounding direction of
\xref{foundations/req:rounding-direction}, and Remark~\ref{rem:unsound} shows
that the operator form of the \termref{term-class-marginal-identity}{class-marginal identity} on the target, which
the fourth consequence below sets aside, is unsound.
}
\begin{definition}[Reduction instance]\label{def:reduction-instance}
Fix the declared \exttermref{../foundations/contract.pdf}{term-lattice-exponent}{lattice exponent} $\kappa$ of
Definition~\ref{def:exactadm}\,(b) and a \termref{term-joint-alphabet}{joint alphabet} with
\termref{term-classes}{classes} and \termref{term-class-weights}{class weights} $\big(\Vset,\{\Vset_t\}_{t\in\Tset},\{q_t\}\big)$
(Definition~\ref{def:joint-family}\,(a)). Let $\SecSet_{\mathrm{dat}}$ be a
finite ordered list of
\termdef{term-data-sector}{data sector}s,
each $s$ with a declared dimension $d_s\ge1$ and a declared
\termdef{term-support}{support}
$\Vset^{(s)}\subseteq\Vset$ that meets every class. A
\termdef{term-datum}{datum}
for these sectors is a family
$A=\{A^{(s)}_v\}_{s\in\SecSet_{\mathrm{dat}},\,v\in\Vset}$ of $d_s\times d_s$
matrices with entries in $\Fldi$ and $A^{(s)}_v=0$ for
$v\notin\Vset^{(s)}$: these elements lie
\termdef{term-outside-the-sector-support}{outside the sector support},
and no step of this contract reads, rounds or repairs them. It is an
\termdef{term-admissible-datum}{admissible datum},
written $A\in\mathcal{D}_{\mathrm{adm}}$, when for every $s$:
(a) $A^{(s)}_v=(A^{(s)}_v)^{\dagger}$ entry for entry; (b)
$A^{(s)}_v\succeq0$, decided by \eqref{eq:psdtest}; (c) it is
\termdef{term-class-complete}{class complete},
\begin{equation}\label{eq:classcomplete}
  \sum_{v\in\Vset_t}A^{(s)}_v\;=\;q_t\,\idon{d_s}
  \qquad\text{exactly, for every }t\in\Tset ;
\end{equation}
and it is an
\termdef{term-exactly-admissible-datum}{exactly admissible datum},
written $A\in\mathcal{D}_{\kappa}$, when in addition (d) every entry
satisfies \eqref{eq:latinv}. These are conditions (b)--(d) of
Definition~\ref{def:exactadm}\,(c) and the lattice invariant, read per data
sector with the class weight $q_t$ as the constant of the completeness
identity.

A
\termdef{term-reduction-instance}{reduction instance}
$\mathcal{J}$ over $\big(\Vset,\{\Vset_t\},\{q_t\}\big)$ consists of the
following four parts.
\begin{enumerate}[label=(I\arabic*),leftmargin=3.0em]
\item \emph{(the handed model and its datum)} A system $\Bsys$ with
  $\Hs{\Bsys}$ separable (Definition~\ref{def:system-setting}); data sectors
  $\SecSet_{\mathrm{dat}}$ with their declared dimensions and supports; a
  \termdef{term-handed-datum}{handed datum}
  $A\in\mathcal{D}_{\mathrm{adm}}$, every entry of which is formed from the
  numbers the instance was handed by finitely many operations in $\Fldi$; a
  finite list of
  \termdef{term-hypotheses-h}{hypotheses (H)}
  about the rest of the model (operators without a finite presentation,
  facts about a part of $\Hs{\Bsys}$ the datum does not describe, a register
  POVM), each of which is either decided exactly on \exttermref{../foundations/contract.pdf}{term-declared-data}{declared data}, the
  request being refused where it fails, or carried as an assumption of the
  emitted ledger; and a
  \termdef{term-model-map}{model map}
  $A'\mapsto M[A']=\{M_v[A']\}_{v\in\Vset}$ that assigns to every admissible
  datum $A'\in\mathcal{D}_{\mathrm{adm}}$ a POVM on $\Bsys$ over $\Vset$,
  under (H). The
  \termdef{term-handed-model}{handed model}
  is $M[A]$, the model map at the handed datum. The
  \termdef{term-stability-clause}{stability clause}
  requires that no hypothesis in (H) read an entry of the datum, so that (H)
  holds of $M[A']$ for every $A'\in\mathcal{D}_{\mathrm{adm}}$ as soon as it
  holds of the handed model; and that parts (I2)--(I4) be established
  for \emph{every} $A'\in\mathcal{D}_{\mathrm{adm}}$ under (H), not only at
  the handed datum.
\item \emph{(the target)} A finite ordered
  \termdef{term-sector-set}{sector set}
  $\SecSet=\SecSet_{\mathrm{dat}}\sqcup\SecSet_{\mathrm{sc}}$, the second
  summand the
  \termdef{term-scalar-sector}{scalar sector}s,
  each $s\in\SecSet$ with a declared dimension $d_s\ge1$; for each scalar
  sector $s$ and each $v\in\Vset$ a
  \termdef{term-scalar-value}{scalar value}
  $\phi_s(v)\in[0,1]$ satisfying \eqref{eq:latinv} with
  $\sum_{v\in\Vset}\phi_s(v)=1$ exactly, each formed from the numbers the
  instance was handed by finitely many operations in $\Fld$ and, where a
  rounding onto the lattice is needed, by the \emph{downward} rounding of a
  quantity the value is a certified lower bound of, any complementary value
  being recomputed from the rounded ones and never rounded on its own; the
  target space $\Hs{\Qsys}\coloneqq\bigoplus_{s\in\SecSet}\mathbb{C}^{d_s}$
  with sector projectors $\Pi_s$, so that $\dim\Hs{\Qsys}=\sum_sd_s$; and the
  \termdef{term-target-family}{target family}
  \begin{equation}\label{eq:instance-target}
    F_v[A']\;\coloneqq\;\bigoplus_{s\in\SecSet}F^{(s)}_v[A'],
    \qquad
    F^{(s)}_v[A']\;\coloneqq\;
    \begin{cases}
      A'^{(s)}_v & s\in\SecSet_{\mathrm{dat}},\\[2pt]
      \phi_s(v)\,\idon{d_s} & s\in\SecSet_{\mathrm{sc}},
    \end{cases}
  \end{equation}
  assigned one sector block at a time, so that $F_v[A']$ is
  sector-block-diagonal \emph{by construction}: no off-sector entry is
  written and none is to be verified. The \exttermref{../foundations/contract.pdf}{term-sector-map}{sector map} of the instance
  is the list $\SecSet$ with its order, each sector's dimension and its kind
  (data or scalar), and it is what a consumer compares
  (Definition~\ref{def:exactadm}\,(c), condition (a)). The
  \termdef{term-default-party-partition}{default party partition}
  of the target assigns every factor of every $\Hs{\Qsys_s}$,
  $s\in\SecSet$, to \emph{Bob}: \emph{every sector of the target is Bob's}.
  That is the assignment the physical model makes: the target describes what
  reaches Bob's receiver, the handed model living on his system $\Bsys$. An
  instance may in addition be declared with \emph{another party's label},
  one label for the whole instance, and then the same assignment is made
  with that label in its place: every factor of every $\Hs{\Qsys_s}$ is that
  party's, because the handed model lives on the system \emph{that} party
  receives. Where no label is declared the label is Bob's. Nothing else in
  the instance depends on the label: it is a declared identifier, compared
  and never computed, and no condition of (I1)--(I4) reads it.
\item \emph{(the existence theorem)} A channel
  $\Lambda[A']=\PinS\circ\Lambda_0[A']\colon\Tcl{\Bsys}\to\Tcl{\Qsys}$,
  given by a constructor and recorded as one, with $\PinS$ the sector
  pinching of Proposition~\ref{prop:pinching} for $\{\Pi_s\}_{s\in\SecSet}$;
  and a finite
  \termdef{term-transported-set}{transported set}
  $\hat O[A']$ of sector-block-diagonal Hermitian operators on $\Qsys$
  containing every sector projector $\Pi_s$, with
  \termref{term-pull-back}{pull-backs} $O[A']$; such that
  \begin{equation}\label{eq:instance-relation}
    \big(M[A'];\,O[A']\big)\ \dsq_{\Lambda[A']}\ \big(F[A'];\,\hat O[A']\big)
    \qquad\text{for every }A'\in\mathcal{D}_{\mathrm{adm}}\text{ under (H)}.
  \end{equation}
  By linearity of $\Lambda[A']^{\dagger}$, every real linear combination of
  transported operators is transported to the same combination of their
  pull-backs.
\item \emph{(the certified entries)} A finite list
  $\CList[A']=\{(\hat O_k,\bowtie_k,c_k)\}_{k}$ of
  \termdef{term-certified-entries}{certified entries},
  each $\hat O_k$ a real linear combination of transported operators of the
  form
  \begin{equation}\label{eq:instance-entry}
    \hat O_k\;=\;\sum_{v\in\Vset}\sum_{s\in\SecSet}
      \alpha_{k,v,s}\,\Pi_sF_v[A']\Pi_s\;+\;Z_k ,
  \end{equation}
  with real coefficients $\alpha_{k,v,s}$ \emph{on the lattice}, each an
  integer multiple of $2^{-\kappa}$, and $Z_k$ a real combination of sector
  projectors with coefficients of the same kind, $\bowtie_k\in\{\preceq,=,\succeq\}$
  and $c_k$ a declared constant on the lattice, whose \termref{term-certificate}{certificate}
  \eqref{eq:certificate},
  $\Lambda[A']^{\dagger}(\hat O_k)\bowtie_kc_k\,\idon{\Bsys}$, is
  \emph{discharged} for every $A'\in\mathcal{D}_{\mathrm{adm}}$ under (H); and
  a
  \termdef{term-class-marginal-flag}{class-marginal flag}
  recording whether the \termref{term-class-marginal-identity}{class-marginal identity}
  \begin{equation}\label{eq:instance-classmarginal}
    \sum_{v\in\Vset_t}M_v[A']\;=\;q_t\,\idon{\Bsys}
    \qquad\text{for every }t\in\Tset
  \end{equation}
  is \emph{discharged}, that is, established by the construction for every
  $A'\in\mathcal{D}_{\mathrm{adm}}$ under (H), or not.
  The modulus bound of \eqref{eq:latinv} is not imposed on a coefficient:
  $|x|\le1$ is a condition on a declared or emitted \emph{entry} (a value of
  the datum, a scalar value, a declared scalar or constant), where a POVM
  element and a completeness identity give it, and not on the coefficients
  of a linear combination, which a branching instance forms as the
  difference $\beta_s-c_k$ of two lattice multiples
  (Theorem~\ref{thm:instrument-instance}), again a lattice multiple;
  \emph{being on the lattice is all that is asked of a coefficient}. What
  \eqref{eq:latinv} does bind is the entries of the operator the combination
  produces (Corollary~\ref{cor:entries}\,(iii)).
\end{enumerate}
The instance \emph{at} a datum $A'\in\mathcal{D}_{\mathrm{adm}}$ is the same
$\Bsys$, (H), sectors, supports, scalar values, constructor and entry rule
with $A'$ in place of the handed datum; by the stability clause it is a
reduction instance whenever $\mathcal{J}$ is.

\emph{Five consequences, each immediate.} \emph{First}, for every
$A'\in\mathcal{D}_{\mathrm{adm}}$ the target $\{F_v[A']\}_{v\in\Vset}$ is a
POVM on $\Qsys$: every sector block is Hermitian and positive semidefinite,
by (a) and (b) on a data sector and by $\phi_s(v)\in[0,1]$ on a scalar one;
and it is complete, on a data sector because
$\sum_{v\in\Vset}A'^{(s)}_v=\sum_{t}\sum_{v\in\Vset_t}A'^{(s)}_v
=\sum_tq_t\idon{d_s}=\idon{d_s}$ by \eqref{eq:classcomplete} and
$\sum_tq_t=1$, on a scalar sector because $\sum_v\phi_s(v)=1$. Its data
sectors \emph{are} the datum, entry by entry, so the emitted target carries
the datum and nothing else of the model. \emph{Second}, every $F_v[A']$,
every $\Pi_s$ and every transported operator is sector-block-diagonal, so
Corollary~\ref{cor:blockdiag}\,(i) applies: a relation established for
$\Lambda_0[A']$ holds for $\PinS\circ\Lambda_0[A']$ with the same target,
the same transported set and the same pull-backs, which is why (I3) may be
stated with the pinching in front, and the squashed analysis may be
restricted to block-diagonal states without excluding a reachable one
(Corollary~\ref{cor:blockdiag}\,(iii)). \emph{Third}, the pull-back of every
sector projector is positive semidefinite, $\Lambda[A']^{\dagger}(\Pi_s)
\succeq0$, because the adjoint of a channel is a positive map, and
$\sum_s\Lambda[A']^{\dagger}(\Pi_s)=\idon{\Bsys}$ by unitality; so
$(\Pi_s,\succeq,0)$ carries the certificate of
Corollary~\ref{cor:pullback}\,(ii) for \emph{every} instance and needs no
part of (I4) to say so. \emph{That certificate travels with the transported
set, whatever the request carries}, which is where (I3) puts every sector
projector and where the artefact records the certificate case per entry; it
becomes an \emph{entry} of $\CList$ only where a caller requests it \emph{as
an entry} (Corollary~\ref{cor:entries}\,(i)), the emitted list being empty
where the request carries no entry. \emph{Fourth}, the relation pulls a class
sum of the target back to the class sum of the model,
$\Lambda[A']^{\dagger}\big(\sum_{v\in\Vset_t}F_v[A']\big)
=\sum_{v\in\Vset_t}M_v[A']$, by linearity; so
\eqref{eq:instance-classmarginal} holds of the model exactly when the
pull-back identity $\Lambda[A']^{\dagger}\big(\sum_{v\in\Vset_t}F_v[A']\big)
=q_t\idon{\Bsys}$ does, which is the expectation-form entry
Corollary~\ref{cor:entries}\,(ii) emits where the flag says it is discharged,
and never the operator identity
$\sum_{v\in\Vset_t}F_v[A']=q_t\idon{\Qsys}$, which fails on every scalar
sector whose values do not sum to $q_t$ over $\Vset_t$. \emph{Fifth},
$\dim\Hs{\Qsys}=\sum_sd_s$ is a sum of declared dimensions, finite, and the
sector map is a comparison of declared labels; both are read off the
declaration with no arithmetic performed.

The link of Definition~\ref{def:link} an instance produces is the relation
\eqref{eq:instance-relation} at the \termref{term-declared-datum}{declared datum}, with the transported set
and the entry list as its slots, and its chain (Definition~\ref{def:chain})
is that link preceded by whatever post-processing the instance's submodule
applies before it.
\end{definition}
\claimmeta{\metalabel{Uses}\mref{Def.}{def:squashing}, \mref{Cor.}{cor:pullback}, \mref{Def.}{def:certificate}, \mref{Def.}{def:joint-family}, \mref{Prop.}{prop:pinching}, \mref{Cor.}{cor:blockdiag}, \mref{Thm.}{thm:licence}, \mref{Def.}{def:exactadm}, \mref{Lem.}{lem:exactpsd}
\metasep \metalabel{Used by}\mref{Rem.}{rem:unsound}, \mref{Cor.}{cor:identity-instance}, \mref{Cor.}{cor:flag-instance}, \mref{Thm.}{thm:instrument-instance}, \mref{Cor.}{cor:register-instance}, \mref{Lem.}{lem:precomposed-instance}, \mref{Def.}{def:emitnorm}, \mref{Thm.}{thm:emit}, \mref{Cor.}{cor:entries}, \mref{Thm.}{thm:emitted-product}, \mref{Cor.}{cor:emitted-identity}, \mref{Rem.}{rem:emitclaim}, \mref{Rem.}{rem:claimscope}, \mref{Rem.}{rem:exactboundary}, \mref{Cor.}{cor:reference-companion}, \mref{Cor.}{cor:companion-factor-enclosure}, \mref{Req.}{req:licensed-entry-lattice}, \mref{Req.}{req:reference-companion}, \mref{Req.}{req:party-label}, \mref{Prop.}{prop:precomposed-blocks}, \mref{Prop.}{prop:flat-instance}, \mref{Cor.}{cor:emitted-flat}, \mref{Cor.}{cor:two-sector-decided}, \mref{Cor.}{cor:flat-companion}, \mref{Def.}{def:domination-link}, \mref{Cor.}{cor:dom-chain}, \extref{../meat-core/contract.pdf}{meat-core/cor:reference-companion-factors}
\metasep newly derived
\metasep implemented (effective: accepted)
\metasep \metalabel{Code}\coderef{https://github.com/fraunhofersg/certiqs-key/blob/v0.1.0/src/certiqs_key/reductions/instance.py}{certiqs_key.reductions.instance:ReductionInstance}}

% claim reductions/cor:identity-instance
\claimsummary{reductions/cor:identity-instance}{A declared block-diagonal measurement, decided exactly admissible on entry, is a reduction instance with the sector pinching as its channel, no scalar sector, no entry and the class-marginal identity discharged. Its chain is the identity link of \hyperlink{claim-reductions/ex:identity}{Example~\ref*{ex:identity}}.}
\claimbookmark{reductions/cor:identity-instance}{Corollary}{cor:identity-instance}{The identity instance}
\claimleadin{%
A receiver that is already finite needs no squasher, only the sector
pinching that every instance's channel carries in front. Its chain is the
identity link of Example~\ref{ex:identity}, and its emission is
Corollary~\ref{cor:emitted-identity}.
}
\begin{corollary}[The identity instance]\label{cor:identity-instance}
Let $\big(\Vset,\{\Vset_t\},\{q_t\}\big)$ be a joint alphabet with classes
and weights, let $\SecSet_0$ be a declared finite ordered sector set with
dimensions $d_s$, and let $\{M_v\}_{v\in\Vset}$ be a POVM on
$\Hs{\Bsys}=\bigoplus_{s\in\SecSet_0}\mathbb{C}^{d_s}$ presented as the
sector-block list $M_v=\bigoplus_sA^{(s)}_v$ of an admissible datum
$A\in\mathcal{D}_{\mathrm{adm}}$ for the sectors $\SecSet_0$ with the default
supports $\Vset^{(s)}=\Vset$, that is, a declared family every block of which
is decided \termref{term-exactly-admissible-on-the-declared-lattice}{exactly admissible on the declared lattice} on entry with
the class sums at $q_t$ (Definition~\ref{def:exactadm}\,(c)). Then the
\emph{identity instance} $\mathcal{J}_{\mathrm{id}}$, given by
\begin{enumerate}[label=(I\arabic*),leftmargin=3.0em]
\item $\Bsys$ as given; $\SecSet_{\mathrm{dat}}=\SecSet_0$; the handed datum
  $A$; (H) empty; and the model map $M_v[A']\coloneqq\bigoplus_{s}A'^{(s)}_v$;
\item $\SecSet_{\mathrm{sc}}=\varnothing$, so that $\Qsys=\Bsys$ and
  $F[A']=M[A']$;
\item $\Lambda=\PinS$ for the decomposition $\SecSet_0$ (that is,
  $\Lambda_0=\idmap$), independent of the datum, with the transported set
  $\{\Pi_s\}_{s\in\SecSet_0}$ and pull-backs $\Pi_s$;
\item an empty entry list, and the class-marginal identity
  discharged,
\end{enumerate}
is a reduction instance. Its chain is the identity link: one receiver-stage
link of kind
\textsf{identity} with $\Lambda_{\Lnk}=\PinS$ (the identity on every state
block diagonal for $\SecSet_0$, and $\idmap$ outright when $\SecSet_0$ has one
sector), with $M^{\mathrm{out}}=M^{\mathrm{in}}$, the sector projectors
carried to themselves, the \termref{term-reachable-set-statement}{reachable-set statement} the declared input
family, pinched, no entries, no licence and
$\textsf{status}=\textsf{exact}$.
\end{corollary}
\claimmeta{\metalabel{Uses}\mref{Def.}{def:system-setting}, \mref{Def.}{def:joint-family}, \mref{Prop.}{prop:pinching}, \mref{Cor.}{cor:blockdiag}, \mref{Def.}{def:link}, \mref{Def.}{def:chain}, \mref{Thm.}{thm:chain}, \mref{Def.}{def:exactadm}, \mref{Def.}{def:reduction-instance}
\metasep \metalabel{Used by}\mref{Cor.}{cor:emitted-identity}, \mref{Cor.}{cor:reference-companion}, \mref{Cor.}{cor:companion-factor-enclosure}
\metasep newly derived
\metasep accepted}

\begin{proof}
(I1) $M[A']$ is a POVM on $\Bsys$ for every $A'\in\mathcal{D}_{\mathrm{adm}}$:
each block is positive semidefinite by admissibility (b), and
$\sum_vM_v[A']=\bigoplus_s\sum_t\sum_{v\in\Vset_t}A'^{(s)}_v
=\bigoplus_s\sum_tq_t\idon{d_s}=\idon{\Bsys}$ by class completeness
\eqref{eq:classcomplete} and $\sum_tq_t=1$. $\Hs{\Bsys}$ is
finite-dimensional, hence separable (Definition~\ref{def:system-setting}).
(H) is
empty, so the stability clause holds vacuously and what follows is
established for every admissible datum.

(I2) The target \eqref{eq:instance-target} with no scalar sector is
$F_v[A']=\bigoplus_sA'^{(s)}_v=M_v[A']$, on $\Qsys=\Bsys$, and the sector map
is $\SecSet_0$ with its dimensions, every sector of data kind.

(I3) Every $F_v[A']$ is block diagonal for $\{\Pi_s\}_{s\in\SecSet_0}$ by
construction, so Proposition~\ref{prop:pinching} gives
$F[A']\dsq_{\PinS}F[A']$, which is $M[A']\dsq_{\PinS}F[A']$; and
$\PinS^{\dagger}(\Pi_s)=\PinS(\Pi_s)=\Pi_s$, so every sector projector is
transported to itself. $\Lambda=\PinS\circ\idmap$ has the form (I3) requires.

(I4) The instance carries no entry of its own: the sector-projector
certificate travels with the transported set, not in the entry list (the
third consequence of Definition~\ref{def:reduction-instance}), and there is
nothing else to certify. The class-marginal identity
\eqref{eq:instance-classmarginal} reads
$\sum_{v\in\Vset_t}M_v[A']=\bigoplus_sq_t\idon{d_s}=q_t\idon{\Bsys}$, which
is class completeness \eqref{eq:classcomplete}; it is therefore discharged
for every admissible datum.

The chain: the link with \termref{term-source-family}{source family} $(\Bsys,M[A'])$, target
$(\Qsys,F[A'])=(\Bsys,M[A'])$, \termref{term-channel-witness}{channel witness} $\PinS$, the carried
pair $(\{\Pi_s\},\{\Pi_s\})$, the \termref{term-reachable-set-statement}{reachable-set statement} ``the
declared input family, pinched'', $\CList=\varnothing$, no licence and
$\textsf{status}=\textsf{exact}$ is a receiver-stage link of
Definition~\ref{def:link} and, with an empty \exttermref{../transmitter/contract.pdf}{term-source-segment}{source segment}, a chain of
Definition~\ref{def:chain}; Theorem~\ref{thm:chain}\,(iv) gives the status,
the identity $\PinS^{\dagger}(M_v[A'])=M_v[A']$ carrying no error term. It is
the identity link of Example~\ref{ex:identity} with $\PinS$ in the channel
slot: $\PinS(\sigma)=\sum_s\Pi_s\sigma\Pi_s=\sigma$ for every $\sigma$ block
diagonal for $\SecSet_0$, so on the block-diagonal states to which
Corollary~\ref{cor:blockdiag}\,(iii) restricts the analysis the two channels
agree, and when $\SecSet_0$ is a single sector $\Pi_s=\idon{\Bsys}$ and
$\PinS=\idmap$.
\end{proof}

% claim reductions/cor:flag-instance
\claimsummary{reductions/cor:flag-instance}{Any POVM on any separable system is a reduction instance with one classical flag per outcome, no data sector and no entry, so there is nothing to normalise. The class-marginal identity is discharged exactly where the branch's declaration supplies it.}
\claimbookmark{reductions/cor:flag-instance}{Corollary}{cor:flag-instance}{The flag instance of an uncharacterised branch}
\claimleadin{%
A branch about which nothing beyond its POVM is declared, such as one an
adversary controls, still squashes, to the flag PVM of
Lemma~\ref{lem:top}; Remark~\ref{rem:flagcost} names the resulting instance
as the one for a branch whose device is unknown or adversarially chosen. The
corollary shows that this squasher is a reduction instance, so that the
branching instance of Theorem~\ref{thm:instrument-instance} can combine such
a branch with characterised ones.
}
\begin{corollary}[The flag instance of an uncharacterised branch]
\label{cor:flag-instance}
Let $\big(\Vset,\{\Vset_t\},\{q_t\}\big)$ be a joint alphabet with classes
and weights and let $\{M_v\}_{v\in\Vset}$ be any POVM on any system $\Bsys$
with $\Hs{\Bsys}$ separable, about which nothing beyond that is declared.
Then the \emph{flag instance} $\mathcal{J}_{\mathrm{flag}}[M]$, given by
\begin{enumerate}[label=(I\arabic*),leftmargin=3.0em]
\item $\Bsys$ as given; no data sector, $\SecSet_{\mathrm{dat}}=\varnothing$,
  so that the datum is empty and $\mathcal{D}_{\mathrm{adm}}=\{\varnothing\}$;
  (H) the single hypothesis that $\{M_v\}$ is a POVM on $\Bsys$; and the
  constant model map $M[\varnothing]\coloneqq M$;
\item one scalar sector $(f)$ of dimension $1$ for each $f\in\Vset$, in the
  order of $\Vset$, with $\phi_f(v)\coloneqq\delta_{vf}$, so that
  $\Hs{\Qsys}=\Hs{\Flag}=\mathbb{C}^{|\Vset|}$ and $F_v=\dyad{v}_{\Flag}$,
  the flag PVM $\Ppvm$ of Lemma~\ref{lem:top};
\item $\Lambda=\PinS\circ\Lmp=\Lmp$, the measure-and-prepare channel
  \eqref{eq:mpsquasher} of $M$, with the transported set
  $\{\Pi_f=\dyad{f}\}_{f\in\Vset}$ and pull-backs $M_f$;
\item an empty entry list, and the class-marginal identity
  discharged exactly when (H) is extended by the declared fact
  $\sum_{v\in\Vset_t}M_v=q_t\idon{\Bsys}$ for every $t$, as it is when the
  branch is declared to be a \termref{term-joint-family}{joint family}
  (Definition~\ref{def:joint-family}\,(b)), and not otherwise,
\end{enumerate}
is a reduction instance. The normalizer of Definition~\ref{def:emitnorm} has
nothing to act on: the declared datum is the handed one, no $\Delta_s$
exists, and the \termref{term-declared-model}{declared model} is $M$ itself.
\end{corollary}
\claimmeta{\metalabel{Uses}\mref{Def.}{def:joint-family}, \mref{Prop.}{prop:pinching}, \mref{Lem.}{lem:top}, \mref{Def.}{def:exactadm}, \mref{Def.}{def:reduction-instance}, \mref{Def.}{def:emitnorm}
\metasep newly derived
\metasep implemented (effective: accepted)
\metasep \metalabel{Code}\coderef{https://github.com/fraunhofersg/certiqs-key/blob/v0.1.0/src/certiqs_key/reductions/sectors.py}{certiqs_key.reductions.sectors:flag_scalar_values}}

\begin{proof}
(I1) There is one datum, the empty family, and it is admissible vacuously;
$M[\varnothing]=M$ is a POVM by (H); the stability clause holds because (H)
reads no datum entry, there being none.

(I2) $\phi_f(v)=\delta_{vf}\in\{0,1\}$ satisfies \eqref{eq:latinv} for every
$\kappa$, is formed with no rounding, and $\sum_v\delta_{vf}=1$;
\eqref{eq:instance-target} gives $F_v=\bigoplus_f\delta_{vf}$, the diagonal
operator $\dyad{v}$ in the basis $\{\ket{f}\}_{f\in\Vset}$ of $\Hs{\Flag}$,
and $\dim\Hs{\Qsys}=|\Vset|$.

(I3) Lemma~\ref{lem:top} gives $M\dsq_{\Lmp}\Ppvm$ with
$\Lmp^{\dagger}(X)=\sum_v\bra{v}X\ket{v}M_v$, so
$\Lmp^{\dagger}(\dyad{f})=M_f$, which is the pull-back of the sector
projector $\Pi_f$; the output of $\Lmp$ is diagonal in the flag basis, so
$\PinS\circ\Lmp=\Lmp$ (Proposition~\ref{prop:pinching}), and $\Lambda$ has
the form (I3) requires with $\Lambda_0=\Lmp$. Since the datum is empty this
holds for every admissible datum.

(I4) The instance carries no entry of its own: the sector-projector
certificate travels with the transported set, not in the entry list (the
third consequence of Definition~\ref{def:reduction-instance}); here that
certificate reads $M_f\succeq0$. The class-marginal identity
\eqref{eq:instance-classmarginal} is a statement about
$\sum_{v\in\Vset_t}M_v$, and (H) alone can supply it, nothing else about $M$
being declared: where the branch is declared to be a joint family it is
\eqref{eq:classmarginal}, and for an arbitrary POVM
$\sum_{v\in\Vset_t}M_v$ is some operator with
$0\preceq\sum_{v\in\Vset_t}M_v\preceq\idon{\Bsys}$ and nothing more. So the
flag is set exactly where the declared fact is part of (H).

The last sentence: Definition~\ref{def:emitnorm} acts on data sectors, of
which there are none, so $A'=A=\varnothing$ and $\Mdec=M[\varnothing]=M$.
\end{proof}

% claim reductions/thm:instrument-instance
\claimsummary{reductions/thm:instrument-instance}{A generalized instrument followed by a reduction instance in each branch is a reduction instance, keeping or forgetting the branch label. A branch entry with constant c becomes a zero-constant entry with c times the branch projector moved into the operator, and the normalizer acts branch by branch.}
\claimbookmark{reductions/thm:instrument-instance}{Theorem}{thm:instrument-instance}{The branching instance: an instrument followed by branch instances}
\claimleadin{%
A receiver described branch by branch, each branch a reduction instance (the
flag instance of Corollary~\ref{cor:flag-instance} for an uncharacterised
branch, for example), is a reduction instance as a whole. The theorem builds
it on Lemma~\ref{lem:master} in two forms: the joint-outcome form keeps the
branch label in the outcome, and the coarse-grained form forgets it. Both
file the branch label as a block index of the target, so they model branches
the adversary may know (Remark~\ref{rem:master-scope}).
Corollary~\ref{cor:register-instance} is the case of a register read by a
POVM.
}
\begin{theorem}[The branching instance: an instrument followed by branch instances]\label{thm:instrument-instance}
Let $J$ be a finite ordered set, let $\{\Ins_j\}_{j\in J}$,
$\Ins_j\colon\Tcl{\Bsys}\to\Tcl{\Bsys'_j}$, be a
\termref{term-generalized-instrument}{generalized instrument} (Definition~\ref{def:generalized-instrument}),
$\sum_j\tr\circ\Ins_j=\tr$, given independently of any datum, and for each
$j$ let $\mathcal{J}_j$ be a reduction instance on $\Bsys'_j$ over one common
joint alphabet with classes and weights $\big(\Vset,\{\Vset_t\},\{q_t\}\big)$,
with handed datum $A_j$, data sectors $\SecSet_{j,\mathrm{dat}}$ with supports
$\Vset^{(j,s)}$, scalar sectors $\SecSet_{j,\mathrm{sc}}$ with values
$\phi_{j,s}$, hypotheses (H$_j$), channel $\Lambda_j[\cdot]$, transported set
$\hat O^{(j)}[\cdot]$ and entry list $\CList_j[\cdot]$. Write
$\SecSet\coloneqq\bigsqcup_{j\in J}\{j\}\times\SecSet_j$, ordered by $j$ and
then as $\SecSet_j$ is, $\Pi_{(j,s)}$ for its sector projectors,
$\Pi_j\coloneqq\sum_{s\in\SecSet_j}\Pi_{(j,s)}$ for the \emph{branch
projectors}, $\hat X|_j$ for the operator equal to $\hat X$ in block $j$ and
zero elsewhere, and $A\coloneqq(A_j)_{j\in J}$ for the joint datum. Then each
of the two constructions (i) and (ii) is a reduction instance, with
\begin{itemize}[leftmargin=2.2em]
\item the channel
  \begin{equation}\label{eq:instrument-channel}
    \Lambda[A']\;\coloneqq\;\PinS\circ\bigoplus_{j\in J}\Lambda_j[A'_j]\circ\Ins_j ,
  \end{equation}
\item the sector set $\SecSet$, the hypotheses (H)~$=\bigcup_j$(H$_j$)
  together with the hypothesis that $\{\Ins_j\}$ is a generalized
  instrument and, for every $j$ in a declared set $J_{\mathrm{w}}\subseteq J$
  of branches with declared lattice constants $q_j$, the
  \emph{branch-weight hypothesis}
  $\Ins_j^{\dagger}(\idon{\Bsys'_j})=q_j\idon{\Bsys}$,
\item the transported set consisting of the sector projectors and, for every
  $j$ and every $\hat X\in\hat O^{(j)}$, the operator $\hat X|_j$, with
  pull-back $\Ins_j^{\dagger}\big(\Lambda_j^{\dagger}(\hat X)\big)$, in
  particular
  \begin{equation}\label{eq:instrument-transport}
    \Lambda^{\dagger}\big(\Pi_{(j,s)}\big)
    =\Ins_j^{\dagger}\big(\Lambda_j^{\dagger}(\Pi_s)\big),
    \qquad
    \Lambda^{\dagger}\big(\Pi_j\big)=\Ins_j^{\dagger}\big(\idon{\Bsys'_j}\big) ,
  \end{equation}
\item the entry list consisting of, for every branch $j$ and every entry
  $(\hat O_k,\bowtie_k,c_k)\in\CList_j$, the entry
  \begin{equation}\label{eq:instrument-entry}
    \big(\hat O_k|_j-c_k\,\Pi_j,\ \bowtie_k,\ 0\big),
  \end{equation}
  together with the \emph{branch-weight entry} $(\Pi_j,=,q_j)$ for every
  $j\in J_{\mathrm{w}}$, where that set is not empty:
\end{itemize}
\begin{enumerate}[label=(\roman*),leftmargin=3.0em]
\item \emph{(the joint-outcome form)} over the joint alphabet
  $\Vset\times J$ of Definition~\ref{def:joint-family}\,(a) whose
  \termref{term-outcome-alphabet}{outcome alphabet} is $\Yset\times J$, ordered with the branch label
  major and the outcome minor, and whose settings are $\Tset$: writing its
  element $(y,j,t)$ as $(v,j)$ with $v=(y,t)\in\Vset$, its
  \exttermref{../foundations/contract.pdf}{term-declared-order}{declared order} has
  $(v,j)<(v',j')$ iff $t<t'$, or $t=t'$ and $j<j'$, or $t=t'$, $j=j'$ and
  $y<y'$, so that every class $\Vset_t\times J$ is a contiguous run and the
  restriction of the order to $\Vset\times\{j\}$ \emph{is} the order of
  $\Vset$; its classes are $\Vset_t\times J$ and its weights $q_t$. On it, the model map
  $M_{(v,j)}[A']\coloneqq\Ins_j^{\dagger}\big(M^{(j)}_v[A'_j]\big)$; data
  sectors $(j,s)$, $s\in\SecSet_{j,\mathrm{dat}}$, of dimension $d_{j,s}$ and
  support $\Vset^{(j,s)}\times\{j\}$, with
  $A'^{(j,s)}_{(v,j')}\coloneqq\delta_{jj'}A'^{(s)}_{j,v}$; scalar sectors
  $(j,s)$, $s\in\SecSet_{j,\mathrm{sc}}$, with
  $\phi_{(j,s)}(v,j')\coloneqq\delta_{jj'}\phi_{j,s}(v)$; and the target
  $F_{(v,j)}[A']=F^{(j)}_v[A'_j]\big|_j$, the branch target in block $j$ and
  zero elsewhere;
\item \emph{(the coarse-grained form)} over $\big(\Vset,\{\Vset_t\},\{q_t\}\big)$
  itself: the model map
  $M_v[A']\coloneqq\sum_{j}\Ins_j^{\dagger}\big(M^{(j)}_v[A'_j]\big)$; the same
  data and scalar sectors with supports $\Vset^{(j,s)}$,
  $A'^{(j,s)}_v\coloneqq A'^{(s)}_{j,v}$ and $\phi_{(j,s)}(v)\coloneqq
  \phi_{j,s}(v)$; and the target $F_v[A']=\bigoplus_jF^{(j)}_v[A'_j]$.
\end{enumerate}
In both forms the class-marginal identity is discharged when every branch
discharges its own, and the flag is set in exactly that case; and the
normalizer of Definition~\ref{def:emitnorm} acts on the composite datum
\emph{branch by branch}: its output at the data sectors of branch $j$ is the
output of the normalizer of $\mathcal{J}_j$ at $A_j$, so that the declared
datum is $A'=(A'_j)_j$ with $A'_j$ the branches' declared data and the
declared model is $M[A']$ at those.
\end{theorem}
\claimmeta{\metalabel{Uses}\mref{Lem.}{lem:sep}, \mref{Cor.}{cor:pullback}, \mref{Lem.}{lem:linear}, \mref{Def.}{def:joint-family}, \mref{Cor.}{cor:blockdiag}, \mref{Def.}{def:generalized-instrument}, \mref{Lem.}{lem:master}, \mref{Thm.}{thm:licence}, \mref{Def.}{def:exactadm}, \mref{Def.}{def:reduction-instance}, \mref{Def.}{def:emitnorm}
\metasep \metalabel{Used by}\mref{Cor.}{cor:register-instance}, \mref{Cor.}{cor:reference-companion}
\metasep newly derived
\metasep accepted}

\begin{proof}
\emph{Idea.} Lemma~\ref{lem:master} combines the branch relations into one;
each part of a reduction instance is then read off branch by branch, and a
branch entry survives the instrument once its constant is moved into the
operator. Throughout, $A'=(A'_j)_j$ denotes a datum of the composite.

\emph{Step 1 (the data).} In the joint form, sector $(j,s)$ has support
$\Vset^{(j,s)}\times\{j\}$, which meets every class $\Vset_t\times J$ because
$\Vset^{(j,s)}$ meets $\Vset_t$; its elements off the support are zero by the
definition of $A'^{(j,s)}$, and its elements on the support are those of
$A'_j$ at sector $s$. Hermiticity and positivity of the elements are those of
$A'_j$; class completeness over $\Vset_t\times J$ reads
$\sum_{(v,j')\in\Vset_t\times J}\delta_{jj'}A'^{(s)}_{j,v}
=\sum_{v\in\Vset_t}A'^{(s)}_{j,v}=q_t\idon{d_{j,s}}$, which is class
completeness of $A'_j$; and the lattice condition is read entry by entry.
In the coarse form the support is $\Vset^{(j,s)}$ and every condition is
literally the branch's. So in both forms a composite datum is an admissible
datum (resp.\ an exactly admissible datum) if and only if every $A'_j$ is:
$\mathcal{D}_{\mathrm{adm}}=\prod_j\mathcal{D}^{(j)}_{\mathrm{adm}}$. The
handed datum $A=(A_j)_j$ is admissible because each $A_j$ is, and its entries
are formed exactly because each branch's are.

\emph{Step 2 ((I1), the model map and stability).} The joint-outcome
alphabet of (i) is a joint alphabet with classes and weights in the sense of
Definition~\ref{def:joint-family}\,(a): its outcome alphabet $\Yset\times J$
is finite, non-empty and ordered, its settings and weights are those of
$\Vset$, and its declared order is that definition's, the outcome varying
fastest and the setting slowest. For every admissible $A'$,
each $M^{(j)}[A'_j]$ is a POVM on $\Bsys'_j$ by (I1) of $\mathcal{J}_j$ under
(H$_j$), which holds by the branch's stability clause. Lemma~\ref{lem:master}\,(i)
with these branch POVMs gives that
$\{\Ins_j^{\dagger}(M^{(j)}_v[A'_j])\}_{(v,j)\in\Vset\times J}$ is a POVM on
$\Bsys$, the joint-form model; coarse-graining $j$ away is
Lemma~\ref{lem:linear}\,(i) for the stochastic matrix that forgets $j$, so
the coarse-form model is a POVM as well. $\Hs{\Bsys}$ is separable by
hypothesis. Stability: (H$_j$) reads no entry of $A'_j$ by the branch's
clause, and the instrument hypothesis and the branch-weight hypotheses read
no datum because the $\Ins_j$ are given independently of any; so (H) reads
no entry of the composite datum,
and the remaining parts, established in Steps 3 to 5 for every admissible $A'$ under
(H), are established as (I1) demands.

\emph{Step 3 ((I2), the target).} The scalar values $\phi_{(j,s)}$ are the
branches' values, \exttermref{../foundations/contract.pdf}{term-lattice-value}{lattice values} in $[0,1]$ formed as (I2) of the branch
requires, with $\sum_{(v,j')}\delta_{jj'}\phi_{j,s}(v)=\sum_v\phi_{j,s}(v)=1$
in the joint form and $\sum_v\phi_{j,s}(v)=1$ in the coarse form. Read
\eqref{eq:instance-target} at the sectors $(j',s)$: in the joint form the
block $(j',s)$ of $F_{(v,j)}[A']$ is $\delta_{jj'}A'^{(s)}_{j,v}$ on a data
sector and $\delta_{jj'}\phi_{j,s}(v)\idon{}$ on a scalar one, so
$F_{(v,j)}[A']$ is the branch target $F^{(j)}_v[A'_j]$ of
\eqref{eq:instance-target} for $\mathcal{J}_j$ in block $j$ and zero
elsewhere, the $F_{v,j}$ of Lemma~\ref{lem:master}\,(ii); in the coarse
form the block $(j,s)$ of $F_v[A']$ is the block $s$ of $F^{(j)}_v[A'_j]$, so
$F_v[A']=\bigoplus_jF^{(j)}_v[A'_j]$, the target of
Lemma~\ref{lem:master}\,(iii). Block diagonality for $\SecSet$ holds by
construction in both forms, and $\dim\Hs{\Qsys}=\sum_j\dim\Hs{\Qsys_j}$ is
finite.

\emph{Step 4 ((I3), the relation).} By (I3) of each branch,
$\big(M^{(j)}[A'_j];O^{(j)}\big)\dsq_{\Lambda_j[A'_j]}
\big(F^{(j)}[A'_j];\hat O^{(j)}\big)$ with $F^{(j)}$ on the
finite-dimensional $\Qsys_j$. Lemma~\ref{lem:master}\,(ii) with these branch
relations gives $\{M_{(v,j)}[A']\}\dsq_{\Lambda_0}\{F_{(v,j)}[A']\}$ for
$\Lambda_0\coloneqq\bigoplus_j\Lambda_j[A'_j]\circ\Ins_j$; its part (v)
transports $\hat X|_j$ to $\Ins_j^{\dagger}(\Lambda_j^{\dagger}(\hat X))$ for
every $\hat X\in\hat O^{(j)}$; its part (iv) gives
$\Lambda_0^{\dagger}(\Pi_j)=\Ins_j^{\dagger}(\idon{\Bsys'_j})$; and
$\Pi_{(j,s)}=\Pi_s|_j$ with $\Pi_s\in\hat O^{(j)}$ by (I3) of the branch, so
$\Lambda_0^{\dagger}(\Pi_{(j,s)})=\Ins_j^{\dagger}(\Lambda_j^{\dagger}(\Pi_s))$,
which is \eqref{eq:instrument-transport}. Its part (iii) gives the
coarse-grained relation with the same channel. Every operator involved is
block diagonal for $\SecSet$ (Step 3; $\hat X|_j$ is block diagonal for
$\SecSet_j$ inside block $j$ by (I3) of the branch), so
Corollary~\ref{cor:blockdiag}\,(i) gives the same relations, transported set
and pull-backs for $\PinS\circ\Lambda_0=\Lambda[A']$, which has the form (I3)
requires; the transported set contains every $\Pi_{(j,s)}$. All of this for
every admissible $A'$ under (H).

\emph{Step 5 ((I4), the entries).} Let $(\hat O_k,\bowtie_k,c_k)\in\CList_j[A'_j]$.
By (I4) of the branch, $\hat O_k$ is a real combination of transported
operators of $\mathcal{J}_j$ of the form \eqref{eq:instance-entry} over
$\SecSet_j$, with coefficients $\alpha_{k,v,s}$ and a projector combination
$Z_k$. Then $\hat O_k|_j$ is the same combination of the operators
$\hat X|_j$, hence transported, and it has the form \eqref{eq:instance-entry}
over $\SecSet$: in the joint form
$\Pi_{(j,s)}F_{(v,j')}[A']\Pi_{(j,s)}=\delta_{jj'}\big(\Pi_sF^{(j)}_v[A'_j]\Pi_s\big)\big|_j$,
so $\hat O_k|_j=\sum_{(v,j')}\sum_{s}\alpha_{k,(v,j'),(j,s)}
\Pi_{(j,s)}F_{(v,j')}\Pi_{(j,s)}+Z_k|_j$ with
$\alpha_{k,(v,j'),(j,s)}\coloneqq\delta_{jj'}\alpha_{k,v,s}$; in the coarse
form $\hat O_k|_j=\sum_v\sum_s\alpha_{k,v,s}\Pi_{(j,s)}F_v\Pi_{(j,s)}+Z_k|_j$
directly. $Z_k|_j$ and $c_k\Pi_j$ are real combinations of sector
projectors of $\SecSet$ with coefficients on the lattice: a coefficient of
$Z_k|_j-c_k\Pi_j$ is $\beta_s-c_k$, a difference of two integer multiples of
$2^{-\kappa}$ and hence one itself, which is all that is asked of it, since
Definition~\ref{def:reduction-instance}\,(I4) puts the band of
\eqref{eq:latinv} on an \emph{entry} and not on a coefficient and this step
uses nothing about the coefficient beyond its being on the lattice. So
$\hat O_k|_j-c_k\Pi_j$ has the form \eqref{eq:instance-entry} with
$Z=Z_k|_j-c_k\Pi_j$, and it is a real combination of transported operators.
Its pull-back, by \eqref{eq:instrument-transport} and linearity, is
\begin{equation*}
  \Lambda^{\dagger}\big(\hat O_k|_j-c_k\Pi_j\big)
  =\Ins_j^{\dagger}\big(\Lambda_j^{\dagger}(\hat O_k)\big)
   -c_k\,\Ins_j^{\dagger}\big(\idon{\Bsys'_j}\big)
  =\Ins_j^{\dagger}\Big(\Lambda_j^{\dagger}(\hat O_k)-c_k\,\idon{\Bsys'_j}\Big).
\end{equation*}
The branch certificate says
$\Lambda_j^{\dagger}(\hat O_k)-c_k\idon{\Bsys'_j}\bowtie_k0$, and
$\Ins_j^{\dagger}$ is a positive map, being the adjoint of a completely
positive one (Lemma~\ref{lem:sep}\,(i) identifies it; positivity is
\eqref{eq:pairing} tested on positive operators): the image of an operator
$\succeq0$ is $\succeq0$, of one $\preceq0$ is $\preceq0$, and of $0$ is $0$.
So the pull-back $\bowtie_k0$, which is \eqref{eq:certificate} for
\eqref{eq:instrument-entry} at the constant $0$, for every admissible $A'$
under (H). The constant moves into the operator because positivity of
$\Ins_j^{\dagger}$ bounds $\Ins_j^{\dagger}(\Lambda_j^{\dagger}(\hat O_k))$
against $c_k\Ins_j^{\dagger}(\idon{})$, the branch weight, and not against
$c_k\idon{\Bsys}$; the entry \eqref{eq:instrument-entry} is the form in
which the branch certificate survives without a hypothesis on that weight.
For $j\in J_{\mathrm{w}}$ the branch-weight hypothesis of (H) makes the
weight a constant, $\Ins_j^{\dagger}(\idon{\Bsys'_j})=q_j\idon{\Bsys}$, and
the branch-weight entry is listed only there; Lemma~\ref{lem:master}\,(iv)
gives
$\Lambda^{\dagger}(\Pi_j)=q_j\idon{\Bsys}$, an equality certificate at
$c=q_j$ for $\Pi_j$, which has the form \eqref{eq:instance-entry} with every
$\alpha=0$; and then also $\Lambda^{\dagger}(\hat O_k|_j)\bowtie_kc_kq_j\idon{\Bsys}$,
the branch entry at the scaled constant, which is
\eqref{eq:instrument-entry} plus $c_k$ times the branch-weight entry.
The class-marginal identity: in the joint form,
$\sum_{(v,j')\in\Vset_t\times J}M_{(v,j')}[A']
=\sum_{j'}\Ins_{j'}^{\dagger}\big(\sum_{v\in\Vset_t}M^{(j')}_v[A'_{j'}]\big)$,
and in the coarse form $\sum_{v\in\Vset_t}M_v[A']$ is the same sum; where
every branch discharges \eqref{eq:instance-classmarginal} this is
$q_t\sum_{j'}\Ins_{j'}^{\dagger}(\idon{\Bsys'_{j'}})=q_t\idon{\Bsys}$ by the
instrument identity $\sum_j\Ins_j^{\dagger}(\idon{\Bsys'_j})=\idon{\Bsys}$
established in the proof of Lemma~\ref{lem:master}, for every admissible
$A'$ under (H). Where some branch does not discharge it, nothing in this
construction establishes the identity, and the flag is not set.

\emph{Step 6 (the normalizer acts per branch).} Definition~\ref{def:emitnorm}
acts per data sector and per class on the elements of the class within the
sector's support. At sector $(j,s)$ of the joint form the support is
$\Vset^{(j,s)}\times\{j\}$ and the class is $\Vset_t\times J$, so the elements
acted on are $\{A^{(s)}_{j,v}:v\in\Vset_t\cap\Vset^{(j,s)}\}$, exactly the
elements the normalizer of $\mathcal{J}_j$ acts on at sector $s$ and class
$t$, with the same residue $q_t\idon{}-\sum_{v\in\Vset_t}A^{(s)}_{j,v}$,
the same exact traces, the same \exttermref{../foundations/contract.pdf}{term-declared-order}{declared order} (the order of
$\Vset\times J$ declared in (i) restricts on $\Vset\times\{j\}$ to the order
of $\Vset$, the setting varying slowest in both), hence the same
designation order, the same roundings to nearest and mixing weights of the fidelity snap
(Definition~\ref{def:emitnorm}\,(2)), the same components and hosts, and the
same last resort (Definition~\ref{def:emitnorm}\,(3)), which read the
elements of the class, their number, the class weight, the designation order
and the \exttermref{../foundations/contract.pdf}{term-structural-support}{structural supports} alone, all the same
in the two runs under one run configuration and its strict option; every
step of the two runs therefore coincides, and so do their outputs, their
displacements and their moved classes. The
elements \termref{term-outside-the-sector-support}{outside the sector support} are touched by neither run. In
the coarse form the support is $\Vset^{(j,s)}$ and the class is $\Vset_t$,
and the identification is immediate. Hence $A'=(A'_j)_j$, and the declared
model is $M[A']$ at the branches' declared data by \eqref{eq:declmodel}.
\end{proof}

% claim reductions/cor:register-instance
\claimsummary{reductions/cor:register-instance}{When a register read by a POVM selects the branch, the branching instance is a reduction instance whose pull-backs are the branches' tensored with the register element. Branch entries hold at constant zero, and a branch-weight entry only where (H) makes that element a multiple of the identity.}
\claimbookmark{reductions/cor:register-instance}{Corollary}{cor:register-instance}{The register instance: register-conditioned squashing of instances}
\claimleadin{%
A register $D$, read by a POVM that need not be projective
(Remark~\ref{rem:povmreg}), selects which branch instance acts; a branch may
be the flag instance of Corollary~\ref{cor:flag-instance}. The corollary is
the branching instance at the instrument that reads the register.
}
\begin{corollary}[The register instance: register-conditioned squashing of instances]\label{cor:register-instance}
Let $\Bsys=\Bsys_0D$, let $\{N_j\}_{j\in J}$ be a POVM on the register $D$
with $J$ finite and ordered, and for each $j$ let $\mathcal{J}_j$ be a
reduction instance on $\Bsys_0$ over one common joint alphabet with classes
and weights, with handed model $\{M_{v|j}\}_{v\in\Vset}$, channel
$\Lambda_j$, transported set $\hat O^{(j)}$ and entry list $\CList_j$. Let
$\Ins_j$ be the instrument of Theorem~\ref{thm:mixture},
$\Ins_j(\rho)=\tr_{D}\big[(\idon{\Bsys_0}\otimes N_j^{1/2})\rho
(\idon{\Bsys_0}\otimes N_j^{1/2})\big]$. Then the branching instance of
Theorem~\ref{thm:instrument-instance} for $\{\Ins_j\}$ and
$\{\mathcal{J}_j\}$, the \emph{register instance}, is a reduction
instance on $\Bsys$ in both forms, with the handed model
\begin{equation}\label{eq:register-model}
  M_v=\sum_{j\in J}M_{v|j}\otimes N_j
  \quad\text{(coarse-grained)},
  \qquad
  M_{(v,j)}=M_{v|j}\otimes N_j
  \quad\text{(joint-outcome)},
\end{equation}
the channel \eqref{eq:regchannel} pinched,
$\Lambda=\PinS\circ\bigoplus_j\Lambda_j\circ\Ins_j$, the target
$\bigoplus_jF^{(j)}_v$ resp.\ $F^{(j)}_v|_j$, the hypotheses
$\bigcup_j$(H$_j$) together with ``$\{N_j\}$ is a POVM on $D$'' and, for every
$j$ in a declared set $J_{\mathrm{w}}\subseteq J$ with declared lattice
constants $q_j$, ``$N_j=q_j\idon{D}$'', and:
\begin{enumerate}[label=(\roman*),leftmargin=3.0em]
\item \emph{(transport)} for every branch $j$ and sector $s$ of
  $\mathcal{J}_j$, and for every $\hat X\in\hat O^{(j)}$,
  \begin{equation}\label{eq:register-transport}
    \Lambda^{\dagger}\big(\Pi_{(j,s)}\big)=\Lambda_j^{\dagger}(\Pi_s)\otimes N_j,
    \qquad
    \Lambda^{\dagger}\big(\Pi_j\big)=\idon{\Bsys_0}\otimes N_j,
    \qquad
    \Lambda^{\dagger}\big(\hat X|_j\big)=\Lambda_j^{\dagger}(\hat X)\otimes N_j ;
  \end{equation}
\item \emph{(the entry clause)} every entry $(\hat O,\bowtie,c)$ of a branch
  $j$ is carried as $\big(\hat O|_j-c\,\Pi_j,\bowtie,0\big)$, with pull-back
  \begin{equation}\label{eq:register-entry}
    \Lambda^{\dagger}\big(\hat O|_j-c\,\Pi_j\big)
    =\big(\Lambda_j^{\dagger}(\hat O)-c\,\idon{\Bsys_0}\big)\otimes N_j ,
  \end{equation}
  a tensor product of the branch's certified operator with $N_j\succeq0$, so
  that the certificate holds at $0$ with the branch's relation; the
  branch-weight entry $(\Pi_j,=,q_j)$ is listed for every $j\in J_{\mathrm{w}}$;
\item \emph{(the class-marginal identity)} it is discharged when every
  branch discharges it, and the flag is set in exactly that case;
\item \emph{(the declared model)} the normalizer acts branch by branch, and
  the declared model is
  \begin{equation}\label{eq:register-declared}
    \Mdec_v=\sum_{j\in J}{\Mdec}^{(j)}_v\otimes N_j
    \quad\text{resp.}\quad
    \Mdec_{(v,j)}={\Mdec}^{(j)}_v\otimes N_j ,
  \end{equation}
  with ${\Mdec}^{(j)}$ the declared model of branch $j$.
\end{enumerate}
\end{corollary}
\claimmeta{\metalabel{Uses}\mref{Def.}{def:generalized-instrument}, \mref{Lem.}{lem:master}, \mref{Thm.}{thm:mixture}, \mref{Def.}{def:reduction-instance}, \mref{Thm.}{thm:instrument-instance}, \mref{Def.}{def:emitnorm}
\metasep \metalabel{Used by}\mref{Cor.}{cor:reference-companion}
\metasep newly derived
\metasep implemented (effective: accepted)
\metasep \metalabel{Code}\coderef{https://github.com/fraunhofersg/certiqs-key/blob/v0.1.0/src/certiqs_key/reductions/instance.py}{certiqs_key.reductions.instance:combine_branches}}

\begin{proof}
The proof of Theorem~\ref{thm:mixture} shows that $\{\Ins_j\}$ is a
\termref{term-generalized-instrument}{generalized instrument} (Definition~\ref{def:generalized-instrument})
with $\Ins_j^{\dagger}(X)=X\otimes N_j$ for $X\in\Lin{\Bsys_0}$, and $\Ins_j$
is given by $N_j$ alone, independently of any datum; the hypothesis that it
is an instrument is the hypothesis that $\{N_j\}$ is a POVM on $D$, the
branch-weight hypothesis of the theorem at $j\in J_{\mathrm{w}}$ is
$N_j=q_j\idon{D}$ (below), and that
$\Hs{\Bsys}$ is separable is inherited from $\Bsys_0$ and $D$.
Theorem~\ref{thm:instrument-instance} therefore applies. Its model maps are
$M_{(v,j)}[A']=\Ins_j^{\dagger}(M^{(j)}_v[A'_j])=M^{(j)}_v[A'_j]\otimes N_j$
and their sum over $j$, which at the handed data are the two models
\eqref{eq:register-model}; its channel \eqref{eq:instrument-channel} is
\eqref{eq:regchannel} composed with the sector pinching, as in
Theorem~\ref{thm:mixture}\,(i) and (iv). (i) is
\eqref{eq:instrument-transport} and the transport clause of the theorem with
$\Ins_j^{\dagger}(X)=X\otimes N_j$ substituted, the second identity being
Theorem~\ref{thm:mixture}\,(iii). (ii) is \eqref{eq:instrument-entry} with
the same substitution:
$\Ins_j^{\dagger}\big(\Lambda_j^{\dagger}(\hat O)-c\idon{\Bsys_0}\big)
=\big(\Lambda_j^{\dagger}(\hat O)-c\idon{\Bsys_0}\big)\otimes N_j$; if
$Y\succeq0$ on $\Bsys_0$ and $N_j\succeq0$ then $Y\otimes N_j\succeq0$, if
$Y\preceq0$ then $Y\otimes N_j\preceq0$, and $Y=0$ gives $0$, so the branch's
relation survives at the constant $0$. The branch weight
$\Ins_j^{\dagger}(\idon{\Bsys_0})=\idon{\Bsys_0}\otimes N_j$ is a multiple
$q_j\idon{\Bsys}$ exactly when $N_j=q_j\idon{D}$, so the theorem's
branch-weight hypothesis at $j\in J_{\mathrm{w}}$ is the hypothesis
$N_j=q_j\idon{D}$ stated here, and the entry is the theorem's (and that is
the reason Theorem~\ref{thm:mixture}\,(iii) calls the general case an
identity rather than a constraint). (iii) and (iv)
are the theorem's two closing sentences at this instrument, the declared
model \eqref{eq:register-declared} being the model map at
$A'=(A'_j)_j$ by \eqref{eq:declmodel}.
\end{proof}

% claim reductions/lem:precomposed-instance
\claimsummary{reductions/lem:precomposed-instance}{A squashing relation precomposed with a channel on the physical system is again one, with the same target and every spectral certificate at the same constant. A reduction instance so precomposed is again one, with the same datum, target, entries and emission.}
\claimbookmark{reductions/lem:precomposed-instance}{Lemma}{lem:precomposed-instance}{Precomposition of a reduction instance with a channel}
\claimleadin{%
Proposition~\ref{prop:preorder} composes two squashing relations whose
middle system is a finite target. The lemma below precomposes with a channel
between two physical systems instead, possibly infinite-dimensional ones, as
when a device is a known channel followed by a model a reduction instance is
built on; the certificates it keeps are those of the form
Theorem~\ref{thm:licence} reads.
}
\begin{lemma}[Precomposition of a reduction instance with a channel]\label{lem:precomposed-instance}
Let $\Bsys$ and $\Bsys'$ be systems with separable Hilbert spaces, and let
$\Phi\colon\Tcl{\Bsys}\to\Tcl{\Bsys'}$ be a channel with adjoint
$\Phi^{\dagger}\colon\Lin{\Bsys'}\to\Lin{\Bsys}$ (Lemma~\ref{lem:sep}).
\begin{enumerate}[label=(\roman*),leftmargin=3.0em]
\item \emph{(relations)} If $(M';O')\dsq_{\Lambda}(F;\hat O)$ with
  $\Lambda\colon\Tcl{\Bsys'}\to\Tcl{\Qsys}$, then
  $\Phi^{\dagger}(M')\coloneqq\{\Phi^{\dagger}(M'_v)\}_{v\in\Vset}$ is a POVM
  on $\Bsys$ and
  \begin{equation}\label{eq:precomposed-relation}
    \big(\Phi^{\dagger}(M');\Phi^{\dagger}(O')\big)\;\dsq_{\Lambda\circ\Phi}\;
    \big(F;\hat O\big),\qquad
    (\Lambda\circ\Phi)^{\dagger}=\Phi^{\dagger}\circ\Lambda^{\dagger}.
  \end{equation}
  If $\Lambda^{\dagger}(\hat O_k)\bowtie c\,\idon{\Bsys'}$ for a relation
  $\bowtie\,\in\{\preceq,=,\succeq\}$ and a real $c$, then
  $(\Lambda\circ\Phi)^{\dagger}(\hat O_k)\bowtie c\,\idon{\Bsys}$, with the
  same relation and the same constant.
\item \emph{(instances)} Let $\mathcal{J}$ be a \termref{term-reduction-instance}{reduction instance} on
  $\Bsys'$ with model map $M[\cdot]$, channel
  $\Lambda[\cdot]=\PinS\circ\Lambda_0[\cdot]$ and pull-backs $O[\cdot]$.
  The \emph{precomposed instance} $\mathcal{J}\circ\Phi$ has the system
  $\Bsys$; the joint alphabet, data sectors, supports,
  \termref{term-handed-datum}{handed datum}, \termref{term-hypotheses-h}{hypotheses (H)}, \termref{term-scalar-value}{scalar values},
  \termref{term-sector-set}{sector set}, \termref{term-target-family}{target family}, \termref{term-transported-set}{transported set}, entry list and
  \termref{term-class-marginal-flag}{class-marginal flag} of $\mathcal{J}$; the model map
  $A'\mapsto\Phi^{\dagger}(M[A'])$; the channel
  $A'\mapsto\Lambda[A']\circ\Phi=\PinS\circ(\Lambda_0[A']\circ\Phi)$; and the
  pull-backs $A'\mapsto\Phi^{\dagger}(O[A'])$. It is a reduction instance
  on $\Bsys$.
\item \emph{(at a datum, and the emission)} For every admissible datum $A'$,
  $\mathcal{J}\circ\Phi$ at $A'$ is $\mathcal{J}$ at $A'$ precomposed with
  $\Phi$. The \termref{term-emission-normalizer}{emission normalizer} produces the same declared datum from
  the handed datum of either, so the emitted target, scalar values, entries
  and constants of $\mathcal{J}\circ\Phi$ are those of $\mathcal{J}$; its
  emitted channel is the emitted channel of $\mathcal{J}$ precomposed with
  $\Phi$, and its declared model is $\Phi^{\dagger}$ of the declared model of
  $\mathcal{J}$.
\end{enumerate}
\end{lemma}
\claimmeta{\metalabel{Uses}\mref{Lem.}{lem:sep}, \mref{Def.}{def:squashing}, \mref{Prop.}{prop:pinching}, \mref{Def.}{def:reduction-instance}, \mref{Def.}{def:emitnorm}
\metasep \metalabel{Used by}\mref{Cor.}{cor:normalized-flat}, \mref{Req.}{req:normalized-link-record}
\metasep newly derived
\metasep implemented (effective: accepted)
\metasep \metalabel{Code}\coderef{https://github.com/fraunhofersg/certiqs-key/blob/v0.1.0/src/certiqs_key/reductions/instance.py}{certiqs_key.reductions.instance:precompose}}

\begin{proof}
(i) $\Phi$ is completely positive and trace preserving, so by
Lemma~\ref{lem:sep} (its trace-non-increasing case, with $\Bsys'$ in place of
the target) its adjoint is completely positive, and it is unital: for every
$\rho\in\Tcl{\Bsys}$,
$\tr[\Phi^{\dagger}(\idon{\Bsys'})\rho]=\tr\Phi(\rho)=\tr\rho$, so
$\Phi^{\dagger}(\idon{\Bsys'})=\idon{\Bsys}$ by Lemma~\ref{lem:sep}\,(i).
Hence every $\Phi^{\dagger}(M'_v)$ is positive semidefinite and
$\sum_v\Phi^{\dagger}(M'_v)=\Phi^{\dagger}(\idon{\Bsys'})=\idon{\Bsys}$. The
composite $\Lambda\circ\Phi\colon\Tcl{\Bsys}\to\Tcl{\Qsys}$ is a channel, and
for $X\in\Lin{\Qsys}$ and $\rho\in\Tcl{\Bsys}$ the pairings
\eqref{eq:pairing} of $\Lambda$ and of $\Phi$, the second applied to the
bounded operator $\Lambda^{\dagger}(X)$, give
$\tr[X\,\Lambda(\Phi(\rho))]=\tr[\Lambda^{\dagger}(X)\,\Phi(\rho)]
=\tr[\Phi^{\dagger}(\Lambda^{\dagger}(X))\,\rho]$; the adjoint fixed by
\eqref{eq:pairing} is unique (Lemma~\ref{lem:sep}), so
$(\Lambda\circ\Phi)^{\dagger}=\Phi^{\dagger}\circ\Lambda^{\dagger}$. Then
$(\Lambda\circ\Phi)^{\dagger}(F_v)=\Phi^{\dagger}(M'_v)$ and
$(\Lambda\circ\Phi)^{\dagger}(\hat O_k)=\Phi^{\dagger}(O'_k)$, which is
\eqref{eq:precomposed-relation} by Definition~\ref{def:squashing}. For the
certificate with $\succeq$: $\Lambda^{\dagger}(\hat O_k)-c\,\idon{\Bsys'}
\succeq0$, and $\Phi^{\dagger}$ is positive and unital, so
$(\Lambda\circ\Phi)^{\dagger}(\hat O_k)-c\,\idon{\Bsys}
=\Phi^{\dagger}\big(\Lambda^{\dagger}(\hat O_k)-c\,\idon{\Bsys'}\big)\succeq0$;
for $\preceq$ apply this to $-\hat O_k$ and $-c$, and for $=$ use both.

(ii) We check the four parts of Definition~\ref{def:reduction-instance}.
(I1): $\Hs{\Bsys}$ is separable; the handed datum and its formation are those
of $\mathcal{J}$; the hypotheses (H) are those of $\mathcal{J}$ and read no
entry of the datum; and for every admissible $A'$ the family
$\Phi^{\dagger}(M[A'])$ is a POVM on $\Bsys$ by (i), so the stability clause
holds wherever it holds for $\mathcal{J}$. (I2): the sector set, the scalar
values and the target family are functions of the datum and the scalar
values alone, which are unchanged. (I3): $\Lambda_0[A']\circ\Phi$ is a channel
from $\Tcl{\Bsys}$, given by the constructor of $\mathcal{J}$ preceded by
$\Phi$, so $\Lambda[A']\circ\Phi$ has the required form with the same sector
pinching of Proposition~\ref{prop:pinching}; the transported set is unchanged
and contains every sector projector, and \eqref{eq:instance-relation} for
$\mathcal{J}\circ\Phi$ at every admissible $A'$ under (H) is (i) applied to
that relation for $\mathcal{J}$. (I4): every entry is the same combination
\eqref{eq:instance-entry} of the same target, so it is an entry of the
required form, and its certificate is discharged at every admissible datum by
the certificate clause of (i). Where $\mathcal{J}$ discharges the
class-marginal identity \eqref{eq:instance-classmarginal},
$\sum_{v\in\Vset_t}\Phi^{\dagger}(M_v[A'])=\Phi^{\dagger}(q_t\idon{\Bsys'})
=q_t\idon{\Bsys}$, so $\mathcal{J}\circ\Phi$ discharges it too, and the flag
is the same.

(iii) The first sentence is the definition in (ii) read at $A'$. The emission
normalizer of Definition~\ref{def:emitnorm} acts on the handed datum, class by
class and sector by sector, with the class weights and the lattice exponent,
and reads nothing else of the instance; the handed datum, the class weights
and the lattice exponent of $\mathcal{J}\circ\Phi$ are those of $\mathcal{J}$,
so it produces the same declared datum $A'$. The emitted target, scalar values
and entries are functions of $A'$, the scalar values and the entry rule, all
unchanged, and the emitted channel and declared model are those of
$\mathcal{J}\circ\Phi$ at $A'$, namely $\Lambda[A']\circ\Phi$ and
$\Phi^{\dagger}(M[A'])$.
\end{proof}

\section{Exact emission}\label{group-reductions/def:emitnorm}
\sectionintro{Where the exact construction of a \termref{term-reduction-instance}{reduction instance} leaves the declared lattice, the \termref{term-emission-normalizer}{emission normalizer} declares a nearby datum, keeping the classes already on the lattice and snapping every other class for fidelity without reviving a zero row or column, or, where no element can host its residue, by a flagged last resort, and decides all six emission conditions last (\hyperlink{claim-reductions/def:emitnorm}{Definition~\ref*{def:emitnorm}}). The main result, \hyperlink{claim-reductions/thm:emit}{Theorem~\ref*{thm:emit}}, is that the emitted family is exactly admissible and squashes the \termref{term-declared-model}{declared model} exactly, and that every claim of the instance holds of that model; where the normalizer moved the datum, nothing proved relates the two models, their distance least of all, and a run either evaluates off the snap or rests on a disclosed premise. Which constraint entries are emitted and at which constants (\hyperlink{claim-reductions/cor:entries}{Corollary~\ref*{cor:entries}}), and how several \termref{term-measuring-parties}{measuring parties}' emissions combine without a product datum (\hyperlink{claim-reductions/thm:emitted-product}{Theorem~\ref*{thm:emitted-product}}), follow; the remarks say what an emitted claim is about, where exact decisions stop, and why the emitted constraints are sound by theorem rather than by assumption.}
\sectionimports{\extref{../foundations/contract.pdf}{foundations/def:exact-reference}, \extref{../foundations/contract.pdf}{foundations/def:snap}}

% claim reductions/def:emitnorm
\claimsummary{reductions/def:emitnorm}{The normalizer puts an exact datum back on the lattice class by class: it keeps the classes already there and snaps the others for fidelity, rounding to nearest and placing the residue only where support allows, or, failing that, by a flagged last resort.}
\claimbookmark{reductions/def:emitnorm}{Definition}{def:emitnorm}{The emission normalizer, the declared datum and the declared model}
\claimleadin{%
A \termref{term-reduction-instance}{reduction instance} forms its \termref{term-handed-datum}{handed datum} by exact
arithmetic; every exact step multiplies, and a product of dyadics adds
significand widths, so the datum can leave the declared lattice on which this
module emits. The normalizer defined here puts it back by a snap of the
module's own emission, never of anything it was handed. Every class of the
handed datum is Hermitian, positive semidefinite and complete at its class
weight with Gaussian-rational entries, so it is an \exttermref{../foundations/contract.pdf}{term-exact-reference}{exact reference} of
itself (\xref{foundations/def:exact-reference}, every radius zero). The
normalizer keeps the classes already on the lattice and snaps every other
class by one procedure that stays as close to the handed datum as it can
while keeping every row and column that the handed datum holds at zero
(\xref{foundations/req:support-safe-snap}): it rounds to nearest, softens the
off-diagonal entries where positivity needs it, and places what completeness
leaves only on elements whose support holds it. Where no such placement
exists, a last resort mixes the class with a multiple of the identity
(\xref{foundations/req:support-safe-snap}\,(iv)); that may give an element a
row the handed datum denies it, so the class is flagged. No claim relates a
snapped class to the handed datum. A run that forms the number its key
length is formed from on the reference companion
(Corollary~\ref{cor:reference-companion}), which the reductions build with no
normalizer from the \exttermref{../foundations/contract.pdf}{term-handed-reference}{handed reference}, or from the instance's own
snapped blocks where none is handed, evaluates every class off its snap; a run that does not, such as one whose
key length is formed from a number certified on the snapped declaration in
the fallback of \xref{meat-core/frank-wolfe/req:certified-evaluation}\,(i),
rests, on the classes the
normalizer moved, on \xref{foundations/ass:snapped-measurement}, and a run
that declares the strict option refuses there instead
(\xref{foundations/req:snap-reference}\,(iv)). The emission theorem
(Theorem~\ref{thm:emit}) is stated about what the normalizer produces. The
six conditions its last step decides are those of
\xref{foundations/req:exact-emission}; when a producer re-establishes its
lattice at an emission, and on what it may snap at all, is
\xref{foundations/req:who-snaps}. Everything the normalizer records is
Requirement~\ref{req:emission-record}, and its refusals are
Requirement~\ref{req:emission-refusal}.
}
\begin{definition}[The emission normalizer, the declared datum and the declared model]
\label{def:emitnorm}
Let $\mathcal{J}$ be a \termref{term-reduction-instance}{reduction instance}
(Definition~\ref{def:reduction-instance}) with \termref{term-handed-datum}{handed datum}
$A\in\mathcal{D}_{\mathrm{adm}}$, formed exactly. Those numbers \emph{are} the
\termref{term-data-sector}{data sectors} of the emitted target
(Definition~\ref{def:reduction-instance}\,(I2)), so they are this module's own
declaration to its \exttermref{../foundations/contract.pdf}{term-consumer}{consumer}. Where the exact arithmetic that formed
them has left the declared lattice \eqref{eq:latinv}, they
are re-established on it, per data sector $s$ and per class $t$, by the
following steps, the
\termdef{term-emission-normalizer}{emission normalizer}.
The steps act on the elements $A^{(s)}_v$ with
$v\in\Vset_t\cap\Vset^{(s)}$, the elements of the class $t$ in the sector's
support; the elements \termref{term-outside-the-sector-support}{outside the sector support} are exactly zero,
lie on the lattice, and are touched by no step. Every positive
semidefiniteness below is decided exactly by \eqref{eq:psdtest}, on matrices
with entries in $\Fldi$. The
\exttermref{../foundations/contract.pdf}{term-structural-support}{structural support} of a matrix is the set of indices $i$ whose row $i$
or column $i$ holds a non-zero entry (\xref{foundations/def:snap}). The
\emph{designation order} ranks the elements of a class in the sector's
support by largest exact trace and, among equal traces, by the
\exttermref{../foundations/contract.pdf}{term-declared-order}{declared order} of the class. \emph{A tie is exact and ordinary here}:
the traces compared are exact rationals, and a sector symmetric under a
relabelling of the class ties every element of it, so ``largest trace''
alone names no element; the declared order is the order in which this
contract reads $\Vset$, the outcome varying fastest, and it is emitted with
the artefact, so the order is a function of the declaration and of nothing
else. \emph{Rounding to nearest} a matrix at $\kappa$ rounds every real and
imaginary part of every entry to the nearest integer multiple of
$2^{-\kappa}$, \emph{ties to even}: a part lying exactly halfway between two
multiples goes to the one whose integer numerator is even. \emph{The tie is
exact}: the rounding is carried out in the substrate of
Lemma~\ref{lem:exactpsd} on exact rationals, so a halfway part is an exact
equality rather than a floating-point near-miss. \emph{Ties to even is
odd-symmetric}: negating a part negates both of the multiples nearest it and
their numerators, which preserves evenness, so the rounding commutes with
$X\mapsto X^{\dagger}$ --- mirror entries have equal real parts and opposite
imaginary parts before it and therefore after it --- and a matrix rounded
from a Hermitian one is Hermitian; a zero entry stays zero.
\begin{enumerate}[label=(\arabic*),leftmargin=2.4em]
\item \emph{Keep.} A class $t$ at a sector $s$ every element of which
  satisfies \eqref{eq:latinv} entry by entry is \emph{kept}:
  $A'^{(s)}_v\coloneqq A^{(s)}_v$ for every $v\in\Vset_t\cap\Vset^{(s)}$, and
  no further step acts on it.
\item \emph{Snap for fidelity.} At every other class and sector, for $w=0$
  and then $w=2^{-\kappa},2^{-\kappa+1},\dots,2^{-1}$ in this order, let
  $P^{w}_v$ be $A^{(s)}_v$ with every off-diagonal entry multiplied by
  $1-w$, that is $(1-w)A^{(s)}_v+w\,T_v$ with $T_v$ the diagonal part of
  $A^{(s)}_v$. Every element of the class in the sector's support receives
  $P^{w}_v$ rounded to nearest at $\kappa$, and the \emph{class residue}
  $E^{w}\coloneqq q_t\idon{d_s}-\sum_{v\in\Vset_t\cap\Vset^{(s)}}A'^{(s)}_v$,
  a Hermitian matrix of \exttermref{../foundations/contract.pdf}{term-lattice-value}{lattice values}, is split by its \emph{components}:
  the classes of the equivalence relation on the indices $i$ whose row $i$ of
  $E^{w}$ is non-zero generated by $i\sim l$ whenever $E^{w}_{il}\ne0$. For a
  component $C$ write $E_C$ for the matrix that agrees with $E^{w}$ on the
  entries whose row and column both lie in $C$ and is zero elsewhere. Each
  $E_C$ is added to the first element $A'^{(s)}_u$ in the designation order
  whose structural support in the handed datum, that of $A^{(s)}_u$,
  contains $C$, its \emph{host}. The first $w$ at which every component has
  a host and every element so formed is decided positive semidefinite gives
  the family, and $w$ is recorded. Where no $w$ does, the class is snapped
  by step~(3).
\item \emph{The last resort.} The support-safe steps of (2) produced no snap
  of the class, and the class is snapped by a uniform lift, applied to the
  class with constant $q_t$ and starting from the handed datum. Let $n$ be
  the number of elements of the class in the sector's support, $v_1$ the
  first of them in the designation order, and, for $w\in(0,1)$,
  \begin{equation}\label{eq:lastresort}
    B^{w}_v\;\coloneqq\;(1-w)\,A^{(s)}_v+w\,\frac{q_t}{n}\,\idon{d_s},
    \qquad v\in\Vset_t\cap\Vset^{(s)},
  \end{equation}
  a family that is Hermitian, complete at $q_t$, with entries in $\Fldi$, and
  positive definite in every element, since $q_t>0$ at a class that is not
  kept (at $q_t=0$ completeness and positive semidefiniteness make every
  element zero, which is on the lattice). Its designation order is that of
  the class, $\tr B^{w}_v$ being an increasing affine function of
  $\tr A^{(s)}_v$. For $w=2^{-\kappa},2^{-\kappa+1},\dots,2^{-1}$ in this
  order, every element other than $v_1$ receives $B^{w}_v$ rounded to
  nearest at $\kappa$, and $v_1$ receives $q_t\idon{d_s}$ minus their sum;
  the first $w$ at which every element so formed is decided positive
  semidefinite gives the family. Where no $w$ does, every element $v\ne v_1$
  receives the first matrix $A'^{(s)}_v$ of the following list that is
  Hermitian, satisfies \eqref{eq:latinv} entry by entry, and satisfies
  \begin{equation}\label{eq:rounddown}
    0\;\preceq\;A'^{(s)}_v\;\preceq\;B^{1/2}_v :
  \end{equation}
  $B^{1/2}_v$ itself; then, for $r=\kappa,\kappa-1,\dots,1$ in this order,
  the shift $B^{1/2}_v-2^{-r}\idon{d_s}$ and the scaling
  $(1-2^{-r})B^{1/2}_v$, each rounded to nearest at $\kappa$; finally the
  zero matrix, which always qualifies, $B^{1/2}_v$ being positive
  semidefinite; and $v_1$ receives $q_t\idon{d_s}$ minus their sum. The
  weight $w$, or the fall-back, is recorded, and the class is flagged as one
  whose support is not preserved: an element may then hold a non-zero entry
  in a row or column in which the handed element holds none, that is,
  outside its structural support (\xref{foundations/def:snap}).
\item \emph{All six conditions are decided last}, after (2) or (3), each of
  which changes partial sums from the touched element onward: (a)--(e) of
  Definition~\ref{def:exactadm} together with the lattice invariant
  \eqref{eq:latinv} of every entry, read on the \termref{term-target-family}{target family}
  \eqref{eq:instance-target} the datum enters. They are the six
  Theorem~\ref{thm:emit}\,(ii) decides; a family for which one of them fails
  is refused.
\end{enumerate}
A class at a sector that step~(1) does not keep is \emph{moved}, whichever
of (2) and (3) snapped it. Write $A'$ for the resulting datum --- the
\termdef{term-declared-datum}{declared datum},
an element of $\mathcal{D}_{\kappa}$ by (1) or the step of (2) and (3) that
snapped the class, each of which makes it Hermitian, on the lattice and
\termref{term-class-complete}{class complete} exactly, positive semidefinite by an exact decision in
(2) and in (3) at a weight of its ladder, and by construction in the fall-back of (3), and
by (4), which decides all six --- and
\begin{equation}\label{eq:declmodel}
  \Mdec\;\coloneqq\;M[A'],
  \qquad\text{that is}\qquad
  \Mdec_v\;=\;M_v[A']\quad(v\in\Vset),
\end{equation}
for the
\termdef{term-declared-model}{declared model}:
the \termref{term-handed-model}{handed model} with the datum replaced by $A'$, as the instance's
\termref{term-model-map}{model map} defines it. Per data sector $s$, those being the sectors the
normalizer touches and the only ones both data have,
\begin{equation}\label{eq:displacement}
  \Delta_s\;\coloneqq\;\max_{v\in\Vset}\
  \mxn{A'^{(s)}_v-A^{(s)}_v}
\end{equation}
is the
\termdef{term-emission-displacement}{emission displacement};
it is recorded per data sector, together with the kept and the moved classes,
for a moved class the step that snapped it, the weight $w$ of (2) or (3) or
the fall-back of (3), the hosts of the components of (2), and the flag of a
class snapped by (3).

The normalizer has five properties, each of which Theorem~\ref{thm:emit}, or
a claim that instantiates it, uses.
\emph{First, $A'=A$ and every $\Delta_s=0$ whenever the exact construction
already satisfies \eqref{eq:latinv}}: every class is then kept, the
normalizer is the identity and nothing below has an antecedent. That is the
ordinary case at a nominal declaration and is not the general one: every
exact step of a construction multiplies, a product of dyadics adds
significand widths, and a lattice established at the instance is closed
under none of them; that is the case in which a producer re-establishes its
lattice at an emission.
\emph{Second, the normalizer reads data sectors and nothing else.} No object
the \termref{term-hypotheses-h}{hypotheses (H)} of the instance are about --- an operator without a finite
presentation, a family on a part of $\Hs{\Bsys}$ the datum does not describe,
a register POVM --- is read, and none is repaired; so (H) holds of $\Mdec$
by the \termref{term-stability-clause}{stability clause} of
Definition~\ref{def:reduction-instance}\,(I1), $\Mdec$ is a POVM on $\Bsys$
by (I1) at $A'$, and every statement the instance makes about its model under
(H) holds of $\Mdec$. For an instance whose datum is a family of preserved
blocks below a cutoff, this is the statement that the declared model agrees
with the handed one above the cutoff and satisfies every above-cutoff fact the
handed one does.
\emph{Third, what the normalizer produces is a datum, and the family the
emitted identity is about is the model the instance builds from it.} Two
objects, two symbols, and every statement below names which it is on: $A'$ is
the declared datum, a finite list of sector blocks carried entry by entry in
the emitted target; $\Mdec$ is the declared model on $\Bsys$, the object an
adjoint landing on $\Bsys$ can equal (Theorem~\ref{thm:emit}\,(iii)). Where
the model map is the identity on a finite $\Bsys$ the two coincide; where the
model has parts the datum does not describe --- an above-cutoff family,
further branches --- they do not, and the instance's model map says how $A'$
sits inside $\Mdec$. For a branching instance
(Theorem~\ref{thm:instrument-instance}) the normalizer acts branch by branch,
so the declared model is the branches' declared models composed.
\emph{Fourth, this is a snap of what this module declares and never of what it
was handed}, which is the condition under
which a producer may snap at all.
\emph{Fifth, no step revives a zero row or column, save the last resort
(3)}: for every data sector $s$ and every $v\in\Vset$ whose class at $s$ is
not snapped by (3) the structural support of $A'^{(s)}_v$ is contained in
that of $A^{(s)}_v$. In (2) the family $\{T_v\}$ of diagonal parts is
complete at $q_t$, the diagonal of $\sum_vA^{(s)}_v=q_t\idon{d_s}$ being
$q_t$ in every place, and $T_v$ is diagonal, with non-negative entries,
non-zero only at indices of the structural support of $A^{(s)}_v$, so
$P^{w}$ is a mixture with a complete diagonal family that adds only on
coordinates of each element's own structural support; $P^{w}_v$ has the
structural support of $A^{(s)}_v$, and a rounding to nearest keeps a zero
entry zero. Each $E_C$ has its non-zero rows and columns in $C$, which the
host's structural support contains, and the entries of $E^{w}$ between two
components are zero by their definition, so $E^{w}=\sum_CE_C$ and every
non-zero entry of the residue goes to a host. In (2) the residue need not be
positive semidefinite, and the family is taken only where every element is
decided positive semidefinite. No element is mixed with a multiple of the
identity outside (3). In (3) the elements of the class in the sector's
support may gain rows and columns, which is what its flag discloses; the
elements outside the support stay zero there too. And (3) always yields a
family whose elements are positive semidefinite: at a weight of its ladder by decision, and
in the fall-back because every element of $B^{1/2}$ is positive definite, the
rounded elements satisfy \eqref{eq:rounddown}, and
$A'^{(s)}_{v_1}=B^{1/2}_{v_1}+\sum_{v\ne v_1}\big(B^{1/2}_v-A'^{(s)}_v\big)
\succeq B^{1/2}_{v_1}\succ0$ by the completeness of $B^{1/2}$.
\end{definition}
\claimmeta{\metalabel{Uses}\mref{Def.}{def:joint-family}, \mref{Def.}{def:exactadm}, \mref{Lem.}{lem:exactpsd}, \mref{Def.}{def:reduction-instance}, \extref{../foundations/contract.pdf}{foundations/def:exact-reference}, \extref{../foundations/contract.pdf}{foundations/def:snap}
\metasep \metalabel{Used by}\mref{Cor.}{cor:flag-instance}, \mref{Thm.}{thm:instrument-instance}, \mref{Cor.}{cor:register-instance}, \mref{Lem.}{lem:precomposed-instance}, \mref{Thm.}{thm:emit}, \mref{Cor.}{cor:entries}, \mref{Thm.}{thm:emitted-product}, \mref{Cor.}{cor:emitted-identity}, \mref{Rem.}{rem:emitclaim}, \mref{Rem.}{rem:claimscope}, \mref{Rem.}{rem:exactboundary}, \mref{Cor.}{cor:reference-companion}, \mref{Req.}{req:emission-refusal}, \mref{Req.}{req:emission-record}, \mref{Req.}{req:reference-companion}, \mref{Prop.}{prop:precomposed-blocks}, \mref{Prop.}{prop:flat-instance}, \mref{Cor.}{cor:emitted-flat}
\metasep newly derived
\metasep implemented (effective: accepted)
\metasep \metalabel{Code}\coderef{https://github.com/fraunhofersg/certiqs-key/blob/v0.1.0/src/certiqs_key/reductions/normalizer.py}{certiqs_key.reductions.normalizer:normalize}}

% claim reductions/thm:emit
\claimsummary{reductions/thm:emit}{The emitted target is computed exactly, decided exactly admissible before emission, and squashes the declared model exactly; nearness to the handed model bounds nothing. On blocks the record names as resting on \hyperref{../foundations/contract.pdf}{}{claim-foundations/ass:snapped-measurement}{A-SNAP}, the run's conclusions about the device rest on that premise.}
\claimbookmark{reductions/thm:emit}{Theorem}{thm:emit}{Exact emission of a reduction instance}
\claimleadin{%
This is the theorem every emission of the module rests on: what is emitted,
the displacement record aside, is computed exactly, is decided before it is
emitted, and is about the declared model of Definition~\ref{def:emitnorm}.
The identity instance applies it through
Corollary~\ref{cor:emitted-identity}, and each submodule's instance through
its own emitted corollary. Part~(v) says what relates the declared model to
the handed one: nothing proved, and in particular not the size of the move,
since a family near the constructed one need not have a key length near its
own. A run that forms the number its key length is formed from on the
reference companion (Corollary~\ref{cor:reference-companion}) needs no such
relation: it reads
the emitted target only to find points, and every data sector is
\exttermref{../foundations/contract.pdf}{term-evaluated-off-its-snap}{evaluated off its snap} (\xref{foundations/def:evaluated-off-snap}).
Elsewhere, where the normalizer moved a class, a run on that emission rests
on the premise that \xref{foundations/ass:snapped-measurement}\,(ii) states
for a producer's snap of its own emission: every conclusion about the device
that the run draws from the emitted target holds with the declared datum in
place of the handed one on the sectors concerned. Since the emitted target
squashes the declared model exactly (part~(iii)), that premise holds in
particular where the device measures the declared model there. A moved
block of the declaring instance's own snap of the measurement the handed
datum is formed from is in the same position; there the premise is
clause~(i) of the same assumption.
Part~(vi) states what a run on the emission then rests on. It is
conditional: it asserts nothing of a run whose record names no such block,
so that the premise is engaged exactly where the record names a block as
resting on it. A run whose key length is formed from a number on the
reference companion (Corollary~\ref{cor:reference-companion}) names no data
sector of this emission among them, since the companion carries none of the
normalizer's data, and of the declaring instance's snap only the blocks on
which the companion carries the snapped block; the bound then reads the
companion, and this theorem describes the target the search ran on.
Where the key length is formed instead from a number certified on the
snapped declaration, because the companion's bound counts but no point of it
can be certified (\xref{meat-core/frank-wolfe/req:certified-evaluation}\,(i)),
the run reads the emitted target there, and names the data sectors with a
moved class among the blocks resting on the premise. The six conditions of part~(ii) are those of
\xref{foundations/req:exact-emission}, and that no status is emitted for an
approximate relation is Requirement~\ref{req:no-approximate-relations}.
}
\begin{mainclaim}
\begin{theorem}[Exact emission of a reduction instance]\label{thm:emit}
Let the \exttermref{../foundations/contract.pdf}{term-declared-data}{declared data} pass every refusal rule of this module's
interface, so that every declared block and scalar it reads is
\termref{term-exactly-admissible-on-the-declared-lattice}{exactly admissible on the declared lattice}
(Definition~\ref{def:exactadm}). Let $\mathcal{J}$ be a reduction instance
(Definition~\ref{def:reduction-instance}) with handed datum $A$, and let $A'$
be the declared datum the normalizer of Definition~\ref{def:emitnorm}
produces from it (where the normalizer refuses, nothing is emitted). Let
$\{F_v\}\coloneqq F[A']$, $\Lambda\coloneqq\Lambda[A']$ and
$\CList\coloneqq\CList[A']$ be the target, the channel and the entry list of
the instance at $A'$, and $\Mdec=M[A']$ the declared model
\eqref{eq:declmodel}. Then:
\begin{enumerate}[label=(\roman*),leftmargin=3.0em]
\item \emph{(everything on the list below is computed exactly, and the
  displacement $\Delta_s$ is the one emitted number that is not on it)} every
  number on the following list is obtained from the numbers the instance was
  handed by finitely many operations in $\Fldi$, floors and comparisons: the
  data-sector values $A'^{(s)}_v$; the \termref{term-scalar-value}{scalar values}
  $\phi_s(v)$; the sector projectors; and every emitted constraint operator
  and constant of Corollary~\ref{cor:entries}. \emph{That list is everything
  this module emits with exactly one exception}: the emission displacement
  $\Delta_s$ of \eqref{eq:displacement}, an entrywise maximum \emph{modulus}
  and hence the square root of an exact rational, in general irrational, which
  is emitted only as a record of what the emission normalizer did and which
  no statement of this contract reads. \emph{No floating-point
  comparison, no eigenvalue routine, no enclosure and no interval occurs
  anywhere in that list}; the intermediate numbers an instance forms on the
  way to its datum and its scalar values are listed by the corollary that
  instantiates this theorem for it;
\item \emph{(the emitted family is \exttermref{../foundations/contract.pdf}{term-exactly-admissible}{exactly admissible}, and is decided so
  before it is emitted)} the family $\{F_v\}$, presented as the sector-block list
  \eqref{eq:instance-target}, satisfies conditions (a)--(e) of
  Definition~\ref{def:exactadm} together with \eqref{eq:latinv} of every
  entry --- six conditions --- and each of the six is \emph{decided}, exactly, on the numbers
  about to be emitted: (a) by comparison of declared labels against the
  instance's \exttermref{../foundations/contract.pdf}{term-sector-map}{sector map}; (b) and (d) by entrywise equality and (c) by
  Lemma~\ref{lem:exactpsd}, per sector block; (e) by Lemma~\ref{lem:latsum}, whose hypothesis
  \eqref{eq:latsumcond} is a declared input condition; and \eqref{eq:latinv}
  entry by entry. \emph{A family for which any of the six cannot be
  established is not emitted}.
  \emph{And the same numbers discharge what the consuming declaration asks of
  a party's side}: the tuple $\big(\Qsys,\SecSet,\{F_v\}_{v\in\Vset}\big)$
  \emph{together with the party partition} of
  Definition~\ref{def:reduction-instance}\,(I2) --- the
  \termref{term-default-party-partition}{default party partition}, every factor Bob's, unless the instance
  declares another party's label, in which case every factor is that party's
  --- satisfies the well-formedness conditions the consuming declaration
  places on \emph{that party's} side (\xref{meat-core/def:admissible-declaration}\,(i) and
  (iii), with the party labels per tensor factor demanded through its (viii))
  ---
  \begin{equation}\label{eq:emitwellformed}
    F_v=F_v^{\dagger},\qquad
    F_v\succeq0,\qquad
    \sum_{v\in\Vset}F_v=\idon{\Qsys},
  \end{equation}
  with every $F_v$ block diagonal for $\SecSet$ and
  $\Hs{\Qsys}=\bigoplus_{s\in\SecSet}\Hs{\Qsys_s}$ finite-dimensional, every
  factor of every $\Hs{\Qsys_s}$ the instance's declared party's and Bob's by
  default. \emph{Those are conditions on a party partition and not on a
  defect}, and they are a \emph{different} condition from the conditions of
  \xref{foundations/def:exact-admissibility} that (a)--(e) discharge, which
  say nothing about who owns which factor;
\item \emph{(the defining identity is about the numbers emitted, and it is an
  identity on $\Bsys$)} for every $v\in\Vset$,
  \begin{equation}\label{eq:emitidentity}
    \Lambda^{\dagger}(F_v)\;=\;\Mdec_v
    \qquad\text{exactly},
  \end{equation}
  \emph{both sides operators on $\Bsys$}, which is where $\Lambda^{\dagger}$
  lands, $\Lambda$ being a channel $\Tcl{\Bsys}\to\Tcl{\Qsys}$. \emph{Only the
  data sectors carry anything the normalizer produced}: $\Mdec$ is the handed
  model with the datum replaced by $A'$ as the instance's model map defines
  it, and everything the hypotheses (H) are about is the declared one,
  unchanged. Moreover every claim of the instance holds verbatim of the
  declared model --- the relation \eqref{eq:instance-relation} with its
  \termref{term-transported-set}{transported set} and pull-backs, every certificate of
  \eqref{eq:instance-entry}, and the \termref{term-class-marginal-identity}{class-marginal identity}
  \eqref{eq:instance-classmarginal} where the instance discharges it --- and
  Theorem~\ref{thm:transfer} and Theorem~\ref{thm:chain} hold of the one-link
  \termref{term-receiver-segment}{receiver segment} (Definition~\ref{def:link}) whose
  \termref{term-source-family}{source family} is $(\Bsys,\Mdec)$, whose target is $(\Qsys,F)$ and
  whose \termref{term-channel-witness}{channel witness} is $\Lambda$, with
  $\textsf{status}=\textsf{exact}$;
\item \emph{(the emitted bounds, and the direction each is rounded in)} every
  scalar value that was rounded to reach the lattice was rounded
  \emph{down}, as the lower bound of the quantity it certifies, and every
  complementary value was \emph{recomputed} from the rounded ones and never
  rounded on its own (Definition~\ref{def:reduction-instance}\,(I2)); and a
  caller-declared constant of modulus at most one off the lattice is rounded
  in the direction that enlarges the feasible set for its declared relation
  (Corollary~\ref{cor:entries}\,(iii)), an equality constant off the lattice
  and any constant of modulus above one being refused. \emph{Every other
  number this module emits is exact, the displacement record $\Delta_s$ of
  (i) aside}; these are all the rounding directions there are among the
  numbers a statement of this contract reads, and the direction of each is
  recorded;
\item \emph{(what relates the declared model to the handed one)} where some
  $\Delta_s>0$ the emitted numbers are a declaration of a family near the
  constructed one, and nearness is not what relates them: \emph{no bound
  relates a key length computed from $\Mdec$ alone to one computed from the
  handed model $M[A]$, or from the supplier's own output}, and this contract
  has no theorem converting an operator perturbation of the target into a
  security penalty, nor any other relation between $A'$ and $A$ on a moved
  class: the consuming contract's impossibility
  (\xref{meat-core/prop:optol-nobound}) is the whole of what is proved
  there; what a run rests on there instead is~(vi);
\item \emph{(what a run on the emission rests on where it reads a moved
  block, and only there)} let $B=B_{\mathrm{snap}}\sqcup B_{\mathrm{em}}$
  be the blocks that the record of the run names, for this artefact, as
  resting on \xref{foundations/ass:snapped-measurement}, kept as two
  disjoint parts, since they are blocks of two different snaps and their
  labels may coincide: $B_{\mathrm{em}}$ holds the data sectors of the
  emission so named, and $B_{\mathrm{snap}}$, where $A$ is formed from
  measurement blocks of a snap (\xref{foundations/def:snap}) that the
  declaring party made of the measurement it states, the blocks of that
  snap so named; the declaring party is the instance that states the
  measurement of the device, or a reference of it, from which $A$ is formed.
  Where $B$ is empty, this part asserts nothing, and no part of this theorem
  rests on that assumption. Where $B$ is not empty, let that assumption
  hold. Then on every block in $B_{\mathrm{snap}}$, the block of the
  measurement the device performs is the snapped block, with the entries
  recorded for that snap, from which $A$ is formed, in place of the block of
  the model the declaring party states; and every conclusion about the
  device that the claims of the run draw from the emitted target, under
  their hypothesis that its data sectors carry the handed datum $A$, holds
  when they carry the declared datum $A'$ on every data sector in
  $B_{\mathrm{em}}$, which by~(iii) is the datum whose model $\Mdec$ the
  emitted target squashes exactly. The run carries no bound relating its key
  length to the declaring party's model on the blocks in
  $B_{\mathrm{snap}}$, nor to $M[A]$ on the data sectors in
  $B_{\mathrm{em}}$: no claim relates that model to the snapped block on a
  block in $B_{\mathrm{snap}}$, nor $A$ to $A'$ on a data sector in
  $B_{\mathrm{em}}$.
\end{enumerate}
\end{theorem}
\end{mainclaim}
\claimmeta{\metalabel{Uses}\mref{Thm.}{thm:transfer}, \mref{Def.}{def:joint-family}, \mref{Cor.}{cor:blockdiag}, \mref{Def.}{def:link}, \mref{Def.}{def:chain}, \mref{Thm.}{thm:chain}, \mref{Def.}{def:exactadm}, \mref{Def.}{def:substrate}, \mref{Lem.}{lem:exactpsd}, \mref{Lem.}{lem:latsum}, \mref{Def.}{def:reduction-instance}, \mref{Def.}{def:emitnorm}, \extref{../foundations/contract.pdf}{foundations/def:snap}
\metasep \metalabel{Assumes}\extaref{../foundations/contract.pdf}{foundations/ass:snapped-measurement}{A-SNAP}
\metasep \metalabel{Used by}\mref{Rem.}{rem:why-exact}, \mref{Cor.}{cor:entries}, \mref{Thm.}{thm:emitted-product}, \mref{Cor.}{cor:emitted-identity}, \mref{Rem.}{rem:emitclaim}, \mref{Rem.}{rem:claimscope}, \mref{Req.}{req:emission-refusal}, \mref{Req.}{req:emission-record}, \mref{Cor.}{cor:emitted-flat}, \mref{Cor.}{cor:two-sector-decided}, \mref{Cor.}{cor:normalized-flat}, \extref{../protocol-compiler/contract.pdf}{protocol-compiler/req:link-row-provenance}
\metasep newly derived
\metasep implemented (effective: accepted)
\metasep \metalabel{Code}\coderef{https://github.com/fraunhofersg/certiqs-key/blob/v0.1.0/src/certiqs_key/reductions/emit.py}{certiqs_key.reductions.emit:reduce}
\metasep \metalabel{Oracles}\coderef{https://github.com/fraunhofersg/certiqs-key/blob/v0.1.0/tests/oracles/registry_reductions.py}{registry_reductions.py}}

\begin{proof}
\emph{Idea.} Parts (i) and (ii) follow what the instance and the normalizer
compute and what the six conditions ask of it; part (iii) is the
\termref{term-stability-clause}{stability clause} of the instance applied at the declared datum; part
(iv) is the rounding rule of the instance's scalar values and of a caller's
constants; part (vi) applies the two clauses of the snapped-measurement
assumption to the blocks the record names.

\emph{Part (i).} Every entry of a declared block and every declared scalar is
a dyadic rational (Definition~\ref{def:exactadm}\,(a),(b)), hence an element
of $\Fldi$ resp.\ $\Fld$. The handed datum is formed from them by finitely
many operations in $\Fldi$ (Definition~\ref{def:reduction-instance}\,(I1));
the scalar values by finitely many operations in $\Fld$ and downward
roundings $\lat{\cdot}$, each a floor and a comparison in $\Fld$ (I2); the
sector projectors are the indicator values $0$ and $1$; each emitted
constraint operator of Corollary~\ref{cor:entries} is a real combination
\eqref{eq:instance-entry} of emitted blocks and sector projectors with
lattice coefficients, and each constant is declared or rounded by a floor or
a ceiling in $\Fld$; and each of the normalizer's steps (1)--(4) is a
rounding to nearest, a comparison, an addition, a subtraction, a product
with a rational number, a positive semidefiniteness decision, or a test of
which entries are zero, and its mixing weights are the dyadic numbers
$2^{-l}$. Field operations, comparisons and floors in $\Fldi$ are exact by the
definition of the substrate (Definition~\ref{def:substrate}), and the only
non-arithmetical predicate that occurs --- positive semidefiniteness, the
form every operator inequality of this contract takes, the L\"owner order
being the only order it uses --- is decided by sign comparisons in $\Fld$
(Lemma~\ref{lem:exactpsd}\,(i), \eqref{eq:psdtest}); the other conditions
of\,(ii) above are equalities, compared entry by entry in $\Fldi$, and
comparisons of declared labels. \emph{No bound in the operator norm is decided
here or anywhere on the emission path}: Lemma~\ref{lem:exactpsd}\,(ii) is
correct as it stands and is carried without a client. \emph{The one emitted
number the list omits is $\Delta_s$}, and it is omitted because
\eqref{eq:displacement} takes an entrywise maximum \emph{modulus} of a complex
difference, which is a square root of an exact rational and in general leaves
$\Fld$. It is emitted only as a record
(Requirement~\ref{req:emission-record}), so the statement above is a
statement about the enumerated list and about everything this module emits
except that one number; no step of this proof, and no statement of this
contract, is proved from its value (Remark~\ref{rem:claimscope}).

\emph{Part (ii).} Each of the six is decidable on a finite list of elements
of $\Fldi$: (a) is a comparison of declared labels; (b) is an entrywise
equality; (c) is \eqref{eq:psdtest}; (d) is an entrywise equality between two
finite sums; \eqref{eq:latinv} is, per entry, integrality of $2^{\kappa}x$
together with $|x|\le1$, two comparisons in $\Fld$; and (e) is
Lemma~\ref{lem:latsum}, whose hypothesis is the declared input condition
\eqref{eq:latsumcond} and whose conclusion covers every entry position, the
accumulation being entrywise.
\emph{That (b), (c) and (d) hold of the emitted family} is, on a data sector,
Definition~\ref{def:emitnorm}, whose step (1) keeps exactly admissible
classes, whose steps (2) and (3) establish Hermiticity, \eqref{eq:latinv}
and class completeness for a moved class and establish or decide its
positivity, and whose step (4) decides all six on the result ---
completeness over $\Vset$ following from class completeness and
$\sum_tq_t=1$ (Definition~\ref{def:reduction-instance}, first consequence);
and, on a \termref{term-scalar-sector}{scalar sector}, the \exttermref{../foundations/contract.pdf}{term-declared-value}{declared values}
$\phi_s(v)\in[0,1]$ on the
lattice with $\sum_v\phi_s(v)=1$ exactly (I2), which give a real, positive
semidefinite and complete block outright.
\emph{So the decision of step~(4) is a defensive exact check}: on declared
data that pass the refusal rules, as the theorem assumes, every path of the
normalizer produces a family the six conditions hold of, and no path reaches
the refusal it guards. The refusal of a class that can occur is the strict
option's refusal of a moved class (Requirement~\ref{req:emission-refusal}).
\emph{The consuming declaration's own conditions follow from what has just
been established}: (b), (c) and (d) are \eqref{eq:emitwellformed} --- the
families are POVMs and every $F_v$ is sector-block-diagonal --- which is
\xref{meat-core/def:admissible-declaration}\,(iii), block diagonality for $\SecSet$ being by
construction; and $\dim\Hs{\Qsys}=\sum_{s\in\SecSet}d_s$ is a finite sum of
declared dimensions (Definition~\ref{def:reduction-instance}\,(I2), its fifth
consequence), so $\Hs{\Qsys}$ is finite-dimensional and the alphabet $\Vset$
finite, which is the part of its (i) this tuple is about. \emph{And the
default party partition assigns every factor of every $\Hs{\Qsys_s}$ to Bob,
which is the assignment the physical model already makes}: the target
describes what reaches Bob's receiver, the instance's system $\Bsys$ being
his. Where the instance declares another party's label
(Definition~\ref{def:reduction-instance}\,(I2)) the same sentence holds of
that party, whose system $\Bsys$ then is; what is emitted is the declared
label either way, no step above reads its value, and an emission of several
\termref{term-measuring-parties}{measuring parties} carries one label per factor
(Theorem~\ref{thm:emitted-product}\,(i)). \emph{The party labels themselves
are neither (i) nor (iii)}: a declaration names a party partition $\Party$
(\xref{protocol-compiler/def:declaration}) and \xref{meat-core/def:admissible-declaration}\,(viii) is
what demands it; what is emitted here is that partition's value on the
target's sectors. Nothing in this paragraph is measured and nothing is
concluded from a returned number.

\emph{Part (iii).} By Definition~\ref{def:emitnorm}\,(1)--(4),
$A'\in\mathcal{D}_{\kappa}\subseteq\mathcal{D}_{\mathrm{adm}}$: it is
\termref{term-class-complete}{class complete} at $q_t$ by (1) or by the step of (2) and (3) that
snapped the class, the last step that changes an element on every path, and
Hermitian, positive semidefinite and on the lattice per data sector by the
decision of (4). By
the second property of that definition the normalizer read no object the
hypotheses (H) are about, so (H) holds of $M[A']$ by the stability clause of
Definition~\ref{def:reduction-instance}\,(I1), and $M[A']=\Mdec$ is a POVM on
$\Bsys$ by (I1). Parts (I3) and (I4) of the instance are established for
\emph{every} \termref{term-admissible-datum}{admissible datum} under (H); at $A'$ they read
$(\Mdec;O[A'])\dsq_{\Lambda}(F;\hat O[A'])$, whose first clause is
\eqref{eq:emitidentity}, together with every certificate of $\CList$ and,
where the flag says so, \eqref{eq:instance-classmarginal} at $A'$. That is
what ``holds verbatim of the declared model'' asserts, and nothing was
re-proved: the instance at $A'$ is the same construction with a different
admissible datum, which is what the stability clause was stated for. The
sector pinching in front of $\Lambda_0[A']$ changes none of it
(Corollary~\ref{cor:blockdiag}\,(i), the second consequence of
Definition~\ref{def:reduction-instance}). Theorem~\ref{thm:transfer} needs
only the relation just established; the one-link segment with source family
$(\Bsys,\Mdec)$, target $(\Qsys,F)$, channel witness $\Lambda$, the transported
pair $(\hat O[A'],O[A'])$, the \termref{term-reachable-set-statement}{reachable-set statement} of that channel
and the entry list $\CList$ is a receiver-stage link of
Definition~\ref{def:link} and a chain of Definition~\ref{def:chain} with an
empty \exttermref{../transmitter/contract.pdf}{term-source-segment}{source segment}, and Theorem~\ref{thm:chain}\,(iv) gives it
$\textsf{status}=\textsf{exact}$ because \eqref{eq:emitidentity} carries no
error term. \emph{Every remaining ingredient is common to the handed model and
the declared one}: $\Bsys$, the alphabet with its classes and weights, the
sector map, the supports, the scalar values, the constructor, and every object
of (H), which the normalizer never reads.

\emph{Part (iv).} The scalar values are (I2) of
Definition~\ref{def:reduction-instance}, which requires the downward direction
and the recomputed complement; the constants are
Corollary~\ref{cor:entries}\,(iii). \emph{That a complement is recomputed
rather than rounded is what makes a scalar sector's completeness an identity
of lattice values}: the rounded values and their complement are formed from
the same finitely many rounded quantities, so $\sum_v\phi_s(v)=1$ holds
exactly, and Lemma~\ref{lem:latsum} makes the accumulation that checks it
exact as well. Rounding the two independently would leave a residue of up to
$2^{-\kappa}$ on the sector and would fail (d) --- which is the case that
\xref{foundations/req:rounding-direction} closes by requiring the
complementary quantity to be recomputed.

\emph{Part (v).} It asserts the absence of a bound and of a relation and
needs no proof; the consuming contract rules the same inference out from its
own side (\xref{meat-core/prop:optol-nobound}).

\emph{Part (vi).} Where $B$ is empty there is nothing to prove, and no step
of the proofs of (i)--(v) reads
\xref{foundations/ass:snapped-measurement}. Otherwise each of
$B_{\mathrm{snap}}$ and $B_{\mathrm{em}}$ is, for its own snap, the set that
the assumption calls $B$: the blocks of that snap the record of the run
names as resting on it. The assumption is applied to the two snaps
separately, and a label in both parts names two different blocks.

\emph{Blocks in $B_{\mathrm{snap}}$.} By the statement of (vi), the
declaring party is the instance that states the measurement $A$ is formed
from, so the snap whose blocks form
$B_{\mathrm{snap}}$ is the declaring instance's snap of the measurement of
the device, or of a reference of it, that it states. Its $D$ is therefore
data an instance states as the measurement of a device, or a reference of
it, and clause~(i) of the assumption applies: on each block in
$B_{\mathrm{snap}}$ the block of the measurement the device performs is the
snapped family with the entries recorded for the snap, which are the entries
$A$ is formed from on that block, and a statement of another assumption of
the run that the declaring party's model describes the device there is not
asserted.

\emph{Data sectors in $B_{\mathrm{em}}$.} A data sector of the emission is a
block of the normalizer's snap of the handed datum $A$ to the declared datum
$A'$ (Definition~\ref{def:emitnorm}), a producer's snap of what it emits
whose exact construction is $A$, so clause~(ii) applies to it with $A'$ as
the snap: every conclusion about the device that the claims of the run draw
from the emitted family, under their hypothesis that it is $A$, holds when
it is $A'^{(s)}$ on every data sector $s\in B_{\mathrm{em}}$. The data
sectors of the emitted target carry the values $A'^{(s)}_v$ (part~(i)), and
the target squashes $\Mdec=M[A']$ exactly (part~(iii)).

\emph{What the run does not carry.} The argument is that of part~(v): a
nearness of the families relates no key lengths
(\xref{meat-core/prop:optol-nobound}), and no claim of these contracts
relates a moved block of a snap to the data it was made from. A block in
$B_{\mathrm{snap}}$ is a moved block of the declaring instance's snap, by the
assumption's definition of $B$, and a data sector in $B_{\mathrm{em}}$ holds
a moved class of the normalizer's; so the run carries no relation on these
blocks. The sentence is confined to them: it says nothing of a relation to
$M[A]$ on a block in $B_{\mathrm{snap}}$, where $A$ is formed from the
snapped block and part~(iii) stands as it is established for the data
sectors, and nothing of a block of either snap that the record does not
name.
\end{proof}

% claim reductions/cor:entries
\claimsummary{reductions/cor:entries}{The emitted list is empty unless the request carries entries; then it holds the instance's entries, the class-marginal entry where discharged and further entries whose certificate a claim decides. Certificates hold of the emitted numbers, and off-lattice inequality constants are rounded outward.}
\claimbookmark{reductions/cor:entries}{Corollary}{cor:entries}{The emitted constraint entries}
\claimleadin{%
A request either carries entries or carries none, and there is no third
case; whether the emitted constraint list is empty, and with it whether the
attack-set licence is present, depends on that alone. The flag-state
submodule's encoded and licensed routes are that submodule's instance of the
two cases, named here for the reader; this module declares no route key of
its own. The entries below are untagged; where an entry carries the
party-local tag instead is Requirement~\ref{req:party-local-tag}, and
Remark~\ref{rem:entry-tags} explains why a one-factor emission carries none.
Remark~\ref{rem:unsound} identifies the expectation form of the
class-marginal identity, read in part (ii), as its sound squashed form.
}
\begin{corollary}[The emitted constraint entries]\label{cor:entries}
Assume the hypotheses of Theorem~\ref{thm:emit} for a \termref{term-reduction-instance}{reduction instance}
$\mathcal{J}$ with declared datum $A'$, target $\{F_v\}$, channel $\Lambda$
and declared model $\Mdec$. A request \emph{carries entries} when it selects
the case in which the instance's own entries are emitted: in the flag-state
submodule that case is the licensed route; in a core-only artefact it is a
caller's request for an entry --- for a default carried operator as an entry,
or for the class-marginal entry. \emph{What is emitted is decided by the
request}: \emph{where the request carries no entry the emitted constraint
list is empty} --- whatever entries the instance's own list $\CList[A']$
contains --- and the artefact carries no attack-set licence; \emph{where the
request carries entries the emitted list is the instance's own entries, the
class-marginal entry on request and any further requested entry that (iii)
admits}, so that
the list is non-empty and the licence of Theorem~\ref{thm:licence} is
present. In \emph{both} cases the
certificate $(\Pi_s,\succeq,0)$ of every sector projector of every instance
is carried \emph{per entry in the transported set}, with the certificate case
of Corollary~\ref{cor:pullback} that makes each usable, and not as a
constraint entry. Then
\emph{(i) and the identity \eqref{eq:entryclass} of (ii) of every request,
and what each of (ii) and (iii) emits of a request that carries entries}:
\begin{enumerate}[label=(\roman*),leftmargin=3.0em]
\item \emph{(the default \termref{term-carried-operators}{carried operators})} every operator of the
  instance's \emph{default carried set} --- every sector projector $\Pi_s$,
  $s\in\SecSet$, which Definition~\ref{def:reduction-instance}\,(I3) puts in
  the transported set of every instance, together with the sums of sector
  projectors that instance transports besides (for a target split at a
  cutoff, the projector onto the blocks below it and the projector onto the
  flag blocks; for a branching instance, the branch projector $\Pi_j$) ---
  is emitted at $c=0$ with $\bowtie\,=\,\succeq$: its entries
  are $0$ and $1$, hence on the lattice at every $\kappa\ge1$, and its
  pull-back
  \begin{equation}\label{eq:entrysector}
    \Lambda^{\dagger}(\Pi_s)\;\succeq\;0
  \end{equation}
  is exactly positive for every instance, as therefore is the pull-back of
  every sum of them. \emph{It is emitted with the transported set}, per entry
  and whatever the request carries, as the certificate case of
  Corollary~\ref{cor:pullback}\,(ii) that entry records, and it becomes an
  entry of $\CList$ only where a caller requests it \emph{as an entry} ---
  a request that makes the emitted list non-empty and so brings the licence
  with it;
\item \emph{(the class-marginal entry)} For $\hat O=\sum_{v\in\Vset_t}F_v$ at
  a declared $t\in\Tset$ --- the expectation form of the physical
  \termref{term-class-marginal-identity}{class-marginal identity} \eqref{eq:classmarginal} ---
  \begin{equation}\label{eq:entryclass}
    \Lambda^{\dagger}\big(\hat O\big)\;=\;q_t\,\idon{\Bsys}
    \qquad\text{exactly}
  \end{equation}
  holds \emph{whenever} the instance discharges the class-marginal identity
  \eqref{eq:instance-classmarginal}, and is established by this contract
  \emph{only} there --- \emph{a fact about the instance and the channel, which
  this part asserts of every request}. \emph{Where the request carries
  entries}, so that the emitted list is non-empty and the licence is present,
  the entry is then emitted, on request, as the single equality entry
  $(\hat O,=,q_t)$ at the declared $q_t$; where the instance does not
  discharge it the request is refused, and the entry is never emitted
  on the strength of a branch that does not establish it;
\item \emph{(any further entry --- the general rule)} \emph{under the same
  hypothesis as (ii)'s emission sentence --- the request carries entries},
  let
  $\hat O=\sum_{v\in\Vset}\sum_{s\in\SecSet}\alpha_{v,s}\Pi_sF_v\Pi_s+Z$ with
  declared real lattice coefficients $\alpha_{v,s}$ and $Z$ a fixed operator
  of the declared sector decomposition (a sector projector or a combination
  of them whose coefficients are declared lattice values, as
  Definition~\ref{def:reduction-instance}\,(I4) requires), the form
  \eqref{eq:instance-entry}, with a relation
  $\bowtie$ and a declared real constant $c$.
  \emph{Where a declared constant of modulus at most one does not satisfy
  \eqref{eq:latinv} it is first rounded onto the lattice in the direction that
  enlarges the feasible set for $\bowtie$} --- down for $\succeq$, up for
  $\preceq$ --- and \emph{an equality entry whose constant is off the lattice
  is refused rather than split into a pair}, an equality having no loosening
  direction; \emph{a declared constant of modulus above one is refused under
  every relation}. Suppose the certificate
  \begin{equation}\label{eq:entrygeneral}
    \Lambda^{\dagger}(\hat O)\ \bowtie\ c\,\idon{\Bsys}
  \end{equation}
  is established for the emitted numbers, at the constant $c$ so obtained, in
  one of two ways: for an entry of the instance's own list $\CList[A']$, by
  (I4) of the instance; for a requested entry, by a \emph{decided
  certificate} --- a condition on the emitted numbers and the request's
  declared constants, decided exactly from them, that a claim stated for the
  instance proves sufficient for \eqref{eq:entrygeneral} under the instance's
  hypotheses (H) --- that holds there. Then the entry is emitted at that $c$.
  \emph{No other requested entry is emitted}: a requested operator of this
  form for which no claim stated for the instance gives a decided
  certificate, or whose decided certificate fails at the emitted numbers, is
  refused, and no certificate is taken from the request on trust.
\end{enumerate}
Every entry that (i)--(iii) emit is \emph{untagged}, $(\hat O,\bowtie,c)$,
and is read in the form of the declared source mode
(Definition~\ref{def:licensed}).

\emph{Each certificate above is discharged for the emitted numbers
themselves}, those numbers being the family every part is read on
(Theorem~\ref{thm:emit}\,(iii)), so no constant is shifted to pay for a
displacement between the family a certificate is asserted of and the numbers
its operator is assembled from. The rounding rule of (iii) applies only to a
constant a caller declares, never to one this module or an instance
computes. \emph{Together, (i)--(iii) cover every entry this module can
emit}: the default carried operators of every instance, the class-marginal
entry a caller may request, and every entry of the form
\eqref{eq:entrygeneral} --- which includes every entry an instance carries in
its own list, such as a flag-sector entry of a weight-preserving squasher,
and the branch entries a branching instance inherits
(Theorem~\ref{thm:instrument-instance},
Corollary~\ref{cor:register-instance}). \emph{A requested operator outside
that form, or of that form without a decided certificate that holds, is not
emitted at some other constant; it is refused}. So the
three parts are
exhaustive over what is \emph{emitted}, which is what the artefact carries.
\emph{Where the request carries no entry there is nothing for (ii)--(iii) to
be exhaustive about}: $\CList=\varnothing$ there, whatever entries the
instance's own list contains, and the certificate of every default carried
operator travels with the transported set alone.
\end{corollary}
\claimmeta{\metalabel{Uses}\mref{Cor.}{cor:pullback}, \mref{Def.}{def:joint-family}, \mref{Def.}{def:licensed}, \mref{Thm.}{thm:licence}, \mref{Def.}{def:exactadm}, \mref{Def.}{def:reduction-instance}, \mref{Def.}{def:emitnorm}, \mref{Thm.}{thm:emit}
\metasep \metalabel{Assumes}\aref{A-CONST}{ass:constdecl}
\metasep \metalabel{Used by}\mref{Rem.}{rem:entry-tags}, \mref{Cor.}{cor:emitted-identity}, \mref{Rem.}{rem:exactboundary}, \mref{Cor.}{cor:reference-companion}, \mref{Req.}{req:licence-record}, \mref{Req.}{req:constraint-list-by-request}, \mref{Req.}{req:transported-set-certificates}, \mref{Req.}{req:class-marginal-entry}, \mref{Req.}{req:caller-constant-rounding}, \mref{Req.}{req:uncertified-operator-refused}, \mref{Req.}{req:reference-companion}, \mref{Cor.}{cor:emitted-flat}, \mref{Cor.}{cor:two-sector-decided}, \extref{../protocol-compiler/contract.pdf}{protocol-compiler/req:link-row-provenance}
\metasep newly derived
\metasep implemented (effective: accepted)
\metasep \metalabel{Code}\coderef{https://github.com/fraunhofersg/certiqs-key/blob/v0.1.0/src/certiqs_key/reductions/constraints.py}{certiqs_key.reductions.constraints:certify_carried_operators}}

\begin{proof}
(i) The sector projectors are determined by the declared dimensions and
order alone (Definition~\ref{def:reduction-instance}\,(I2)); each is in the
transported set by (I3), and its pull-back is positive semidefinite because
$\Lambda^{\dagger}$ is a positive map --- the third consequence of
Definition~\ref{def:reduction-instance} --- which is
\eqref{eq:certificate} with $c=0$ and $\bowtie\,=\,\succeq$, the certificate
case of Corollary~\ref{cor:pullback}\,(ii). A sum of distinct sector
projectors is again a projector of the declared decomposition, with entries
$0$ and $1$ and with the sum of their pull-backs as its own, positive by
additivity; those sums are the remaining operators of the default carried
set. \emph{Where each is emitted is the transported set}, which (I3) is
about and which the artefact carries in either case with the certificate
case per entry; nothing in this part makes it an entry of $\CList$, and
where the request carries no entry $\CList$ is empty. A caller who asks for
one of them \emph{as an entry} asks for an entry, and gets it at this
constant --- the list being non-empty then, and the licence present.
\emph{Nothing in this part uses which of the two cases holds}, so it is
asserted of every request.

(ii) By Theorem~\ref{thm:emit}\,(iii) and linearity,
$\Lambda^{\dagger}(\hat O)=\sum_{v\in\Vset_t}\Mdec_v$, the class sum of the
declared model (the fourth consequence of
Definition~\ref{def:reduction-instance}). Where the instance discharges
\eqref{eq:instance-classmarginal}, it does so for every admissible datum
under (H), hence at $A'$ by Theorem~\ref{thm:emit}\,(iii), so the class sum
is $q_t\idon{\Bsys}$ and \eqref{eq:entryclass} is \eqref{eq:certificate}
with $\bowtie\,=\,=$ at $c=q_t$; the constant is declared before the run
(Assumption~\ref{ass:constdecl}). That derivation reads the instance and the
channel alone: \emph{nothing in it uses which of the two cases holds}, so
\eqref{eq:entryclass} is asserted of every request, exactly as (i) is. What
the request decides is whether the entry is \emph{emitted} --- where it
carries entries the list is non-empty and the licence present, and the entry
joins the list on request. That is the whole of the implication, and
\emph{no converse is asserted}: where the instance does not discharge
\eqref{eq:instance-classmarginal}, \eqref{eq:entryclass} may still hold of
the device and nothing in this contract establishes it --- an
uncharacterised branch contributes some operator between $0$ and its branch
weight, not $q_t$ times it (Corollary~\ref{cor:flag-instance}) --- and
emitting the entry on an undischarged identity would impose
$\tr[\hat O\sigma]=q_t$ on reachable states for which it may be false, the
security break of Remark~\ref{rem:unsound}; so the request is refused, which
is the safe direction and the only one this corollary licenses.

(iii) The entry is emitted at the constant its certificate is established
at, which is the hypothesis; that the constant is declared and not observed
is Assumption~\ref{ass:constdecl}. An instance's own entries carry their
certificates for every admissible datum
(Definition~\ref{def:reduction-instance}\,(I4)), hence at $A'$. A requested
entry's decided certificate holds at the emitted numbers by hypothesis, and
the claim that gives it proves \eqref{eq:entrygeneral} from it under (H),
which holds of the instance at $A'$ by the stability clause of
Definition~\ref{def:reduction-instance}, $A'$ being admissible
(Theorem~\ref{thm:emit}\,(iii)); so \eqref{eq:entrygeneral} holds at the
emitted numbers, and nothing about it is taken from the request. A requested
operator without a decided certificate that holds has nothing that
establishes \eqref{eq:entrygeneral} for it, and emitting it would impose on
reachable states a condition that may be false of them, the security break
of Remark~\ref{rem:unsound}; refusing it is the safe direction, nothing
being emitted. \emph{The rounding rule}: let $c'$ be a declared constant off
the lattice with $|c'|\le1$ and $c$ its rounding in the stated direction, so
$c\le c'$ when $\bowtie$ is $\succeq$ and $c\ge c'$ when $\bowtie$ is
$\preceq$; since $-1$ and $1$ are lattice values, $c$ lies in $[-1,1]$ and
satisfies \eqref{eq:latinv}. The certificate is established at $c$, the
constant emitted; and wherever it holds at $c'$ it holds at $c$, since
$\Lambda^{\dagger}(\hat O)\succeq c'\idon{\Bsys}\succeq c\,\idon{\Bsys}$ in
the first case and $\preceq$ in the second, so the rounding loses no entry
whose certificate holds at the declared constant, while $\AttSet_{\CList}$
only grows --- which is the safe direction, a smaller admissible set being
the one in which Theorem~\ref{thm:licence}'s inclusion can fail. \emph{For an
equality entry neither direction survives}, which is why such an entry is
refused rather than rounded. A constant of modulus above one is refused under
every relation: on its enlarging side there need be no value satisfying
\eqref{eq:latinv} at all ($\preceq$ at $c'>1$, $\succeq$ at $c'<-1$), and
refusing is the safe direction, nothing being emitted.
\end{proof}

% claim reductions/rem:entry-tags
\claimsummary{reductions/rem:entry-tags}{The licence covers a party-local entry only through a certificate on its own factor's channel; a joint one can hold while the operator inequality fails. On a one-factor emission the untagged entry is implied by the party-local reading, and rate is lost only in untrusted-source mode at d\_A \textgreater{} 1.}
\claimbookmark{reductions/rem:entry-tags}{Remark}{rem:entry-tags}{Party-local tags and the one-factor emission}
\claimleadin{%
The entries of Corollary~\ref{cor:entries} are untagged; the party-local tag
exists for a \termref{term-measured-system}{measured system} that is a declared product of party-scoped
factors (Requirement~\ref{req:party-local-tag}). This remark says what
licenses the tagged reading and what a one-factor emission gives up by
emitting the untagged one.
}
\begin{remark}[Party-local tags and the one-factor emission]\label{rem:entry-tags}
An entry \termref{term-scoped-to-a-factor}{scoped to a factor} $f$ of a declared product
$\Qsys=\bigotimes_f\Qsys_f$ is read as the operator inequality
\eqref{eq:partylocal-entry} on the other systems of $\omega$, with $\hat O$
an operator on $\Qsys_f$ alone. \emph{What licenses that reading is a
certificate discharged on the factor's own channel} $\Lambda_f$, which is the
hypothesis of Theorem~\ref{thm:licence}\,(v); a certificate on the joint
channel does not license it, and does not almost license it,
Lemma~\ref{lem:party-local}\,(iv) exhibiting a channel that satisfies the
scalar bound and violates the operator one.

\emph{On a one-factor emission the untagged entry is safe, whatever the
case}: it is \emph{implied} by the
\termref{term-party-local-form}{party-local reading} of the same
$(\hat O,\bowtie,c)$ --- in
\exttermref{../transmitter/contract.pdf}{term-untrusted-source-mode}{untrusted-source mode} it is that reading's trace
(Lemma~\ref{lem:party-local}\,(iii)), and in \exttermref{../transmitter/contract.pdf}{term-trusted-source-mode}{trusted-source mode} it is
the same operator inequality under the marginal condition $\omega_A=\tau_A$
--- and the licence covers it on the very same certificate, by
Theorem~\ref{thm:licence}\,(i)--(iv). \emph{What it gives up does depend on
the case, and Definition~\ref{def:licensed} states all three}: in
trusted-source mode under the consumer's marginal condition
$\omega_A=\tau_A$, and in untrusted-source mode at $d_A=1$, the untagged entry
\emph{is} the party-local reading and nothing is given up; in untrusted-source
mode with $d_A>1$ the party-local reading is an operator inequality on $A$
against the marginal that is actually there, of which the untagged entry is
only the trace, so what is emitted is strictly weaker --- a rate loss and
never a soundness loss. Emitting the stronger form there would need a tag
naming a scope that no declared product provides; widening it is an amendment
of Definition~\ref{def:licensed} and of the consuming \exttermref{../foundations/contract.pdf}{term-round-model}{round model}.
Theorem~\ref{thm:emitted-product}\,(vi) says the same of the one-factor
emission as a whole.
\end{remark}
\claimmeta{\metalabel{Uses}\mref{Def.}{def:licensed}, \mref{Lem.}{lem:party-local}, \mref{Thm.}{thm:licence}, \mref{Cor.}{cor:entries}
\metasep \metalabel{Used by}\mref{Req.}{req:party-local-tag}
\metasep newly derived
\metasep accepted}

% claim reductions/thm:emitted-product
\claimsummary{reductions/thm:emitted-product}{Several measuring parties are emitted as one exactly admissible artefact per party with a combination rule, never a product matrix. The product target is a POVM squashing the product of the declared models, and every entry keeps its constant, scoped to its factor.}
\claimbookmark{reductions/thm:emitted-product}{Theorem}{thm:emitted-product}{The emission of several measuring parties}
\claimleadin{%
Where several parties measure, each party's measurement is its own
reduction instance and is emitted as one; the theorem says what the product
of the emitted factors is, without forming a product datum. At one factor its
parts (i), (ii), (iv) and (v) are Theorem~\ref{thm:emit} again and its part
(iii) is Theorem~\ref{thm:licence}\,(v) (part~(vi)). A class-marginal
identity is emitted in expectation form, as in Remark~\ref{rem:unsound}, and
what a factor's emission records is Requirement~\ref{req:emission-record}.
}
\begin{theorem}[The emission of several measuring parties]
\label{thm:emitted-product}
Let $\mathbb{F}$ be a finite ordered set of \emph{factors}, each $f$ with a
party label, and for each $f$ let $\mathcal{J}^{f}$ be a
\termref{term-reduction-instance}{reduction instance} (Definition~\ref{def:reduction-instance}) on a
system $\Bsys_f$ over
$\big(\Vset^{f},\{\Vset^{f}_{t}\}_{t\in\Tset^{f}},\{q^{f}_{t}\}\big)$ to
which Theorem~\ref{thm:emit} applies, with emitted datum $A'_f$, emitted
target $F^{f}\coloneqq F^{f}[A'_f]$ on
$\Qsys_f=\bigoplus_{s\in\SecSet^{f}}\Qsys_{f,s}$, channel
$\Lambda_f\coloneqq\Lambda_f[A'_f]$, declared model ${M'}^{\mathrm{ph},f}=M^{f}[A'_f]$
and entry list $\CList^{f}$. Put $\Bsys\coloneqq\bigotimes_f\Bsys_f$,
$\Qsys\coloneqq\bigotimes_f\Qsys_f$, $\Vset\coloneqq\prod_f\Vset^{f}$,
$\Lambda\coloneqq\bigotimes_f\Lambda_f$, all products in the order of
$\mathbb{F}$, and for $v=(v_f)_f\in\Vset$
\begin{equation}\label{eq:producttarget}
  F_v\coloneqq\bigotimes_{f}F^{f}_{v_f},\qquad
  \Mdec_v\coloneqq\bigotimes_{f}{M'}^{\mathrm{ph},f}_{v_f} .
\end{equation}
Suppose that \emph{what is emitted} is the ordered list of the factors'
artefacts, each with its party label, and the rule \eqref{eq:producttarget},
and that no entry of any $F_v$ is emitted, declared or rounded. Then:
\begin{enumerate}[label=(\roman*),leftmargin=3.0em]
\item \emph{(a measurement, block diagonal for the product sector map)}
  $\{F_v\}_{v\in\Vset}$ is a POVM on $\Qsys$: every $F_v$ is Hermitian and
  positive semidefinite and $\sum_{v\in\Vset}F_v=\idon{\Qsys}$, as identities
  of the exact rational numbers the factors emitted. With the \emph{product
  sector map} --- index set $\prod_f\SecSet^{f}$ in the product order, the
  last factor fastest, identifiers the factors' identifiers joined in that
  order, dimensions the products, projectors
  $\Pi_{s}=\bigotimes_f\Pi^{f}_{s_f}$ --- every $F_v$ is block diagonal, its
  block at $s=(s_f)_f$ being the Kronecker product
  $\bigotimes_fF^{f,(s_f)}_{v_f}$ of the factors' blocks. The party
  partition assigns every tensor factor of $\Hs{\Qsys_{f,s_f}}$ to the party
  of $f$.
\item \emph{(the \termref{term-squashing-relation}{squashing relation} and the declared model)}
  $\Lambda^{\dagger}(F_v)=\Mdec_v$ for every $v$, exactly, so
  $\Mdec\dsq_{\Lambda}F$; the declared joint model is the \emph{product} of
  the factors' declared models. $\Lambda$ is of the form the artefact
  records for one factor: with $\PinS$ the pinching for the product
  projectors, $\PinS\circ\Lambda=\Lambda$, and the squashed
  analysis may be restricted to states block diagonal for the product sector
  map without excluding a reachable one. The transported set contains every
  product projector $\Pi_s$, every lifted projector
  $\idon{}\otimes\Pi^{f}_{s_f}\otimes\idon{}$ and every lifted transported
  operator $\idon{}\otimes\hat O\otimes\idon{}$, $\hat O$ transported by
  $\mathcal{J}^{f}$, with pull-back the lift of $\hat O$'s pull-back.
\item \emph{(entries keep their constants)} for every factor $f$ and every
  entry $(\hat O_k,\bowtie_k,c_k)\in\CList^{f}$, the certificate
  $\Lambda_f^{\dagger}(\hat O_k)\bowtie_kc_k\,\idon{\Bsys_f}$ which
  $\mathcal{J}^{f}$ discharges at $A'_f$ is a certificate on that factor's
  channel in the sense of Theorem~\ref{thm:licence}\,(v). So the joint list
  $\CList\coloneqq\bigsqcup_f\{(\hat O_k,\bowtie_k,c_k;f)\}$, each entry
  scoped to its factor and carrying the same operator, the same relation and
  the same constant, is licensed with its entries read party-locally
  (Definition~\ref{def:licensed}), or as scalars on the lifted operators
  $\idon{}\otimes\hat O_k\otimes\idon{}$, or any mixture; every squashed
  \exttermref{../foundations/contract.pdf}{term-attack-sequence}{attack sequence} against $\Mdec$ is admissible for it and reproduces
  the protocol's output state; and the lifted operators are block diagonal
  for the product sector map with entries that are the factor's emitted
  entries and zeros and ones. No constant is re-derived and none is rounded.
\item \emph{(class-marginal identities)} with classes
  $\Vset_{t}=\prod_f\Vset^{f}_{t_f}$ and weights $q_t=\prod_fq^{f}_{t_f}$,
  $\sum_{v\in\Vset_t}\Mdec_v=q_t\,\idon{\Bsys}$ holds for every $t$ if every
  factor's \termref{term-class-marginal-identity}{class-marginal identity} is discharged. Where it is
  discharged for the factors in a subset $\mathbb{G}\subseteq\mathbb{F}$ only,
  the identity holds for the classes that are full in every factor outside
  $\mathbb{G}$:
  $\sum_{v:\,v_f\in\Vset^{f}_{t_f}\,(f\in\mathbb{G})}\Mdec_v
  =\prod_{f\in\mathbb{G}}q^{f}_{t_f}\,\idon{\Bsys}$. What is emitted for it
  is the expectation-form
  entry on the pulled-back class sum and never an operator identity on
  $\Qsys$. The weight $q_t$ is a product of declared numbers and is formed
  exactly by whoever reads it; it is not emitted as a number of this module,
  and it could not be: each $q^{f}_{t_f}$ lies on its factor's declared
  lattice (Definition~\ref{def:exactadm}\,(b),
  Definition~\ref{def:joint-family}\,(a)), and a product of two lattice
  numbers lies in general on a finer lattice only and need not be a binary64
  number, exactly as for the blocks of \eqref{eq:producttarget}.
\item \emph{(what was rounded)} the numbers behind \eqref{eq:producttarget}
  are the factors' emitted numbers and nothing else. Every displacement,
  kept or moved class and rounding direction is a factor's and is
  recorded with that factor's emission,
  and Theorem~\ref{thm:emit}\,(v) applies to each factor as it stands: no
  bound relates a key length computed from $\Mdec$ alone to one computed
  from the product of the handed models. The product introduces no further
  displacement.
\item \emph{(one factor)} for $|\mathbb{F}|=1$ parts (i), (ii), (iv) and (v)
  are the corresponding statements of Theorem~\ref{thm:emit} for that
  instance, and what is emitted is that instance's artefact, every entry
  untagged. Part (iii) at one factor is Theorem~\ref{thm:licence}\,(v) at one
  factor, which says \emph{more} than (i)--(iv) of that theorem in
  \exttermref{../transmitter/contract.pdf}{term-untrusted-source-mode}{untrusted-source mode} with $d_A>1$ and the same otherwise
  (Definition~\ref{def:licensed}); the untagged entry that is emitted is
  implied by it (Lemma~\ref{lem:party-local}\,(iii)), so the one-factor
  emission asserts nothing beyond Theorem~\ref{thm:licence}\,(i)--(iv).
\end{enumerate}
These statements are about the factors' instance links and their emitted
families.
\end{theorem}
\claimmeta{\metalabel{Uses}\mref{Lem.}{lem:sep}, \mref{Prop.}{prop:pinching}, \mref{Cor.}{cor:blockdiag}, \mref{Prop.}{prop:tensor}, \mref{Def.}{def:licensed}, \mref{Lem.}{lem:party-local}, \mref{Thm.}{thm:licence}, \mref{Def.}{def:reduction-instance}, \mref{Def.}{def:emitnorm}, \mref{Thm.}{thm:emit}
\metasep \metalabel{Used by}\mref{Req.}{req:party-local-tag}, \mref{Req.}{req:factorised-emission}, \mref{Cor.}{cor:normalized-flat}, \extref{../protocol-compiler/contract.pdf}{protocol-compiler/req:link-row-provenance}
\metasep newly derived
\metasep implemented (effective: accepted)
\metasep \metalabel{Code}\coderef{https://github.com/fraunhofersg/certiqs-key/blob/v0.1.0/src/certiqs_key/reductions/factors.py}{certiqs_key.reductions.factors:reduce_factors}}

\begin{proof}
(i) By Theorem~\ref{thm:emit}\,(ii) each $F^{f}_{v_f}$ is Hermitian and
positive semidefinite with $\sum_{v_f}F^{f}_{v_f}=\idon{\Qsys_f}$, exactly.
A tensor product of Hermitian operators is Hermitian and of positive
semidefinite operators positive semidefinite, and
$\sum_{v\in\Vset}\bigotimes_fF^{f}_{v_f}
=\bigotimes_f\sum_{v_f}F^{f}_{v_f}=\idon{\Qsys}$ by multilinearity. These
are identities between polynomials in the emitted entries, which are dyadic
rationals. Each $F^{f}_{v_f}=\bigoplus_{s_f}F^{f,(s_f)}_{v_f}$ is block
diagonal for $\{\Pi^{f}_{s_f}\}$, so
$\bigotimes_fF^{f}_{v_f}=\bigoplus_{s}\bigotimes_fF^{f,(s_f)}_{v_f}$ along
$\Hs{\Qsys}=\bigoplus_s\bigotimes_f\Hs{\Qsys_{f,s_f}}$, the distributive law
for tensor products over direct sums; the stated order of the index set is a
convention for listing the summands. The party partition is each factor's
own (Definition~\ref{def:reduction-instance}\,(I2)) with its party's label
in place of the default one.

(ii) $\Lambda^{\dagger}(\bigotimes_fF^{f}_{v_f})
=\bigotimes_f\Lambda_f^{\dagger}(F^{f}_{v_f})=\bigotimes_f{M'}^{\mathrm{ph},f}_{v_f}$ by
Lemma~\ref{lem:sep}\,(ii) and Theorem~\ref{thm:emit}\,(iii) for each factor;
this is Proposition~\ref{prop:tensor} iterated. Each
$\Lambda_f=\mathcal{P}_f\circ\Lambda_{0,f}$ with $\mathcal{P}_f$ the
pinching for $\{\Pi^{f}_{s_f}\}$
(Definition~\ref{def:reduction-instance}\,(I3)), and
$\bigotimes_f\mathcal{P}_f$ is the pinching $\PinS$ of
Proposition~\ref{prop:pinching} for the product projectors, because
$\big(\bigotimes_f\mathcal{P}_f\big)(\xi)=\sum_{s}\Pi_s\,\xi\,\Pi_s$ by
multilinearity; a pinching is idempotent, so $\PinS\circ\Lambda=\Lambda$.
Every $F_v$, every $\Pi_s$ and every lifted operator named is block
diagonal for the product projectors by (i), so
Corollary~\ref{cor:blockdiag}\,(iii) applies to $\Lambda$. For a transported
$\hat O$ of factor $f$ with pull-back $O$,
$\Lambda^{\dagger}(\idon{}\otimes\hat O\otimes\idon{})
=\idon{}\otimes\Lambda_f^{\dagger}(\hat O)\otimes\idon{}
=\idon{}\otimes O\otimes\idon{}$ by the same lemma and unitality of the
other factors' adjoints; $\Pi_s$ is a product of lifted projectors and its
pull-back is $\bigotimes_f\Lambda_f^{\dagger}(\Pi^{f}_{s_f})$.

(iii) The certificate of an entry of $\CList^{f}$ is discharged by
$\mathcal{J}^{f}$ at every admissible datum, in particular at $A'_f$
(Definition~\ref{def:reduction-instance}\,(I4), Theorem~\ref{thm:emit}\,(iii)),
and it is a statement about $\Lambda_f$ and $\hat O_k$ alone. That is the
hypothesis of Theorem~\ref{thm:licence}\,(v), whose conclusion is the
admissibility in the party-local set, in the scalar set for the lifted
operators, and for every mixture, together with the equality of output
states and the security implication, the relation $\Mdec\dsq_{\Lambda}F$
being (ii). The lifted operator is a Kronecker product of $\hat O_k$ with
identity matrices, so its entries are entries of $\hat O_k$, zeros and ones,
and it is block diagonal because $\hat O_k$ is for $\{\Pi^{f}_{s_f}\}$.

(iv) $\sum_{v\in\Vset_t}\Mdec_v
=\bigotimes_f\sum_{v_f\in\Vset^{f}_{t_f}}{M'}^{\mathrm{ph},f}_{v_f}$. Where factor $f$
discharges its identity the $f$-th factor of this product is
$q^{f}_{t_f}\idon{\Bsys_f}$ (Theorem~\ref{thm:emit}\,(iii)); summing over all
$t_f$ for $f\notin\mathbb{G}$ replaces those factors by $\idon{\Bsys_f}$, by
completeness of ${M'}^{\mathrm{ph},f}$. The remark on the emitted form is
Definition~\ref{def:reduction-instance}, fourth consequence, applied to the
pulled-back class sum.

(v) \eqref{eq:producttarget} is a function of the factors' emitted numbers,
and by hypothesis nothing else is emitted; so there is no number to which a
further displacement could attach.

(vi) With one factor, \eqref{eq:producttarget} is $F_v=F^{f}_v$, the product
sector map is the factor's and $\Lambda=\Lambda_f$, so (i), (ii), (iv) and
(v) are Theorem~\ref{thm:emit}'s. For (iii): by
Definition~\ref{def:licensed} the party-local reading at one factor is the
entry of the source mode in \exttermref{../transmitter/contract.pdf}{term-trusted-source-mode}{trusted-source mode} and at $d_A=1$, and in
untrusted-source mode with $d_A>1$ it is an operator inequality on $A$ whose
trace is the scalar entry, Lemma~\ref{lem:party-local}\,(iii); the emitted
entry is the untagged one in every case, which
Theorem~\ref{thm:licence}\,(i)--(iv) license on the same certificate.
\end{proof}

% claim reductions/cor:emitted-identity
\claimsummary{reductions/cor:emitted-identity}{For the identity instance the normalizer moves nothing, so the emitted family is the handed datum itself, squashed by the sector pinching with no error term on an exact link. A requested entry's certificate is the operator inequality itself, decided exactly on the emitted numbers.}
\claimbookmark{reductions/cor:emitted-identity}{Corollary}{cor:emitted-identity}{The emitted family of the identity instance}
\claimleadin{%
The identity instance is the one a protocol with an already-finite receiver
would use; this corollary is Theorem~\ref{thm:emit} read at it.
}
\begin{corollary}[The emitted family of the identity instance]
\label{cor:emitted-identity}
Let the declared data pass every refusal rule of this module's interface, let
$\mathcal{J}_{\mathrm{id}}$ be the identity instance of
Corollary~\ref{cor:identity-instance} with handed datum $A$, and let $A'$ be
the datum the normalizer of Definition~\ref{def:emitnorm} declares. Then
Theorem~\ref{thm:emit} holds of $\mathcal{J}_{\mathrm{id}}$; the normalizer
moves nothing, $A'=A$ and $\Delta_s=0$ for every data sector $s$; and, for
every $v\in\Vset$,
\begin{equation}\label{eq:emitted-identity}
  \Lambda^{\dagger}(F_v)\;=\;\PinS(F_v)\;=\;F_v\;=\;\Mdec_v
  \;=\;\bigoplus_{s\in\SecSet_0}A'^{(s)}_v
  \qquad\text{exactly, on }\Bsys=\Qsys ,
\end{equation}
so the emitted family \emph{is} the declared datum, which is the handed
datum, exactly admissible and decided so before emission; the emitted
constraint list is \emph{empty where the request carries no entry} --- this
instance's own list being empty in any case
(Corollary~\ref{cor:identity-instance}\,(I4)) --- and where the request
carries entries it carries, on request, the class-marginal entries
$\big(\sum_{v\in\Vset_t}F_v,=,q_t\big)$, which the instance discharges, and
a further requested entry $(\hat O,\bowtie,c)$ of the form
\eqref{eq:instance-entry}, which Corollary~\ref{cor:entries}\,(iii) emits
exactly where $\hat O\bowtie c\,\idon{\Bsys}$ holds at the constant $c$ it
emits, this instance carrying none of its own: $\Lambda^{\dagger}(\hat O)=\hat O$, so the
certificate \eqref{eq:entrygeneral} of such an entry is that operator
inequality on the emitted numbers, decided exactly by finitely many field
operations and sign comparisons, a decided certificate in the sense of
Corollary~\ref{cor:entries}\,(iii) for every operator of that form; its
default carried operators --- here the sector projectors and nothing else,
the transported set being $\{\Pi_s\}_{s\in\SecSet_0}$ --- are emitted with
the transported set at $c=0$ with $\succeq$ in either case
(Corollary~\ref{cor:entries}\,(i)); no scalar value is emitted and no bound
of this module's own is rounded; the chain is the identity link with
$\textsf{status}=\textsf{exact}$. These statements are about the instance's
link and its emitted family.
The instance is selectable only once a declared input carries an
already-finite POVM.
\end{corollary}
\claimmeta{\metalabel{Uses}\mref{Prop.}{prop:pinching}, \mref{Thm.}{thm:chain}, \mref{Def.}{def:exactadm}, \mref{Lem.}{lem:exactpsd}, \mref{Def.}{def:reduction-instance}, \mref{Cor.}{cor:identity-instance}, \mref{Def.}{def:emitnorm}, \mref{Thm.}{thm:emit}, \mref{Cor.}{cor:entries}
\metasep newly derived
\metasep accepted}

\begin{proof}
Corollary~\ref{cor:identity-instance} makes $\mathcal{J}_{\mathrm{id}}$ a
reduction instance, so Theorem~\ref{thm:emit} applies to it with the data
passing every refusal rule. Its handed datum is the declared family, every
block of which is decided exactly admissible on the declared lattice on
entry, so every entry satisfies \eqref{eq:latinv}, and the first of the
properties of Definition~\ref{def:emitnorm} gives $A'=A$ and every
$\Delta_s=0$. Part (iii) of that theorem at this instance reads
$\Lambda^{\dagger}(F_v)=M_v[A']=\bigoplus_sA'^{(s)}_v$; and $\Lambda=\PinS$ is
self-adjoint with $\PinS(F_v)=F_v$ for the block-diagonal $F_v$
(Proposition~\ref{prop:pinching}), which gives the two middle equalities of
\eqref{eq:emitted-identity}. That the emitted family is exactly admissible
and decided so is part (ii); that it is the declared datum is (I2) of the
instance, there being no scalar sector. The default carried operators and
the entries are Corollary~\ref{cor:entries}\,(i) and, the identity instance
discharging the class-marginal identity, (ii). For a further requested
entry, $\hat O$ is a real combination of the blocks $\Pi_sF_v\Pi_s$ and of
sector projectors, hence block diagonal for $\SecSet_0$ and Hermitian, so
$\Lambda^{\dagger}(\hat O)=\PinS(\hat O)=\hat O$ (Proposition~\ref{prop:pinching});
its entries are formed from emitted lattice entries and declared lattice
coefficients by finitely many field operations, so they lie in $\Fldi$; and
the certificate \eqref{eq:entrygeneral} is then $\hat O-c\,\idon{\Bsys}\succeq0$,
$c\,\idon{\Bsys}-\hat O\succeq0$ or an entrywise equality, the first two
decided by \eqref{eq:psdtest} of Lemma~\ref{lem:exactpsd} and the last by
comparing entries. The instance's (H) being empty, that condition is a decided
certificate for every operator of the form \eqref{eq:instance-entry}, and
Corollary~\ref{cor:entries}\,(iii) emits the entry exactly where it holds.
Parts (ii) and (iii) are statements
about a request that carries entries; where the request carries none the
emitted list is empty by that corollary's rule, whatever the instance's own
list contains --- and this instance's is empty in any case
(Corollary~\ref{cor:identity-instance}\,(I4)) --- so the artefact then
carries no entry and no licence, which is what the last sentence of that
corollary says of its chain. With no scalar sector the rounding clause of
Theorem~\ref{thm:emit}\,(iv) has nothing of this module's to round, and only
a caller's constant can be. The chain and its status are the last sentence
of Corollary~\ref{cor:identity-instance} with Theorem~\ref{thm:chain}\,(iv),
now at the declared datum. The last sentence is
Example~\ref{ex:identity}.
\end{proof}

% claim reductions/rem:emitclaim
\claimsummary{reductions/rem:emitclaim}{Where the normalizer moves the datum, the claims are about the declared model, and nothing proved relates it to the handed one. A run can evaluate off the snap on the reference companion; elsewhere a disclosed premise relates them.}
\claimbookmark{reductions/rem:emitclaim}{Remark}{rem:emitclaim}{What the reduction's claim is about after an emission snap}
\begin{remark}[What the reduction's claim is about after an emission snap]
\label{rem:emitclaim}
Where the exact construction stays on the declared lattice --- every
$\Delta_s=0$ --- there is nothing to say: the emitted target is the
constructed one, the defining identity is about the handed model itself, and
what the artefact asserts is what the instance's existence theorem proves
(Definition~\ref{def:reduction-instance}\,(I3)). \emph{Where some
$\Delta_s>0$ the emitted numbers are a declaration and not the constructed
ones}, and what the artefact asserts is then this. There is a datum $A'$ ---
exhibited rather than posited, and carried entry by entry in the emitted
target's data sectors --- which is exactly admissible, class complete at the
declared weights, and differs from the handed datum by at most $\Delta_s$
entrywise per data sector. \emph{It is the datum of a whole declared model}
$\Mdec=M[A']$ \eqref{eq:declmodel}: the handed model with the datum replaced
as the instance's model map defines it, every object the hypotheses (H) are
about --- an above-cutoff family, a register POVM, a further branch --- being
the declared one, which the normalizer never reads. \emph{The emitted channel
$\Lambda$ squashes that model, exactly and with no error term} ---
$\Lambda^{\dagger}(F_v)=\Mdec_v$ \eqref{eq:emitidentity}, an identity on
$\Bsys$ and not on the data sectors alone; every hypothesis (H) holds of it by
the stability clause; every certificate of the instance is discharged for it;
and every claim of this contract that reads a declared number reads it.
\emph{The physical content of the artefact alone is therefore a statement
about $\Mdec$, of which the normalizer moved the data sectors and nothing
else}, and $A'$ stands one modelling step further from the supplier's own
output than the instance's declaration already stood. Both steps are
recorded, the first in the instance's record of its snap and the second in
this module's emission record (Requirement~\ref{req:emission-record}).

\emph{Nothing proved relates $\Mdec$ to the handed model, and distance is not
such a relation.} The size $\Delta_s$ of the move bounds nothing that
matters to a key rate: the consuming contract shows that two exactly
admissible families within any declared distance can have rate bounds far
apart (\xref{meat-core/prop:optol-nobound}), so no bound relates a key length
computed from $\Mdec$ alone to one computed from the handed model, and none
is asserted. The snap keeps every row and column that the handed datum holds
at zero (Definition~\ref{def:emitnorm}, fifth property), save on a class
snapped by the last resort, which the record flags, and it may move support
along other directions. A run therefore does not rest on the declared model
where it can avoid it: where it forms the number its key length is formed
from on the reference companion (Corollary~\ref{cor:reference-companion}),
which the module builds with no normalizer from the handed reference, or from
the instance's own snapped blocks where none is handed, every data sector is \exttermref{../foundations/contract.pdf}{term-evaluated-off-its-snap}{evaluated off its snap}
(\xref{foundations/def:evaluated-off-snap}), and the emitted target serves
only the search. Elsewhere, and so also where the key length is formed from
a number certified on the snapped declaration in place of the companion's
(\xref{meat-core/frank-wolfe/req:certified-evaluation}\,(i)), the run rests on
\xref{foundations/ass:snapped-measurement}\,(ii), the premise for a
producer's snap of its own emission: what the run concludes about the device
from the emitted target holds with $A'$ in place of $A$ on the sectors
concerned. The artefact being an exact statement about $\Mdec$, the premise
holds in particular where the device measures $\Mdec$ there. The record
names the moved sectors (Requirement~\ref{req:emission-record}).

\emph{Two things this does not license.} It does not license repairing what
this module was handed: the normalizer runs on the datum an instance computed,
and on nothing else (Definition~\ref{def:emitnorm}, fourth property). And it
does not license emitting $\textsf{status}=\textsf{exact}$ about a family the
identity does not hold of: the value is truthful precisely because the
identity is re-established for $\Mdec$, whose data sectors are $A'$, rather
than asserted of a construction the artefact does not carry.
\end{remark}
\claimmeta{\metalabel{Uses}\mref{Def.}{def:reduction-instance}, \mref{Def.}{def:emitnorm}, \mref{Thm.}{thm:emit}
\metasep \metalabel{Used by}\mref{Rem.}{rem:claimscope}
\metasep newly derived
\metasep accepted}

% claim reductions/rem:claimscope
\claimsummary{reductions/rem:claimscope}{The squashing, transfer, calculus and core instance results read no declared number; a submodule's instance reads declared data decided exactly, the licence exactly Hermitian operators and declared constants. Where the normalizer acts, every statement is read on the declared model.}
\claimbookmark{reductions/rem:claimscope}{Remark}{rem:claimscope}{The claims that read a declared number}
\begin{remark}[The claims that read a declared number]\label{rem:claimscope}
\emph{The \termref{term-squashing-relation}{squashing relation} and its soundness, the factorization and
transfer theorems, the calculus of the pre-order, the branching rule, the
flags-enough corollary, register-conditioned squashing and the obstruction
read no declared number} and assert nothing about the emitted artefact: each
is a statement about an arbitrary squashing relation, an arbitrary
instrument, an arbitrary register-conditioned mixture, or an explicit
counterexample. \emph{The instance theorems of the core} --- the identity,
flag, branching and register instances
(Corollaries~\ref{cor:identity-instance}, \ref{cor:flag-instance}
and~\ref{cor:register-instance}, Theorem~\ref{thm:instrument-instance}) ---
are statements about arbitrary admissible data and arbitrary branch instances
and read no declared number either. \emph{A submodule's instance} builds its
datum out of declared blocks and declared classical data; its hypotheses are
conditions on a POVM family and on declared scalars and matrices, and this
module's refusal rules decide every one of those conditions \emph{exactly},
so it applies to the declared data verbatim. \emph{The attack-set licence}
(Theorem~\ref{thm:licence}) needs its emitted constraint operators exactly
Hermitian and its constants declared: the first holds by construction and is
decided by the refusal rule that says so, the second by
Assumption~\ref{ass:constdecl}. \emph{The chain theorem}
(Theorem~\ref{thm:chain}) asserts $\textsf{status}=\textsf{exact}$ of the
artefact, which is true of the numbers emitted
(Theorem~\ref{thm:emit}\,(iii)). \emph{The emission theorem}
(Theorem~\ref{thm:emit}) is the claim about all of the above, stated once over
the reduction instance (Definition~\ref{def:reduction-instance}).

\emph{One conditioning survives and it is the only one}: where the emission
normalizer acts, each of those statements is read on the declared model of
Definition~\ref{def:emitnorm} --- on $\Bsys$, the declared model
$\Mdec=M[A']$; in the data sectors, the declared datum $A'$ --- and not on
the handed datum, and Remark~\ref{rem:emitclaim} is what that means.
\emph{Every residual is recorded}: the displacement and the rounding
directions are (Requirement~\ref{req:emission-record}), and the consumer's
key length is Theorem~\ref{thm:emit}\,(v) --- not established, and no refusal
rule can establish it.
\end{remark}
\claimmeta{\metalabel{Uses}\mref{Def.}{def:reduction-instance}, \mref{Def.}{def:emitnorm}, \mref{Thm.}{thm:emit}, \mref{Rem.}{rem:emitclaim}
\metasep \metalabel{Assumes}\aref{A-CONST}{ass:constdecl}
\metasep newly derived
\metasep accepted}

% claim reductions/rem:exactboundary
\claimsummary{reductions/rem:exactboundary}{Data without a finite presentation, the objects of an instance's hypotheses, and combinatorial data offer no numerical predicate: the first stay hypotheses, the second are settled by inspection. The normalizer never touches the first, so the hypotheses hold of the declared model.}
\claimbookmark{reductions/rem:exactboundary}{Remark}{rem:exactboundary}{The boundary of the exact decisions}
\begin{remark}[The boundary of the exact decisions]\label{rem:exactboundary}
Every operator predicate this module checks is decided exactly. \emph{Two kinds
of declared datum carry no such predicate.}

\emph{First: data with no declared finite presentation.} The hypotheses (H) of
a reduction instance (Definition~\ref{def:reduction-instance}\,(I1)) are about
exactly these: an operator family on a part of $\Hs{\Bsys}$ that no finite
list of matrices describes --- $\Hs{\Bsys}$ being separable and possibly
infinite-dimensional (Definition~\ref{def:system-setting}) --- such as the blocks of a
family above a photon-number cutoff; a POVM on a register whose dimension is
nowhere declared or bounded; the POVM of an uncharacterised branch, which the
instance that files it deliberately declines to constrain
(Corollary~\ref{cor:flag-instance}). \emph{None of these has a finite matrix
presentation anywhere in this contract, so none offers anything to decide}:
there is no list of entries for Lemma~\ref{lem:exactpsd} to run on. Their
conditions --- that a family is a POVM, that an above-cutoff block satisfies
a certified inequality, that a register family is complete --- are therefore
\emph{hypotheses of the instance}, carried in the ledger and decided by this
module's refusal rules only to the extent that the declaration presents
something to decide. Which objects these are for a given instance is that
instance's to say; of the core's leaf instances only the flag instance has
one, its branch POVM, and the branching and register instances add the
hypothesis that their instrument is a \termref{term-generalized-instrument}{generalized instrument}, and for
the register form that the register family is a POVM
(Theorem~\ref{thm:instrument-instance},
Corollary~\ref{cor:register-instance}).

\emph{Second: data that is combinatorial rather than numerical.} A
coarse-graining map with its totality, surjectivity and fibre conditions; the
alphabets $\Yset$, $\Tset$, $\Vset$ and their declared orders; the sector
identifiers, their dimensions, their supports and their order; a branch index
set; a cutoff; and the \exttermref{../foundations/contract.pdf}{term-lattice-exponent}{lattice exponent} $\kappa$ itself. \emph{None of these is a
number a condition could be decided about}, and a predicate on them is settled
by finite inspection --- which is the same reason the \exttermref{../foundations/contract.pdf}{term-exact-island}{exact island}
(\xref{foundations/def:exact-island}) is drawn where it is.

\emph{Two consequences.} First, the objects of the first kind are also the
ones the emission normalizer never touches: it acts on the data sectors and on
nothing else (Definition~\ref{def:emitnorm}, second property), so none of them
enters any $\Delta_s$ and none is repaired --- which is what makes every
hypothesis (H) hold of the declared model by the stability clause, and with it
Corollary~\ref{cor:entries}\,(ii)'s identity available where the instance
discharges it. Second, \emph{extending an exact decision to such an object
would first require declaring a finite representation for it}, and for an
uncharacterised branch that means constraining the very model the instance
leaves free; that is a widening of the instance's model rather than a change
of number model. The declared classical data an instance's route reads --- a
stochastic matrix and its descended form, the \termref{term-class-weights}{class weights}, the
\termref{term-certified-weights}{certified weights} --- \emph{is} finitely presented, so it is decided here,
exactly.
\end{remark}
\claimmeta{\metalabel{Uses}\mref{Def.}{def:system-setting}, \mref{Lem.}{lem:exactpsd}, \mref{Def.}{def:reduction-instance}, \mref{Def.}{def:emitnorm}, \mref{Cor.}{cor:entries}
\metasep newly derived
\metasep accepted}

% claim reductions/rem:sound-and-atk
\claimsummary{reductions/rem:sound-and-atk}{Every entry this module emits holds on every reachable squashed state, as a conclusion of \hyperlink{claim-reductions/thm:chain}{Theorem~\ref*{thm:chain}}, \hyperlink{claim-reductions/cor:entries}{Corollary~\ref*{cor:entries}} and \hyperlink{claim-reductions/thm:licence}{Theorem~\ref*{thm:licence}} and not as an assumption. Inside the engine's security chain the whole list is covered by the stronger \hyperref{../protocol-compiler/contract.pdf}{}{claim-protocol-compiler/ass:atk}{A-ATK} through the engine's range lemma.}
\claimbookmark{reductions/rem:sound-and-atk}{Remark}{rem:sound-and-atk}{Soundness of the emitted constraint list}
\begin{remark}[Soundness of the emitted constraint list]\label{rem:sound-and-atk}
Every entry of the constraint list $\CList$ this module emits is sound for
the \termref{term-reduction-chain}{reduction chain} that produced the declaration, in the sense of
Definition~\ref{def:sound}: every reachable squashed state satisfies it. For
the entries this module emits this is a conclusion and not an assumption:
the entries are the ones the chain theorem, the emission corollary and the
licence theorem put into $\CList$ (Theorem~\ref{thm:chain},
Corollary~\ref{cor:entries}, Theorem~\ref{thm:licence}), under their
hypotheses. \emph{Inside} the consuming engine's security chain, soundness of
the whole list, entries a declaration brings itself included, is implied by
the stronger A-ATK (\xref{protocol-compiler/ass:atk}) through the engine's
range lemma (\xref{meat-core/lem:range}\,(i)), which is declared there.
Of an entry this module did not emit, nothing here asserts soundness outside
that chain.
\end{remark}
\claimmeta{newly derived
\metasep accepted}

\section{The reference companion}\label{group-reductions/cor:reference-companion}
\sectionintro{Where the declaring instance holds a reference of the measurement it snapped, the same instance run on the reference in place of the snap is again a \termref{term-reduction-instance}{reduction instance} with the emitted alphabet, sectors and \termref{term-scalar-value}{scalar values}: its target is a POVM that the reference model squashes to exactly, with its own \termref{term-certified-entries}{certified entries}, known through enclosures formed from the reference's (\hyperlink{claim-reductions/cor:reference-companion}{Corollary~\ref*{cor:reference-companion}}). This needs only that the instance's datum, formed from any family with the properties of a reference, is admissible, which holds for the identity instance, passes through branching, and holds for the flat instance. An entry built from the reference is emitted in binary64 by rounding its operator and moving its constant outward by a bound on the rounding, which keeps a certificate and enlarges the set the entry cuts out (\hyperlink{claim-reductions/lem:row-relaxation}{Lemma~\ref*{lem:row-relaxation}}). Where the instance forms its datum as a non-negative combination of the blocks, as the identity and the flat instance do, factor enclosures of the reference's blocks pass to the companion's datum and target, columns side by side and cores scaled (\hyperlink{claim-reductions/cor:companion-factor-enclosure}{Corollary~\ref*{cor:companion-factor-enclosure}}).}
\sectionimports{\extref{../foundations/contract.pdf}{foundations/def:evaluated-off-snap}, \extref{../foundations/contract.pdf}{foundations/def:exact-reference}, \extref{../foundations/contract.pdf}{foundations/def:factor-enclosure}, \extref{../foundations/contract.pdf}{foundations/def:snap}, \extref{../receiver/contract.pdf}{receiver/def:complex-enclosure}, \extref{../receiver/contract.pdf}{receiver/lem:enclosure-arithmetic}}

% claim reductions/cor:reference-companion
\claimsummary{reductions/cor:reference-companion}{The same instance, run on the reference in place of the snap, or on its own snapped blocks where none is handed, is a reduction instance of the emitted shape that squashes that model exactly. Its entries are enclosed from the reference, requested ones where their certificates hold there.}
\claimbookmark{reductions/cor:reference-companion}{Corollary}{cor:reference-companion}{The reference companion of a reduction instance}
\claimleadin{%
An emission is computed on the snapped measurement a declaring instance
states, and the normalizer may move the instance's datum once more
(Theorem~\ref{thm:emit}). Where the declaring instance holds a reference of
the measurement and hands it beside its snap, the same construction can be
run on the reference instead: the instance's datum formed from the reference
blocks, with no normalizer. Where it hands none, the same construction run on
the snapped blocks themselves gives the producer's exact construction, the
\termref{term-handed-datum}{handed datum}, which the normalizer has not touched: the snap of the
emission is then bypassed though the instance's own snap is not. Nothing has
to be proved again for that, because
every part of a \termref{term-reduction-instance}{reduction instance} is established at every \termref{term-admissible-datum}{admissible datum};
the one thing to supply is that the datum formed from the reference is
admissible, which is hypothesis (H-REF) below, and the substrate in which the
construction is read, since reference entries are in general irrational.
What results is no declaration a consumer solves in its own right
(Requirement~\ref{req:reference-companion}); its entries, which are formed
from the reference, are emitted in binary64 by
Lemma~\ref{lem:row-relaxation}. Part~(iv) discharges (H-REF) for the identity
instance and passes it through branching; the flat instance discharges it in
Corollary~\ref{cor:flat-companion}. Part~(vi) carries to the companion the
entries a caller requests beside the instance's own: their operators are
formed at the companion's datum by the same formula, and their certificates
are re-established there, since a certificate decided on the emitted
numbers says nothing about the reference; for the operator of
\eqref{eq:two-sector-entry} the flat instance does so in
Corollary~\ref{cor:flat-companion}.
}
\begin{mainclaim}
\begin{corollary}[The reference companion of a reduction instance]\label{cor:reference-companion}
Let $\mathcal{J}$ be a \termref{term-reduction-instance}{reduction instance}
(Definition~\ref{def:reduction-instance}) whose \termref{term-handed-datum}{handed datum} $A$ is
formed, by the rule of its part (I1), from a family $K=\{K^{(b)}_j\}$ of
measurement blocks, over blocks $b$ of a \exttermref{../foundations/contract.pdf}{term-sector-map}{sector map} and outcomes $j$ of an
alphabet (one such family per setting where the blocks are given per
setting), that a declaring instance states as a \exttermref{../foundations/contract.pdf}{term-snap}{snap}
(\xref{foundations/def:snap}) of the measurement of a device. Let
$G=\{G^{(b)}_j\}$ be a family over the same blocks, alphabet and dimensions,
every block of which is Hermitian and positive semidefinite and satisfies
exactly every completeness identity that the declaration of $K$ states of
$K$ on that block. In the use this corollary is for, $G^{(b)}$ is the block
of the \exttermref{../foundations/contract.pdf}{term-handed-reference}{handed reference} (\xref{foundations/def:evaluated-off-snap})
where the declaring instance hands one, and $K^{(b)}$ elsewhere, so that
$G=K$ where it hands none. Suppose
\begin{itemize}[leftmargin=4.2em]
\item[(H-REF)] $\mathcal{J}$ is a reduction instance over the substrate
  $\Fld=\mathbb{R}$ (Definition~\ref{def:substrate}), and the rule of its
  part (I1) forms, from every family with the properties just listed for $G$,
  a datum in $\mathcal{D}_{\mathrm{adm}}$.
\end{itemize}
Let $A^{\mathrm{ref}}$ be the datum that rule forms from $G$ in place of
$K$, and call the instance at $A^{\mathrm{ref}}$ --- the same $\Bsys$, (H),
sectors, supports, \termref{term-scalar-value}{scalar values}, constructor and entry rule, with
$A^{\mathrm{ref}}$ in place of the handed datum --- the
\emph{reference companion} of $\mathcal{J}$ at $G$. Then:
\begin{enumerate}[label=(\roman*),leftmargin=3.0em]
\item \emph{(an admissible target over the emitted shape)}
  $A^{\mathrm{ref}}\in\mathcal{D}_{\mathrm{adm}}$, the reference companion
  is a reduction instance, and its target $\{F_v[A^{\mathrm{ref}}]\}_{v\in\Vset}$
  is a POVM on $\Qsys$, sector-block-diagonal by construction. Its alphabet
  with classes and weights, its sector map (the sector list, its order, each
  sector's dimension and kind), its scalar values and its party partition are
  those of the target the instance has at every admissible datum, the
  \termref{term-declared-datum}{declared datum} the normalizer produces included, and it differs
  from such a target on the \termref{term-data-sector}{data sectors} alone, where it
  carries $A^{\mathrm{ref}}$. Where $G=K$, $A^{\mathrm{ref}}=A$: the
  reference companion is $\mathcal{J}$ at its handed datum;
\item \emph{(an exact squash of the companion's model)}
  $\big(M[A^{\mathrm{ref}}];O[A^{\mathrm{ref}}]\big)\dsq_{\Lambda[A^{\mathrm{ref}}]}
  \big(F[A^{\mathrm{ref}}];\hat O[A^{\mathrm{ref}}]\big)$, so that
  $\Lambda[A^{\mathrm{ref}}]^{\dagger}(F_v[A^{\mathrm{ref}}])=M_v[A^{\mathrm{ref}}]$
  exactly for every $v$; every certificate of the entry list
  $\CList[A^{\mathrm{ref}}]$ is discharged, and the \termref{term-class-marginal-identity}{class-marginal identity}
  \eqref{eq:instance-classmarginal} holds at $A^{\mathrm{ref}}$ where the
  instance's flag says it is discharged; and Theorem~\ref{thm:transfer} and
  Theorem~\ref{thm:chain} hold of the one-link \termref{term-receiver-segment}{receiver segment}
  (Definition~\ref{def:link}) whose \termref{term-source-family}{source family} is
  $(\Bsys,M[A^{\mathrm{ref}}])$, whose target is
  $(\Qsys,F[A^{\mathrm{ref}}])$ and whose \termref{term-channel-witness}{channel witness} is
  $\Lambda[A^{\mathrm{ref}}]$, with $\textsf{status}=\textsf{exact}$;
\item \emph{(the companion is known through enclosures)} let every block of
  $G$ be read with an \exttermref{../foundations/contract.pdf}{term-entrywise-enclosure}{entrywise enclosure}
  (\xref{foundations/def:exact-reference}), with radius zero on every block
  where $G^{(b)}=K^{(b)}$. Where the rule of (I1) forms each entry of the
  datum as a finite expression in entries of the blocks and in numbers the
  instance was handed, and every reciprocal node of that expression receives
  an enclosure by part (iv) of \xref{receiver/lem:enclosure-arithmetic}, as
  it does where the expression has none, the rules of that lemma applied to
  that expression give a \exttermref{../receiver/contract.pdf}{term-complex-enclosure}{complex enclosure} of every element of
  $A^{\mathrm{ref}}$ whose centre is formed by the same expression from the
  centres and whose radius is zero at every entry whose inputs are all
  enclosed with radius zero; part (v) of that lemma, applied element by
  element, makes it an entrywise enclosure of $A^{\mathrm{ref}}$, whose
  centre is the Hermitian part of that centre and whose radius is zero at an
  entry $(i,l)$ where the radii at $(i,l)$ and $(l,i)$ both are. Every entry operator
  $\hat O_k[A^{\mathrm{ref}}]$ of \eqref{eq:instance-entry}, a combination
  of the blocks $\Pi_sF_v[A^{\mathrm{ref}}]\Pi_s$ with real coefficients
  $\alpha_{k,v,s}$ plus a combination $Z_k$ of sector projectors, is then
  enclosed entrywise by the same combination of the centres, with radius
  $\sum_{v}|\alpha_{k,v,s}|\,R^{(s)}_{v,il}$ at the entry $(i,l)$ of each
  data sector $s$, where $(C,R)$ is any entrywise enclosure of the datum and
  $R^{(s)}_{v,il}$ is the radius of its element $v$ at that entry, and
  radius zero on every \termref{term-scalar-sector}{scalar sector};
\item \emph{(where (H-REF) holds)} (H-REF) holds for the identity instance of
  Corollary~\ref{cor:identity-instance}; and it holds for the branching
  instance of Theorem~\ref{thm:instrument-instance}, and for the register
  instance of Corollary~\ref{cor:register-instance}, whenever it holds for
  every branch instance whose datum is formed from blocks of $K$, the data of
  the other branches being formed from no block of $K$;
\item \emph{(what the companion rests on, and only there)} no part of
  (i)--(iv) rests on \xref{foundations/ass:snapped-measurement}. Let $B$ be
  the blocks on which $G^{(b)}=K^{(b)}$ that the record of the run names as
  resting on that assumption. Where $B$ is empty this part asserts nothing.
  Where it is not, let the assumption hold: then on every block in $B$ the
  block of the measurement the device performs is the snapped block
  $K^{(b)}$, with the entries recorded for the snap, in place of the block
  of the declaring instance's model, so that the companion is formed there
  from the block the device then measures; and no claim relates a key length
  computed from the companion to that model on the blocks in $B$;
\item \emph{(entries requested beside the instance's own)} let
  $(\hat O,\bowtie,c)$ be an entry emitted at a request that carries entries
  (Corollary~\ref{cor:entries}), written in the form
  \eqref{eq:instance-entry} with coefficients $\alpha_{v,s}$ and a
  combination $Z$ of sector projectors,
  $\hat O=\sum_{v\in\Vset}\sum_{s\in\SecSet}\alpha_{v,s}\Pi_sF_v[A']\Pi_s+Z$
  at the \termref{term-declared-datum}{declared datum} $A'$, and call
  $(\hat O^{\mathrm{ref}},\bowtie,c)$, with
  \begin{equation}\label{eq:companion-form}
    \hat O^{\mathrm{ref}}\;\coloneqq\;\sum_{v\in\Vset}\sum_{s\in\SecSet}
      \alpha_{v,s}\,\Pi_sF_v[A^{\mathrm{ref}}]\Pi_s\;+\;Z ,
  \end{equation}
  its \emph{companion form} at that writing. Part~(iii) encloses
  $\hat O^{\mathrm{ref}}$. Where $\hat O$ is a sum of sector projectors
  (Corollary~\ref{cor:entries}\,(i)), written with every $\alpha_{v,s}=0$,
  $\hat O^{\mathrm{ref}}=\hat O$ and
  $\Lambda[A^{\mathrm{ref}}]^{\dagger}(\hat O)\succeq0$. Where it is the
  class-marginal entry $\big(\sum_{v\in\Vset_t}F_v,=,q_t\big)$ of
  Corollary~\ref{cor:entries}\,(ii), written with $\alpha_{v,s}=1$ for
  $v\in\Vset_t$, $0$ otherwise, and $Z=0$, $\hat O^{\mathrm{ref}}=\hat O$,
  and where the instance discharges \eqref{eq:instance-classmarginal},
  $\Lambda[A^{\mathrm{ref}}]^{\dagger}(\hat O^{\mathrm{ref}})=q_t\,\idon{\Bsys}$
  exactly. For any other entry this part asserts no certificate at
  $A^{\mathrm{ref}}$. In every case, where
  $\Lambda[A^{\mathrm{ref}}]^{\dagger}(\hat O^{\mathrm{ref}})\bowtie
  c\,\idon{\Bsys}$ holds,
  $\big(M[A^{\mathrm{ref}}];O[A^{\mathrm{ref}}],
  \Lambda[A^{\mathrm{ref}}]^{\dagger}(\hat O^{\mathrm{ref}})\big)
  \dsq_{\Lambda[A^{\mathrm{ref}}]}
  \big(F[A^{\mathrm{ref}}];\hat O[A^{\mathrm{ref}}],\hat O^{\mathrm{ref}}\big)$,
  and the companion form is an entry carrying the spectral certificate
  \eqref{eq:certificate} at its constant.
\end{enumerate}
\end{corollary}
\end{mainclaim}
\claimmeta{\metalabel{Uses}\mref{Def.}{def:squashing}, \mref{Thm.}{thm:transfer}, \mref{Cor.}{cor:blockdiag}, \mref{Thm.}{thm:licence}, \mref{Def.}{def:link}, \mref{Def.}{def:chain}, \mref{Thm.}{thm:chain}, \mref{Def.}{def:substrate}, \mref{Def.}{def:reduction-instance}, \mref{Cor.}{cor:identity-instance}, \mref{Thm.}{thm:instrument-instance}, \mref{Cor.}{cor:register-instance}, \mref{Def.}{def:emitnorm}, \mref{Cor.}{cor:entries}, \extref{../foundations/contract.pdf}{foundations/def:evaluated-off-snap}, \extref{../foundations/contract.pdf}{foundations/def:exact-reference}, \extref{../foundations/contract.pdf}{foundations/def:snap}, \extref{../receiver/contract.pdf}{receiver/def:complex-enclosure}, \extref{../receiver/contract.pdf}{receiver/lem:enclosure-arithmetic}
\metasep \metalabel{Assumes}\extaref{../foundations/contract.pdf}{foundations/ass:snapped-measurement}{A-SNAP}
\metasep \metalabel{Used by}\mref{Cor.}{cor:companion-factor-enclosure}, \mref{Req.}{req:reference-companion}, \mref{Cor.}{cor:flat-companion}, \mref{Cor.}{cor:normalized-flat}, \extref{../meat-core/contract.pdf}{meat-core/cor:reference-companion-factors}, \extref{../meat-core/contract.pdf}{meat-core/prop:reference-companion}, \extref{../meat-core/contract.pdf}{meat-core/req:companion-admission}, \extref{../protocol-compiler/contract.pdf}{protocol-compiler/req:reference-companion-compiled}
\metasep newly derived
\metasep accepted}

\begin{proof}
\emph{Idea.} The \termref{term-stability-clause}{stability clause} of Definition~\ref{def:reduction-instance}
was stated so that the instance at any admissible datum is again an
instance; (H-REF) supplies an admissible datum formed from $G$, and every
conclusion is a part of the instance read at that datum.

(i) The family $G$ has the properties (H-REF) lists, so
$A^{\mathrm{ref}}\in\mathcal{D}_{\mathrm{adm}}$, the set of admissible data
over the substrate $\Fld=\mathbb{R}$, which Definition~\ref{def:substrate}
admits: it leaves the field open between $\mathbb{Q}$ and $\mathbb{R}$, and
entries of a reference may be irrational. The definition's closing sentence
says that the instance at an admissible datum is a reduction instance
whenever $\mathcal{J}$ is, by the stability clause: (H) reads no entry of
the datum, so it holds of $M[A^{\mathrm{ref}}]$ as it holds of the \termref{term-handed-model}{handed
model}, and (I2)--(I4) are established at every admissible datum. By the
first consequence of Definition~\ref{def:reduction-instance},
$\{F_v[A^{\mathrm{ref}}]\}$ is a POVM on $\Qsys$, its data sectors carrying
$A^{\mathrm{ref}}$ and its scalar sectors the scalar values; it is assigned
one sector block at a time by \eqref{eq:instance-target}, so it is
sector-block-diagonal by construction. The alphabet, the classes and
weights, the sector list with its order, dimensions and kinds, the scalar
values and the declared party label are data of the instance that the
instance at a datum keeps by definition; only the datum changes, and the
datum occupies the data sectors alone. Where $G=K$, the rule of (I1) is
applied to the family it formed the handed datum from, so
$A^{\mathrm{ref}}=A$.

(ii) The relation is (I3) of the reference companion, which (i) shows to be
a reduction instance; its first clause is
$\Lambda[A^{\mathrm{ref}}]^{\dagger}(F_v[A^{\mathrm{ref}}])
=M_v[A^{\mathrm{ref}}]$ by Definition~\ref{def:squashing}. Every
certificate of $\CList[A^{\mathrm{ref}}]$ and, where flagged, the
class-marginal identity are (I4) at $A^{\mathrm{ref}}$, discharged there by
the stability clause. Theorem~\ref{thm:transfer} needs only the relation.
The one-link segment with source family $(\Bsys,M[A^{\mathrm{ref}}])$,
target $(\Qsys,F[A^{\mathrm{ref}}])$, channel witness
$\Lambda[A^{\mathrm{ref}}]$, the transported pair
$(\hat O[A^{\mathrm{ref}}],O[A^{\mathrm{ref}}])$, the \termref{term-reachable-set-statement}{reachable-set statement}
of that channel and the entry list $\CList[A^{\mathrm{ref}}]$ is a
receiver-stage link of Definition~\ref{def:link} and a chain of
Definition~\ref{def:chain} with an empty \exttermref{../transmitter/contract.pdf}{term-source-segment}{source segment}, and
Theorem~\ref{thm:chain}\,(iv) gives it $\textsf{status}=\textsf{exact}$,
the relation carrying no error term. The sector pinching in front of
$\Lambda_0[A^{\mathrm{ref}}]$ changes none of this
(Corollary~\ref{cor:blockdiag}\,(i)).

(iii) Each entry of $A^{\mathrm{ref}}$ is the value of the expression the
rule of (I1) uses, at the entries of $G$ and at the handed numbers. Those
handed numbers that are not entries of a block are the same as for $A$ and
are known exactly; an exactly known number is enclosed by itself with radius
zero (\xref{receiver/def:complex-enclosure}), and an entrywise enclosure of a
block is a complex enclosure of each of its entries. The rules (i)--(iii) of
\xref{receiver/lem:enclosure-arithmetic} enclose sums, products, conjugates
and, where the condition of its part (iii) is met, reciprocals of enclosed
numbers, with centres formed by the same operation on the centres, and keep
an enclosure exact where the inputs are; by its part (iv), applied along the
expression under the hypothesis that every reciprocal node receives an
enclosure, they give a complex enclosure of every entry of
$A^{\mathrm{ref}}$, with the centre and the radii stated. Every element of
$A^{\mathrm{ref}}$ is a Hermitian matrix, the datum being admissible (i), so
part (v) of that lemma turns the complex enclosure of each element into an
entrywise enclosure in the sense of \xref{foundations/def:exact-reference},
with Hermitian centre $\tfrac12(X_c+X_c^{\dagger})$ and symmetric radii
$\tfrac12(r+r^{\transpose})$, exact at an entry $(i,l)$ where
$r_{il}=r_{li}=0$. For an entry
operator, $\hat O_k[A^{\mathrm{ref}}]$ restricted to a data sector $s$ is
$\sum_v\alpha_{k,v,s}A^{\mathrm{ref},(s)}_v$ plus the coefficient of $\Pi_s$
in $Z_k$ times $\idon{d_s}$, exactly known; at each entry the triangle
inequality gives $\big|\sum_v\alpha_{k,v,s}(A^{\mathrm{ref},(s)}_v
-C^{(s)}_v)_{il}\big|\le\sum_v|\alpha_{k,v,s}|\,R^{(s)}_{v,il}$, with
$(C,R)$ any entrywise enclosure of the datum. On a scalar sector the operator is a
combination of the scalar values, which are the instance's and exact.

(iv) For the identity instance the rule of (I1) is the identity: the datum
is the declared family read as a sector-block list, and it is admissible
exactly when each block is Hermitian and positive semidefinite and the class
sums are $q_t$ times the identity, the completeness identities its
declaration states. A family $G$ with the properties listed has them, so the
datum it forms is admissible. That the identity instance is an instance over
$\mathbb{R}$ is the proof of Corollary~\ref{cor:identity-instance}, which
reads of the substrate its field operations and its order and of a datum
only admissibility: its channel, the sector pinching, does not depend on
the datum. For the branching and the register instances,
$\mathcal{D}_{\mathrm{adm}}=\prod_j\mathcal{D}^{(j)}_{\mathrm{adm}}$, a
composite datum being admissible exactly when every branch datum is (Step~1
of the proof of Theorem~\ref{thm:instrument-instance}, which
Corollary~\ref{cor:register-instance} applies); replacing $K$ by $G$ changes
only the branch data formed from blocks of $K$, each admissible by (H-REF)
for its branch, and leaves the others as they are. Both results establish
the composite instance from its branch instances by arguments that read of
the substrate its field operations and its order, so the composite is an
instance over $\mathbb{R}$ when its branches are.

(v) No step of (i)--(iv) reads the assumption. Where $B$ is not empty, each
block in $B$ is a block of the declaring instance's snap of the measurement
of a device, the data $D$ of that snap being what the instance states as
that measurement or a reference of it, so clause (i) of
\xref{foundations/ass:snapped-measurement} applies to it: the block of the
measurement the device performs is the snapped family there, with the
entries recorded for the snap, in place of the model. On such a block
$G^{(b)}=K^{(b)}$, so the companion is formed there from that snapped block.
A nearness of the snapped block and the model relates no key lengths, and
no claim relates the two on a block in $B$, which is a \exttermref{../foundations/contract.pdf}{term-moved}{moved} block of
the snap, as the assumption's $B$ is defined; so no claim relates the
companion's key length to the model there.

(vi) The companion form is an operator of the form
\eqref{eq:instance-entry} at the datum $A^{\mathrm{ref}}$, which
(iii) encloses. Where every $\alpha_{v,s}$ is zero, \eqref{eq:companion-form}
reads no datum, so $\hat O^{\mathrm{ref}}=Z=\hat O$; $Z$ is a sum of sector
projectors, each in the \termref{term-transported-set}{transported set} of the companion by (I3) at
$A^{\mathrm{ref}}$, and $\Lambda[A^{\mathrm{ref}}]^{\dagger}$ is a positive
map (the third consequence of Definition~\ref{def:reduction-instance}), so
$\Lambda[A^{\mathrm{ref}}]^{\dagger}(Z)\succeq0$, as
Corollary~\ref{cor:entries}\,(i) argues at $A'$. For the class-marginal
entry, $F_v[A]$ is sector-block-diagonal at every admissible datum (i), so
$\sum_s\Pi_sF_v[A]\Pi_s=F_v[A]$ and
$\hat O^{\mathrm{ref}}=\sum_{v\in\Vset_t}F_v[A^{\mathrm{ref}}]$. On a scalar
sector $s$ the block of $\sum_{v\in\Vset_t}F_v$ is formed from the scalar
values, which are the same at both data (i). On a data sector $s$ it is
$\sum_{v\in\Vset_t}A^{(s)}_v$ at the datum it is read at, the elements
\termref{term-outside-the-sector-support}{outside the sector support} being zero; an admissible datum is
\termref{term-class-complete}{class complete}, so this sum is $q_t\idon{d_s}$ at $A^{\mathrm{ref}}$,
admissible by (i), and at $A'$, admissible as an element of
$\mathcal{D}_{\kappa}$ (Definition~\ref{def:emitnorm}). The two operators
therefore agree block by block, and $\hat O^{\mathrm{ref}}=\hat O$. Where
the instance discharges \eqref{eq:instance-classmarginal}, (ii) gives that
identity at $A^{\mathrm{ref}}$ and
$\Lambda[A^{\mathrm{ref}}]^{\dagger}(F_v[A^{\mathrm{ref}}])
=M_v[A^{\mathrm{ref}}]$, so by linearity
$\Lambda[A^{\mathrm{ref}}]^{\dagger}(\hat O^{\mathrm{ref}})
=\sum_{v\in\Vset_t}M_v[A^{\mathrm{ref}}]=q_t\idon{\Bsys}$, which is
\eqref{eq:certificate} with $\bowtie$ the equality at $c=q_t$. For any
other entry the certificate was established at $A'$, for the numbers
emitted there (Corollary~\ref{cor:entries}\,(iii)), and a certificate of an operator
formed from one datum is no statement about the operator formed from
another; so none is asserted. A writing of $\hat O$ in the form
\eqref{eq:instance-entry} need not be unique, and two writings may give
different companion forms; each is an entry wherever its own certificate
holds, which is all this part claims of it. Finally, where
$\Lambda[A^{\mathrm{ref}}]^{\dagger}(\hat O^{\mathrm{ref}})\bowtie
c\,\idon{\Bsys}$ holds, (ii) is a \termref{term-squashing-relation}{squashing relation} at
$A^{\mathrm{ref}}$, and adjoining the pair
$\big(\Lambda[A^{\mathrm{ref}}]^{\dagger}(\hat O^{\mathrm{ref}}),
\hat O^{\mathrm{ref}}\big)$ to its \termref{term-carried-operators}{carried operators} keeps
\eqref{eq:squashing}, whose clause for that pair holds by the choice of its
physical operator (Definition~\ref{def:squashing}); the certificate is then
\eqref{eq:certificate} for the adjoined carried operator.
\end{proof}

% claim reductions/cor:companion-factor-enclosure
\claimsummary{reductions/cor:companion-factor-enclosure}{Where an instance forms its datum as a non-negative rational combination of measurement blocks, factor enclosures of the reference's blocks give one of each element of the companion's datum, columns side by side and cores scaled, and so of the target's block on that data sector.}
\claimbookmark{reductions/cor:companion-factor-enclosure}{Corollary}{cor:companion-factor-enclosure}{Factor enclosures carried to the reference companion}
\claimleadin{%
Corollary~\ref{cor:reference-companion}\,(iii) carries the
\exttermref{../foundations/contract.pdf}{term-entrywise-enclosure}{entrywise enclosures} of the handed reference to
the companion's datum through the expression the instance's rule of (I1) is.
A \exttermref{../foundations/contract.pdf}{term-factor-enclosure}{factor enclosure} passes through a narrower class of rules. A sum, with non-negative weights, of members of
factor enclosures is a member of the enclosure whose columns are theirs side
by side and whose core is theirs scaled by the weights, and no arithmetic on
the columns is needed for it. The corollary names that class of rules and
carries the factor enclosures through it. The identity instance is in it, and
so is the flat instance (Corollary~\ref{cor:flat-companion}). A rule outside
it, such as one that maps a block by a conjugation before it combines blocks,
carries no factor enclosure here, and the elements it forms are known
through their entrywise enclosures alone
(Requirement~\ref{req:reference-companion}).
}
\begin{corollary}[Factor enclosures carried to the reference companion]\label{cor:companion-factor-enclosure}
Let $\mathcal{J}$, $K$, $G$ and $A^{\mathrm{ref}}$ be as in
Corollary~\ref{cor:reference-companion}, the blocks $K^{(b)}_j$ indexed by a
block $b$ and an outcome $j$, together with the setting where the blocks are
given per setting. Call the rule of part (I1) of $\mathcal{J}$
(Definition~\ref{def:reduction-instance}) a \emph{non-negative combination}
when for every \termref{term-data-sector}{data sector} $s$ and every $v\in\Vset$ there are
non-negative rationals $\lambda^{(s,v)}_{b,j}$ such that, for every family
$K'$ with the properties that Corollary~\ref{cor:reference-companion} lists
for $G$, the datum $A'$ the rule forms from $K'$ satisfies
\begin{equation}\label{eq:datum-combination}
  A'^{(s)}_v\;=\;\sum_{b,j}\lambda^{(s,v)}_{b,j}\,K'^{(b)}_j ,
\end{equation}
every block $K'^{(b)}_j$ with $\lambda^{(s,v)}_{b,j}>0$ being of dimension
$d_s$. Let the rule be a non-negative combination, and let some blocks
$G^{(b)}_j$ carry \exttermref{../foundations/contract.pdf}{term-factor-enclosure}{factor enclosures}
$\big(L^{(b)}_{j},\varrho^{(b)}_{j},W^{(b)}_{j}\big)$
(\xref{foundations/def:factor-enclosure}).
\begin{enumerate}[label=(\roman*),leftmargin=3.0em]
\item Let $s$ be a data sector and $v\in\Vset$ such that every block
  $G^{(b)}_j$ with $\lambda^{(s,v)}_{b,j}>0$ carries one, and list the pairs
  $(b,j)$ with $\lambda^{(s,v)}_{b,j}>0$ in a fixed order. The triple whose
  \exttermref{../foundations/contract.pdf}{term-matrix-box}{matrix box} has as its centre and its radii those of the boxes
  $\big(L^{(b)}_{j},\varrho^{(b)}_{j}\big)$ side by side in that order, and
  whose core is the diagonal matrix of the
  $\lambda^{(s,v)}_{b,j}W^{(b)}_{j}$ in that order, is a factor enclosure of
  $A^{\mathrm{ref},(s)}_v$, its number of columns the sum of theirs; where
  every $\lambda^{(s,v)}_{b,j}$ is zero, it has no column.
\item The block of the companion's target $F_v[A^{\mathrm{ref}}]$ on the data
  sector $s$ is $A^{\mathrm{ref},(s)}_v$, so the triple of (i) is a factor
  enclosure of that block.
\item The rule of (I1) of the identity instance of
  Corollary~\ref{cor:identity-instance} is a non-negative combination: each
  element of its datum is the block of the family it is read from at the
  same sector and outcome, with the coefficient $1$, every other coefficient
  being $0$.
\end{enumerate}
\end{corollary}
\claimmeta{\metalabel{Uses}\mref{Def.}{def:reduction-instance}, \mref{Cor.}{cor:identity-instance}, \mref{Cor.}{cor:reference-companion}, \extref{../foundations/contract.pdf}{foundations/def:factor-enclosure}
\metasep \metalabel{Used by}\mref{Req.}{req:reference-companion}, \mref{Cor.}{cor:flat-companion}, \extref{../meat-core/contract.pdf}{meat-core/cor:reference-companion-factors}
\metasep newly derived
\metasep implemented (effective: accepted)
\metasep \metalabel{Code}\coderef{https://github.com/fraunhofersg/certiqs-key/blob/v0.1.0/src/certiqs_key/reductions/companion.py}{certiqs_key.reductions.companion:combined_factor}}

\begin{proof}
\emph{(i).} By Corollary~\ref{cor:reference-companion}, $A^{\mathrm{ref}}$ is
the datum the rule of (I1) forms from $G$, a family with the properties that
corollary lists, so \eqref{eq:datum-combination} with $K'=G$ gives
$A^{\mathrm{ref},(s)}_v=\sum_{(b,j)}\lambda^{(s,v)}_{b,j}G^{(b)}_j$, the sum
over the listed pairs, the others having coefficient $0$. Each listed block
carries its factor enclosure, so there is a matrix $L^{(b)}_{j,*}$ that the
box $\big(L^{(b)}_{j},\varrho^{(b)}_{j}\big)$ contains with
$G^{(b)}_j=L^{(b)}_{j,*}W^{(b)}_{j}L^{(b)\dagger}_{j,*}$
(\xref{foundations/def:factor-enclosure}); every such matrix has $d_s$ rows,
its block being of dimension $d_s$. Let $L_*$ be the matrix with the
$L^{(b)}_{j,*}$ side by side in the listed order and $W$ the block-diagonal
matrix of the $\lambda^{(s,v)}_{b,j}W^{(b)}_{j}$ in the same order. Each
column of the box of the triple is a column of one of the boxes it is made
of, with its centre and radii, so the box contains $L_*$; its centre has
entries in $\mathbb{Q}(i)$ and its radii are non-negative rationals, as
those of the boxes it is made of; and $W$ is diagonal, each
$W^{(b)}_j$ being diagonal, with positive rational diagonal entries, each a
product of a positive rational $\lambda^{(s,v)}_{b,j}$ and a positive
rational entry of $W^{(b)}_j$. Multiplying by blocks,
\begin{equation*}
  L_*WL_*^{\dagger}
  =\sum_{(b,j)}L^{(b)}_{j,*}\big(\lambda^{(s,v)}_{b,j}W^{(b)}_{j}\big)L^{(b)\dagger}_{j,*}
  =\sum_{(b,j)}\lambda^{(s,v)}_{b,j}\,G^{(b)}_j
  =A^{\mathrm{ref},(s)}_v ,
\end{equation*}
which is what a factor enclosure of $A^{\mathrm{ref},(s)}_v$ asks
(\xref{foundations/def:factor-enclosure}). The number of columns of $L_*$ is
the sum of those of the $L^{(b)}_{j,*}$. Where every coefficient is zero, no
pair is listed, $L_*$ has no column, and $A^{\mathrm{ref},(s)}_v=0$, the one
member of a triple with no column.

\emph{(ii).} By Corollary~\ref{cor:reference-companion}\,(i) the reference
companion is a reduction instance at the datum $A^{\mathrm{ref}}$, and its
target is \eqref{eq:instance-target} of
Definition~\ref{def:reduction-instance} at that datum, whose block on a data
sector $s$ is $A^{\mathrm{ref},(s)}_v$.

\emph{(iii).} The rule of (I1) of the identity instance is the identity: its
datum is the declared family read as a sector-block list,
$A^{(s)}_v$ being the block of $M_v$ on the sector $s$
(Corollary~\ref{cor:identity-instance}), and the same rule reads any family
$K'$ over the same sectors so. That is \eqref{eq:datum-combination} with
$\lambda^{(s,v)}_{s,v}=1$ and every other coefficient $0$, and the block
with the coefficient $1$ is of dimension $d_s$.
\end{proof}

% claim reductions/lem:row-relaxation
\claimsummary{reductions/lem:row-relaxation}{Rounding a certified entry's operator and moving its constant outward by a bound on the rounding's norm gives an entry that is again certified and that every attack and every state satisfying the exact entry satisfies. An equality becomes two inequalities.}
\claimbookmark{reductions/lem:row-relaxation}{Lemma}{lem:row-relaxation}{A rounded constraint entry with its constant relaxed}
\claimleadin{%
A constraint entry built in exact arithmetic from data off the binary64 grid
has an operator that binary64 cannot carry. Rounding the operator moves the
set the entry cuts out, in either direction. Moving the constant outward by a
bound on the size of the rounding makes the rounded entry weaker than the
exact one, so that whatever satisfies the exact entry satisfies the rounded
one, and the rounded entry still carries a certificate of its own. The bound
is read off the entries, since a Hermitian matrix has operator norm at most
its largest absolute row sum.
}
\begin{lemma}[A rounded constraint entry with its constant relaxed]\label{lem:row-relaxation}
Let $\Lambda\colon\Tcl{\Bsys}\to\Tcl{\Qsys}$ be a channel with
$(M;O)\dsq_{\Lambda}(F;\hat O)$ (Definition~\ref{def:squashing}), and let
$(\hat O_k,\bowtie_k,c_k)$ be an entry with $\hat O_k$ a carried operator,
$\bowtie_k\in\{\preceq,=,\succeq\}$, $c_k\in\mathbb{R}$ and the certificate
$\Lambda^{\dagger}(\hat O_k)\bowtie_kc_k\,\idon{\Bsys}$ of
\eqref{eq:certificate}. Let $\hat O'_k$ be a Hermitian operator on $\Qsys$,
let $E_k\coloneqq\hat O'_k-\hat O_k$, and let $\rho_{il}\ge0$ be numbers with
$|(E_k)_{il}|\le\rho_{il}$ for every entry in a fixed orthonormal basis of
$\Hs{\Qsys}$, and $\delta_k\ge\max_i\sum_l\rho_{il}$. Define the
\emph{relaxed entries} of $(\hat O_k,\bowtie_k,c_k)$ as
$(\hat O'_k,\succeq,c'_k)$ with $c'_k\le c_k-\delta_k$ where $\bowtie_k$ is
$\succeq$; $(\hat O'_k,\preceq,c''_k)$ with $c''_k\ge c_k+\delta_k$ where
$\bowtie_k$ is $\preceq$; and both of these where $\bowtie_k$ is $=$. Then:
\begin{enumerate}[label=(\roman*),leftmargin=2.4em]
\item $-\delta_k\idon{\Qsys}\preceq E_k\preceq\delta_k\idon{\Qsys}$;
\item \emph{(certificate)} $(M;O,\Lambda^{\dagger}(\hat O'_k))\dsq_{\Lambda}
  (F;\hat O,\hat O'_k)$, and every relaxed entry carries the spectral
  certificate \eqref{eq:certificate} at its own constant:
  $\Lambda^{\dagger}(\hat O'_k)\succeq c'_k\idon{\Bsys}$, respectively
  $\Lambda^{\dagger}(\hat O'_k)\preceq c''_k\idon{\Bsys}$;
\item \emph{(\exttermref{../foundations/contract.pdf}{term-attack-set}{attack sets})} for all finite-dimensional $E',E''$, every
  $\Aev$ in the set \eqref{eq:attset}, \eqref{eq:attsetun} or, for an entry
  tagged party-local with the factor $f$, \eqref{eq:attsetpl}, of the list
  holding the exact entry lies in the set of the list in which the exact
  entry is replaced by its relaxed entries, the other entries unchanged;
\item \emph{(states)} for every finite-dimensional system $X$ and every
  positive semidefinite $\omega$ on $\Hs{X}\otimes\Hs{\Qsys}$,
  $\tr_{\Qsys}[(\idon{X}\otimes\hat O_k)\,\omega]\succeq c_k\tr_{\Qsys}\omega$
  implies $\tr_{\Qsys}[(\idon{X}\otimes\hat O'_k)\,\omega]\succeq
  c'_k\tr_{\Qsys}\omega$, and the same with $\preceq$, $c''_k$ in place of
  $\succeq$, $c'_k$; with $X$ trivial these are the scalar readings
  $\tr[\hat O_k\omega]\ge c_k\tr\omega$ and the like.
\end{enumerate}
\end{lemma}
\claimmeta{\metalabel{Uses}\mref{Lem.}{lem:sep}, \mref{Def.}{def:squashing}, \mref{Def.}{def:certificate}, \mref{Def.}{def:attset}, \mref{Thm.}{thm:licence}
\metasep \metalabel{Used by}\mref{Req.}{req:reference-companion}, \extref{../meat-core/contract.pdf}{meat-core/prop:reference-companion}, \extref{../protocol-compiler/contract.pdf}{protocol-compiler/req:reference-companion-compiled}
\metasep newly derived
\metasep implemented (effective: accepted)
\metasep \metalabel{Code}\coderef{https://github.com/fraunhofersg/certiqs-key/blob/v0.1.0/src/certiqs_key/reductions/companion.py}{certiqs_key.reductions.companion:relaxed_entries}}

\begin{proof}
(i) $E_k$ is Hermitian, as a difference of Hermitian operators, so its
operator norm is the largest modulus of an eigenvalue. Let $E_kx=\lambda x$
with $x\neq0$ and let $i$ be an index with $|x_i|=\max_l|x_l|$. Then
$|\lambda||x_i|=|\sum_l(E_k)_{il}x_l|\le\sum_l\rho_{il}|x_l|
\le\big(\sum_l\rho_{il}\big)|x_i|$, so $|\lambda|\le\delta_k$, and the
spectral theorem gives $-\delta_k\idon{}\preceq E_k\preceq\delta_k\idon{}$.

(ii) $\Lambda^{\dagger}$ is linear, positive and unital (Lemma~\ref{lem:sep}),
so (i) gives $-\delta_k\idon{\Bsys}\preceq\Lambda^{\dagger}(E_k)\preceq
\delta_k\idon{\Bsys}$. Where $\bowtie_k$ is $\succeq$ or $=$,
$\Lambda^{\dagger}(\hat O_k)\succeq c_k\idon{\Bsys}$, hence
$\Lambda^{\dagger}(\hat O'_k)=\Lambda^{\dagger}(\hat O_k)+
\Lambda^{\dagger}(E_k)\succeq(c_k-\delta_k)\idon{\Bsys}\succeq c'_k
\idon{\Bsys}$; where it is $\preceq$ or $=$, symmetrically
$\Lambda^{\dagger}(\hat O'_k)\preceq(c_k+\delta_k)\idon{\Bsys}\preceq
c''_k\idon{\Bsys}$. The relation with $\hat O'_k$ carried is
\eqref{eq:squashing} with the one further carried operator $\hat O'_k$ and
its pull-back $\Lambda^{\dagger}(\hat O'_k)$, which holds by the definition
of that pull-back; the conditions on $F$, $M$, $\hat O$ and $O$ are the given
relation. Each relaxed entry's certificate is the operator inequality just
shown, a case of \eqref{eq:certificate}.

(iii) The adjoint of a channel, and of $\idmap_{R'}\otimes\Aev$, is unital and
completely positive (Lemma~\ref{lem:sep}). For \eqref{eq:attset},
$-\delta_k\idon{}\preceq E_k\otimes\idon{E''}\preceq\delta_k\idon{}$ by (i),
so $\Aev^{\dagger}(\hat O'_k\otimes\idon{E''})=\Aev^{\dagger}(\hat O_k\otimes
\idon{E''})+\Aev^{\dagger}(E_k\otimes\idon{E''})$ is $\succeq(c_k-\delta_k)
\idon{E'}\succeq c'_k\idon{E'}$ where $\Aev^{\dagger}(\hat O_k\otimes
\idon{E''})\succeq c_k\idon{E'}$, and $\preceq(c_k+\delta_k)\idon{E'}\preceq
c''_k\idon{E'}$ where $\Aev^{\dagger}(\hat O_k\otimes\idon{E''})\preceq
c_k\idon{E'}$; an equality gives both. The same with
$\idon{A}\otimes\hat O_k\otimes\idon{E''}$ proves it for \eqref{eq:attsetun}.
For \eqref{eq:attsetpl}, with $X\succeq0$ on $R'R\,\Qsys_{f^c}$,
\begin{equation*}
  X\otimes(\hat O'_k-c'_k\idon{})\otimes\idon{E''}
  =X\otimes(\hat O_k-c_k\idon{})\otimes\idon{E''}
  +X\otimes\big(E_k+(c_k-c'_k)\idon{}\big)\otimes\idon{E''},
\end{equation*}
and $E_k+(c_k-c'_k)\idon{}\succeq E_k+\delta_k\idon{}\succeq0$ by (i), so
the last term is a tensor product of positive semidefinite operators and its
image under the completely positive adjoint is positive semidefinite; the
first term's image is $\succeq0$ where the exact entry holds. For
$\preceq$ the second bracket is $E_k-(c''_k-c_k)\idon{}\preceq
E_k-\delta_k\idon{}\preceq0$, and the images are $\preceq0$. An equality
entry gives both. The other entries of the list are unchanged, so their
conditions are the same in both sets.

(iv) Put $P\coloneqq E_k+\delta_k\idon{\Qsys}\succeq0$ by (i). For
$\omega\succeq0$,
$\tr_{\Qsys}[(\idon{X}\otimes P)\omega]
=\tr_{\Qsys}[(\idon{X}\otimes\sqrt P)\,\omega\,(\idon{X}\otimes\sqrt P)]
\succeq0$, by cyclicity of the partial trace for an operator acting on
$\Qsys$ alone and because a partial trace of a positive semidefinite operator
is positive semidefinite. Hence
$\tr_{\Qsys}[(\idon{X}\otimes\hat O'_k)\omega]\succeq
\tr_{\Qsys}[(\idon{X}\otimes\hat O_k)\omega]-\delta_k\tr_{\Qsys}\omega
\succeq(c_k-\delta_k)\tr_{\Qsys}\omega\succeq c'_k\tr_{\Qsys}\omega$, the
last step because $\tr_{\Qsys}\omega\succeq0$ and $c_k-\delta_k\ge c'_k$. The
case $\preceq$ uses $\delta_k\idon{}-E_k\succeq0$ in the same way. With $X$
trivial the operators are numbers and the inequalities are the scalar ones.
\end{proof}

\section{Assumptions}\label{group-reductions/ass:meas}
\sectionintro{These assumptions are the premises the reductions rest on and do not prove: a measuring party's device entering only through one POVM per round, the same physical measurement in every round with no memory across rounds, a supplier's artefact being its module's own output at an accepted version and true of the device, a measuring party's private setting being its own declared randomness, output-state test statistics and declared constants. The dimensions of the systems are the setting's, and the adversary's registers are finite-dimensional as the round model's are; no assumption is made about the adversary's memory, which the consuming engine reaches at any dimension by a limit (\xref{meat-core/cor:composed-security}). No assumption about the distance of the declared numbers from a true operator family is made, because their exact admissibility is decided on entry; what remains a hypothesis is that the supplier's facts hold of the family declared (\hyperlink{claim-reductions/ass:recv}{Assumption~\ref*{ass:recv}}, \hyperlink{claim-reductions/rem:recv-scope}{Remark~\ref*{rem:recv-scope}}), which for an exact, rigorously evaluated receiver artefact the run built itself holds under the receiver's own ledger (\hyperlink{claim-reductions/req:recv-row-provenance}{Requirement~\ref*{req:recv-row-provenance}}).}
\sectionimports{\extref{../foundations/contract.pdf}{foundations/def:round-homogeneity}}

% claim reductions/ass:meas
\claimsummary{reductions/ass:meas}{In every round each measuring party's device acts as one POVM: it emits a classical outcome, retains no quantum system, and everything after it acts on classical data. It is the hypothesis about the devices behind the output-state equality of the physical and the squashed protocol.}
\claimbookmark{reductions/ass:meas}{Assumption}{ass:meas}{Each measuring party's device as one POVM per round}
\claimleadin{%
The assumption is what lets a round's processing factor through the
\termref{term-measuring-parties}{measuring parties}' \termref{term-measurement-channel}{measurement channel}: once each device is one
POVM whose outcome is all that leaves it, the rest of the round acts on
classical data. Where several parties measure, it is asserted of each of
their devices, and where the declaration declares several measured factors
of the \termref{term-measured-system}{measured system}, of the device at each measured factor, two
factors of one party included; where one party
measures (Bob, in a prepare-and-measure declaration), it is asserted of that
party's device. The round's order is
not part of it; classical detector memory across rounds is excluded by the
round-homogeneity assumption.
}
\SetShortId{A-MEAS}
\begin{assumption}[Each measuring party's device as one POVM per round]\label{ass:meas}
In every round, each measuring party's interaction with the quantum system it
receives satisfies:
(i) it is described by a single POVM on that system;
(ii) it emits the classical outcome and nothing else, and no
post-measurement quantum system is retained or reused;
(iii) all downstream processing (announcements, sifting, acceptance, key map,
error correction, privacy amplification) acts on classical data;
(iv) the round's \termref{term-measurement-channel}{measurement channel} is the product of one
measurement channel per factor, each acting on its own factor alone, for each
of the following factorizations of the \termref{term-measured-system}{measured system} $\Bsys_j$:
where the declaration declares several measured factors
(\xref{protocol-compiler/def:spec}), the factorization
$\Bsys_j=\bigotimes_f\Bsys_{f,j}$ with one factor per measured factor,
whichever parties carry their labels and two factors of one party included;
and where several parties measure, the factorization into the parties'
parts of Definition~\ref{def:protocol}.
\end{assumption}
\claimverification{The engine does not verify this assumption; a run that rests on it records it as unverified.}
\claimreading{Each measuring device, at every measuring station of either party, turns the incoming quantum signal into a classical result in every round and keeps nothing quantum afterwards: no post-measurement system is stored, re-used or passed on. Everything after detection (announcements, sifting, acceptance, the key map, error correction, privacy amplification) uses classical data only. Where several measuring stations are declared, of different parties or of one, or where several parties measure inside one declared station, each device measures independently of the others, so that the round's measurement is the product of theirs: no shared randomness or shared ancilla couples them. A receiver with a quantum memory that keeps the measured light, or passes it on, violates this.}
\claimmeta{\metalabel{Uses}\mref{Def.}{def:protocol}
\metasep \metalabel{Category}receiver
\metasep \metalabel{Affects}secrecy
\metasep \metalabel{Assumed by}\mref{Lem.}{lem:opreduction-transfer}, \mref{Thm.}{thm:factorization}, \mref{Rem.}{rem:coherent}, \mref{Rem.}{rem:criteria}, \extref{../protocol-compiler/contract.pdf}{protocol-compiler/req:link-row-provenance}
\metasep \metalabel{Machine key}\iface{measurement_only_receiver}
\metasep newly derived
\metasep accepted}

% claim reductions/ass:round
\claimsummary{reductions/ass:round}{One physical POVM describes the measurement of every round, and each round's measurement acts as it describes whatever happened in other rounds: no detector memory and no drift. Nothing is assumed about the attacks being independent across rounds.}
\claimbookmark{reductions/ass:round}{Assumption}{ass:round}{Round homogeneity of the physical measurement}
\claimleadin{%
The measurement this assumption holds fixed is the physical one, the POVM by
which the devices of the \termref{term-measuring-parties}{measuring parties} measure the \termref{term-measured-system}{measured
system} of a round; where several parties measure, the assumption fixes each
party's factor across the $n$ rounds. Rounds are not averaged, and per-round
variation is not modelled. The \termref{term-declared-model}{declared model} and the
\termref{term-squashing-channel}{squashing channel} are objects of the analysis, not of the device, and
are not part of this assumption; neither is the round's \emph{order}, every
message of the round after every measurement of the round.
}
\SetShortId{A-ROUND-SQ}
\begin{assumption}[Round homogeneity of the physical measurement]\label{ass:round}
Let $M$ be the physical POVM: the POVM on the \termref{term-measured-system}{measured system} of a
round that the analysis takes as the measurement the devices of the
\termref{term-measuring-parties}{measuring parties} perform on it.
The family $\{M\}$ is \exttermref{../foundations/contract.pdf}{term-round-homogeneous}{round-homogeneous}
(\xref{foundations/def:round-homogeneity}) in both clauses of that
definition:
(i) the same $M$ describes the measurement of every round;
(ii) a round's measurement acts as $M$ describes, whatever the inputs and
outputs of the other rounds, with no detector memory across rounds.
\end{assumption}
\claimverification{The engine does not verify this assumption; a run that rests on it records it as unverified.}
\claimreading{Each measuring device, at every measuring station of either party, performs the same measurement in every round: the same optics, detectors and dark-count behaviour, settings drawn with the same probabilities, and, where you declared detector-response modes, the same response in each mode. What a device does in one round does not depend on the inputs or outputs of other rounds. Drift of efficiency, alignment or dark counts violates this. Only parts of a run fixed before it starts, each processed as its own execution and composed under \hyperref{../meat-core/contract.pdf}{}{claim-meat-core/ass:separate-executions}{A-SEPARATE}, can be analysed separately; a split chosen after seeing the data is not covered. Detector memory (afterpulsing, a dead time, any dependence on earlier rounds) violates it too, and nothing here repairs that. Attacks may still correlate the rounds. The engine's finite model and squashing map are its own construction, not part of this premise.}
\claimmeta{\metalabel{Uses}\mref{Def.}{def:protocol}, \extref{../foundations/contract.pdf}{foundations/def:round-homogeneity}
\metasep \metalabel{Category}receiver
\metasep \metalabel{Affects}secrecy
\metasep \metalabel{Refines}\extaref{../protocol-compiler/contract.pdf}{protocol-compiler/ass:rhom}{A-RHOM}
\metasep \metalabel{Assumed by}\mref{Thm.}{thm:factorization}, \mref{Rem.}{rem:variants}, \mref{Rem.}{rem:coherent}
\metasep \metalabel{Machine key}\iface{round_homogeneity_squashing}
\metasep newly derived
\metasep accepted}

% claim reductions/ass:recv
\claimsummary{reductions/ass:recv}{A supplier artefact read by a reduction instance is that module's own output at the accepted version it declares, and the device facts its claims state hold of the numbers exactly as declared. It is the premise behind a reduction on supplier numbers, save an exact, rigorous artefact the run built.}
\claimbookmark{reductions/ass:recv}{Assumption}{ass:recv}{Supplier artefact whose claims hold of the device}
\claimleadin{%
For the receiver supplier the numbers in question are the photon-number
blocks, the high-photon weights with their quantifier domain and the
\exttermref{../receiver/contract.pdf}{term-dark-count-channel}{dark-count channel}. Which of a
supplier's claims are consumed, and so which of its assumptions the imported
facts rest on, is stated where the supplier's artefact is consumed; the
assumption says that those facts hold of the numbers declared. Where the run
built a receiver artefact itself, from a declaration with a model admissible
exactly and with every weight evaluated rigorously, both clauses hold
without it, by the receiver's own theorems under its own assumptions, and
the engine omits its row; at a weight evaluated in floating point, or at a
declaration admitted only within the receiver's tolerances, clause~(ii)
remains a premise.
}
\SetShortId{A-RECV}
\begin{assumption}[Supplier artefact whose claims hold of the device]\label{ass:recv}
Every declared block, weight, matrix and scalar that a
\termref{term-reduction-instance}{reduction instance} of this module reads from a supplier's artefact
satisfies:
(i) the artefact is the output of the supplying module, at the module version
it declares, which is an \emph{accepted} version of the supplying module, and
on the inputs it declares, and the validation results it carries are that
module's own;
(ii) the facts the supplier's claims state hold of those numbers as declared,
under that module's own assumption ledger.
\end{assumption}
\claimverification{The engine does not verify this assumption; a run that rests on it records it as unverified.}
\claimreading{The receiver data your reduction uses (blocks, high-photon weights with their quantifier domain, dark-count channel) come from a receiver artefact produced by an accepted version of the receiver module, the version the artefact names, on the inputs it names, with its own validation results attached, and the device facts that module's claims state are true of those numbers exactly as declared, under its own assumptions. Where this run built it from a declaration describing a valid receiver exactly, every weight evaluated rigorously, both halves follow and the report omits this premise (\hyperlink{claim-reductions/req:recv-row-provenance}{Requirement~\ref*{req:recv-row-provenance}}). For another artefact this run built, acceptance is the status the report gives the receiver's claims and you rely on the second half, floating-point weights included; a receiver file from outside the run needs your own justification for both halves.}
\claimmeta{\metalabel{Category}supplied artefact
\metasep \metalabel{Affects}secrecy
\metasep \metalabel{Assumed by}\mref{Rem.}{rem:notlicensed}, \mref{Prop.}{prop:route}, \mref{Rem.}{rem:absorbneeded}, \mref{Rem.}{rem:minweight}, \mref{Rem.}{rem:darkcompose}, \mref{Rem.}{rem:receiver-imports}, \mref{Prop.}{prop:flat-instance}, \mref{Cor.}{cor:emitted-flat}, \mref{Cor.}{cor:two-sector-decided}, \mref{Cor.}{cor:flat-companion}, \mref{Cor.}{cor:normalized-flat}
\metasep \metalabel{Machine key}\iface{receiver_artifact_accepted}
\metasep newly derived
\metasep accepted}

% claim reductions/rem:recv-scope
\claimsummary{reductions/rem:recv-scope}{\hyperlink{claim-reductions/ass:recv}{A-RECV} is stated of the exactly admissible declaration, not of an idealisation of it, and of one declared receiver. No claim relates a key length from the declaration to one from the supplier's own floating-point output; an exact, rigorous receiver artefact the run built has its row omitted.}
\claimbookmark{reductions/rem:recv-scope}{Remark}{rem:recv-scope}{What the supplier assumption is stated of}
\begin{remark}[What the supplier assumption is stated of]\label{rem:recv-scope}
Assumption~\ref{ass:recv} is stated of the declaration rather than of an
idealisation of it: what this module is handed is \exttermref{../foundations/contract.pdf}{term-exactly-admissible}{exactly admissible}
(Definition~\ref{def:exactadm}) and is decided so on entry, and what the
assumption adds is that the device facts the supplier states are true of
it. Where the declaration and the supplier's own floating-point output
differ, the difference is a modelling act taken at the \exttermref{../foundations/contract.pdf}{term-instance}{instance} and
recorded there, with that party's \exttermref{../foundations/contract.pdf}{term-declared-tolerance}{declared tolerance} as its ceiling
(\xref{foundations/req:declared-tolerance}); no claim of this contract
relates a key length computed from the declaration to one computed from
that output, exactly as Theorem~\ref{thm:emit}\,(v) says of this module's
own emission. Every one of the imported facts is inherited here unverified,
and a single device family, uncertainty set or characterisation protocol
appears nowhere: this contract reduces \emph{one} declared receiver. Which
of the supplier's claims are imported is stated where its artefact is
consumed (for the receiver, Remark~\ref{rem:receiver-imports}).
For a receiver artefact the run built itself, from a declaration with a
model admissible exactly and with every weight evaluated rigorously,
Requirement~\ref{req:recv-row-provenance} omits the row of
Assumption~\ref{ass:recv}: both its clauses hold there by the receiver's
theorems under the receiver's own ledger, with no premise about arithmetic.
Assumption~\ref{ass:recv} is the assumption on which the quantitative
content of every other reduction instance built on a supplier's numbers
rests (it
is part of the \termref{term-hypotheses-h}{hypotheses (H)} of such an instance,
Definition~\ref{def:reduction-instance}\,(I1)), and it is the reason
Theorem~\ref{thm:licence} is stated with hypothesis \eqref{eq:certificate}
rather than with a claim about optics.
\end{remark}
\claimmeta{newly derived
\metasep accepted}

% claim reductions/ass:probe
\claimsummary{reductions/ass:probe}{A station's private per-round setting is drawn from its party's own randomness with the declared class weights, independently of the adversary and the round's quantum input, never inferred from an outcome. It lets the setting enter a joint outcome under one common squashing channel.}
\claimbookmark{reductions/ass:probe}{Assumption}{ass:probe}{The station's private setting drawn from its own declared randomness}
\claimleadin{%
The setting of the optical instance is the probe setting. Drawn this way,
the setting enters as part of a joint outcome, the weight inside the
element, and one squashing channel serves every value of it; a setting the
adversary could read or steer would instead have to be handed to her.
}
\SetShortId{A-PROBE}
\begin{assumption}[The station's private setting drawn from its own declared randomness]\label{ass:probe}
The private per-round setting $t\in\Tset$ of a measuring party's device
satisfies, in every round:
(i) it is drawn from that party's own randomness with the
declared \termref{term-class-weights}{class weights} $q_t$;
(ii) it is drawn independently of the adversary's actions and of the round's
quantum input;
(iii) it is not inferred from any outcome.
\end{assumption}
\claimverification{The engine does not verify this assumption; a run that rests on it records it as unverified.}
\claimreading{The station's private setting in each round (for the optical instance, how its probe is switched) is chosen by a random number generator of the station's own party with exactly the declared probabilities. It is independent of anything the attacker does and of the light arriving in that round, and it is never derived from a detection result. An attacker who can predict or influence the setting violates this.}
\claimmeta{\metalabel{Category}local randomness
\metasep \metalabel{Affects}secrecy
\metasep \metalabel{Assumed by}\mref{Prop.}{prop:flat-instance}, \mref{Cor.}{cor:emitted-flat}, \mref{Cor.}{cor:flat-companion}, \mref{Cor.}{cor:normalized-flat}
\metasep \metalabel{Machine key}\iface{probe_setting_private}
\metasep newly derived
\metasep accepted}

% claim reductions/ass:tsm
\claimsummary{reductions/ass:tsm}{Every per-round public choice of known bias at the measurement the emitted family describes, its private setting included where it has one, is a component of the emitted outcome and enters only through its statistics. No derivation then assumes a state-independent test fraction.}
\claimbookmark{reductions/ass:tsm}{Assumption}{ass:tsm}{Output-state test statistics}
\claimleadin{%
For the optical instance the private per-round setting is the probe
setting. A choice that is a component of the outcome is analysed through the
statistics of the emitted family like every other component, so no fraction
of rounds of one kind is fixed independently of the state.
}
\SetShortId{A-TSM}
\begin{assumption}[Output-state test statistics]\label{ass:tsm}
Every per-round public choice of known bias in the physical protocol that is
applied at the measurement the emitted family $\{F_v\}$ describes, a private
per-round setting of that measurement and any round-type designation the
consumer adds by relabelling its outcomes included,
(i) is a component of the emitted outcome $v$, and
(ii) enters the analysis only through the statistics of $\{F_v\}$.
\end{assumption}
\claimverification{The engine does not verify this assumption; a run that rests on it records it as unverified.}
\claimreading{Every per-round random choice with a known bias made at this station (for example its probe setting, where it has one, or whether a round is a test round, where that is added by relabelling its outcomes) is recorded as part of that round's outcome, and it enters the analysis only through the observed statistics. No derivation of this contract then assumes a state-independent test fraction.}
\claimmeta{\metalabel{Category}protocol structure
\metasep \metalabel{Affects}secrecy
\metasep \metalabel{Assumed by}\mref{Rem.}{rem:intertwining}, \mref{Prop.}{prop:flat-instance}, \mref{Cor.}{cor:emitted-flat}, \mref{Cor.}{cor:flat-companion}, \mref{Cor.}{cor:normalized-flat}, \extref{../meat-core/contract.pdf}{meat-core/rem:not-licensed}, \extref{../statistics/contract.pdf}{statistics/binomial-box/rem:guard-activation}, \extref{../statistics/contract.pdf}{statistics/binomial-box/rem:sound-indep}
\metasep \metalabel{Machine key}\iface{test_statistics_model}
\metasep newly derived
\metasep accepted}

% claim reductions/ass:constdecl
\claimsummary{reductions/ass:constdecl}{Every constant of an emitted constraint entry is fixed before the run by the device fact that certifies it and is no function of an observed statistic. Observed data reach the analysis only through the acceptance set, which belongs to another module.}
\claimbookmark{reductions/ass:constdecl}{Assumption}{ass:constdecl}{Declared constants of the constraint list}
\claimleadin{%
The only sanctioned carrier of observed data in the analysis is the
acceptance set, which belongs to a different module and appears nowhere in
this contract.
}
\SetShortId{A-CONST}
\begin{assumption}[Declared constants of the constraint list]\label{ass:constdecl}
Every constant $c_k$ in an emitted constraint entry
(i) is fixed before the protocol run by the device fact that certifies it,
and
(ii) is not a function of any statistic observed during the run.
\end{assumption}
\claimverification{The engine does not verify this assumption; a run that rests on it records it as unverified.}
\claimreading{Applies only where the reduction emits constraints on the attacker; a run that emits none does not rest on it. Every number used as a constant in such a constraint (for example a bound on a detector's multi-photon response) is fixed before the run from the device characterisation that justifies it. It is never adjusted using data observed in the run; observed data enter only through the acceptance test.}
\claimmeta{\metalabel{Category}supplied artefact
\metasep \metalabel{Affects}secrecy
\metasep \metalabel{Assumed by}\mref{Rem.}{rem:unsound}, \mref{Thm.}{thm:licence}, \mref{Rem.}{rem:notlicensed}, \mref{Cor.}{cor:entries}, \mref{Rem.}{rem:claimscope}, \mref{Rem.}{rem:published}
\metasep \metalabel{Machine key}\iface{declared_constants}
\metasep newly derived
\metasep accepted}

\clearpage
\part{Engine requirements}
\sectionintro{What an implementation of this module must do, record or refuse; each requirement names the result that motivates it.}

\section{Requirements: channel witnesses, post-processing and relations}\label{group-reductions/req:channel-witness}
\sectionintro{These requirements fix what the artefact carries: an explicit \termref{term-channel-witness}{channel witness} for every emitted link, the artefact's witness being the \termref{term-squashing-channel}{squashing channel} followed by the sector pinching, and the carried tuple every emitted constraint is built from. They also fix what the engine refuses: a relation with an error term, a noise or coarse-graining map after a squashing link, a declared relation that fails a necessary condition, and a link of a reserved kind. Only the round-homogeneous model is emitted, and branch blocks are emitted as the branches supply them, a padded variant only where a submodule proves it and declares it in its interface.}

% claim reductions/req:channel-witness
\claimsummary{reductions/req:channel-witness}{Every link the engine emits carries an explicit channel witness, and a request naming a reduction that comes with none is refused. Without a squashing channel the transfer of \hyperlink{claim-reductions/thm:transfer}{Theorem~\ref*{thm:transfer}} is not available (\hyperlink{claim-reductions/rem:notchannel}{Remark~\ref*{rem:notchannel}}).}
\claimbookmark{reductions/req:channel-witness}{Requirement}{req:channel-witness}{A channel witness for every emitted link}
\claimleadin{%
A reduction that exhibits no \termref{term-squashing-channel}{squashing channel} cannot be fed into
Theorem~\ref{thm:transfer} (Remark~\ref{rem:notchannel}). Where the witness
sits in the artefact and which rule refuses a reduction without one are the
outputs table and the refusal rules of the interface
(\href{../../../spec/modules/reductions/interface.yaml}{\nolinkurl{spec/modules/reductions/interface.yaml}}).
}
\begin{requirement}[A channel witness for every emitted link]\label{req:channel-witness}
The engine emits every link with an explicit \termref{term-channel-witness}{channel witness} and
refuses a request that names a reduction coming with none.
\end{requirement}
\claimmeta{\metalabel{Motivated by}\mref{Rem.}{rem:notchannel}, \mref{Thm.}{thm:transfer}
\metasep accepted}

% claim reductions/req:pinched-channel-witness
\claimsummary{reductions/req:pinched-channel-witness}{The artefact's channel witness is the squashing channel followed by the sector pinching. The block-diagonality restriction and every other emitted constraint are then sound for one channel, which a block-diagonal optimisation of the emitted target needs.}
\claimbookmark{reductions/req:pinched-channel-witness}{Requirement}{req:pinched-channel-witness}{The pinched squashing channel as the emitted witness}
\claimleadin{%
Part (iii) of Corollary~\ref{cor:blockdiag} is what makes the emitted target
usable by a block-diagonal optimisation, and it holds for the pinched
channel; Remark~\ref{rem:channel-choice} requires every constraint imposed
together to be sound for one and the same channel.
}
\begin{requirement}[The pinched squashing channel as the emitted witness]\label{req:pinched-channel-witness}
The engine emits $\PinS\circ\Lambda$, the \termref{term-squashing-channel}{squashing channel} followed by the
pinching onto the sectors of Corollary~\ref{cor:blockdiag}, as the
artefact's channel witness (output \iface{squashing_channel}), so that the
block-diagonality restriction of that corollary's part (iii) and every other
emitted constraint are sound for \emph{one} channel.
\end{requirement}
\claimmeta{\metalabel{Motivated by}\mref{Rem.}{rem:channel-choice}, \mref{Cor.}{cor:blockdiag}
\metasep accepted}

% claim reductions/req:carried-tuple-emitted
\claimsummary{reductions/req:carried-tuple-emitted}{Beside each link's channel witness the engine emits the carried tuple of the link's squashing relation, and it builds every emitted constraint from an operator in that tuple. Such a constraint belongs to the relation, not to one realizing channel.}
\claimbookmark{reductions/req:carried-tuple-emitted}{Requirement}{req:carried-tuple-emitted}{Emission of the carried tuple, and constraints built from it}
\claimleadin{%
A constraint built from one of the \termref{term-carried-operators}{carried operators} has the same
pull-back for every channel realizing the \termref{term-squashing-relation}{squashing relation}, while
one read off a particular channel is tied to it
(Remark~\ref{rem:channel-choice}).
}
\begin{requirement}[Emission of the carried tuple, and constraints built from it]\label{req:carried-tuple-emitted}
Beside the channel witness of each link it emits, the engine emits the
carried tuple of the link's squashing relation
(Definition~\ref{def:squashing}; \iface{transported_operators}), and it
builds every constraint it emits from an operator in that tuple.
\end{requirement}
\claimmeta{\metalabel{Motivated by}\mref{Def.}{def:squashing}, \mref{Rem.}{rem:channel-choice}
\metasep accepted}

% claim reductions/req:no-approximate-relations
\claimsummary{reductions/req:no-approximate-relations}{A declared relation carrying an error term is refused, and no link is emitted with the status reserved for approximate squashing. No theorem of this contract converts an error term into a security penalty (\hyperlink{claim-reductions/rem:why-exact}{Remark~\ref*{rem:why-exact}}).}
\claimbookmark{reductions/req:no-approximate-relations}{Requirement}{req:no-approximate-relations}{Refusal of approximate squashing relations}
\claimleadin{%
Theorem~\ref{thm:transfer} rests on an equality of output states and has no
penalty slot, so an error term in a declared relation has nowhere to enter
(Remark~\ref{rem:why-exact}). The refusal of a chain containing a link that
is not exact is Requirement~\ref{req:mixed-status-chain}.
}
\begin{requirement}[Refusal of approximate squashing relations]\label{req:no-approximate-relations}
The engine refuses a declared \termref{term-squashing-relation}{squashing relation} that carries an error
term, with the reason \iface{approximate_squashing_unsupported}, and emits no
link with the status \iface{certified_outer}, which is reserved for
approximate squashing.
\end{requirement}
\claimmeta{\metalabel{Motivated by}\mref{Rem.}{rem:why-exact}, \mref{Thm.}{thm:transfer}, \mref{Req.}{req:mixed-status-chain}
\metasep accepted}

% claim reductions/req:round-homogeneous-emission
\claimsummary{reductions/req:round-homogeneous-emission}{The engine emits one declared model and one squashing channel for every round, built from the one physical POVM that \hyperlink{claim-reductions/ass:round}{A-ROUND-SQ} holds fixed. A round-dependent squasher, for which \hyperlink{claim-reductions/rem:variants}{Remark~\ref*{rem:variants}} proves the factorization, is not emitted.}
\claimbookmark{reductions/req:round-homogeneous-emission}{Requirement}{req:round-homogeneous-emission}{Emission of the round-homogeneous case only}
\claimleadin{%
The factorization holds for round-dependent squashers as well
(Remark~\ref{rem:variants}), but no declared input describes more than one
round.
}
\begin{requirement}[Emission of the round-homogeneous case only]\label{req:round-homogeneous-emission}
The engine emits one \termref{term-declared-model}{declared model} and one squashing channel for
every round, both built from the one physical POVM that
Assumption~\ref{ass:round} holds fixed (the interface convention
\iface{round_homogeneity}). It emits no round-dependent squasher.
\end{requirement}
\claimmeta{\metalabel{Motivated by}\mref{Rem.}{rem:variants}, \mref{Ass.}{ass:round}
\metasep accepted}

% claim reductions/req:post-processing-order
\claimsummary{reductions/req:post-processing-order}{A noise or coarse-graining map on the physical POVM is applied before the first squashing link, and a request to apply one after a link is refused. The refusal reaches these two kinds of map and no other, so a relabelling of both sides after a link stays available.}
\claimbookmark{reductions/req:post-processing-order}{Requirement}{req:post-processing-order}{Noise and coarse-graining before the first squashing link}
\claimleadin{%
A noise or coarse-graining map changes the physical family a squasher is
built from, and a squasher built for one physical family is not a squasher
to the same target for the family the map produces from it
(Remark~\ref{rem:intertwining}).
}
\begin{requirement}[Noise and coarse-graining before the first squashing link]\label{req:post-processing-order}
The engine applies a noise or coarse-graining map on the physical POVM
before the first squashing link, and refuses a request to apply one after a
squashing link (input \iface{post_processing_order_declaration}, refusal
\iface{post_processing_order_violated}). The refusal reaches exactly those
two kinds of map: a stochastic relabelling applied to both sides after a
link, depending on the private component of the outcome only, is not
refused by it (Remark~\ref{rem:intertwining}, \eqref{eq:relabel}).
\end{requirement}
\claimmeta{\metalabel{Motivated by}\mref{Rem.}{rem:intertwining}
\metasep accepted}

% claim reductions/req:declared-relation-check
\claimsummary{reductions/req:declared-relation-check}{Every relation a caller declares is tested against both necessary conditions of \hyperlink{claim-reductions/lem:obstructions}{Lemma~\ref*{lem:obstructions}}, decided exactly on the declared numbers, and refused where one fails. No eigenvalue is computed and no threshold is used.}
\claimbookmark{reductions/req:declared-relation-check}{Requirement}{req:declared-relation-check}{Exact check of a declared relation}
\claimleadin{%
Lemma~\ref{lem:obstructions} rules out an impossible relation by inspection,
and Remark~\ref{rem:obstructiondecide} gives the exact form in which each of
its conditions is decided on declared numbers.
}
\begin{requirement}[Exact check of a declared relation]\label{req:declared-relation-check}
The engine tests every declared relation (input \iface{declared_relations})
against both necessary conditions of Lemma~\ref{lem:obstructions} and
refuses one that fails either, with \iface{declared_relation_impossible}.
With $X$ and $X'$ as in that lemma, it decides (i) in operator form on the
rational enclosures $a\le b$ the declaration exhibits: the pair $a\idon{\Qsys}\preceq X\preceq b\idon{\Qsys}$
exactly by \eqref{eq:psdtest} and, where it holds, the same pair for $X'$;
and (ii) as the exact rank over $\Fldi$ of the matrix whose columns are the
vectorized declared elements, by elimination in the field, together with its
linear relations: an exact basis of the solutions $(a,c)$ of
$\sum_va_vF_v=c\,\idon{\Qsys}$, obtained by elimination in the field, each
basis solution tested on the declared physical elements as
$\sum_va_vM_v=c\,\idon{\Bsys}$. No eigenvalue is computed and no threshold is
used.
\end{requirement}
\claimmeta{\metalabel{Motivated by}\mref{Lem.}{lem:obstructions}, \mref{Rem.}{rem:obstructiondecide}, \mref{Lem.}{lem:exactpsd}
\metasep accepted}

% claim reductions/req:reserved-link-kinds
\claimsummary{reductions/req:reserved-link-kinds}{The engine emits no link of the reserved label-space kind and refuses a request that declares one. The family-dominating kind is the domination submodule's, and no requirement emits it.}
\claimbookmark{reductions/req:reserved-link-kinds}{Requirement}{req:reserved-link-kinds}{Refusal of the reserved receiver-stage link kind}
\claimleadin{%
Proposition~\ref{prop:common} merges branches whose targets squash to one
common target; with the device label read as a branch register, it reduces a
family of devices to a single squashed model
(Remark~\ref{rem:dominating}). The receiver-stage link kind
(Definition~\ref{def:link}) for a general device-label space is reserved. The
kind \texttt{family-dominating}, whose first instance is the \termref{term-domination-link}{domination link}
of Definition~\ref{def:domination-link}, belongs to the domination submodule,
and no requirement emits it.
}
\begin{requirement}[Refusal of the reserved receiver-stage link kind]\label{req:reserved-link-kinds}
The engine emits no receiver-stage link of the reserved kind
\iface{label-space}, and refuses a request that declares a link of that kind
with \iface{link_kind_deferred}.
\end{requirement}
\claimmeta{\metalabel{Motivated by}\mref{Prop.}{prop:common}, \mref{Rem.}{rem:dominating}, \mref{Def.}{def:link}, \mref{Def.}{def:domination-link}
\metasep implemented (effective: accepted)
\metasep \metalabel{Code}\coderef{https://github.com/fraunhofersg/certiqs-key/blob/v0.1.0/src/certiqs_key/reductions/validation.py}{certiqs_key.reductions.validation:rule_link_kind}}

% claim reductions/req:unpadded-branch-blocks
\claimsummary{reductions/req:unpadded-branch-blocks}{The engine does not pad the branch blocks of a composite target: a reduction instance emits the blocks its branches supply. A padded variant a submodule proves from \hyperlink{claim-reductions/cor:flagsenough}{Corollary~\ref*{cor:flagsenough}} is emitted only where that submodule's interface declares it.}
\claimbookmark{reductions/req:unpadded-branch-blocks}{Requirement}{req:unpadded-branch-blocks}{Branch blocks emitted as supplied}
\claimleadin{%
Corollary~\ref{cor:flagsenough} allows the branch blocks of a composite
target to be padded to a common size; this requirement says when the engine
does so.
}
\begin{requirement}[Branch blocks emitted as supplied]\label{req:unpadded-branch-blocks}
The engine does not pad the branch blocks of a composite target: a
\termref{term-reduction-instance}{reduction instance} emits the blocks its branches supply, and a padded variant a
submodule proves from Corollary~\ref{cor:flagsenough} is not emitted unless
that submodule's interface declares it.
\end{requirement}
\claimmeta{\metalabel{Motivated by}\mref{Cor.}{cor:flagsenough}
\metasep accepted}

\section{Requirements: chains, constraint slots and the licence}\label{group-reductions/req:constraint-and-licence-slots}
\sectionintro{Every artefact carries a constraint list and a licence slot, even when both are empty, and its \termref{term-licence-record}{licence record} names the licensing theorem, the certificate of each entry, the claim that discharges it and, where that claim reads a supplier's fact, the supplier artefact. The engine builds no entry whose constant depends on data observed during the run, refuses a chain containing a link that is not exact, and records an identity link where nothing was reduced.}

% claim reductions/req:constraint-and-licence-slots
\claimsummary{reductions/req:constraint-and-licence-slots}{Every artefact carries a constraint list and a licence slot, even when both are empty. A record without them would exclude device knowledge that no classical target can carry.}
\claimbookmark{reductions/req:constraint-and-licence-slots}{Requirement}{req:constraint-and-licence-slots}{Constraint and licence slots on every artefact}
\begin{requirement}[Constraint and licence slots on every artefact]\label{req:constraint-and-licence-slots}
The engine emits every artefact with a constraint list
\iface{constraint_list} and a licence slot \iface{attack_set_license}, even
when both are empty: a record without them excludes a class of device
knowledge that no classical target can carry
(Proposition~\ref{prop:obstruction}, Remark~\ref{rem:obstruction-scope}\,(a)),
and the licence slot is where the licence of Theorem~\ref{thm:licence} is
carried.
\end{requirement}
\claimmeta{\metalabel{Motivated by}\mref{Prop.}{prop:obstruction}, \mref{Rem.}{rem:obstruction-scope}, \mref{Thm.}{thm:licence}
\metasep accepted}

% claim reductions/req:licence-record
\claimsummary{reductions/req:licence-record}{The licence record names the licensing theorem, the certificate per entry, the claim that discharges it and, where that claim reads a supplier's fact, the supplier artefact by version; it is emitted only when every entry has all of these. It also names the proof technique it was issued for.}
\claimbookmark{reductions/req:licence-record}{Requirement}{req:licence-record}{Content and presence of the licence record}
\claimleadin{%
Where the consumer's mechanical check degenerates, its gate rests on this
record alone (Remark~\ref{rem:gateweight}); so the record names a proved
theorem, its hypothesis and the claim that discharges it, never a bare
attestation. Every emitted certificate is a conclusion of a claim
(Corollary~\ref{cor:entries}); an instance with no supplier, such as the
identity instance, has no supplier artefact to name.
}
\begin{requirement}[Content and presence of the licence record]\label{req:licence-record}
The engine emits the \termref{term-licence-record}{licence record} \iface{attack_set_license} naming
the theorem that licenses the constraint list (Theorem~\ref{thm:licence});
the hypothesis it was discharged on, which is the certificate
\eqref{eq:certificate} of each entry; and, per entry, what discharges that
hypothesis: the claim that establishes the certificate under
Corollary~\ref{cor:entries} --- the instance's own claim for an entry of its
list, that corollary for a default carried operator and for the
class-marginal entry, the claim giving the decided certificate for a
requested entry --- and, where that claim reads a fact a supplier's claims
state, the supplier's artefact under Assumption~\ref{ass:recv}, named by its
version. It emits the record only when every entry has its certificate and
what discharges it. The artefact names the proof technique the licence was
issued for, since under another technique the licence must be re-established
(Remark~\ref{rem:notlicensed}\,(e)).
\end{requirement}
\claimmeta{\metalabel{Motivated by}\mref{Thm.}{thm:licence}, \mref{Rem.}{rem:notlicensed}, \mref{Rem.}{rem:gateweight}, \mref{Cor.}{cor:entries}, \mref{Ass.}{ass:recv}
\metasep accepted
\metasep \metalabel{Code}\coderef{https://github.com/fraunhofersg/certiqs-key/blob/v0.1.0/src/certiqs_key/reductions/constraints.py}{certiqs_key.reductions.constraints:issue_licence}}

% claim reductions/req:declared-constants-only
\claimsummary{reductions/req:declared-constants-only}{A constraint entry whose constant is a function of a statistic observed during the run is refused, never built. Such an entry would be unsound regardless of provenance, the licence resting on constants declared before the run.}
\claimbookmark{reductions/req:declared-constants-only}{Requirement}{req:declared-constants-only}{Refusal of data-dependent constants}
\begin{requirement}[Refusal of data-dependent constants]\label{req:declared-constants-only}
The engine builds no constraint entry whose constant $c_k$ is a function of a
statistic observed during the run; a request for one is refused with
\iface{data_dependent_constant}. Such an entry would be unsound regardless of
provenance (Remark~\ref{rem:notlicensed}\,(c)), the constants being declared
before the run (Assumption~\ref{ass:constdecl}).
\end{requirement}
\claimmeta{\metalabel{Motivated by}\mref{Rem.}{rem:notlicensed}, \mref{Ass.}{ass:constdecl}
\metasep accepted}

% claim reductions/req:mixed-status-chain
\claimsummary{reductions/req:mixed-status-chain}{A chain containing a link that is not exact is refused rather than given a composite status. No theorem of this contract converts an error term into a security penalty.}
\claimbookmark{reductions/req:mixed-status-chain}{Requirement}{req:mixed-status-chain}{Refusal of a chain with a link that is not exact}
\begin{requirement}[Refusal of a chain with a link that is not exact]\label{req:mixed-status-chain}
The engine assigns a chain no composite status other than the one of
Theorem~\ref{thm:chain}\,(iv): a chain containing a link that is not
\textsf{exact} is refused, with
\iface{approximate_squashing_unsupported}, rather than given a composite
status, because no theorem of this contract converts an error term into a
security penalty. The refusal is a rule of the engine and not a converse
this contract proves.
\end{requirement}
\claimmeta{\metalabel{Motivated by}\mref{Thm.}{thm:chain}
\metasep implemented (effective: accepted)
\metasep \metalabel{Code}\coderef{https://github.com/fraunhofersg/certiqs-key/blob/v0.1.0/src/certiqs_key/reductions/validation.py}{certiqs_key.reductions.validation:rule_approximate}}

% claim reductions/req:identity-link-recorded
\claimsummary{reductions/req:identity-link-recorded}{Where nothing is reduced, the artefact carries the identity link rather than an empty chain. An empty chain could not say whether a reduction was performed and found trivial or never considered.}
\claimbookmark{reductions/req:identity-link-recorded}{Requirement}{req:identity-link-recorded}{Identity link recorded for an already-finite receiver}
\begin{requirement}[Identity link recorded for an already-finite receiver]\label{req:identity-link-recorded}
The engine records the reduction of an already-finite receiver as the
identity link of Example~\ref{ex:identity}, never as an empty chain: an empty
chain would leave the artefact unable to say whether a reduction was performed
and found trivial or never considered, and the composition order the chain
carries as data (Definition~\ref{def:chain}) has to be carried even when there
is nothing to order.
\end{requirement}
\claimmeta{\metalabel{Motivated by}\mref{Def.}{def:chain}, \mref{Thm.}{thm:chain}, \mref{Ex.}{ex:identity}
\metasep accepted}

\section{Requirements: constraint entries}\label{group-reductions/req:certified-entries-only}
\sectionintro{The engine emits only certified constraint entries, each with its recorded certificate case, as the request decides, and carries the certificates of the \termref{term-transported-set}{transported set} beside them. The \termref{term-class-marginal-identity}{class-marginal identity} is emitted only in expectation form and only where the instance discharges it; its operator form, and any requested operator whose certificate no claim decides on the emitted numbers, are refused. A caller's inequality constant off the lattice is rounded outward and an equality constant off it is refused; a party-local entry carries its tag, and every entry of an emitted constraint operator satisfies the declared lattice condition, its modulus bound included.}

% claim reductions/req:certified-entries-only
\claimsummary{reductions/req:certified-entries-only}{Every constraint entry the engine emits carries exactly one certificate, named by its case of \hyperlink{claim-reductions/cor:pullback}{Corollary~\ref*{cor:pullback}}, and is of the certified kind of \hyperlink{claim-reductions/cor:transported}{Corollary~\ref*{cor:transported}}, never an assumption on the input. The artefact records each entry's certificate case.}
\claimbookmark{reductions/req:certified-entries-only}{Requirement}{req:certified-entries-only}{Emission of certified constraint entries only}
\claimleadin{%
A certified family is sound for every family of \termref{term-physical-inputs}{physical inputs},
while an assumed one is sound only for the family it was assumed on
(Corollary~\ref{cor:transported}); an emitted entry therefore never rests on
an assumption about the input (Remark~\ref{rem:unsound} gives the standard
way a constraint fails to be sound).
}
\begin{requirement}[Emission of certified constraint entries only]\label{req:certified-entries-only}
Every constraint entry the engine emits carries exactly one
\termref{term-certificate}{certificate}, named by its case of Corollary~\ref{cor:pullback}
(identity, positivity or spectral bounds), and is of the certified kind, case
(i) of Corollary~\ref{cor:transported}; the engine emits no entry of the
assumed kind, case (ii). Each emitted entry's certificate case is recorded
with the \termref{term-transported-set}{transported set}, as
Requirement~\ref{req:transported-set-certificates} states.
\end{requirement}
\claimmeta{\metalabel{Motivated by}\mref{Cor.}{cor:pullback}, \mref{Cor.}{cor:transported}, \mref{Rem.}{rem:unsound}, \mref{Req.}{req:transported-set-certificates}
\metasep accepted}

% claim reductions/req:class-marginal-operator-form
\claimsummary{reductions/req:class-marginal-operator-form}{The engine refuses a request to impose the class-marginal identity on the target as an operator identity, a form \hyperlink{claim-reductions/rem:unsound}{Remark~\ref*{rem:unsound}} shows unsound. Its sound squashed form is the expectation identity.}
\claimbookmark{reductions/req:class-marginal-operator-form}{Requirement}{req:class-marginal-operator-form}{Refusal of the operator form of the class-marginal identity}
\claimleadin{%
Where the physical \termref{term-class-marginal-identity}{class-marginal identity} holds, its squashed counterpart
holds in expectation on every reachable state but not, in general, as an
operator identity on the target, and imposing the operator form where the
target does not satisfy it would exclude reachable states
(Remark~\ref{rem:unsound}).
}
\begin{requirement}[Refusal of the operator form of the class-marginal identity]\label{req:class-marginal-operator-form}
The engine refuses, with the reason \iface{unsound_constraint_requested}, a
request to impose on the squashed states the operator form
$\sum_{v\in\Vset_t}F_v=q_t\,\idon{\Qsys}$ of the class-marginal identity
(Remark~\ref{rem:unsound}).
\end{requirement}
\claimmeta{\metalabel{Motivated by}\mref{Rem.}{rem:unsound}
\metasep accepted}

% claim reductions/req:constraint-list-by-request
\claimsummary{reductions/req:constraint-list-by-request}{Where the request carries no entry the engine emits an empty constraint list and no licence; where it carries entries it emits the instance's entries and the requested ones its claims admit, with the licence. A licence is present exactly when the list is non-empty.}
\claimbookmark{reductions/req:constraint-list-by-request}{Requirement}{req:constraint-list-by-request}{The constraint list and the licence decided by the request}
\claimleadin{%
A non-empty list without a licence would carry no key-length claim, and a
licence without a list would attest to nothing; which case holds is the
request's, as Corollary~\ref{cor:entries} states.
}
\begin{requirement}[The constraint list and the licence decided by the request]\label{req:constraint-list-by-request}
Where the request carries no entry (Corollary~\ref{cor:entries}), the engine
emits an empty \iface{constraint_list}, whatever entries the instance's own
list contains, and no \iface{attack_set_license}. Where the request carries
entries, it emits in \iface{constraint_list} the instance's own entries, the
class-marginal entry on request and any further requested entry that
Corollary~\ref{cor:entries}\,(iii) admits
(Requirement~\ref{req:uncertified-operator-refused}), and it emits
the licence of Theorem~\ref{thm:licence} in \iface{attack_set_license}. The
licence is present if and only if the list is non-empty
(\iface{licence_presence}): a non-empty list without a licence is refused
with \iface{licence_missing}, and a licence without a list with
\iface{licence_without_constraints}.
\end{requirement}
\claimmeta{\metalabel{Motivated by}\mref{Thm.}{thm:licence}, \mref{Cor.}{cor:entries}, \mref{Req.}{req:uncertified-operator-refused}
\metasep implemented (effective: accepted)
\metasep \metalabel{Code}\coderef{https://github.com/fraunhofersg/certiqs-key/blob/v0.1.0/src/certiqs_key/reductions/validation.py}{certiqs_key.reductions.validation:rule_licence_presence}}

% claim reductions/req:transported-set-certificates
\claimsummary{reductions/req:transported-set-certificates}{Whatever the request carries, the engine emits the transported set with the certificate case of each operator. The positivity certificate of every sector projector travels there, and becomes a constraint entry only where a caller requests it as one.}
\claimbookmark{reductions/req:transported-set-certificates}{Requirement}{req:transported-set-certificates}{Certificates carried with the transported set}
\claimleadin{%
Every instance transports its sector projectors, and their pull-backs are
positive for every instance (Corollary~\ref{cor:entries}\,(i)); that
certificate is a property of the channel and does not wait for a request.
}
\begin{requirement}[Certificates carried with the transported set]\label{req:transported-set-certificates}
Whatever the request carries, the engine emits in
\iface{transported_operators}, per operator of the \termref{term-transported-set}{transported set},
the certificate case of Corollary~\ref{cor:pullback} that makes it usable. In
particular the certificate $(\Pi_s,\succeq,0)$ of every sector projector of
every instance travels there and not as a constraint entry; such an operator
becomes an entry of \iface{constraint_list} only where a caller requests it
as an entry (Corollary~\ref{cor:entries}\,(i)).
\end{requirement}
\claimmeta{\metalabel{Motivated by}\mref{Cor.}{cor:pullback}, \mref{Cor.}{cor:entries}
\metasep implemented (effective: accepted)
\metasep \metalabel{Code}\coderef{https://github.com/fraunhofersg/certiqs-key/blob/v0.1.0/src/certiqs_key/reductions/channels.py}{certiqs_key.reductions.channels:TransportedOperator}}

% claim reductions/req:class-marginal-entry
\claimsummary{reductions/req:class-marginal-entry}{The engine emits the class-marginal entry only in expectation form, on request, and only where the instance discharges the class-marginal identity; a request for it anywhere else is refused. The entry is never emitted on the strength of a branch that does not establish it.}
\claimbookmark{reductions/req:class-marginal-entry}{Requirement}{req:class-marginal-entry}{The class-marginal entry only where the instance discharges it}
\claimleadin{%
Emitting the entry where the identity is not discharged would impose it on
reachable states for which it may be false, which is the security break of
Remark~\ref{rem:unsound}; its operator form is refused outright
(Requirement~\ref{req:class-marginal-operator-form}).
}
\begin{requirement}[The class-marginal entry only where the instance discharges it]\label{req:class-marginal-entry}
The engine emits the class-marginal entry $\big(\sum_{v\in\Vset_t}F_v,=,q_t\big)$
at a declared $t$, in the expectation form of
Corollary~\ref{cor:entries}\,(ii), only on request and only where the
instance discharges the \termref{term-class-marginal-identity}{class-marginal identity}
\eqref{eq:instance-classmarginal}. A request for it of an instance that does
not discharge the identity is refused with
\iface{constraint_without_certificate}, and the entry is never emitted on the
strength of a branch that does not establish it.
\end{requirement}
\claimmeta{\metalabel{Motivated by}\mref{Rem.}{rem:unsound}, \mref{Def.}{def:reduction-instance}, \mref{Cor.}{cor:entries}, \mref{Req.}{req:class-marginal-operator-form}
\metasep implemented (effective: accepted)
\metasep \metalabel{Code}\coderef{https://github.com/fraunhofersg/certiqs-key/blob/v0.1.0/src/certiqs_key/reductions/validation.py}{certiqs_key.reductions.validation:rule_constraint_certificate}}

% claim reductions/req:caller-constant-rounding
\claimsummary{reductions/req:caller-constant-rounding}{A caller's inequality constant off the declared lattice is rounded onto it in the direction that enlarges the feasible set, and the direction is recorded. An equality constant off the lattice, or any constant of modulus above one, is refused rather than rounded.}
\claimbookmark{reductions/req:caller-constant-rounding}{Requirement}{req:caller-constant-rounding}{Outward rounding of a caller's constant off the lattice}
\claimleadin{%
Rounding toward the enlarging side keeps the certificate true at the rounded
constant, while the other side could exclude a reachable state
(Corollary~\ref{cor:entries}\,(iii)); an equality has no enlarging side.
}
\begin{requirement}[Outward rounding of a caller's constant off the lattice]\label{req:caller-constant-rounding}
Where the declared constant of a requested entry does not satisfy
\eqref{eq:latinv}, the engine first rounds it onto the lattice in the
direction that enlarges the feasible set for the entry's relation, down for
$\succeq$ and up for $\preceq$, and records the direction taken per entry in
\iface{emission_record} (\xref{foundations/req:rounding-direction}). An
equality entry whose constant is off the lattice is refused with
\iface{constraint_without_certificate} rather than split into a pair, and a
declared constant of modulus above one is refused with
\iface{constraint_without_certificate} under every relation, before any
rounding.
\end{requirement}
\claimmeta{\metalabel{Motivated by}\mref{Def.}{def:exactadm}, \mref{Cor.}{cor:entries}, \extref{../foundations/contract.pdf}{foundations/req:rounding-direction}
\metasep implemented (effective: accepted)
\metasep \metalabel{Code}\coderef{https://github.com/fraunhofersg/certiqs-key/blob/v0.1.0/src/certiqs_key/reductions/constraints.py}{certiqs_key.reductions.constraints:rounded_constant}}

% claim reductions/req:uncertified-operator-refused
\claimsummary{reductions/req:uncertified-operator-refused}{A requested operator is emitted only where a claim gives a condition on the emitted numbers implying its certificate and the engine decides that condition exactly. Every other requested operator is refused, and none is emitted at some other constant or on the caller's word.}
\claimbookmark{reductions/req:uncertified-operator-refused}{Requirement}{req:uncertified-operator-refused}{Refusal of a requested operator without a decided certificate}
\claimleadin{%
Parts (i)--(iii) of Corollary~\ref{cor:entries} cover every entry the module
can emit. A requested operator they do not cover has no certificate to be
emitted on, and a certificate the caller declares is not one: nothing in this
contract decides it, and no assumption of the ledger names it as a premise,
while an entry on a false certificate excludes reachable states, the
security break of Remark~\ref{rem:unsound}.
}
\begin{requirement}[Refusal of a requested operator without a decided certificate]\label{req:uncertified-operator-refused}
The engine emits an operator that a request carries as an entry only where
Corollary~\ref{cor:entries} covers it: as a default carried operator (part
(i)), as the class-marginal entry of an instance that discharges the
class-marginal identity (part (ii)), or, under part (iii), where a claim
stated for the instance gives a decided certificate for the operator's form
and the engine decides that condition exactly on the emitted numbers and the
request's declared constants, finding that it holds. It refuses, with
\iface{constraint_without_certificate}, every other requested operator: one
outside the form of Corollary~\ref{cor:entries}\,(i)--(iii), one of that form
for which no claim stated for the instance gives a decided certificate, and
one whose decided certificate fails at the emitted numbers. It takes no
certificate \eqref{eq:certificate} from the request on trust, and it never
emits a requested operator at some other constant.
\end{requirement}
\claimmeta{\metalabel{Motivated by}\mref{Rem.}{rem:unsound}, \mref{Thm.}{thm:licence}, \mref{Cor.}{cor:entries}
\metasep implemented (effective: accepted)
\metasep \metalabel{Code}\coderef{https://github.com/fraunhofersg/certiqs-key/blob/v0.1.0/src/certiqs_key/reductions/validation.py}{certiqs_key.reductions.validation:rule_constraint_certificate}}

% claim reductions/req:party-local-tag
\claimsummary{reductions/req:party-local-tag}{Every emitted entry carries its reading as a tag: untagged in the form of the source mode, or party-local with a factor scope only on a declared product with a certificate on that factor's channel. A party-local request on a one-factor emission gets the untagged entry.}
\claimbookmark{reductions/req:party-local-tag}{Requirement}{req:party-local-tag}{The party-local tag and its scope}
\claimleadin{%
The tag makes a consumer's gate a type check that compares no numbers, and a
party-local reading is covered by the licence only through a certificate on
its own factor's channel (Remark~\ref{rem:entry-tags}).
}
\begin{requirement}[The party-local tag and its scope]\label{req:party-local-tag}
The engine emits every constraint entry with the reading it is to be read in.
An entry is emitted untagged, as $(\hat O,\bowtie,c)$, and is read in the
form of the declared source mode (Definition~\ref{def:licensed}); or, only
where the \termref{term-measured-system}{measured system} is a declared product
$\Qsys=\bigotimes_f\Qsys_f$ of party-scoped factors and the entry's
certificate is discharged on that factor's own channel $\Lambda_f$
(Theorem~\ref{thm:licence}\,(v)), it is emitted with the tag
\iface{party_local} and a scope, as $(\hat O,\bowtie,c;f)$, with $\hat O$ an
operator on $\Qsys_f$ alone carried as a sector-block list against that
factor's \exttermref{../foundations/contract.pdf}{term-sector-map}{sector map}. In this module the tag is emitted only with the
entries of several \termref{term-measuring-parties}{measuring parties}
(Theorem~\ref{thm:emitted-product}\,(iii)).
A request for the \termref{term-party-local-form}{party-local form} on a one-factor emission is answered
with the untagged entry, at the same operator, relation and constant, and is
not refused.
\end{requirement}
\claimmeta{\metalabel{Motivated by}\mref{Def.}{def:licensed}, \mref{Thm.}{thm:licence}, \mref{Rem.}{rem:entry-tags}, \mref{Thm.}{thm:emitted-product}
\metasep implemented (effective: accepted)
\metasep \metalabel{Code}\coderef{https://github.com/fraunhofersg/certiqs-key/blob/v0.1.0/src/certiqs_key/reductions/constraints.py}{certiqs_key.reductions.constraints:party_local_entry}}

% claim reductions/req:licensed-entry-lattice
\claimsummary{reductions/req:licensed-entry-lattice}{Where the licence is present, every entry of every emitted constraint operator is decided to satisfy the lattice condition before the licence is issued, and the request is refused otherwise. The modulus bound binds these entries, not the coefficients that form them.}
\claimbookmark{reductions/req:licensed-entry-lattice}{Requirement}{req:licensed-entry-lattice}{Lattice condition on the entries of emitted constraint operators}
\claimleadin{%
A certified entry is a combination of transported operators whose
coefficients need only be integer multiples of $2^{-\kappa}$
(Definition~\ref{def:reduction-instance}\,(I4)). Condition
\eqref{eq:latinv}, the modulus bound $|x|\le1$ included, binds the entries of
the operator the combination produces, and those are emitted numbers.
}
\begin{requirement}[Lattice condition on the entries of emitted constraint operators]\label{req:licensed-entry-lattice}
Where the licence is present, the engine decides, on the emitted numbers and
before the licence is issued, that every entry of every emitted constraint
operator satisfies \eqref{eq:latinv}, and refuses the request
otherwise, with the reason \iface{constraint_without_certificate} (the lattice
half of the rule \iface{licensed_hermiticity};
Corollary~\ref{cor:entries}\,(iii)).
\end{requirement}
\claimmeta{\metalabel{Motivated by}\mref{Def.}{def:exactadm}, \mref{Def.}{def:reduction-instance}, \mref{Cor.}{cor:entries}
\metasep implemented (effective: accepted)
\metasep \metalabel{Code}\coderef{https://github.com/fraunhofersg/certiqs-key/blob/v0.1.0/src/certiqs_key/reductions/validation.py}{certiqs_key.reductions.validation:rule_licensed_hermiticity}}

\section{Requirements: exact admission and emission}\label{group-reductions/req:exact-admission-refusal}
\sectionintro{The engine refuses a declaration that is not exactly representable in the field of its arithmetic, and an emitted family it cannot decide exactly admissible. A class the normalizer moves, the flagged last resort included, is disclosed rather than refused, unless the run declares the strict option and would rest on it. It records every displacement, kept or moved class and rounding direction of an emission, emits the reference companion beside the artefact where a reference is handed, with no normalizer and its entries relaxed, emits each \termref{term-reduction-instance}{reduction instance}'s party partition with its sector map, and emits several \termref{term-measuring-parties}{measuring parties}' artefacts with their combination rule rather than a product.}
\sectionimports{\extref{../foundations/contract.pdf}{foundations/def:evaluated-off-snap}, \extref{../foundations/contract.pdf}{foundations/def:exact-island}, \extref{../foundations/contract.pdf}{foundations/def:factor-enclosure}, \extref{../foundations/contract.pdf}{foundations/lem:recentred-enclosure}, \extref{../foundations/contract.pdf}{foundations/req:snap-reference}}

% claim reductions/req:exact-admission-refusal
\claimsummary{reductions/req:exact-admission-refusal}{A declaration whose entries are not exactly representable in the field chosen for the arithmetic is refused, never rounded into it. A declaration on the declared lattice is never refused on this ground.}
\claimbookmark{reductions/req:exact-admission-refusal}{Requirement}{req:exact-admission-refusal}{Refusal of a declaration outside the substrate}
\claimleadin{%
Declared operator data cross every \exttermref{../foundations/contract.pdf}{term-module-boundary}{module boundary} \exttermref{../foundations/contract.pdf}{term-exactly-admissible}{exactly admissible}
(\xref{foundations/req:boundary-admission}), and a consumer that finds one of
its admission predicates violated refuses rather than repairs
(\xref{foundations/req:consumer-decides-exactly}). For this module one
such predicate is that the declared entries lie in the field that carries its
arithmetic.
}
\begin{requirement}[Refusal of a declaration outside the substrate]\label{req:exact-admission-refusal}
The engine refuses a declaration whose entries are not exactly representable
in the field $\Fld$ chosen for the arithmetic
(Definition~\ref{def:substrate}), with the reason
\iface{declaration_not_exactly_admissible}, rather than rounding the entries
into that field. No declaration whose entries satisfy \eqref{eq:latinv} is
refused on this ground.
\end{requirement}
\claimmeta{\metalabel{Motivated by}\mref{Def.}{def:exactadm}, \mref{Def.}{def:substrate}, \extref{../foundations/contract.pdf}{foundations/req:boundary-admission}, \extref{../foundations/contract.pdf}{foundations/req:consumer-decides-exactly}
\metasep accepted}

% claim reductions/req:emission-refusal
\claimsummary{reductions/req:emission-refusal}{The engine emits a family only after deciding the six emission conditions on the numbers about to be emitted, and otherwise refuses, repairing nothing it was handed. A moved class, the flagged last resort included, is refused only under the strict option, where it would rest on the snap.}
\claimbookmark{reductions/req:emission-refusal}{Requirement}{req:emission-refusal}{Refusal of an emitted family that is not exactly admissible}
\claimleadin{%
The decision on the six conditions is taken on what comes out of the
normalizer (Definition~\ref{def:emitnorm}\,(4), Theorem~\ref{thm:emit}\,(ii)).
A class the normalizer moves is not refused by default: it is snapped without
reviving a zero row or column and disclosed
(Requirement~\ref{req:emission-record}), as
\xref{foundations/req:snap-reference}\,(iii) asks of any snap of a
measurement family. A class whose residue no single element can take
without reviving a zero row or column has no support-safe snap in the
normalizer; it is snapped by the last resort, which always yields a positive
semidefinite family and is flagged in the record
(\xref{foundations/req:support-safe-snap}\,(iv)), so it is not refused by
default either. The usual cause of a last resort is a class whose elements
are rank-deficient on the sector, which a model with the device's loss and
dark counts can remove where the conditions the receiver module states for it
hold (\xref{receiver/detector-imperfections/prop:noise-floor}). One refusal
remains for a moved class: a run that declares the strict option refuses
wherever the class would rest on the snapped-measurement premise, which is
certain at emission where no reference companion that bypasses the
normalizer is formed (Requirement~\ref{req:reference-companion}).
}
\begin{requirement}[Refusal of an emitted family that is not exactly admissible]\label{req:emission-refusal}
The engine emits a family only if the six conditions of
Theorem~\ref{thm:emit}\,(ii) --- (a)--(e) of
Definition~\ref{def:exactadm} together with \eqref{eq:latinv} of every entry
--- hold of the numbers about to be emitted, decided exactly after the
\termref{term-emission-normalizer}{emission normalizer} where that acted; otherwise it refuses with the
reason \iface{emitted_family_not_admissible}, naming the predicate that
failed, and repairs nothing it was handed. Under the same reason, where the
run declares the strict option \iface{strict_snaps}
(\xref{foundations/req:snap-reference}\,(iv)) and no reference companion of
Requirement~\ref{req:reference-companion} is formed beside the artefact, it
refuses every moved class of Definition~\ref{def:emitnorm}, one snapped by
the last resort (3) included, naming the \termref{term-data-sector}{data sector}, the class, the
cause and the fix; where a companion is formed, whether the class rests is
decided by the run after its search, and the run refuses there under the
same option. Where the run does not declare the option, no class is refused
for being moved. The cause of a moved class that would rest is that no
reference companion bypasses the normalizer, which, by
Requirement~\ref{req:reference-companion}, is where no claim of this module
establishes (H-REF) for the instance; for a class snapped by (3) it also
names that a component of its residue has no host. The fix named is
declared numbers on which the exact construction stays on the lattice, so
that the class is kept; for a class
snapped by (3) also a model that includes the device's loss and dark counts,
so that the elements of the class are positive definite on the sector. A
moved class is mixed only with the family of its diagonal parts in (2), or
with a multiple of the identity in the last resort (3).
\end{requirement}
\claimmeta{\metalabel{Motivated by}\mref{Def.}{def:exactadm}, \mref{Def.}{def:emitnorm}, \mref{Thm.}{thm:emit}, \mref{Req.}{req:emission-record}, \mref{Req.}{req:reference-companion}, \extref{../foundations/contract.pdf}{foundations/req:snap-reference}, \extref{../foundations/contract.pdf}{foundations/req:support-safe-snap}, \extref{../receiver/contract.pdf}{receiver/detector-imperfections/prop:noise-floor}
\metasep implemented (effective: accepted)
\metasep \metalabel{Code}\coderef{https://github.com/fraunhofersg/certiqs-key/blob/v0.1.0/src/certiqs_key/reductions/validation.py}{certiqs_key.reductions.validation:rule_emitted_family}}

% claim reductions/req:emission-record
\claimsummary{reductions/req:emission-record}{Every emission records the lattice exponent and the strict option, each displacement, the kept and moved classes with the step that snapped each, and every rounding direction. It names each moved sector, which rests on the snapped-measurement premise unless evaluated off its snap.}
\claimbookmark{reductions/req:emission-record}{Requirement}{req:emission-record}{The record of an emission}
\claimleadin{%
Where the normalizer moved the datum, nothing proved relates the
\termref{term-declared-model}{declared model} to the \termref{term-handed-model}{handed model}, and the size of the move is
not such a relation (Theorem~\ref{thm:emit}\,(v),
Remark~\ref{rem:emitclaim}); the record carries the move so that it is visible
to the consumer. A run that forms the number its key length is formed
from on the reference companion of Requirement~\ref{req:reference-companion},
which carries no data of this normalizer, evaluates off its snap every data
sector that holds a moved class
(\xref{foundations/def:evaluated-off-snap}); elsewhere a moved sector rests on
the premise \xref{foundations/ass:snapped-measurement}, and the record says so
in the form \xref{foundations/req:snap-reference}\,(iii) asks, so that a
report on the key rate can flag it, as the record of a snap names every block
that rests on the premise (\xref{foundations/req:who-snaps}). A class snapped
by the last resort carries its flag, since its elements may treat as possible
an outcome the \termref{term-handed-datum}{handed datum} makes impossible
(\xref{foundations/req:support-safe-snap}\,(iv)).
}
\begin{requirement}[The record of an emission]\label{req:emission-record}
The engine emits with every artefact, in \iface{emission_record}: the
declared \exttermref{../foundations/contract.pdf}{term-lattice-exponent}{lattice exponent} $\kappa$; whether the run declared the strict
option of \xref{foundations/req:snap-reference}\,(iv); the reference kind of
every snap the normalizer makes, \emph{exact}, every radius zero, the
reference being the \termref{term-handed-datum}{handed datum} itself
(Definition~\ref{def:emitnorm}), so that the enclosure setting, recorded as
\xref{foundations/req:who-snaps} asks of the record of every snap, does not
bear on it; per \termref{term-data-sector}{data sector}, the \termref{term-emission-displacement}{emission displacement}
$\Delta_s$ of \eqref{eq:displacement}; per data sector and class, whether the
class was kept (Definition~\ref{def:emitnorm}\,(1)) or moved, and for a moved
class the step (2) or (3) of Definition~\ref{def:emitnorm} that snapped it,
the weight $w$ taken or the fall-back of (3), the hosts of the components of
the residue under (2), and for (3) the flag \iface{support_not_preserved}; and
the direction of every bound it rounds, each \termref{term-scalar-value}{scalar value} rounded
down and each complementary value recomputed (Theorem~\ref{thm:emit}\,(iv)).
A data sector \emph{rests on} \xref{foundations/ass:snapped-measurement}
where some class of it is moved, save where it is
\exttermref{../foundations/contract.pdf}{term-evaluated-off-its-snap}{evaluated off its snap} (\xref{foundations/def:evaluated-off-snap}):
in a run whose key length is formed from a number on the reference
companion of Requirement~\ref{req:reference-companion}, its first-order
bound or a number certified on it, which carries no data of this emission's
normalizer, every data sector that holds a moved class is, a kept sector
being named as kept and not counted among them, and none rests on the
assumption. Where the key length is formed instead from a number certified
on the snapped declaration in the fallback of
\xref{meat-core/frank-wolfe/req:certified-evaluation}\,(i), none is, and
every data sector that holds a moved class rests on the assumption, with the
cause \iface{companion_not_certified}
(\xref{foundations/req:snap-reference}\,(iii)(b)). The record names every moved sector, with the flag of any class
of it snapped by (3), and carries
\xref{foundations/ass:snapped-measurement} among the assumptions the
artefact may rest on exactly where there is one. This record is emitted with
the artefact, before the run knows whether the number its key length is
formed from is formed on the reference companion; what it names is therefore provisional, and the
record of the run decides which sectors rest on the assumption, each with
its cause and group (\xref{foundations/req:snap-reference}\,(iii);
Theorem~\ref{thm:emit}\,(vi)).
\end{requirement}
\claimmeta{\metalabel{Motivated by}\mref{Def.}{def:emitnorm}, \mref{Thm.}{thm:emit}, \mref{Rem.}{rem:emitclaim}, \mref{Req.}{req:reference-companion}, \extref{../foundations/contract.pdf}{foundations/def:evaluated-off-snap}, \extref{../foundations/contract.pdf}{foundations/req:snap-reference}, \extref{../foundations/contract.pdf}{foundations/req:support-safe-snap}, \extref{../foundations/contract.pdf}{foundations/req:who-snaps}, \extref{../meat-core/contract.pdf}{meat-core/frank-wolfe/req:certified-evaluation}
\metasep \metalabel{Assumes}\extaref{../foundations/contract.pdf}{foundations/ass:snapped-measurement}{A-SNAP}
\metasep implemented (effective: accepted)
\metasep \metalabel{Code}\coderef{https://github.com/fraunhofersg/certiqs-key/blob/v0.1.0/src/certiqs_key/reductions/artifact.py}{certiqs_key.reductions.artifact:EmissionRecord}}

% claim reductions/req:reference-companion
\claimsummary{reductions/req:reference-companion}{Where a reference is handed or the normalizer moved a class, the engine runs the same construction exactly on the reference, or on the snapped blocks, with no normalizer, and emits the companion's enclosures, factor ones where carried, and its entries, requested ones where certified there.}
\claimbookmark{reductions/req:reference-companion}{Requirement}{req:reference-companion}{Emitting the reference companion beside the artefact}
\claimleadin{%
The companion is what a run evaluates its first-order bound on
(\xref{meat-core/frank-wolfe/req:bracket-on-reference}). It is the same
chain run on other data, so it is built where the artefact is built, by the
same code path, and carries the same shape by construction. It keeps the
reference's exact data where the reference has them, as Gaussian-rational
centres and rational radii, and is rounded to binary64 only by whoever
evaluates it; its entries, which a solver's oracle reads as constraints, are
the one part emitted in binary64 here, and Lemma~\ref{lem:row-relaxation}
says how they are rounded without enlarging what they assert. Where the
instance is one whose (H-REF) a claim establishes, Corollary~\ref{cor:reference-companion}\,(iv)
or Corollary~\ref{cor:flat-companion}, the companion is emitted. It is
emitted also where the declaring instance hands no reference but the
normalizer moved a class: the companion is then the instance at its handed
datum, the module's exact construction, which no step of the normalizer has
touched, so that the classes the normalizer moved are evaluated off their
snap while the instance's own blocks stay snapped. An entry a caller
requested beside the instance's own is carried to the companion in its
companion form where its certificate is established at the companion's
datum, and is left out where it is not: leaving an entry out removes one
constraint and so only enlarges the companion's feasible set, and no bound
formed on the companion is raised by it. Where the reference also carries
\exttermref{../foundations/contract.pdf}{term-factor-enclosure}{factor enclosures} and the instance's rule combines blocks with non-negative
weights, the companion carries them to its datum and target
(Corollary~\ref{cor:companion-factor-enclosure}), unaltered but for the
weights.
}
\begin{requirement}[Emitting the reference companion beside the artefact]\label{req:reference-companion}
Where a claim of this module establishes (H-REF) of
Corollary~\ref{cor:reference-companion} for the \termref{term-reduction-instance}{reduction instance} the
artefact is emitted from, and either the declaring instance hands the
reference beside its snap of the measurement it states
(\xref{foundations/def:evaluated-off-snap}\,(i)) on at least one block or
the \termref{term-emission-normalizer}{emission normalizer} moves at least one class
(Definition~\ref{def:emitnorm}), the engine emits beside the artefact, in
\iface{reference_companion}, the reference companion of that instance at the
family $G$ that carries the \exttermref{../foundations/contract.pdf}{term-handed-reference}{handed reference} on every handed block and the
snapped block the instance declares on every other block, so that, where no
block is handed, $G$ is the instance's snapped family and the companion's
datum is the \termref{term-handed-datum}{handed datum}
(Corollary~\ref{cor:reference-companion}\,(i)):
\begin{enumerate}[label=(\roman*),leftmargin=2.4em]
\item \emph{The datum.} Every enclosure of a handed block is first
  recentred by \xref{foundations/lem:recentred-enclosure}, so that every
  entry whose enclosure contains $0$ has centre $0$. The engine then forms,
  in exact arithmetic on the \exttermref{../foundations/contract.pdf}{term-exact-island}{exact island}
  (\xref{foundations/def:exact-island}), the enclosure of the companion's
  datum that Corollary~\ref{cor:reference-companion}\,(iii) gives, from these
  enclosures and from the snapped blocks with radius zero, by the
  construction it applies to the snapped family: the same rule of (I1), the
  same \termref{term-weight-route}{weight route}, scalar values, sectors and supports. It then
  recentres that enclosure again, element by element and block by block,
  by \xref{foundations/lem:recentred-enclosure}, after making its centre
  Hermitian and its radii symmetric by part (iv) of that lemma where they
  are not, every block of the datum being Hermitian
  (Corollary~\ref{cor:reference-companion}\,(i)); where the rule of (I1)
  combines entries of several blocks, a sum that is exactly zero may so
  have a centre that is not, and this second recentring sets it to zero
  wherever its enclosure contains zero. The entries of (iii) are formed
  from the datum's enclosure so recentred. It applies no
  normalizer (Definition~\ref{def:emitnorm}), decides no lattice condition,
  rounds nothing, and decides nothing of the companion's admissibility,
  which is Corollary~\ref{cor:reference-companion}\,(i) on the claims the
  handed reference names (\xref{foundations/req:snap-reference}\,(i)). It
  emits the centres, entries of $\mathbb{Q}(i)$, and the radii, rationals,
  entry by entry, so that an entry the reference knows exactly is emitted
  exactly, with radius zero.
\item \emph{The shape.} The companion is emitted with the \exttermref{../foundations/contract.pdf}{term-sector-map}{sector map}, the
  alphabet with its classes and weights, the scalar values and the party
  partition of the artefact's target, which are equal by
  Corollary~\ref{cor:reference-companion}\,(i), and with the artefact's
  \termref{term-channel-witness}{channel witness} and \termref{term-transported-set}{transported set} read at the companion's datum.
\item \emph{The entries.} On a route whose request carries entries
  (Corollary~\ref{cor:entries}), the companion's entry list holds the
  instance's own entries at the companion's datum and, for every entry
  requested beside them that the artefact carries, its companion form
  (Corollary~\ref{cor:reference-companion}\,(vi)) at the writing by which
  the engine recognised it, where its certificate at the companion's datum
  is established: for a sum of sector projectors and for the class-marginal
  entry, by Corollary~\ref{cor:reference-companion}\,(vi), the operator
  being the artefact's, enclosed with radius zero; for another entry, where
  a claim of this module states for the instance a condition on an
  enclosure of the companion's datum under which its certificate holds,
  by deciding that condition exactly on the datum's enclosure recentred as
  (i) states (for the operator \eqref{eq:two-sector-entry} on the flat instance's
  licensed route, \eqref{eq:companion-two-sector} of
  Corollary~\ref{cor:flat-companion}). A requested entry whose certificate
  at the companion's datum is not so established is not formed on the
  companion and has no entry there in its place. For each entry
  $(\hat O_k,\bowtie_k,c_k)$ of that list, with the enclosure $(C_k,R_k)$
  of its operator that Corollary~\ref{cor:reference-companion}\,(iii)
  gives, or radius zero where the operator is the artefact's: the operator
  $\hat O'_k$ whose every real and imaginary part is the binary64 number
  nearest to that of $C_k$, ties to even, which is Hermitian and has the
  blocks of $\hat O_k$; the rationals
  $\rho_{il}\coloneqq|\operatorname{Re}(\hat O'_k-C_k)_{il}|
  +|\operatorname{Im}(\hat O'_k-C_k)_{il}|+(R_k)_{il}$ and
  $\delta_k\coloneqq\max_i\sum_l\rho_{il}$, computed exactly; and the
  relaxed entries of Lemma~\ref{lem:row-relaxation} at $\delta_k$, with
  $c'_k$ the largest binary64 number at most $c_k-\delta_k$ and $c''_k$ the
  smallest at least $c_k+\delta_k$. Where $\delta_k=0$ the entry is
  $\hat O_k$ itself and is emitted as it is, an equality included. Each
  emitted entry carries the certificate case of the entry it relaxes, its
  $\delta_k$, and the mark that it is relaxed, its certificate being
  Lemma~\ref{lem:row-relaxation}\,(ii); its operator is added to the
  companion's transported set with its pull-back. The lattice condition of
  Requirement~\ref{req:licensed-entry-lattice} binds the artefact's own
  entries and not these.
\item \emph{No normalizer.} No step of the \termref{term-emission-normalizer}{emission normalizer}
  (Definition~\ref{def:emitnorm}) acts on the companion's datum: the
  companion carries the reference, not a snap of it.
\item \emph{The record.} The companion carries, per block of $G$, whether
  it is the handed reference or the snapped block, the kind of every handed
  block (\emph{exact}, every radius zero, or \emph{enclosed}, with the
  largest radius of the block after recentring), the entries recentred in
  the handed blocks and in the datum, the largest radius of the datum's
  enclosure after its recentring, each $\delta_k$, every requested entry
  formed on the companion with what established its certificate there, and
  every requested entry not formed with the reason; and it is
  marked evaluation-only: it has no run of its own, no consumer admitting
  it as the declaration of a run, and what a consumer decides or computes
  of it beside the declaration it accompanies is the consumer's
  (\xref{meat-core/req:companion-admission}). The artefact
  itself, its target, normalizer, emission record
  (Requirement~\ref{req:emission-record}) and entries, is emitted as it is
  without a companion.
\item \emph{No companion.} Where no block is handed and the normalizer
  moves no class, or no claim of this module establishes (H-REF) for the
  instance, or a reciprocal node of the rule of (I1) receives no enclosure,
  so that Corollary~\ref{cor:reference-companion}\,(iii) gives none, no
  companion is emitted, and the record says which of the three holds.
\item \emph{Factor enclosures.} Where the handed reference carries
  \exttermref{../foundations/contract.pdf}{term-factor-enclosure}{factor enclosures}
  (\xref{foundations/def:factor-enclosure}) on some of its blocks
  $G^{(b)}_j$, $b$ a handed block, as the receiver's does
  (\xref{receiver/detector-imperfections/req:dark-composed-reference}), and a
  claim of this module establishes that the rule of (I1) of the instance is a
  non-negative combination (Corollary~\ref{cor:companion-factor-enclosure},
  part (iii) for the identity instance and
  Corollary~\ref{cor:flat-companion} for the flat instance), the engine
  forms, for every element of every data sector of the companion's datum
  every combining block of which carries one, the factor enclosure of
  Corollary~\ref{cor:companion-factor-enclosure}\,(i): the boxes of the
  combining blocks side by side in a fixed order, their centres and radii
  unaltered, and their cores multiplied by the coefficients in exact rational
  arithmetic. It emits each beside the \exttermref{../foundations/contract.pdf}{term-entrywise-enclosure}{entrywise enclosure} of the
  datum of (i), which it does not alter, as the factor enclosure of the target's block
  on that data sector (part (ii) of that corollary). It forms none on any other
  element: none where a combining block carries none, and none for an
  instance no claim of this module shows to have such a rule, in particular
  none through a rule that maps the blocks of a data sector by a conjugation
  before it combines them, whose factor enclosure this module does not form.
  The record of (v) carries, per element of the datum, whether a factor
  enclosure is carried, with its number of columns, or why none is, and it
  names the claim that shows the rule of (I1) to be a non-negative
  combination, Corollary~\ref{cor:companion-factor-enclosure}\,(iii) or
  Corollary~\ref{cor:flat-companion}.
\end{enumerate}
\end{requirement}
\claimmeta{\metalabel{Motivated by}\mref{Def.}{def:reduction-instance}, \mref{Def.}{def:emitnorm}, \mref{Cor.}{cor:entries}, \mref{Cor.}{cor:reference-companion}, \mref{Cor.}{cor:companion-factor-enclosure}, \mref{Lem.}{lem:row-relaxation}, \mref{Req.}{req:licensed-entry-lattice}, \mref{Req.}{req:emission-record}, \mref{Rem.}{rem:published}, \mref{Cor.}{cor:flat-companion}, \extref{../foundations/contract.pdf}{foundations/def:evaluated-off-snap}, \extref{../foundations/contract.pdf}{foundations/def:exact-island}, \extref{../foundations/contract.pdf}{foundations/def:factor-enclosure}, \extref{../foundations/contract.pdf}{foundations/lem:recentred-enclosure}, \extref{../foundations/contract.pdf}{foundations/req:snap-reference}, \extref{../meat-core/contract.pdf}{meat-core/frank-wolfe/req:bracket-on-reference}, \extref{../meat-core/contract.pdf}{meat-core/req:companion-admission}, \extref{../receiver/contract.pdf}{receiver/detector-imperfections/req:dark-composed-reference}
\metasep \metalabel{Used by}\extref{../meat-core/contract.pdf}{meat-core/req:companion-admission}
\metasep implemented (effective: accepted)
\metasep \metalabel{Code}\coderef{https://github.com/fraunhofersg/certiqs-key/blob/v0.1.0/src/certiqs_key/reductions/companion.py}{certiqs_key.reductions.companion:reference_companion}}

% claim reductions/req:party-label
\claimsummary{reductions/req:party-label}{Each reduction instance's party partition is emitted with its sector map: the default one, every factor Bob's, unless a party label is declared. Where several parties measure, each factor's artefact carries that factor's own label.}
\claimbookmark{reductions/req:party-label}{Requirement}{req:party-label}{Party partition emitted with the sector map}
\claimleadin{%
The consuming declaration places its well-formedness conditions on a
party's side (Theorem~\ref{thm:emit}\,(ii)), so the party partition of
Definition~\ref{def:reduction-instance}\,(I2) is part of what a consumer
reads with the target.
}
\begin{requirement}[Party partition emitted with the sector map]\label{req:party-label}
The engine emits the party partition of each \termref{term-reduction-instance}{reduction instance} with
its \exttermref{../foundations/contract.pdf}{term-sector-map}{sector map} (\iface{target_descriptor}): the
\termref{term-default-party-partition}{default party partition} where the instance declares no party label,
and otherwise the declared label (\iface{party_label}) on every factor of
every $\Hs{\Qsys_s}$. Where several parties measure, each factor's artefact
carries that factor's own label with its own sector map
(Theorem~\ref{thm:emitted-product}\,(i)).
\end{requirement}
\claimmeta{\metalabel{Motivated by}\mref{Def.}{def:reduction-instance}, \mref{Thm.}{thm:emit}, \mref{Thm.}{thm:emitted-product}
\metasep implemented (effective: accepted)
\metasep \metalabel{Code}\coderef{https://github.com/fraunhofersg/certiqs-key/blob/v0.1.0/src/certiqs_key/reductions/sectors.py}{certiqs_key.reductions.sectors:party_partition_statement}}

% claim reductions/req:factorised-emission
\claimsummary{reductions/req:factorised-emission}{Where several parties measure, the engine emits one artefact per factor, each with its party label, in the declared order, together with the combination rule. It emits, declares and rounds no entry of any product element.}
\claimbookmark{reductions/req:factorised-emission}{Requirement}{req:factorised-emission}{Emission of several measuring parties as factors}
\claimleadin{%
A Kronecker product of two families on the declared lattice lies in general
only on a finer lattice, and a normalised product would declare a non-product
measurement for separated parties; the product is therefore emitted as a
rule, and Theorem~\ref{thm:emitted-product} states what that rule gives.
}
\begin{requirement}[Emission of several measuring parties as factors]\label{req:factorised-emission}
Where several \termref{term-measuring-parties}{measuring parties} measure, the engine emits in
\iface{factor_emission} the ordered list of the factors' artefacts, each with
its party label, in the \exttermref{../foundations/contract.pdf}{term-declared-order}{declared order} of the factors, and the combination
rule \eqref{eq:producttarget}; it emits, declares and rounds no entry of any
product element $F_v$ (Theorem~\ref{thm:emitted-product}).
\end{requirement}
\claimmeta{\metalabel{Motivated by}\mref{Thm.}{thm:emitted-product}
\metasep implemented (effective: accepted)
\metasep \metalabel{Code}\coderef{https://github.com/fraunhofersg/certiqs-key/blob/v0.1.0/src/certiqs_key/reductions/factors.py}{certiqs_key.reductions.factors:reduce_factors}}

\section{Requirements: the assumption ledger}\label{group-reductions/req:assumption-ledger}
\sectionintro{The emitted assumption ledger carries the rows of the emission's assumptions under their presence rules and marks those that are never verified. It refuses a supplier artefact of a version the caller has not declared accepted, passes the supplier's ledger through entire, refuses one with a repeated or missing key, and relies only on supplier presence rules that omit no premise an imported fact rests on. It omits the row of the supplier assumption only where the run attests that it built every receiver artefact read and each is exact with rigorous weights (\hyperlink{claim-reductions/req:recv-row-provenance}{Requirement~\ref*{req:recv-row-provenance}}). The engine declares the output-state model of the test statistics.}
\sectionimports{\extref{../receiver/contract.pdf}{receiver/cor:unconditioned-family}, \extref{../receiver/contract.pdf}{receiver/def:setting-family}, \extref{../receiver/contract.pdf}{receiver/def:slot}, \extref{../receiver/contract.pdf}{receiver/thm:click-povm}, \extref{../receiver/contract.pdf}{receiver/thm:coarse-graining}, \extref{../receiver/contract.pdf}{receiver/thm:enclosed-blocks}, \extref{../receiver/contract.pdf}{receiver/thm:enclosed-weight}, \extref{../receiver/contract.pdf}{receiver/thm:representation-independence}}

% claim reductions/req:assumption-ledger
\claimsummary{reductions/req:assumption-ledger}{Every artefact carries the assumption ledger: a row per assumption of this contract or taken from another module, under its presence rule. The rows cover every such key of the producing claims' closure but those that hold vacuously and \hyperlink{claim-reductions/ass:recv}{A-RECV} for an exact, rigorous receiver artefact the run attests.}
\claimbookmark{reductions/req:assumption-ledger}{Requirement}{req:assumption-ledger}{Rows of the emitted assumption ledger}
\claimleadin{%
A consumer inherits from this module the assumptions its artefact names, and
it can match an assumption against its own ledger only by machine key.
}
\begin{requirement}[Rows of the emitted assumption ledger]\label{req:assumption-ledger}
The engine emits with every artefact the assumption ledger
\iface{assumption_ledger}: one row, carrying its machine key, for each
assumption of this contract and for each assumption the interface lists as
taken from another module, each under its declared presence rule. The row
of A-SRC (\xref{transmitter/ass:src}) is present only where the emitted chain's
\exttermref{../transmitter/contract.pdf}{term-source-segment}{source segment} is non-empty, and there under the presence rule the
interface declares for it, the row of Assumption~\ref{ass:constdecl} only
where the emitted constraint list is non-empty, the row of
Assumption~\ref{ass:probe} only where the declared setting set $\Tset$ has
at least two values, the row of Assumption~\ref{ass:tsm} only where $\Tset$
has at least two values or a public choice is applied at the measurement the
emitted family describes (and also where the emission is not told whether
one is), a submodule's own assumption
where that submodule's instance is emitted and any narrower presence rule the
submodule declares for it holds, the row of Assumption~\ref{ass:recv} on
every emission save where Requirement~\ref{req:recv-row-provenance} omits
it, and every other row on every emission.
These rows include every key
of the core's own assumptions, and every key the interface lists as taken
from another module, that the closure of the claims producing the artefact
contains, save the row of Assumption~\ref{ass:constdecl} where the emitted
constraint list is empty, a condition on the constants of emitted entries
holding vacuously of an empty list, the row of Assumption~\ref{ass:probe}
where $\Tset=\{\ast\}$, whose clauses hold of the one setting drawn with
weight one, and the row of Assumption~\ref{ass:tsm} where in addition no
public choice is applied at that measurement, whose clauses then concern no
choice, and the row of Assumption~\ref{ass:recv} where
Requirement~\ref{req:recv-row-provenance} omits it, whose clauses that
requirement shows to hold there (the supplier's keys ride the row of
Assumption~\ref{ass:recv}
instead, under the supplier's own presence rules,
Requirement~\ref{req:supplier-ledger-pass-through}); a row may be present that
no such closure contains.
\end{requirement}
\claimmeta{\metalabel{Motivated by}\mref{Ass.}{ass:recv}, \mref{Ass.}{ass:probe}, \mref{Ass.}{ass:tsm}, \mref{Ass.}{ass:constdecl}, \mref{Req.}{req:supplier-ledger-pass-through}, \mref{Req.}{req:recv-row-provenance}, \extref{../transmitter/contract.pdf}{transmitter/ass:src}
\metasep implemented
\metasep \metalabel{Code}\coderef{https://github.com/fraunhofersg/certiqs-key/blob/v0.1.0/src/certiqs_key/reductions/emit.py}{certiqs_key.reductions.emit:emitted_assumption_rows}}

% claim reductions/req:never-verified-assumptions
\claimsummary{reductions/req:never-verified-assumptions}{The ledger marks as never verifiable the row of every assumption the engine never verifies, as its verification field says, whichever module states it. Validation decides the declared data the assumptions are read with, and no code path decides more of a never-verified one.}
\claimbookmark{reductions/req:never-verified-assumptions}{Requirement}{req:never-verified-assumptions}{Never-verified rows and validation of declared data}
\claimleadin{%
An assumption that is a statement about the physical devices, the source,
the adversary or the run, rather than about \exttermref{../foundations/contract.pdf}{term-declared-data}{declared data}, has a
clause that no code path can decide; what validation decides beside that
clause does not change its mark. Which assumptions the engine never verifies
is recorded in each one's verification field
(\href{../../../spec/claim-format.md\#32-whether-the-engine-verifies-an-assumption}{\nolinkurl{spec/claim-format.md}\texttt{\#32-whether-the-engine-verifies-an-assumption}}),
which the contract prints beside the statement. A row that this module takes
from another module carries the mark its owner gives it, as the row of the
transmitter's retention assumption (\xref{transmitter/ass:src}) does.
}
\begin{requirement}[Never-verified rows and validation of declared data]\label{req:never-verified-assumptions}
The engine marks as never verifiable (\iface{never_verified}) the row, in the
assumption ledger, of every assumption that the engine never verifies, as
that assumption's verification field says, whether this contract, one of its
submodules or another module states it; no code path decides what any of
them says beyond its clauses on the declared data. The engine decides by
validation the conditions on the declared data that the assumptions of this
contract are read with, such as the declared dimensions of
Definition~\ref{def:system-setting} and the declared \termref{term-class-weights}{class weights} of
Definition~\ref{def:joint-family}.
\end{requirement}
\claimmeta{\metalabel{Motivated by}\mref{Def.}{def:system-setting}, \mref{Def.}{def:joint-family}, \extref{../transmitter/contract.pdf}{transmitter/ass:src}
\metasep implemented
\metasep \metalabel{Code}\coderef{https://github.com/fraunhofersg/certiqs-key/blob/v0.1.0/src/certiqs_key/reductions/artifact.py}{certiqs_key.reductions.artifact:assumption_ledger}}

% claim reductions/req:supplier-ledger-pass-through
\claimsummary{reductions/req:supplier-ledger-pass-through}{The row of \hyperlink{claim-reductions/ass:recv}{A-RECV} carries every machine key of the supplied ledger, as supplied and in the supplier's order, beside the artefact's version. The engine keeps no list of the supplier's keys and adds no presence rule of its own.}
\claimbookmark{reductions/req:supplier-ledger-pass-through}{Requirement}{req:supplier-ledger-pass-through}{Pass-through of the supplier's assumption ledger}
\claimleadin{%
The facts a \termref{term-reduction-instance}{reduction instance} imports from a supplier rest on the
supplier's assumptions (Assumption~\ref{ass:recv}); a consumer inherits them
only if the artefact names them by the supplier's own machine keys.
}
\begin{requirement}[Pass-through of the supplier's assumption ledger]\label{req:supplier-ledger-pass-through}
On the row of Assumption~\ref{ass:recv} in the assumption ledger
\iface{assumption_ledger}, the engine carries, as machine keys and beside the
artefact's module version \iface{module_version}, every key the supplied
ledger carries, as supplied and in the supplier's order; a key beyond those
the supplier's presence rules require is carried rather than refused. The
engine holds no list of the supplier's keys and adds no presence rule of its
own: what it reads from the supplying module, as it is installed where the
check runs, is that module's machine keys and the set of keys its rules make
conditional, and it evaluates the condition itself on the artefact's own
data.
\end{requirement}
\claimmeta{\metalabel{Motivated by}\mref{Ass.}{ass:recv}
\metasep accepted}

% claim reductions/req:supplier-ledger-refusals
\claimsummary{reductions/req:supplier-ledger-refusals}{A supplied ledger that repeats a machine key is refused, and so is an artefact whose ledger lacks a key the supplier's own presence rules make present. A shortened list is never passed on in silence.}
\claimbookmark{reductions/req:supplier-ledger-refusals}{Requirement}{req:supplier-ledger-refusals}{Refusal of a supplier ledger with a repeated or missing key}
\claimleadin{%
The pass-through of Requirement~\ref{req:supplier-ledger-pass-through}
carries the supplied keys as they stand, so a repeated key would be carried
twice and a missing one would silently shorten what a consumer inherits.
}
\begin{requirement}[Refusal of a supplier ledger with a repeated or missing key]\label{req:supplier-ledger-refusals}
The engine refuses a supplied ledger that carries one machine key more than
once, and it refuses an artefact whose ledger lacks a key the supplier's own
presence rules make present, rather than passing it on as a silently
shortened list; both refusals carry the reason
\iface{receiver_spec_not_accepted}.
\end{requirement}
\claimmeta{\metalabel{Motivated by}\mref{Req.}{req:supplier-ledger-pass-through}
\metasep accepted}

% claim reductions/req:supplier-version-accepted
\claimsummary{reductions/req:supplier-version-accepted}{A supplier artefact of a version the caller does not declare accepted, or without its supplier's validation results, is refused. Outside the run, that the set names accepted versions and the artefact is genuine stays \hyperlink{claim-reductions/ass:recv}{A-RECV}; in the run, acceptance is the claims' effective status.}
\claimbookmark{reductions/req:supplier-version-accepted}{Requirement}{req:supplier-version-accepted}{Refusal of a supplier artefact of an unaccepted version}
\claimleadin{%
The version an artefact declares and the presence of its supplier's
validation results are declared data, decided on entry. The check is against
the caller's declared set of accepted versions
(\iface{accepted_receiver_versions}); that every version in that set is an
accepted version of the supplying module and that the artefact is what it
declares are Assumption~\ref{ass:recv}\,(i) for an artefact from outside the
run, and that its facts hold of the device is Assumption~\ref{ass:recv}\,(ii)
wherever the ledger carries its row. For every receiver artefact the run
built itself, whether or not Requirement~\ref{req:recv-row-provenance} omits
that row, acceptance is no premise but the effective status of the
receiver's claims in the run's closure, which the record of the run
reports.
}
\begin{requirement}[Refusal of a supplier artefact of an unaccepted version]\label{req:supplier-version-accepted}
The engine refuses a supplier artefact whose declared module version
\iface{module_version} is not in the caller's declared set of accepted
versions \iface{accepted_receiver_versions}, and one that does not carry the
supplying module's validation results; both refusals carry the reason
\iface{receiver_spec_not_accepted}, the reason the ledger refusals of
Requirement~\ref{req:supplier-ledger-refusals} carry as well.
\end{requirement}
\claimmeta{\metalabel{Motivated by}\mref{Ass.}{ass:recv}, \mref{Req.}{req:supplier-ledger-refusals}, \mref{Req.}{req:recv-row-provenance}
\metasep implemented
\metasep \metalabel{Code}\coderef{https://github.com/fraunhofersg/certiqs-key/blob/v0.1.0/src/certiqs_key/reductions/validation.py}{certiqs_key.reductions.validation:rule_receiver_spec}}

% claim reductions/req:supplier-presence-rules
\claimsummary{reductions/req:supplier-presence-rules}{The supplier's presence rules, read as installed, may omit a key of the imported facts' closure only where the facts imported through it hold without it, and a change of the rules must meet this too. The pass-through then never drops a premise an imported fact rests on.}
\claimbookmark{reductions/req:supplier-presence-rules}{Requirement}{req:supplier-presence-rules}{Supplier presence rules that omit no imported premise}
\claimleadin{%
The row of Assumption~\ref{ass:recv} carries the supplier's ledger under the
supplier's own presence rules (Requirement~\ref{req:supplier-ledger-pass-through}),
and those rules are read from the supplying module as it is installed where
the check runs, so a change of them takes effect here without an edit.
}
\begin{requirement}[Supplier presence rules that omit no imported premise]\label{req:supplier-presence-rules}
The supplier's presence rules that the engine reads omit a key of the closure
of the facts a \termref{term-reduction-instance}{reduction instance} imports from that supplier only where
the facts imported through it hold of the declaration without it; a change
of those rules, which takes effect in the engine's check without an edit of
this module, must itself meet this requirement.
\end{requirement}
\claimmeta{\metalabel{Motivated by}\mref{Ass.}{ass:recv}, \mref{Req.}{req:supplier-ledger-pass-through}
\metasep accepted}

% claim reductions/req:recv-row-provenance
\claimsummary{reductions/req:recv-row-provenance}{The ledger omits \hyperlink{claim-reductions/ass:recv}{A-RECV} only where every supplier artefact read is a receiver artefact this run built, with a model admissible exactly and every weight evaluated rigorously. Both clauses then hold by the receiver's theorems under its own ledger, with no premise on arithmetic.}
\claimbookmark{reductions/req:recv-row-provenance}{Requirement}{req:recv-row-provenance}{Omission of the supplier row for a receiver artefact the run built}
\claimleadin{%
Assumption~\ref{ass:recv} names what a \termref{term-reduction-instance}{reduction instance} cannot see
in a supplier artefact: that it is its module's output at an accepted
version, on the inputs it declares, and that the facts the supplier's claims
state hold of its numbers. Version, validation results and ledger are data
that an artefact written by hand can carry equally; whether the run built
the artefact is known to the part of the engine that ran the supplier, as
for the chain of \xref{protocol-compiler/req:link-row-provenance}. Where the
run built it, clause~(i) holds by construction: the artefact is the
installed receiver's output on the declaration the run handed it, with that
receiver's own validation results, at the version it states. For every
receiver artefact the run built, whether or not the row is omitted,
acceptance is no premise: it is the effective status of the receiver's
claims in the run's closure, which the run's record reports, as for every
module the run executes; the acceptance part of clause~(i) is a premise only
for an artefact from outside the run. Clause~(ii) concerns the facts the instance imports, those the
consuming claims cite (Remark~\ref{rem:receiver-imports};
Proposition~\ref{prop:route}\,(iv)), and condition~(iv) below makes each a
theorem of the receiver with no premise on arithmetic. A declaration with a
model admissible exactly (\xref{receiver/req:exact-route-decision}) has
the port Grams of a \exttermref{../receiver/contract.pdf}{term-receiver-instrument}{receiver instrument}, which determine its
click-pattern measurement (\xref{receiver/thm:representation-independence}),
so the theorems stated of a receiver instrument apply to it under the
receiver's ledger, \xref{receiver/ass:trust} saying that it is the physical
receiver: the block diagonality in the photon number of the
elements of its \exttermref{../receiver/contract.pdf}{term-setting-family}{setting family} and of their coarse-graining
(\xref{receiver/def:setting-family}, \xref{receiver/thm:coarse-graining},
\xref{receiver/thm:click-povm}, \xref{receiver/cor:unconditioned-family}); the
\exttermref{../receiver/contract.pdf}{term-dark-count-channel}{dark-count channel}, the declared matrix or the product form, of which
(D1) and (D2) hold by construction and which this module decides exactly on
entry (\xref{receiver/detector-imperfections/def:darkadm},
\xref{receiver/detector-imperfections/thm:dark-counts}); at a
\exttermref{../receiver/contract.pdf}{term-normalized-station}{normalized station} the efficiency normalization
(\xref{receiver/detector-imperfections/prop:efficiency-normalization}); and
the \exttermref{../receiver/contract.pdf}{term-strengthening-slot}{strengthening slot} (\xref{receiver/def:slot}). Its weights carry
the rigorous label (\xref{receiver/req:evaluation-label}), so each
\exttermref{../receiver/contract.pdf}{term-per-setting-weight}{per-setting weight} is the largest binary64 number at or below an
\exttermref{../receiver/contract.pdf}{term-enclosed-weight}{enclosed weight} (\xref{receiver/req:rigorous-weight}\,(iv)), and
\xref{receiver/thm:enclosed-weight}\,(iii) gives its guarantee on every
photon-number block above the declared cutoff at its own setting, the
quantifier domain being the declared cutoff and settings; at a setting
family a receiver submodule constructs, that submodule states the same of
its rigorously evaluated weight (\xref{receiver/thm:exact-and-evaluated},
lead-in). No computed eigenvalue enters such a weight, so
\xref{receiver/ass:eig} is invoked at no setting, and the non-negativity and
the sum of the weights at a setting are decided on entry (the slot row of
\href{../../../spec/modules/reductions/wpfss/interface.yaml}{\nolinkurl{spec/modules/reductions/wpfss/interface.yaml}}). The photon-number
blocks below the cutoff are no imported fact: they reach the instance as its
declaration (\xref{foundations/def:declared-data}\,(ii)), beside rational
enclosures that contain the model's blocks
(\xref{receiver/thm:enclosed-blocks}), against which the declaring party
snaps them, a moved block resting on
\xref{foundations/ass:snapped-measurement} (Remark~\ref{rem:recv-scope}).
Where a weight is evaluated in floating point, the receiver asserts no
bound between the emitted number and the weight its theorems speak of
(\xref{receiver/thm:exact-and-evaluated}\,(ii)), and where the declaration
is admitted only within the \exttermref{../receiver/contract.pdf}{term-receiver-tolerances}{receiver tolerances} no receiver instrument has
its declared numbers; there clause~(ii) stays a premise.
}
\begin{requirement}[Omission of the supplier row for a receiver artefact the run built]\label{req:recv-row-provenance}
The engine omits the row of Assumption~\ref{ass:recv} from the assumption
ledger \iface{assumption_ledger} of an emission only where the
\termref{term-reduction-instance}{reduction instance} read at least one supplier artefact and, for every
supplier artefact it read, all of the following hold:
(i) the supplying module is the receiver, and the artefact is one that the
part of the engine that ran this module obtained in the same run by running
the receiver module, as installed where the run runs, on the receiver
declaration the run holds, and handed to this module as it was emitted;
(ii) the module version \iface{module_version} the artefact declares is the
one that installed module states for the artefacts it emits, and the
artefact passes Requirements~\ref{req:supplier-version-accepted}
and~\ref{req:supplier-ledger-refusals};
(iii) the record of the run carries that artefact's own assumption ledger;
(iv) the artefact records that its declaration has a model admissible
exactly (\xref{receiver/req:exact-route-decision}\,(iv)), and it carries the
evaluation label \iface{rigorous} (\xref{receiver/req:evaluation-label}),
every per-setting weight having been formed by the rigorous evaluation of
\xref{receiver/req:rigorous-weight}.
There clause~(i) of Assumption~\ref{ass:recv} holds by construction, the
acceptance of the version being the effective status of the receiver's
claims in the run's closure, which the record of the run reports, and
clause~(ii) holds of the facts the instance imports, each a theorem of the
receiver about the receiver instrument the declaration is, holding of the
numbers emitted under the ledger~(iii) carries and resting on no premise
about arithmetic. That (i) and (iii) hold is established by the part of the
engine that ran the receiver module and docked its artefact into the
instance, and is passed to this module beside the artefact
(\iface{receiver_provenance}); the engine never reads it from the artefact.
The engine decides (ii) and (iv) on the artefact. Where it omits the row, it
records the omission with the artefact's version
(\iface{receiver_row_omitted_on_provenance}), and the supplier's keys, which
ride that row where it is emitted
(Requirements~\ref{req:supplier-ledger-pass-through}
and~\ref{req:supplier-eig-record}), reach the record of the run from the
artefact's own ledger by~(iii). Where \iface{receiver_provenance} is not
passed, or a condition fails for any artefact read, the row is emitted as
Requirement~\ref{req:assumption-ledger} says.
\end{requirement}
\claimreading{Where this run built the receiver artefact itself from your declared receiver, your declaration describes a valid receiver exactly, and every high-photon weight was computed rigorously, the report does not list the premise that a supplied receiver file is genuine and true of your device (\hyperlink{claim-reductions/ass:recv}{A-RECV}). What that premise says then follows from the receiver's theorems under the receiver's own assumptions, which the report lists, and whether those theorems are accepted is shown by their status in the report. Wherever a weight was evaluated in floating point, the receiver was admitted only within its tolerances, or the receiver file came from outside the run, the report keeps \hyperlink{claim-reductions/ass:recv}{A-RECV}.}
\claimmeta{\metalabel{Motivated by}\mref{Ass.}{ass:recv}, \mref{Rem.}{rem:recv-scope}, \mref{Req.}{req:assumption-ledger}, \mref{Req.}{req:supplier-ledger-pass-through}, \mref{Req.}{req:supplier-ledger-refusals}, \mref{Req.}{req:supplier-version-accepted}, \mref{Prop.}{prop:route}, \mref{Rem.}{rem:receiver-imports}, \mref{Req.}{req:supplier-eig-record}, \extref{../foundations/contract.pdf}{foundations/ass:snapped-measurement}, \extref{../foundations/contract.pdf}{foundations/def:declared-data}, \extref{../protocol-compiler/contract.pdf}{protocol-compiler/req:link-row-provenance}, \extref{../receiver/contract.pdf}{receiver/ass:eig}, \extref{../receiver/contract.pdf}{receiver/ass:trust}, \extref{../receiver/contract.pdf}{receiver/cor:unconditioned-family}, \extref{../receiver/contract.pdf}{receiver/def:setting-family}, \extref{../receiver/contract.pdf}{receiver/def:slot}, \extref{../receiver/contract.pdf}{receiver/detector-imperfections/def:darkadm}, \extref{../receiver/contract.pdf}{receiver/detector-imperfections/prop:efficiency-normalization}, \extref{../receiver/contract.pdf}{receiver/detector-imperfections/thm:dark-counts}, \extref{../receiver/contract.pdf}{receiver/req:evaluation-label}, \extref{../receiver/contract.pdf}{receiver/req:exact-route-decision}, \extref{../receiver/contract.pdf}{receiver/req:rigorous-weight}, \extref{../receiver/contract.pdf}{receiver/thm:click-povm}, \extref{../receiver/contract.pdf}{receiver/thm:coarse-graining}, \extref{../receiver/contract.pdf}{receiver/thm:enclosed-blocks}, \extref{../receiver/contract.pdf}{receiver/thm:enclosed-weight}, \extref{../receiver/contract.pdf}{receiver/thm:exact-and-evaluated}, \extref{../receiver/contract.pdf}{receiver/thm:representation-independence}
\metasep implemented (effective: accepted)
\metasep \metalabel{Code}\coderef{https://github.com/fraunhofersg/certiqs-key/blob/v0.1.0/src/certiqs_key/reductions/emit.py}{certiqs_key.reductions.emit:receiver_row_omitted_on_provenance}}

% claim reductions/req:output-state-statistics
\claimsummary{reductions/req:output-state-statistics}{The engine declares the output-state model of the test statistics and emits no state-independent test fraction. Every public choice of known bias enters through the statistics of the emitted family, as \hyperlink{claim-reductions/ass:tsm}{A-TSM} states.}
\claimbookmark{reductions/req:output-state-statistics}{Requirement}{req:output-state-statistics}{Output-state model of the test statistics}
\claimleadin{%
No derivation of this contract assumes a state-independent test fraction:
every per-round public choice of known bias enters only through the
statistics of the emitted family (Assumption~\ref{ass:tsm}).
}
\begin{requirement}[Output-state model of the test statistics]\label{req:output-state-statistics}
The engine emits no state-independent test fraction, and it declares the
test-statistics model \iface{test_statistics_model} to be
\iface{output_state}.
\end{requirement}
\claimmeta{\metalabel{Motivated by}\mref{Ass.}{ass:tsm}
\metasep accepted}

\clearpage
\part{Flag-state and weight-preserving squashers}
\sectionintro{This submodule consumes a receiver artefact and turns it into a finite reduction instance. It keeps the photon-number blocks below a cutoff and replaces everything above by flags: the flag-state target, which puts each outcome on its own flag, or the weight-preserving target, which puts each outcome's certified high-photon weight on every flag and one minus the total weight on its own. Its first main result is that, whenever the certified weights total less than one, the physical measurement squashes to both targets, the two squashers compose exactly, and a weight the weight-preserving target encodes holds on every squashed state (\hyperlink{claim-reductions/wpfss/prop:chain}{Proposition~\ref*{prop:chain}}). The weight route first applies the dark-count channel on the emitted alphabet, which requires it to descend through the coarse-graining, then forms the joint family of the measuring party's private setting and routes the receiver's per-setting multiclick weights into certified weights rounded down onto the lattice. The resulting flat instance, on the encoded or the licensed route, is a reduction instance, and the second main result is its emission: the emitted family is exactly admissible and pulls back exactly to the declared model on the channel system, with its weights rounded down and still certified (\hyperlink{claim-reductions/wpfss/cor:emitted-flat}{Corollary~\ref*{cor:emitted-flat}}). Where the emission moves a block below the cutoff, that statement is about the declared model, which nothing proved relates to the unrounded family; a run evaluates off the snap on the reference companion of the flat instance, or the moved class is disclosed. Where the receiver is compiled at its normalized station, the instance is formed on that station and precomposed with the factored loss, which changes the channel witness and nothing else emitted (\hyperlink{claim-reductions/wpfss/cor:normalized-flat}{Corollary~\ref*{cor:normalized-flat}}). The receiver's facts enter by named claims under A-RECV, and the receiver assumptions those claims rest on enter the assumption closures with them.}

\section{Flag-state targets and their chain}\label{group-reductions/wpfss/def:targets}
\sectionintro{Above a cutoff the receiver is replaced by flags: the \termref{term-flag-state-target}{flag-state target} keeps the blocks below the cutoff and puts each outcome on its own flag, and the \termref{term-weight-preserving-target}{weight-preserving target} puts each outcome's certified high-photon weight on every flag and one minus the \termref{term-total-weight}{total weight} on its own. The main result, \hyperlink{claim-reductions/wpfss/prop:chain}{Proposition~\ref*{prop:chain}}, is that the physical measurement squashes to both targets whenever the \termref{term-certified-weights}{certified weights} total less than one, that the chain from the physical measurement through the weight-preserving target to the flag-state target composes exactly, and that a weight the weight-preserving target encodes holds on every squashed state, whereas under the flag-state target it holds only on reachable states. An artefact therefore takes one of two routes, and a lowered certified weight stays certified and, while the total weight is below one, gives a cleaner target (\hyperlink{claim-reductions/wpfss/prop:weight-monotone}{Proposition~\ref*{prop:weight-monotone}}).}

% claim reductions/wpfss/def:targets
\claimsummary{reductions/wpfss/def:targets}{Certified weights are per-outcome lower bounds on a POVM's blocks above a cutoff. Both targets keep the blocks below the cutoff; the flag-state target puts each outcome on its own flag, the weight-preserving target puts its certified weight on every flag and one minus the total weight on its own.}
\claimbookmark{reductions/wpfss/def:targets}{Definition}{def:targets}{Certified weights and the flag-state and weight-preserving targets}
\claimleadin{%
Above a photon-number cutoff the receiver is replaced by one classical flag
per outcome, and the only device knowledge kept there is a lower bound on each
outcome's block. The weight-preserving target keeps these bounds in its flags,
so that the corresponding inequality holds on every squashed state
(Proposition~\ref{prop:chain}\,(iii)); both targets are known from the
literature (Remark~\ref{rem:published}).
}
\begin{definition}[Certified weights and the flag-state and weight-preserving targets]
\label{def:targets}
Throughout this definition and the claims of this submodule that use it,
$\Hs{\Bsys}=\Hs{\Bsys_{\leqN}}\dsum\Hs{\Bsys_{\gtN}}$
is an orthogonal two-block decomposition with projections $\Pleq$ and $\Pgt$
--- a photon-number cutoff in the optical instance, where
$\Hs{\Bsys_{\leqN}}=\bigoplus_{m\le N}\Hs{\Chan_m}$ is finite-dimensional ---
and every $M_v$ is block diagonal for it,
$M_v=\blk{M}{\leqN}_v\dsum\blk{M}{\gtN}_v$. $\Hs{\Flag}=\mathbb{C}^{|\Vset|}$
carries the orthonormal basis $\{\ket{v}\}_{v\in\Vset}$, and
$\Hs{\Qsys}=\Hs{\Bsys_{\leqN}}\dsum\Hs{\Flag}$ with block projections
$\Pleq$ (the same symbol on either space) and $\Pflag$.

A family $(p_v)_{v\in\Vset}$ of
\termdef{term-certified-weights}{certified weights}
for $M$ at the cutoff consists of non-negative numbers $p_v$ --- on the
emission path, integer multiples of $2^{-\kappa}$, the $p_{(y,t)}$ of
\eqref{eq:weightroute} --- with
\begin{equation}\label{eq:certweights}
  \blk{M}{\gtN}_v\;\succeq\;p_v\,\idon{\Bsys_{\gtN}}
  \qquad\text{for every }v\in\Vset ,
\end{equation}
and we write $p\coloneqq\sum_vp_v$ for their
\termdef{term-total-weight}{total weight}.
Define
\begin{align}
  \Ffss_v&\coloneqq\blk{M}{\leqN}_v\dsum\dyad{v}_{\Flag},
  \label{eq:FSS}\\
  \Fwp_v&\coloneqq\blk{M}{\leqN}_v\dsum
     \big(p_v\,\idon{\Flag}+(1-p)\,\dyad{v}_{\Flag}\big).
  \label{eq:WPFSS}
\end{align}
$\Ffss$ is the
\termdef{term-flag-state-target}{flag-state target}
and $\Fwp$ the
\termdef{term-weight-preserving-target}{weight-preserving target},
the multi-bound weight-preserving flag-state-squasher target, with one
certified weight per outcome. The
\termdef{term-single-bound-case}{single-bound case}
is the one in which only one $p_v$ is non-zero, and $\Ffss$ is the case
$p_v\equiv0$.
\end{definition}
\claimmeta{\metalabel{Uses}\mref{Def.}{def:weight-route}
\metasep \metalabel{Used by}\mref{Lem.}{lem:ple1}, \mref{Rem.}{rem:peqone}, \mref{Prop.}{prop:chain}, \mref{Def.}{def:routes}, \mref{Rem.}{rem:tworoutes}, \mref{Prop.}{prop:weight-monotone}, \mref{Rem.}{rem:weightfree}, \mref{Rem.}{rem:published}, \mref{Prop.}{prop:route}, \mref{Prop.}{prop:precomposed-blocks}, \mref{Rem.}{rem:absorbneeded}, \mref{Cor.}{cor:fssinstance}, \mref{Prop.}{prop:flat-instance}, \mref{Cor.}{cor:emitted-flat}, \mref{Cor.}{cor:two-sector-decided}, \mref{Cor.}{cor:flat-companion}
\metasep newly derived
\metasep accepted}

% claim reductions/wpfss/lem:ple1
\claimsummary{reductions/wpfss/lem:ple1}{Certified weights are non-negative, and their total is at most one whenever the space above the cutoff is non-trivial. On a trivial space above the cutoff the certified-weight condition is vacuous and bounds the total not at all.}
\claimbookmark{reductions/wpfss/lem:ple1}{Lemma}{lem:ple1}{Bound on the total certified weight}
\claimleadin{%
The hypothesis on the space above the cutoff is needed: where that block is
trivial, \eqref{eq:certweights} is vacuous and constrains $p$ not at all.
}
\begin{technicalclaim}
\begin{lemma}[Bound on the total certified weight]\label{lem:ple1}
For any \termref{term-certified-weights}{certified weights}, $p\ge0$; and $p\le1$ whenever
$\Hs{\Bsys_{\gtN}}\ne\{0\}$.
\end{lemma}
\end{technicalclaim}
\claimmeta{\metalabel{Uses}\mref{Def.}{def:targets}
\metasep \metalabel{Used by}\mref{Rem.}{rem:peqone}, \mref{Req.}{req:total-weight-refusal}
\metasep newly derived
\metasep accepted}

\begin{proof}
Summing \eqref{eq:certweights} over $v$ and using block completeness
$\sum_v\blk{M}{\gtN}_v=\idon{\Bsys_{\gtN}}$ gives
$\idon{\Bsys_{\gtN}}\succeq p\,\idon{\Bsys_{\gtN}}$, hence $p\le1$ whenever
$\Hs{\Bsys_{\gtN}}\ne\{0\}$; and $p\ge0$ termwise.
\end{proof}

% claim reductions/wpfss/rem:peqone
\claimsummary{reductions/wpfss/rem:peqone}{At total weight one the certified weights are exact and the relation still holds, through a map that keeps the block below the cutoff and replaces the block above it by a fixed flag state. Refusing the case is therefore conservative rather than forced by an obstruction.}
\claimbookmark{reductions/wpfss/rem:peqone}{Remark}{rem:peqone}{The boundary case of total weight one}
\claimleadin{%
Lemma~\ref{lem:ple1} settles two of the failure conditions of a squashing
request (\href{../../../spec/architecture.md\#54-failure-policy}{\nolinkurl{spec/architecture.md}\texttt{\#54-failure-policy}}) to the extent it
proves and no further: $p_v<0$ cannot be produced by a valid weight family,
and neither can $\sum_vp_v>1$ \emph{whenever $\Hs{\Bsys_{\gtN}}\ne\{0\}$}, so
on that hypothesis they are impossibilities rather than error paths. Where the
block above the cutoff is trivial the second is possible, the certified-weight
condition being vacuous there, and the construction of
Proposition~\ref{prop:chain} divides by $1-p$; a \termref{term-total-weight}{total weight} of one
or more is therefore refused (Requirement~\ref{req:total-weight-refusal}). The
remark shows that the boundary case $p=1$ is no obstruction.
}
\begin{remark}[The boundary case of total weight one]
\label{rem:peqone}
At $p=1$ the certified weights are forced to be exact: summing
\eqref{eq:certweights} gives
$\sum_v\big(\blk{M}{\gtN}_v-p_v\idon{}\big)=0$ with every summand positive
semidefinite, so every summand vanishes and
$\blk{M}{\gtN}_v=p_v\idon{}$ for each $v$. The relation then still holds, with
the target $F_v=\blk{M}{\leqN}_v\dsum p_v\idon{\Flag}$ --- which is
\eqref{eq:WPFSS} at $1-p=0$ --- realised by
$\Lambda(\rho)=\Pleq\rho\Pleq\dsum\tr[\Pgt\rho\Pgt]\,\mu$ with
$\mu=\sum_vp_v\dyad{v}$ a state: that map is a channel, its adjoint sends
$X_{\leqN}\dsum X_{\Flag}$ to
$X_{\leqN}\dsum\tr[\mu X_{\Flag}]\idon{\Bsys_{\gtN}}$, and evaluating at $F_v$
returns $\blk{M}{\leqN}_v\dsum p_v\idon{}=M_v$. So $p=1$ is not an
obstruction. Refusing it instead of constructing that map
(Requirement~\ref{req:total-weight-refusal}) costs nothing that a declaration
reaching $p=1$ could not also obtain by lowering one weight by the smallest
step the declared lattice allows, $2^{-\kappa}$.
\end{remark}
\claimmeta{\metalabel{Uses}\mref{Def.}{def:targets}, \mref{Lem.}{lem:ple1}
\metasep newly derived
\metasep accepted}

% claim reductions/wpfss/prop:chain
\claimsummary{reductions/wpfss/prop:chain}{With total certified weight below one, the physical POVM squashes to both targets, and the squashers compose exactly through a flag-noise channel. A weight the weight-preserving target encodes holds on every squashed state; under the flag-state target it holds only on reachable ones.}
\claimbookmark{reductions/wpfss/prop:chain}{Proposition}{prop:chain}{Flag-state squashers: existence, chain and encoded weights}
\claimleadin{%
Both targets are reached by measuring the part above the cutoff and preparing
a flag, the weight-preserving one after the certified weights have been taken
out of the measured POVM. In the weight-preserving target the weights are
then an operator inequality on the target itself (part (iii)), which is the
difference between the two routes (Remark~\ref{rem:tworoutes}).
}
\begin{mainclaim}
\begin{proposition}[Flag-state squashers: existence, chain and encoded weights]
\label{prop:chain}
Let $(p_v)$ be \termref{term-certified-weights}{certified weights} with $p<1$ and define the channels
\begin{align}
  \Lwp(\rho)
   &\coloneqq\Pleq\rho\,\Pleq\ \dsum\
      \sum_{v}\tr\!\big[\Mwp_v\,\Pgt\rho\,\Pgt\big]\,\dyad{v}_{\Flag},
   \qquad
   \Mwp_v\coloneqq\frac{\blk{M}{\gtN}_v-p_v\,\idon{\Bsys_{\gtN}}}{1-p},
   \label{eq:Lwp}\\
  \Lfss(\rho)
   &\coloneqq\Pleq\rho\,\Pleq\ \dsum\
      \sum_{v}\tr\!\big[\blk{M}{\gtN}_v\,\Pgt\rho\,\Pgt\big]\,\dyad{v}_{\Flag},
   \label{eq:Lfss}\\
  \Lnoise(\sigma)
   &\coloneqq(1-p)\,\sigma+p\,\Pleq\sigma\,\Pleq
      +\tr\!\big[\Pflag\sigma\big]\,\mu,
   \qquad \mu\coloneqq\sum_vp_v\dyad{v}_{\Flag}.
   \label{eq:Lnoise}
\end{align}
Then:
\begin{enumerate}[label=(\roman*),leftmargin=3.0em]
\item \emph{(existence, with the cutoff projectors transported)}
  \begin{equation}\label{eq:transportproj}
    \big(M;\Pleq,\Pgt\big)\dsq_{\Lwp}\big(\Fwp;\Pleq,\Pflag\big),
    \qquad
    \big(M;\Pleq,\Pgt\big)\dsq_{\Lfss}\big(\Ffss;\Pleq,\Pflag\big);
  \end{equation}
\item \emph{(the chain)} $\Fwp\dsq_{\Lnoise}\Ffss$, and the channels compose
  exactly, $\Lfss=\Lnoise\circ\Lwp$ on all of $\Tcl{\Bsys}$, so that
  \begin{equation}\label{eq:ladder}
    M\ \dsq_{\Lwp}\ \Fwp\ \dsq_{\Lnoise}\ \Ffss\ \dsq_{\Lmp}\ \Ppvm ;
  \end{equation}
\item \emph{(the encoded constraint)} $\Fwp_v\succeq p_v\,\Pflag$ as operators
  on $\Qsys$, so $\tr[\Fwp_v\sigma]\ge p_v\tr[\Pflag\sigma]$ holds on
  \emph{every} squashed state, reachable or not. Under $\Ffss$ the same
  inequality is only a pull-back constraint: it holds on every reachable state,
  since
  $\Lfss^{\dagger}\big(\Ffss_v-p_v\Pflag\big)
  =\blk{M}{\leqN}_v\dsum\big(\blk{M}{\gtN}_v-p_v\idon{\Bsys_{\gtN}}\big)
  \succeq0$, and it \emph{fails} at $\sigma=\dyad{v'}_{\Flag}$ for any $v'\ne v$
  with $p_v>0$;
\item \emph{(no reverse relation in general)} $\Ffss\dsq\Fwp$ can fail: at
  $\Hs{\Bsys_{\leqN}}=\Hs{\Bsys_{\gtN}}=\mathbb{C}$, $\Vset=\{1,2\}$,
  $M_1=\tfrac12\dsum1$, $M_2=\tfrac12\dsum0$, $p_1=\tfrac12$, $p_2=0$, one has
  $\|\Ffss_2\|=1>\tfrac12=\|\Fwp_2\|$, which
  Lemma~\ref{lem:obstructions}\,(i) forbids.
\end{enumerate}
\end{proposition}
\end{mainclaim}
\claimmeta{\metalabel{Uses}\mref{Lem.}{lem:sep}, \mref{Cor.}{cor:pullback}, \mref{Lem.}{lem:obstructions}, \mref{Lem.}{lem:top}, \mref{Lem.}{lem:master}, \mref{Cor.}{cor:instances}, \mref{Def.}{def:targets}
\metasep \metalabel{Used by}\mref{Rem.}{rem:tworoutes}, \mref{Prop.}{prop:weight-monotone}, \mref{Rem.}{rem:weightfree}, \mref{Rem.}{rem:published}, \mref{Prop.}{prop:precomposed-blocks}, \mref{Cor.}{cor:fssinstance}, \mref{Prop.}{prop:flat-instance}, \mref{Cor.}{cor:emitted-flat}, \mref{Cor.}{cor:two-sector-decided}, \mref{Cor.}{cor:flat-companion}, \mref{Req.}{req:total-weight-refusal}
\metasep newly derived
\metasep accepted
\metasep \metalabel{Oracles}\coderef{https://github.com/fraunhofersg/certiqs-key/blob/v0.1.0/tests/oracles/registry_reductions.py}{registry_reductions.py}}

\begin{proof}
(i) \emph{$\{\Mwp_v\}$ is a POVM on $\Bsys_{\gtN}$}: $\Mwp_v\succeq0$ by
\eqref{eq:certweights}, and
$\sum_v\Mwp_v=\big(\idon{\Bsys_{\gtN}}-p\,\idon{\Bsys_{\gtN}}\big)/(1-p)
=\idon{\Bsys_{\gtN}}$ by block completeness.

\emph{$\Lwp$ is a channel.} It is the map \eqref{eq:master} of
Lemma~\ref{lem:master} for the pinching instrument
$\Ins_1(\rho)=\Pleq\rho\Pleq$, $\Ins_2(\rho)=\Pgt\rho\Pgt$
(Corollary~\ref{cor:instances}\,(ii)), the identity channel on branch~1, and
the measure-and-prepare channel $\Lmp$ of the POVM $\Mwp$
(Lemma~\ref{lem:top}) on branch~2; part~(ii) of that lemma gives complete
positivity and trace preservation. The POVM that $\Lwp$ squashes is
computed directly below rather than read off
Lemma~\ref{lem:master}\,(iii): for these branch data that part names the
coarse-grained target $\blk{M}{\leqN}_v\dsum\dyad{v}$, which is the flag-state
target of the \emph{modified} POVM $\Mwp$, whereas $\Fwp$ is a different target
on the same space, and the weights $p_v$ are what the difference records.

\emph{The relation.} The output of $\Lwp$ is block
diagonal, so only the diagonal blocks of an observable contribute to
\eqref{eq:pairing}, and for $X=X_{\leqN}\dsum X_{\Flag}$,
\begin{equation}\label{eq:WPadjoint}
  \Lwp^{\dagger}\big(X_{\leqN}\dsum X_{\Flag}\big)
  = X_{\leqN}\ \dsum\ \sum_v\bra{v}X_{\Flag}\ket{v}\;\Mwp_v .
\end{equation}
Applying \eqref{eq:WPadjoint} to $\Fwp_v$, whose flag block is
$p_v\idon{\Flag}+(1-p)\dyad{v}$ with diagonal entries
$\bra{v'}\cdot\ket{v'}=p_v+(1-p)\delta_{vv'}$,
\begin{align*}
  \Lwp^{\dagger}\big(\Fwp_v\big)
  &= \blk{M}{\leqN}_v\ \dsum\
    \Big(p_v\sum_{v'}\Mwp_{v'}+(1-p)\,\Mwp_v\Big)\\
  &= \blk{M}{\leqN}_v\ \dsum\
    \big(p_v\,\idon{\Bsys_{\gtN}}+\blk{M}{\gtN}_v-p_v\,\idon{\Bsys_{\gtN}}\big)
  \;=\; M_v .
\end{align*}
For the transported projectors, $X=\Pleq$ means $X_{\leqN}=\idon{}$,
$X_{\Flag}=0$, and \eqref{eq:WPadjoint} returns $\Pleq$; $X=\Pflag$ means
$X_{\leqN}=0$, $X_{\Flag}=\idon{\Flag}$, and it returns
$0\dsum\sum_v\Mwp_v=\Pgt$. The statement for $\Lfss$ is the case
$p_v\equiv0$ of the same computation, with $\Mwp_v=\blk{M}{\gtN}_v$; note that
$\Ffss$ is a POVM because $\sum_v\Ffss_v=\idon{\Qsys}$, and so is $\Fwp$
because $\sum_v\Fwp_v=\idon{\Bsys_{\leqN}}\dsum
\big(p\,\idon{\Flag}+(1-p)\idon{\Flag}\big)=\idon{\Qsys}$.

(ii) $\Lnoise$ is a channel: it is the convex combination, with weights $1-p$
and $p$, of the identity and of the channel that pinches onto the two blocks
and replaces the flag block by the fixed state $\mu/p$ (for $p=0$,
$\Lnoise=\idmap$); trace preservation is
$\tr\Lnoise(\sigma)=(1-p)\tr\sigma+p\,\tr[\Pleq\sigma]
+\tr[\Pflag\sigma]\,\tr\mu
=(1-p)\tr\sigma+p\big(\tr[\Pleq\sigma]+\tr[\Pflag\sigma]\big)=\tr\sigma$,
using $\tr\mu=p$ and $\Pleq+\Pflag=\idon{\Qsys}$. Its adjoint is
$\Lnoise^{\dagger}(X)=(1-p)X+p\,\Pleq X\Pleq+\tr[\mu X]\,\Pflag$, so on block
observables
\begin{equation}\label{eq:L2adjoint}
  \Lnoise^{\dagger}\big(X_{\leqN}\dsum X_{\Flag}\big)
  = X_{\leqN}\ \dsum\
    \big((1-p)\,X_{\Flag}+\tr[\mu X_{\Flag}]\,\idon{\Flag}\big),
\end{equation}
whence $\Lnoise^{\dagger}(\Ffss_v)
=\blk{M}{\leqN}_v\dsum\big((1-p)\dyad{v}+p_v\,\idon{\Flag}\big)=\Fwp_v$, which
is the relation. For the composition, $\Lwp(\rho)$ is block diagonal, so
\begin{align*}
  \Lnoise\big(\Lwp(\rho)\big)
  &=\Pleq\rho\Pleq\ \dsum\
    \sum_v\Big[(1-p)\tr\big[\Mwp_v\Pgt\rho\Pgt\big]
      +p_v\,\tr\big[\Pgt\rho\Pgt\big]\Big]\dyad{v}\\
  &=\Pleq\rho\Pleq\ \dsum\
    \sum_v\tr\!\Big[\big((1-p)\Mwp_v+p_v\,\idon{\Bsys_{\gtN}}\big)\,
      \Pgt\rho\Pgt\Big]\dyad{v}
   \;=\;\Lfss(\rho),
\end{align*}
since $(1-p)\Mwp_v+p_v\idon{}=\blk{M}{\gtN}_v$; this holds for every $\rho$,
including states with coherences across the cutoff, which both sides discard.
The last relation in \eqref{eq:ladder} is Lemma~\ref{lem:top}.

(iii) $\Fwp_v-p_v\Pflag
=\blk{M}{\leqN}_v\dsum(1-p)\dyad{v}\succeq0$. The pull-back identity is
\eqref{eq:WPadjoint} at $p_v\equiv0$ applied to $\Ffss_v-p_v\Pflag$, whose
flag block has diagonal entries $\delta_{vv'}-p_v$, giving
$\blk{M}{\leqN}_v\dsum\big(\blk{M}{\gtN}_v-p_v\idon{}\big)$, which is
$\succeq0$ by \eqref{eq:certweights}; so
Corollary~\ref{cor:pullback}\,(ii) applies on every reachable state. For the
failure, $\tr\big[(\Ffss_v-p_v\Pflag)\dyad{v'}\big]=0-p_v<0$ when $v'\ne v$ and
$p_v>0$.

(iv) With the stated data, \eqref{eq:certweights} holds ($1\succeq\tfrac12$ and
$0\succeq0$) and $p=\tfrac12<1$. Then
$\Ffss_2=\tfrac12\dsum\dyad{2}$ has norm $1$, while
$\Fwp_2=\tfrac12\dsum\big(0\cdot\idon{\Flag}+\tfrac12\dyad{2}\big)$ has norm
$\tfrac12$. If $\Ffss\dsq\Fwp$ held, Lemma~\ref{lem:obstructions}\,(i) applied
with $\Ffss$ in the physical role would give $\|\Ffss_2\|\le\|\Fwp_2\|$, i.e.
$1\le\tfrac12$.
\end{proof}

% claim reductions/wpfss/def:routes
\claimsummary{reductions/wpfss/def:routes}{The encoded route uses the weight-preserving target and emits no constraint; the licensed route uses the flag-state target and emits the licensed entry.}
\claimbookmark{reductions/wpfss/def:routes}{Definition}{def:routes}{The encoded and the licensed route}
\claimleadin{%
An artefact carries the device knowledge above the cutoff either inside its
target or as a constraint entry beside it; the two choices are named here.
}
\begin{definition}[The encoded and the licensed route]\label{def:routes}
The \emph{encoded} route uses the \termref{term-weight-preserving-target}{weight-preserving target} $\Fwp$ and
emits $\CList=\varnothing$; the \emph{licensed} route uses the
\termref{term-flag-state-target}{flag-state target} $\Ffss$ and emits the entry of
Theorem~\ref{thm:licence}.
\end{definition}
\claimmeta{\metalabel{Uses}\mref{Def.}{def:targets}
\metasep \metalabel{Used by}\mref{Rem.}{rem:tworoutes}, \mref{Prop.}{prop:flat-instance}
\metasep newly derived
\metasep accepted}

% claim reductions/wpfss/rem:tworoutes
\claimsummary{reductions/wpfss/rem:tworoutes}{Both routes are sound, and the encoded one, which needs no attestation, is preferable where it is available. Pairing the licensed entry with the encoded route's channel is unsound, so an artefact takes one route.}
\claimbookmark{reductions/wpfss/rem:tworoutes}{Remark}{rem:tworoutes}{Soundness of the two routes and one route per artefact}
\claimleadin{%
Part~(iii) of Proposition~\ref{prop:chain} is the whole difference between the
two targets. Under $\Fwp$ the device fact is an operator inequality \emph{on
the target}, true of every state the adversary could possibly prepare, so it
needs no constraint, no licence and no attestation. Under $\Ffss$ the same
fact is true only on the image, so it must travel as a declared constraint
entry, and something must certify that the adversary cannot leave the image,
which is exactly Theorem~\ref{thm:licence}. Part~(ii) of the proposition says
which of the two hands the adversary more: by Corollary~\ref{cor:monotone} a
cleaner target never helps, and $\Ffss$ is cleaner than $\Fwp$, in general
strictly so: part~(iv) gives data on which the reverse relation fails.
}
\begin{remark}[Soundness of the two routes and one route per artefact]\label{rem:tworoutes}
Both routes of Definition~\ref{def:routes} are sound, and an artefact takes
exactly one of them (Requirement~\ref{req:one-route-per-artifact}). The encoded route is
preferable wherever it is available, because it requires no attestation a
program cannot check. It is \emph{not} always available:
Proposition~\ref{prop:obstruction} exhibits a device fact, a bound on a
\emph{union} of outcomes that must not be merged, that no target classical on
the whole of $\Qsys$ can carry and that $\Fwp$, one weight per outcome, does
not carry, and for which the licensed route is the only route; this is why
every artefact carries a constraint slot and a licence slot
(Requirement~\ref{req:constraint-and-licence-slots}).

\emph{Why one route per artefact.} The two routes run on different channels,
$\Lwp$ and $\Lfss$, and by Remark~\ref{rem:channel-choice}\,(c) a constraint
sound for one channel need not be sound for another. Here it is not: pairing
the licensed entry with the encoded route's channel gives, by
\eqref{eq:WPadjoint},
\begin{equation}\label{eq:crosspairing}
  \Lwp^{\dagger}\big(\Ffss_o-p_o\Pflag\big)
  =\blk{M}{\leqN}_o\ \dsum\
   \Big(\tfrac{1}{1-p}\big(\blk{M}{\gtN}_o-p_o\idon{}\big)-p_o\idon{}\Big),
\end{equation}
whose flag block is \emph{not} implied positive by \eqref{eq:certweights}: at
the tight end $\blk{M}{\gtN}_o=p_o\idon{}$ it equals $-p_o\idon{}\prec0$ for
$p_o>0$. So that pairing would impose a constraint some reachable state
violates, voiding Theorem~\ref{thm:transfer}, which is exactly the failure
Remark~\ref{rem:unsound} describes. Since an artefact's payload is
single-valued (Requirement~\ref{req:one-route-per-artifact}) and
\eqref{eq:link} binds each link's constraint list to \emph{that link's}
\termref{term-channel-witness}{channel witness}, the remedy is structural: one route per artefact, and
two requests for a caller who wants both bundles, with no schema in which the
unsound pairing can be expressed.
\end{remark}
\claimmeta{\metalabel{Uses}\mref{Thm.}{thm:transfer}, \mref{Rem.}{rem:unsound}, \mref{Rem.}{rem:channel-choice}, \mref{Cor.}{cor:monotone}, \mref{Def.}{def:targets}, \mref{Prop.}{prop:chain}, \mref{Def.}{def:routes}
\metasep \metalabel{Used by}\mref{Req.}{req:one-route-per-artifact}
\metasep newly derived
\metasep accepted}

% claim reductions/wpfss/prop:weight-monotone
\claimsummary{reductions/wpfss/prop:weight-monotone}{Certified weights stay certified when lowered but not below zero, and while the total weight is below one the weight-preserving target of the larger weights squashes to that of the smaller ones. By \hyperlink{claim-reductions/cor:monotone}{Corollary~\ref*{cor:monotone}} the smaller weights give the cleaner target, which hands the adversary more.}
\claimbookmark{reductions/wpfss/prop:weight-monotone}{Proposition}{prop:weight-monotone}{Monotonicity of the weight-preserving target in its weights}
\claimleadin{%
Nothing in Definition~\ref{def:targets} or Proposition~\ref{prop:chain}
requires $p_v$ to be the largest admissible value. The proposition says what
lowering the weights does: they stay valid, and the target becomes cleaner, a
statement about the pre-order that is proved rather than asserted.
}
\begin{proposition}[Monotonicity of the weight-preserving target in its weights]
\label{prop:weight-monotone}
Let $(p_v)_{v\in\Vset}$ be \termref{term-certified-weights}{certified weights} for $M$ with
\termref{term-total-weight}{total weight} $p$, let $p'_v\in[0,p_v]$ for every $v\in\Vset$, and
write $p'\coloneqq\sum_vp'_v$. Write $\Fwp[p]$ and $\Fwp[p']$ for the
\termref{term-weight-preserving-target}{weight-preserving targets} \eqref{eq:WPFSS}
built from $(p_v)$ and from $(p'_v)$.
\begin{enumerate}[label=(\roman*),leftmargin=3.0em]
\item $(p'_v)$ are certified weights for $M$, and $p'\le p$; in particular,
  where $p<1$ the construction of Proposition~\ref{prop:chain} runs
  unchanged with $(p'_v)$.
\item If $p<1$, put
  \begin{equation}\label{eq:weightmono}
    \Psi(\sigma)\coloneqq\lambda\,\sigma+(1-\lambda)\,\Pleq\sigma\Pleq
     +\tr\!\big[\Pflag\sigma\big]\,\nu,
    \qquad
    \lambda\coloneqq\frac{1-p}{1-p'},
    \qquad
    \nu\coloneqq\sum_v\frac{p_v-p'_v}{1-p'}\dyad{v}_{\Flag}.
  \end{equation}
  Then $\Psi$ is a channel and $\Fwp[p]\dsq_{\Psi}\Fwp[p']$, so by
  Corollary~\ref{cor:monotone} the smaller-weight target is the cleaner one.
\end{enumerate}
\end{proposition}
\claimmeta{\metalabel{Uses}\mref{Cor.}{cor:monotone}, \mref{Def.}{def:targets}, \mref{Prop.}{prop:chain}
\metasep \metalabel{Used by}\mref{Rem.}{rem:weightfree}, \mref{Prop.}{prop:route}, \mref{Rem.}{rem:minweight}, \mref{Req.}{req:largest-certified-weight}
\metasep newly derived
\metasep accepted}

\begin{proof}
(i) By \eqref{eq:certweights},
$\blk{M}{\gtN}_v\succeq p_v\idon{\Bsys_{\gtN}}\succeq p'_v\idon{\Bsys_{\gtN}}$
with $p'_v\ge0$, for every $v$, and summing $p'_v\le p_v$ gives $p'\le p$; so
where $p<1$ also $p'<1$, and Proposition~\ref{prop:chain} applies to
$(p'_v)$.

(ii) \emph{$\Psi$ is a channel.} $\lambda\in[0,1]$ and $\nu\succeq0$ because
$p'\le p<1$ and $p'_v\le p_v$, so $\Psi$ is CP; and with
$\tr\nu=(p-p')/(1-p')$ and $\Pleq+\Pflag=\idon{\Qsys}$,
$\tr\Psi(\sigma)=\tr[\Pleq\sigma]+\big(\lambda+\tr\nu\big)\tr[\Pflag\sigma]
=\tr\sigma$, since
$\lambda+\tr\nu=\big((1-p)+(p-p')\big)/(1-p')=1$.

\emph{The adjoint.} On block observables
$\Psi^{\dagger}(X_{\leqN}\dsum X_{\Flag})
=X_{\leqN}\dsum\big(\lambda X_{\Flag}+\tr[\nu X_{\Flag}]\idon{\Flag}\big)$, and
evaluating at $X_{\Flag}=p'_v\idon{\Flag}+(1-p')\dyad{v}$ gives flag block
\begin{equation*}
  \Big[\tfrac{(1-p)p'_v}{1-p'}+\tfrac{p'_v(p-p')}{1-p'}+(p_v-p'_v)\Big]\idon{\Flag}
  +(1-p)\dyad{v}
  \;=\;p_v\,\idon{\Flag}+(1-p)\dyad{v},
\end{equation*}
because the first two brackets sum to $p'_v$; with the unchanged $\leqN$ block
this is $\Fwp[p]_v$. The conclusion is Corollary~\ref{cor:monotone} applied
to $\Fwp[p]\dsq_{\Psi}\Fwp[p']$.

Taking $p'\equiv0$ recovers the flag-noise channel $\Lnoise$ of
\eqref{eq:Lnoise}, as it must.
\end{proof}

% claim reductions/wpfss/rem:weightfree
\claimsummary{reductions/wpfss/rem:weightfree}{Every smaller non-negative weight is also certified, but it gives a cleaner target, which hands the adversary more (\hyperlink{claim-reductions/wpfss/prop:weight-monotone}{Proposition~\ref*{prop:weight-monotone}}). The largest weight the artefact certifies is therefore the one used, and a smaller one stays valid.}
\claimbookmark{reductions/wpfss/rem:weightfree}{Remark}{rem:weightfree}{Why the largest certified weight is used}
\begin{remark}[Why the largest certified weight is used]
\label{rem:weightfree}
Nothing in Definition~\ref{def:targets} or Proposition~\ref{prop:chain}
requires $p_v$ to be the largest admissible value: every
$p'_v\in[0,p_v]$ is again a certified weight, and the construction runs
unchanged with $p'=\sum_vp'_v\le p$
(Proposition~\ref{prop:weight-monotone}\,(i)). A smaller weight is not free,
however: it gives a \emph{cleaner} target, which by
Corollary~\ref{cor:monotone} hands the adversary more
(Proposition~\ref{prop:weight-monotone}\,(ii)). The largest weight the
artefact certifies is therefore the best choice, and any smaller one is valid
but weaker (Requirement~\ref{req:largest-certified-weight}).
\end{remark}
\claimmeta{\metalabel{Uses}\mref{Cor.}{cor:monotone}, \mref{Def.}{def:targets}, \mref{Prop.}{prop:chain}, \mref{Prop.}{prop:weight-monotone}
\metasep \metalabel{Used by}\mref{Req.}{req:largest-certified-weight}
\metasep newly derived
\metasep accepted}

% claim reductions/wpfss/rem:published
\claimsummary{reductions/wpfss/rem:published}{The flag-state target is the flag-state squasher of Zhang et al., and the single-bound weight-preserving target that of Nahar et al. Their usual scalar flag-weight constraint is a carried-operator entry stronger than the default licensed one, and only its operator form is ever emitted.}
\claimbookmark{reductions/wpfss/rem:published}{Remark}{rem:published}{Relation to the published flag-state squashers}
\begin{remark}[Relation to the published flag-state squashers]
\label{rem:published}
$\Ffss$ is the flag-state squasher of \cite[Thm.~1, Eq.~(8)]{zhang2021security},
which states exactly \eqref{eq:FSS}: $\widetilde M_y=M_{y,n\le k}\oplus
\dyad{y}$ for a POVM block diagonal at a photon-number cutoff. The adjoint
form is stated there in part: \cite[Eqs.~(149)--(150)]{tupkary2026decoy}
gives, for that map, the defining identity $\Lambda^{\dagger}(F_i)=M_i$ and the
\emph{flag} half of \eqref{eq:transportproj},
$\Lambda^{\dagger}(\Pflag)=\Pgt$, and states no second projector identity. The
preserved half follows from it in one step and is written out here rather than
attributed: $\Lambda$ is trace preserving, so $\Lambda^{\dagger}$ is unital, and
$\Pleq+\Pflag=\idon{\Qsys}$ with $\Pleq+\Pgt=\idon{\Bsys}$ give
$\Lambda^{\dagger}(\Pleq)=\idon{\Bsys}-\Lambda^{\dagger}(\Pflag)
=\idon{\Bsys}-\Pgt=\Pleq$. This contract proves both halves at
Proposition~\ref{prop:chain}\,(i) and the citation is corroboration, which is
why the narrower attribution costs nothing. The \termref{term-single-bound-case}{single-bound case} of
the \termref{term-weight-preserving-target}{weight-preserving target} $\Fwp$
is the weight-preserving flag-state squasher of
\cite[Lemma~7, Eqs.~(28)--(29)]{nahar2024postselection}, whose target places
$f_N$ on the diagonal of the flag block of one distinguished element and
$(1-f_N)$ on the others; the multi-bound form \eqref{eq:WPFSS} is the
symmetric statement with one weight per outcome.

The scalar flag-weight constraint these targets are usually paired with is
\cite[Eq.~(27)]{nahar2024postselection}, and the comparison with what this
module emits runs in the direction opposite to the naive one. Writing
$a\coloneqq\lmin\big(\blk{M}{\leqN}_e\big)$ and
$b\coloneqq\lmin\big(\blk{M}{\gtN}_e\big)$ for a designated event $e$
(declared device constants, fixed before the run and never recomputed from
the declaration, Assumption~\ref{ass:constdecl}), that
equation reads, in full, $W\le1-p(m\le N)\le\big(p(e)-a\big)/(b-a)$ for the
flag weight $W=\tr[\Pflag\sigma]$ and the observed
$p(e)=\tr[\Ffss_e\sigma]$;
the right-hand member is the operative one and the only one compared with
below. It is
\emph{a} member of the family Corollary~\ref{cor:transported} generates,
namely the carried operator $\hat O$, which carries the certificate
\begin{equation}\label{eq:two-sector-entry}
\begin{split}
  \hat O&=\Ffss_e-a\,\Pleq-b\,\Pflag,\\
  \Lfss^{\dagger}\big(\hat O\big)&=M_e-a\,\Pleq-b\,\Pgt
  =\big(\blk{M}{\leqN}_e-a\idon{}\big)
    \dsum\big(\blk{M}{\gtN}_e-b\idon{}\big)
  \;\succeq\;0
\end{split}
\end{equation}
by the definitions of $a$ and $b$ and the transported projectors of
\eqref{eq:transportproj}. The emitted constant is $0$, with no shift
(Corollary~\ref{cor:entries}\,(iii)).
Corollary~\ref{cor:pullback}\,(ii) gives
$p(e)-a(1-W)-bW\ge0$, which rearranges to Eq.~(27) when $b>a$ (the published
division is by $b-a$; at $b\le a$ it reverses or is undefined, and only the
operator form above is used here). It is \emph{not} the member
Proposition~\ref{prop:chain}\,(iii) supplies and not the one
Corollary~\ref{cor:fssinstance} emits: those use $a=0$, where the same
corollary gives only $p(e)-bW\ge0$.

Which of the two is stronger is settled at the level of the \emph{constraint
sets}, not by comparing the two caps on $W$: since
$\hat O=\hat O\big|_{a=0}-a\,\Pleq$ with $a\ge0$ and $\Pleq\succeq0$, one has
$\tr[\hat O|_{a=0}\,\sigma]\ge\tr[\hat O\,\sigma]$ for every $\sigma\succeq0$,
so $\{\sigma:\tr[\hat O\sigma]\ge0\}\subseteq
\{\sigma:\tr[\hat O|_{a=0}\sigma]\ge0\}$: the $a>0$ entry cuts a set no
larger and is at least as strong a constraint, unconditionally. The implication
therefore runs from Eq.~(27) to the $a=0$ form and not the other way. (The
comparison of the two \emph{caps} $\big(p(e)-a\big)/(b-a)$ and $p(e)/b$ is not
the right one and does not run in a fixed direction: the former is the smaller
exactly when $p(e)<b$, because $p(e)$ is itself a function of $\sigma$ rather
than a fixed number, so the two inequalities cut the state set differently.)

Two consequences follow. \emph{(a)} The default licensed entry of
Corollary~\ref{cor:fssinstance} is weaker than standard practice, not
stronger; a declaration may request the entry \eqref{eq:two-sector-entry}
as a carried operator with declared constants $a$ and $b$, which need not be
the least eigenvalues above. Its certificate is then not taken from the
request: it is decided on the emitted numbers, where every emitted block of
$e$ below the cutoff clears $a$ and $b$ is at most the certified weight
$p_e$ (Corollary~\ref{cor:two-sector-decided}), and the request is refused
where either fails. Both constants must be declared on the lattice of
\eqref{eq:latinv}, as every declared constant must: an $a$ or $b$ off the
lattice puts the requested operator's entries off the lattice, and the
request is refused (Requirement~\ref{req:licensed-entry-lattice}). Smaller
lattice values of $a$ and $b$ keep both conditions. \emph{(b)} The \emph{cap} in
the displayed form is never emitted: its right-hand side contains the run's
observed $p(e)$, whereas Assumption~\ref{ass:constdecl} admits only constants
declared before the run, and an entry with such a constant is refused
(Requirement~\ref{req:declared-constants-only}). The admissible object is the
operator entry \eqref{eq:two-sector-entry} with $\bowtie\,=\,\succeq$ and $c=0$, of
which the cap is a run-time consequence and not a declaration. Every statement
of this remark is proved in the claims it names
(Definition~\ref{def:targets}, Proposition~\ref{prop:chain},
Corollary~\ref{cor:fssinstance}, Corollary~\ref{cor:two-sector-decided} and
Corollary~\ref{cor:entries}); the
citations record where a reader finds the correspondence.
\end{remark}
\claimmeta{\metalabel{Uses}\mref{Cor.}{cor:pullback}, \mref{Cor.}{cor:transported}, \mref{Def.}{def:targets}, \mref{Prop.}{prop:chain}
\metasep \metalabel{Assumes}\aref{A-CONST}{ass:constdecl}
\metasep newly derived
\metasep accepted}

\section{The weight route}\label{group-reductions/wpfss/lem:descent}
\sectionintro{The receiver supplies coarse-grained blocks, a dark-count channel on raw click patterns and per-setting weights for raw outcome sets. The dark-count channel can act on the coarse blocks only if it descends through the coarse-graining (\hyperlink{claim-reductions/wpfss/lem:descent}{Lemma~\ref*{lem:descent}}); the \termref{term-joint-family}{joint family} of the private setting is then formed, and each certified weight is routed to its bounded letter, which must be absorbing, sharpened by the dark counts and rounded down onto the lattice (\hyperlink{claim-reductions/wpfss/prop:route}{Proposition~\ref*{prop:route}}). Blocks supplied with the dark counts already composed form the same instance wherever no snap intervenes (\hyperlink{claim-reductions/wpfss/prop:precomposed-blocks}{Proposition~\ref*{prop:precomposed-blocks}}). The remarks show why the constancy, absorbing and fibre conditions cannot be dropped and which receiver claims are imported.}
\sectionimports{\extref{../foundations/contract.pdf}{foundations/def:exact-reference}, \extref{../foundations/contract.pdf}{foundations/def:snap}, \extref{../receiver/contract.pdf}{receiver/cor:unconditioned-family}, \extref{../receiver/contract.pdf}{receiver/def:setting-family}, \extref{../receiver/contract.pdf}{receiver/def:slot}, \extref{../receiver/contract.pdf}{receiver/detector-imperfections/def:darkadm}, \extref{../receiver/contract.pdf}{receiver/detector-imperfections/def:darktol}, \extref{../receiver/contract.pdf}{receiver/detector-imperfections/prop:efficiency-normalization}, \extref{../receiver/contract.pdf}{receiver/detector-imperfections/thm:dark-counts}, \extref{../receiver/contract.pdf}{receiver/prop:slotvalid}, \extref{../receiver/contract.pdf}{receiver/thm:click-povm}, \extref{../receiver/contract.pdf}{receiver/thm:coarse-graining}, \extref{../receiver/contract.pdf}{receiver/thm:emitted-bound}}

% claim reductions/wpfss/lem:descent
\claimsummary{reductions/wpfss/lem:descent}{R descends through a coarse-graining when each fibre sum depends on a column only through its letter; then R and the coarse-graining commute via a column-stochastic R\textasciicircum{}cg. A monotone R descends if the coarse-graining is injective off the multiclick set and maps exactly that set to one letter.}
\claimbookmark{reductions/wpfss/lem:descent}{Lemma}{lem:descent}{Descent of a dark-count matrix through a coarse-graining}
\claimleadin{%
The receiver emits its blocks on the coarse alphabet $\Yset$ but supplies its
\exttermref{../receiver/contract.pdf}{term-dark-count-channel}{dark-count channel} on the raw click patterns $2^{[K]}$. This lemma says
when the channel can be moved onto the coarse alphabet through the
\exttermref{../receiver/contract.pdf}{term-declared-coarse-graining}{declared coarse-graining} $\cg$, and gives the sufficient condition the
canonical threshold-detector declaration meets; the constancy clause of that
condition cannot be dropped (Remark~\ref{rem:descentneedsconstant}).
}
\begin{lemma}[Descent of a dark-count matrix through a coarse-graining]
\label{lem:descent}
Let $R$ be column-stochastic on $2^{[K]}$ and $\cg\colon2^{[K]}\to\Yset$
surjective. Say that \emph{$R$ descends through $\cg$} if the quantity
$P(y\,|\,S')\coloneqq\sum_{S\in\cg^{-1}(y)}R_{S|S'}$ depends on $S'$ only
through $\cg(S')$; the resulting matrix
$R^{\cg}_{y|y'}\coloneqq P(y\,|\,S')$ for any $S'\in\cg^{-1}(y')$ is then well
defined and column-stochastic. In that case, for every family of operators
$\{X_{S}\}$,
\begin{equation}\label{eq:descend}
  \sum_{S\in\cg^{-1}(y)}\ \sum_{S'}R_{S|S'}X_{S'}
  \;=\;\sum_{y'\in\Yset}R^{\cg}_{y|y'}\ \sum_{S'\in\cg^{-1}(y')}X_{S'},
\end{equation}
so applying $R$ and then coarse-graining agrees with coarse-graining and then
applying $R^{\cg}$. A sufficient condition, which the canonical
threshold-detector declaration satisfies, is that $R$ be \emph{monotone}
($R_{S|S'}>0\Rightarrow S'\subseteq S$) and that, with
$\MCK\coloneqq\{S:|S|\ge2\}$,
\begin{equation}\label{eq:canonicalcg}
  \cg\ \text{is injective on }2^{[K]}\setminus\MCK,
  \qquad
  \cg\ \text{is \emph{constant} on }\MCK,\ \text{say }\cg(\MCK)=\{y_{\mc}\},
  \qquad
  \cg^{-1}(y_{\mc})=\MCK ;
\end{equation}
the case $\cg=\mathrm{id}$ is trivial.

On declared data both the hypothesis and the conclusion are read exactly. The
declared $R$ is exactly column-stochastic, and the descent
condition, that $P(y|S')=P(y|S'')$ for every $y$ and every pair $S',S''$ with
$\cg(S')=\cg(S'')$, is an identity between the dyadic rationals the
declaration carries and is decided as one
(Definition~\ref{def:substrate}).
The representative is declared, not chosen:
$R^{\cg}_{y|y'}\coloneqq P(y|S'_{y'})$ with $S'_{y'}$ the first element of
$\cg^{-1}(y')$ in the \exttermref{../foundations/contract.pdf}{term-declared-order}{declared order} of $2^{[K]}$, so the descended
matrix is a determinate object and not one of several. Column-stochasticity of
$R^{\cg}$ needs no second decision: the identity
$\sum_yR^{\cg}_{y|y'}=\sum_SR_{S|S'}$ of the proof regroups one finite sum,
and in the exact substrate a regrouped sum \emph{is} the sum it regroups. So \eqref{eq:descend} holds of the declared
numbers as an exact identity.
\end{lemma}
\claimmeta{\metalabel{Used by}\mref{Rem.}{rem:descentneedsconstant}, \mref{Def.}{def:weight-route}, \mref{Prop.}{prop:route}, \mref{Prop.}{prop:precomposed-blocks}, \mref{Rem.}{rem:colstoch}, \mref{Rem.}{rem:darkcompose}, \mref{Rem.}{rem:receiver-imports}, \mref{Prop.}{prop:flat-instance}, \mref{Cor.}{cor:emitted-flat}, \mref{Cor.}{cor:flat-companion}, \mref{Req.}{req:descent-decided-exactly}, \mref{Req.}{req:column-stochastic-refusal}
\metasep newly derived
\metasep accepted}

\begin{proof}
Well-definedness is the stated hypothesis, and
$\sum_yR^{\cg}_{y|y'}=\sum_yP(y|S')=\sum_SR_{S|S'}=1$ gives
column-stochasticity. For \eqref{eq:descend}, exchange the two sums on the left
and group the columns by $\cg(S')$:
$\sum_{S'}P(y|S')X_{S'}
=\sum_{y'}\sum_{S'\in\cg^{-1}(y')}R^{\cg}_{y|y'}X_{S'}$.

For the sufficient condition \eqref{eq:canonicalcg}: if $S'\in\MCK$ and
$R_{S|S'}>0$ then $S\supseteq S'$ by monotonicity, so $|S|\ge2$ and $S\in\MCK$;
hence the whole column $S'$ is supported on $\MCK$, which by the constancy and
preimage clauses is exactly $\cg^{-1}(y_{\mc})$. Column-stochasticity then gives
$P(y_{\mc}|S')=\sum_{S\in\MCK}R_{S|S'}=1$ and $P(y|S')=0$ for $y\ne y_{\mc}$,
the same values for \emph{every} $S'\in\MCK$. This is where constancy of
$\cg$ on $\MCK$ is needed: without it, $\MCK$ would split across several
letters and the individual $P(y|S')$ could differ inside $\MCK$
(Remark~\ref{rem:descentneedsconstant}). If $S'\notin\MCK$ then $\cg(S')$
determines $S'$: by the preimage clause $\cg(S')\ne y_{\mc}$, so by the
constancy clause no pattern of $\MCK$ has the letter $\cg(S')$, and $\cg$ is
injective on $2^{[K]}\setminus\MCK$. So $P(y|S')$ depends on $S'$
only through $\cg(S')$ in both cases.
\end{proof}

% claim reductions/wpfss/rem:descentneedsconstant
\claimsummary{reductions/wpfss/rem:descentneedsconstant}{With four ports and independent dark counts, one coarse-graining injective off the multiclick set that splits that set over two letters has two columns in one fibre with different fibre sums, so R\textasciicircum{}cg is undefined. The constancy clause of \hyperlink{claim-reductions/wpfss/lem:descent}{Lemma~\ref*{lem:descent}} cannot be left unchecked.}
\claimbookmark{reductions/wpfss/rem:descentneedsconstant}{Remark}{rem:descentneedsconstant}{Necessity of the constancy clause}
\begin{remark}[Necessity of the constancy clause]
\label{rem:descentneedsconstant}
The constancy clause of \eqref{eq:canonicalcg} cannot be dropped: without it
the conclusion of Lemma~\ref{lem:descent} fails. Take $K=4$, independent
per-port dark counts of common probability $p\in(0,1)$ (the
\exttermref{../receiver/contract.pdf}{term-independent-product-form}{independent product form} with the same probability at every port), and
let $\cg$ send $\{1,2\},\{1,2,4\},\{1,2,3,4\}\mapsto y_a$, every other
multiclick pattern to $y_b$, and be injective with distinct fresh letters off
$\MCK$. Then $\cg$ is injective on $2^{[K]}\setminus\MCK$ and
$\cg^{-1}(\{y_a,y_b\})=\MCK$, and $R$ is monotone, but
\begin{equation*}
  P\big(y_a\,\big|\,\{1,2\}\big)=(1-p)^2+p(1-p)+p^2=1-p+p^2,
  \qquad
  P\big(y_a\,\big|\,\{1,2,4\}\big)=(1-p)+p=1,
\end{equation*}
two columns with the same $\cg$-image and different values, so $R^{\cg}$ is not
well defined. The first count is over the supersets of $\{1,2\}$ lying in
$\cg^{-1}(y_a)$, namely $\{1,2\}$ itself, $\{1,2,4\}$ and $\{1,2,3,4\}$; the
second over the supersets of $\{1,2,4\}$ there, namely itself and
$\{1,2,3,4\}$. The canonical declaration maps all of $\MCK$ to the single
letter $\mc$ and is unaffected; what the counterexample rules out is testing
the other two clauses of \eqref{eq:canonicalcg} in place of the descent
property itself
(Requirement~\ref{req:descent-decided-exactly}).
\end{remark}
\claimmeta{\metalabel{Uses}\mref{Lem.}{lem:descent}
\metasep \metalabel{Used by}\mref{Req.}{req:descent-decided-exactly}
\metasep newly derived
\metasep accepted}

% claim reductions/wpfss/def:weight-route
\claimsummary{reductions/wpfss/def:weight-route}{The receiver's blocks pass through the descended dark-count channel, unless supplied already composed, and join the private setting. Each joint outcome at a bounded letter gets its class weight times its setting's certified weight, sharpened by the dark counts and rounded down onto the lattice.}
\claimbookmark{reductions/wpfss/def:weight-route}{Definition}{def:weight-route}{The weight route}
\claimleadin{%
One granularity fact shapes this definition. The supplier emits the
\emph{coarse-grained} blocks $\blk{M}{=m}_{y|t}$, one per emitted letter
$y\in\Yset$, together with a \exttermref{../receiver/contract.pdf}{term-dark-count-channel}{dark-count channel} $R$ that acts on
\emph{raw} click patterns $2^{[K]}$ and weights $p^{(t)}_i$, one per setting,
certified for \emph{raw} outcome sets $\mathcal{O}_i\subseteq2^{[K]}$. A
consumer holding only the coarse blocks cannot in general apply $R$ to them,
because $S\mapsto\cg(S)$ is many-to-one and $R$ need not respect it. This
definition therefore requires $R$ to \emph{descend} through $\cg$
(Lemma~\ref{lem:descent}; decided exactly,
Requirement~\ref{req:descent-decided-exactly}). The case $\cg=\mathrm{id}$, in
which the emitted blocks are the raw ones, satisfies the condition
trivially. Remarks~\ref{rem:absorbneeded} and~\ref{rem:fibre} show that
neither the absorbing condition nor the fibre condition below can be dropped
(without the absorbing condition the bound \eqref{eq:certroute} for the
weights of \eqref{eq:weightroute} fails in general, in the \emph{unsafe}
direction), and Remark~\ref{rem:minweight} shows why the weight at $(y_i,t)$
is formed from the weight of its own setting, and why a slot carrying one
weight at every setting, as every one-setting slot does, gives the numbers of
the uniform route, which consumes each entry's minimum over the settings.
The order in which steps~1--2 precede the first squashing link is that of
Remark~\ref{rem:intertwining} and Requirement~\ref{req:post-processing-order};
the column-stochasticity of $R^{\cg}$, the fibre condition and the absorbing
condition are decided by Requirements~\ref{req:column-stochastic-refusal},
\ref{req:fibre-equality} and~\ref{req:absorbing-bounded-letters}; and the
downward rounding of the routed weights is that of
\xref{foundations/req:rounding-direction}, which rounds every emitted lower
bound down. A declaring party may compose the dark counts into the
receiver's model before it snaps the blocks, so that the snap is made on a
model whose elements are positive definite where the receiver module's
conditions for that hold (loss in the receiver and a dark-count probability
strictly between zero and one in every port,
\xref{receiver/detector-imperfections/prop:noise-floor}); it then supplies
the blocks already composed, and possibly snapped. The dark-composed form
below admits them, and Proposition~\ref{prop:precomposed-blocks} shows that
it forms the same instance as composing afterwards wherever no snap
intervenes.
}
\begin{definition}[The weight route]\label{def:weight-route}
Let a receiver artifact supply, for each setting $t\in\Tset$, preserved
photon-number blocks, $m\le N$, in one of the two input forms below, for a
conditional POVM $\{M_{y|t}\}_{y\in\Yset}$ on $\Chan$ block diagonal at the
cutoff; a
dark-count matrix $R$ on $2^{[K]}$ that descends through the
\exttermref{../receiver/contract.pdf}{term-declared-coarse-graining}{declared coarse-graining} $\cg$ with descended matrix $R^{\cg}$
(Lemma~\ref{lem:descent}); \termref{term-class-weights}{class weights} $\{q_t\}_{t\in\Tset}$ as in
Definition~\ref{def:joint-family}; and a slot
$\{(\mathcal{O}_i,(p^{(t)}_i)_{t\in\Tset})\}_{i=1}^{r}$ of certified
high-photon weights, one per entry and per setting of the same $\Tset$;
every declared entry and scalar lies on the declared lattice
\eqref{eq:latinv}, and every one of the supplier's facts about them is held
under Assumption~\ref{ass:recv} (Proposition~\ref{prop:route}\,(iv)). The
\termdef{term-weight-route}{weight route}
is the following three steps, in this order, each with the claim that
licenses it. Steps~1--2 are classical post-processing of the physical outcome
and precede the first squashing link.
\begin{enumerate}[leftmargin=2.2em]
\item \emph{Dark counts, on the emitted alphabet.} Replace the supplied
  conditional POVM by
  \begin{equation}\label{eq:darkstep}
    \widetilde M_{y|t}\;\coloneqq\;\sum_{y'\in\Yset}R^{\cg}_{y|y'}\,M_{y'|t},
    \qquad y\in\Yset,\ t\in\Tset .
  \end{equation}
  This is post-processing by a column-stochastic matrix
  (Lemma~\ref{lem:linear}\,(i)), so the result is a conditional POVM at each
  $t$ with the same photon-number block structure, and
  Lemma~\ref{lem:descent} shows it equals the coarse-graining of the raw-level
  composition.
\item \emph{\termref{term-joint-alphabet}{Joint alphabet}.} Form the
  \termref{term-joint-family}{joint family} of Definition~\ref{def:joint-family}\,(b),
  $M_{(y,t)}\coloneqq q_t\,\widetilde M_{y|t}$ on the joint alphabet
  $\Vset=\Yset\times\Tset$;
  this is a POVM by \eqref{eq:jointfamily} (Lemma~\ref{lem:linear}\,(ii);
  the weights are the measuring party's own, Assumption~\ref{ass:probe}), block diagonal at
  the cutoff, and the \termref{term-class-marginal-identity}{class-marginal identity} \eqref{eq:classmarginal}
  holds of it block by block. Its preserved blocks $\blk{M}{=m}_{(y,t)}$,
  $m\le N$, are the
  \termdef{term-post-processed-datum}{post-processed datum}.
\item \emph{The instance.} The post-processed joint family, with the routed
  weights below, is the datum, the hypotheses and the target data of the
  \emph{flat instance} (Proposition~\ref{prop:flat-instance}), and its
  emission is that instance's.
\end{enumerate}

\emph{The two input forms.} The supplied preserved blocks are declared in one
of two forms. In the
\termdef{term-bare-form}{bare form}
they are the blocks $\blk{M}{=m}_{y|t}$, $m\le N$, of the conditional POVM
above, and step~1 acts on them. In the
\termdef{term-dark-composed-form}{dark-composed form}
they are a declared family $\{\blk{\widetilde M}{=m}_{y|t}\}_{y\in\Yset}$,
$m\le N$, per setting $t$, each block Hermitian, positive semidefinite and
complete on its block, which the model takes below the cutoff after the
dark counts: in this form the family $\widetilde M_{y|t}$ of step~1 is
defined as these blocks below the cutoff and by \eqref{eq:darkstep} above
it, step~1 is not applied to the supplied blocks, and the post-processed
datum is formed from them by step~2 alone. Where no snap intervened between
composing and supplying, this form is the bare form composed in advance
(Proposition~\ref{prop:precomposed-blocks}\,(ii)); where the declaring party
snapped the composed blocks, they are its snapped family, and this
definition asserts of them no relation to a bare family. In either form $R^{\cg}$ is
formed and its conditions are decided as below, \eqref{eq:darkstep} defines
the family above the cutoff, $\blk{\widetilde M}{\gtN}_{y|t}
=\sum_{y'}R^{\cg}_{y|y'}\blk{M}{\gtN}_{y'|t}$, and the routed weights below
read $R^{\cg}$ as they stand. In either form $\{\widetilde M_{y|t}\}_y$ is a
POVM at each $t$, block diagonal at the cutoff: below it by step~1 in the
bare form and by the declared properties of the supplied blocks in the
dark-composed form, and above it by step~1.

\emph{Exactness of steps~1 and~2.} The two steps are written on the declared
numbers, and their POVM conclusions hold \emph{exactly} of them. The declared
$R^{\cg}$ is exactly column-stochastic with non-negative entries, decided on
entry, and $\{q_t\}$ is
exactly a probability distribution, as Definition~\ref{def:joint-family}
requires of \termref{term-class-weights}{class weights}; and the combination is formed in the exact
substrate of Lemma~\ref{lem:exactpsd} (Definition~\ref{def:substrate}), so
$\{\widetilde M_{y|t}\}_y$ and $\{M_{(y,t)}\}$ are POVMs and every claim of
this submodule reads them as such. What the two steps can nevertheless do is
leave the lattice: a product of two dyadic rationals adds significand widths,
so the post-processed blocks $\blk{M}{=m}_{(y,t)}$ need not be binary64
numbers even though every factor is. That is the one place in this submodule
where exact arithmetic outruns the representation; it is where the emission
normaliser of Definition~\ref{def:emitnorm} acts, and it acts on numbers this
module computed and is about to declare, never on numbers it was handed.

\emph{The routed weights.} They follow the same route, one bounded letter at a
time. Fix an entry $i$ whose bounded set is exactly a fibre,
$\mathcal{O}_i=\cg^{-1}(y_i)$, and whose letter is absorbing,
$R^{\cg}_{y_i|y_i}=1$; both are declared input validity conditions, and both
are decided exactly. Write
\begin{equation}\label{eq:weightroute}
\begin{gathered}
  \hat p_{(y_i,t)}\ \coloneqq\
  q_t\Big(\big(1-r_{\min,i}\big)p^{(t)}_i+r_{\min,i}\Big),
  \qquad
  p_{(y_i,t)}\ \coloneqq\ \lat{\hat p_{(y_i,t)}},\\
  r_{\min,i}\coloneqq\min_{y'\ne y_i}R^{\cg}_{y_i|y'}
  \ \ \big(\coloneqq0\ \text{when that minimum is over the empty set}\big),
\end{gathered}
\end{equation}
with $p_{(y,t)}\coloneqq0$ for every other $(y,t)$; the family
$(p_v)_{v\in\Vset}$ is the
\termdef{term-routed-weights}{routed weights},
and $p\coloneqq\sum_{v\in\Vset}p_v$. Here
\begin{equation}\label{eq:latdown}
  \lat{x}\;\coloneqq\;2^{-\kappa}\big\lfloor2^{\kappa}x\big\rfloor
\end{equation}
is the largest integer multiple of $2^{-\kappa}$ that is at most $x$: the
\emph{downward} rounding onto the declared lattice, computed exactly on the
rational $\hat p_{(y_i,t)}$ and equal to it whenever it is already a lattice
number. The routed weights are \emph{certified} when the element
$M_{(y_i,t)}$ of step~2 satisfies, on every channel block above the cutoff,
\begin{equation}\label{eq:certroute}
  \blk{M}{\gtN}_{(y_i,t)}\ \succeq\ p_{(y_i,t)}\,\idon{}
  \qquad\text{for every }i\text{ and every }t\in\Tset ,
\end{equation}
which Proposition~\ref{prop:route} proves from the declared input validity
conditions together with the absorbing hypothesis those conditions impose,
and from nothing else.

\emph{The direction of the rounding.} The emitted weight is $p_{(y_i,t)}$ and
not $\hat p_{(y_i,t)}$, and the direction of that rounding is the whole of
what makes it safe: it is the one emitted number of this submodule that a
downstream inequality consumes as a \emph{bound}, the certified-weight
inequality of Definition~\ref{def:targets} is preserved when a certified
weight is lowered and, when one is raised, is no longer proved and can be
false, and $\lat{\cdot}$ only lowers (Proposition~\ref{prop:route};
Proposition~\ref{prop:weight-monotone}\,(i)). Four features of the route are relied on, and none of them may
be dropped.
\emph{(a)} $r_{\min}$ is indexed by the bounded letter: it is a property of row
$y_i$ of $R^{\cg}$ read across the columns $y'\ne y_i$, and one entry's value is
not another's.
\emph{(b)} The absorbing condition is a hypothesis, not a consequence:
without it the bound \eqref{eq:certroute} for the weights of
\eqref{eq:weightroute} is no longer proved.
\emph{(c)} A non-zero $r_{\min,i}$ can occur only at $r=1$
(Corollary~\ref{cor:oneabsorbing}), so for $r\ge2$ \emph{every} entry
contributes the un-improved weight $\lat{q_tp^{(t)}_i}$, which is sound,
merely not an improvement.
\emph{(d)} $p_{(y_i,t)}\ge0$ needs no clamp: $\hat p_{(y_i,t)}$ is $q_t>0$
times $(1-r_{\min,i})p^{(t)}_i+r_{\min,i}$, a convex combination of
$p^{(t)}_i\ge0$ and $1$, since $r_{\min,i}\in[0,1]$ as an entry of the
column-stochastic $R^{\cg}$ or as $0$; and the largest lattice number at most
a non-negative rational is itself non-negative.
\end{definition}
\claimmeta{\metalabel{Uses}\mref{Lem.}{lem:linear}, \mref{Def.}{def:joint-family}, \mref{Def.}{def:exactadm}, \mref{Def.}{def:substrate}, \mref{Lem.}{lem:exactpsd}, \mref{Lem.}{lem:descent}
\metasep \metalabel{Used by}\mref{Def.}{def:targets}, \mref{Prop.}{prop:route}, \mref{Prop.}{prop:precomposed-blocks}, \mref{Cor.}{cor:oneabsorbing}, \mref{Rem.}{rem:colstoch}, \mref{Rem.}{rem:minweight}, \mref{Rem.}{rem:darkcompose}, \mref{Rem.}{rem:receiver-imports}, \mref{Prop.}{prop:flat-instance}, \mref{Cor.}{cor:emitted-flat}, \mref{Cor.}{cor:flat-companion}, \mref{Cor.}{cor:normalized-flat}, \mref{Req.}{req:supplier-version}
\metasep newly derived
\metasep accepted}

% claim reductions/wpfss/prop:route
\claimsummary{reductions/wpfss/prop:route}{At absorbing bounded letters the receiver's per-setting high-photon bound survives the dark counts, sharpened by r\_min, on every block of the joint family above the cutoff. Rounding it down onto the lattice keeps it certified, so the routed weights are certified weights for the flag-state targets.}
\claimbookmark{reductions/wpfss/prop:route}{Proposition}{prop:route}{Certified weights routed from the receiver artifact}
\claimleadin{%
The routed weight of Definition~\ref{def:weight-route} is certified for two
reasons: the receiver's per-setting bound survives the dark counts at an
absorbing letter, and the emitted weight is rounded down. Had that rounding
been to nearest, or upwards, the second step would fail: an emitted weight
above $\hat p_{(y_i,t)}$ leaves \eqref{eq:certweights} unproved, and it can
make it false, the failure Assumption~\ref{ass:recv}'s ledger row calls a
security break.
\xref{foundations/req:rounding-direction} names the same direction for every
emitted lower bound. The input conditions part~(iv) leaves to the declaration
are decided by Requirements~\ref{req:column-stochastic-refusal},
\ref{req:descent-decided-exactly}, \ref{req:fibre-equality}
and~\ref{req:absorbing-bounded-letters}.
}
\begin{proposition}[Certified weights routed from the receiver artifact]\label{prop:route}
Let $\{M_{y|t}\}_{y\in\Yset}$ be, for each $t\in\Tset$, a POVM on $\Chan$ block
diagonal at the cutoff, and let $R$ descend through $\cg$ with descended matrix
$R^{\cg}$ (Lemma~\ref{lem:descent}). Let $\{q_t\}$ be a probability
distribution, and let $\{(\mathcal{O}_i,(p^{(t)}_i)_{t\in\Tset})\}_{i=1}^{r}$
be the declared slot, with $\mathcal{O}_i=\cg^{-1}(y_i)$ for pairwise distinct
letters $y_1,\dots,y_r$, with $0\le p^{(t)}_i\le1$, and with the per-setting
guarantee $\blk{M}{\gtN}_{y_i|t}\succeq p^{(t)}_i\,\idon{}$ for every
$t\in\Tset$ and every $i$. Suppose each bounded letter is \emph{absorbing},
\begin{equation}\label{eq:absorbing}
  R^{\cg}_{y_i|y_i}=1
  \qquad\text{for every }i=1,\dots,r,
  \tag{absorbing}
\end{equation}
and put $r_{\min,i}\coloneqq\min_{y'\ne y_i}R^{\cg}_{y_i|y'}$, with the
convention $r_{\min,i}\coloneqq0$ when the minimum is over the empty set (that
is, when $\Yset=\{y_i\}$), the conservative value, which reduces
\eqref{eq:routebound} to the un-improved bound. Then:
\begin{enumerate}[label=(\roman*),leftmargin=2.4em]
\item The joint-alphabet POVM
  $M_{(y,t)}\coloneqq q_t\sum_{y'}R^{\cg}_{y|y'}M_{y'|t}$ of
  Definitions~\ref{def:joint-family} and~\ref{def:weight-route} satisfies, for
  every $i$ and every $t\in\Tset$,
  \begin{equation}\label{eq:routebound}
    \blk{M}{\gtN}_{(y_i,t)}\ \succeq\
    q_t\Big(\big(1-r_{\min,i}\big)\,p^{(t)}_i+r_{\min,i}\Big)\,\idon{}
    \;=\;\hat p_{(y_i,t)}\,\idon{} .
  \end{equation}
  A slot whose weights do not depend on the setting, $p^{(t)}_i=p_i$ for
  every $t$ (the $t$-uniform form, which every one-setting slot has), is the
  special case, and there \eqref{eq:routebound} is the bound in $p_i$.
\item The \emph{emitted} weight of \eqref{eq:weightroute} is certified, because
  the rounding runs downwards. Every number on the right of
  \eqref{eq:routebound} lies on the declared lattice: $q_t$ and $p^{(t)}_i$ by
  declaration, and $r_{\min,i}$ as an entry of $R^{\cg}$, which is a sum of
  declared entries and lies in $[0,1]$. So $\hat p_{(y_i,t)}$ is an exact
  element of $\Fld$ formed from them; and
  $p_{(y_i,t)}=\lat{\hat p_{(y_i,t)}}\le\hat p_{(y_i,t)}$ by
  \eqref{eq:latdown}, whence
  \begin{equation*}
    \blk{M}{\gtN}_{(y_i,t)}\ \succeq\ p_{(y_i,t)}\,\idon{} ,
  \end{equation*}
  which is \eqref{eq:certroute}; so the emitted $(p_v)$ are
  \termref{term-certified-weights}{certified weights} in the sense of Definition~\ref{def:targets} and they lie on the declared
  lattice.
\item Under the sufficient condition \eqref{eq:canonicalcg} of
  Lemma~\ref{lem:descent} and with the single entry $\mathcal{O}_1=\MCK$,
  \eqref{eq:absorbing} holds at $y_1=y_{\mc}$, and its per-setting hypothesis
  is the supplier's slot guarantee
  $\sum_{S\in\MCK}\blk{M}{\gtN}_{S|t}\succeq p^{(t)}_1\,\idon{}$ at each $t$,
  read through $\cg^{-1}(y_{\mc})=\MCK$.
\item On the emission path the hypotheses about the supplier's objects hold
  under Assumption~\ref{ass:recv}: the conditional families, and their
  block diagonality in the photon number, block by block and so in
  particular at the cutoff, are at every setting those of the coarse-grained
  elements of the \exttermref{../receiver/contract.pdf}{term-setting-family}{setting family} on which the supplier states its
  weights, by item~(a) of \xref{receiver/def:setting-family} and
  \xref{receiver/thm:coarse-graining}; for a receiver without probe inputs
  that family is the one-setting family of
  \xref{receiver/cor:unconditioned-family}, and the conditional families and
  their block diagonality are given by the receiver's
  \xref{receiver/thm:click-povm}\,(iii) and
  \xref{receiver/thm:coarse-graining}; the
  dark-count matrix $R$, and the block diagonality in the photon number of
  the elements with the dark counts composed, by
  \xref{receiver/detector-imperfections/def:darkadm} and
  \xref{receiver/detector-imperfections/thm:dark-counts}; so every element of
  the joint family of (i), and of the post-processed joint family that
  Definition~\ref{def:weight-route} forms in either of its forms, is block
  diagonal in the photon number in every block, below the cutoff as above
  it; and the slot, its
  per-setting guarantee (the first conclusion of
  \xref{receiver/thm:emitted-bound}) and $0\le p^{(t)}_i\le1$,
  $\sum_ip^{(t)}_i\le1$ at every $t$ by \xref{receiver/def:slot},
  \xref{receiver/thm:emitted-bound} and \xref{receiver/prop:slotvalid}. The
  other hypotheses are not supplier facts: that $R$ is column-stochastic and descends through $\cg$, that
  $\mathcal{O}_i=\cg^{-1}(y_i)$ and \eqref{eq:absorbing} are decided exactly
  on the declared data, a declaration failing one being
  refused, and $\{q_t\}$ are the declared \termref{term-class-weights}{class weights}
  (Definition~\ref{def:joint-family}).
\end{enumerate}
\end{proposition}
\claimmeta{\metalabel{Uses}\mref{Def.}{def:joint-family}, \mref{Def.}{def:substrate}, \mref{Lem.}{lem:exactpsd}, \mref{Def.}{def:targets}, \mref{Prop.}{prop:weight-monotone}, \mref{Lem.}{lem:descent}, \mref{Def.}{def:weight-route}, \extref{../receiver/contract.pdf}{receiver/cor:unconditioned-family}, \extref{../receiver/contract.pdf}{receiver/def:setting-family}, \extref{../receiver/contract.pdf}{receiver/def:slot}, \extref{../receiver/contract.pdf}{receiver/detector-imperfections/def:darkadm}, \extref{../receiver/contract.pdf}{receiver/detector-imperfections/thm:dark-counts}, \extref{../receiver/contract.pdf}{receiver/prop:slotvalid}, \extref{../receiver/contract.pdf}{receiver/thm:click-povm}, \extref{../receiver/contract.pdf}{receiver/thm:coarse-graining}, \extref{../receiver/contract.pdf}{receiver/thm:emitted-bound}
\metasep \metalabel{Assumes}\aref{A-RECV}{ass:recv}
\metasep \metalabel{Used by}\mref{Cor.}{cor:oneabsorbing}, \mref{Rem.}{rem:absorbneeded}, \mref{Rem.}{rem:fibre}, \mref{Rem.}{rem:colstoch}, \mref{Rem.}{rem:minweight}, \mref{Rem.}{rem:darkcompose}, \mref{Rem.}{rem:receiver-imports}, \mref{Prop.}{prop:flat-instance}, \mref{Cor.}{cor:emitted-flat}, \mref{Req.}{req:absorbing-bounded-letters}, \mref{Req.}{req:supplier-version}
\metasep newly derived
\metasep accepted
\metasep \metalabel{Oracles}\coderef{https://github.com/fraunhofersg/certiqs-key/blob/v0.1.0/tests/oracles/registry_reductions.py}{registry_reductions.py}}

\begin{proof}
(i) Fix $i$ and $t$. Split the sum over columns and use \eqref{eq:absorbing}
at that letter:
\begin{align*}
  \sum_{y'}R^{\cg}_{y_i|y'}M_{y'|t}
  &=M_{y_i|t}+\sum_{y'\ne y_i}R^{\cg}_{y_i|y'}M_{y'|t}
  \;\succeq\;M_{y_i|t}+r_{\min,i}\sum_{y'\ne y_i}M_{y'|t}\\
  &=(1-r_{\min,i})\,M_{y_i|t}+r_{\min,i}\,\idon{},
\end{align*}
using $M_{y'|t}\succeq0$ and $R^{\cg}_{y_i|y'}\ge r_{\min,i}$ in the
inequality and POVM completeness in the last step. All operators are block
diagonal at the cutoff, so the inequality restricts to the $\gtN$ block, where
$\blk{M}{\gtN}_{y_i|t}\succeq p^{(t)}_i\,\idon{}$ and $1-r_{\min,i}\ge0$, every
entry of $R^{\cg}$ lying in $[0,1]$ by non-negativity and unit columns
(Lemma~\ref{lem:descent}). Multiplying by $q_t\ge0$ gives
\eqref{eq:routebound}. Nothing in the argument couples the entries: each $i$
uses only its own row of $R^{\cg}$ and its own weight, which is why $r_{\min}$
must be indexed (Corollary~\ref{cor:oneabsorbing} says what that costs).

(ii) $\hat p_{(y_i,t)}$ is formed from $q_t$, $p^{(t)}_i$ and $r_{\min,i}$ by
finitely many field operations in $\Fld$ (Lemma~\ref{lem:exactpsd},
Definition~\ref{def:substrate}), so it \emph{is} the real value of
\eqref{eq:weightroute}'s expression and \eqref{eq:routebound} is a statement
about that value. Monotonicity of $\succeq$ in a scalar multiple of the
identity, with $p_{(y_i,t)}\le\hat p_{(y_i,t)}$, gives \eqref{eq:certroute};
and $p_{(y_i,t)}\ge0$ because $\hat p_{(y_i,t)}\ge0$ and $\lat{\cdot}$ carries
a non-negative rational to a non-negative lattice value. No rounding-error
budget is tracked and none is needed: the inequality is exact, and the only
inexactness anywhere near it is a rounding that weakens it
(Proposition~\ref{prop:weight-monotone}\,(i)).

(iii) The proof of Lemma~\ref{lem:descent} under \eqref{eq:canonicalcg} gives
$P(y_{\mc}|S')=1$ for every $S'\in\MCK$, hence
$R^{\cg}_{y_{\mc}|y_{\mc}}=1$; and
$M_{y_{\mc}|t}=\sum_{S\in\cg^{-1}(y_{\mc})}M_{S|t}=\sum_{S\in\MCK}M_{S|t}$ by
the definition of the coarse-grained POVM.

(iv) Each receiver claim cited there states the hypothesis it is cited for,
of the objects the receiver artifact supplies, and
Assumption~\ref{ass:recv} makes the artifact's facts hold of the numbers as
declared. The supplier states its weights on a setting family
(\xref{receiver/thm:emitted-bound}), whose elements at every setting commute
with the photon number by item~(a) of \xref{receiver/def:setting-family},
and a coarse-grained element is a sum of them
(\xref{receiver/thm:coarse-graining}); for a receiver without probe inputs
\xref{receiver/cor:unconditioned-family} identifies that family as the
click-pattern POVM of \xref{receiver/thm:click-povm}, block diagonal by its
item~(iii). The remaining hypotheses are the input decisions and the
declaration named there. An element of the joint family of (i), and an
element of the post-processed joint family in either form, is a combination
with non-negative coefficients ($q_t$ and the entries of $R^{\cg}$, or $q_t$
alone) of the receiver's elements, conditional or with the dark counts
composed, which the claims cited give block diagonal in the photon number;
and a sum of operators each block diagonal in the photon number is again
block diagonal in it, block by block.
\end{proof}

% claim reductions/wpfss/prop:precomposed-blocks
\claimsummary{reductions/wpfss/prop:precomposed-blocks}{Supplying the preserved blocks already composed with the dark counts forms the same flat instance, and so the same emission, as composing them here, wherever no snap moves a block. A snap made before or after the composition changes only the data sectors, and the composed reference is a reference.}
\claimbookmark{reductions/wpfss/prop:precomposed-blocks}{Proposition}{prop:precomposed-blocks}{Preserved blocks supplied with the dark counts already composed}
\claimleadin{%
The dark counts are a trusted classical channel, and composing them before or
after the supplier's blocks are handed over is the same operation; what the
order changes is only where a snap, if one is needed, is made. A declaring
party snaps the composed model when it wants the noise floor of the dark
counts in place before its snap. This proposition says that the flat instance
does not notice the order (part (ii)) and what changes when a snap
intervenes (part (iii)), the composed reference being what the receiver
hands beside its snap of the composed blocks.
In the \termref{term-dark-composed-form}{dark-composed form} the dark-count channel is not applied to the
preserved blocks again, but it is still read: it defines the family above the
cutoff, and the routed weights read its minima and its absorbing letters, as
in the \termref{term-bare-form}{bare form}.
}
\begin{proposition}[Preserved blocks supplied with the dark counts already composed]\label{prop:precomposed-blocks}
Let the data of Definition~\ref{def:weight-route} be given in the
\termref{term-bare-form}{bare form}, with supplied preserved blocks $B^{(m)}_{y|t}$, $m\le N$,
and put
\begin{equation}\label{eq:precomposed}
  \widetilde B^{(m)}_{y|t}\;\coloneqq\;\sum_{y'\in\Yset}R^{\cg}_{y|y'}\,B^{(m)}_{y'|t}.
\end{equation}
\begin{enumerate}[label=(\roman*),leftmargin=3.0em]
\item For every $t$ and $m$, $\{\widetilde B^{(m)}_{y|t}\}_{y\in\Yset}$ is a
  family of Hermitian positive semidefinite matrices complete on its block.
\item The flat instance of Proposition~\ref{prop:flat-instance} formed from
  the blocks $\widetilde B$ in the \termref{term-dark-composed-form}{dark-composed form}, and the one
  formed from the blocks $B$ in the bare form, have the same
  \termref{term-handed-datum}{handed datum}, $A^{(m)}_{(y,t)}=q_t\widetilde B^{(m)}_{y|t}$, the same
  \termref{term-hypotheses-h}{hypotheses (H)}, \termref{term-model-map}{model map}, \termref{term-routed-weights}{routed weights},
  \termref{term-scalar-value}{scalar values}, channel, \termref{term-transported-set}{transported set} and
  entries; so they are one \termref{term-reduction-instance}{reduction instance}, and the
  \termref{term-emission-normalizer}{emission normalizer} gives them the same \termref{term-declared-datum}{declared datum},
  \termref{term-declared-model}{declared model} and record.
\item Let $\{G^{(m)}_{y|t}\}_{y\in\Yset}$ be a \exttermref{../foundations/contract.pdf}{term-reference}{reference} of the bare
  blocks of setting $t$ with constant $1$, read through the
  \exttermref{../foundations/contract.pdf}{term-entrywise-enclosure}{entrywise enclosures}
  $\big(C^{(m)}_{y|t},r^{(m)}_{y|t}\big)$
  (\xref{foundations/def:exact-reference}), and put
  $\widetilde G^{(m)}_{y|t}\coloneqq\sum_{y'}R^{\cg}_{y|y'}G^{(m)}_{y'|t}$.
  Then $\widetilde G$ is a reference with constant $1$, read through the
  enclosures
  $\big(\sum_{y'}R^{\cg}_{y|y'}C^{(m)}_{y'|t},\,\sum_{y'}R^{\cg}_{y|y'}r^{(m)}_{y'|t}\big)$,
  and it is an \exttermref{../foundations/contract.pdf}{term-exact-reference}{exact reference} where $G$ is. If $B$ is the
  snapped family of a \exttermref{../foundations/contract.pdf}{term-snap}{snap} of $G$ and $B'$ that of a snap of
  $\widetilde G$ (\xref{foundations/def:snap}), then the handed data
  $A^{\mathrm{b}}_{(y,t)}=q_t\sum_{y'}R^{\cg}_{y|y'}B_{y'|t}$ of the bare
  form and $A^{\mathrm{d}}_{(y,t)}=q_tB'_{y|t}$ of the dark-composed form
  differ in general. The two flat instances agree in everything but their
  data, so their emissions differ only in their
  \termref{term-data-sector}{data sectors}.
\end{enumerate}
\end{proposition}
\claimmeta{\metalabel{Uses}\mref{Def.}{def:reduction-instance}, \mref{Def.}{def:emitnorm}, \mref{Def.}{def:targets}, \mref{Prop.}{prop:chain}, \mref{Lem.}{lem:descent}, \mref{Def.}{def:weight-route}, \mref{Prop.}{prop:flat-instance}, \extref{../foundations/contract.pdf}{foundations/def:exact-reference}, \extref{../foundations/contract.pdf}{foundations/def:snap}
\metasep newly derived
\metasep implemented (effective: accepted)
\metasep \metalabel{Code}\coderef{https://github.com/fraunhofersg/certiqs-key/blob/v0.1.0/src/certiqs_key/reductions/wpfss/weights.py}{certiqs_key.reductions.wpfss.weights:post_processed_preserved_blocks}}

\begin{proof}
(i) Each $\widetilde B^{(m)}_{y|t}$ is a combination of Hermitian positive
semidefinite matrices with the non-negative coefficients $R^{\cg}_{y|y'}$,
and $\sum_y\widetilde B^{(m)}_{y|t}=\sum_{y'}\big(\sum_yR^{\cg}_{y|y'}\big)
B^{(m)}_{y'|t}=\sum_{y'}B^{(m)}_{y'|t}=\idon{}$ by the column-stochasticity
of $R^{\cg}$ (Lemma~\ref{lem:descent}) and the completeness of the supplied
family on its block.

(ii) In the bare form the handed datum is formed by steps~1--2 of
Definition~\ref{def:weight-route},
$A^{(m)}_{(y,t)}=q_t\sum_{y'}R^{\cg}_{y|y'}B^{(m)}_{y'|t}$
(Proposition~\ref{prop:flat-instance}\,(I1)), which is
$q_t\widetilde B^{(m)}_{y|t}$ by \eqref{eq:precomposed}; in the dark-composed
form it is formed by step~2 alone from the supplied $\widetilde B$, which
gives the same matrices, entry for entry, in exact arithmetic. Every other
part of the flat instance reads no supplied block below the cutoff: the
family above the cutoff is
$q_t\sum_{y'}R^{\cg}_{y|y'}\blk{M}{\gtN}_{y'|t}$ in either form
(Definition~\ref{def:weight-route}); the hypotheses (H1)--(H3) are about that
family and the routed weights; the routed weights read $R^{\cg}$, the class
weights and the declared slot; the scalar values
\eqref{eq:flat-scalars} read the routed weights; the channels \eqref{eq:Lwp}
and \eqref{eq:Lfss} read the cutoff projectors, the family above the cutoff
and the weights; and the transported set and the entries are built from
these and from the sector projectors (Proposition~\ref{prop:flat-instance}\,(I3)
and (I4)). So the two instances coincide. The normalizer of
Definition~\ref{def:emitnorm} is a function of the handed datum, the
declared order, the lattice exponent and the run option, so it produces the
same declared datum and records, and the declared model, the model map at
that datum, is the same.

(iii) $\widetilde G$ is the image of $G$ under the same non-negative
combination as in (i), so its blocks are Hermitian, positive semidefinite and
complete at $1$. Its centres $\sum_{y'}R^{\cg}_{y|y'}C^{(m)}_{y'|t}$ are
Hermitian with Gaussian-rational entries and its radii are non-negative
rationals, the entries of $R^{\cg}$ being non-negative dyadic rationals; and
for every entry $(i,l)$,
\begin{equation*}
  \Big|\big(\widetilde G^{(m)}_{y|t}\big)_{il}
    -\Big(\sum_{y'}R^{\cg}_{y|y'}C^{(m)}_{y'|t}\Big)_{il}\Big|
  \;\le\;\sum_{y'}R^{\cg}_{y|y'}\Big|\big(G^{(m)}_{y'|t}-C^{(m)}_{y'|t}\big)_{il}\Big|
  \;\le\;\sum_{y'}R^{\cg}_{y|y'}\big(r^{(m)}_{y'|t}\big)_{il},
\end{equation*}
by the triangle inequality and $R^{\cg}\ge0$. That is
\xref{foundations/def:exact-reference}, and every radius is zero where every
radius of $G$ is. That the two handed data differ in general
is that rounding and composing do not commute: a snap of $G$ followed by
$R^{\cg}$ and a snap of $R^{\cg}G$ are different lattice families. The last
sentence is (ii)'s list of what the instance reads besides its handed datum.
\end{proof}

% claim reductions/wpfss/cor:oneabsorbing
\claimsummary{reductions/wpfss/cor:oneabsorbing}{With two or more absorbing bounded letters every r\_min,i vanishes, so each routed weight is the un-improved lat(q\_t p\_i\textasciicircum{}(t)). The dark counts sharpen a weight only when the slot has a single entry.}
\claimbookmark{reductions/wpfss/cor:oneabsorbing}{Corollary}{cor:oneabsorbing}{Dark-count sharpening at a single bounded letter only}
\begin{corollary}[Dark-count sharpening at a single bounded letter only]
\label{cor:oneabsorbing}
Under \eqref{eq:absorbing}: if $r\ge2$ then $r_{\min,i}=0$ for \emph{every}
$i$, so \eqref{eq:routebound} reads
$\blk{M}{\gtN}_{(y_i,t)}\succeq q_tp^{(t)}_i\,\idon{}$ throughout and the
emitted weight is $\lat{q_tp^{(t)}_i}$. Consequently a positive
$r_{\min,i}$, the only case in which the dark counts sharpen a weight, can
occur for at most one bounded letter, and only when $r=1$.
\end{corollary}
\claimmeta{\metalabel{Uses}\mref{Lem.}{lem:exactpsd}, \mref{Def.}{def:weight-route}, \mref{Prop.}{prop:route}
\metasep \metalabel{Used by}\mref{Rem.}{rem:colstoch}, \mref{Rem.}{rem:minweight}
\metasep newly derived
\metasep accepted}

\begin{proof}
Let $r\ge2$ and fix any $i$; choose $j\ne i$, which is possible. Column $y_j$
of $R^{\cg}$ is a probability vector and $R^{\cg}_{y_j|y_j}=1$ by
\eqref{eq:absorbing}, so every other entry of that column vanishes; in
particular $R^{\cg}_{y_i|y_j}=0$. Since $y_j\ne y_i$, the minimum defining
$r_{\min,i}$ ranges over a set containing $y_j$, so
$0\le r_{\min,i}\le R^{\cg}_{y_i|y_j}=0$. As $i$ was arbitrary, every
$r_{\min,i}$ vanishes; substituting into \eqref{eq:routebound} gives the
displayed form, which is the un-improved bound and is sound, and
\eqref{eq:weightroute} then reads $\hat p_{(y_i,t)}=q_tp^{(t)}_i$. The last
sentence follows: a positive $r_{\min,i}$ is incompatible with the existence
of a second absorbing bounded letter. Every step is an identity or a
comparison between declared values and is decided as such in the substrate of
Lemma~\ref{lem:exactpsd}, so the conclusion holds exactly.
\end{proof}

% claim reductions/wpfss/rem:absorbneeded
\claimsummary{reductions/wpfss/rem:absorbneeded}{Without the absorbing condition the routed bound is unproved and fails in general, unsafely: a single-click letter loses weight to multiclick patterns under a dark count at another port, and no other column restores it. The receiver's multiclick entry is absorbing, so its slot is admitted.}
\claimbookmark{reductions/wpfss/rem:absorbneeded}{Remark}{rem:absorbneeded}{Necessity of the absorbing hypothesis}
\begin{remark}[Necessity of the absorbing hypothesis]\label{rem:absorbneeded}
Without \eqref{eq:absorbing}, \eqref{eq:routebound} is no longer proved, and
it fails in general, in the direction that certifies a weight the device does
not have. Under the supplier's \exttermref{../receiver/contract.pdf}{term-independent-product-form}{independent product form}
(\xref{receiver/detector-imperfections/def:darkadm}), a bounded
\emph{single-click} letter $y_i=\cg(\{k\})$ whose fibre $\cg^{-1}(y_i)$ is
$\{\{k\}\}$, as it is when $\cg$ is injective off the \exttermref{../receiver/contract.pdf}{term-multiclick-set}{multiclick set}
and maps no multiclick pattern to $y_i$, has $R^{\cg}_{y_i|y_i}=R_{\{k\}|\{k\}}
=\prod_{k'\ne k}(1-p_{\dc|k'})<1$
whenever some $p_{\dc|k'}>0$, because a dark count at another port moves the
pattern $\{k\}$ into the multiclick set and never back. And, for any
monotone dark-count channel,
$r_{\min,i}=0$, because monotonicity kills \emph{at least one} other column:
any other single-click letter $y'=\cg(\{k''\})$, $k''\ne k$, has
$R^{\cg}_{y_i|y'}=0$ since $\{k''\}\not\subseteq\{k\}$, which is all a minimum
needs, and $K\ge2$ supplies such a letter. Formula \eqref{eq:routebound} would
then assert $\blk{\widetilde M}{\gtN}_{y_i|t}\succeq p^{(t)}_i\idon{}$, while
in fact the inequality that does hold is
$\widetilde M_{y_i|t}\succeq R^{\cg}_{y_i|y_i}M_{y_i|t}$ with
$R^{\cg}_{y_i|y_i}<1$, and the columns that remain do not compensate: under
the independent product form the only other column
feeding $y_i$ is the no-click one,
$R^{\cg}_{y_i|\cg(\varnothing)}=p_{\dc|k}\prod_{k'\ne k}(1-p_{\dc|k'})$, which
multiplies $M_{\cg(\varnothing)|t}$ and not $M_{y_i|t}$ and carries no
high-photon weight of its own. In the sub-case $p_{\dc|k}=0$ with some
$p_{\dc|k'}>0$ that column vanishes too and the loss is exactly the factor
$R^{\cg}_{y_i|y_i}=\prod_{k'\ne k}(1-p_{\dc|k'})<1$: for a receiver whose
certified weight is attained and positive,
$\lmin\big(\blk{M}{\gtN}_{y_i|t}\big)=p^{(t)}_i>0$, the smallest eigenvalue of
$\blk{\widetilde M}{\gtN}_{y_i|t}$ is $R^{\cg}_{y_i|y_i}\,p^{(t)}_i<p^{(t)}_i$,
and \eqref{eq:routebound} is false. A weight asserted above
the true one makes \eqref{eq:certweights} false, so $\Fwp$ is not a squashing
of the real receiver and the emitted licence certifies an inclusion that does
not hold, the failure Assumption~\ref{ass:recv}'s ledger row calls a security
break. Condition \eqref{eq:absorbing} is therefore a declared input validity
condition at every bounded letter
(Requirement~\ref{req:absorbing-bounded-letters}), and no bound is stated for
a letter that fails it. The multiclick letter of the canonical declaration is
absorbing (Proposition~\ref{prop:route}\,(iii)), which is why the slot the
receiver emits, the single entry $(\MCK,(p^{(t)}_{\mc})_{t\in\Tset})$ of its
\xref{receiver/def:slot}, written $(\MCK,p_{\mc})$ at a probe-free receiver,
is admitted.
\end{remark}
\claimmeta{\metalabel{Uses}\mref{Def.}{def:targets}, \mref{Prop.}{prop:route}, \extref{../receiver/contract.pdf}{receiver/def:slot}
\metasep \metalabel{Assumes}\aref{A-RECV}{ass:recv}
\metasep \metalabel{Used by}\mref{Req.}{req:absorbing-bounded-letters}
\metasep newly derived
\metasep accepted}

% claim reductions/wpfss/rem:fibre
\claimsummary{reductions/wpfss/rem:fibre}{A certified bound on a sum of positive operators says nothing about its summands, so a certified set split over two letters gives neither a bound, and one merged with other outcomes attaches the weight to a larger fibre. Each certified set must therefore be exactly one fibre of the coarse-graining.}
\claimbookmark{reductions/wpfss/rem:fibre}{Remark}{rem:fibre}{Necessity of the fibre condition}
\begin{remark}[Necessity of the fibre condition]\label{rem:fibre}
Proposition~\ref{prop:route} needs the certified set to be \emph{exactly} a
fibre, $\mathcal{O}=\cg^{-1}(y_{*})$. The reason is that the certified quantity
is the \emph{sum} $\sum_{S\in\mathcal{O}}M_{S|t}$, and a bound on a sum of
positive operators says nothing about the individual summands. Two failures
follow, in opposite directions. If $\cg$ \emph{splits} $\mathcal{O}$ across two
emitted letters, neither letter inherits a bound, and the natural repair,
bounding each piece, is unavailable because the supplier certified only the
sum. If $\cg$ merges $\mathcal{O}$ with outcomes \emph{outside} it, the bound
still holds for the merged letter (adding positive operators only helps), but
the coarse blocks the supplier emits no longer let this module recover
$\sum_{S\in\mathcal{O}}M_{S|t}$ separately, so the emitted weight is attached
to a letter whose fibre is larger than the certified set: sound, and
recorded as such, but only if the containment
$\mathcal{O}\subseteq\cg^{-1}(y_{*})$ is validated rather than assumed.
Equality is therefore required, and decided at input
(Requirement~\ref{req:fibre-equality}). Where the pieces carry key
information and must not be merged, the fact does not become encodable by
re-declaring $\cg$; it becomes an event-level bound, and
Proposition~\ref{prop:obstruction} shows that a target classical on the
whole of $\Qsys$ cannot in general encode it.
\end{remark}
\claimmeta{\metalabel{Uses}\mref{Prop.}{prop:route}
\metasep \metalabel{Used by}\mref{Req.}{req:fibre-equality}
\metasep newly derived
\metasep accepted}

% claim reductions/wpfss/rem:colstoch
\claimsummary{reductions/wpfss/rem:colstoch}{Exact column-stochasticity of the declared R gives that of R\textasciicircum{}cg by regrouping one finite sum, so one exact decision on R serves both; the route, the post-processing step and the one-letter corollary rely on it. Every result here is proved for an exactly column-stochastic R only.}
\claimbookmark{reductions/wpfss/rem:colstoch}{Remark}{rem:colstoch}{Column-stochasticity of R and of its descended matrix}
\begin{remark}[Column-stochasticity of R and of its descended matrix]\label{rem:colstoch}
\emph{One matrix suffices.} The declared $R$ is decided exactly
column-stochastic, $R_{S|S'}\ge0$ for every entry and $\sum_SR_{S|S'}=1$ for
every column (Requirement~\ref{req:column-stochastic-refusal}). The fibres of
$\cg$ partition $2^{[K]}$, $\cg$ being a total surjection onto $\Yset$
(Lemma~\ref{lem:descent}). So
$R^{\cg}_{y|y'}=\sum_{S\in\cg^{-1}(y)}R_{S|S'}$ is non-negative as a sum of
non-negative entries, and
$\sum_{y\in\Yset}R^{\cg}_{y|y'}=\sum_{S}R_{S|S'}=1$, which is exactly the
column-stochasticity Lemma~\ref{lem:descent} states of $R^{\cg}$, there
derived from the same hypothesis on $R$. In the exact substrate that inference
carries: both identities regroup one finite sum in $\Fld$
(Lemma~\ref{lem:exactpsd}, Definition~\ref{def:substrate}), so the decision on
$R$ settles $R^{\cg}$ as well and no second decision is needed. A regrouped
binary64 sum need not be the sum it regroups; under exact arithmetic the
question does not arise.

\emph{Where the property is used.} Three results rely on it.
\emph{(i)} Proposition~\ref{prop:route}, whose proof restricts
$(1-r_{\min,i})\blk{M}{\gtN}_{y_i|t}\succeq(1-r_{\min,i})p^{(t)}_i\idon{}$ and
needs $1-r_{\min,i}\ge0$, which is non-negativity of the entries together with
unit columns. \emph{(ii)} Lemma~\ref{lem:linear}\,(i) at step~1 of
Definition~\ref{def:weight-route}: post-processing by $R^{\cg}$ yields a POVM
because the coefficients are non-negative and each column sums to one; without
both, $\{\widetilde M_{y|t}\}$ need not be one. \emph{(iii)}
Corollary~\ref{cor:oneabsorbing}, whose proof reads a column of $R^{\cg}$ as a
probability vector.

\emph{What the decision closes: coverage and attribution, not soundness.} A
violation that moves an emitted number in the unsafe direction can also be
caught by the \termref{term-total-weight}{total weight}: with $\sum_tq_t=1$ (Definition~\ref{def:joint-family})
and the weights read as the values of \eqref{eq:weightroute}'s expression,
the routed total is $\sum_i\big(r_{\min,i}(1-\bar p_i)+\bar p_i\big)$ with
$\bar p_i\coloneqq\sum_tq_tp^{(t)}_i$ and $0\le\bar p_i\le1$ from the slot's
validity condition, which gives $0\le p^{(t)}_i\le1$ at every $t$ (at a slot
with one weight at every setting $\bar p_i$ is that weight), so any
$r_{\min,i}\ge1$ makes the $i$-th summand at least $1$, and the total at least
$1$ as long as no other summand is negative
(Requirement~\ref{req:total-weight-refusal}); a negative $r_{\min,j}$, a defect
of the same class, can make one negative. A defect outside that path,
one that neither lifts some $r_{\min,i}$ to $1$ nor shows up in the total,
would meet no other rule. That path is not relied on: the decision on the
declared matrix settles the two conditions directly and exactly, and the
absorbing condition fixes $R^{\cg}_{y_i|y_i}=1$
(Requirement~\ref{req:absorbing-bounded-letters}), so every admitted
$r_{\min,i}$ lies in $[0,1]$. Every result of this contract is proved for an
exactly column-stochastic $R$ and asserts nothing whatever about a
declaration that is not one. Beyond diagnosis, the decision is used in one
place: every admitted $R$ satisfies (D1) and (D2) of
\xref{receiver/detector-imperfections/def:darkadm} exactly, which is half of
Remark~\ref{rem:darkcompose}.
\end{remark}
\claimmeta{\metalabel{Uses}\mref{Lem.}{lem:linear}, \mref{Def.}{def:joint-family}, \mref{Lem.}{lem:descent}, \mref{Def.}{def:weight-route}, \mref{Prop.}{prop:route}, \mref{Cor.}{cor:oneabsorbing}
\metasep \metalabel{Used by}\mref{Req.}{req:column-stochastic-refusal}
\metasep newly derived
\metasep accepted}

% claim reductions/wpfss/rem:minweight
\claimsummary{reductions/wpfss/rem:minweight}{The route forms each joint outcome's weight from its own setting's certified weight; consuming each entry's minimum over the settings instead would be sound but weakly smaller everywhere. A slot with one weight at every setting, as at every probe-free receiver, gives the same numbers either way.}
\claimbookmark{reductions/wpfss/rem:minweight}{Remark}{rem:minweight}{Per-setting weights against the minimum over the settings}
\begin{remark}[Per-setting weights against the minimum over the settings]
\label{rem:minweight}
\emph{What is consumed.} The hypothesis Proposition~\ref{prop:route} runs on
is the guarantee the supplier's contract \emph{states} at each setting: one
weight $p^{(t)}_i$ per slot entry and per setting $t\in\Tset$, with
$\sum_{S\in\mathcal{O}_i}\blk{M}{\gtN}_{S|t}\succeq p^{(t)}_i\idon{}$ at that
setting (\xref{receiver/def:slot} and the first conclusion of
\xref{receiver/thm:emitted-bound}). The weight route forms the weight of a
joint outcome $(y_i,t)$ from $p^{(t)}_i$, the weight of that outcome's own
setting (\eqref{eq:weightroute}). Assumption~\ref{ass:recv} names the
supplier's weights generically, and which claims state them is
Proposition~\ref{prop:route}\,(iv).

\emph{Consuming the minimum would also be sound, and is weakly smaller at
every joint outcome.} The \exttermref{../receiver/contract.pdf}{term-uniform-weight-of-an-entry}{uniform weight of an entry},
$p_i=\min_tp^{(t)}_i$, satisfies the supplier's guarantee at every setting, so
a route consuming it in place of $p^{(t)}_i$, the \emph{uniform route}, would
be certified too. Since
$1-r_{\min,i}\ge0$ and $p_i\le p^{(t)}_i$,
\begin{equation*}
  q_t\Big(\big(1-r_{\min,i}\big)p_i+r_{\min,i}\Big)
  \;\le\;
  q_t\Big(\big(1-r_{\min,i}\big)p^{(t)}_i+r_{\min,i}\Big)=\hat p_{(y_i,t)}
\end{equation*}
at every $(y_i,t)$, and $\lat{\cdot}$ is non-decreasing, so every routed weight
of that route would be at most the one \eqref{eq:weightroute} emits.
\emph{Where the per-setting total $p$ is below one},
Proposition~\ref{prop:weight-monotone}\,(ii) and Corollary~\ref{cor:monotone} then
make that route's target the cleaner one, which grants the adversary at least
as much. At a slot with a single entry, $r=1$, that holds whenever the
uniform route's total is below one. Each of the two totals is at most
$\sum_tq_t\big((1-r_{\min,1})w_t+r_{\min,1}\big)\le1$, $w_t\in[0,1]$ being the
weight it reads at $t$, and reaches one exactly when $r_{\min,1}=1$ or $w_t=1$
at every $t$, because every $q_t$ is positive and a lattice number, which
$\lat{\cdot}$ leaves unchanged (Definitions~\ref{def:joint-family}
and~\ref{def:weight-route}); and $p_1=\min_tp^{(t)}_1$ is $1$ exactly when
every $p^{(t)}_1$ is, so the two totals reach one together. At $r\ge2$, where
every $r_{\min,i}$ vanishes (Corollary~\ref{cor:oneabsorbing}), the
per-setting total can reach one where the uniform one is below it: with two
settings $a,b$, $q_a=q_b=\tfrac12$, $p^{(a)}_1=p^{(b)}_2=1$ and the other two
weights $0$, the totals are $1$ and $0$. The per-setting route is then refused
for its total (Requirement~\ref{req:total-weight-refusal}), a conservative
refusal (Remark~\ref{rem:peqone}), and there is no second target to compare it
with. What the per-setting route gains at entry $i$ is, before rounding,
$\sum_tq_t(1-r_{\min,i})\big(p^{(t)}_i-p_i\big)\ge0$ of routed weight: nothing
at a setting whose weight is the minimum, and at any other setting the
difference its own floor makes, for instance at a bright probe setting, where
the high-photon multiclick floor is expected to be the larger one, an
expectation no claim states.

\emph{The uniform case is the special case, number for number.} If
$p^{(t)}_i=p_i$ for every $t$, which holds at every one-setting family
$\Tset=\{\off\}$, whose one weight is its
\exttermref{../receiver/contract.pdf}{term-uniform-weight-of-an-entry}{uniform weight} (\xref{receiver/def:slot}),
and so at every probe-free receiver, \eqref{eq:weightroute} is evaluated on
the same declared numbers $q_t$, $r_{\min,i}$ and $p_i$ as the uniform route
evaluates, in the same exact substrate and with the same downward rounding. The routed weights,
their total $p$, the complement $1-p$ and everything the flat instance builds
from them are then the numbers that route gives, so at a declaration without
a probe the two routes emit the same numbers.
\end{remark}
\claimmeta{\metalabel{Uses}\mref{Cor.}{cor:monotone}, \mref{Def.}{def:joint-family}, \mref{Prop.}{prop:weight-monotone}, \mref{Def.}{def:weight-route}, \mref{Prop.}{prop:route}, \mref{Cor.}{cor:oneabsorbing}, \extref{../receiver/contract.pdf}{receiver/def:slot}, \extref{../receiver/contract.pdf}{receiver/thm:emitted-bound}
\metasep \metalabel{Assumes}\aref{A-RECV}{ass:recv}
\metasep newly derived
\metasep accepted}

% claim reductions/wpfss/rem:darkcompose
\claimsummary{reductions/wpfss/rem:darkcompose}{This submodule composes with neither of the supplier's two dark-count inequalities and proves its own bound from R itself. Every R it admits satisfies (D1) and (D2) exactly, so the supplier's cap is inactive, its correction is slack, and its exact inequality holds as stated here.}
\claimbookmark{reductions/wpfss/rem:darkcompose}{Remark}{rem:darkcompose}{Composition with the supplier's dark-count bounds}
\begin{remark}[Composition with the supplier's dark-count bounds]\label{rem:darkcompose}
\emph{Two supplier inequalities.} The supplier is the receiver at the version
the artifact names (Assumption~\ref{ass:recv}\,(i)), and it separates two
inequalities about the same object. The first, part of
\xref{receiver/detector-imperfections/thm:dark-counts}\,(iv), reads
$\blk{\widetilde M}{\gtN}_{\MCK}\succeq
\big((1-r_{\min})p_{\mc}+r_{\min}\big)\id^{(>N)}$; the second, in
\xref{receiver/detector-imperfections/prop:dark-slack-bound}\,(iii), is the
same statement with the capped coefficient $\bar r=\min\{r_{\min},1\}$ in
place of $r_{\min}$ and a subtracted correction $\delta_{\dc}$, both defined
there. The supplier says which is which: the second is the inequality a
downstream module composes with
(\xref{receiver/detector-imperfections/req:dark-channel-record}), and the
first is stated for an $R$ satisfying (D1) and (D2) exactly and asserted for
no other. A consumer that used the first where the second is required would
overstate by $(r_{\min}-\bar r)(1-p_{\mc})+\delta_{\dc}$
(\xref{receiver/detector-imperfections/rem:d2slack}\,(e)), first order in the
\exttermref{../receiver/contract.pdf}{term-receiver-tolerances}{receiver tolerances} and in the \emph{unsafe} direction. Which case
holds here is therefore a question this contract must answer rather than
inherit.

\emph{First: this module composes with neither inequality.} It consumes the
matrix $R$ itself, forms $R^{\cg}$ by Lemma~\ref{lem:descent} and its own
per-letter constants $r_{\min,i}=\min_{y'\ne y_i}R^{\cg}_{y_i|y'}$, and proves
\eqref{eq:routebound} inline in Proposition~\ref{prop:route} from POVM
positivity and completeness, the absorbing hypothesis \eqref{eq:absorbing} and
the exact column-stochasticity of $R^{\cg}$, and from nothing else. It reads
none of $r_{\min}$, $\bar r$ or $\delta_{\dc}$ from the artifact. Those three
are the supplier's \emph{raw}-level scalars, minima and masses over $\MCK$;
this module's $r_{\min,i}$ is a different object, indexed by a \emph{coarse}
bounded letter and read across the columns $y'\ne y_i$ of $R^{\cg}$, and
Definitions~\ref{def:joint-family} and~\ref{def:weight-route} say so.

\emph{Second, and this settles the question rather than leaving it to a
reading: every declaration this module admits satisfies (D1) and (D2)
exactly.} \xref{receiver/detector-imperfections/def:darkadm} calls $R$ a
dark-count channel when it is (D1) column-stochastic and (D2) monotone, and
both are decided on the declared matrix \emph{exactly}: (D1) by the
column-stochasticity rule (Requirement~\ref{req:column-stochastic-refusal}),
and (D2) by the descent rule, whose monotonicity half is the implication
$R_{S|S'}>0\Rightarrow S'\subseteq S$
(Requirement~\ref{req:descent-decided-exactly}). The two rules are needed
together and neither covers the other: a strictly positive entry in a
(D2)-forbidden position fails monotonicity while leaving column-stochasticity
untouched, and a \emph{negative} entry in such a position does the reverse, it
satisfies the monotonicity implication vacuously and fails non-negativity.
Consequently, on the set of declarations this module emits for,
$R_{\MCK|S'}\le\sum_SR_{S|S'}=1$ for every column, so the supplier's minimum
$r_{\min}$ of \xref{receiver/detector-imperfections/thm:dark-counts}\,(iv) is
at most $1$ and the cap in $\bar r$ is inactive, $\bar r=r_{\min}$; and
\xref{receiver/detector-imperfections/thm:dark-counts}\,(iii) holds with no
defect, so the $\delta_{\dc}$ that contract emits is slack in a bound rather
than a defect of the datum. The first term of the overstatement is therefore
nil, and the second, $\delta_{\dc}$, is slack: for such an $R$ the first
inequality holds as stated
(\xref{receiver/detector-imperfections/thm:dark-counts}\,(iv)), so using it
overstates nothing. That first inequality, the case of (D1) and (D2) holding
exactly, is the statement available here.

\emph{What that means route by route.} On the supplier's
\exttermref{../receiver/contract.pdf}{term-generated-dark-count-route}{generated dark-count route}, where $R$ is its
\exttermref{../receiver/contract.pdf}{term-independent-product-form}{independent product form}, (D1) and (D2) are a theorem
(\xref{receiver/detector-imperfections/thm:dark-counts}\,(i)) and
$\delta_{\dc}=0$ on that route by its definition in
\xref{receiver/detector-imperfections/prop:dark-slack-bound}, so the two
inequalities coincide there outright. On its \exttermref{../receiver/contract.pdf}{term-declared-dark-count-route}{declared dark-count route}, where
\xref{receiver/detector-imperfections/def:darktol} admits entries down to
$-\varepsilon_4$ and column sums anywhere in
$[1-\varepsilon_4,1+\varepsilon_4]$, a matrix reaches a claim of this contract
only if it satisfies (D1) and (D2) exactly, and is \emph{refused here}
otherwise, under whichever of the two rules it fails. Nothing in
\eqref{eq:weightroute} is restated in $\bar r$ or $\delta_{\dc}$, and nothing
needs to be: that formula is proved for an $R$ satisfying (D1) and (D2)
exactly, and this module admits no other. What is \emph{not} settled here is
the composition in general: what a consumer that admits the supplier's widened
set must do. That question is not this contract's, and this contract does not
answer it: it is a question about the supplier's own receiver tolerances
$\varepsilon_2,\varepsilon_4$, which this module neither reads nor converts,
and \xref{foundations/req:norm-conversion} puts any such conversion at the
instance, explicitly and on that party's own authority
(Assumption~\ref{ass:recv}).
\end{remark}
\claimmeta{\metalabel{Uses}\mref{Def.}{def:joint-family}, \mref{Lem.}{lem:descent}, \mref{Def.}{def:weight-route}, \mref{Prop.}{prop:route}, \extref{../receiver/contract.pdf}{receiver/detector-imperfections/def:darkadm}, \extref{../receiver/contract.pdf}{receiver/detector-imperfections/def:darktol}, \extref{../receiver/contract.pdf}{receiver/detector-imperfections/thm:dark-counts}
\metasep \metalabel{Assumes}\aref{A-RECV}{ass:recv}
\metasep newly derived
\metasep accepted}

% claim reductions/wpfss/rem:receiver-imports
\claimsummary{reductions/wpfss/rem:receiver-imports}{The receiver facts imported here are the claims \hyperlink{claim-reductions/wpfss/prop:route}{Proposition~\ref*{prop:route}} (iv) cites, and the emitted ledger names every supplier assumption they rest on and may name more. A-EIG enters only at an exact window evaluated in floating point; acyclicity needs no supplier assumption.}
\claimbookmark{reductions/wpfss/rem:receiver-imports}{Remark}{rem:receiver-imports}{Supplier assumptions behind the receiver imports}
\begin{remark}[Supplier assumptions behind the receiver imports]
\label{rem:receiver-imports}
The receiver claims this submodule imports are those
Proposition~\ref{prop:route}\,(iv) cites for the hypotheses on the
supplier's objects, read under Assumption~\ref{ass:recv} or, for an
artifact the run built itself from a declaration with a model admissible
exactly and with every weight evaluated rigorously, under
Requirement~\ref{req:recv-row-provenance}, which shows them to hold of its
numbers under the supplier's own ledger. They are read at
the module version the artifact names, an accepted version: for an
artifact from outside the run by Assumption~\ref{ass:recv}\,(i) and
Requirement~\ref{req:supplier-version-accepted}, and for one the run built
as the effective status of the receiver's claims in the run's closure, which
the record of the run reports (Requirement~\ref{req:recv-row-provenance}),
whether or not the row is omitted; there each resolves and
reads as Proposition~\ref{prop:route}\,(iv) cites it; for the
\exttermref{../receiver/contract.pdf}{term-per-setting-weight}{per-setting weights} that version must state them
(Requirement~\ref{req:supplier-version}). Where the receiver is compiled at
its \exttermref{../receiver/contract.pdf}{term-normalized-station}{normalized station}, Corollary~\ref{cor:normalized-flat} imports in
addition \xref{receiver/detector-imperfections/prop:efficiency-normalization},
read under the same assumption. The claims named here are
dependencies of this remark too, so the supplier assumptions the imported
facts rest on are the assumption closure of those dependencies.

\emph{A-EIG is invoked conditionally.} The supplier's \exttermref{../receiver/contract.pdf}{term-evaluation-model}{evaluation model}
(\xref{receiver/def:evaluation-model}), a restriction the
supplier places on itself and not an assumption, governs
\xref{receiver/thm:exact-and-evaluated} and the evaluation of the weights of
\xref{receiver/thm:emitted-bound}, hence the weights $p^{(t)}_i$: each is
evaluated in floating point or, where the supplier encloses its port Grams,
at a setting where the needed enclosures are formed (the artifact records
which), rigorously, and its artifact says which
(\xref{receiver/req:evaluation-label}). A-EIG
supports \xref{receiver/thm:exact-and-evaluated}\,(v) and \emph{the
exact-window branch} of \xref{receiver/thm:emitted-bound}, hence the emitted
$p^{(t)}_{\mc}$ itself at a setting carrying a depth whose weight the supplier
evaluated in floating point. The supplier's
\xref{receiver/ass:eig} is asserted only at a setting whose declaration
carries an \exttermref{../receiver/contract.pdf}{term-exact-window-depth}{exact-window depth} $m^{\ast}_t$ (its
\xref{receiver/def:pmcex}) and whose weight is so evaluated; where no depth is
declared, or the weight is the supplier's \exttermref{../receiver/contract.pdf}{term-enclosed-weight}{enclosed weight}
(\xref{receiver/thm:enclosed-weight}), the emitted value does not depend on
it. \emph{This submodule can neither verify nor improve
it}, and its declared \exttermref{../receiver/contract.pdf}{term-eigenvalue-slack}{eigenvalue slack} $\varepsilon_{\mathrm{eig}}$,
which A-EIG assumes bounds the overestimate in the supplier's floating-point
evaluation of an exact block minimum eigenvalue, is declared separately and
never derived from the \exttermref{../foundations/contract.pdf}{term-declared-tolerance}{declared tolerance}
(\xref{foundations/req:tolerance-scope}\,(iii)), nor from anything else this
module reads. Its invocation record travels with its key
(Requirement~\ref{req:supplier-eig-record}), and a caller that wants it out of
the chain declares no exact window, or a receiver whose port Grams the
supplier encloses and at every setting of which the weight is evaluated
rigorously, as the artifact records.

\emph{The emitted ledger names every assumption of that closure, and may name
more.} The entry of the emitted assumption ledger for
Assumption~\ref{ass:recv} is the supplier's own ledger, not that closure
(Requirements~\ref{req:supplier-ledger-pass-through}
and~\ref{req:supplier-presence-rules}). The two answer two questions: the
closure says which of the supplier's assumptions the facts imported here rest
on; the ledger passes on what the supplier asserts of the artifact it
supplied. Where Requirement~\ref{req:recv-row-provenance} omits that row,
the supplier's ledger reaches the record of the run as the artifact's own,
which the run holds. Under the supplier's presence rules every row is emitted on every
artifact, so no key is conditional. The ledger may name an assumption the
closure is without, which names an assumption more than the imported facts
need. \emph{No assumption that a fact imported here rests on at the
declaration is missing from it}: the supplier's rules omit no key, and an
artifact whose ledger lacks a key they make present is refused
(Requirement~\ref{req:supplier-ledger-refusals}).

\emph{Acyclicity needs no supplier assumption.} The acyclicity of a supplier's element graph is
condition (E2) of \xref{receiver/def:graph}, a well-formedness condition on a
supplier declaration that presents its receiver as an element graph, its
route (iii), the route on which \xref{receiver/thm:compilation-well-defined}
and \xref{receiver/component-catalogue/thm:catalogue-admissibility} produce
the blocks consumed here, so a declaration whose element graph has a cycle is
refused before any block exists and produces no blocks to consume. A
declaration on the supplier's route (i) or (ii) declares no element graph at
all. None of the claims imported here rests on a graph: each is stated of the
supplier's \exttermref{../receiver/contract.pdf}{term-receiver-instrument}{receiver instrument} tuple (\xref{receiver/def:instrument}),
the object on which its three input routes converge, or of an object more
abstract still: a \exttermref{../receiver/contract.pdf}{term-setting-family}{setting family} for \xref{receiver/thm:emitted-bound}
and for the slot claims that read it, a dark-count channel for
\xref{receiver/detector-imperfections/def:darkadm} and
\xref{receiver/detector-imperfections/thm:dark-counts}. The supplier's
\xref{receiver/cor:compiled-instrument}, which states that a well-formed
\exttermref{../receiver/contract.pdf}{term-declared-element-graph}{declared element graph} with a \exttermref{../receiver/contract.pdf}{term-declared-mode-space}{declared mode space}, an input
split and a port map \emph{is} a receiver instrument whose matrix is the
compiled one, is deliberately not a dependency of the supplier's emitting
claims, which are true of every receiver instrument and therefore on every one
of its three routes.

\emph{What is deliberately not imported: the supplier's own admission
procedure for a declared $R$.} The receiver decides (D1) and (D2) for a
\emph{declared} matrix at its own \exttermref{../receiver/contract.pdf}{term-receiver-tolerances}{receiver tolerances} $\varepsilon_4$
and $\varepsilon_2$ (\xref{receiver/detector-imperfections/def:darktol}) and
carries what that widening costs in the capped coefficient $\bar r$ and the
correction $\delta_{\dc}$ of its composed inequality. What this submodule
imports is the \emph{definition} and not that procedure: it decides (D1) and
(D2) on the declared matrix \emph{exactly} and refuses everything else
(Requirements~\ref{req:column-stochastic-refusal}
and~\ref{req:descent-decided-exactly}), so every declaration reaching a claim
of this submodule satisfies (D1) and (D2) outright and
\xref{receiver/detector-imperfections/thm:dark-counts} applies to it verbatim
(Remark~\ref{rem:darkcompose}). \emph{Its residual budgets are different
declared numbers and this submodule reads none of them}: the supplier's
$r_1,\dots,r_6$ and the Hermiticity half of its
\exttermref{../receiver/contract.pdf}{term-hermiticity-and-positivity-defect}{Hermiticity-and-positivity defect} are stated in the \emph{spectral}
norm (\xref{receiver/def:residuals}) and are decided
against its own $\varepsilon_2,\varepsilon_4$, while this module decides each
of its own conditions at exactly zero (Definition~\ref{def:exactadm}) and
carries no number to compare them with. \emph{No conversion between the two is
performed here and none is needed}; an instance that sizes a supplier budget
performs one explicitly, with its dimension factor written down, on its own
authority (\xref{foundations/req:norm-conversion}).
\end{remark}
\claimmeta{\metalabel{Uses}\mref{Lem.}{lem:descent}, \mref{Def.}{def:weight-route}, \mref{Prop.}{prop:route}, \extref{../receiver/contract.pdf}{receiver/cor:unconditioned-family}, \extref{../receiver/contract.pdf}{receiver/def:setting-family}, \extref{../receiver/contract.pdf}{receiver/def:slot}, \extref{../receiver/contract.pdf}{receiver/detector-imperfections/def:darkadm}, \extref{../receiver/contract.pdf}{receiver/detector-imperfections/prop:efficiency-normalization}, \extref{../receiver/contract.pdf}{receiver/detector-imperfections/thm:dark-counts}, \extref{../receiver/contract.pdf}{receiver/prop:slotvalid}, \extref{../receiver/contract.pdf}{receiver/thm:click-povm}, \extref{../receiver/contract.pdf}{receiver/thm:coarse-graining}, \extref{../receiver/contract.pdf}{receiver/thm:emitted-bound}
\metasep \metalabel{Assumes}\aref{A-RECV}{ass:recv}
\metasep \metalabel{Used by}\mref{Req.}{req:supplier-eig-record}
\metasep newly derived
\metasep accepted}

\section{The flat instance and its emission}\label{group-reductions/wpfss/cor:fssinstance}
\sectionintro{On the licensed route a certified weight becomes a zero-constant entry for the \termref{term-flag-state-target}{flag-state target} (\hyperlink{claim-reductions/wpfss/cor:fssinstance}{Corollary~\ref*{cor:fssinstance}}). The post-processed \termref{term-joint-family}{joint family} with its \termref{term-routed-weights}{routed weights} and either target is a \termref{term-reduction-instance}{reduction instance} (\hyperlink{claim-reductions/wpfss/prop:flat-instance}{Proposition~\ref*{prop:flat-instance}}). The main result, \hyperlink{claim-reductions/wpfss/cor:emitted-flat}{Corollary~\ref*{cor:emitted-flat}}, is the core's emission theorem applied to it: on each route the emitted family is exactly admissible and pulls back exactly to the \termref{term-declared-model}{declared model} on the channel system, with its weights rounded down and still certified. A requested two-sector entry is emitted where its certificate is decided on the emitted numbers (\hyperlink{claim-reductions/wpfss/cor:two-sector-decided}{Corollary~\ref*{cor:two-sector-decided}}). Run on a reference of the receiver blocks in place of their snap, the flat instance has a reference companion, whose datum combines those blocks with non-negative weights, so that their factor enclosures reach it (\hyperlink{claim-reductions/wpfss/cor:flat-companion}{Corollary~\ref*{cor:flat-companion}}). Formed on a normalized station and precomposed with the pure-loss channel of the factored loss, the flat instance is a reduction instance for the receiver itself, with the same emission (\hyperlink{claim-reductions/wpfss/cor:normalized-flat}{Corollary~\ref*{cor:normalized-flat}}).}
\sectionimports{\extref{../receiver/contract.pdf}{receiver/detector-imperfections/def:dark-composed}, \extref{../receiver/contract.pdf}{receiver/detector-imperfections/def:normalized-station}, \extref{../receiver/contract.pdf}{receiver/detector-imperfections/def:pure-loss}, \extref{../receiver/contract.pdf}{receiver/detector-imperfections/prop:efficiency-normalization}, \extref{../receiver/contract.pdf}{receiver/thm:click-povm}, \extref{../receiver/contract.pdf}{receiver/thm:coarse-graining}}

% claim reductions/wpfss/cor:fssinstance
\claimsummary{reductions/wpfss/cor:fssinstance}{Under the flag-state target, the certified weight of an outcome becomes a constraint entry with constant 0 whose certificate holds, so the one-entry list is licensed. Read on the family the artefact carries, that constant is emitted unshifted.}
\claimbookmark{reductions/wpfss/cor:fssinstance}{Corollary}{cor:fssinstance}{Licensed flag-sector entry of a certified weight}
\claimleadin{%
On the licensed route a certified weight travels as a constraint entry rather
than inside the target (Definition~\ref{def:routes}). The corollary is read on
the family the artefact carries: its $c=0$ is the constant for the family the
target is built from, and on the emission path that is the
\termref{term-declared-model}{declared model} $M'$ of the flat instance (Proposition~\ref{prop:flat-instance},
Corollary~\ref{cor:emitted-flat}\,(iii); Definition~\ref{def:emitnorm},
Theorem~\ref{thm:emit}\,(iii)). So the constant emitted is $0$, unshifted
(Corollary~\ref{cor:entries}\,(iii)), and nothing in the statement changes; a
branching instance that files the flat instance as a branch carries this entry
as the core's register instance says
(Corollary~\ref{cor:register-instance}\,(ii)). The corollary is about the
flag-state squasher alone, which no further link follows on any route the
engine emits.
}
\begin{corollary}[Licensed flag-sector entry of a certified weight]\label{cor:fssinstance}
In the situation of Proposition~\ref{prop:chain} with the
\termref{term-flag-state-target}{flag-state target} $\Ffss$, fix an outcome $o\in\Vset$ with certified
weight $p_o$ and put
\begin{equation}\label{eq:fssentry}
  \hat O\coloneqq\Ffss_o-p_o\,\Pflag,
  \qquad \bowtie\;=\;\succeq,\qquad c=0 .
\end{equation}
Then \eqref{eq:certificate} holds, so Theorem~\ref{thm:licence} applies and the
single-entry list $\CList=\{(\hat O,\succeq,0)\}$ is licensed with the
admissible set
\begin{equation}\label{eq:fssattset}
  \AttSet_{\CList}(E',\Qsys E'')
  =\Big\{\Aev:\ \Aev^{\dagger}\big(\Ffss_o\otimes\idon{E''}\big)
  \;\succeq\;p_o\,\Aev^{\dagger}\big(\Pflag\otimes\idon{E''}\big)\Big\}.
\end{equation}
\end{corollary}
\claimmeta{\metalabel{Uses}\mref{Thm.}{thm:licence}, \mref{Def.}{def:targets}, \mref{Prop.}{prop:chain}
\metasep \metalabel{Used by}\mref{Prop.}{prop:flat-instance}, \mref{Cor.}{cor:emitted-flat}, \mref{Cor.}{cor:flat-companion}
\metasep newly derived
\metasep accepted}

\begin{proof}
By Proposition~\ref{prop:chain}\,(i),
$\Lfss^{\dagger}(\Ffss_o)=M_o$ and $\Lfss^{\dagger}(\Pflag)=\Pgt$, so by
linearity $\Lfss^{\dagger}(\hat O)=M_o-p_o\Pgt
=\blk{M}{\leqN}_o\dsum\big(\blk{M}{\gtN}_o-p_o\idon{}\big)\succeq0$, the first
block being positive because $M_o$ is and the second by
\eqref{eq:certweights}. That is \eqref{eq:certificate} with $c=0$ and
$\bowtie\,=\,\succeq$. The displayed form of the set is
$\Aev^{\dagger}\big((\Ffss_o-p_o\Pflag)\otimes\idon{E''}\big)\succeq0$
rearranged by linearity of $\Aev^{\dagger}$.
\end{proof}

% claim reductions/wpfss/prop:flat-instance
\claimsummary{reductions/wpfss/prop:flat-instance}{The post-processed receiver family with its routed weights defines a reduction instance, with the weight-preserving target on the encoded route and the flag-state target on the licensed route. Its declared model agrees with the handed one above the cutoff and keeps every certified weight.}
\claimbookmark{reductions/wpfss/prop:flat-instance}{Proposition}{prop:flat-instance}{The flat instance of the post-processed receiver}
\claimleadin{%
The flat instance is the \termref{term-reduction-instance}{reduction instance} of an optical receiver
with no branch register: the \termref{term-joint-family}{joint family} the \termref{term-weight-route}{weight route} forms, with
its blocks below the cutoff as the finite datum, and the blocks above the
cutoff with the \termref{term-routed-weights}{routed weights} entering only through hypotheses that
read no entry of the datum.
}
\begin{proposition}[The flat instance of the post-processed receiver]
\label{prop:flat-instance}
Let the data of Definition~\ref{def:weight-route} pass every refusal rule of
this module's interface; let
$\{M_v\}_{v\in\Vset}$ be the post-processed joint family of its step~2 on
$\Chan$, with $\Hs{\Chan}=\bigoplus_{m\ge0}\Hs{\Chan_m}$, $\Hs{\Chan_{\leqN}}$
finite-dimensional, every $M_v$ block diagonal at the cutoff $N$; let $(p_v)$
be its routed weights with $p=\sum_vp_v<1$; and let a route, encoded or
licensed, be declared (Definition~\ref{def:routes}). Then the \emph{flat instance}
$\mathcal{J}_{\mathrm{flat}}$, given by
\begin{enumerate}[label=(I\arabic*),leftmargin=3.0em]
\item $\Bsys=\Chan$; the \termref{term-data-sector}{data sectors} $(m)$, $0\le m\le N$, in that order,
  with $d_m=\dim\Hs{\Chan_m}$ and the default supports $\Vset^{(m)}=\Vset$;
  the \termref{term-handed-datum}{handed datum} $A^{(m)}_v\coloneqq\blk{M}{=m}_v$, formed from the declared
  blocks, $R^{\cg}$ and $\{q_t\}$ by steps~1--2 of
  Definition~\ref{def:weight-route} in the \termref{term-bare-form}{bare form}, and from the
  declared blocks and $\{q_t\}$ by step~2 alone in the
  \termref{term-dark-composed-form}{dark-composed form}; the \termref{term-hypotheses-h}{hypotheses (H)}: \emph{(H1)} for
  each $t\in\Tset$ the above-cutoff blocks $\{\blk{M}{\gtN}_{y|t}\}_{y\in\Yset}$
  of the supplied conditional POVM form a POVM on $\Chan_{\gtN}$
  (the supplier's fact, Assumption~\ref{ass:recv},
  Proposition~\ref{prop:route}\,(iv)); \emph{(H2)} the routed weights are
  certified, \eqref{eq:certroute} (Proposition~\ref{prop:route});
  \emph{(H3)} $p<1$, decided on entry; and the \termref{term-model-map}{model map}
  \begin{equation}\label{eq:flat-model}
    M_v[A']\;\coloneqq\;\bigoplus_{m\le N}A'^{(m)}_v\ \dsum\ \blk{M}{\gtN}_v,
    \qquad
    \blk{M}{\gtN}_{(y,t)}=q_t\sum_{y'\in\Yset}R^{\cg}_{y|y'}\blk{M}{\gtN}_{y'|t} ;
  \end{equation}
\item the \termref{term-scalar-sector}{scalar sectors} $(f)$, $f\in\Vset$, in the order of $\Vset$, each of
  dimension $1$, with
  \begin{equation}\label{eq:flat-scalars}
    \phi_f(v)\;\coloneqq\;p_v+(1-p)\,\delta_{vf}\quad\text{(encoded route)},
    \qquad
    \phi_f(v)\;\coloneqq\;\delta_{vf}\quad\text{(licensed route)},
  \end{equation}
  so that $\Hs{\Qsys}=\Hs{\Chan_{\leqN}}\dsum\Hs{\Flag}$ and the target
  \eqref{eq:instance-target} is $F_v[A']=\Fwp_v$ resp.\ $\Ffss_v$ of
  Definition~\ref{def:targets} for the model $M[A']$;
\item $\Lambda=\PinS\circ\Lwp$ resp.\ $\PinS\circ\Lfss$ of
  Proposition~\ref{prop:chain}, independent of the datum, with the
  \termref{term-transported-set}{transported set} $\{\Pleq,\Pflag\}\cup\{\Pi_{(m)}\}_{m\le N}\cup
  \{\Pi_{(f)}\}_{f\in\Vset}$ and pull-backs $\Pleq$, $\Pgt$, the
  photon-number projectors $\Pnumon{m}{\Chan}$, and $\Mwp_f$ resp.\
  $\blk{M}{\gtN}_f$ read as operators on $\Chan$ supported in $\Chan_{\gtN}$;
\item an empty entry list on the encoded route, and on the
  licensed route the flag-sector entries
  $\big(\Ffss_o-p_o\Pflag,\succeq,0\big)$ of Corollary~\ref{cor:fssinstance}
  for every bounded letter $o=(y_i,t)$; the \termref{term-class-marginal-identity}{class-marginal identity}
  discharged,
\end{enumerate}
is a reduction instance over the \termref{term-joint-alphabet}{joint alphabet} of
Definition~\ref{def:weight-route}. Its \termref{term-handed-model}{handed model} $M[A]$ is the
post-processed joint family $\{M_v\}$ itself where every $M_v$ is also block
diagonal for the decomposition
$\Hs{\Chan_{\leqN}}=\bigoplus_{m\le N}\Hs{\Chan_m}$, the photon-number blocks
the datum reads, and otherwise the family's pinching onto those blocks below
the cutoff; on the emission path the first case holds, the receiver's
elements being block diagonal in every photon number
(Proposition~\ref{prop:route}\,(iv)), so that there the handed model is the
post-processed joint family itself. Its declared model, at the
\termref{term-declared-datum}{declared datum} $A'$ the normalizer of
Definition~\ref{def:emitnorm} produces from $A$, is
\begin{equation}\label{eq:flat-declared}
  M'_v\;\coloneqq\;M_v[A']\;=\;\bigoplus_{m\le N}A'^{(m)}_v\ \dsum\ \blk{M}{\gtN}_v ,
\end{equation}
a POVM on $\Chan$ block diagonal at the cutoff, agreeing with the handed one
above it and satisfying every certified-weight inequality \eqref{eq:certroute}
the handed one does.
\end{proposition}
\claimmeta{\metalabel{Uses}\mref{Def.}{def:system-setting}, \mref{Lem.}{lem:linear}, \mref{Def.}{def:joint-family}, \mref{Cor.}{cor:blockdiag}, \mref{Def.}{def:exactadm}, \mref{Def.}{def:reduction-instance}, \mref{Def.}{def:emitnorm}, \mref{Def.}{def:targets}, \mref{Prop.}{prop:chain}, \mref{Def.}{def:routes}, \mref{Lem.}{lem:descent}, \mref{Def.}{def:weight-route}, \mref{Prop.}{prop:route}, \mref{Cor.}{cor:fssinstance}
\metasep \metalabel{Assumes}\aref{A-RECV}{ass:recv}, \aref{A-PROBE}{ass:probe}, \aref{A-TSM}{ass:tsm}
\metasep \metalabel{Used by}\mref{Prop.}{prop:precomposed-blocks}, \mref{Cor.}{cor:emitted-flat}, \mref{Cor.}{cor:two-sector-decided}, \mref{Cor.}{cor:flat-companion}, \mref{Cor.}{cor:normalized-flat}
\metasep newly derived
\metasep implemented (effective: accepted)
\metasep \metalabel{Code}\coderef{https://github.com/fraunhofersg/certiqs-key/blob/v0.1.0/src/certiqs_key/reductions/wpfss/construction.py}{certiqs_key.reductions.wpfss.construction:flat_instance}}

\begin{proof}
(I1) \emph{Formation and admissibility of the handed datum.} For $v=(y,t)$,
$A^{(m)}_v=\blk{M}{=m}_{(y,t)}=q_t\sum_{y'}R^{\cg}_{y|y'}\blk{M}{=m}_{y'|t}$
by \eqref{eq:darkstep} and \eqref{eq:jointfamily} read block by block ---
finitely many products and sums of declared entries in $\Fldi$. It is
Hermitian as a real combination of Hermitian blocks; positive semidefinite as
a combination of positive semidefinite blocks with non-negative coefficients
($R^{\cg}\ge0$ entrywise, $q_t>0$); and \termref{term-class-complete}{class complete}:
$\sum_{y\in\Yset}A^{(m)}_{(y,t)}=q_t\sum_{y'}\big(\sum_yR^{\cg}_{y|y'}\big)
\blk{M}{=m}_{y'|t}=q_t\sum_{y'}\blk{M}{=m}_{y'|t}=q_t\idon{\Chan_m}$ by
column-stochasticity of $R^{\cg}$ (Lemma~\ref{lem:descent}) and completeness
of the conditional POVM at $t$ read on block $m$. In the dark-composed form
$A^{(m)}_{(y,t)}=q_t\blk{\widetilde M}{=m}_{y|t}$ with the supplied blocks,
a product of declared entries with $q_t$; it is Hermitian and positive
semidefinite because the supplied blocks are, and
$\sum_yA^{(m)}_{(y,t)}=q_t\idon{\Chan_m}$ because the supplied family is
complete on its block (Definition~\ref{def:weight-route}). The default
supports meet every class. So $A\in\mathcal{D}_{\mathrm{adm}}$ in either
form.
\emph{Stability.} (H1) is about the blocks $m>N$, which are not in the datum;
(H2) is an inequality between above-cutoff blocks and scalars --- the proof
of Proposition~\ref{prop:route} restricts to the $\gtN$ block and reads no
preserved block; (H3) reads the weights. None reads an entry of the datum.
\emph{The model map.} For every $A'\in\mathcal{D}_{\mathrm{adm}}$, $M_v[A']$
is Hermitian and positive semidefinite block by block; it is complete below
the cutoff by class completeness and $\sum_tq_t=1$ (the first consequence
of Definition~\ref{def:reduction-instance}), and above it by (H1) and
column-stochasticity: $\sum_{(y,t)}\blk{M}{\gtN}_{(y,t)}
=\sum_tq_t\sum_{y'}\big(\sum_yR^{\cg}_{y|y'}\big)\blk{M}{\gtN}_{y'|t}
=\sum_tq_t\idon{\Chan_{\gtN}}=\idon{\Chan_{\gtN}}$. It is block diagonal at
the cutoff by construction, and $\Hs{\Chan}$ is separable
(Definition~\ref{def:system-setting}). So $M[A']$ is a POVM on $\Chan$ under (H), for
every \termref{term-admissible-datum}{admissible datum}. At the handed datum,
$M_v[A]=\bigoplus_{m\le N}\blk{M}{=m}_v\dsum\blk{M}{\gtN}_v$ keeps the
diagonal blocks of $M_v$ for the blocks $m\le N$ and $\gtN$ and drops its
off-diagonal blocks, of which only those between two blocks $m\ne m'\le N$
can be non-zero, $M_v$ being block diagonal at the cutoff. So $M_v[A]$ is the
pinching of $M_v$ onto those blocks, and it is $M_v$ exactly where those
off-diagonal blocks vanish. On the emission path they do: every $M_v$ is
block diagonal in every photon number by Proposition~\ref{prop:route}\,(iv),
read under Assumption~\ref{ass:recv}.

(I2) On the encoded route every $p_v$ is a lattice value rounded
\emph{down} from $\hat p_v$ (\eqref{eq:weightroute}, \eqref{eq:latdown}), a
certified lower bound by (H2), and $1-p$ is recomputed from the rounded
$p_v$ and is a lattice value as their complement; $\phi_f(v)\ge0$ because
$p_v\ge0$ and $p<1$, and $\phi_f(v)\le p_v+(1-p)\le p+(1-p)=1$; and
$\sum_v\phi_f(v)=p+(1-p)=1$ exactly. On the licensed route
$\phi_f(v)=\delta_{vf}\in\{0,1\}$ with $\sum_v\delta_{vf}=1$. Reading
\eqref{eq:instance-target} at these sectors,
$F_v[A']=\bigoplus_{m\le N}A'^{(m)}_v\dsum\bigoplus_{f}\phi_f(v)
=\blk{M[A']}{\leqN}_v\dsum\big(p_v\idon{\Flag}+(1-p)\dyad{v}_{\Flag}\big)$,
the operator $p_v\idon{\Flag}+(1-p)\dyad{v}$ being diagonal in the flag basis
with entries $\phi_f(v)$ --- that is $\Fwp_v$ of \eqref{eq:WPFSS} for the
model $M[A']$; on the licensed route the flag block is $\dyad{v}$ and the
target is $\Ffss_v$ of \eqref{eq:FSS}.

(I3) Fix $A'\in\mathcal{D}_{\mathrm{adm}}$ under (H). By (I1), $M[A']$ is a
POVM on $\Chan$ block diagonal at the cutoff; $(p_v)$ are \termref{term-certified-weights}{certified weights}
for it in the sense of Definition~\ref{def:targets}, because
\eqref{eq:certweights} concerns the above-cutoff blocks only, which are the
handed ones, and at a bounded letter it is (H2) while at every other $v$ it
reads $\blk{M}{\gtN}_v\succeq0$; and $p<1$ by (H3). So
Proposition~\ref{prop:chain}\,(i) applies to $M[A']$:
$(M[A'];\Pleq,\Pgt)\dsq_{\Lwp}(\Fwp;\Pleq,\Pflag)$, resp.\ with $\Lfss$ and
$\Ffss$. The channels \eqref{eq:Lwp} and \eqref{eq:Lfss} read $\Pleq$,
$\Pgt$, the above-cutoff blocks and the weights, and no entry of the datum,
so they are independent of $A'$. The sector projectors: with
$\Pi_{(m)}=\Pnumon{m}{\Chan}\dsum0$ and $\Pi_{(f)}=0\dsum\dyad{f}$, the adjoint
\eqref{eq:WPadjoint} gives $\Lwp^{\dagger}(\Pi_{(m)})=\Pnumon{m}{\Chan}$ and
$\Lwp^{\dagger}(\Pi_{(f)})=\sum_v\bra{v}\dyad{f}\ket{v}\Mwp_v=\Mwp_f$, and at
$p_v\equiv0$ the same computation gives $\blk{M}{\gtN}_f$ for $\Lfss$. Every
transported operator is block diagonal for $\SecSet$, $\Pleq$ and $\Pflag$
being sums of sector projectors, so Corollary~\ref{cor:blockdiag}\,(i) gives
the relation \eqref{eq:instance-relation} for $\PinS\circ\Lwp$ resp.\
$\PinS\circ\Lfss$ with the same pull-backs.

(I4) On the encoded route the list is empty; the certificate every instance
carries for a sector projector is not an entry of it but travels with the
transported set (the third consequence of
Definition~\ref{def:reduction-instance},
Corollary~\ref{cor:entries}\,(i)). On the licensed route, for a
bounded letter $o$, $\hat O=\Ffss_o-p_o\Pflag=\sum_{s\in\SecSet}\Pi_sF_o\Pi_s
-p_o\sum_{f}\Pi_{(f)}$ has the form \eqref{eq:instance-entry} with
$\alpha_{o,s}=1$ for every $s$, every other $\alpha=0$, $Z=-p_o\Pflag$ and
the lattice coefficient $p_o$; it is a combination of transported operators;
and its certificate is the computation of Corollary~\ref{cor:fssinstance},
$\Lfss^{\dagger}(\hat O)=\bigoplus_{m\le N}A'^{(m)}_o\dsum
\big(\blk{M}{\gtN}_o-p_o\idon{}\big)\succeq0$, the first summand positive by
admissibility of $A'$ and the second by (H2) --- for every admissible datum.
The class-marginal identity: $\sum_{v\in\Vset_t}M_v[A']
=\bigoplus_{m\le N}q_t\idon{\Chan_m}\dsum\sum_{y}\blk{M}{\gtN}_{(y,t)}$ by
class completeness, and $\sum_y\blk{M}{\gtN}_{(y,t)}
=q_t\sum_{y'}\big(\sum_yR^{\cg}_{y|y'}\big)\blk{M}{\gtN}_{y'|t}
=q_t\idon{\Chan_{\gtN}}$ by column-stochasticity and (H1) --- which is
Lemma~\ref{lem:linear}\,(ii) read block by block; so it is discharged for
every admissible datum.

The declared model is \eqref{eq:declmodel} at the model map
\eqref{eq:flat-model}, which is \eqref{eq:flat-declared}; it agrees with the
\termref{term-handed-model}{handed model} above the cutoff because the normaliser touches data sectors
only, and \eqref{eq:certroute} is a statement about those blocks alone.
\end{proof}

% claim reductions/wpfss/cor:emitted-flat
\claimsummary{reductions/wpfss/cor:emitted-flat}{The flat instance's emitted family is exactly admissible and pulls back exactly to the declared model on the channel system, its weights rounded down yet certified. Where the normalizer moved a block below the cutoff, nearness to the unrounded family bounds nothing.}
\claimbookmark{reductions/wpfss/cor:emitted-flat}{Corollary}{cor:emitted-flat}{Exact emission of the flat instance}
\claimleadin{%
The core's emission theorem, specialised to the flat instance: it adds the
numbers of the weight route to the list of exact ones, names the rounding
directions the emission record carries (Requirement~\ref{req:emission-record}),
and says, route by route, which constraint entries the artefact carries.
The six conditions of part~(ii) are those of
\xref{foundations/req:exact-emission}, and a family that fails one of them is
refused by Requirement~\ref{req:emission-refusal}.
}
\begin{mainclaim}
\begin{corollary}[Exact emission of the flat instance]\label{cor:emitted-flat}
Let the \exttermref{../foundations/contract.pdf}{term-declared-data}{declared data} pass every refusal rule of this module's
interface, let
$\mathcal{J}_{\mathrm{flat}}$ be the flat instance of
Proposition~\ref{prop:flat-instance} at the declared route with handed datum
$A$, and let $A'$ be the datum the normaliser of
Definition~\ref{def:emitnorm} declares. Then Theorem~\ref{thm:emit} holds of
$\mathcal{J}_{\mathrm{flat}}$, and in particular:
\begin{enumerate}[label=(\roman*),leftmargin=3.0em]
\item every number on the list of Theorem~\ref{thm:emit}\,(i) and, in
  addition, the descended matrix $R^{\cg}$, the constants $r_{\min,i}$, the
  exact weights $\hat p_v$ of \eqref{eq:weightroute}, the total $p$ and its
  complement $1-p$ are formed from the declared entries and scalars by
  finitely many operations in $\Fldi$, floors and comparisons; the
  \termref{term-emission-displacement}{emission displacement} $\Delta_m$ per preserved block is the one
  emitted number that is not;
\item the emitted family, $\Fwp$ resp.\ $\Ffss$ of the declared model
  \eqref{eq:flat-declared}, satisfies conditions (a)--(e) of
  Definition~\ref{def:exactadm} together with \eqref{eq:latinv} of every
  entry and is decided so before it is emitted;
\item for every $v\in\Vset$,
  \begin{equation}\label{eq:emitted-flat}
    \Lambda^{\dagger}(F_v)\;=\;M'_v\;=\;\bigoplus_{m\le N}A'^{(m)}_v\ \dsum\ \blk{M}{\gtN}_v
    \qquad\text{exactly, on }\Chan ,
  \end{equation}
  both sides operators on $\Chan$ itself: the physical system of this
  instance is the channel system and the identity is on it, not on an
  ampliation. $M'$ is a POVM on $\Chan$ block diagonal at the cutoff for which
  $(p_v)$ are \termref{term-certified-weights}{certified weights} in the sense of
  Definition~\ref{def:targets}, and Proposition~\ref{prop:chain},
  Proposition~\ref{prop:route} and Corollary~\ref{cor:fssinstance} hold
  verbatim of it;
\item the per-outcome weights are $p_v=\lat{\hat p_v}$, rounded \emph{down},
  and \eqref{eq:certroute} holds for $M'$ at those values
  (Proposition~\ref{prop:route}); the complement $1-p$ is \emph{recomputed}
  from the rounded weights and is never rounded on its own; a caller's
  constant off the lattice is rounded toward the enlarging side
  (Corollary~\ref{cor:entries}\,(iii)); and there are no other rounding
  directions;
\item where some $\Delta_m>0$ no bound relates a key length computed from
  $M'$ alone to one computed from the pre-snap post-processed family, or from
  the receiver's own output, and nothing else relates them
  (Theorem~\ref{thm:emit}\,(v)).
\end{enumerate}
The emitted constraint list is \emph{empty on the encoded route}; on the
licensed route it carries the flag-sector entries
$(\Ffss_o-p_o\Pflag,\succeq,0)$ of Corollary~\ref{cor:fssinstance} at $c=0$,
unshifted, and, on request, the class-marginal entries
$\big(\sum_{v\in\Vset_t}F_v,=,q_t\big)$, which the flat instance discharges,
a default \termref{term-carried-operators}{carried operator} as an entry
(Corollary~\ref{cor:entries}\,(i)),
and any further requested entry that Corollary~\ref{cor:entries}\,(iii)
admits, which for this instance is the two-sector entry
$(\Ffss_e-a\,\Pleq-b\,\Pflag,\succeq,0)$ where its decided certificate of
Corollary~\ref{cor:two-sector-decided} holds, all under the licence of that
route (Corollary~\ref{cor:entries}). The instance's default
\termref{term-carried-operators}{carried operators},
the sector projectors and the two projectors $\Pleq$ and $\Pflag$ it
transports besides, \emph{are emitted with the transported set at $c=0$ with
$\succeq$ on either route} (Corollary~\ref{cor:entries}\,(i),
Proposition~\ref{prop:flat-instance}\,(I3)). These statements are about the
instance's link and its emitted family.
\end{corollary}
\end{mainclaim}
\claimmeta{\metalabel{Uses}\mref{Def.}{def:exactadm}, \mref{Def.}{def:substrate}, \mref{Lem.}{lem:exactpsd}, \mref{Def.}{def:reduction-instance}, \mref{Def.}{def:emitnorm}, \mref{Thm.}{thm:emit}, \mref{Cor.}{cor:entries}, \mref{Def.}{def:targets}, \mref{Prop.}{prop:chain}, \mref{Lem.}{lem:descent}, \mref{Def.}{def:weight-route}, \mref{Prop.}{prop:route}, \mref{Cor.}{cor:fssinstance}, \mref{Prop.}{prop:flat-instance}, \mref{Cor.}{cor:two-sector-decided}
\metasep \metalabel{Assumes}\aref{A-RECV}{ass:recv}, \aref{A-PROBE}{ass:probe}, \aref{A-TSM}{ass:tsm}
\metasep \metalabel{Used by}\mref{Cor.}{cor:normalized-flat}, \mref{Req.}{req:normalized-link-record}
\metasep newly derived
\metasep implemented (effective: accepted)
\metasep \metalabel{Code}\coderef{https://github.com/fraunhofersg/certiqs-key/blob/v0.1.0/src/certiqs_key/reductions/wpfss/api.py}{certiqs_key.reductions.wpfss.api:reduce_flat}}

\begin{proof}
Proposition~\ref{prop:flat-instance} makes $\mathcal{J}_{\mathrm{flat}}$ a
reduction instance, so Theorem~\ref{thm:emit} applies to it with the data
passing every refusal rule.

(i) The theorem's list, and in addition: every entry of
$R^{\cg}$ is a finite sum of declared entries of $R$ (Lemma~\ref{lem:descent});
$r_{\min,i}$ is a minimum over finitely many of them, decided by comparisons;
$\hat p_{(y_i,t)}$ is two products and two sums of lattice values;
$p$ is a finite sum and $1-p$ a difference; and $\lat{\cdot}$ is a floor and
a comparison in $\Fld$, all exact in the substrate
(Lemma~\ref{lem:exactpsd}, Definition~\ref{def:substrate}). The displacement
\eqref{eq:displacement} at the data sector $(m)$ is $\Delta_m$.

(ii) is the theorem's part (ii), the emitted family being
\eqref{eq:instance-target} at the flat instance, which
Proposition~\ref{prop:flat-instance}\,(I2) identifies with $\Fwp$ resp.\
$\Ffss$ of the declared model.

(iii) is the theorem's part (iii) with the model map \eqref{eq:flat-model}:
$\Mdec_v=M_v[A']=M'_v$ of \eqref{eq:flat-declared}, an operator on $\Bsys=\Chan$.
That $(p_v)$ are certified weights for $M'$ is (H2) at $A'$, available by the
\termref{term-stability-clause}{stability clause} because \eqref{eq:certroute} is a statement about the
above-cutoff blocks, which the normaliser does not touch; that $M'$ is a POVM
block diagonal at the cutoff is (I1) at $A'$; and $p<1$ is (H3). Those are
exactly the hypotheses Proposition~\ref{prop:chain} and
Corollary~\ref{cor:fssinstance} place on their data, so each applies to $M'$
word for word, and Proposition~\ref{prop:route} is a statement about the
above-cutoff blocks and the declared scalars, which are common to the handed
and the declared model.

(iv) The weights are Proposition~\ref{prop:route}
through Definition~\ref{def:weight-route}; the complement and the constant
rule are the theorem's part (iv), the \termref{term-scalar-value}{scalar values}
\eqref{eq:flat-scalars}
being the only rounded numbers of this instance's own.

(v) is the theorem's part (v).

The constraint list. The default carried operators, with the transported
set and, where a caller requests one as an entry, in the list, and the
entries are Corollary~\ref{cor:entries}\,(i), (ii), the
flat instance discharging the class-marginal identity, and (iii) for the
instance's own list, which on the licensed route is the flag-sector entry of
Corollary~\ref{cor:fssinstance} at its constant $0$; on the encoded route
that corollary emits nothing, the instance's own list being empty there
(Proposition~\ref{prop:flat-instance}\,(I4)). A further requested entry is
emitted by (iii) only on a decided certificate that a claim stated for the
flat instance gives; the one such claim is
Corollary~\ref{cor:two-sector-decided}, whose part (iii) is that decided
certificate for the two-sector entry at the declared datum.
\end{proof}

% claim reductions/wpfss/cor:two-sector-decided
\claimsummary{reductions/wpfss/cor:two-sector-decided}{On the licensed route the requested two-sector entry keeps its certificate wherever every block of its outcome below the cutoff clears its first constant and its second constant is at most the outcome's certified weight. Both conditions are decided exactly on the emitted numbers.}
\claimbookmark{reductions/wpfss/cor:two-sector-decided}{Corollary}{cor:two-sector-decided}{The requested two-sector entry, decided on the emitted numbers}
\claimleadin{%
Remark~\ref{rem:published} compares the scalar flag-weight constraint of
the literature with the operator entry \eqref{eq:two-sector-entry}, whose two
constants are the least eigenvalues of the two halves of an element. A
request declares the constants instead, and the entry is emitted only where
its certificate is decided on the emitted numbers. The above-cutoff half has
no finite presentation, so it is decided through the certified weight; the
below-cutoff half is decided on the emitted blocks themselves.
}
\begin{corollary}[The requested two-sector entry, decided on the emitted numbers]\label{cor:two-sector-decided}
Let $\mathcal{J}_{\mathrm{flat}}$ be the flat instance of
Proposition~\ref{prop:flat-instance} on the licensed route, with channel
$\Lambda$ and routed weights $(p_v)$, and let $A''$ be an admissible datum of
it over any substrate of Definition~\ref{def:substrate} over which it is a
reduction instance, $\Fld=\mathbb{R}$ included. Let $e\in\Vset$, let $a$ and $b$ be declared real constants, and put
\begin{equation}\label{eq:two-sector-at}
  \hat O[A'']\;\coloneqq\;\Ffss_e[A'']-a\,\Pleq-b\,\Pflag ,
\end{equation}
written in the form \eqref{eq:instance-entry} with coefficient $1$ for $e$
on every sector, every other coefficient $0$, and $Z=-a\,\Pleq-b\,\Pflag$.
Then:
\begin{enumerate}[label=(\roman*),leftmargin=3.0em]
\item its pull-back is
  \begin{equation}\label{eq:two-sector-pullback}
    \Lambda^{\dagger}\big(\hat O[A'']\big)
    \;=\;\bigoplus_{m\le N}\big(A''^{(m)}_e-a\,\idon{d_m}\big)
    \ \dsum\ \big(\blk{M}{\gtN}_e-b\,\idon{}\big);
  \end{equation}
\item if
  \begin{equation}\label{eq:two-sector-condition}
    A''^{(m)}_e-a\,\idon{d_m}\;\succeq\;0\quad(m\le N)
    \qquad\text{and}\qquad b\;\le\;p_e ,
  \end{equation}
  then $\Lambda^{\dagger}(\hat O[A''])\succeq0$, which is the certificate
  \eqref{eq:certificate} of the entry $(\hat O[A''],\succeq,0)$;
\item the first condition of \eqref{eq:two-sector-condition} is necessary for
  that certificate. At the declared datum $A'$ of Theorem~\ref{thm:emit},
  whose blocks and routed weights are emitted lattice values, and for
  constants $a$ and $b$ that are lattice values, both conditions are decided
  by finitely many field operations and sign comparisons on these numbers,
  and \eqref{eq:two-sector-condition} is a decided certificate for the entry
  $(\hat O[A'],\succeq,0)$ in the sense of
  Corollary~\ref{cor:entries}\,(iii).
\end{enumerate}
\end{corollary}
\claimreading{A constraint you request in the published two-constant form is added to the model only after the engine has checked exactly, on the numbers it emits, that the receiver blocks below the photon-number cutoff clear the first constant and that the second constant does not exceed the certified high-photon weight the receiver supplies. Where either check fails the request is refused; the constraint is never taken on your word.}
\claimmeta{\metalabel{Uses}\mref{Thm.}{thm:licence}, \mref{Def.}{def:substrate}, \mref{Lem.}{lem:exactpsd}, \mref{Def.}{def:reduction-instance}, \mref{Thm.}{thm:emit}, \mref{Cor.}{cor:entries}, \mref{Def.}{def:targets}, \mref{Prop.}{prop:chain}, \mref{Prop.}{prop:flat-instance}
\metasep \metalabel{Assumes}\aref{A-RECV}{ass:recv}
\metasep \metalabel{Used by}\mref{Cor.}{cor:emitted-flat}, \mref{Cor.}{cor:flat-companion}
\metasep newly derived
\metasep implemented (effective: accepted)
\metasep \metalabel{Code}\coderef{https://github.com/fraunhofersg/certiqs-key/blob/v0.1.0/src/certiqs_key/reductions/constraints.py}{certiqs_key.reductions.constraints:_two_sector_certificate}}

\begin{proof}
No step of (i) and (ii) reads a property of the substrate beyond its
field operations and its order, so they hold over each substrate the
statement admits. (i) The instance at $A''$ is a reduction instance by the stability clause
of Definition~\ref{def:reduction-instance}, and (I3) of
Proposition~\ref{prop:flat-instance} holds there: its channel
$\Lambda=\PinS\circ\Lfss$ does not depend on the datum,
$\Lambda^{\dagger}(\Ffss_e[A''])=M_e[A'']$, and $\Pleq$ and $\Pflag$ are
transported with the pull-backs $\Pleq$ and $\Pgt$
(Proposition~\ref{prop:chain}\,(i)). By linearity of the adjoint and the
model map \eqref{eq:flat-model},
\begin{equation*}
  \Lambda^{\dagger}\big(\hat O[A'']\big)
  \;=\;M_e[A'']-a\,\Pleq-b\,\Pgt
  \;=\;\bigoplus_{m\le N}\big(A''^{(m)}_e-a\,\idon{d_m}\big)
  \ \dsum\ \big(\blk{M}{\gtN}_e-b\,\idon{}\big),
\end{equation*}
since $\Pleq$ is the identity on every photon-number block $m\le N$ and zero
above the cutoff, and $\Pgt$ the reverse.

(ii) The blocks below the cutoff are positive semidefinite by the first
condition. The block above the cutoff is the handed instance's, the model
map reading no entry of the datum there; the routed weights are
\termref{term-certified-weights}{certified weights} by (H2) of Proposition~\ref{prop:flat-instance}, so
\eqref{eq:certweights} gives $\blk{M}{\gtN}_e\succeq p_e\idon{}\succeq
b\,\idon{}$ by the second condition. A direct sum of positive semidefinite
blocks is positive semidefinite, so $\Lambda^{\dagger}(\hat O[A''])\succeq0$,
which is \eqref{eq:certificate} with $c=0$ and $\bowtie\,=\,\succeq$.

(iii) A direct summand of a positive semidefinite operator is positive
semidefinite, so the certificate implies the first condition. At $A'$ the
blocks are lattice values, $A'\in\mathcal{D}_{\kappa}$ by
Theorem~\ref{thm:emit}\,(ii), and so are the routed weights, $a$ and $b$;
lattice values are dyadic rationals. Each $A'^{(m)}_e-a\,\idon{d_m}$ is
therefore Hermitian, $A'$ being admissible and $a$ real, with entries in
$\mathbb{Q}(i)$, and its positivity is decided by \eqref{eq:psdtest} of
Lemma~\ref{lem:exactpsd} over the substrate $\Fld=\mathbb{Q}$, which
Definition~\ref{def:substrate} admits; $b\le p_e$ is one comparison of two
rationals. The condition is thus one on the emitted numbers and the declared
constants, decided exactly from them, and by (ii) it is sufficient for
\eqref{eq:certificate} under the instance's hypotheses (H), which is what a
decided certificate is.
\end{proof}

% claim reductions/wpfss/cor:flat-companion
\claimsummary{reductions/wpfss/cor:flat-companion}{Run on any Hermitian, positive receiver blocks complete per setting, the flat instance forms, as a non-negative combination of them, an admissible datum over the reals, as the core companion asks. A requested two-sector entry keeps its certificate where the centres clear its constant by the radius.}
\claimbookmark{reductions/wpfss/cor:flat-companion}{Corollary}{cor:flat-companion}{The flat instance has a reference companion}
\claimleadin{%
The core states the reference companion of any reduction instance under one
hypothesis, (H-REF): that the instance's rule forms an admissible datum from
any family of blocks with the properties a reference has, and that the
instance holds over the reals. For the flat instance both follow from the
proof that it is an instance at all (Proposition~\ref{prop:flat-instance}),
which reads of the receiver blocks only those properties and of an
admissible datum only its admissibility. The companion's entries on the
licensed route are therefore the flag-sector entries of
Corollary~\ref{cor:fssinstance} at the datum formed from the reference, and its emission is the flat emission
(Corollary~\ref{cor:emitted-flat}) with the datum formed from the reference
in its data sectors and no normalizer. A caller may request besides, on the
licensed route, the operator \eqref{eq:two-sector-entry} of
Remark~\ref{rem:published}. Its certificate on the emitted numbers says
nothing about the reference, so the corollary re-establishes it at the
companion: the below-cutoff half from the centres of the companion's datum,
which must clear the declared constant by the radius, and the above-cutoff
half by deciding the declared constant against the certified weight of that
block \eqref{eq:certweights}, which no block of the receiver enters.
}
\begin{corollary}[The flat instance has a reference companion]\label{cor:flat-companion}
Let the flat instance $\mathcal{J}_{\mathrm{flat}}$ of
Proposition~\ref{prop:flat-instance} be given, at either route and in either
form, with $K$ the declared low-photon blocks from which its handed datum is
formed by Definition~\ref{def:weight-route} (the conditional POVM blocks
$\blk{M}{=m}_{y'|t}$, $m\le N$, of each setting $t$ in the bare form, the
supplied dark-composed blocks $\blk{\widetilde M}{=m}_{y|t}$ in the
dark-composed form). Then $\mathcal{J}_{\mathrm{flat}}$ satisfies (H-REF) of
Corollary~\ref{cor:reference-companion}: for every family $G$ over the same
blocks, settings, outcomes and dimensions whose blocks are Hermitian and
positive semidefinite, with $\sum_{y'}G^{(m)}_{y'|t}=\idon{\Chan_m}$ exactly
for every $t$ and $m\le N$, the datum the same rule forms from $G$ is
admissible, and $\mathcal{J}_{\mathrm{flat}}$ is a reduction instance over
the substrate $\Fld=\mathbb{R}$. Its rule of (I1) is a non-negative
combination in the sense of Corollary~\ref{cor:companion-factor-enclosure}:
the element $(y,t)$ of the datum it forms from a family $K'$ with the
properties just listed for $G$ is, on the data sector $(m)$, $q_t\sum_{y'}R^{\cg}_{y|y'}K'^{(m)}_{y'|t}$ in the bare
form and $q_tK'^{(m)}_{y|t}$ in the dark-composed form. At the datum $A^{\mathrm{ref}}$ so formed,
its target is $\Fwp[A^{\mathrm{ref}}]$ on the encoded route and
$\Ffss[A^{\mathrm{ref}}]$ on the licensed route, and on the licensed route its
entries are $\big(\Ffss_o[A^{\mathrm{ref}}]-p_o\Pflag,\succeq,0\big)$ for
every bounded letter $o$, with the routed weights $p_o$ of the handed
instance. On the licensed route, let the request carry besides an entry
$(\hat O,\succeq,0)$ with $\hat O=\Ffss_e-a\,\Pleq-b\,\Pflag$, at an outcome
$e\in\Vset$ and declared real constants $a$ and $b$, written in the form
\eqref{eq:instance-entry} with coefficient $1$ for $e$ on every sector, every
other coefficient $0$, and $Z=-a\,\Pleq-b\,\Pflag$, so that its companion form
(Corollary~\ref{cor:reference-companion}\,(vi)) is
$\hat O^{\mathrm{ref}}=\Ffss_e[A^{\mathrm{ref}}]-a\,\Pleq-b\,\Pflag$. Let
$(C,R)$ be an entrywise enclosure of $A^{\mathrm{ref}}$ whose centres are
Hermitian, and put $\rho_m\coloneqq\max_i\sum_lR^{(m)}_{e,il}$ for $m\le N$.
If
\begin{equation}\label{eq:companion-two-sector}
  C^{(m)}_e-(a+\rho_m)\,\idon{d_m}\;\succeq\;0\quad(m\le N)
  \qquad\text{and}\qquad b\;\le\;p_e ,
\end{equation}
then $\Lambda^{\dagger}(\hat O^{\mathrm{ref}})\succeq0$, with $\Lambda$ the
instance's channel, which is the same at every datum: the companion form
carries the certificate of the entry it is formed from.
\end{corollary}
\claimmeta{\metalabel{Uses}\mref{Thm.}{thm:licence}, \mref{Def.}{def:substrate}, \mref{Def.}{def:reduction-instance}, \mref{Cor.}{cor:reference-companion}, \mref{Cor.}{cor:companion-factor-enclosure}, \mref{Def.}{def:targets}, \mref{Prop.}{prop:chain}, \mref{Lem.}{lem:descent}, \mref{Def.}{def:weight-route}, \mref{Cor.}{cor:fssinstance}, \mref{Prop.}{prop:flat-instance}, \mref{Cor.}{cor:two-sector-decided}
\metasep \metalabel{Assumes}\aref{A-RECV}{ass:recv}, \aref{A-PROBE}{ass:probe}, \aref{A-TSM}{ass:tsm}
\metasep \metalabel{Used by}\mref{Cor.}{cor:normalized-flat}
\metasep newly derived
\metasep accepted}

\begin{proof}
\emph{Admissible datum.} Part (I1) of the proof of
Proposition~\ref{prop:flat-instance} forms the datum as
$A^{(m)}_{(y,t)}=q_t\sum_{y'}R^{\cg}_{y|y'}\blk{M}{=m}_{y'|t}$ in the bare
form and $q_t\blk{\widetilde M}{=m}_{y|t}$ in the dark-composed form, and
shows it Hermitian, positive semidefinite and class complete using, of the
blocks, only that they are Hermitian, positive semidefinite and complete on
their block at each setting; of $R^{\cg}$, that it is entrywise non-negative
and column-stochastic (Lemma~\ref{lem:descent}); and $q_t>0$. Read with $G$
in place of $K$, the same argument gives
$A^{\mathrm{ref}}\in\mathcal{D}_{\mathrm{adm}}$: $R^{\cg}$, $q_t$ and the
supports are the handed instance's and unchanged.

\emph{An instance over $\mathbb{R}$.} The rest of the proof of
Proposition~\ref{prop:flat-instance} establishes the model map, (I2), (I3)
and (I4) for every admissible datum under (H1)--(H3), and reads of a datum
only its Hermiticity, positivity and class completeness: the model map is a
direct sum of the datum and the above-cutoff blocks; the scalar values
\eqref{eq:flat-scalars} are formed from the routed weights alone; the
channels of Proposition~\ref{prop:chain} read the above-cutoff blocks, the
projectors and the weights and no entry of the datum; and the certificate of
each licensed entry is the computation of
Corollary~\ref{cor:fssinstance}, positive on the data sectors by
admissibility and on the flag sectors by (H2). (H1)--(H3) are about the
above-cutoff blocks and the weights, none of which is a block of $K$, so
they hold unchanged. Every step is a field operation, a comparison in the
Loewner order or an identity of operators, so the argument holds over the
real substrate as over any other that
Definition~\ref{def:substrate} admits, and the flat instance is a
reduction instance over $\mathbb{R}$ (Definition~\ref{def:reduction-instance}).

\emph{A non-negative combination.} The expressions of the first paragraph,
read for any family $K'$ with the properties listed, are the rule of (I1).
Their coefficients $q_tR^{\cg}_{y|y'}$ and $q_t$ are non-negative rationals:
$q_t>0$ is a declared class weight, every entry of $R^{\cg}$ is a sum of
entries of the declared $R$ (Lemma~\ref{lem:descent}), and every declared
entry and scalar lies on the declared lattice of
Definition~\ref{def:weight-route}, so is a dyadic rational; $R^{\cg}$ is
entrywise non-negative. Neither $q_t$ nor $R^{\cg}$ is formed from a block,
so the coefficients are the same for every family. Every block combined into
the element on the data sector $(m)$ is a block of photon number $m$, of the
dimension $d_m$ of that sector. That is \eqref{eq:datum-combination} of
Corollary~\ref{cor:companion-factor-enclosure}.

\emph{The target and the entries.} The target at $A^{\mathrm{ref}}$ is
\eqref{eq:instance-target} at the flat instance's sectors, which (I2) of the
proof of Proposition~\ref{prop:flat-instance} identifies with $\Fwp$ and
$\Ffss$ for the model $M[A^{\mathrm{ref}}]$; the entry list of its (I4) is
empty on the encoded route and on the licensed route consists of the
flag-sector entries of Corollary~\ref{cor:fssinstance} at the datum, whose
weights $p_o$ are the routed weights and are not formed from $K$.

\emph{The requested entry.} The companion is the instance at the
admissible datum $A^{\mathrm{ref}}$, and $\hat O^{\mathrm{ref}}$ is the
operator \eqref{eq:two-sector-at} of Corollary~\ref{cor:two-sector-decided}
at $A''=A^{\mathrm{ref}}$, with the instance's channel $\Lambda$, which is
the same at every datum. By part (ii) of that corollary it suffices that
$A^{\mathrm{ref},(m)}_e\succeq a\,\idon{d_m}$ for every $m\le N$ and
$b\le p_e$; the second is the second condition of
\eqref{eq:companion-two-sector}. Fix $m\le N$ and put $E\coloneqq A^{\mathrm{ref},(m)}_e-C^{(m)}_e$. It is
Hermitian, $A^{\mathrm{ref}}$ being an admissible datum and the centre
Hermitian by hypothesis, and $|E_{il}|\le R^{(m)}_{e,il}$ for every entry.
Let $Ex=\lambda x$ with $x\ne0$ and let $i$ be an index with
$|x_i|=\max_l|x_l|$; then
$|\lambda||x_i|=|\sum_lE_{il}x_l|\le\sum_lR^{(m)}_{e,il}|x_l|\le\rho_m|x_i|$,
so every eigenvalue of $E$ is at least $-\rho_m$ and
$E\succeq-\rho_m\idon{d_m}$. With the first condition of
\eqref{eq:companion-two-sector},
$A^{\mathrm{ref},(m)}_e=C^{(m)}_e+E\succeq(a+\rho_m)\idon{d_m}-\rho_m\idon{d_m}
=a\,\idon{d_m}$, and Corollary~\ref{cor:two-sector-decided}\,(ii) gives
$\Lambda^{\dagger}(\hat O^{\mathrm{ref}})\succeq0$, which is
\eqref{eq:certificate} with $c=0$ and $\bowtie\,=\,\succeq$.
\end{proof}

% claim reductions/wpfss/cor:normalized-flat
\claimsummary{reductions/wpfss/cor:normalized-flat}{The flat instance built on a normalized station, precomposed with the pure-loss channel of the factored loss, is a reduction instance for the receiver itself. Its emission is the normalized one with that channel in front; nothing else emitted changes.}
\claimbookmark{reductions/wpfss/cor:normalized-flat}{Corollary}{cor:normalized-flat}{The flat instance of a normalized station}
\claimleadin{%
The normalization changes the device a reduction is built on, not the
reduction: the receiver is the \exttermref{../receiver/contract.pdf}{term-normalized-station}{normalized station} preceded by a known loss
(\xref{receiver/detector-imperfections/prop:efficiency-normalization}), and
a known channel in front of a reduction instance is absorbed by
Lemma~\ref{lem:precomposed-instance}.
}
\begin{corollary}[The flat instance of a normalized station]\label{cor:normalized-flat}
Let $\Recv$ be a \exttermref{../receiver/contract.pdf}{term-receiver-instrument}{receiver instrument} without \exttermref{../receiver/contract.pdf}{term-declared-source-modes}{declared source modes}, with
efficiencies $\eta_k$ and $\eta_{\max}>0$, let $t\in[\eta_{\max},1]$, and let
$\Recv'$ be a normalized station of $\Recv$ at the
\exttermref{../receiver/contract.pdf}{term-factored-loss}{factored loss} $t$
(\xref{receiver/detector-imperfections/def:normalized-station}), both with
the channel system $\Chan$. Let the data of
Definition~\ref{def:weight-route} be supplied for $\Recv'$, with
$\Tset=\{\off\}$ and $q_{\off}=1$, pass every refusal rule of this module's
interface, and let $\mathcal{J}'_{\mathrm{flat}}$ be the flat instance of
Proposition~\ref{prop:flat-instance} at the declared route. Put
$\Phi\coloneqq\Lambda^{\mathrm{loss}}_t$, the \exttermref{../receiver/contract.pdf}{term-pure-loss-channel}{pure-loss channel} of
transmittance $t$ on $\Chan$
(\xref{receiver/detector-imperfections/def:pure-loss}), and
$\mathcal{J}_{\mathrm{flat}}\coloneqq\mathcal{J}'_{\mathrm{flat}}\circ\Phi$.
\begin{enumerate}[label=(\roman*),leftmargin=3.0em]
\item \emph{(the instance)} $\mathcal{J}_{\mathrm{flat}}$ is a reduction
  instance on $\Chan$ with the datum, sectors, scalar values, target,
  transported set, entries, constants and \termref{term-class-marginal-flag}{class-marginal flag} of
  $\mathcal{J}'_{\mathrm{flat}}$; its channel is that of
  $\mathcal{J}'_{\mathrm{flat}}$ precomposed with $\Phi$, the
  \termref{term-pull-back}{pull-backs} of its transported set are $\Phi^{\dagger}$ of
  those of $\mathcal{J}'_{\mathrm{flat}}$, and its reachable states
  (Definition~\ref{def:sound}) are those of that composite channel. Where the supplied blocks are those of
  $\Recv'$'s coarse-grained click-pattern POVM, in the bare form, or of its
  \exttermref{../receiver/contract.pdf}{term-dark-composed-family}{dark-composed family} with the declared dark-count matrix, in the
  dark-composed form, its handed model is the post-processed joint family of
  steps~1--2 of Definition~\ref{def:weight-route} formed from $\Recv$'s own
  measurement with the same dark-count matrix.
\item \emph{(the emission)} Theorem~\ref{thm:emit} holds of
  $\mathcal{J}_{\mathrm{flat}}$: its declared datum, emitted family, routed
  weights, entries and constants are those Corollary~\ref{cor:emitted-flat}
  gives for $\mathcal{J}'_{\mathrm{flat}}$, its channel is
  $\Lambda\circ\Phi$ with $\Lambda$ that of $\mathcal{J}'_{\mathrm{flat}}$,
  and for every $v\in\Vset$
  \begin{equation}\label{eq:normalized-flat}
    (\Lambda\circ\Phi)^{\dagger}(F_v)\;=\;\Phi^{\dagger}(M'_v)
    \qquad\text{exactly, on }\Chan,
  \end{equation}
  with $M'$ the declared model \eqref{eq:flat-declared} of
  $\mathcal{J}'_{\mathrm{flat}}$. Where moreover the normalizer kept every
  class, the right-hand side is the family of (i) formed from $\Recv$'s
  measurement.
\item \emph{(several measuring factors)} Where each factor $f$ of
  Theorem~\ref{thm:emitted-product} is such an instance
  $\mathcal{J}^{f}_{\mathrm{flat}}=\mathcal{J}'^{f}_{\mathrm{flat}}\circ
  \Lambda^{\mathrm{loss}}_{t_f}$, each at its own factored loss $t_f$, that
  theorem applies, with channel
  $\bigotimes_f\big(\Lambda_f\circ\Lambda^{\mathrm{loss}}_{t_f}\big)$ and every
  entry at the constant of its factor.
\item \emph{(the reference companion)} $\mathcal{J}_{\mathrm{flat}}$
  satisfies (H-REF) of Corollary~\ref{cor:reference-companion}, and its
  reference companion at a family $G$ is that of
  $\mathcal{J}'_{\mathrm{flat}}$ at $G$ precomposed with $\Phi$: the same
  target and entries, the companion's channel precomposed with $\Phi$, and
  $\Phi^{\dagger}$ of the companion's pull-backs.
\end{enumerate}
\end{corollary}
\claimreading{When the receiver is analysed at its normalized station, the finite model the key rate is computed on is exactly the one computed for the normalized station. It is a valid model of your actual receiver where the receiver facts the run carries as \hyperlink{claim-reductions/ass:recv}{A-RECV} hold and where the largest detector efficiency may be factored out as loss in front of the receiver, which the reading of \xref{receiver/detector-imperfections/prop:efficiency-normalization} conditions on the receiver's device assumptions and which an efficiency mismatch the adversary can exploit breaks. The factored loss is then counted as part of the channel the adversary controls. The certified high-photon weights, constraints and their constants are those of the normalized station.}
\claimmeta{\metalabel{Uses}\mref{Def.}{def:system-setting}, \mref{Def.}{def:sound}, \mref{Lem.}{lem:precomposed-instance}, \mref{Thm.}{thm:emit}, \mref{Thm.}{thm:emitted-product}, \mref{Cor.}{cor:reference-companion}, \mref{Def.}{def:weight-route}, \mref{Prop.}{prop:flat-instance}, \mref{Cor.}{cor:emitted-flat}, \mref{Cor.}{cor:flat-companion}, \extref{../receiver/contract.pdf}{receiver/detector-imperfections/def:dark-composed}, \extref{../receiver/contract.pdf}{receiver/detector-imperfections/def:normalized-station}, \extref{../receiver/contract.pdf}{receiver/detector-imperfections/def:pure-loss}, \extref{../receiver/contract.pdf}{receiver/detector-imperfections/prop:efficiency-normalization}, \extref{../receiver/contract.pdf}{receiver/thm:click-povm}, \extref{../receiver/contract.pdf}{receiver/thm:coarse-graining}
\metasep \metalabel{Assumes}\aref{A-RECV}{ass:recv}, \aref{A-PROBE}{ass:probe}, \aref{A-TSM}{ass:tsm}
\metasep \metalabel{Used by}\mref{Req.}{req:normalized-link-record}, \mref{Req.}{req:normalized-branch-refusal}
\metasep newly derived
\metasep implemented (effective: accepted)
\metasep \metalabel{Code}\coderef{https://github.com/fraunhofersg/certiqs-key/blob/v0.1.0/src/certiqs_key/reductions/wpfss/api.py}{certiqs_key.reductions.wpfss.api:FlatInstanceBuilder}}

\begin{proof}
$\Hs{\Chan}$ is separable (Definition~\ref{def:system-setting}) and $\Phi$ is
a channel on it, so Lemma~\ref{lem:precomposed-instance} applies to
$\mathcal{J}'_{\mathrm{flat}}$, which is a reduction instance by
Proposition~\ref{prop:flat-instance}.

(i) The first sentence is Lemma~\ref{lem:precomposed-instance}\,(ii), the
reachable states being by Definition~\ref{def:sound} the images of the
\termref{term-physical-inputs}{physical inputs} under the instance's channel. For the
second: every click-pattern element commutes with the photon number
(\xref{receiver/thm:click-povm}\,(iii)), so do the coarse-grained elements
(\xref{receiver/thm:coarse-graining}) and the dark-composed ones
(\xref{receiver/detector-imperfections/def:dark-composed}), and so does the
post-processed joint family, a non-negative combination of them; it is
therefore block diagonal for $\bigoplus_{m\le N}\Hs{\Chan_m}$, and the handed
model of $\mathcal{J}'_{\mathrm{flat}}$ is that family itself
(Proposition~\ref{prop:flat-instance}). Steps~1--2 of
Definition~\ref{def:weight-route} form each element as a fixed linear
combination, with coefficients from the dark-count matrix and $q_{\off}=1$, of
the coarse-grained click-pattern elements, which are fibre sums of the click
patterns; by
\xref{receiver/detector-imperfections/prop:efficiency-normalization}\,(ii) the
same combination formed from $\Recv$'s measurement is $\Phi^{\dagger}$ of the
one formed from $\Recv'$'s. That is the model map of
$\mathcal{J}_{\mathrm{flat}}$ at the handed datum
(Lemma~\ref{lem:precomposed-instance}\,(ii)).

(ii) $\mathcal{J}_{\mathrm{flat}}$ is a reduction instance whose handed datum,
\termref{term-class-weights}{class weights} and lattice exponent are those of
$\mathcal{J}'_{\mathrm{flat}}$, so Theorem~\ref{thm:emit} applies to it, and
by Lemma~\ref{lem:precomposed-instance}\,(iii) its declared datum, emitted
family, entries and constants are those of $\mathcal{J}'_{\mathrm{flat}}$,
which Corollary~\ref{cor:emitted-flat} describes, the routed weights being
scalar values of the target. Its emitted channel is $\Lambda\circ\Phi$ and
its declared model $\Phi^{\dagger}(M')$, so \eqref{eq:normalized-flat} is
\eqref{eq:emitidentity} for $\mathcal{J}_{\mathrm{flat}}$. Where the
normalizer kept every class, the declared datum is the handed one, $M'$ is the
handed model of $\mathcal{J}'_{\mathrm{flat}}$, and (i) identifies
$\Phi^{\dagger}(M')$.

(iii) Theorem~\ref{thm:emitted-product} asks of each factor only that it be a
reduction instance to which Theorem~\ref{thm:emit} applies, which (ii) gives,
with channel $\Lambda_f\circ\Lambda^{\mathrm{loss}}_{t_f}$ and the entries
and constants of $\mathcal{J}'^{f}_{\mathrm{flat}}$.

(iv) $\mathcal{J}'_{\mathrm{flat}}$ satisfies (H-REF)
(Corollary~\ref{cor:flat-companion}); $\mathcal{J}_{\mathrm{flat}}$ has the
same rule forming its datum and is a reduction instance over the same
substrate by Lemma~\ref{lem:precomposed-instance}\,(ii), so it satisfies
(H-REF) as well. Its reference companion at $G$ is
$\mathcal{J}_{\mathrm{flat}}$ at the datum $A^{\mathrm{ref}}$ formed from
$G$, which is $\mathcal{J}'_{\mathrm{flat}}$ at $A^{\mathrm{ref}}$
precomposed with $\Phi$ by Lemma~\ref{lem:precomposed-instance}\,(iii).
\end{proof}

\section{Requirements}\label{group-reductions/wpfss/req:one-route-per-artifact}
\sectionintro{The engine emits one route per artefact, refuses \termref{term-certified-weights}{certified weights} that total one or more, and uses and records the largest certified weight. It decides descent, absorbing letters, fibres and column-stochasticity exactly on the declared data and refuses a declaration that fails one of them. It reads per-setting weights only from a receiver version that states them, and carries the supplier's eigenvalue-slack assumption beside the settings at which it was invoked. For a receiver compiled at a factored loss below one it records the loss in front of the link's \termref{term-channel-witness}{channel witness}, and refuses such a receiver as a branch of a register instance.}

% claim reductions/wpfss/req:one-route-per-artifact
\claimsummary{reductions/wpfss/req:one-route-per-artifact}{Every artefact is emitted on exactly one route, encoded or licensed, and a request for both is refused. A caller wanting both issues two requests, since pairing one route's entries with the other's channel is unsound (\hyperlink{claim-reductions/wpfss/rem:tworoutes}{Remark~\ref*{rem:tworoutes}}).}
\claimbookmark{reductions/wpfss/req:one-route-per-artifact}{Requirement}{req:one-route-per-artifact}{One route per artefact}
\claimleadin{%
The two routes run on different channels, and the licensed entry paired with
the encoded route's channel is not sound (Remark~\ref{rem:tworoutes},
\eqref{eq:crosspairing}); an artefact that carried both would admit that
pairing.
}
\begin{requirement}[One route per artefact]
\label{req:one-route-per-artifact}
The engine emits every artefact on exactly one route: the request's
\iface{route} is \iface{encoded} or \iface{licensed}, the emitted payload is
single-valued, and a request asking for both is refused with
\iface{route_not_single_valued}. A caller wanting both issues two requests and
receives two artefacts.
\end{requirement}
\claimmeta{\metalabel{Motivated by}\mref{Rem.}{rem:tworoutes}
\metasep accepted}

% claim reductions/wpfss/req:total-weight-refusal
\claimsummary{reductions/wpfss/req:total-weight-refusal}{Certified weights whose total is one or more are refused, whether or not the space above the cutoff is trivial, and no squasher is constructed at total weight one. The construction of \hyperlink{claim-reductions/wpfss/prop:chain}{Proposition~\ref*{prop:chain}} divides by one minus the total.}
\claimbookmark{reductions/wpfss/req:total-weight-refusal}{Requirement}{req:total-weight-refusal}{Refusal of a total weight of one or more}
\claimleadin{%
On a non-trivial space above the cutoff a \termref{term-total-weight}{total weight} above one is
impossible (Lemma~\ref{lem:ple1}), but on a trivial one the certified-weight
condition is vacuous, so the test is the only thing that bounds the total
there. At total weight one the relation still holds
(Remark~\ref{rem:peqone}); the refusal is the conservative choice.
}
\begin{requirement}[Refusal of a total weight of one or more]
\label{req:total-weight-refusal}
The engine refuses \termref{term-certified-weights}{certified weights} whose \termref{term-total-weight}{total weight}
$p=\sum_vp_v$ is not below one, with \iface{total_weight_not_below_one}. It
decides this for every declaration, whether or not $\Hs{\Bsys_{\gtN}}$ is
trivial, and it constructs no squasher at $p=1$, where the construction of
Proposition~\ref{prop:chain} would divide by $1-p=0$.
\end{requirement}
\claimmeta{\metalabel{Motivated by}\mref{Lem.}{lem:ple1}, \mref{Rem.}{rem:peqone}, \mref{Prop.}{prop:chain}
\metasep accepted}

% claim reductions/wpfss/req:largest-certified-weight
\claimsummary{reductions/wpfss/req:largest-certified-weight}{The engine uses the largest weight the artefact certifies for each outcome and records the weights it used. It never silently substitutes a larger one.}
\claimbookmark{reductions/wpfss/req:largest-certified-weight}{Requirement}{req:largest-certified-weight}{Use and record of the largest certified weight}
\claimleadin{%
A smaller certified weight is valid but gives a cleaner target
(Proposition~\ref{prop:weight-monotone}, Remark~\ref{rem:weightfree}), so the
choice of weight is recorded rather than implied.
}
\begin{requirement}[Use and record of the largest certified weight]
\label{req:largest-certified-weight}
The engine uses, for each outcome, the largest weight the artefact certifies
and records the weights it used in \iface{wpfss_data}, and it never silently
substitutes a larger one.
\end{requirement}
\claimmeta{\metalabel{Motivated by}\mref{Prop.}{prop:weight-monotone}, \mref{Rem.}{rem:weightfree}
\metasep accepted}

% claim reductions/wpfss/req:descent-decided-exactly
\claimsummary{reductions/wpfss/req:descent-decided-exactly}{The declared dark-count matrix is decided exactly to be monotone and to descend through the declared coarse-graining, by testing descent itself; a declaration that fails either is refused. The canonical sufficient condition serves only as a shortcut checked in full.}
\claimbookmark{reductions/wpfss/req:descent-decided-exactly}{Requirement}{req:descent-decided-exactly}{Exact decision of monotonicity and descent}
\claimleadin{%
The descent condition is an identity between declared dyadic rationals, so
it can be decided exactly. The sufficient condition \eqref{eq:canonicalcg}
with its constancy clause left out would admit a declaration whose descended
matrix is not defined (Remark~\ref{rem:descentneedsconstant}), so a shortcut
through that condition is checked in full. Its monotonicity half is clause
(D2) of the supplier's \exttermref{../receiver/contract.pdf}{term-dark-count-channel}{dark-count channel}, which every admitted matrix
therefore satisfies exactly (Remark~\ref{rem:darkcompose}).
}
\begin{requirement}[Exact decision of monotonicity and descent]\label{req:descent-decided-exactly}
The engine decides exactly, on the declared dark-count matrix $R$, that $R$
is monotone, $R_{S|S'}>0\Rightarrow S'\subseteq S$, and that it descends
through the \exttermref{../receiver/contract.pdf}{term-declared-coarse-graining}{declared coarse-graining} $\cg$ in the sense of
Lemma~\ref{lem:descent}, $P(y|S')=P(y|S'')$ for every $y$ and every pair
$S',S''$ with $\cg(S')=\cg(S'')$, each an identity or a comparison between the
dyadic rationals the declaration carries (Definition~\ref{def:substrate}). It
tests the descent property itself, and treats the sufficient condition
\eqref{eq:canonicalcg} only as a shortcut that must itself be checked in
full. A declaration that fails either condition is refused with the reason
\iface{dark_counts_do_not_descend}.
\end{requirement}
\claimmeta{\metalabel{Motivated by}\mref{Def.}{def:substrate}, \mref{Lem.}{lem:descent}, \mref{Rem.}{rem:descentneedsconstant}, \mref{Rem.}{rem:darkcompose}
\metasep accepted}

% claim reductions/wpfss/req:absorbing-bounded-letters
\claimsummary{reductions/wpfss/req:absorbing-bounded-letters}{Every bounded letter of the slot must be absorbing under the descended dark-count matrix, decided exactly on input. A declaration with one non-absorbing letter is refused as a whole; no entry's weight is zeroed in its place.}
\claimbookmark{reductions/wpfss/req:absorbing-bounded-letters}{Requirement}{req:absorbing-bounded-letters}{Absorbing bounded letters as an input validity condition}
\claimleadin{%
Without \eqref{eq:absorbing} the routed bound \eqref{eq:routebound} is no
longer proved and fails in general, in the unsafe direction
(Remark~\ref{rem:absorbneeded}), and no bound is stated for a letter that
fails it.
}
\begin{requirement}[Absorbing bounded letters as an input validity condition]\label{req:absorbing-bounded-letters}
The engine takes \eqref{eq:absorbing} as a declared input validity condition:
it decides $R^{\cg}_{y_i|y_i}=1$ exactly at every bounded letter $y_i$ of the
slot, and refuses a declaration in which some bounded letter fails it, with
the reason \iface{bounded_letter_not_absorbing}. The refusal is of the whole
artifact: the failing entry's weight is not set to zero and the rest kept.
\end{requirement}
\claimmeta{\metalabel{Motivated by}\mref{Prop.}{prop:route}, \mref{Rem.}{rem:absorbneeded}
\metasep accepted}

% claim reductions/wpfss/req:fibre-equality
\claimsummary{reductions/wpfss/req:fibre-equality}{Every certified set of the slot must be exactly one fibre of the declared coarse-graining; the equality is decided on input and a slot that fails it is refused. A split set leaves neither letter with a bound.}
\claimbookmark{reductions/wpfss/req:fibre-equality}{Requirement}{req:fibre-equality}{Certified sets equal to fibres of the coarse-graining}
\claimleadin{%
A bound on a sum of positive operators says nothing about its summands, so a
certified set split over two letters leaves neither with a bound, and one
merged with other outcomes attaches the weight to a larger fibre
(Remark~\ref{rem:fibre}).
}
\begin{requirement}[Certified sets equal to fibres of the coarse-graining]\label{req:fibre-equality}
The engine requires every bounded set of the slot to be exactly a fibre of
the \exttermref{../receiver/contract.pdf}{term-declared-coarse-graining}{declared coarse-graining}, $\mathcal{O}_i=\cg^{-1}(y_i)$, validates
the equality at input, and refuses a slot that fails it with the reason
\iface{bound_set_not_a_fibre}.
\end{requirement}
\claimmeta{\metalabel{Motivated by}\mref{Rem.}{rem:fibre}
\metasep accepted}

% claim reductions/wpfss/req:column-stochastic-refusal
\claimsummary{reductions/wpfss/req:column-stochastic-refusal}{The declared dark-count matrix is decided exactly to have non-negative entries and unit column sums, and is refused otherwise. No second rule decides the descended matrix, which inherits the property exactly.}
\claimbookmark{reductions/wpfss/req:column-stochastic-refusal}{Requirement}{req:column-stochastic-refusal}{Exact column-stochasticity of the declared dark-count matrix}
\claimleadin{%
Every result of this contract is proved for an exactly column-stochastic $R$,
and the descended matrix inherits the property by regrouping one finite sum
(Remark~\ref{rem:colstoch}).
}
\begin{requirement}[Exact column-stochasticity of the declared dark-count matrix]\label{req:column-stochastic-refusal}
The engine decides \emph{exactly} that the declared dark-count matrix $R$ is
column-stochastic, $R_{S|S'}\ge0$ for every entry and $\sum_SR_{S|S'}=1$ for
every column, and refuses a declaration that is not, with the reason
\iface{dark_counts_not_column_stochastic}. No second rule decides the
property of the descended matrix $R^{\cg}$ of Lemma~\ref{lem:descent}.
\end{requirement}
\claimmeta{\metalabel{Motivated by}\mref{Lem.}{lem:descent}, \mref{Rem.}{rem:colstoch}
\metasep accepted}

% claim reductions/wpfss/req:supplier-version
\claimsummary{reductions/wpfss/req:supplier-version}{Per-setting weights are read only from an artifact whose receiver version states them as the content of a claim. An artifact of any other version is refused, a probe-free one included, and nothing it reports is read per setting.}
\claimbookmark{reductions/wpfss/req:supplier-version}{Requirement}{req:supplier-version}{Receiver versions from which per-setting weights are read}
\claimleadin{%
A hypothesis of Proposition~\ref{prop:route}, whose part (iv) discharges the
per-setting guarantee by \xref{receiver/thm:emitted-bound}, is discharged
only by a claim that states it. At a receiver
version whose claims state only the minimum over the settings, the
per-setting values an artifact reports beside it are the content of no
claim, and the \termref{term-weight-route}{weight route} reads every slot per setting
(Definition~\ref{def:weight-route}), so such an artifact supplies nothing it
can read.
}
\begin{requirement}[Receiver versions from which per-setting weights are read]\label{req:supplier-version}
The engine reads the \exttermref{../receiver/contract.pdf}{term-per-setting-weight}{per-setting weights}
$p^{(t)}_i$ of a receiver artifact (\iface{high_photon_bounds}) only where
the artifact's module version resolves to an accepted receiver version that
states them: one whose \xref{receiver/thm:emitted-bound} carries the
per-setting inequality as its first conclusion and whose
\xref{receiver/def:slot} carries a weight per entry and setting. An artifact
whose version does not state them, a probe-free one included, is refused with
\iface{high_photon_slot_ill_formed}: the per-setting values it reports beside
the minimum are not read from it, and its minimum is not read at every
setting in their place.
\end{requirement}
\claimmeta{\metalabel{Motivated by}\mref{Def.}{def:weight-route}, \mref{Prop.}{prop:route}, \extref{../receiver/contract.pdf}{receiver/def:slot}, \extref{../receiver/contract.pdf}{receiver/thm:emitted-bound}
\metasep accepted}

% claim reductions/wpfss/req:supplier-eig-record
\claimsummary{reductions/wpfss/req:supplier-eig-record}{Where the emitted ledger carries \hyperlink{claim-reductions/ass:recv}{A-RECV}, it carries the supplier's \hyperref{../receiver/contract.pdf}{}{claim-receiver/ass:eig}{A-EIG} there beside the settings at which the supplier says it was invoked. It is the one supplier record carried beside its key; every other key is carried alone.}
\claimbookmark{reductions/wpfss/req:supplier-eig-record}{Requirement}{req:supplier-eig-record}{Invocation record of the eigenvalue-slack assumption in the ledger}
\begin{technicalclaim}
\begin{requirement}[Invocation record of the eigenvalue-slack assumption in the ledger]\label{req:supplier-eig-record}
On the row of Assumption~\ref{ass:recv}, wherever the emitted assumption
ledger \iface{assumption_ledger} carries it
(Requirement~\ref{req:recv-row-provenance}), the engine records the
supplier's eigenvalue-slack assumption A-EIG (\xref{receiver/ass:eig})
beside the settings at which the supplier says it was invoked. Of the supplier's record behind a key, only A-EIG's is carried
beside it; every other key is carried alone, and its record stays in the
artifact, which the entry names by version.
\end{requirement}
\end{technicalclaim}
\claimmeta{\metalabel{Motivated by}\mref{Ass.}{ass:recv}, \mref{Req.}{req:recv-row-provenance}, \mref{Rem.}{rem:receiver-imports}, \extref{../receiver/contract.pdf}{receiver/ass:eig}
\metasep accepted}

% claim reductions/wpfss/req:normalized-link-record
\claimsummary{reductions/wpfss/req:normalized-link-record}{At a factored loss below one, four slots of the emitted link change: the witness and the carried pull-backs gain the pure-loss channel in front, the source family is the receiver's, the reachable set the composite's. Target, weights, entries and constants do not change.}
\claimbookmark{reductions/wpfss/req:normalized-link-record}{Requirement}{req:normalized-link-record}{The link of a normalized station records the factored loss}
\claimleadin{%
Corollary~\ref{cor:normalized-flat} makes the precomposed instance, and not
the flat instance of the \exttermref{../receiver/contract.pdf}{term-normalized-station}{normalized station}, the reduction of the receiver.
By Lemma~\ref{lem:precomposed-instance}\,(ii) the two links of
Definition~\ref{def:link} differ in four slots: the \termref{term-channel-witness}{channel witness},
the \termref{term-source-family}{source family}, the pull-backs of the carried tuple and the
\termref{term-reachable-set-statement}{reachable-set statement}. Every slot is part of what a link carries to
be re-checked (Requirement~\ref{req:channel-witness}), so the record states
all four.
}
\begin{requirement}[The link of a normalized station records the factored loss]\label{req:normalized-link-record}
Where the receiver artefact's \iface{efficiency_record} carries a
\exttermref{../receiver/contract.pdf}{term-factored-loss}{factored loss} $t<1$, the engine emits the artefact of the precomposed
instance $\mathcal{J}'_{\mathrm{flat}}\circ\Lambda^{\mathrm{loss}}_t$ of
Corollary~\ref{cor:normalized-flat}, and its link records, in place of the
normalized station's: (i) the channel witness as the flat instance's witness
preceded by the \exttermref{../receiver/contract.pdf}{term-pure-loss-channel}{pure-loss channel} of transmittance $t$ on the
\exttermref{../receiver/contract.pdf}{term-channel-modes}{channel modes}, with $t$ recorded under \iface{factored_loss}; (ii) the
source family as the receiver's measurement; (iii) each pull-back of the
carried tuple as the pull-back through that loss of the flat instance's
pull-back, not as the flat instance's pull-back itself; and (iv) the
reachable-set statement as the one for that composite channel
(Definition~\ref{def:sound}). The emitted target, \termref{term-routed-weights}{routed weights},
transported operators, entries, constants, certificate cases and
\termref{term-licence-record}{licence record} are those emitted for the normalized station, unchanged. The reference
companion's link is recorded the same way: its channel and pull-backs carry
the same loss in front, and its target and entries are unchanged.
\end{requirement}
\claimmeta{\metalabel{Motivated by}\mref{Def.}{def:sound}, \mref{Def.}{def:link}, \mref{Lem.}{lem:precomposed-instance}, \mref{Req.}{req:channel-witness}, \mref{Cor.}{cor:emitted-flat}, \mref{Cor.}{cor:normalized-flat}
\metasep implemented (effective: accepted)
\metasep \metalabel{Code}\coderef{https://github.com/fraunhofersg/certiqs-key/blob/v0.1.0/src/certiqs_key/reductions/wpfss/construction.py}{certiqs_key.reductions.wpfss.construction:normalized_precomposition}}

% claim reductions/wpfss/req:normalized-branch-refusal
\claimsummary{reductions/wpfss/req:normalized-branch-refusal}{An artefact that embeds the flat instance as a branch of a branching instance, the register instance included, reads only a receiver artefact compiled without a factored loss, and refuses one compiled at a factor below one.}
\claimbookmark{reductions/wpfss/req:normalized-branch-refusal}{Requirement}{req:normalized-branch-refusal}{A normalized receiver is not read as a branch of a branching instance}
\claimleadin{%
Corollary~\ref{cor:normalized-flat} absorbs the factored loss for the flat
instance on the channel system. A branching instance
(Theorem~\ref{thm:instrument-instance}), of which the register instance of
Corollary~\ref{cor:register-instance} is one, places the flat instance beside
other branches, which need not be the pull-back of anything through a loss on
the channel modes, so no claim here moves the loss in front of it.
}
\begin{requirement}[A normalized receiver is not read as a branch of a branching instance]\label{req:normalized-branch-refusal}
An artefact that embeds the flat instance as a branch of a branching
instance, the register instance included, reads only a receiver artefact whose \iface{efficiency_record} carries the
factored loss $t=1$, and refuses one with $t<1$ with
\iface{normalized_receiver_in_branch}.
\end{requirement}
\claimmeta{\metalabel{Motivated by}\mref{Thm.}{thm:instrument-instance}, \mref{Cor.}{cor:register-instance}, \mref{Cor.}{cor:normalized-flat}
\metasep implemented (effective: accepted)
\metasep \metalabel{Code}\coderef{https://github.com/fraunhofersg/certiqs-key/blob/v0.1.0/src/certiqs_key/reductions/wpfss/validation.py}{certiqs_key.reductions.wpfss.validation:rule_normalized_branch}}

\clearpage
\part{Domination links}
\sectionintro{This submodule relates two measurement families not by the size of their difference but by domination: one family dominates the other at one weight per sector. Its main result, \hyperlink{claim-reductions/domination/thm:noise-channel}{Theorem~\ref*{thm:noise-channel}}, derives from the branching calculus that the dominating family is then a noise channel followed by the dominated one together with classical flags, the channel keeping each sector with probability one minus its weight. That relation is a typed link of its own, the domination link, placed on the target system after a squashing link, and a squashing link followed by it is an exact chain whose security transfers with the same \ensuremath{\varepsilon} (\hyperlink{claim-reductions/domination/cor:dom-chain}{Corollary~\ref*{cor:dom-chain}}). No requirement of the engine uses these claims: a snapped measurement is not related to its model by domination, and a run evaluates its key length on the model instead where it holds one. They are kept for a declared device weight (A-DOM), and the remark on rotated projective elements says why an idealised model needs its noise floor before it dominates a lattice family at a small weight.}

\section{The domination link and its noise channel}\label{group-reductions/domination/def:domination-link}
\sectionintro{A \termref{term-domination-link}{domination link} runs from a finite \termref{term-target-family}{target family} to a family it dominates at one weight per sector, its flags prepared into sectors added for the purpose or into \termref{term-indicator-sectors}{indicator sectors} the target already has. The main result, \hyperlink{claim-reductions/domination/thm:noise-channel}{Theorem~\ref*{thm:noise-channel}}, derives its channel from the \termref{term-generalized-instrument}{generalized instrument} that keeps or flags each sector and from the common-target rule. The closing remark says why a rotated projective element dominates no lattice element, so that a model needs its noise floor before domination can relate a lattice family to it at a small weight.}
\sectionimports{\extref{../foundations/contract.pdf}{foundations/def:domination}, \extref{../foundations/contract.pdf}{foundations/lem:absorbing-domination}, \extref{../foundations/contract.pdf}{foundations/lem:domination}, \extref{../foundations/contract.pdf}{foundations/lem:rational-obstruction}}

% claim reductions/domination/def:domination-link
\claimsummary{reductions/domination/def:domination-link}{A domination link runs from a target family G to a family F it dominates at one weight per sector, through a noise channel that keeps each sector with probability one minus its weight and otherwise measures a residual family and prepares a flag, in added sectors or in sectors F already has.}
\claimbookmark{reductions/domination/def:domination-link}{Definition}{def:domination-link}{The domination link}
\claimleadin{%
A family $G$ that dominates a family $F$ at weights below one is a noise
channel followed by $F$, together with classical flags
\cite[Thm.~3, Eqs.~(20)--(24)]{nahar2025imperfect}: by
\xref{foundations/def:domination} each element splits as
$G_v=(1-\vartheta_s)F_v+\vartheta_sR_v$ on sector $s$, with a residual family $R$
that is again complete, and a channel can keep the sector with probability
$1-\vartheta_s$ and otherwise measure $R$ and prepare a flag
(Theorem~\ref{thm:noise-channel}, which also shows that the tuple defined
here is a receiver-stage link with an exact identity). The link defined here
records that channel. It is placed on the \emph{target} side of a chain: it runs from the
\termref{term-target-family}{target family} a squashing link produces at the exact datum to the target
family at a datum it dominates, both on one finite target system with
one \exttermref{../foundations/contract.pdf}{term-sector-map}{sector map}. The relation could equally be placed on the channel
system before the squasher, from the model at the exact datum to the model at
the other datum together with flags; the composite would be the same exact
relation, but the intermediate descriptor would then be the
infinite-dimensional channel system with flags, whereas on the target side
every descriptor of the chain after the first link is the finite target
system with its sector map (Corollary~\ref{cor:dom-chain}). The flags are
prepared either into $|\Vset|$ one-dimensional sectors added to the target
for the purpose, or, where the target already has one one-dimensional sector
per outcome carrying that outcome's indicator, as the \termref{term-flag-state-target}{flag-state
target} does, into those sectors, and the target system does not change
\cite[Technical aside~1]{nahar2025imperfect}.
}
\begin{definition}[The domination link]\label{def:domination-link}
Let $\SecSet$ be a finite ordered set of sectors with dimensions $d_s\ge1$,
$\Hs{\Qsys}=\bigoplus_{s\in\SecSet}\mathbb{C}^{d_s}$ with sector projectors
$\Pi_s$, and let $G=\{G_v\}_{v\in\Vset}$ and $F=\{F_v\}_{v\in\Vset}$ be POVMs
on $\Qsys$ over one finite alphabet $\Vset$, both block diagonal for
$\{\Pi_s\}$, $G_v=\bigoplus_sG^{(s)}_v$ and $F_v=\bigoplus_sF^{(s)}_v$. Let
$\vartheta=(\vartheta_s)_{s\in\SecSet}$ with $\vartheta_s\in[0,1)$ be weights at which
$G$ dominates $F$, every block complete at $1$
(\xref{foundations/def:domination}):
\begin{equation}\label{eq:dom-link}
  G^{(s)}_v-(1-\vartheta_s)\,F^{(s)}_v\;\succeq\;0
  \qquad\text{for every }v\in\Vset\text{ and }s\in\SecSet ;
\end{equation}
they are the
\termdef{term-domination-weights}{domination weights}
of the pair. Let $R_v\coloneqq\bigoplus_sR^{(s)}_v$ with $R^{(s)}$ the
\exttermref{../foundations/contract.pdf}{term-residual-family}{residual family} of \xref{foundations/def:domination} on sector $s$.
A family $F$ on $\Qsys$ has
\termdef{term-indicator-sectors}{indicator sectors}
when there are pairwise distinct sectors $(f)\in\SecSet$, one for each
$f\in\Vset$, with $d_{(f)}=1$ and $F^{((f))}_v=\delta_{vf}$ for every $v$;
write $\Pi_{(f)}$ for the projector onto the sector $(f)$, a rank-one
projector.
\begin{enumerate}[label=(\alph*),leftmargin=2.4em]
\item \emph{(added flag sectors)} Let $\Qsys^{+}$ be the system with
  $\Hs{\Qsys^{+}}=\Hs{\Qsys}\dsum\Hs{\Flag^{+}}$, $\Hs{\Flag^{+}}=\mathbb{C}^{|\Vset|}$
  with orthonormal basis $\{\ket{v}\}_{v\in\Vset}$, whose \termref{term-sector-set}{sector set} is
  $\SecSet$ followed by $|\Vset|$ one-dimensional sectors $(f)^{+}$, $f\in\Vset$,
  in the order of $\Vset$, the sector $(f)^{+}$ spanned by $\ket{f}$. Put
  $F^{+}_v\coloneqq F_v\dsum\dyad{v}_{\Flag^{+}}$, so that the added sectors are
  indicator sectors of $F^{+}$, and
  \begin{equation}\label{eq:noise-fresh}
    \Phi^{+}_{\vartheta}(\sigma)\;\coloneqq\;
    \sum_{s\in\SecSet}(1-\vartheta_s)\,\Pi_s\sigma\Pi_s\ \dsum\
    \sum_{v\in\Vset}\Big(\sum_{s\in\SecSet}\vartheta_s
      \tr\big[R_v\,\Pi_s\sigma\Pi_s\big]\Big)\dyad{v}_{\Flag^{+}}
  \end{equation}
  for $\sigma\in\Tcl{\Qsys}$.
\item \emph{(shared flag sectors)} Suppose $F$ has indicator sectors $(f)$ and
  $\vartheta_{(f)}=0$ for every $f\in\Vset$. Put
  \begin{equation}\label{eq:noise-shared}
    \Phi_{\vartheta}(\sigma)\;\coloneqq\;
    \sum_{s\in\SecSet}(1-\vartheta_s)\,\Pi_s\sigma\Pi_s\ +\
    \sum_{f\in\Vset}\Big(\sum_{s\in\SecSet}\vartheta_s
      \tr\big[R_f\,\Pi_s\sigma\Pi_s\big]\Big)\Pi_{(f)} .
  \end{equation}
\end{enumerate}
A
\termdef{term-domination-link}{domination link}
is the tuple \eqref{eq:link} of a receiver-stage \termref{term-typed-link}{typed link}
(Definition~\ref{def:link}) with $\textsf{stage}=\textsf{receiver}$, kind
$\textsf{family-dominating}$ and \termref{term-source-family}{source family} $(\Qsys,G)$; with
target $(\Qsys^{+},F^{+})$ and \termref{term-channel-witness}{channel witness}
$\Phi^{+}_{\vartheta}$ in case (a), or target $(\Qsys,F)$ and channel witness
$\Phi_{\vartheta}$ in case (b); with \termref{term-carried-operators}{carried operators} the projectors
$\Pi_s\dsum0$, $s\in\SecSet$, and the flag projector
$\Pflag^{+}\coloneqq0\dsum\idon{\Flag^{+}}$ in case (a), and the projectors
$\Pi_s$ of the sectors that are not indicator sectors and the flag projector
$\Pflag\coloneqq\sum_{f}\Pi_{(f)}$ in case (b), each with its pull-back under
the channel witness; with the \termref{term-reachable-set-statement}{reachable-set statement} of
Definition~\ref{def:sound} for that channel; with the constraint entries and
\termref{term-licence-record}{licence record} of the chain it ends; and with
$\textsf{status}=\textsf{exact}$.
\end{definition}
\claimmeta{\metalabel{Uses}\mref{Def.}{def:squashing}, \mref{Def.}{def:sound}, \mref{Def.}{def:link}, \mref{Def.}{def:reduction-instance}, \extref{../foundations/contract.pdf}{foundations/def:domination}
\metasep \metalabel{Used by}\mref{Thm.}{thm:noise-channel}, \mref{Rem.}{rem:rotated-projective}, \mref{Cor.}{cor:dom-chain}
\metasep newly derived
\metasep accepted}

% claim reductions/domination/thm:noise-channel
\claimsummary{reductions/domination/thm:noise-channel}{If G dominates F sector by sector, G squashes exactly to F with added flag sectors, or to F itself where F has flag sectors, through a noise channel that keeps each sector with probability one minus its weight. The adjoint returns each sector projector scaled by that probability.}
\claimbookmark{reductions/domination/thm:noise-channel}{Theorem}{thm:noise-channel}{The noise channel of a domination}
\claimleadin{%
The channel is the one of \cite[Thm.~3, Eqs.~(20)--(24)]{nahar2025imperfect},
derived here from the branching calculus of this contract: a \termref{term-generalized-instrument}{generalized
instrument} decides, sector by sector, whether to keep the input
(probability $1-\vartheta_s$) or to flag it (probability $\vartheta_s$); the keep
branch passes the input on unchanged, and the flag branch measures the
residual family and prepares the flag of its outcome. The domination
$G=(1-\vartheta)F+\vartheta R$ sector by sector is exactly the statement that the
two branches together reproduce $G$. Remark~\ref{rem:dominating} reads the
same split as a dominated detector. Part~(iii), in which the flags land in
sectors the target already has, is the use of an existing flag space that
\cite[Technical aside~1]{nahar2025imperfect} describes.
}
\begin{mainclaim}
\begin{theorem}[The noise channel of a domination]\label{thm:noise-channel}
In the setting of Definition~\ref{def:domination-link}, let $G$ dominate $F$
at the \termref{term-domination-weights}{domination weights} $\vartheta$ and let $R$ be the residual family.
\begin{enumerate}[label=(\roman*),leftmargin=3.0em]
\item $R=\{R_v\}_{v\in\Vset}$ is a POVM on $\Qsys$, block diagonal for
  $\{\Pi_s\}$, and
  $G^{(s)}_v=(1-\vartheta_s)F^{(s)}_v+\vartheta_sR^{(s)}_v$ for every $v$ and $s$.
\item \emph{(added flag sectors)} $\Phi^{+}_{\vartheta}$ of \eqref{eq:noise-fresh} is a
  channel $\Tcl{\Qsys}\to\Tcl{\Qsys^{+}}$, and
  \begin{equation}\label{eq:noise-fresh-relation}
    \big(G;\,(1-\vartheta_s)\Pi_s\,(s\in\SecSet),\,\textstyle\sum_s\vartheta_s\Pi_s\big)
    \ \dsq_{\Phi^{+}_{\vartheta}}\
    \big(F^{+};\,\Pi_s\dsum0\,(s\in\SecSet),\,0\dsum\idon{\Flag^{+}}\big).
  \end{equation}
\item \emph{(shared flag sectors)} If $F$ has \termref{term-indicator-sectors}{indicator sectors} $(f)$ with
  $\vartheta_{(f)}=0$ for every $f$, then
  $\Phi_{\vartheta}$ of \eqref{eq:noise-shared} is a channel
  $\Tcl{\Qsys}\to\Tcl{\Qsys}$ with $G\dsq_{\Phi_{\vartheta}}F$, and its adjoint
  takes the sector projectors and the flag projector to
  \begin{equation}\label{eq:noise-shared-pullbacks}
  \begin{gathered}
    \Phi_{\vartheta}^{\dagger}(\Pi_s)=(1-\vartheta_s)\Pi_s\quad
    (s\text{ not an indicator sector}),\\
    \Phi_{\vartheta}^{\dagger}(\Pi_{(f)})=\Pi_{(f)}+\sum_{s}\vartheta_s\,
      \Pi_sR_f\Pi_s,\qquad
    \Phi_{\vartheta}^{\dagger}(\Pflag)=\Pflag+\sum_{s}\vartheta_s\Pi_s .
  \end{gathered}
  \end{equation}
\item \emph{(output)} The outputs of both channels are block diagonal for the
  sectors of their target, so either channel is unchanged when the sector
  pinching of its target is applied after it.
\item \emph{(weight zero)} If $\vartheta_s=0$ for every $s$, then $F=G$, and
  $\Phi_{\vartheta}$, where defined, is the sector pinching
  $\sigma\mapsto\sum_s\Pi_s\sigma\Pi_s$.
\end{enumerate}
In particular the domination link of Definition~\ref{def:domination-link}, in
either case, is a receiver-stage link of Definition~\ref{def:link}: its
channel witness relates its source family and carried operators to its
target and carried operators exactly, as \eqref{eq:link} requires with
$\textsf{status}=\textsf{exact}$.
\end{theorem}
\end{mainclaim}
\claimmeta{\metalabel{Uses}\mref{Def.}{def:squashing}, \mref{Lem.}{lem:top}, \mref{Def.}{def:generalized-instrument}, \mref{Lem.}{lem:master}, \mref{Prop.}{prop:common}, \mref{Def.}{def:link}, \mref{Def.}{def:domination-link}, \extref{../foundations/contract.pdf}{foundations/def:domination}, \extref{../foundations/contract.pdf}{foundations/lem:domination}
\metasep \metalabel{Used by}\mref{Cor.}{cor:dom-chain}
\metasep newly derived
\metasep accepted}

\begin{proof}
\emph{Idea.} The two channels are the composite channels of
Lemma~\ref{lem:master} and Proposition~\ref{prop:common} for one generalized
instrument with a keep branch and a flag branch.

\emph{(i)} On every sector $s$, $G^{(s)}$ and $F^{(s)}$ are Hermitian, positive
semidefinite and complete at $1$, $G$ and $F$ being POVMs block diagonal for
$\{\Pi_s\}$. \xref{foundations/lem:domination}\,(vi) gives that $R^{(s)}$ has
positive semidefinite elements, is complete at $1$ and satisfies the stated
split; so $R_v=\bigoplus_sR^{(s)}_v\succeq0$ and
$\sum_vR_v=\bigoplus_s\idon{d_s}=\idon{\Qsys}$.

\emph{The instrument.} Put
$\Ins_{\mathrm{k}}(\sigma)\coloneqq\sum_s(1-\vartheta_s)\Pi_s\sigma\Pi_s$ and
$\Ins_{\mathrm{f}}(\sigma)\coloneqq\sum_s\vartheta_s\Pi_s\sigma\Pi_s$, both maps
$\Tcl{\Qsys}\to\Tcl{\Qsys}$. Each is completely positive, a sum of maps
$\sigma\mapsto c\,\Pi_s\sigma\Pi_s$ with $c\ge0$, and
$\tr\Ins_{\mathrm{k}}(\sigma)+\tr\Ins_{\mathrm{f}}(\sigma)
=\sum_s\tr[\Pi_s\sigma\Pi_s]=\tr\sigma$, since $\sum_s\Pi_s=\idon{\Qsys}$.
So $\{\Ins_{\mathrm{k}},\Ins_{\mathrm{f}}\}$ is a generalized instrument
(Definition~\ref{def:generalized-instrument}). Its adjoints are
$\Ins_{\mathrm{k}}^{\dagger}(X)=\sum_s(1-\vartheta_s)\Pi_sX\Pi_s$ and
$\Ins_{\mathrm{f}}^{\dagger}(X)=\sum_s\vartheta_s\Pi_sX\Pi_s$. In the keep branch
take the POVM $F$ on $\Qsys$ with the identity channel, so
$F\dsq_{\idmap}F$; in the flag branch take the POVM $R$ on $\Qsys$ with the
measure-and-prepare channel
$\Lmp(\tau)=\sum_v\tr[R_v\tau]\dyad{v}_{\Flag^{+}}$, so $R\dsq_{\Lmp}\Ppvm$ by
Lemma~\ref{lem:top}. By Lemma~\ref{lem:master}\,(i) the branch elements are
$M_{v,\mathrm{k}}=\Ins_{\mathrm{k}}^{\dagger}(F_v)=\sum_s(1-\vartheta_s)\Pi_sF_v\Pi_s$
and $M_{v,\mathrm{f}}=\Ins_{\mathrm{f}}^{\dagger}(R_v)=\sum_s\vartheta_s\Pi_sR_v\Pi_s$,
and since $F_v$ and $R_v$ are block diagonal,
$M_{v,\mathrm{k}}+M_{v,\mathrm{f}}
=\bigoplus_s\big((1-\vartheta_s)F^{(s)}_v+\vartheta_sR^{(s)}_v\big)=G_v$ by (i).

\emph{(ii)} Lemma~\ref{lem:master}\,(iii) gives
$\{M_{v,\mathrm{k}}+M_{v,\mathrm{f}}\}_v\dsq_{\Lambda}\{F_v\dsum\dyad{v}\}_v$
with $\Lambda(\sigma)=\Ins_{\mathrm{k}}(\sigma)\dsum\Lmp\big(\Ins_{\mathrm{f}}(\sigma)\big)$,
a channel by Lemma~\ref{lem:master}\,(ii). The left side is $G$, the right side
is $F^{+}$, and
$\Lmp(\Ins_{\mathrm{f}}(\sigma))=\sum_v\sum_s\vartheta_s\tr[R_v\Pi_s\sigma\Pi_s]\dyad{v}$,
so $\Lambda=\Phi^{+}_{\vartheta}$ of \eqref{eq:noise-fresh}. For the carried
operators, \eqref{eq:masteradjoint} gives
$(\Phi^{+}_{\vartheta})^{\dagger}(X\dsum Y)=\Ins_{\mathrm{k}}^{\dagger}(X)
+\Ins_{\mathrm{f}}^{\dagger}\big(\Lmp^{\dagger}(Y)\big)$, with
$\Lmp^{\dagger}(Y)=\sum_v\bra{v}Y\ket{v}R_v$. At $X=\Pi_s$, $Y=0$ this is
$(1-\vartheta_s)\Pi_s$; at $X=0$, $Y=\idon{\Flag^{+}}$ it is
$\Ins_{\mathrm{f}}^{\dagger}(\sum_vR_v)=\Ins_{\mathrm{f}}^{\dagger}(\idon{\Qsys})
=\sum_s\vartheta_s\Pi_s$, which is Lemma~\ref{lem:master}\,(iv) for the flag
branch. That is \eqref{eq:noise-fresh-relation} in the notation of
Definition~\ref{def:squashing}.

\emph{(iii)} Let $W\colon\mathbb{C}^{|\Vset|}\to\Hs{\Qsys}$ be the isometry with
$W\ket{f}$ the unit vector spanning the indicator sector $(f)$, and
$\Psi(\tau)\coloneqq W\tau W^{\dagger}$, a channel
$\Tcl{\Flag^{+}}\to\Tcl{\Qsys}$ with adjoint $\Psi^{\dagger}(X)=W^{\dagger}XW$.
Because every $F_v$ is block diagonal and equals $\delta_{vf}$ on the sector
$(f)$, $\Psi^{\dagger}(F_v)=\sum_f\delta_{vf}\dyad{f}=\dyad{v}$, so
$\Ppvm\dsq_{\Psi}F$. With $F\dsq_{\idmap}F$ in the keep branch, $F$ is a
common target of the two branch targets $F$ and $\Ppvm$, and
Proposition~\ref{prop:common} gives
$\{M_{v,\mathrm{k}}+M_{v,\mathrm{f}}\}_v\dsq_{\Lambda'}F$ with
$\Lambda'=\idmap\circ\idmap\circ\Ins_{\mathrm{k}}+\Psi\circ\Lmp\circ\Ins_{\mathrm{f}}$,
a channel by that proposition. The left side is $G$, and
$\Psi(\Lmp(\Ins_{\mathrm{f}}(\sigma)))=\sum_f\sum_s\vartheta_s\tr[R_f\Pi_s\sigma\Pi_s]\,W\dyad{f}W^{\dagger}$
with $W\dyad{f}W^{\dagger}=\Pi_{(f)}$, so $\Lambda'=\Phi_{\vartheta}$ of
\eqref{eq:noise-shared}. Its adjoint is
$\Phi_{\vartheta}^{\dagger}(X)=\Ins_{\mathrm{k}}^{\dagger}(X)
+\Ins_{\mathrm{f}}^{\dagger}\big(\sum_f\bra{f}W^{\dagger}XW\ket{f}R_f\big)$,
by linearity and the adjoints of the three maps. For a sector $s$ that is not
an indicator sector, $W^{\dagger}\Pi_sW=0$, so
$\Phi_{\vartheta}^{\dagger}(\Pi_s)=(1-\vartheta_s)\Pi_s$. For an indicator sector,
$\vartheta_{(f)}=0$ gives $\Ins_{\mathrm{k}}^{\dagger}(\Pi_{(f)})=\Pi_{(f)}$, and
$W^{\dagger}\Pi_{(f)}W=\dyad{f}$ gives the second term
$\Ins_{\mathrm{f}}^{\dagger}(R_f)=\sum_s\vartheta_s\Pi_sR_f\Pi_s$. Summing over
$f$, with $\sum_fR_f=\idon{\Qsys}$ by (i),
$\Phi_{\vartheta}^{\dagger}(\Pflag)=\Pflag+\sum_s\vartheta_s\Pi_s$.

\emph{(iv)} Every term of \eqref{eq:noise-fresh} and \eqref{eq:noise-shared}
is $\Pi_s\sigma\Pi_s$ for a sector $s$, a multiple of $\dyad{v}_{\Flag^{+}}$, the
projector onto the added sector $(v)^{+}$ of $\Qsys^{+}$, or a multiple
of $\Pi_{(f)}$; each is block diagonal for the sectors of the target, and the
sector pinching fixes it.

\emph{(v)} If every $\vartheta_s=0$, \eqref{eq:dom-link} reads
$G^{(s)}_v-F^{(s)}_v\succeq0$ for all $v$ and $s$, and both families are
complete at $1$ on every sector, so $F^{(s)}=G^{(s)}$ by
\xref{foundations/lem:domination}\,(iii). The second sum of
\eqref{eq:noise-shared} vanishes and the first is the sector pinching.

\emph{The closing sentence.} The carried operators of the link in case (a)
are those of \eqref{eq:noise-fresh-relation}, and in case (b) the
projectors of the sectors that are not indicator sectors and $\Pflag$, whose
pull-backs are \eqref{eq:noise-shared-pullbacks}; so (ii) and (iii) are the
relation the channel-witness slot of \eqref{eq:link} asks for, with no error
term.
\end{proof}

% claim reductions/domination/rem:rotated-projective
\claimsummary{reductions/domination/rem:rotated-projective}{A projective element on a direction with no Gaussian-rational vector dominates no non-zero lattice element, and completeness then forces a weight of one, so no domination link runs from that block. A noise floor in the model can remove it; after another channel it applies anew.}
\claimbookmark{reductions/domination/rem:rotated-projective}{Remark}{rem:rotated-projective}{Rotated projective elements and the noise floor}
\claimleadin{%
The remark is about domination only, the relation the link of
Definition~\ref{def:domination-link} carries, whose channel acts on each
sector as the identity followed by a branching. A reduction that first
applies another channel to the reference changes what is compared: the
obstruction then applies to the image of the reference, compared with a
lattice family written in the new basis, where the supports may have
rational directions that they lacked. Such reductions are planned beyond
these contracts (\href{../../../spec/future-scope.md}{\nolinkurl{spec/future-scope.md}}).
}
\begin{remark}[Rotated projective elements and the noise floor]\label{rem:rotated-projective}
Let a reference family $G$, complete at $1$ on a block, have an element
$G_j=\dyad{\psi}$ that is a rank-one projector on a direction containing no
non-zero vector of $\mathbb{Q}(i)^{d}$, as happens for an ideal analyser
whose analysed polarisation $\cos\alpha\ket{H}+\sin\alpha\ket{V}$ has an
irrational $\tan\alpha$ on its one-photon block, or for the
two-photon block of a diagonal-basis analyser written in the orthonormal
occupation basis. Any family $F$ with Gaussian-rational entries that $G$
dominates at a weight below one has
$\mathrm{supp}\,F_j\subseteq\mathrm{supp}\,G_j$ by
\xref{foundations/lem:domination}\,(ii), so $F_j=\lambda\dyad{\psi}$ with
$\lambda\ge0$, and a non-zero such $F_j$ would have a column that is a
non-zero Gaussian-rational vector on the direction $\psi$; hence $F_j=0$:
$G_j$ dominates no non-zero lattice element. Completeness then forces the
weight up to one: \xref{foundations/lem:rational-obstruction}\,(v) gives
$\vartheta\ge\lmax(G_j)/1=1$, so no weight below one exists, at any lattice
exponent. No domination link runs from the reference on a block that holds
such an element, and no lattice family is related to it by domination there;
this is the obstruction \xref{foundations/rem:noise-floor} explains. The
remedy within domination is in the model and not in the lattice. With the
device's loss and dark counts composed into the reference, every element is
positive definite on every block under the conditions that the receiver
module states (\xref{receiver/detector-imperfections/prop:noise-floor}:
loss in the receiver and a dark-count probability strictly between zero and
one in every port); loss alone or dark counts alone need not suffice. Where
those conditions hold, the obstruction disappears: a lattice family whose
elements but one lie between zero and their reference elements, the last
taking what completeness leaves, is dominated at a weight of at most
$\delta/(\lambda+\delta)$, with $\delta$ the norm of the residue that last
element absorbs and $\lambda$ the least eigenvalue of the reference's element
at the same outcome, and not of the lattice element, which
the residue raises above it
(\xref{foundations/lem:absorbing-domination}\,(iv);
\xref{foundations/rem:noise-floor}).
After another channel the obstruction is to be read in the new basis, where
it may not arise: an ideal analyser at an angle with irrational
$\tan\alpha$ is the
image, under one unitary applied before every outcome, of the same analyser
at an angle where its elements are lattice values, and the isometry from the
two-photon block into $\mathbb{C}^{2}\otimes\mathbb{C}^{2}$ carries the
diagonal-basis element above to a projector on a rational vector; neither
route is a claim of this contract. The \termref{term-emission-normalizer}{emission normalizer} does not meet the
obstruction in the form above: the datum it snaps has Gaussian-rational
entries, so every support it meets is spanned by Gaussian-rational vectors
(\xref{foundations/lem:rational-obstruction}\,(ii)).
\end{remark}
\claimmeta{\metalabel{Uses}\mref{Def.}{def:domination-link}, \extref{../foundations/contract.pdf}{foundations/def:domination}, \extref{../foundations/contract.pdf}{foundations/lem:absorbing-domination}, \extref{../foundations/contract.pdf}{foundations/lem:domination}, \extref{../foundations/contract.pdf}{foundations/lem:rational-obstruction}
\metasep newly derived
\metasep accepted}

\section{Domination in a reduction chain}\label{group-reductions/domination/cor:dom-chain}
\sectionintro{A squashing link at an exact datum followed by the \termref{term-domination-link}{domination link} to another datum it dominates is an exact chain, so security transfers through it with the same \ensuremath{\varepsilon} and every entry is decided on the composite pull-back (\hyperlink{claim-reductions/domination/cor:dom-chain}{Corollary~\ref*{cor:dom-chain}}). No requirement emits such a chain; it is the relation a declared device weight would use.}
\sectionimports{\extref{../foundations/contract.pdf}{foundations/def:domination}}

% claim reductions/domination/cor:dom-chain
\claimsummary{reductions/domination/cor:dom-chain}{If an exact datum dominates another datum, the instance's squashing link at the exact datum followed by the domination link is an exact chain from the model at the exact datum to the other target with flags. Security of that target transfers to the model with the same \ensuremath{\varepsilon}.}
\claimbookmark{reductions/domination/cor:dom-chain}{Corollary}{cor:dom-chain}{A squashing link followed by a domination link}
\claimleadin{%
This is Theorem~\ref{thm:chain} for a chain of two receiver-stage links: the
squashing link of a \termref{term-reduction-instance}{reduction instance} at an exact datum, and the
\termref{term-domination-link}{domination link} from its target to the target at another datum. Where the
exact datum is formed from a reference by trusted classical post-processing
--- the dark counts and the \termref{term-class-weights}{class weights} --- and the domination is
established after that post-processing, every map
on the physical outcome acts before the first squashing link, and the
domination link is a squashing link on the target system with an exact
identity, not a post-processing map on an outcome; the order of
Requirement~\ref{req:post-processing-order} is therefore kept
(Remark~\ref{rem:intertwining}). Hypothesis (a) below is the existence
theorem (I3) of the instance where the exact datum is admissible, with
$M^{\circ}=M[A^{\circ}]$ and $\Lambda^{\circ}=\Lambda[A^{\circ}]$; it is
stated on its own so that an exact datum whose entries are not in the
substrate, such as a reference with square-root entries, can be used where an
instance proves the relation for it directly.
}
\begin{corollary}[A squashing link followed by a domination link]\label{cor:dom-chain}
Let $\mathcal{J}$ be a reduction instance with \termref{term-sector-set}{sector set}
$\SecSet=\SecSet_{\mathrm{dat}}\sqcup\SecSet_{\mathrm{sc}}$, \termref{term-target-family}{target family}
$F[\cdot]$ and \termref{term-model-map}{model map} $M[\cdot]$
(Definition~\ref{def:reduction-instance}), and let $A^{\circ}$ and $A'$ be
data for its \termref{term-data-sector}{data sectors} such that
\begin{enumerate}[label=(\alph*),leftmargin=2.4em]
\item $A'\in\mathcal{D}_{\mathrm{adm}}$; the blocks $A^{\circ(s)}_v$ are
  Hermitian, positive semidefinite and \termref{term-class-complete}{class complete}, so that the
  \termref{term-target-family}{target family} $F[A^{\circ}]$ of \eqref{eq:instance-target} is a POVM
  on $\Qsys$; and there are a POVM $M^{\circ}$ on $\Bsys$, a channel
  $\Lambda^{\circ}\colon\Tcl{\Bsys}\to\Tcl{\Qsys}$ and a finite set
  $\hat O^{\circ}$ of operators on $\Qsys$ containing every sector projector,
  with pull-backs $O^{\circ}$, such that
  $\big(M^{\circ};O^{\circ}\big)\dsq_{\Lambda^{\circ}}
  \big(F[A^{\circ}];\hat O^{\circ}\big)$;
\item $A^{\circ}$ dominates $A'$ at weights
  $\vartheta_s\in[0,1)$, $s\in\SecSet_{\mathrm{dat}}$:
  $A^{\circ(s)}_v-(1-\vartheta_s)A'^{(s)}_v\succeq0$ for every $v\in\Vset$
  (\xref{foundations/def:domination}).
\end{enumerate}
Put $\vartheta_s\coloneqq0$ on every \termref{term-scalar-sector}{scalar sector}. Then:
\begin{enumerate}[label=(\roman*),leftmargin=3.0em]
\item $F[A^{\circ}]$ dominates $F[A']$ at the \termref{term-domination-weights}{domination weights}
  $\vartheta=(\vartheta_s)_{s\in\SecSet}$, as POVMs on $\Qsys$ block diagonal for
  the sectors.
\item Let $\Lnk_{\mathrm{dom}}$ be the domination link
  (Definition~\ref{def:domination-link}) with \termref{term-source-family}{source family}
  $(\Qsys,F[A^{\circ}])$ and target family $F[A']$, with added flag sectors
  (case (a) of that definition), or with shared flag sectors (case (b)) where
  $F[A']$ has \termref{term-indicator-sectors}{indicator sectors} among its scalar sectors; write $F^{\mathrm{dom}}$ and $\Phi$ for its target and \termref{term-channel-witness}{channel
  witness}. Let $\Lnk_1$ be the receiver-stage link with source family
  $(\Bsys,M^{\circ})$, target $(\Qsys,F[A^{\circ}])$, channel witness
  $\Lambda^{\circ}$ and target-side \termref{term-carried-operators}{carried operators} the pull-backs
  under $\Phi$ of the carried operators of $\Lnk_{\mathrm{dom}}$, with
  $\textsf{status}=\textsf{exact}$. Then $(\Lnk_1,\Lnk_{\mathrm{dom}})$ is a
  composable and tuple-chained \termref{term-receiver-segment}{receiver segment}, and with an empty source
  segment a \termref{term-reduction-chain}{reduction chain} (Definition~\ref{def:chain}).
\item Its \termref{term-receiver-composite}{receiver composite} $\Phi\circ\Lambda^{\circ}$ satisfies
  $M^{\circ}\dsq_{\Phi\circ\Lambda^{\circ}}F^{\mathrm{dom}}$ exactly,
  its composite status is $\textsf{exact}$, and
  Theorem~\ref{thm:factorization} and Theorem~\ref{thm:transfer} apply to
  it: the protocol whose measurement is $M^{\circ}$ and the \termref{term-squashed-protocol}{squashed
  protocol} with target $F^{\mathrm{dom}}$ have the same output state, and
  security of the squashed protocol over analysed inputs that contain the
  composite's image transfers to the protocol with $M^{\circ}$ with the
  same $\varepsilon$.
\item An entry $(\hat O,\bowtie,c)$ declared on $F^{\mathrm{dom}}$ is sound for
  the composite, and Theorem~\ref{thm:licence} applies to the list of such
  entries with the composite channel, when
  $(\Lambda^{\circ})^{\dagger}\big(\Phi^{\dagger}(\hat O)\big)\bowtie
  c\,\idon{\Bsys}$. Every projector onto a sector of the target has a
  positive semidefinite composite pull-back.
\item Where $\sum_{v\in\Vset_t}M^{\circ}_v=q_t\idon{\Bsys}$ for every
  $t\in\Tset$, the \termref{term-class-marginal-identity}{class-marginal identity}
  \eqref{eq:instance-classmarginal} of $M^{\circ}$, the composite pulls
  $\sum_{v\in\Vset_t}F^{\mathrm{dom}}_v$ back to $q_t\idon{\Bsys}$ for every
  $t\in\Tset$, so the class-marginal entry of
  Definition~\ref{def:reduction-instance}\,(I4) is certified for the
  composite.
\end{enumerate}
\end{corollary}
\claimmeta{\metalabel{Uses}\mref{Cor.}{cor:pullback}, \mref{Thm.}{thm:factorization}, \mref{Thm.}{thm:transfer}, \mref{Thm.}{thm:licence}, \mref{Def.}{def:link}, \mref{Def.}{def:chain}, \mref{Thm.}{thm:chain}, \mref{Def.}{def:reduction-instance}, \mref{Def.}{def:domination-link}, \mref{Thm.}{thm:noise-channel}, \extref{../foundations/contract.pdf}{foundations/def:domination}
\metasep newly derived
\metasep accepted}

\begin{proof}
(i) Both families are POVMs on $\Qsys$, $F[A^{\circ}]$ by (a) and $F[A']$ by
the first consequence of Definition~\ref{def:reduction-instance}, and both
are block diagonal by \eqref{eq:instance-target}. On a \termref{term-data-sector}{data sector} $s$ the
blocks are $A^{\circ(s)}_v$ and $A'^{(s)}_v$, and \eqref{eq:dom-link} at $s$ is
(b). On a scalar sector both blocks are $\phi_s(v)\idon{d_s}$, the \termref{term-scalar-value}{scalar
values} reading no datum, so the difference at weight $0$ is zero.

(ii) Composability (C1): the target of $\Lnk_1$ is the source family of
$\Lnk_{\mathrm{dom}}$. Tuple chaining (C2): by
Theorem~\ref{thm:noise-channel}\,(ii) and (iii) the pull-back under $\Phi$ of
every carried operator of $\Lnk_{\mathrm{dom}}$ (Definition~\ref{def:domination-link})
is a real linear combination of sector projectors of $\Qsys$, namely
$(1-\vartheta_s)\Pi_s$, $\sum_s\vartheta_s\Pi_s$ or $\Pflag+\sum_s\vartheta_s\Pi_s$;
every sector projector is in the \termref{term-transported-set}{transported set} of (I3), and by the last
sentence of (I3), read for the relation of (a), such a combination is
transported by $\Lambda^{\circ}$ to the same combination of the pull-backs
of the sector projectors, by linearity of $(\Lambda^{\circ})^{\dagger}$. So
$\Lnk_1$ carries exactly the operators $\Lnk_{\mathrm{dom}}$ needs on its
source side, and its relation, which is (a) with those operators, carries no
error term. Both links are receiver-stage links of
Definition~\ref{def:link}: $\Lnk_1$ by (a), $\Lnk_{\mathrm{dom}}$ by
Definition~\ref{def:domination-link} and Theorem~\ref{thm:noise-channel}.
Stage order (C3) is structural, the source segment being empty.

(iii) Theorem~\ref{thm:chain}\,(i) gives the relation for the composite,
with $(\Phi\circ\Lambda^{\circ})^{\dagger}
=(\Lambda^{\circ})^{\dagger}\circ\Phi^{\dagger}$; both links are exact, so
Theorem~\ref{thm:chain}\,(iv) gives the composite status \textsf{exact}; and
Theorem~\ref{thm:chain}\,(ii) is the application of
Theorem~\ref{thm:factorization} and Theorem~\ref{thm:transfer}.

(iv) is Theorem~\ref{thm:chain}\,(iii) with the composite pull-back of (iii).
For a sector projector $\Pi$ of the target, $\Phi^{\dagger}(\Pi)\succeq0$
because the adjoint of a channel is a positive map, and
$(\Lambda^{\circ})^{\dagger}$ is positive for the same reason; so the
certificate $(\Pi,\succeq,0)$ of Corollary~\ref{cor:pullback}\,(ii) holds for
the composite.

(v) The domination link relates $F[A^{\circ}]$ to $F^{\mathrm{dom}}$ exactly,
$\Phi^{\dagger}(F^{\mathrm{dom}}_v)=F_v[A^{\circ}]$, so by linearity
$\Phi^{\dagger}\big(\sum_{v\in\Vset_t}F^{\mathrm{dom}}_v\big)
=\sum_{v\in\Vset_t}F_v[A^{\circ}]$, and by (a)
$(\Lambda^{\circ})^{\dagger}$ of that is $\sum_{v\in\Vset_t}M^{\circ}_v$,
which is $q_t\idon{\Bsys}$ by hypothesis.
\end{proof}
\clearpage
\bookmarksetupnext{level=part}
\section*{Interface}\addcontentsline{toc}{section}{Interface}
\sectionintro{The declared inputs and outputs of the module and its submodules: each key with its symbol, when it is present and the claims that fix it. The domains, types and refusal rules are in the interface files themselves.}
\subsection*{Inputs: Proof reductions by squashing (reductions)}
\noindent{\small From \href{interface.yaml}{\nolinkurl{interface.yaml}}.}\par\smallskip
{\footnotesize\begin{xltabular}{\linewidth}{@{}>{\raggedright\arraybackslash}p{4.4cm}>{\raggedright\arraybackslash}X>{\raggedright\arraybackslash}p{2.9cm}@{}}
\toprule
Key and symbol & Presence & Fixed by\\
\midrule
\endhead
\iface{lattice_exponent}\newline kappa & always; REQUIRED, no default and no fallback & \mref{Def.}{def:exactadm}, \mref{Lem.}{lem:latsum}\\
\addlinespace
\iface{strict_snaps}\newline bool & read from the run configuration; off when the run does not declare it, and recorded either way & \mref{Def.}{def:emitnorm}, \mref{Req.}{req:emission-refusal}, \mref{Req.}{req:emission-record}, \extref{../foundations/contract.pdf}{foundations/req:snap-reference}, \extref{../foundations/contract.pdf}{foundations/req:support-safe-snap}\\
\addlinespace
\iface{detection_alphabet}\newline Y & always & \mref{Def.}{def:joint-family}\\
\addlinespace
\iface{probe_declaration}\newline T, \{q\_t\} & always & \mref{Def.}{def:joint-family}, \mref{Ass.}{ass:probe}, \mref{Lem.}{lem:latsum}\\
\addlinespace
\iface{joint_alphabet}\newline V = Y x T with its classes \{V\_t\} and weights \{q\_t\} & always & \mref{Def.}{def:joint-family}\\
\addlinespace
\iface{sector_declaration}\newline S = S\_dat + S\_sc, with dimensions d\_s, supports V\textasciicircum{}(s), and the order & always & \mref{Def.}{def:reduction-instance}, \mref{Def.}{def:exactadm}\\
\addlinespace
\iface{declared_datum}\newline A = \{A\textasciicircum{}(s)\_v\} & always & \mref{Def.}{def:reduction-instance}, \mref{Def.}{def:exactadm}\\
\addlinespace
\iface{reduction_instance}\newline J = (I1)-(I4) & always & \mref{Def.}{def:reduction-instance}\\
\addlinespace
\iface{receiver_artifact} & always for an instance that reads a supplier artifact & \mref{Ass.}{ass:recv}, \mref{Req.}{req:supplier-version-accepted}, \mref{Req.}{req:supplier-ledger-refusals}\\
\addlinespace
\iface{handed_reference}\newline the reference G of the declaring instance's snapped receiver blocks, per setting and block: centres C (Gaussian rationals) and radii r (rationals) & where the declaring instance hands the reference beside its snap of the receiver blocks (foundations/def:evaluated-off-snap (i)), on every block it hands & \mref{Cor.}{cor:reference-companion}, \mref{Req.}{req:reference-companion}, \extref{../foundations/contract.pdf}{foundations/def:evaluated-off-snap}, \extref{../foundations/contract.pdf}{foundations/lem:recentred-enclosure}\\
\addlinespace
\iface{accepted_receiver_versions} & always & \mref{Req.}{req:supplier-version-accepted}, \mref{Ass.}{ass:recv}\\
\addlinespace
\iface{receiver_provenance} & optional; passed only by the part of the engine that ran the receiver module in the same run and docked its artifact into the instance, and absent on every other route & \mref{Req.}{req:recv-row-provenance}, \mref{Ass.}{ass:recv}, \extref{../receiver/contract.pdf}{receiver/req:exact-route-decision}, \extref{../receiver/contract.pdf}{receiver/req:evaluation-label}\\
\addlinespace
\iface{link_declarations}\newline l\_1, ..., l\_k & always & \mref{Def.}{def:link}, \mref{Def.}{def:chain}, \mref{Rem.}{rem:ordering}, \mref{Thm.}{thm:two-sided}\\
\addlinespace
\iface{carried_operators}\newline hat O\_1, ..., hat O\_m & always; the DEFAULT is the instance's DEFAULT CARRIED SET & \mref{Def.}{def:reduction-instance}, \mref{Cor.}{cor:entries}, \mref{Cor.}{cor:pullback}, \mref{Req.}{req:transported-set-certificates}, \mref{Req.}{req:uncertified-operator-refused}\\
\addlinespace
\iface{requested_assurance} & always & \extref{../foundations/contract.pdf}{foundations/req:tiers}, \extref{../foundations/contract.pdf}{foundations/def:tiers}\\
\addlinespace
\iface{source_mode}\newline tau\_A or nomarg & always & \mref{Def.}{def:licensed}, \mref{Def.}{def:attset}\\
\addlinespace
\iface{declared_relations}\newline (M, F) pairs & optional; empty by default & \mref{Lem.}{lem:obstructions}, \mref{Rem.}{rem:obstructiondecide}, \mref{Req.}{req:declared-relation-check}\\
\addlinespace
\iface{post_processing_order_declaration} & always & \mref{Rem.}{rem:intertwining}, \mref{Def.}{def:joint-family}, \mref{Req.}{req:post-processing-order}\\
\addlinespace
\iface{declared_orders} & always & \mref{Def.}{def:reduction-instance}, \mref{Def.}{def:joint-family}\\
\addlinespace
\iface{public_choice_applied} & OPTIONAL, one per call (per station). ABSENT MEANS NOT TOLD, and is the default. & \mref{Req.}{req:assumption-ledger}, \mref{Ass.}{ass:tsm}\\
\bottomrule
\end{xltabular}}
\subsection*{Outputs: Proof reductions by squashing (reductions)}
\noindent{\small From \href{interface.yaml}{\nolinkurl{interface.yaml}}.}\par\smallskip
{\footnotesize\begin{xltabular}{\linewidth}{@{}>{\raggedright\arraybackslash}p{4.4cm}>{\raggedright\arraybackslash}X>{\raggedright\arraybackslash}p{2.9cm}@{}}
\toprule
Key and symbol & Presence & Fixed by\\
\midrule
\endhead
\iface{source_descriptor}\newline B & always & \mref{Def.}{def:reduction-instance}, \mref{Def.}{def:system-setting}\\
\addlinespace
\iface{target_descriptor}\newline Q, S, the sector map and the party partition & always & \mref{Def.}{def:reduction-instance}, \mref{Def.}{def:exactadm}, \mref{Thm.}{thm:emit}, \mref{Thm.}{thm:emitted-product}, \mref{Req.}{req:party-label}\\
\addlinespace
\iface{target_povm}\newline \{F\_v\}\_\{v in V\} & always & \mref{Def.}{def:reduction-instance}, \mref{Thm.}{thm:emit}\\
\addlinespace
\iface{factor_emission}\newline the ordered factor list, with the combination rule & PRESENT EXACTLY WHEN THE CALL DECLARES MORE THAN ONE INSTANCE -- one per measuring party. ABSENT FOR A ONE-FACTOR EMISSION, and nothing else in this file changes there. & \mref{Thm.}{thm:emitted-product}, \mref{Def.}{def:reduction-instance}, \mref{Prop.}{prop:tensor}, \mref{Cor.}{cor:blockdiag}, \mref{Req.}{req:factorised-emission}\\
\addlinespace
\iface{outcome_correspondence}\newline id\_V & always & \mref{Def.}{def:joint-family}, \mref{Rem.}{rem:intertwining}\\
\addlinespace
\iface{squashing_channel}\newline Lambda = Pin\_S o Lambda\_0 & always & \mref{Def.}{def:reduction-instance}, \mref{Prop.}{prop:pinching}, \mref{Cor.}{cor:blockdiag}\\
\addlinespace
\iface{exact_identity}\newline Lambda\textasciicircum{}dagger(F\_v) = M\textasciicircum{}dec\_v & always & \mref{Thm.}{thm:emit}, \mref{Def.}{def:emitnorm}\\
\addlinespace
\iface{transported_operators}\newline (hat O\_k, O\_k) & always & \mref{Def.}{def:squashing}, \mref{Cor.}{cor:pullback}, \mref{Cor.}{cor:transported}, \mref{Def.}{def:reduction-instance}, \mref{Req.}{req:transported-set-certificates}\\
\addlinespace
\iface{reachable_set_statement}\newline S\textasciicircum{}rch & always & \mref{Def.}{def:sound}, \mref{Cor.}{cor:sound}\\
\addlinespace
\iface{constraint_list}\newline CList = \{(hat O\_k, bowtie\_k, c\_k)\} & always & \mref{Def.}{def:licensed}, \mref{Def.}{def:attset}, \mref{Cor.}{cor:entries}, \mref{Rem.}{rem:entry-tags}, \mref{Req.}{req:constraint-list-by-request}, \mref{Req.}{req:party-local-tag}, \mref{Req.}{req:declared-constants-only}, \mref{Thm.}{thm:licence}, \mref{Ass.}{ass:constdecl}\\
\addlinespace
\iface{attack_set_license} & always as a slot (req:constraint-and-licence-slots); its value is the licence record EXACTLY WHEN the constraint list is non-empty, and null otherwise & \mref{Thm.}{thm:licence}, \mref{Def.}{def:attset}, \mref{Rem.}{rem:gateweight}, \mref{Cor.}{cor:entries}, \mref{Req.}{req:licence-record}, \mref{Req.}{req:constraint-list-by-request}\\
\addlinespace
\iface{reduction_chain}\newline Chain = (Chain\_src, Chain\_rec) & always & \mref{Def.}{def:link}, \mref{Def.}{def:chain}, \mref{Thm.}{thm:chain}, \mref{Thm.}{thm:two-sided}, \mref{Rem.}{rem:ordering}\\
\addlinespace
\iface{status} & always & \mref{Def.}{def:link}, \mref{Thm.}{thm:chain}, \mref{Thm.}{thm:emit}, \mref{Req.}{req:mixed-status-chain}, \mref{Req.}{req:no-approximate-relations}\\
\addlinespace
\iface{verification_report} & always & \mref{Ass.}{ass:recv}\\
\addlinespace
\iface{assumption_ledger}\newline A-... & always & \mref{Ass.}{ass:recv}, \mref{Rem.}{rem:sound-and-atk}, \mref{Lem.}{lem:sep}, \extref{../meat-core/contract.pdf}{meat-core/lem:memory-lift}, \mref{Rem.}{rem:receiver-imports}, \mref{Req.}{req:assumption-ledger}, \mref{Req.}{req:supplier-ledger-pass-through}, \mref{Req.}{req:supplier-ledger-refusals}, \mref{Req.}{req:supplier-version-accepted}, \mref{Req.}{req:supplier-presence-rules}, \mref{Req.}{req:never-verified-assumptions}\\
\addlinespace
\iface{receiver_row_omitted_on_provenance} & always & \mref{Req.}{req:recv-row-provenance}\\
\addlinespace
\iface{emission_record}\newline kappa, (Delta\_s) & ALWAYS PRESENT & \mref{Def.}{def:emitnorm}, \mref{Thm.}{thm:emit}, \mref{Rem.}{rem:emitclaim}, \mref{Req.}{req:emission-record}\\
\addlinespace
\iface{reference_companion}\newline the instance at the datum A\textasciicircum{}ref formed from the handed reference, or at the handed datum A where no block is handed & exactly where a claim of this module establishes (H-REF) of cor:reference-companion for the instance (cor:reference-companion (iv); wpfss/cor:flat-companion), and either the declaring instance hands the reference on at least one block or the emission normalizer moves at least one class; otherwise absent, the record saying which of the two failed (req:reference-companion (vi)). Where no block is handed, the companion carries the snapped block on every block, so its datum is the handed datum A, the module's exact construction before the normalizer (cor:reference-companion (i)): the moved classes are evaluated off their snap, and the instance's moved blocks still rest on A-SNAP. & \mref{Cor.}{cor:reference-companion}, \mref{Cor.}{cor:companion-factor-enclosure}, \mref{Lem.}{lem:row-relaxation}, \mref{Req.}{req:reference-companion}, \mref{Cor.}{cor:flat-companion}\\
\addlinespace
\iface{mathematical_spec_version} & always & ---\\
\bottomrule
\end{xltabular}}
\subsection*{Inputs: Flag-state and weight-preserving squashers}
\noindent{\small From \href{wpfss/interface.yaml}{\nolinkurl{wpfss/interface.yaml}}.}\par\smallskip
{\footnotesize\begin{xltabular}{\linewidth}{@{}>{\raggedright\arraybackslash}p{4.4cm}>{\raggedright\arraybackslash}X>{\raggedright\arraybackslash}p{2.9cm}@{}}
\toprule
Key and symbol & Presence & Fixed by\\
\midrule
\endhead
\iface{photon_blocks}\newline blk M\textasciicircum{}\{=m\}\_\{y\textbar{}t\}, 0 \textless{}= m \textless{}= N & always & \mref{Def.}{def:weight-route}, \mref{Def.}{def:exactadm}, \mref{Prop.}{prop:route}, \extref{../receiver/contract.pdf}{receiver/thm:click-povm}, \extref{../receiver/contract.pdf}{receiver/thm:coarse-graining}\\
\addlinespace
\iface{photon_blocks_form}\newline bare \textbar{} dark\_composed & always; NO DEFAULT -- an artefact that does not state the form is refused & \mref{Def.}{def:weight-route}, \mref{Prop.}{prop:flat-instance}, \mref{Prop.}{prop:precomposed-blocks}\\
\addlinespace
\iface{high_photon_bounds}\newline \{(O\_i, (p\_i\textasciicircum{}(t))\_\{t in T\})\}\_\{i=1\}\textasciicircum{}\{r\} & always & \mref{Def.}{def:weight-route}, \mref{Prop.}{prop:route}, \extref{../receiver/contract.pdf}{receiver/thm:emitted-bound}, \extref{../receiver/contract.pdf}{receiver/def:slot}, \extref{../receiver/contract.pdf}{receiver/prop:slotvalid}, \mref{Rem.}{rem:minweight}, \mref{Req.}{req:supplier-version}\\
\addlinespace
\iface{bound_quantifier_domain}\newline N, T & always & \mref{Def.}{def:weight-route}, \mref{Prop.}{prop:route}\\
\addlinespace
\iface{dark_count_channel}\newline R & always & \mref{Def.}{def:weight-route}, \mref{Lem.}{lem:descent}, \mref{Prop.}{prop:route}, \extref{../receiver/contract.pdf}{receiver/detector-imperfections/def:darkadm}, \extref{../receiver/contract.pdf}{receiver/detector-imperfections/thm:dark-counts}, \mref{Rem.}{rem:absorbneeded}, \mref{Req.}{req:column-stochastic-refusal}, \mref{Req.}{req:descent-decided-exactly}, \mref{Req.}{req:absorbing-bounded-letters}\\
\addlinespace
\iface{outcome_coarse_graining}\newline cg & always & \mref{Def.}{def:weight-route}, \mref{Lem.}{lem:descent}, \mref{Rem.}{rem:fibre}, \mref{Req.}{req:fibre-equality}\\
\addlinespace
\iface{bound_evaluation} & always & \mref{Rem.}{rem:receiver-imports}, \mref{Ass.}{ass:recv}, \mref{Def.}{def:exactadm}, \extref{../receiver/contract.pdf}{receiver/req:evaluation-label}\\
\addlinespace
\iface{factored_loss}\newline t & always & \mref{Cor.}{cor:normalized-flat}, \mref{Req.}{req:normalized-link-record}, \mref{Req.}{req:normalized-branch-refusal}, \extref{../receiver/contract.pdf}{receiver/detector-imperfections/prop:efficiency-normalization}\\
\addlinespace
\iface{photon_cutoff}\newline N & always & \mref{Def.}{def:weight-route}, \mref{Prop.}{prop:flat-instance}, \mref{Def.}{def:targets}\\
\addlinespace
\iface{route} & always & \mref{Def.}{def:routes}, \mref{Rem.}{rem:tworoutes}, \mref{Def.}{def:targets}, \mref{Prop.}{prop:flat-instance}, \mref{Req.}{req:one-route-per-artifact}\\
\bottomrule
\end{xltabular}}
\subsection*{Outputs: Flag-state and weight-preserving squashers}
\noindent{\small From \href{wpfss/interface.yaml}{\nolinkurl{wpfss/interface.yaml}}.}\par\smallskip
{\footnotesize\begin{xltabular}{\linewidth}{@{}>{\raggedright\arraybackslash}p{4.4cm}>{\raggedright\arraybackslash}X>{\raggedright\arraybackslash}p{2.9cm}@{}}
\toprule
Key and symbol & Presence & Fixed by\\
\midrule
\endhead
\iface{wpfss_data}\newline \{p\_v\}, flag basis, N, p = sum\_v p\_v & always, for an artifact emitted from the flat instance & \mref{Def.}{def:targets}, \mref{Def.}{def:weight-route}, \mref{Lem.}{lem:ple1}, \mref{Rem.}{rem:peqone}, \mref{Rem.}{rem:weightfree}, \mref{Req.}{req:largest-certified-weight}, \mref{Req.}{req:total-weight-refusal}\\
\addlinespace
\iface{route} & always, for an artifact emitted from the flat instance & \mref{Def.}{def:routes}, \mref{Rem.}{rem:tworoutes}, \mref{Prop.}{prop:flat-instance}, \mref{Req.}{req:one-route-per-artifact}\\
\bottomrule
\end{xltabular}}
\clearpage
\bookmarksetupnext{level=part}
\section*{Claim index}\addcontentsline{toc}{section}{Claim index}
{\small\begin{xltabular}{\linewidth}{@{}>{\raggedright\arraybackslash}p{2.7cm}>{\raggedright\arraybackslash}Xr@{}}
\toprule
Claim & Title and identifier & Page\\
\midrule
\endhead
\mref{Definition}{def:system-setting} & The system setting of the reductions\newline{\footnotesize\color{metaGrey}\claimid{reductions/def:system-setting}} & \pageref{def:system-setting}\\
\addlinespace[2pt]
\mref{Lemma}{lem:sep} & Separation and tensor adjoints in the infinite-dimensional case\newline{\footnotesize\color{metaGrey}\claimid{reductions/lem:sep}} & \pageref{lem:sep}\\
\addlinespace[2pt]
\mref{Definition}{def:squashing} & Squashing relation\newline{\footnotesize\color{metaGrey}\claimid{reductions/def:squashing}} & \pageref{def:squashing}\\
\addlinespace[2pt]
\mref{Lemma}{lem:statistics} & Squashing as statistics equivalence\newline{\footnotesize\color{metaGrey}\claimid{reductions/lem:statistics}} & \pageref{lem:statistics}\\
\addlinespace[2pt]
\mref{Remark}{rem:dictionary} & Direction of the squashing symbol in the literature\newline{\footnotesize\color{metaGrey}\claimid{reductions/rem:dictionary}} & \pageref{rem:dictionary}\\
\addlinespace[2pt]
\mref{Remark}{rem:carried-operators} & Carried operators and the image of the squashing channel\newline{\footnotesize\color{metaGrey}\claimid{reductions/rem:carried-operators}} & \pageref{rem:carried-operators}\\
\addlinespace[2pt]
\mref{Corollary}{cor:pullback} & Pull-back of operator bounds\newline{\footnotesize\color{metaGrey}\claimid{reductions/cor:pullback}} & \pageref{cor:pullback}\\
\addlinespace[2pt]
\mref{Definition}{def:certificate} & Certificates of transported constraints\newline{\footnotesize\color{metaGrey}\claimid{reductions/def:certificate}} & \pageref{def:certificate}\\
\addlinespace[2pt]
\mref{Remark}{rem:notchannel} & Reductions that are not squashings\newline{\footnotesize\color{metaGrey}\claimid{reductions/rem:notchannel}} & \pageref{rem:notchannel}\\
\addlinespace[2pt]
\mref{Definition}{def:msm} & Measurement channel\newline{\footnotesize\color{metaGrey}\claimid{reductions/def:msm}} & \pageref{def:msm}\\
\addlinespace[2pt]
\mref{Lemma}{lem:msmfact} & Measurement-channel factorization\newline{\footnotesize\color{metaGrey}\claimid{reductions/lem:msmfact}} & \pageref{lem:msmfact}\\
\addlinespace[2pt]
\mref{Definition}{def:protocol} & Measurement-factorizing $n$-round protocol\newline{\footnotesize\color{metaGrey}\claimid{reductions/def:protocol}} & \pageref{def:protocol}\\
\addlinespace[2pt]
\mref{Definition}{def:criterion} & Security criterion\newline{\footnotesize\color{metaGrey}\claimid{reductions/def:criterion}} & \pageref{def:criterion}\\
\addlinespace[2pt]
\mref{Definition}{def:opreduction} & Exact operational reduction\newline{\footnotesize\color{metaGrey}\claimid{reductions/def:opreduction}} & \pageref{def:opreduction}\\
\addlinespace[2pt]
\mref{Lemma}{lem:opreduction-transfer} & Transfer and composition of exact operational reductions\newline{\footnotesize\color{metaGrey}\claimid{reductions/lem:opreduction-transfer}} & \pageref{lem:opreduction-transfer}\\
\addlinespace[2pt]
\mref{Remark}{rem:why-exact} & Exactness of the reductions and generality of the transfer\newline{\footnotesize\color{metaGrey}\claimid{reductions/rem:why-exact}} & \pageref{rem:why-exact}\\
\addlinespace[2pt]
\mref{Theorem}{thm:factorization} & Equality of the physical and squashed output states\newline{\footnotesize\color{metaGrey}\claimid{reductions/thm:factorization}} & \pageref{thm:factorization}\\
\addlinespace[2pt]
\mref{Definition}{def:roundmodel} & Round model and attack sequences, typed by source mode\newline{\footnotesize\color{metaGrey}\claimid{reductions/def:roundmodel}} & \pageref{def:roundmodel}\\
\addlinespace[2pt]
\mref{Corollary}{cor:sequential} & Output-state equality against adaptive attack sequences\newline{\footnotesize\color{metaGrey}\claimid{reductions/cor:sequential}} & \pageref{cor:sequential}\\
\addlinespace[2pt]
\mref{Proposition}{prop:platform} & Measurement factorization of the platform's protocol channel\newline{\footnotesize\color{metaGrey}\claimid{reductions/prop:platform}} & \pageref{prop:platform}\\
\addlinespace[2pt]
\mref{Remark}{rem:platform-published} & Published forms of the sequential reduction\newline{\footnotesize\color{metaGrey}\claimid{reductions/rem:platform-published}} & \pageref{rem:platform-published}\\
\addlinespace[2pt]
\mref{Remark}{rem:variants} & Variants of the factorization\newline{\footnotesize\color{metaGrey}\claimid{reductions/rem:variants}} & \pageref{rem:variants}\\
\addlinespace[2pt]
\mref{Remark}{rem:coherent} & Coherent attacks covered by the output-state equality\newline{\footnotesize\color{metaGrey}\claimid{reductions/rem:coherent}} & \pageref{rem:coherent}\\
\addlinespace[2pt]
\mref{Remark}{rem:criteria} & Correctness, secrecy and completeness as criteria\newline{\footnotesize\color{metaGrey}\claimid{reductions/rem:criteria}} & \pageref{rem:criteria}\\
\addlinespace[2pt]
\mref{Remark}{rem:completeness} & Completeness statements in the physical and squashed models\newline{\footnotesize\color{metaGrey}\claimid{reductions/rem:completeness}} & \pageref{rem:completeness}\\
\addlinespace[2pt]
\mref{Definition}{def:sound} & Reachable states and sound constraints\newline{\footnotesize\color{metaGrey}\claimid{reductions/def:sound}} & \pageref{def:sound}\\
\addlinespace[2pt]
\mref{Theorem}{thm:transfer} & Security transfer from a constrained squashed model\newline{\footnotesize\color{metaGrey}\claimid{reductions/thm:transfer}} & \pageref{thm:transfer}\\
\addlinespace[2pt]
\mref{Corollary}{cor:sound} & Imposing and intersecting sound constraints\newline{\footnotesize\color{metaGrey}\claimid{reductions/cor:sound}} & \pageref{cor:sound}\\
\addlinespace[2pt]
\mref{Corollary}{cor:transported} & Sound constraints from transported operators\newline{\footnotesize\color{metaGrey}\claimid{reductions/cor:transported}} & \pageref{cor:transported}\\
\addlinespace[2pt]
\mref{Remark}{rem:unsound} & The standard way to make a squashed constraint unsound\newline{\footnotesize\color{metaGrey}\claimid{reductions/rem:unsound}} & \pageref{rem:unsound}\\
\addlinespace[2pt]
\mref{Remark}{rem:channel-choice} & Soundness relative to the channel, fixed by the carried tuple\newline{\footnotesize\color{metaGrey}\claimid{reductions/rem:channel-choice}} & \pageref{rem:channel-choice}\\
\addlinespace[2pt]
\mref{Proposition}{prop:preorder} & The squashing relation as a pre-order\newline{\footnotesize\color{metaGrey}\claimid{reductions/prop:preorder}} & \pageref{prop:preorder}\\
\addlinespace[2pt]
\mref{Corollary}{cor:monotone} & Monotonicity of the worst-case criterion along the pre-order\newline{\footnotesize\color{metaGrey}\claimid{reductions/cor:monotone}} & \pageref{cor:monotone}\\
\addlinespace[2pt]
\mref{Lemma}{lem:linear} & Linearity: post-processing, private settings and mixtures\newline{\footnotesize\color{metaGrey}\claimid{reductions/lem:linear}} & \pageref{lem:linear}\\
\addlinespace[2pt]
\mref{Definition}{def:joint-family} & Joint alphabet, classes and the joint family of a private setting\newline{\footnotesize\color{metaGrey}\claimid{reductions/def:joint-family}} & \pageref{def:joint-family}\\
\addlinespace[2pt]
\mref{Remark}{rem:intertwining} & Classical maps before a squasher, and relabelling after one\newline{\footnotesize\color{metaGrey}\claimid{reductions/rem:intertwining}} & \pageref{rem:intertwining}\\
\addlinespace[2pt]
\mref{Proposition}{prop:pinching} & Sector pinching as a squashing\newline{\footnotesize\color{metaGrey}\claimid{reductions/prop:pinching}} & \pageref{prop:pinching}\\
\addlinespace[2pt]
\mref{Corollary}{cor:blockdiag} & Restriction to block-diagonal squashed states\newline{\footnotesize\color{metaGrey}\claimid{reductions/cor:blockdiag}} & \pageref{cor:blockdiag}\\
\addlinespace[2pt]
\mref{Proposition}{prop:tensor} & Tensor products of squashing relations\newline{\footnotesize\color{metaGrey}\claimid{reductions/prop:tensor}} & \pageref{prop:tensor}\\
\addlinespace[2pt]
\mref{Lemma}{lem:obstructions} & Necessary conditions for a squashing relation\newline{\footnotesize\color{metaGrey}\claimid{reductions/lem:obstructions}} & \pageref{lem:obstructions}\\
\addlinespace[2pt]
\mref{Remark}{rem:obstructiondecide} & Exact decision of the necessary conditions\newline{\footnotesize\color{metaGrey}\claimid{reductions/rem:obstructiondecide}} & \pageref{rem:obstructiondecide}\\
\addlinespace[2pt]
\mref{Lemma}{lem:top} & Flag PVM as a top element\newline{\footnotesize\color{metaGrey}\claimid{reductions/lem:top}} & \pageref{lem:top}\\
\addlinespace[2pt]
\mref{Definition}{def:generalized-instrument} & Generalized instrument\newline{\footnotesize\color{metaGrey}\claimid{reductions/def:generalized-instrument}} & \pageref{def:generalized-instrument}\\
\addlinespace[2pt]
\mref{Lemma}{lem:master} & Instrument-conditioned squashing\newline{\footnotesize\color{metaGrey}\claimid{reductions/lem:master}} & \pageref{lem:master}\\
\addlinespace[2pt]
\mref{Remark}{rem:reading-master} & The operator form of the branch-weight constraint\newline{\footnotesize\color{metaGrey}\claimid{reductions/rem:reading-master}} & \pageref{rem:reading-master}\\
\addlinespace[2pt]
\mref{Remark}{rem:master-scope} & Branch labels as information the adversary may hold\newline{\footnotesize\color{metaGrey}\claimid{reductions/rem:master-scope}} & \pageref{rem:master-scope}\\
\addlinespace[2pt]
\mref{Remark}{rem:branch-blocks} & Conditionally characterized receivers under the branching rule\newline{\footnotesize\color{metaGrey}\claimid{reductions/rem:branch-blocks}} & \pageref{rem:branch-blocks}\\
\addlinespace[2pt]
\mref{Remark}{rem:branch-probability} & The branch probability as an operator\newline{\footnotesize\color{metaGrey}\claimid{reductions/rem:branch-probability}} & \pageref{rem:branch-probability}\\
\addlinespace[2pt]
\mref{Corollary}{cor:instances} & Three special cases of instrument-conditioned squashing\newline{\footnotesize\color{metaGrey}\claimid{reductions/cor:instances}} & \pageref{cor:instances}\\
\addlinespace[2pt]
\mref{Proposition}{prop:common} & Common and dominating targets\newline{\footnotesize\color{metaGrey}\claimid{reductions/prop:common}} & \pageref{prop:common}\\
\addlinespace[2pt]
\mref{Remark}{rem:dominating} & Dominated detectors and device families\newline{\footnotesize\color{metaGrey}\claimid{reductions/rem:dominating}} & \pageref{rem:dominating}\\
\addlinespace[2pt]
\mref{Corollary}{cor:flagsenough} & Arbitrary unpopulated block in a flag-extended target\newline{\footnotesize\color{metaGrey}\claimid{reductions/cor:flagsenough}} & \pageref{cor:flagsenough}\\
\addlinespace[2pt]
\mref{Remark}{rem:flagcost} & Rate cost of a flag-extended target\newline{\footnotesize\color{metaGrey}\claimid{reductions/rem:flagcost}} & \pageref{rem:flagcost}\\
\addlinespace[2pt]
\mref{Theorem}{thm:mixture} & Register-conditioned squashing\newline{\footnotesize\color{metaGrey}\claimid{reductions/thm:mixture}} & \pageref{thm:mixture}\\
\addlinespace[2pt]
\mref{Remark}{rem:povmreg} & A general POVM on the register, and its price\newline{\footnotesize\color{metaGrey}\claimid{reductions/rem:povmreg}} & \pageref{rem:povmreg}\\
\addlinespace[2pt]
\mref{Proposition}{prop:obstruction} & Obstruction to encoding event bounds in classical targets\newline{\footnotesize\color{metaGrey}\claimid{reductions/prop:obstruction}} & \pageref{prop:obstruction}\\
\addlinespace[2pt]
\mref{Remark}{rem:obstruction-scope} & Scope of the obstruction, and what it forces\newline{\footnotesize\color{metaGrey}\claimid{reductions/rem:obstruction-scope}} & \pageref{rem:obstruction-scope}\\
\addlinespace[2pt]
\mref{Definition}{def:attset} & Admissible attack sets in the consumer's format\newline{\footnotesize\color{metaGrey}\claimid{reductions/def:attset}} & \pageref{def:attset}\\
\addlinespace[2pt]
\mref{Remark}{rem:relation-symbols} & Relation symbols at the channel and declaration levels\newline{\footnotesize\color{metaGrey}\claimid{reductions/rem:relation-symbols}} & \pageref{rem:relation-symbols}\\
\addlinespace[2pt]
\mref{Definition}{def:licensed} & Licensed entries and the party-local form\newline{\footnotesize\color{metaGrey}\claimid{reductions/def:licensed}} & \pageref{def:licensed}\\
\addlinespace[2pt]
\mref{Lemma}{lem:party-local} & Party-local transport\newline{\footnotesize\color{metaGrey}\claimid{reductions/lem:party-local}} & \pageref{lem:party-local}\\
\addlinespace[2pt]
\mref{Theorem}{thm:licence} & The attack-set licence\newline{\footnotesize\color{metaGrey}\claimid{reductions/thm:licence}} & \pageref{thm:licence}\\
\addlinespace[2pt]
\mref{Remark}{rem:notlicensed} & Scope of the attack-set licence\newline{\footnotesize\color{metaGrey}\claimid{reductions/rem:notlicensed}} & \pageref{rem:notlicensed}\\
\addlinespace[2pt]
\mref{Remark}{rem:gateweight} & Weight of the licence in the consumer's gate\newline{\footnotesize\color{metaGrey}\claimid{reductions/rem:gateweight}} & \pageref{rem:gateweight}\\
\addlinespace[2pt]
\mref{Remark}{rem:publishedlicence} & The published squashed-attack-set format\newline{\footnotesize\color{metaGrey}\claimid{reductions/rem:publishedlicence}} & \pageref{rem:publishedlicence}\\
\addlinespace[2pt]
\mref{Definition}{def:link} & Typed link\newline{\footnotesize\color{metaGrey}\claimid{reductions/def:link}} & \pageref{def:link}\\
\addlinespace[2pt]
\mref{Definition}{def:chain} & Reduction chain\newline{\footnotesize\color{metaGrey}\claimid{reductions/def:chain}} & \pageref{def:chain}\\
\addlinespace[2pt]
\mref{Remark}{rem:chain-admission} & The admission rule and what this module emits\newline{\footnotesize\color{metaGrey}\claimid{reductions/rem:chain-admission}} & \pageref{rem:chain-admission}\\
\addlinespace[2pt]
\mref{Theorem}{thm:chain} & Composition of a chain\newline{\footnotesize\color{metaGrey}\claimid{reductions/thm:chain}} & \pageref{thm:chain}\\
\addlinespace[2pt]
\mref{Example}{ex:identity} & The identity link\newline{\footnotesize\color{metaGrey}\claimid{reductions/ex:identity}} & \pageref{ex:identity}\\
\addlinespace[2pt]
\mref{Lemma}{lem:precomposition} & Licensed admissibility survives precomposition with a source channel\newline{\footnotesize\color{metaGrey}\claimid{reductions/lem:precomposition}} & \pageref{lem:precomposition}\\
\addlinespace[2pt]
\mref{Remark}{rem:precomposition-scope} & Scope of the precomposition lemma\newline{\footnotesize\color{metaGrey}\claimid{reductions/rem:precomposition-scope}} & \pageref{rem:precomposition-scope}\\
\addlinespace[2pt]
\mref{Lemma}{lem:source-transfer} & Exactness of the source map in trusted-source mode\newline{\footnotesize\color{metaGrey}\claimid{reductions/lem:source-transfer}} & \pageref{lem:source-transfer}\\
\addlinespace[2pt]
\mref{Remark}{rem:source-transfer-registers} & The released signals as systems in transit\newline{\footnotesize\color{metaGrey}\claimid{reductions/rem:source-transfer-registers}} & \pageref{rem:source-transfer-registers}\\
\addlinespace[2pt]
\mref{Remark}{rem:source-transfer-scope} & Scope of the source-map lemma\newline{\footnotesize\color{metaGrey}\claimid{reductions/rem:source-transfer-scope}} & \pageref{rem:source-transfer-scope}\\
\addlinespace[2pt]
\mref{Theorem}{thm:two-sided} & The two-sided reduction in trusted-source mode\newline{\footnotesize\color{metaGrey}\claimid{reductions/thm:two-sided}} & \pageref{thm:two-sided}\\
\addlinespace[2pt]
\mref{Remark}{rem:two-sided-scope} & Scope of the two-sided reduction\newline{\footnotesize\color{metaGrey}\claimid{reductions/rem:two-sided-scope}} & \pageref{rem:two-sided-scope}\\
\addlinespace[2pt]
\mref{Remark}{rem:ordering} & Recorded stage order, immaterial to soundness for licensed lists\newline{\footnotesize\color{metaGrey}\claimid{reductions/rem:ordering}} & \pageref{rem:ordering}\\
\addlinespace[2pt]
\mref{Definition}{def:exactadm} & The number model, the declared lattice and exact admissibility\newline{\footnotesize\color{metaGrey}\claimid{reductions/def:exactadm}} & \pageref{def:exactadm}\\
\addlinespace[2pt]
\mref{Definition}{def:substrate} & The arithmetic substrate\newline{\footnotesize\color{metaGrey}\claimid{reductions/def:substrate}} & \pageref{def:substrate}\\
\addlinespace[2pt]
\mref{Lemma}{lem:exactpsd} & Exact decision of positivity and of an operator-norm bound\newline{\footnotesize\color{metaGrey}\claimid{reductions/lem:exactpsd}} & \pageref{lem:exactpsd}\\
\addlinespace[2pt]
\mref{Remark}{rem:psdalgorithm} & Algorithms that decide the positivity test\newline{\footnotesize\color{metaGrey}\claimid{reductions/rem:psdalgorithm}} & \pageref{rem:psdalgorithm}\\
\addlinespace[2pt]
\mref{Lemma}{lem:latsum} & Partial sums on the declared lattice are exact\newline{\footnotesize\color{metaGrey}\claimid{reductions/lem:latsum}} & \pageref{lem:latsum}\\
\addlinespace[2pt]
\mref{Definition}{def:reduction-instance} & Reduction instance\newline{\footnotesize\color{metaGrey}\claimid{reductions/def:reduction-instance}} & \pageref{def:reduction-instance}\\
\addlinespace[2pt]
\mref{Corollary}{cor:identity-instance} & The identity instance\newline{\footnotesize\color{metaGrey}\claimid{reductions/cor:identity-instance}} & \pageref{cor:identity-instance}\\
\addlinespace[2pt]
\mref{Corollary}{cor:flag-instance} & The flag instance of an uncharacterised branch\newline{\footnotesize\color{metaGrey}\claimid{reductions/cor:flag-instance}} & \pageref{cor:flag-instance}\\
\addlinespace[2pt]
\mref{Theorem}{thm:instrument-instance} & The branching instance: an instrument followed by branch instances\newline{\footnotesize\color{metaGrey}\claimid{reductions/thm:instrument-instance}} & \pageref{thm:instrument-instance}\\
\addlinespace[2pt]
\mref{Corollary}{cor:register-instance} & The register instance: register-conditioned squashing of instances\newline{\footnotesize\color{metaGrey}\claimid{reductions/cor:register-instance}} & \pageref{cor:register-instance}\\
\addlinespace[2pt]
\mref{Lemma}{lem:precomposed-instance} & Precomposition of a reduction instance with a channel\newline{\footnotesize\color{metaGrey}\claimid{reductions/lem:precomposed-instance}} & \pageref{lem:precomposed-instance}\\
\addlinespace[2pt]
\mref{Definition}{def:emitnorm} & The emission normalizer, the declared datum and the declared model\newline{\footnotesize\color{metaGrey}\claimid{reductions/def:emitnorm}} & \pageref{def:emitnorm}\\
\addlinespace[2pt]
\mref{Theorem}{thm:emit} & Exact emission of a reduction instance\newline{\footnotesize\color{metaGrey}\claimid{reductions/thm:emit}} & \pageref{thm:emit}\\
\addlinespace[2pt]
\mref{Corollary}{cor:entries} & The emitted constraint entries\newline{\footnotesize\color{metaGrey}\claimid{reductions/cor:entries}} & \pageref{cor:entries}\\
\addlinespace[2pt]
\mref{Remark}{rem:entry-tags} & Party-local tags and the one-factor emission\newline{\footnotesize\color{metaGrey}\claimid{reductions/rem:entry-tags}} & \pageref{rem:entry-tags}\\
\addlinespace[2pt]
\mref{Theorem}{thm:emitted-product} & The emission of several measuring parties\newline{\footnotesize\color{metaGrey}\claimid{reductions/thm:emitted-product}} & \pageref{thm:emitted-product}\\
\addlinespace[2pt]
\mref{Corollary}{cor:emitted-identity} & The emitted family of the identity instance\newline{\footnotesize\color{metaGrey}\claimid{reductions/cor:emitted-identity}} & \pageref{cor:emitted-identity}\\
\addlinespace[2pt]
\mref{Remark}{rem:emitclaim} & What the reduction's claim is about after an emission snap\newline{\footnotesize\color{metaGrey}\claimid{reductions/rem:emitclaim}} & \pageref{rem:emitclaim}\\
\addlinespace[2pt]
\mref{Remark}{rem:claimscope} & The claims that read a declared number\newline{\footnotesize\color{metaGrey}\claimid{reductions/rem:claimscope}} & \pageref{rem:claimscope}\\
\addlinespace[2pt]
\mref{Remark}{rem:exactboundary} & The boundary of the exact decisions\newline{\footnotesize\color{metaGrey}\claimid{reductions/rem:exactboundary}} & \pageref{rem:exactboundary}\\
\addlinespace[2pt]
\mref{Remark}{rem:sound-and-atk} & Soundness of the emitted constraint list\newline{\footnotesize\color{metaGrey}\claimid{reductions/rem:sound-and-atk}} & \pageref{rem:sound-and-atk}\\
\addlinespace[2pt]
\mref{Corollary}{cor:reference-companion} & The reference companion of a reduction instance\newline{\footnotesize\color{metaGrey}\claimid{reductions/cor:reference-companion}} & \pageref{cor:reference-companion}\\
\addlinespace[2pt]
\mref{Corollary}{cor:companion-factor-enclosure} & Factor enclosures carried to the reference companion\newline{\footnotesize\color{metaGrey}\claimid{reductions/cor:companion-factor-enclosure}} & \pageref{cor:companion-factor-enclosure}\\
\addlinespace[2pt]
\mref{Lemma}{lem:row-relaxation} & A rounded constraint entry with its constant relaxed\newline{\footnotesize\color{metaGrey}\claimid{reductions/lem:row-relaxation}} & \pageref{lem:row-relaxation}\\
\addlinespace[2pt]
\mref{Assumption}{ass:meas} & Each measuring party's device as one POVM per round\newline{\footnotesize\color{metaGrey}\claimid{reductions/ass:meas}} & \pageref{ass:meas}\\
\addlinespace[2pt]
\mref{Assumption}{ass:round} & Round homogeneity of the physical measurement\newline{\footnotesize\color{metaGrey}\claimid{reductions/ass:round}} & \pageref{ass:round}\\
\addlinespace[2pt]
\mref{Assumption}{ass:recv} & Supplier artefact whose claims hold of the device\newline{\footnotesize\color{metaGrey}\claimid{reductions/ass:recv}} & \pageref{ass:recv}\\
\addlinespace[2pt]
\mref{Remark}{rem:recv-scope} & What the supplier assumption is stated of\newline{\footnotesize\color{metaGrey}\claimid{reductions/rem:recv-scope}} & \pageref{rem:recv-scope}\\
\addlinespace[2pt]
\mref{Assumption}{ass:probe} & The station's private setting drawn from its own declared randomness\newline{\footnotesize\color{metaGrey}\claimid{reductions/ass:probe}} & \pageref{ass:probe}\\
\addlinespace[2pt]
\mref{Assumption}{ass:tsm} & Output-state test statistics\newline{\footnotesize\color{metaGrey}\claimid{reductions/ass:tsm}} & \pageref{ass:tsm}\\
\addlinespace[2pt]
\mref{Assumption}{ass:constdecl} & Declared constants of the constraint list\newline{\footnotesize\color{metaGrey}\claimid{reductions/ass:constdecl}} & \pageref{ass:constdecl}\\
\addlinespace[2pt]
\mref{Requirement}{req:channel-witness} & A channel witness for every emitted link\newline{\footnotesize\color{metaGrey}\claimid{reductions/req:channel-witness}} & \pageref{req:channel-witness}\\
\addlinespace[2pt]
\mref{Requirement}{req:pinched-channel-witness} & The pinched squashing channel as the emitted witness\newline{\footnotesize\color{metaGrey}\claimid{reductions/req:pinched-channel-witness}} & \pageref{req:pinched-channel-witness}\\
\addlinespace[2pt]
\mref{Requirement}{req:carried-tuple-emitted} & Emission of the carried tuple, and constraints built from it\newline{\footnotesize\color{metaGrey}\claimid{reductions/req:carried-tuple-emitted}} & \pageref{req:carried-tuple-emitted}\\
\addlinespace[2pt]
\mref{Requirement}{req:no-approximate-relations} & Refusal of approximate squashing relations\newline{\footnotesize\color{metaGrey}\claimid{reductions/req:no-approximate-relations}} & \pageref{req:no-approximate-relations}\\
\addlinespace[2pt]
\mref{Requirement}{req:round-homogeneous-emission} & Emission of the round-homogeneous case only\newline{\footnotesize\color{metaGrey}\claimid{reductions/req:round-homogeneous-emission}} & \pageref{req:round-homogeneous-emission}\\
\addlinespace[2pt]
\mref{Requirement}{req:post-processing-order} & Noise and coarse-graining before the first squashing link\newline{\footnotesize\color{metaGrey}\claimid{reductions/req:post-processing-order}} & \pageref{req:post-processing-order}\\
\addlinespace[2pt]
\mref{Requirement}{req:declared-relation-check} & Exact check of a declared relation\newline{\footnotesize\color{metaGrey}\claimid{reductions/req:declared-relation-check}} & \pageref{req:declared-relation-check}\\
\addlinespace[2pt]
\mref{Requirement}{req:reserved-link-kinds} & Refusal of the reserved receiver-stage link kind\newline{\footnotesize\color{metaGrey}\claimid{reductions/req:reserved-link-kinds}} & \pageref{req:reserved-link-kinds}\\
\addlinespace[2pt]
\mref{Requirement}{req:unpadded-branch-blocks} & Branch blocks emitted as supplied\newline{\footnotesize\color{metaGrey}\claimid{reductions/req:unpadded-branch-blocks}} & \pageref{req:unpadded-branch-blocks}\\
\addlinespace[2pt]
\mref{Requirement}{req:constraint-and-licence-slots} & Constraint and licence slots on every artefact\newline{\footnotesize\color{metaGrey}\claimid{reductions/req:constraint-and-licence-slots}} & \pageref{req:constraint-and-licence-slots}\\
\addlinespace[2pt]
\mref{Requirement}{req:licence-record} & Content and presence of the licence record\newline{\footnotesize\color{metaGrey}\claimid{reductions/req:licence-record}} & \pageref{req:licence-record}\\
\addlinespace[2pt]
\mref{Requirement}{req:declared-constants-only} & Refusal of data-dependent constants\newline{\footnotesize\color{metaGrey}\claimid{reductions/req:declared-constants-only}} & \pageref{req:declared-constants-only}\\
\addlinespace[2pt]
\mref{Requirement}{req:mixed-status-chain} & Refusal of a chain with a link that is not exact\newline{\footnotesize\color{metaGrey}\claimid{reductions/req:mixed-status-chain}} & \pageref{req:mixed-status-chain}\\
\addlinespace[2pt]
\mref{Requirement}{req:identity-link-recorded} & Identity link recorded for an already-finite receiver\newline{\footnotesize\color{metaGrey}\claimid{reductions/req:identity-link-recorded}} & \pageref{req:identity-link-recorded}\\
\addlinespace[2pt]
\mref{Requirement}{req:certified-entries-only} & Emission of certified constraint entries only\newline{\footnotesize\color{metaGrey}\claimid{reductions/req:certified-entries-only}} & \pageref{req:certified-entries-only}\\
\addlinespace[2pt]
\mref{Requirement}{req:class-marginal-operator-form} & Refusal of the operator form of the class-marginal identity\newline{\footnotesize\color{metaGrey}\claimid{reductions/req:class-marginal-operator-form}} & \pageref{req:class-marginal-operator-form}\\
\addlinespace[2pt]
\mref{Requirement}{req:constraint-list-by-request} & The constraint list and the licence decided by the request\newline{\footnotesize\color{metaGrey}\claimid{reductions/req:constraint-list-by-request}} & \pageref{req:constraint-list-by-request}\\
\addlinespace[2pt]
\mref{Requirement}{req:transported-set-certificates} & Certificates carried with the transported set\newline{\footnotesize\color{metaGrey}\claimid{reductions/req:transported-set-certificates}} & \pageref{req:transported-set-certificates}\\
\addlinespace[2pt]
\mref{Requirement}{req:class-marginal-entry} & The class-marginal entry only where the instance discharges it\newline{\footnotesize\color{metaGrey}\claimid{reductions/req:class-marginal-entry}} & \pageref{req:class-marginal-entry}\\
\addlinespace[2pt]
\mref{Requirement}{req:caller-constant-rounding} & Outward rounding of a caller's constant off the lattice\newline{\footnotesize\color{metaGrey}\claimid{reductions/req:caller-constant-rounding}} & \pageref{req:caller-constant-rounding}\\
\addlinespace[2pt]
\mref{Requirement}{req:uncertified-operator-refused} & Refusal of a requested operator without a decided certificate\newline{\footnotesize\color{metaGrey}\claimid{reductions/req:uncertified-operator-refused}} & \pageref{req:uncertified-operator-refused}\\
\addlinespace[2pt]
\mref{Requirement}{req:party-local-tag} & The party-local tag and its scope\newline{\footnotesize\color{metaGrey}\claimid{reductions/req:party-local-tag}} & \pageref{req:party-local-tag}\\
\addlinespace[2pt]
\mref{Requirement}{req:licensed-entry-lattice} & Lattice condition on the entries of emitted constraint operators\newline{\footnotesize\color{metaGrey}\claimid{reductions/req:licensed-entry-lattice}} & \pageref{req:licensed-entry-lattice}\\
\addlinespace[2pt]
\mref{Requirement}{req:exact-admission-refusal} & Refusal of a declaration outside the substrate\newline{\footnotesize\color{metaGrey}\claimid{reductions/req:exact-admission-refusal}} & \pageref{req:exact-admission-refusal}\\
\addlinespace[2pt]
\mref{Requirement}{req:emission-refusal} & Refusal of an emitted family that is not exactly admissible\newline{\footnotesize\color{metaGrey}\claimid{reductions/req:emission-refusal}} & \pageref{req:emission-refusal}\\
\addlinespace[2pt]
\mref{Requirement}{req:emission-record} & The record of an emission\newline{\footnotesize\color{metaGrey}\claimid{reductions/req:emission-record}} & \pageref{req:emission-record}\\
\addlinespace[2pt]
\mref{Requirement}{req:reference-companion} & Emitting the reference companion beside the artefact\newline{\footnotesize\color{metaGrey}\claimid{reductions/req:reference-companion}} & \pageref{req:reference-companion}\\
\addlinespace[2pt]
\mref{Requirement}{req:party-label} & Party partition emitted with the sector map\newline{\footnotesize\color{metaGrey}\claimid{reductions/req:party-label}} & \pageref{req:party-label}\\
\addlinespace[2pt]
\mref{Requirement}{req:factorised-emission} & Emission of several measuring parties as factors\newline{\footnotesize\color{metaGrey}\claimid{reductions/req:factorised-emission}} & \pageref{req:factorised-emission}\\
\addlinespace[2pt]
\mref{Requirement}{req:assumption-ledger} & Rows of the emitted assumption ledger\newline{\footnotesize\color{metaGrey}\claimid{reductions/req:assumption-ledger}} & \pageref{req:assumption-ledger}\\
\addlinespace[2pt]
\mref{Requirement}{req:never-verified-assumptions} & Never-verified rows and validation of declared data\newline{\footnotesize\color{metaGrey}\claimid{reductions/req:never-verified-assumptions}} & \pageref{req:never-verified-assumptions}\\
\addlinespace[2pt]
\mref{Requirement}{req:supplier-ledger-pass-through} & Pass-through of the supplier's assumption ledger\newline{\footnotesize\color{metaGrey}\claimid{reductions/req:supplier-ledger-pass-through}} & \pageref{req:supplier-ledger-pass-through}\\
\addlinespace[2pt]
\mref{Requirement}{req:supplier-ledger-refusals} & Refusal of a supplier ledger with a repeated or missing key\newline{\footnotesize\color{metaGrey}\claimid{reductions/req:supplier-ledger-refusals}} & \pageref{req:supplier-ledger-refusals}\\
\addlinespace[2pt]
\mref{Requirement}{req:supplier-version-accepted} & Refusal of a supplier artefact of an unaccepted version\newline{\footnotesize\color{metaGrey}\claimid{reductions/req:supplier-version-accepted}} & \pageref{req:supplier-version-accepted}\\
\addlinespace[2pt]
\mref{Requirement}{req:supplier-presence-rules} & Supplier presence rules that omit no imported premise\newline{\footnotesize\color{metaGrey}\claimid{reductions/req:supplier-presence-rules}} & \pageref{req:supplier-presence-rules}\\
\addlinespace[2pt]
\mref{Requirement}{req:recv-row-provenance} & Omission of the supplier row for a receiver artefact the run built\newline{\footnotesize\color{metaGrey}\claimid{reductions/req:recv-row-provenance}} & \pageref{req:recv-row-provenance}\\
\addlinespace[2pt]
\mref{Requirement}{req:output-state-statistics} & Output-state model of the test statistics\newline{\footnotesize\color{metaGrey}\claimid{reductions/req:output-state-statistics}} & \pageref{req:output-state-statistics}\\
\addlinespace[2pt]
\mref{Definition}{def:targets} & Certified weights and the flag-state and weight-preserving targets\newline{\footnotesize\color{metaGrey}\claimid{reductions/wpfss/def:targets}} & \pageref{def:targets}\\
\addlinespace[2pt]
\mref{Lemma}{lem:ple1} & Bound on the total certified weight\newline{\footnotesize\color{metaGrey}\claimid{reductions/wpfss/lem:ple1}} & \pageref{lem:ple1}\\
\addlinespace[2pt]
\mref{Remark}{rem:peqone} & The boundary case of total weight one\newline{\footnotesize\color{metaGrey}\claimid{reductions/wpfss/rem:peqone}} & \pageref{rem:peqone}\\
\addlinespace[2pt]
\mref{Proposition}{prop:chain} & Flag-state squashers: existence, chain and encoded weights\newline{\footnotesize\color{metaGrey}\claimid{reductions/wpfss/prop:chain}} & \pageref{prop:chain}\\
\addlinespace[2pt]
\mref{Definition}{def:routes} & The encoded and the licensed route\newline{\footnotesize\color{metaGrey}\claimid{reductions/wpfss/def:routes}} & \pageref{def:routes}\\
\addlinespace[2pt]
\mref{Remark}{rem:tworoutes} & Soundness of the two routes and one route per artefact\newline{\footnotesize\color{metaGrey}\claimid{reductions/wpfss/rem:tworoutes}} & \pageref{rem:tworoutes}\\
\addlinespace[2pt]
\mref{Proposition}{prop:weight-monotone} & Monotonicity of the weight-preserving target in its weights\newline{\footnotesize\color{metaGrey}\claimid{reductions/wpfss/prop:weight-monotone}} & \pageref{prop:weight-monotone}\\
\addlinespace[2pt]
\mref{Remark}{rem:weightfree} & Why the largest certified weight is used\newline{\footnotesize\color{metaGrey}\claimid{reductions/wpfss/rem:weightfree}} & \pageref{rem:weightfree}\\
\addlinespace[2pt]
\mref{Remark}{rem:published} & Relation to the published flag-state squashers\newline{\footnotesize\color{metaGrey}\claimid{reductions/wpfss/rem:published}} & \pageref{rem:published}\\
\addlinespace[2pt]
\mref{Lemma}{lem:descent} & Descent of a dark-count matrix through a coarse-graining\newline{\footnotesize\color{metaGrey}\claimid{reductions/wpfss/lem:descent}} & \pageref{lem:descent}\\
\addlinespace[2pt]
\mref{Remark}{rem:descentneedsconstant} & Necessity of the constancy clause\newline{\footnotesize\color{metaGrey}\claimid{reductions/wpfss/rem:descentneedsconstant}} & \pageref{rem:descentneedsconstant}\\
\addlinespace[2pt]
\mref{Definition}{def:weight-route} & The weight route\newline{\footnotesize\color{metaGrey}\claimid{reductions/wpfss/def:weight-route}} & \pageref{def:weight-route}\\
\addlinespace[2pt]
\mref{Proposition}{prop:route} & Certified weights routed from the receiver artifact\newline{\footnotesize\color{metaGrey}\claimid{reductions/wpfss/prop:route}} & \pageref{prop:route}\\
\addlinespace[2pt]
\mref{Proposition}{prop:precomposed-blocks} & Preserved blocks supplied with the dark counts already composed\newline{\footnotesize\color{metaGrey}\claimid{reductions/wpfss/prop:precomposed-blocks}} & \pageref{prop:precomposed-blocks}\\
\addlinespace[2pt]
\mref{Corollary}{cor:oneabsorbing} & Dark-count sharpening at a single bounded letter only\newline{\footnotesize\color{metaGrey}\claimid{reductions/wpfss/cor:oneabsorbing}} & \pageref{cor:oneabsorbing}\\
\addlinespace[2pt]
\mref{Remark}{rem:absorbneeded} & Necessity of the absorbing hypothesis\newline{\footnotesize\color{metaGrey}\claimid{reductions/wpfss/rem:absorbneeded}} & \pageref{rem:absorbneeded}\\
\addlinespace[2pt]
\mref{Remark}{rem:fibre} & Necessity of the fibre condition\newline{\footnotesize\color{metaGrey}\claimid{reductions/wpfss/rem:fibre}} & \pageref{rem:fibre}\\
\addlinespace[2pt]
\mref{Remark}{rem:colstoch} & Column-stochasticity of R and of its descended matrix\newline{\footnotesize\color{metaGrey}\claimid{reductions/wpfss/rem:colstoch}} & \pageref{rem:colstoch}\\
\addlinespace[2pt]
\mref{Remark}{rem:minweight} & Per-setting weights against the minimum over the settings\newline{\footnotesize\color{metaGrey}\claimid{reductions/wpfss/rem:minweight}} & \pageref{rem:minweight}\\
\addlinespace[2pt]
\mref{Remark}{rem:darkcompose} & Composition with the supplier's dark-count bounds\newline{\footnotesize\color{metaGrey}\claimid{reductions/wpfss/rem:darkcompose}} & \pageref{rem:darkcompose}\\
\addlinespace[2pt]
\mref{Remark}{rem:receiver-imports} & Supplier assumptions behind the receiver imports\newline{\footnotesize\color{metaGrey}\claimid{reductions/wpfss/rem:receiver-imports}} & \pageref{rem:receiver-imports}\\
\addlinespace[2pt]
\mref{Corollary}{cor:fssinstance} & Licensed flag-sector entry of a certified weight\newline{\footnotesize\color{metaGrey}\claimid{reductions/wpfss/cor:fssinstance}} & \pageref{cor:fssinstance}\\
\addlinespace[2pt]
\mref{Proposition}{prop:flat-instance} & The flat instance of the post-processed receiver\newline{\footnotesize\color{metaGrey}\claimid{reductions/wpfss/prop:flat-instance}} & \pageref{prop:flat-instance}\\
\addlinespace[2pt]
\mref{Corollary}{cor:emitted-flat} & Exact emission of the flat instance\newline{\footnotesize\color{metaGrey}\claimid{reductions/wpfss/cor:emitted-flat}} & \pageref{cor:emitted-flat}\\
\addlinespace[2pt]
\mref{Corollary}{cor:two-sector-decided} & The requested two-sector entry, decided on the emitted numbers\newline{\footnotesize\color{metaGrey}\claimid{reductions/wpfss/cor:two-sector-decided}} & \pageref{cor:two-sector-decided}\\
\addlinespace[2pt]
\mref{Corollary}{cor:flat-companion} & The flat instance has a reference companion\newline{\footnotesize\color{metaGrey}\claimid{reductions/wpfss/cor:flat-companion}} & \pageref{cor:flat-companion}\\
\addlinespace[2pt]
\mref{Corollary}{cor:normalized-flat} & The flat instance of a normalized station\newline{\footnotesize\color{metaGrey}\claimid{reductions/wpfss/cor:normalized-flat}} & \pageref{cor:normalized-flat}\\
\addlinespace[2pt]
\mref{Requirement}{req:one-route-per-artifact} & One route per artefact\newline{\footnotesize\color{metaGrey}\claimid{reductions/wpfss/req:one-route-per-artifact}} & \pageref{req:one-route-per-artifact}\\
\addlinespace[2pt]
\mref{Requirement}{req:total-weight-refusal} & Refusal of a total weight of one or more\newline{\footnotesize\color{metaGrey}\claimid{reductions/wpfss/req:total-weight-refusal}} & \pageref{req:total-weight-refusal}\\
\addlinespace[2pt]
\mref{Requirement}{req:largest-certified-weight} & Use and record of the largest certified weight\newline{\footnotesize\color{metaGrey}\claimid{reductions/wpfss/req:largest-certified-weight}} & \pageref{req:largest-certified-weight}\\
\addlinespace[2pt]
\mref{Requirement}{req:descent-decided-exactly} & Exact decision of monotonicity and descent\newline{\footnotesize\color{metaGrey}\claimid{reductions/wpfss/req:descent-decided-exactly}} & \pageref{req:descent-decided-exactly}\\
\addlinespace[2pt]
\mref{Requirement}{req:absorbing-bounded-letters} & Absorbing bounded letters as an input validity condition\newline{\footnotesize\color{metaGrey}\claimid{reductions/wpfss/req:absorbing-bounded-letters}} & \pageref{req:absorbing-bounded-letters}\\
\addlinespace[2pt]
\mref{Requirement}{req:fibre-equality} & Certified sets equal to fibres of the coarse-graining\newline{\footnotesize\color{metaGrey}\claimid{reductions/wpfss/req:fibre-equality}} & \pageref{req:fibre-equality}\\
\addlinespace[2pt]
\mref{Requirement}{req:column-stochastic-refusal} & Exact column-stochasticity of the declared dark-count matrix\newline{\footnotesize\color{metaGrey}\claimid{reductions/wpfss/req:column-stochastic-refusal}} & \pageref{req:column-stochastic-refusal}\\
\addlinespace[2pt]
\mref{Requirement}{req:supplier-version} & Receiver versions from which per-setting weights are read\newline{\footnotesize\color{metaGrey}\claimid{reductions/wpfss/req:supplier-version}} & \pageref{req:supplier-version}\\
\addlinespace[2pt]
\mref{Requirement}{req:supplier-eig-record} & Invocation record of the eigenvalue-slack assumption in the ledger\newline{\footnotesize\color{metaGrey}\claimid{reductions/wpfss/req:supplier-eig-record}} & \pageref{req:supplier-eig-record}\\
\addlinespace[2pt]
\mref{Requirement}{req:normalized-link-record} & The link of a normalized station records the factored loss\newline{\footnotesize\color{metaGrey}\claimid{reductions/wpfss/req:normalized-link-record}} & \pageref{req:normalized-link-record}\\
\addlinespace[2pt]
\mref{Requirement}{req:normalized-branch-refusal} & A normalized receiver is not read as a branch of a branching instance\newline{\footnotesize\color{metaGrey}\claimid{reductions/wpfss/req:normalized-branch-refusal}} & \pageref{req:normalized-branch-refusal}\\
\addlinespace[2pt]
\mref{Definition}{def:domination-link} & The domination link\newline{\footnotesize\color{metaGrey}\claimid{reductions/domination/def:domination-link}} & \pageref{def:domination-link}\\
\addlinespace[2pt]
\mref{Theorem}{thm:noise-channel} & The noise channel of a domination\newline{\footnotesize\color{metaGrey}\claimid{reductions/domination/thm:noise-channel}} & \pageref{thm:noise-channel}\\
\addlinespace[2pt]
\mref{Remark}{rem:rotated-projective} & Rotated projective elements and the noise floor\newline{\footnotesize\color{metaGrey}\claimid{reductions/domination/rem:rotated-projective}} & \pageref{rem:rotated-projective}\\
\addlinespace[2pt]
\mref{Corollary}{cor:dom-chain} & A squashing link followed by a domination link\newline{\footnotesize\color{metaGrey}\claimid{reductions/domination/cor:dom-chain}} & \pageref{cor:dom-chain}\\
\bottomrule
\end{xltabular}}

\clearpage
\bookmarksetupnext{level=part}
\printindex

\clearpage\phantomsection\bookmarksetupnext{level=part}%
\addcontentsline{toc}{section}{References}
\bibliographystyle{plain}
\bibliography{../../../refs/references}

\end{document}
